\documentclass[10pt,oneside]{book}
\pdftrailerid{}

\usepackage[T1]{fontenc}
\usepackage[utf8]{inputenc}
\usepackage[spanish]{babel}
\usepackage{amsmath,amssymb,mathrsfs}
\usepackage[all]{xy}
\usepackage[a4paper,margin=22mm]{geometry}
\setlength{\emergencystretch}{2em}
\usepackage[
  colorlinks=true,
  linkcolor=blue,
  hypertexnames=false,
  unicode=true,
  pdfencoding=auto
]{hyperref}
\hypersetup{
  pdftitle={Elementos de geometría algebraica IV-1 -- Capítulo 0 -- edición española actual},
  pdfauthor={Alexander Grothendieck y Jean Dieudonné},
  pdfsubject={Capítulo 0: preliminares de EGA IV-1, secciones 14 a 23}
}

\renewcommand{\thesection}{\arabic{section}}
\renewcommand{\thesubsection}{\thesection.\arabic{subsection}}
\newcommand{\gr}{\operatorname{gr}}
\newcommand{\sref}[1]{(\ref{#1})}
\makeatletter
\newcommand{\setmanualcurrentlabel}[1]{\def\@currentlabel{#1}}
\makeatother
\newcommand{\oldpage}[2][]{%
  \setcounter{footnote}{0}%
  \marginpar{\scriptsize\textbf{#1}~|~#2}\ignorespaces}
\newenvironment{env}[1][]{\setmanualcurrentlabel{#1}\par\medskip\noindent\textbf{(#1)}\ }{\par}
\newenvironment{definition}[1][]{%
  \setmanualcurrentlabel{#1}\par\medskip\noindent\itshape Definición (#1). ---\ }{\par}
\newenvironment{proposition}[1][]{%
  \setmanualcurrentlabel{#1}\par\medskip\noindent\itshape Proposición (#1). ---\ }{\par}
\newenvironment{theorem}[1][]{%
  \setmanualcurrentlabel{#1}\par\medskip\noindent\itshape Teorema (#1). ---\ }{\par}
\newenvironment{corollary}[1][]{%
  \setmanualcurrentlabel{#1}\par\medskip\noindent\itshape Corolario (#1). ---\ }{\par}
\newenvironment{remark}[1][]{%
  \setmanualcurrentlabel{#1}\par\medskip\noindent\itshape Observación (#1). ---\ }{\par}
\newenvironment{remarks}[1][]{%
  \setmanualcurrentlabel{#1}\par\medskip\noindent\itshape Observaciones (#1). ---\ }{\par}
\newenvironment{lemma}[1][]{%
  \setmanualcurrentlabel{#1}\par\medskip\noindent\itshape Lema (#1). ---\ }{\par}
\newenvironment{example}[1][]{%
  \setmanualcurrentlabel{#1}\par\medskip\noindent\textit{Ejemplo (#1). ---}\ }{\par}
\newenvironment{examples}[1][]{%
  \setmanualcurrentlabel{#1}\par\medskip\noindent\textit{Ejemplos (#1). ---}\ }{\par}
\newenvironment{proof}{%
  \par\medskip\noindent\textit{Demostración. ---\ }}{\par}

\setlength{\emergencystretch}{6em}
\begin{document}
% BEGIN RIGHTS AND AI TRANSLATION NOTICE
\begingroup
\small

% BEGIN INHERITED HUMAN SOURCE CREDIT
\noindent\textbf{Fuentes y contribuciones}

La obra original es de Alexandre Grothendieck, con la colaboración de Jean Dieudonné: Éléments de géométrie algébrique IV, première partie, capítulo 0, §§ 14–23, Publications Mathématiques de l’IHÉS 20 (1964).

Las fuentes inglesas utilizadas en la preparación de esta traducción española incorporan la traducción humana de \href{https://github.com/ryankeleti/ega}{ryankeleti/ega}, versión \href{https://github.com/ryankeleti/ega/tree/3d7afe28c99d1887ccdb85acf2ca1afd93fb1971}{3d7afe28}, en el § 14 del capítulo 0. El proyecto nombra a Mateo Carmona, Tim Holzschuh, Tim Hosgood, Ryan Keleti, Thiago Solovera \& Nery y Wayne Yang como colaboradores; sus \href{https://github.com/ryankeleti/ega/blob/3d7afe28c99d1887ccdb85acf2ca1afd93fb1971/README.md}{créditos por sección} y su \href{https://github.com/ryankeleti/ega/commits/3d7afe28c99d1887ccdb85acf2ca1afd93fb1971}{historial de contribuciones} precisan la autoría.

La traducción inglesa y la estructura TeX heredadas conservan esa procedencia. La traducción al español, las adiciones posteriores y las correcciones y composición producidas por IA se distinguen de esas contribuciones humanas. No se atribuye a los colaboradores de la traducción inglesa la traducción española ni el material posterior que no figura en su versión. La atribución no implica autorización ni aval.

\par\medskip
% END INHERITED HUMAN SOURCE CREDIT
\noindent\textbf{Derechos, disponibilidad y autoría}

Los derechos sobre las obras originales y sobre las contribuciones protegidas a estas traducciones y transcripciones derivadas corresponden a sus respectivos titulares. En la medida en que el editor tenga derechos sobre la traducción, la transcripción, las correcciones y la composición producidas por IA, los dedica al dominio público bajo CC0 1.0 Universal (https://creativecommons.org/publicdomain/zero/1.0/deed.es); esto no afecta a los derechos sobre las obras originales ni concede permiso para usarlas. Las menciones de derechos de terceros y las autorizaciones ya concedidas siguen siendo aplicables.

Estas ediciones se ofrecen gratuitamente, con fines no comerciales y exclusivamente educativos, bajo la presunción de que los respectivos titulares de derechos no se opondrían a que estas traducciones y transcripciones estuvieran disponibles gratuitamente junto a los originales franceses en NUMDAM (www.numdam.org/). Esta presunción no significa que se haya obtenido autorización; el acceso gratuito a los originales no autoriza por sí solo su traducción, adaptación o republicación.

Si usted es titular de derechos sobre algún material presentado aquí y se opone a su disponibilidad, puede abrir una incidencia en GitHub (github.com/KokunoYumeto/ega-fga-es/issues). El responsable de la publicación retirará el material afectado a petición del titular de los derechos.

La traducción al español está producida por IA. Las correcciones, la edición y la composición también están producidas por IA. El responsable de la publicación no tradujo el texto. Se conservan la autoría y las atribuciones de las obras originales y de las contribuciones heredadas; no se atribuyen a la IA. Los modelos documentados se identifican en cada edición.
\par
\endgroup
\clearpage
% END RIGHTS AND AI TRANSLATION NOTICE

\clearpage
\thispagestyle{plain}
\markboth{PRELIMINARES}{PRELIMINARES}
\oldpage[IV-1]{5}
\phantomsection
\label{chapter:0IV1-es}

\begin{center}
{\large CAPÍTULO 0 (continuación)\footnotemark\par}
\vspace{1.5em}
{\Large\bfseries PRELIMINARES\par}
\vspace{2em}
{\large\bfseries Sumario\par}
\end{center}

\footnotetext{Recordemos que el símbolo $0_N$ (resp. $\operatorname{Err}_N$)
remite a la parte del capítulo 0 (resp. a la lista de Erratas) adjunta al
capítulo $N$.}

\begin{center}
\begin{tabular}{@{}r l@{}}
\S~14. & Dimensión combinatoria de un espacio topológico.\\
\S~15. & Sucesiones M-regulares y sucesiones $\mathscr F$-regulares.\\
\S~16. & Dimensión y profundidad en los anillos locales noetherianos.\\
\S~17. & Anillos regulares.\\
\S~18. & Complementos sobre las extensiones de álgebras.\\
\S~19. & Álgebras formalmente lisas y anillos de Cohen.\\
\S~20. & Derivaciones y diferenciales.\\
\S~21. & Diferenciales en los anillos de característica $p$.\\
\S~22. & Criterios diferenciales de lisitud formal y de regularidad.\\
\S~23. & Anillos japoneses.
\end{tabular}
\end{center}

Casi la totalidad de los parágrafos precedentes está dedicada a la exposición
de nociones de Álgebra conmutativa que se utilizarán a lo largo del capítulo
IV. Aunque buena parte de estas nociones figura ya en varias obras ([1], [12],
[13], [17], [30]), se ha considerado que sería más cómodo para el lector
disponer de una exposición continua y aproximadamente independiente. Junto con
los \S\S~5, 6 y 7 del capítulo IV (donde se utiliza el lenguaje de los
esquemas), estos parágrafos constituyen, dentro de nuestro Tratado, un pequeño
Tratado especial, aproximadamente independiente de los capítulos I a III, que
pretende presentar de forma coherente las propiedades de los anillos que «~se
comportan bien~» respecto de operaciones como la completación o la clausura
integral, vinculando sistemáticamente estas propiedades con algunas
concepciones generales\footnote{La mayoría de las propiedades de que se trata
fueron descubiertas por Chevalley, Zariski, Nagata y Serre. El método utilizado
aquí fue desarrollado por primera vez en el otoño de 1961, en un curso impartido
en Harvard University por A.~Grothendieck.}.

\setcounter{section}{13}
\section{Dimensión combinatoria de un espacio topológico}
\label{section:0.14-es}

\setcounter{subsection}{0}
\subsection{Dimensión combinatoria de un espacio topológico}
\label{subsection:0.14.1-es}

\oldpage[0\textsubscript{IV-1}]{6}
\begin{env}[14.1.1]
\label{0.14.1.1-es}
Sea $I$ un conjunto ordenado; una \emph{cadena} de elementos de $I$ es, por
definición, una sucesión finita estrictamente creciente
$i_0<i_1<\cdots<i_n$ de elementos de $I$ ($n\geq 0$); por definición, la
\emph{longitud} de esta cadena es $n$. Si $X$ es un espacio topológico, el
conjunto de sus subconjuntos \emph{cerrados irreducibles} está ordenado por
inclusión, de donde se obtiene la noción de \emph{cadena} de subconjuntos
cerrados irreducibles de $X$.
\end{env}

\begin{definition}[14.1.2]
\label{0.14.1.2-es}
Sea $X$ un espacio topológico. Se llama \emph{dimensión combinatoria de $X$}
(o simplemente \emph{dimensión de $X$}, si no puede haber confusión), y se
denota por $\dim_c(X)$, o simplemente por $\dim(X)$, a la cota superior de las
longitudes de las cadenas de subconjuntos cerrados irreducibles de $X$. Para
todo $x\in X$, se llama \emph{dimensión combinatoria de $X$ en $x$} (o
simplemente \emph{dimensión de $X$ en $x$}), y se denota por $\dim_x(X)$, al
número $\inf_U(\dim(U))$, donde $U$ recorre el conjunto de los entornos
abiertos de $x$ en $X$.
\end{definition}

De esta definición resulta que
\[
  \dim(\varnothing)=-\infty
\]
(la cota superior en $\overline{\mathbb R}$ del conjunto vacío es $-\infty$).
Si $(X_\alpha)$ es la familia de las componentes irreducibles de $X$, se tiene
\[
  \dim(X)=\sup_\alpha\dim(X_\alpha)
  \tag{14.1.2.1}\label{0.14.1.2.1-es}
\]
porque toda cadena de subconjuntos cerrados irreducibles de $X$ está, por
definición, contenida en una componente irreducible de $X$ y, recíprocamente,
estas componentes son cerradas en $X$, por lo que todo subconjunto cerrado
irreducible de un $X_\alpha$ es un subconjunto cerrado irreducible de $X$.

\begin{definition}[14.1.3]
\label{0.14.1.3-es}
Se dice que un espacio topológico $X$ es \emph{equidimensional} si todas sus
componentes irreducibles tienen la misma dimensión (igual entonces a $\dim(X)$
por (14.1.2.1)).
\end{definition}

\begin{proposition}[14.1.4]
\label{0.14.1.4-es}
\begin{enumerate}
  \item[\rm(i)] Para todo subconjunto $Y$ de un espacio topológico $X$, se tiene
  $\dim(Y)\leq\dim(X)$.
  \item[\rm(ii)] Si un espacio topológico $X$ es unión finita de cerrados
  $X_i$, se tiene $\dim(X)=\sup_i\dim(X_i)$.
\end{enumerate}
\end{proposition}

\begin{proof}
Para todo subconjunto cerrado irreducible $Z$ de $Y$, la adherencia $\overline Z$
de $Z$ en $X$ es irreducible ($0_{\mathrm I}$, 2.1.2) y
$\overline Z\cap Y=Z$, de donde se sigue (i). Por otro lado, si
$X=\bigcup_{i=1}^n X_i$, donde los $X_i$ son cerrados, todo subconjunto cerrado
irreducible de $X$ está contenido en uno de los $X_i$ ($0_{\mathrm I}$, 2.1.1),
y, por consiguiente, toda cadena de subconjuntos cerrados irreducibles de $X$
está contenida en uno de los $X_i$~; de donde se sigue (ii).
\end{proof}

De (14.1.4, (i)) se deduce que, para todo $x\in X$, también se puede escribir
\[
  \dim_x(X)=\lim_U\dim(U)
  \tag{14.1.4.1}\label{0.14.1.4.1-es}
\]
donde el límite se toma con respecto al conjunto filtrante decreciente de los
entornos abiertos de $x$ en $X$.

\oldpage[0\textsubscript{IV-1}]{7}
\begin{corollary}[14.1.5]
\label{0.14.1.5-es}
Sean $X$ un espacio topológico, $x$ un punto de $X$, $U$ un entorno de
$x$, e $Y_i$ ($1\leq i\leq n$) subconjuntos cerrados de $U$ tales que
$x\in Y_i$ para todo $i$ y que $U$ sea la reunión de los $Y_i$. Entonces se tiene
\[
  \dim_x(X)=\sup_i\bigl(\dim_x(Y_i)\bigr).
  \tag{14.1.5.1}\label{0.14.1.5.1-es}
\]
\end{corollary}

\begin{proof}
De (14.1.4, (ii)) resulta que
$\dim_x(X)=\inf_V(\sup_i(\dim(Y_i\cap V)))$ cuando $V$ recorre el conjunto
de los entornos abiertos de $x$ contenidos en $U$~; del mismo modo, se tiene
$\dim_x(Y_i)=\inf_V(\dim(Y_i\cap V))$ para todo $i$. El corolario es, pues,
evidente si
\[
  \sup_i\bigl(\dim_x(Y_i)\bigr)=+\infty~;
\]
en caso contrario, existe un entorno abierto $V_0\subset U$ de $x$ tal
que $\dim(Y_i\cap V)=\dim_x(Y_i)$ para $1\leq i\leq n$ y para todo
$V\subset V_0$, de donde se sigue la conclusión.
\end{proof}

\begin{proposition}[14.1.6]
\label{0.14.1.6-es}
Para todo espacio topológico $X$, se tiene
$\dim(X)=\sup_{x\in X}\dim_x(X)$.
\end{proposition}

\begin{proof}
De la definición (14.1.2) y de la prop. (14.1.4) resulta que
$\dim_x(X)\leq\dim(X)$ para todo $x\in X$. Por otro lado, sea
$Z_0\subset Z_1\subset\cdots\subset Z_n$ una cadena de subconjuntos cerrados
irreducibles de $X$, y sea $x\in Z_0$~; para todo abierto $U\subset X$
que contenga a $x$, $U\cap Z_i$ es irreducible ($0_{\mathrm I}$, 2.1.6) y
cerrado en $U$, y como se tiene $\overline{U\cap Z_i}=Z_i$ en $X$, los
$U\cap Z_i$ son distintos dos a dos~; por tanto, $\dim(U)\geq n$, lo que
termina la demostración.
\end{proof}

\begin{corollary}[14.1.7]
\label{0.14.1.7-es}
Si $(X_\alpha)$ es un recubrimiento abierto de $X$, o un recubrimiento cerrado
localmente finito de $X$, se tiene
$\dim(X)=\sup_\alpha(\dim(X_\alpha))$.
\end{corollary}

\begin{proof}
Si $X_\alpha$ es un entorno de $x\in X$, se tiene
$\dim_x(X)\leq\dim(X_\alpha)$, de donde se sigue la primera afirmación. Por otro
lado, si los $X_\alpha$ son cerrados y si $U$ es un entorno de $x\in X$ que
solo interseca un número finito de conjuntos $X_\alpha$, se tiene
\[
  \dim_x(X)\leq\dim(U)=\sup_\alpha\bigl(\dim(U\cap X_\alpha)\bigr)
  \leq\sup_\alpha\bigl(\dim(X_\alpha)\bigr)
\]
por (14.1.4), de donde se sigue la segunda afirmación.
\end{proof}

\begin{corollary}[14.1.8]
\label{0.14.1.8-es}
Sea $X$ un espacio de Kolmogórov ($0_{\mathrm I}$, 2.1.3) noetheriano, y
sea $F$ el conjunto de los puntos cerrados de $X$. Entonces
$\dim(X)=\sup_{x\in F}\dim_x(X)$.
\end{corollary}

\begin{proof}
Con la notación de la demostración de (14.1.6), basta observar
que existe en $Z_0$ un punto cerrado ($0_{\mathrm I}$, 2.1.3).
\end{proof}

\begin{proposition}[14.1.9]
\label{0.14.1.9-es}
Sea $X$ un espacio de Kolmogórov noetheriano no vacío. Para que
$\dim(X)=0$, es necesario y suficiente que $X$ sea finito y discreto.
\end{proposition}

\begin{proof}
Si un espacio $X$ es separado (y \emph{a fortiori} si $X$ es un espacio
discreto), los únicos subconjuntos cerrados irreducibles de $X$ son los
reducidos a un punto, por lo que $\dim(X)=0$. Recíprocamente, supongamos que
$X$ es un espacio de Kolmogórov noetheriano tal que $\dim(X)=0$~; como toda
componente irreducible de $X$ contiene un punto cerrado ($0_{\mathrm I}$, 2.1.3),
ha de reducirse a dicho punto. Como $X$ solo tiene un número finito de
componentes irreducibles, es finito y discreto.
\end{proof}

\begin{corollary}[14.1.10]
\label{0.14.1.10-es}
Sea $X$ un espacio de Kolmogórov noetheriano. Para que un punto $x\in X$ sea
aislado, es necesario y suficiente que $\dim_x(X)=0$.
\end{corollary}

\begin{proof}
La condición es evidentemente necesaria (sin hipótesis sobre $X$). También es
suficiente,
\oldpage[0\textsubscript{IV-1}]{8}
pues de ella resulta que $\dim(U)=0$ para un entorno abierto $U$ de $x$, y
como $U$ es un espacio de Kolmogórov noetheriano, $U$ es finito y discreto.
\end{proof}

\begin{proposition}[14.1.11]
\label{0.14.1.11-es}
La función $x\mapsto\dim_x(X)$ es semicontinua superiormente en $X$.
\end{proposition}

\begin{proof}
Es claro que esta función es semicontinua superiormente en todo
punto en el que vale $+\infty$. Supongamos, pues, que $\dim_x(X)=n<+\infty$~;
entonces, la fórmula (14.1.4.1) muestra que existe un entorno abierto $U_0$ de
$x$ tal que $\dim(U)=n$ para todo entorno abierto $U\subset U_0$ de $x$.
Dicho esto, para todo $y\in U_0$ y todo entorno abierto $V\subset U_0$ de $y$,
se tiene $\dim(V)\leq\dim(U_0)=n$ (14.1.4)~; se deduce, por tanto, de
(14.1.4.1) que $\dim_y(X)\leq n$.
\end{proof}

\begin{remark}[14.1.12]
\label{0.14.1.12-es}
Si $X$ e $Y$ son dos espacios topológicos, y $f:X\to Y$ una aplicación
continua, puede ocurrir que $\dim(f(X))>\dim(X)$~; se obtiene un ejemplo tomando
como $X$ un espacio discreto de dos elementos $a$, $b$, y como $Y$ el conjunto
$\{a,b\}$ dotado de la topología cuyos conjuntos cerrados son
$\varnothing$, $\{a\}$ y $\{a,b\}$~; si $f:X\to Y$ es la aplicación
identidad, se tiene $\dim(Y)=1$ y $\dim(X)=0$. Obsérvese que $Y$ es el espectro
de un anillo de valoración discreta $A$; su único punto cerrado es $a$ y su
punto genérico es $b$~; si $K$ y $k$ son, respectivamente, el cuerpo de
fracciones y el cuerpo residual de $A$, $X$ es el espectro del anillo $k\times K$
y $f$ la aplicación continua correspondiente al homomorfismo
$(\varphi,\psi):A\to k\times K$, donde $\varphi:A\to k$ y $\psi:A\to K$ son
los homomorfismos canónicos (cf. IV, 5.4.3).
\end{remark}

\subsection{Codimensión de una parte cerrada}
\label{subsection:0.14.2-es}

\begin{definition}[14.2.1]
\label{0.14.2.1-es}
Dada una parte cerrada irreducible $Y$ de un espacio topológico $X$, se llama
\emph{codimensión combinatoria} (o simplemente \emph{codimensión}) de $Y$ en
$X$, y se denota por $\operatorname{codim}(Y,X)$, al supremo de las longitudes
de las cadenas de partes cerradas irreducibles de $X$ cuyo menor elemento es
$Y$. Si $Y$ es una parte cerrada cualquiera de $X$, se llama
\emph{codimensión de $Y$ en $X$}, y se denota asimismo por
$\operatorname{codim}(Y,X)$, al ínfimo de las codimensiones en $X$ de las
componentes irreducibles de $Y$. Se dice que $X$ es
\emph{equicodimensional} si todas las partes cerradas irreducibles mínimas de
$X$ tienen la misma codimensión en $X$.
\end{definition}

De esta definición resulta que
$\operatorname{codim}(\varnothing,X)=+\infty$, pues el ínfimo de la parte vacía
de $\overline{\mathbb R}$ es $+\infty$. Si $Y$ es cerrada en $X$ y si
$(X_\alpha)$ (resp. $(Y_\beta)$) es la familia de las componentes irreducibles
de $X$ (resp. de $Y$), todo $Y_\beta$ está contenido en algún $X_\alpha$ y,
más generalmente, toda cadena de partes cerradas irreducibles de $X$ cuyo menor
elemento es $Y_\beta$ está formada por partes de algún $X_\alpha$; por tanto,
se tiene
\[
  \operatorname{codim}(Y,X)=
  \inf_\beta\left(\sup_\alpha
  \bigl(\operatorname{codim}(Y_\beta,X_\alpha)\bigr)\right)
  \tag{14.2.1.1}\label{0.14.2.1.1-es}
\]
donde, para cada $\beta$, $\alpha$ recorre el conjunto de los índices tales que
$Y_\beta\subset X_\alpha$.

\begin{proposition}[14.2.2]
\label{0.14.2.2-es}
Sea $X$ un espacio topológico.
\begin{enumerate}
  \item[\rm(i)] Si $\Phi$ es el conjunto de las partes cerradas irreducibles de
  $X$, se tiene
  \[
    \dim(X)=\sup_{Y\in\Phi}\bigl(\operatorname{codim}(Y,X)\bigr).
    \tag{14.2.2.1}\label{0.14.2.2.1-es}
  \]
  \oldpage[0\textsubscript{IV-1}]{9}
  \item[\rm(ii)] Para toda parte cerrada no vacía $Y$ de $X$, se tiene
  \[
    \dim(Y)+\operatorname{codim}(Y,X)\leq\dim(X).
    \tag{14.2.2.2}\label{0.14.2.2.2-es}
  \]
  \item[\rm(iii)] Si $Y$, $Z$ y $T$ son tres partes cerradas de $X$ tales que
  $Y\subset Z\subset T$, se tiene
  \[
    \operatorname{codim}(Y,Z)+\operatorname{codim}(Z,T)
    \leq\operatorname{codim}(Y,T).
    \tag{14.2.2.3}\label{0.14.2.2.3-es}
  \]
  \item[\rm(iv)] Para que una parte cerrada $Y$ de $X$ satisfaga
  $\operatorname{codim}(Y,X)=0$, es necesario y suficiente que $Y$ contenga
  una componente irreducible de $X$.
\end{enumerate}
\end{proposition}

\begin{proof}
Las afirmaciones (i) y (iv) son consecuencias inmediatas de la definición
(14.2.1). Para demostrar (ii), puede uno reducirse al caso en que $Y$ es
irreducible, y entonces la fórmula resulta de las definiciones (14.1.1) y
(14.2.1). Por último, para demostrar (iii), puede uno, en virtud de la
definición (14.2.1), reducirse primero al caso en que $Y$ es irreducible;
entonces $\operatorname{codim}(Y,Z)=
\sup_\alpha(\operatorname{codim}(Y,Z_\alpha))$ para las componentes
irreducibles $Z_\alpha$ de $Z$ que contienen $Y$; está claro que
$\operatorname{codim}(Y,T)\geq\operatorname{codim}(Y,Z)$, por lo que la
desigualdad es cierta si $\operatorname{codim}(Y,Z)=+\infty$; en caso
contrario, existe un $\alpha$ tal que
$\operatorname{codim}(Y,Z)=\operatorname{codim}(Y,Z_\alpha)$ y, en virtud de
(14.2.1), puede uno reducirse al caso en que $Z$ es también irreducible; pero
entonces la desigualdad (14.2.2.3) es una consecuencia evidente de la
definición (14.2.1).
\end{proof}

\begin{proposition}[14.2.3]
\label{0.14.2.3-es}
Sean $X$ un espacio topológico y $Y$ una parte cerrada de $X$. Para todo
abierto $U$ de $X$, se tiene
\[
  \operatorname{codim}(Y\cap U,U)\geq\operatorname{codim}(Y,X).
  \tag{14.2.3.1}\label{0.14.2.3.1-es}
\]
Además, para que los dos miembros de (14.2.3.1) sean iguales, es necesario y
suficiente que, si $(Y_\alpha)$ es la familia de las componentes irreducibles
de $Y$ que tienen intersección no vacía con $U$, se tenga
$\operatorname{codim}(Y,X)=
\inf_\alpha(\operatorname{codim}(Y_\alpha,X))$.
\end{proposition}

\begin{proof}
Se sabe ($0_{\mathrm I}$, 2.1.6) que $Z\mapsto\overline Z$ es una biyección
del conjunto de las partes cerradas irreducibles de $U$ sobre el conjunto de
las partes cerradas irreducibles de $X$ que tienen intersección no vacía con
$U$ y, en particular, hace corresponder a las componentes irreducibles de
$Y\cap U$ las componentes irreducibles de $Y$ que tienen intersección no vacía
con $U$; si $Y_\alpha$ es una de estas últimas,
se tiene, por tanto,
$\operatorname{codim}(Y_\alpha,X)=
\operatorname{codim}(Y_\alpha\cap U,U)$, y la proposición resulta entonces de
la definición (14.2.1).
\end{proof}

\begin{definition}[14.2.4]
\label{0.14.2.4-es}
Sean $X$ un espacio topológico, $Y$ una parte cerrada de $X$ y $x$ un punto de
$X$. Se llama \emph{codimensión de $Y$ en $X$ en el punto $x$}, y se denota
por $\operatorname{codim}_x(Y,X)$, al número
$\sup_U(\operatorname{codim}(Y\cap U,U))$, donde $U$ recorre el conjunto de
los entornos abiertos de $x$ en $X$.
\end{definition}

En virtud de (14.2.3), también puede escribirse
\[
  \operatorname{codim}_x(Y,X)=
  \lim_U\bigl(\operatorname{codim}(Y\cap U,U)\bigr)
  \tag{14.2.4.1}\label{0.14.2.4.1-es}
\]
donde el límite se toma sobre el conjunto filtrante decreciente de los
entornos abiertos de $x$ en $X$. Obsérvese que se tiene
\[
  \operatorname{codim}_x(Y,X)=+\infty\quad\text{si}\quad x\in X-Y.
\]

\oldpage[0\textsubscript{IV-1}]{10}
\begin{proposition}[14.2.5]
\label{0.14.2.5-es}
Si $(Y_i)_{1\leq i\leq n}$ es una familia finita de partes cerradas de un
espacio topológico $X$, y $Y$ es la unión de esta familia, se tiene
\[
  \operatorname{codim}(Y,X)=
  \inf_i\bigl(\operatorname{codim}(Y_i,X)\bigr).
  \tag{14.2.5.1}\label{0.14.2.5.1-es}
\]
\end{proposition}

\begin{proof}
En efecto, toda componente irreducible de uno de los $Y_i$ está contenida en
una componente irreducible de $Y$ y, recíprocamente, toda componente
irreducible de $Y$ es también componente irreducible de uno de los $Y_i$
($0_{\mathrm I}$, 2.1.1); la conclusión resulta, por tanto, de la definición
(14.2.1) y de la desigualdad (14.2.2.3).
\end{proof}

\begin{corollary}[14.2.6]
\label{0.14.2.6-es}
Sean $X$ un espacio topológico y $Y$ un subespacio cerrado localmente
noetheriano de $X$.
\begin{enumerate}
  \item[\rm(i)] Para todo $x\in X$, solo existe un número finito de
  componentes irreducibles $Y_i$ ($1\leq i\leq n$) de $Y$ que contienen $x$,
  y se tiene
  $\operatorname{codim}_x(Y,X)=
  \inf_i(\operatorname{codim}(Y_i,X))$.
  \item[\rm(ii)] La función
  $x\mapsto\operatorname{codim}_x(Y,X)$ es semicontinua inferiormente en $X$.
\end{enumerate}
\end{corollary}

\begin{proof}
En efecto, por hipótesis existe un entorno abierto $U_0$ de $x$ en $X$ tal
que $Y\cap U_0$ es noetheriano y, por tanto, solo tiene un número finito de
componentes irreducibles, que son las trazas sobre $U_0$ de componentes
irreducibles de $Y$; \emph{a fortiori}, solo existe un número finito de
componentes irreducibles $Y_i$ ($1\leq i\leq n$) de $Y$ que contienen $x$, y
se puede, sustituyendo $U_0$ por un entorno abierto $U\subset U_0$ de $x$ que
no tenga intersección con ninguno de los $Y_j$ que no contienen $x$, suponer
que los $Y_i\cap U$ son las componentes irreducibles de $Y\cap U$; para todo
abierto $V\subset U$ de $x$ en $X$, los $Y_i\cap V$ son entonces las
componentes irreducibles de $Y\cap V$, y (14.2.3) muestra entonces que
$\operatorname{codim}(Y_i,X)=\operatorname{codim}(Y_i\cap V,V)$, lo que
prueba (i). Además, para todo $x'\in U$, las componentes irreducibles de $Y$
que contienen $x'$ son algunas de las $Y_i$, por lo que
$\operatorname{codim}_{x'}(Y,X)\geq\operatorname{codim}_x(Y,X)$, lo que
demuestra (ii).
\end{proof}

\subsection{La condición de las cadenas}
\label{subsection:0.14.3-es}

\begin{env}[14.3.1]
\label{0.14.3.1-es}
En un espacio topológico $X$, diremos que una cadena
$Z_0\subset Z_1\subset\cdots\subset Z_n$ de subconjuntos cerrados irreducibles es
\emph{saturada} si no existe ningún subconjunto cerrado irreducible $Z'$ distinto
de los $Z_i$ y tal que $Z_k\subset Z'\subset Z_{k+1}$ para algún índice $k$.
\end{env}

\begin{proposition}[14.3.2]
\label{0.14.3.2-es}
Sea $X$ un espacio topológico tal que, para cualesquiera dos subconjuntos
cerrados irreducibles $Y$, $Z$ de $X$ tales que $Y\subset Z$, se tenga
$\operatorname{codim}(Y,Z)<+\infty$. Las dos condiciones siguientes son
equivalentes:
\begin{enumerate}
  \item[\rm(a)] Dos cadenas saturadas de subconjuntos cerrados irreducibles de $X$,
  con los mismos extremos, tienen la misma longitud.
  \item[\rm(b)] Si $Y$, $Z$, $T$ son tres subconjuntos cerrados irreducibles de
  $X$ tales que $Y\subset Z\subset T$, se tiene
  \[
    \operatorname{codim}(Y,T)=
    \operatorname{codim}(Y,Z)+\operatorname{codim}(Z,T).
    \tag{14.3.2.1}\label{0.14.3.2.1-es}
  \]
\end{enumerate}
\end{proposition}

\begin{proof}
Es inmediato que a) implica b). Recíprocamente, supongamos que se verifica b) y
demostremos que, si dos cadenas saturadas con los mismos extremos tienen
longitudes $m$ y $n\leq m$, se tiene necesariamente $m=n$. Razonemos por
inducción sobre $n$, siendo evidente la proposición para $n=1$. Supongamos,
pues, que $n>1$, $n<m$, y sea $Z_0\subset Z_1\subset\cdots\subset Z_n$ una
\oldpage[0\textsubscript{IV-1}]{11}
cadena saturada tal que existe otra cadena saturada con extremos
$Z_0$, $Z_n$ y de longitud $m$. Como
$\operatorname{codim}(Z_0,Z_n)\geq m>n$ y
$\operatorname{codim}(Z_0,Z_1)=1$, se deduce de b) que
$\operatorname{codim}(Z_1,Z_n)=\operatorname{codim}(Z_0,Z_n)-1>n-1$, lo que
contradice la hipótesis de inducción.
\end{proof}

Cuando se cumplen las condiciones de (14.3.2), se dice que $X$ satisface la
\emph{condición de las cadenas}, o también que es un \emph{espacio catenario}.
Es claro que todo subespacio cerrado de un espacio catenario es catenario.

\begin{proposition}[14.3.3]
\label{0.14.3.3-es}
Sea $X$ un espacio de Kolmogórov noetheriano de dimensión finita. Las
condiciones siguientes son equivalentes:
\begin{enumerate}
  \item[\rm(a)] Dos cadenas maximales de subconjuntos cerrados irreducibles de
  $X$ tienen la misma longitud.
  \item[\rm(b)] $X$ es equidimensional, equicodimensional y catenario.
  \item[\rm(c)] $X$ es equidimensional y, para dos subconjuntos cerrados
  irreducibles $Y$, $Z$ de $X$ tales que $Y\subset Z$, se tiene
  \[
    \dim(Z)=\dim(Y)+\operatorname{codim}(Y,Z).
    \tag{14.3.3.1}\label{0.14.3.3.1-es}
  \]
  \item[\rm(d)] $X$ es equicodimensional y, para dos subconjuntos cerrados
  irreducibles $Y$, $Z$ de $X$ tales que $Y\subset Z$, se tiene
  \[
    \operatorname{codim}(Y,X)=
    \operatorname{codim}(Y,Z)+\operatorname{codim}(Z,X).
    \tag{14.3.3.2}\label{0.14.3.3.2-es}
  \]
\end{enumerate}
\end{proposition}

\begin{proof}
Las hipótesis sobre $X$ implican que los extremos de una cadena maximal de
subconjuntos cerrados irreducibles de $X$ son necesariamente un punto cerrado y
una componente irreducible de $X$ ($0_{\mathrm I}$, 2.1.3); además, toda
cadena saturada con extremos $Y$, $Z$ (con $Y\subset Z$) está contenida en una
cadena maximal cuyos elementos distintos de los de la cadena dada o bien están
contenidos en $Y$, o bien contienen a $Z$. Estas observaciones establecen de
inmediato la equivalencia de a) y b), y prueban también que, si se verifica a),
se tiene, para todo subconjunto cerrado irreducible $Y$ de $X$,
\[
  \dim(Y)+\operatorname{codim}(Y,X)=\dim(X)~;
  \tag{14.3.3.3}\label{0.14.3.3.3-es}
\]
como se cumple (14.3.2.1), se deducen inmediatamente (14.3.3.1) y (14.3.3.2)
de (14.3.3.3). Recíprocamente, (14.3.3.1) implica (14.3.2.1), luego
(14.3.3.1) implica la condición de las cadenas en virtud de (14.3.2); además,
aplicando (14.3.3.1) al caso en que $Y$ se reduce a un punto cerrado $x$ de $X$
y $Z$ a una componente irreducible de $X$, se obtiene
$\operatorname{codim}(\{x\},X)=\dim(Z)$; se concluye que c) implica b).
Análogamente, (14.3.3.2) implica (14.3.2.1), luego la condición de las cadenas;
además, con la misma elección de $Y$ y $Z$ que antes, (14.3.3.2) implica
también $\operatorname{codim}(\{x\},X)=\dim(Z)$; por tanto (toda componente
irreducible de $X$ contiene un punto cerrado en virtud de ($0_{\mathrm I}$,
2.1.3)), d) implica b).
\end{proof}

Se dice que un espacio de Kolmogórov noetheriano es \emph{biequidimensional}
si es de \emph{dimensión finita} y verifica las condiciones equivalentes de
(14.3.3).

\begin{corollary}[14.3.4]
\label{0.14.3.4-es}
Sea $X$ un espacio de Kolmogórov noetheriano biequidimensional; entonces, para
todo punto cerrado $x$ de $X$ y toda componente irreducible $Z$ de $X$, se tiene
\[
  \dim(X)=\dim(Z)=\operatorname{codim}(\{x\},X)=\dim_x(X).
  \tag{14.3.4.1}\label{0.14.3.4.1-es}
\]
\end{corollary}

\oldpage[0\textsubscript{IV-1}]{12}
\begin{proof}
La última igualdad se deduce del hecho de que, si
$Y_0=\{x\}\subset Y_1\subset\cdots\subset Y_m$ es una cadena maximal de
subconjuntos cerrados irreducibles de $X$ y $U$ un entorno abierto de $x$, los
$U\cap Y_i$ son subconjuntos cerrados irreducibles de $U$ distintos dos a dos
(pues $\overline{U\cap Y_i}=Y_i$), de donde
$\dim(U)=\dim(X)$ en virtud de (14.1.4).
\end{proof}

\begin{corollary}[14.3.5]
\label{0.14.3.5-es}
Sea $X$ un espacio de Kolmogórov noetheriano; si $X$ es biequidimensional, lo
mismo ocurre con toda unión de componentes irreducibles de $X$ y con todo
subconjunto cerrado irreducible de $X$. Además, para todo subconjunto cerrado
$Y$ de $X$, se tiene entonces
\[
  \dim(Y)+\operatorname{codim}(Y,X)=\dim(X).
  \tag{14.3.5.1}\label{0.14.3.5.1-es}
\]
\end{corollary}

\begin{proof}
Como toda cadena de subconjuntos cerrados irreducibles de $X$ está contenida en
una componente irreducible de $X$, la primera afirmación se deduce
inmediatamente de (14.3.3). Además, si $X'$ es un subconjunto cerrado
irreducible de $X$, $X'$ verifica trivialmente las condiciones (14.3.3, c)), de
donde se sigue la segunda afirmación.

Por último, para demostrar (14.3.5.1), observemos que se ha visto en la
demostración de (14.3.3) que esta relación se verifica cuando $Y$ es
irreducible; si $Y_i$ ($1\leq i\leq m$) son las componentes irreducibles de
$Y$, aquella entre las $Y_i$ para la cual $\dim(Y_i)$ es mayor es también
aquella para la cual $\operatorname{codim}(Y_i,X)$ es menor; por tanto,
(14.3.5.1) se deduce de las definiciones de $\dim(Y)$ y de
$\operatorname{codim}(Y,X)$.
\end{proof}

\begin{remark}[14.3.6]
\label{0.14.3.6-es}
El lector observará que la demostración de (14.3.2) se aplica a un conjunto
ordenado cualquiera, pues el hecho de que se trate de subconjuntos cerrados
irreducibles de un espacio topológico no interviene en esta demostración. Lo
mismo ocurre con la demostración de (14.3.3) en un conjunto ordenado $E$ tal
que, para todo $x\in E$, existe un $z\leq x$ que es \emph{minimal} en $E$, y en
el que la longitud de las cadenas de elementos de $E$ está acotada.
\end{remark}

\section{Sucesiones \texorpdfstring{$M$}{M}-regulares y sucesiones
\texorpdfstring{$\mathscr F$}{F}-regulares}
\label{section:0.15-es}

\subsection{Sucesiones \texorpdfstring{$M$}{M}-regulares y sucesiones
\texorpdfstring{$M$}{M}-cuasirregulares}
\label{subsection:0.15.1-es}
\label{0.15.1-es}

\begin{env}[15.1.1]
\label{0.15.1.1-es}
\oldpage[0\textsubscript{IV-1}]{12}
Sean $A$ un anillo y $N$ un $A$-módulo~; en lo que sigue, denotaremos
abreviadamente por $N[T_1,\ldots,T_n]$ el $A$-módulo graduado
$N\otimes_A A[T_1,\ldots,T_n]$ (siendo los $T_i$ indeterminadas), donde $N$
lleva la graduación trivial y $A[T_1,\ldots,T_n]$ la graduación por el grado total.

Sean $M$ un $A$-módulo y $f_i$ ($1\leq i\leq n$) elementos de $A$;
sea $\mathfrak{J}=\sum_{i=1}^n f_iA$ el ideal de $A$ engendrado por los $f_i$,
y dotemos a $M$ de la filtración $\mathfrak{J}$-preádica~; se define entonces un
\emph{homomorfismo graduado sobreyectivo de $A$-módulos graduados}, de grado $0$,
\[
  \varphi : (\gr_0(M))[T_1,\ldots,T_n]\longrightarrow\gr_\bullet(M)
  \tag{15.1.1.1}\label{0.15.1.1.1-es}
\]
de la manera siguiente~: si $\xi\in\gr_0(M)$ es la clase módulo
$\mathfrak{J}M$ de un elemento $x\in M$, al par formado por $\xi$ y un
polinomio $P[T_1,\ldots,T_n]$ homogéneo de grado $k$ se le hace corresponder la
clase módulo $\mathfrak{J}^{k+1}M$ de $P(f_1,\ldots,f_n)x$~; es inmediato
que esta clase
\oldpage[0\textsubscript{IV-1}]{13}
solo depende de $\xi$ y de $P$, y, si $S_k$ es el conjunto de los polinomios
homogéneos de grado $k$ de $A[T_1,\ldots,T_n]$, queda así definida una
aplicación $A$-lineal
\[
  \varphi_k:\gr_0(M)\otimes_A S_k\longrightarrow\gr_k(M),
\]
que es sobreyectiva por la definición de $\mathfrak{J}^kM$.

Diremos en lo que sigue que un elemento $f\in A$ es \emph{no divisor de
cero en un $A$-módulo $M$} si la homotecia $f_M:x\mapsto fx$ de $M$ es
\emph{inyectiva}.
\end{env}

\begin{lemma}[15.1.2]
\label{0.15.1.2-es}
Sean $M$ un $A$-módulo y $f$ un elemento de $A$, y dotemos a $M$ de la
filtración $(fA)$-preádica. Si $f$ es no divisor de cero en $M$,
el homomorfismo canónico definido en (15.1.1)~:
\[
  \varphi:(M/fM)[T]\longrightarrow\gr_\bullet(M)
  \tag{15.1.2.1}\label{0.15.1.2.1-es}
\]
es biyectivo. El recíproco es cierto si $M$ es separado para la topología
$(fA)$-preádica.
\end{lemma}

\begin{proof}
Se tiene aquí $\gr_k(M)=f^kM/f^{k+1}M$~; decir que $\varphi$ es inyectiva
significa que, para todo $x\in M$ y todo $k$, la relación
$f^kx\in f^{k+1}M$ implica $x\in fM$~; es claro que así ocurre si $f$ es no
divisor de cero en $M$. Recíprocamente, obsérvese que la homotecia
$f_M:x\mapsto fx$ define, por paso a los cocientes, homomorfismos
$\gr_k(M)\to\gr_{k+1}(M)$, y, por tanto, un homomorfismo graduado $g$ de grado
$1$ de $\gr_\bullet(M)$ en sí mismo~; decir que $\varphi$ es inyectivo expresa
también que $g$ es inyectivo. Si se supone separada la filtración de $M$, de
ello resulta que $f_M$ es inyectiva (Bourbaki, \emph{Alg.~comm.}, cap.~III,
\S~2, núm.~8, cor.~1 del teorema~1).
\end{proof}

\begin{corollary}[15.1.3]
\label{0.15.1.3-es}
Sean $A$ un anillo noetheriano y $M$ un $A$-módulo de tipo finito. Para que
$f$ sea no divisor de cero en $M$, es necesario y suficiente que $\varphi$ sea
biyectivo.
\end{corollary}

\begin{proof}
Solo hay que demostrar que la condición es suficiente~; para probar que $f_M$
es inyectiva, basta mostrar que, para todo ideal primo $\mathfrak p$ de $A$,
la homotecia de razón $f'=f/1\in A'=A_{\mathfrak p}$ en el
$A_{\mathfrak p}$-módulo $M_{\mathfrak p}=M'$ es inyectiva (Bourbaki,
\emph{Alg.~comm.}, cap.~II, \S~3, núm.~4, teorema~1). Ahora bien, se tiene
$f'^kM'/f'^{k+1}M'=(f^kM/f^{k+1}M)\otimes_A A'$, y es inmediato que
el homomorfismo canónico $\varphi':(M'/f'M')[T]\to\gr_\bullet(M')$ es igual a
$\varphi\otimes 1_{A'}$~; si $\varphi$ es inyectivo, lo mismo vale, por tanto,
para $\varphi'$ ($0_{\mathrm I}$, 1.3.2). Queda así reducido al caso en que $A$
es un anillo noetheriano \emph{local}, cuyo ideal maximal denotaremos por
$\mathfrak m$. Si $f\notin\mathfrak m$, entonces $f$ es invertible y no hay
nada que demostrar. Si $f\in\mathfrak m$, se sabe que $M$ es separado para la
topología $(fA)$-preádica ($0_{\mathrm I}$, 7.3.5), y basta, por tanto,
aplicar (15.1.2).
\end{proof}

\begin{definition}[15.1.4]
\label{0.15.1.4-es}
Sean $A$ un anillo y $M$ un $A$-módulo. Se dice que un elemento $f\in A$ es
\emph{$M$-regular} (resp. \emph{$M$-cuasirregular}) si $f$ es no divisor de
cero en $M$ (resp. si el homomorfismo (15.1.2.1) es biyectivo).
\end{definition}

Acabamos, pues, de demostrar que todo elemento $M$-regular es
$M$-cuasirregular, y que el recíproco es cierto si $A$ es noetheriano y $M$
un $A$-módulo de tipo finito. Cuando ya no se supone que $M$ sea de tipo
finito, este recíproco puede fallar~: por ejemplo, si $A$ es un anillo
principal que no es un cuerpo, $K$ su cuerpo de fracciones y $M$ el
$A$-módulo $K/A$, se tiene $fM=M$ para todo $f\neq0$ en $A$, por lo que los
dos miembros de (15.1.2.1) se reducen a $0$, pero $f$ es entonces divisor de
cero en $M$.

\begin{env}[15.1.5]
\label{0.15.1.5-es}
\oldpage[0\textsubscript{IV-1}]{14}
Supongamos ahora que $M$ está dotado de una \emph{filtración} $(M_k)$ tal que
$M_0=M$, y denotemos por
$\gr_\bullet(M)=\bigoplus_{k\geq0}M_k/M_{k+1}$ el $A$-módulo graduado asociado.
Con las notaciones de (15.1.1), pongamos, para todo $k\geq0$,
\[
  M'_k=M_k+\mathfrak{J}M_{k-1}+\cdots+
  \mathfrak{J}^{k-1}M_1+\mathfrak{J}^kM_0.
  \tag{15.1.5.1}\label{0.15.1.5.1-es}
\]
Es claro que $(M'_k)$ es una filtración sobre $M$. Cuando
$M_k=\mathfrak{Q}^kM$, donde $\mathfrak Q$ es un ideal de $A$, la filtración
$(M'_k)$ no es otra que la filtración
$(\mathfrak{J}+\mathfrak{Q})$-preádica. Denotaremos por
$\gr'_\bullet(M)$ el $A$-módulo graduado asociado a $(M'_k)$. Denotemos asimismo
\[
  (\gr_\bullet(M)\otimes_A(A/\mathfrak{J}))[T_1,\ldots,T_n]
\]
el producto tensorial de $A$-módulos graduados ($\mathrm{II}$, 2.1.2)
\[
  (\gr_\bullet(M)\otimes_A(A/\mathfrak{J}))
  \otimes_A A[T_1,\ldots,T_n].
\]
Se define un \emph{homomorfismo graduado sobreyectivo de $A$-módulos graduados},
de grado $0$,
\[
  \psi:(\gr_\bullet(M)\otimes_A(A/\mathfrak{J}))[T_1,\ldots,T_n]
  \longrightarrow\gr'_\bullet(M)
  \tag{15.1.5.2}\label{0.15.1.5.2-es}
\]
de la manera siguiente~: si $\alpha\in A/\mathfrak{J}$ es la clase de un
elemento $a\in A$, $\xi\in M_k/M_{k+1}$ la clase de un elemento $x\in M_k$ y
$P[T_1,\ldots,T_n]$ un polinomio homogéneo de grado $h$, se hace corresponder
a la terna $(\alpha,\xi,P)$ la clase módulo $M'_{h+k+1}$ de
$P(f_1,\ldots,f_n)ax$~; se comprueba de nuevo inmediatamente que este elemento
solo depende de $\xi$, $\alpha$ y $P$, y esto define el homomorfismo graduado
(15.1.5.2), que es sobreyectivo en virtud de la definición (15.1.5.1). Se
recupera el homomorfismo (15.1.1.1) considerando la filtración sobre $M$ tal
que $M_1=0$.
\end{env}

\begin{lemma}[15.1.6]
\label{0.15.1.6-es}
Se conservan las notaciones de (15.1.5) y se supone $n=1$, poniendo $f=f_1$.
Si $f$ es $\gr_\bullet(M)$-regular, el homomorfismo (15.1.5.2) es biyectivo.
\end{lemma}

\begin{proof}
Sea $Q_k$ (resp. $Q'_k$) el sub-$A$-módulo de los términos de grado $k$ en
el primer (resp. segundo) miembro de (15.1.5.2). Dotemos a $Q_k$ de la filtración
\[
  (Q_k)_i=\sum_{j\leq k-i}
  (\gr_{k-j}(M)\otimes_A(A/fA))T^j
  \tag{15.1.6.1}\label{0.15.1.6.1-es}
\]
de modo que $(Q_k)_0=Q_k$ y $(Q_k)_{k+1}=0$, y, para esta filtración, se tiene
\[
  \gr_i(Q_k)=(\gr_i(M)\otimes_A(A/fA))T^{k-i}.
\]
Dotemos a $Q'_k$ de la filtración formada por los
$\psi((Q_k)_i)=(Q'_k)_i$. Como estas filtraciones son \emph{finitas}, basta
probar que los homomorfismos $\psi_{ki}:\gr_i(Q_k)\to\gr_i(Q'_k)$ deducidos de
$\psi_k$ son inyectivos (Bourbaki, \emph{Alg.~comm.}, cap.~III, \S~3,
núm.~8, cor.~1 del teorema~1). Ahora bien, se tiene
\[
  \gr_i(Q_k)=((M_i/M_{i+1})\otimes_A(A/fA))T^{k-i}~;
\]
por otro lado, $(Q'_k)_{i+1}$ es la imagen de
\[
  M_k+fM_{k-1}+\cdots+f^{k-i-1}M_{i+1}
\]
\oldpage[0\textsubscript{IV-1}]{15}
en $M'_k/M'_{k+1}$. Nos reducimos, por tanto, a demostrar que, para $x\in M_i$,
la relación
\[
  f^{k-i}x\in M_k+fM_{k-1}+\cdots+
  f^{k-i-1}M_{i+1}+M'_{k+1}
  \tag{15.1.6.2}\label{0.15.1.6.2-es}
\]
implica $x\in fM_i+M_{i+1}$. Ahora bien, el segundo miembro de (15.1.6.2) está
contenido en $M_{i+1}+M'_{k+1}$, y, en virtud de (15.1.5.1), $M'_{k+1}$ está
contenido en $M_{i+1}+f^{k+1-i}M$. La hipótesis sobre $f$ implica que $f$ es
no divisor de cero en $M/M_h$ para todo $h$ (Bourbaki, \emph{loc.~cit.}). La
relación $f^{k-i}x\in f^{k+1-i}M+M_{i+1}$ implica, pues (razonando en
$M/M_{i+1}$), la existencia de un $y\in M$ tal que $x-fy\in M_{i+1}$, y como
$x\in M_i$, se tiene $fy\in M_i$, de donde (razonando en $M/M_i$) se sigue
$y\in M_i$, y finalmente $x\in fM_i+M_{i+1}$. C.Q.D.
\end{proof}

\begin{definition}[15.1.7]
\label{0.15.1.7-es}
Sea $(f_i)_{1\leq i\leq n}$ una sucesión de elementos de $A$. Se dice que es
\emph{$M$-regular} si, para $1\leq i\leq n$, $f_i$ es
$\left(M/(\sum_{1\leq j\leq i-1}f_jM)\right)$-regular. Se dice que la sucesión
$(f_i)$ es \emph{$M$-cuasirregular} si el homomorfismo canónico $\varphi$
definido en (15.1.1.1) es biyectivo.

Cuando $M=A$, se dice simplemente «~sucesión regular~» y
«~sucesión cuasirregular~».
\end{definition}

\begin{proposition}[15.1.8]
\label{0.15.1.8-es}
Con las notaciones de (15.1.5), se supone que la sucesión
$(f_i)_{1\leq i\leq n}$ es $\gr_\bullet(M)$-regular. Entonces
el homomorfismo canónico (15.1.5.2) es biyectivo.
\end{proposition}

\begin{proof}
La proposición no es sino (15.1.6) para $n=1$~; razonemos por inducción sobre
$n$. Pongamos
\[
  \mathfrak{J}''=\sum_{i=1}^{n-1}f_iA,
\]
de modo que $\mathfrak{J}=\mathfrak{J}''+f_nA$, y sea
\[
  M''_k=M_k+\mathfrak{J}''M_{k-1}+\cdots+
  \mathfrak{J}''^{k-1}M_1+\mathfrak{J}''^kM_0~;
\]
se puede escribir, por tanto,
\[
  M'_k=M''_k+f_nM''_{k-1}+\cdots+f_n^{k-1}M''_1+f_n^kM''_0.
\]
Denotemos por $\gr''_\bullet(M)$ el $A$-módulo graduado asociado a la filtración
$(M''_k)$ de $M$. Salvo isomorfismos canónicos, se puede escribir
\[
  (\gr_\bullet(M)\otimes(A/\mathfrak{J}))\otimes_A A[T_1,\ldots,T_n]
  =
  (\gr_\bullet(M)\otimes(A/\mathfrak{J}''))[T_1,\ldots,T_{n-1}]
  \otimes(A/f_nA)\otimes_A A[T_n]
\]
y, por consiguiente, (15.1.5.2) se factoriza en
\[
  (\gr''_\bullet(M)\otimes(A/f_nA))[T_n]\longrightarrow\gr'_\bullet(M)
  \tag{15.1.8.1}\label{0.15.1.8.1-es}
\]
y
\[
  ((\gr_\bullet(M)\otimes(A/\mathfrak{J}''))[T_1,\ldots,T_{n-1}])
  \otimes(A/f_nA)\otimes_A A[T_n]
  \xrightarrow{\psi'\otimes1}
  \gr''_\bullet(M)\otimes(A/f_nA)\otimes_A A[T_n]
\]
donde $\psi'$ es la aplicación (15.1.5.2) en la que se reemplazan $n$ por
$n-1$, $\mathfrak{J}$ por $\mathfrak{J}''$ y $\gr'_\bullet(M)$ por
$\gr''_\bullet(M)$. La hipótesis de inducción implica que $\psi'$ es
biyectiva, y queda, por tanto, comprobar que lo mismo vale para (15.1.8.1).
Aplicando (15.1.6), basta mostrar que $f_n$ es
$\gr''_\bullet(M)$-regular. Pero $\gr''_\bullet(M)$ se identifica con
\[
  (\gr_\bullet(M)\otimes(A/\mathfrak{J}''))
  [T_1,\ldots,T_{n-1}]~;
\]
por hipótesis, la homotecia de razón $f_n$ en
$\gr_\bullet(M)\otimes(A/\mathfrak{J}'')$ es inyectiva~; por consiguiente, lo
mismo vale para la homotecia de razón $f_n$ en $\gr''_\bullet(M)$, puesto que
$A[T_1,\ldots,T_{n-1}]$ es un $A$-módulo libre.
\end{proof}
\begin{proposition}[15.1.9]
\label{0.15.1.9-es}
Sean $M$ un $A$-módulo y $(f_i)_{1\leq i\leq n}$ una sucesión de elementos de
$A$. Si la sucesión $(f_i)$ es $M$-regular, es $M$-cuasirregular. El recíproco
es cierto cuando se supone además que los $A$-módulos
\[
  M,\quad M/f_1M,\quad\ldots,\quad
  M/\left(\sum_{1\leq j\leq n-1}f_jM\right)
\]
son separados para la topología $\mathfrak{J}$-preádica (donde
$\mathfrak{J}=\sum_{i=1}^n f_iA$).
\end{proposition}

\begin{proof}
La primera afirmación resulta de (15.1.8), tomando $M_1=0$, y, por
consiguiente, el homomorfismo (15.1.5.2) se reduce a (15.1.1.1).
Recíprocamente, supongamos que el homomorfismo $\varphi$ de (15.1.1.1) es
biyectivo. Obsérvese que, si se pone
\[
  \mathfrak{J}''=\sum_{i=1}^{n-1}f_iA,
\]
\oldpage[0\textsubscript{IV-1}]{16}
y si se denota por $\gr''_\bullet(M)$ el $A$-módulo graduado asociado a $M$
filtrado por la filtración $\mathfrak{J}''$-preádica, $\varphi$ se factoriza en
\[
  (\gr''_\bullet(M)\otimes(A/f_nA))[T_n]\longrightarrow\gr_\bullet(M)
  \tag{15.1.9.1}\label{0.15.1.9.1-es}
\]
y en $\varphi'\otimes1$, donde $\varphi'$ es el homomorfismo canónico
\[
  (\gr''_0(M))[T_1,\ldots,T_{n-1}]
  \longrightarrow\gr''_\bullet(M).
\]
Como estos dos homomorfismos son sobreyectivos, son biyectivos si se supone que
$\varphi$ es biyectivo; por tanto, $\varphi'$ es necesariamente inyectivo y,
por consiguiente, biyectivo, puesto que es sobreyectivo. Se puede entonces
razonar por recurrencia sobre $n$ (el caso $n=1$ resulta de (15.1.2)), pues
sobre
\[
  M,\quad M/f_1M,\quad\ldots,\quad
  M/\left(\sum_{j=1}^{n-2}f_jM\right)
\]
la topología $\mathfrak{J}''$-preádica es más fina que la topología
$\mathfrak{J}$-preádica y, por consiguiente, es separada. Se ve, pues, que la
sucesión $(f_i)_{1\leq i\leq n-1}$ es $M$-regular. Por otra parte, demostremos
que la hipótesis implica que, si $y\in M/\mathfrak{J}''M$ es tal que
\[
  f_n^ky\in f_n^{k+1}(M/\mathfrak{J}''M),
\]
entonces $y\in f_n(M/\mathfrak{J}''M)$. En efecto, si $x$ es un representante
de $y$, demostremos que $f_n^kx\in\mathfrak{J}^{k+1}M$; nuestra afirmación
resultará entonces del hecho de que (15.1.9.1) es inyectivo. La hipótesis
implica
\[
  f_n^kx\in f_n^{k+1}M+\mathfrak{J}''M\subset\mathfrak{J}M,
\]
de donde, como $\varphi$ es inyectivo,
$f_n^{k-1}x\in\mathfrak{J}^kM$; por recurrencia se obtiene así
$x\in\mathfrak{J}M$, y de ahí, finalmente,
$f_n^kx\in\mathfrak{J}^{k+1}M$. Se termina el razonamiento como en
(15.1.2): sobre $M/\mathfrak{J}''M$, la filtración
$\mathfrak{J}$-preádica es idéntica a la filtración
$(f_nA)$-preádica, luego es separada por hipótesis, y se concluye que la
homotecia de razón $f_n$ en $M/\mathfrak{J}''M$ es inyectiva, lo que termina
de demostrar que la sucesión $(f_i)_{1\leq i\leq n}$ es $M$-regular.
\end{proof}

\begin{corollary}[15.1.10]
\label{0.15.1.10-es}
Con las notaciones de (15.1.9), supongamos que, para todo $A$-módulo $N$ de
tipo finito, $M\otimes_A N$ es separado para la topología
$\mathfrak{J}$-preádica. Entonces, para que la sucesión $(f_i)$ sea
$M$-regular, es necesario y suficiente que sea $M$-cuasirregular.
\end{corollary}

Obsérvese que la hipótesis de (15.1.10) implica que, si la sucesión
$(f_i)_{1\leq i\leq n}$ es $M$-regular, también lo es la sucesión
$(f_{\sigma(i)})$ para toda permutación $\sigma$ del intervalo $[1,n]$.

\begin{corollary}[15.1.11]
\label{0.15.1.11-es}
Sean $A$ un anillo noetheriano, $M$ un $A$-módulo de tipo finito y
$(f_i)_{1\leq i\leq n}$ una sucesión finita de elementos contenidos en el
radical de $A$. Entonces, para que la sucesión $(f_i)$ sea $M$-regular, es
necesario y suficiente que sea $M$-cuasirregular.
\end{corollary}

\begin{proof}
Es un caso particular de (15.1.10), puesto que todo $A$-módulo de tipo finito
es entonces separado para la topología $\mathfrak{J}$-preádica
($0_{\mathrm I}$, 7.3.5).
\end{proof}

\begin{remarks}[15.1.12]
\label{0.15.1.12-es}
\mbox{}
\begin{enumerate}
  \item[(i)] Cuando $A$ es un anillo local noetheriano, $M$ un $A$-módulo de
  tipo finito, y cuando los $f_i$ están contenidos en el ideal maximal de $A$,
  la noción de sucesión $M$-regular fue introducida por J.-P.~Serre con el
  nombre de $M$-\emph{sucesión} [17].
  \item[(ii)] Incluso si $A$ es noetheriano y $M=A$, una sucesión $(f,g)$ de
  dos elementos de $A$ que es cuasirregular no es necesariamente regular. Por
  ejemplo, la sucesión $(0,1)$ es cuasirregular, puesto que entonces
  $\mathfrak{J}=fA+gA=A$, luego $\gr_\bullet(A)=0$; pero, si $A\neq0$, no es
  regular; obsérvese, no obstante, que en este caso la sucesión $(1,0)$ es
  regular, pues $1$ no es divisor de $0$ en $A$, y como $A/A=0$, $0$ no es
  divisor de cero en este módulo (cf. (15.1.1)).
  \item[(iii)]
  \oldpage[0\textsubscript{IV-1}]{17}
  En (15.1.6), en general no se puede sustituir la hipótesis de que $f$ es
  $\gr_\bullet(M)$-regular por la hipótesis de que es
  $\gr_\bullet(M)$-cuasirregular. Para verlo, retomemos el ejemplo (15.1.4),
  donde $A$ es un anillo principal de cuerpo de fracciones $K\neq A$, y
  tomemos $M=K$, definiéndose la filtración $(M_k)$ de $M$ por
  $M_0=M=K$, $M_1=A$, $M_k=0$ para $k\geq2$, de donde
  $\gr_0(M)=K/A$, $\gr_1(M)=A$, $\gr_k(M)=0$ para $k\geq2$. Si $f\neq0$ en
  $A$, se ha visto que $f$ no es $\gr_\bullet(M)$-regular, pero, como se tiene
  \[
    f^k\gr_\bullet(M)/f^{k+1}\gr_\bullet(M)=f^kA/f^{k+1}A,
  \]
  es inmediato que el homomorfismo (15.1.2.1) (donde $M$ se sustituye por
  $\gr_\bullet(M)$) es inyectivo; por tanto, $f$ es
  $\gr_\bullet(M)$-cuasirregular. Pero, con las notaciones de (15.1.6), se
  tiene entonces $M'_k=K$ para todo $k\geq0$, luego
  $\gr'_\bullet(M)=0$, mientras que el miembro izquierdo de (15.1.5.2) no se
  reduce a $0$.
\end{enumerate}
\end{remarks}

\begin{proposition}[15.1.13]
\label{0.15.1.13-es}
Sean $A$ un anillo, $M$ y $N$ dos $A$-módulos, y
$(f_i)_{1\leq i\leq n}$ una sucesión de elementos de $A$. Si la sucesión
$(f_i)$ es $M$-regular y si $N$ es un $A$-módulo plano, entonces la sucesión
$(f_i)$ es $M\otimes_A N$-regular. El recíproco es cierto si $N$ es un
$A$-módulo fielmente plano.
\end{proposition}

\begin{proof}
En efecto (sin hipótesis sobre $N$),
\[
  (M\otimes_A N)/
  \left(\sum_{j=1}^{i-1}f_j(M\otimes_A N)\right)
\]
es isomorfo a
\[
  \left(M/\left(\sum_{j=1}^{i-1}f_jM\right)\right)\otimes_A N~;
\]
basta entonces aplicar ($0_{\mathrm I}$, 6.1.1) y
($0_{\mathrm I}$, 6.4.1) a la homotecia de razón $f_i$ en
$M/(\sum_{j=1}^{i-1}f_jM)$.
\end{proof}

\begin{corollary}[15.1.14]
\label{0.15.1.14-es}
Sean $A$ un anillo, $M$ un $A$-módulo y $(f_i)_{1\leq i\leq n}$ una sucesión
finita de elementos de $A$. Sean $\rho:A\to B$ un homomorfismo de anillos,
$M'=M_{(B)}=M\otimes_A B$ y $f'_i=\rho(f_i)$ para $1\leq i\leq n$. Si $B$ es
un $A$-módulo plano y si la sucesión $(f_i)$ es $M$-regular, entonces la
sucesión $(f'_i)$ es $M'$-regular; el recíproco es cierto si $B$ es un
$A$-módulo fielmente plano.
\end{corollary}

\begin{proof}
En efecto,
\[
  M'/\left(\sum_{j=1}^{i-1}f'_jM'\right)
\]
se identifica con
\[
  \left(M/\left(\sum_{j=1}^{i-1}f_jM\right)\right)\otimes_A B,
\]
y la homotecia de razón $f'_i$ en el $B$-módulo
$M'/(\sum_{j=1}^{i-1}f'_jM')$ se identifica con la homotecia de razón $f_i$
en el $A$-módulo
\[
  \left(M/\left(\sum_{j=1}^{i-1}f_jM\right)\right)\otimes_A B~;
\]
la afirmación resulta, por tanto, de (15.1.13).
\end{proof}

\begin{proposition}[15.1.15]
\label{0.15.1.15-es}
Sean $A$ un anillo, $B$ una $A$-álgebra conmutativa,
$(f_i)_{1\leq i\leq n}$ una sucesión finita de elementos de $B$ y $M$ un
$B$-módulo. Supongamos que la sucesión $(f_i)$ es $M$-regular y que cada uno
de los $A$-módulos
\[
  M/\left(\sum_{j=1}^{i-1}f_jM\right)\quad(1\leq i\leq n)
\]
es plano. Entonces, para todo homomorfismo de anillos $\rho:A\to A'$, si se
pone $B'=B\otimes_A A'$, $M'=M\otimes_A A'$ y $f'_i=f_i\otimes1$
($1\leq i\leq n$), la sucesión $(f'_i)$ es $M'$-regular y los
$A'$-módulos
\[
  M'/\left(\sum_{j=1}^{i-1}f'_jM'\right)\quad(1\leq i\leq n)
\]
son planos.
\end{proposition}

\begin{proof}
Consideremos la sucesión (exacta por hipótesis)
\[
  0\longrightarrow M\xrightarrow{f_1}M\longrightarrow M/f_1M
  \longrightarrow0
\]
(siendo la flecha $M\xrightarrow{f_1}M$ la homotecia de razón $f_1$ en
$M$). La hipótesis de que $M/f_1M$ es $A$-plano implica que la sucesión
\[
  0\longrightarrow M\otimes_A A'
  \xrightarrow{f_1\otimes1}M\otimes_A A'
  \longrightarrow(M/f_1M)\otimes_A A'\longrightarrow0
\]
\oldpage[0\textsubscript{IV-1}]{18}
es exacta ($0_{\mathrm I}$, 6.1.2). Esto muestra que la homotecia de razón
$f'_1$ en $M'$ es inyectiva.

Pongamos a continuación, para $1\leq i\leq n-1$,
\[
  M_i=M/\left(\sum_{j=1}^i f_jM\right),\qquad
  M'_i=M'/\left(\sum_{j=1}^i f'_jM'\right)~;
\]
se tiene $M'_i=M_i\otimes_A A'$, $M_i/f_{i+1}M_i=M_{i+1}$,
$M'_i/f'_{i+1}M'_i=M'_{i+1}$; basta sustituir en el razonamiento anterior
$M$ y $f_1$ por $M_i$ y $f_{i+1}$ para concluir que la homotecia de razón
$f'_{i+1}$ en $M'_i$ es inyectiva, en virtud de la planitud de
$M_i/f_{i+1}M_i$. La última afirmación resulta de ($0_{\mathrm I}$, 6.2.1).
\end{proof}
\begin{proposition}[15.1.16]
\label{0.15.1.16-es}
Sean $A$, $B$ dos anillos locales noetherianos, $k$ el cuerpo residual de
$A$, $\varphi:A\to B$ un homomorfismo local y $M$ un $B$-módulo de tipo finito.
Sea $(f_i)_{1\leq i\leq n}$ una sucesión de elementos del ideal máximo de $B$
y, para cada $i$, sea $g_i$ la imagen de $f_i$ en $B\otimes_A k$. Las
condiciones siguientes son equivalentes:
\begin{enumerate}
  \item[\rm(a)] La sucesión $(f_i)$ es $M$-regular y los cocientes
  \[
    M_i=M/\left(\sum_{j=1}^i f_jM\right)
  \]
  son $A$-módulos planos para $1\leq i\leq n$.
  \item[\rm(b)] La sucesión $(f_i)$ es $M$-regular y el $A$-módulo $M_n$ es
  plano.
  \item[\rm(c)] El $A$-módulo $M$ es plano y la sucesión
  $(g_i)_{1\leq i\leq n}$ es $(M\otimes_A k)$-regular.
  \item[\rm(d)] El $A$-módulo $M$ es plano y, para todo homomorfismo de
  anillos $\rho:A\to A'$, si se pone $B'=B\otimes_A A'$,
  $M'=M\otimes_A A'$, $f'_i=f_i\otimes1$ ($1\leq i\leq n$), la sucesión
  $(f'_i)$ es $M'$-regular.
\end{enumerate}
\end{proposition}

\begin{proof}
Es claro que a) implica b); para demostrar que b) implica d), basta, en virtud
de (15.1.15), demostrar que b) implica que los $M_i$ son $A$-módulos planos
para $0\leq i\leq n-1$ (poniendo $M_0=M$). Como
$M_{i+1}=M_i/f_{i+1}M_i$, por inducción descendente nos reducimos a probar que,
si $h$ es un elemento del ideal máximo de $B$ tal que la multiplicación por
$h$ en $M$ sea inyectiva y que $M/hM$ sea un $A$-módulo plano, entonces $M$
es un $A$-módulo plano; pero esto no es sino
($0_{\mathrm{III}}$, 10.2.7). Es claro que c) es el caso particular de d)
obtenido tomando $A'=k$. Queda, pues, probar que c) implica a). Como $M$ es un
$A$-módulo plano y un $B$-módulo de tipo finito, se deduce de c) y de
($0_{\mathrm{III}}$, 10.2.4) que la multiplicación por $f_1$ en $M$ es
inyectiva y que su conúcleo $M/f_1M=M_1$ es un $A$-módulo plano. Razonemos por
inducción sobre $i$ y supongamos demostrado que $M_i$ es un $A$-módulo plano;
como la multiplicación por $g_{i+1}$ en $M_i\otimes_A k$ es inyectiva por
hipótesis, se deduce de nuevo de ($0_{\mathrm{III}}$, 10.2.4) que la
multiplicación por $f_{i+1}$ en $M_i$ es inyectiva y que
$M_i/f_{i+1}M_i=M_{i+1}$ es un $A$-módulo plano, lo que termina la
demostración. Se observará que la demostración de que c) implica a) no utiliza
el hecho de que los $f_i$ pertenezcan al ideal máximo de $B$.
\end{proof}

\begin{corollary}[15.1.17]
\label{0.15.1.17-es}
Bajo las hipótesis sobre $A$ y $B$ de (15.1.16), sean $M$ un $B$-módulo de
tipo finito, plano sobre $A$, y $R$ un sub-$B$-módulo de $M$ tal que $N=M/R$
sea un $A$-módulo plano. Sea $(f_i)_{1\leq i\leq n}$ una sucesión de elementos
del ideal máximo de $B$ que tenga las propiedades siguientes:
\begin{enumerate}
  \item[(i)] $f_iM\subset R$ para $1\leq i\leq n$.
  \item[(ii)] Si $g_i$ es la imagen de $f_i$ en $B\otimes_A k$, la sucesión
  $(g_i)_{1\leq i\leq n}$ es $(M\otimes_A k)$-regular.
  \item[(iii)] La suma de los sub-$(B\otimes_A k)$-módulos
  $g_i(M\otimes_A k)$ de $M\otimes_A k$ es la imagen canónica de
  $R\otimes_A k$ en $M\otimes_A k$ (igual al núcleo del homomorfismo canónico
  $M\otimes_A k\to N\otimes_A k$).
\end{enumerate}
Entonces la sucesión $(f_i)$ es $M$-regular y se tiene
$R=\sum_{i=1}^n f_iM$.
\end{corollary}

\begin{proof}
\oldpage[0\textsubscript{IV-1}]{19}
Como $N$ es un $A$-módulo plano, se deduce de la exactitud de la sucesión
$0\to R\to M\to N\to0$ que la sucesión
\[
  0\longrightarrow R\otimes_A k\longrightarrow M\otimes_A k
  \longrightarrow N\otimes_A k\longrightarrow0
\]
es exacta ($0_{\mathrm I}$, 6.1.2); se tiene, pues,
$R\otimes_A k=\sum_i g_i(M\otimes_A k)$. Como $M$ es un $B$-módulo de tipo
finito y $B$ es noetheriano, $R$ es también un $B$-módulo de tipo finito, de
donde $R=\sum_i f_iM$ en virtud del lema de Nakayama, pues los $f_i$ pertenecen
al ideal máximo de $B$.
\end{proof}

\begin{lemma}[15.1.18]
\label{0.15.1.18-es}
Sean $A$ un anillo, $M$ un $A$-módulo provisto de una filtración separada y
exhaustiva $(M_k)_{k\geq0}$, y sea $f\in A$ un elemento
$\gr_\bullet(M)$-regular. Entonces:
\begin{enumerate}
  \item[(i)] $f$ es $M$-regular.
  \item[(ii)] Para todo $k$, se tiene $M_k\cap fM=fM_k$, y la imagen canónica
  $M'_k$ de $M_k$ en $M'=M/fM$ es, por tanto, canónicamente isomorfa a
  $M_k/fM_k$.
  \item[(iii)] La sucesión
  \[
    0\longrightarrow M_{k+1}/fM_{k+1}\longrightarrow M_k/fM_k
    \longrightarrow\gr_k(M)/f\cdot\gr_k(M)\longrightarrow0
  \]
  es exacta y, si se dota a $M'$ de la filtración cociente $(M'_k)$,
  $\gr_\bullet(M')$ es, por tanto, canónicamente isomorfo a
  $\gr_\bullet(M)/f\cdot\gr_\bullet(M)$.
\end{enumerate}
\end{lemma}

\begin{proof}
La afirmación (i) se deduce de la hipótesis sobre $f$ y de Bourbaki,
\emph{Alg.~comm.}, cap.~III, \S~2, n.º~8, cor.~1 del teo.~1. Como
$\gr_\bullet(M/M_k)$ es un sumando directo de $\gr_\bullet(M)$, $f$ es
$\gr_\bullet(M/M_k)$-regular para todo $k$, luego $(M/M_k)$-regular por (i),
lo que establece (ii). Para probar (iii), basta ver que la aplicación
$M_{k+1}/fM_{k+1}\to M_k/fM_k$ es inyectiva: ahora bien, si
$x\in M_{k+1}$ es tal que $x\in fM_k$, se deduce de (ii) que
$x\in fM_{k+1}$, lo que termina de probar el lema.
\end{proof}

\begin{proposition}[15.1.19]
\label{0.15.1.19-es}
Sean $A$ un anillo, $M$ un $A$-módulo provisto de una filtración
$(M_k)_{k\geq0}$ tal que $M_0=M$, y $(f_i)_{1\leq i\leq n}$ una sucesión
$\gr_\bullet(M)$-regular de elementos de $A$. Supongamos además que, para
$0\leq i\leq n-1$, la filtración cociente de $(M_k)$ sobre
\[
  M/\left(\sum_{j\leq i}f_jM\right)
\]
es separada. Entonces la sucesión $(f_i)$ es $M$-regular y, si se pone
\[
  M'=M/\left(\sum_{i\leq n}f_iM\right)
\]
y se dota a $M'$ de la filtración cociente de $(M_k)$,
$\gr_\bullet(M')$ es canónicamente isomorfo a
\[
  \gr_\bullet(M)/
  \left(\sum_{i\leq n}f_i\gr_\bullet(M)\right).
\]
\end{proposition}

\begin{proof}
Para $n=1$, la proposición se deduce del lema (15.1.18). Razonemos por
inducción sobre $n$, poniendo
\[
  N=M/\left(\sum_{i\leq n-1}f_iM\right)
\]
y denotando por $(N_k)$ la filtración cociente de $(M_k)$ sobre $N$, que es
separada por hipótesis; además, $\gr_\bullet(N)$ es isomorfo a
\[
  \gr_\bullet(M)/
  \left(\sum_{i\leq n-1}f_i\gr_\bullet(M)\right).
\]
Por hipótesis, $f_n$ es, pues, $\gr_\bullet(N)$-regular, y se deduce del lema
(15.1.18), aplicado a $N$, que $f_n$ es $N$-regular y $\gr_\bullet(M')$ es
isomorfo a $\gr_\bullet(N)/f_n\gr_\bullet(N)$, lo que demuestra la
proposición.
\end{proof}

La proposición (15.1.19) se aplicará, en particular, a las filtraciones que
verifiquen una u otra de las hipótesis siguientes:
\begin{enumerate}
  \item[$1^\circ$] La filtración $(M_k)$ es finita y separada (pues esto
  implica que $M_k=0$ para $k$ suficientemente grande);
  \item[$2^\circ$] $A$ es un anillo noetheriano, $\mathfrak{J}$ un ideal
  contenido en el radical de $A$, $M$ un $A$-módulo de tipo finito y $(M_k)$
  la filtración $\mathfrak{J}$-preádica ($0_{\mathrm I}$, 7.3.5).
\end{enumerate}

\begin{corollary}[15.1.20]
\label{0.15.1.20-es}
\oldpage[0\textsubscript{IV-1}]{20}
Sean $A$ un anillo, $M$ un $A$-módulo y $(f_i)_{1\leq i\leq n}$ una sucesión
$M$-regular de elementos de $A$. Entonces, cualesquiera que sean los enteros
$\nu_i\geq1$, la sucesión $(f_i^{\nu_i})$ es $M$-regular y
$M/(\sum_{i\leq n}f_i^{\nu_i}M)$ posee una serie de composición cuyos
cocientes son isomorfos a $M/(\sum_{i\leq n}f_iM)$.
\end{corollary}

\begin{proof}
Razonemos por inducción sobre $n$; para $n=1$, las afirmaciones son inmediatas,
pues $f^jM/f^{j+1}M$ es isomorfo a $M/fM$, ya que la multiplicación por $f$ en
$M$ es inyectiva por hipótesis. Pongamos entonces
\[
  N=M/\left(\sum_{i\leq n-1}f_i^{\nu_i}M\right)~;
\]
por la hipótesis de inducción, $N$ admite una filtración finita $(N_k)$ para la
cual los $\gr_k(N)$ son isomorfos a
$M/(\sum_{i\leq n-1}f_iM)$; por hipótesis, $f_n$ es, pues,
$\gr_\bullet(N)$-regular; se deduce entonces de (15.1.19) que $f_n$ es
$N$-regular. Aplicando el caso $n=1$ demostrado al principio, se concluye que
$f_n^{\nu_n}$ es $N$-regular y que $N/f_n^{\nu_n}N$ admite una serie de
composición cuyos cocientes son isomorfos a $N/f_nN$; pero, para la filtración
de $N/f_nN$ cociente de la filtración $(N_k)$,
$\gr_k(N/f_nN)$ es isomorfo a $\gr_k(N)/f_n\gr_k(N)$ por (15.1.19) para todo
$k$, luego a $M/(\sum_{i\leq n}f_iM)$, lo que prueba el corolario.
\end{proof}

\begin{proposition}[15.1.21]
\label{0.15.1.21-es}
Sean $A$, $B$ dos anillos locales noetherianos, $\varphi:A\to B$ un
homomorfismo local, $M$ un $B$-módulo de tipo finito y
$(f_i)_{1\leq i\leq n}$ una sucesión $A$-regular formada por elementos del
ideal máximo de $A$. Las propiedades siguientes son equivalentes:
\begin{enumerate}
  \item[a)] $M$ es un $A$-módulo plano.
  \item[b)] Si $\mathfrak{J}$ es el ideal $\sum_{i\leq n}f_iA$,
  $M/\mathfrak{J}M$ es un $A/\mathfrak{J}$-módulo plano y la sucesión $(f_i)$
  es $M$-regular.
\end{enumerate}
\end{proposition}

\begin{proof}
El hecho de que a) implica b) se deduce de (15.1.13) y de
($0_{\mathrm{III}}$, 10.2.2). Recíprocamente, si la sucesión $(f_i)$ es
$M$-regular, es $M$-cuasi-regular (15.1.9), luego el homomorfismo
(15.1.1.1) es biyectivo. Pero, como también se supone que la sucesión
$(f_i)_{1\leq i\leq n}$ es $A$-regular, es $A$-cuasi-regular (15.1.9), y el
homomorfismo canónico (15.1.1.1) correspondiente
$(A/\mathfrak{J})[T_1,\ldots,T_n]\to\gr_\bullet(A)$ es biyectivo. Se concluye
que el homomorfismo canónico ($0_{\mathrm{III}}$, 10.1.1.2)
\[
  \gr_0(M)\otimes_{A/\mathfrak{J}}\gr_\bullet(A)
  \longrightarrow\gr_\bullet(M)
\]
es biyectivo. La conclusión se deduce entonces de ($0_{\mathrm{III}}$,
10.2.2), pues los $f_i$ están en el ideal máximo de $A$ y el homomorfismo
$\varphi$ es local.
\end{proof}

\subsection{Sucesiones \texorpdfstring{$\mathscr{F}$}{F}-regulares}
\label{subsection:0.15.2-es}

\begin{env}[15.2.1]
\label{0.15.2.1-es}
Sean $X$ un espacio anillado, $\mathscr{F}$ un $\mathcal{O}_X$-Módulo,
$f_i$ ($1\leq i\leq n$) una sucesión de secciones de $\mathcal{O}_X$ sobre
$X$, y $\mathscr{J}$ el Ideal de $\mathcal{O}_X$ generado por los $f_i$
(imagen de $\mathcal{O}_X^n$ por el homomorfismo
$\mathcal{O}_X^n\to\mathcal{O}_X$ definido canónicamente por los $f_i$
($0_{\mathrm I}$, 5.1.1)). Los sub-$\mathcal{O}_X$-Módulos
$\mathscr{J}^k\mathscr{F}$ definen una filtración sobre $\mathscr{F}$, y
se pone también
$\gr_k(\mathscr{F})=\mathscr{J}^k\mathscr{F}/\mathscr{J}^{k+1}\mathscr{F}$~;
si $S_k$ es el conjunto de los polinomios homogéneos de grado total $k$ en
el anillo $\mathbb{Z}[T_1,\ldots,T_n]$, $S_k$ puede considerarse como un
haz simple sobre $X$, y se define, como en (15.1.1), un homomorfismo
canónico
\[
\label{0.15.2.1.1-es}
  \varphi_k:\gr_0(\mathscr{F})\otimes_{\mathbb Z}S_k
  \longrightarrow\gr_k(\mathscr{F}).
\tag{15.2.1.1}
\]
\oldpage[0\textsubscript{IV-1}]{21}
de la manera siguiente: para todo abierto $U$ de $X$ (o de una base de la
topología de $X$), se considera el homomorfismo (15.1.1.1)
\[
  \varphi_{k,U}:\gr_0(\Gamma(U,\mathscr{F}))\otimes_{\mathbb Z}S_k
  \longrightarrow\gr_k(\Gamma(U,\mathscr{F}))
\]
definido por las $n$ secciones $f_i|U\in\Gamma(U,\mathcal{O}_X)$, estando el
$\Gamma(U,\mathcal{O}_X)$-módulo $\Gamma(U,\mathscr{F})$ filtrado por los
submódulos $(\Gamma(U,\mathscr{J}))^k\Gamma(U,\mathscr{F})$. Como el
$\mathcal{O}_X$-Módulo $\gr_k(\mathscr{F})$ está asociado al prehaz
$U\mapsto\gr_k(\Gamma(U,\mathscr{F}))$ para todo $k$, los $\varphi_{k,U}$
definen un homomorfismo de $\mathcal{O}_X$-Módulos (15.2.1.1). En virtud
de las propiedades de exactitud de los límites inductivos de módulos y de su
conmutación con los productos tensoriales, de lo anterior resulta que, para
todo $x\in X$, el homomorfismo $(\varphi_k)_x$ de los tallos en el punto $x$
de los dos miembros de (15.2.1.1) se identifica con el homomorfismo de (15.1.1)
\[
\label{0.15.2.1.2-es}
  \gr_0(\mathscr{F}_x)\otimes_{\mathbb Z}S_k
  \longrightarrow\gr_k(\mathscr{F}_x)
\tag{15.2.1.2}
\]
definido por la sucesión $((f_i)_x)_{1\leq i\leq n}$ de elementos de
$\mathcal{O}_x$.
\end{env}

\begin{env}[15.2.2]
\label{0.15.2.2-es}
Con la notación de (15.2.1), diremos que la sucesión
$(f_i)_{1\leq i\leq n}$ es $\mathscr{F}$-regular si, denotando por
$f_j\mathscr{F}$ el sub-$\mathcal{O}_X$-Módulo de $\mathscr{F}$ imagen del
endomorfismo de $\mathscr{F}$ definido por $f_j$, el endomorfismo de
$\mathscr{F}/(\sum_{j=1}^{i-1}f_j\mathscr{F})$ definido por $f_i$ es inyectivo
para $1\leq i\leq n$; diremos que la sucesión $(f_i)$ es
$\mathscr{F}$-cuasirregular si los homomorfismos (15.2.1.1) son inyectivos.
De estas definiciones y de las observaciones de (15.2.1) resulta que, para que
la sucesión $(f_i)$ sea $\mathscr{F}$-regular (resp.
$\mathscr{F}$-cuasirregular), es necesario y suficiente que, para todo
$x\in X$, la sucesión $((f_i)_x)_{1\leq i\leq n}$ sea
$\mathscr{F}_x$-regular (resp. $\mathscr{F}_x$-cuasirregular). Toda sucesión
$\mathscr{F}$-regular es, pues, $\mathscr{F}$-cuasirregular (15.1.9); el
recíproco es cierto si $\mathscr{F}$ es un $\mathcal{O}_X$-Módulo de tipo
finito, si los $\mathcal{O}_x$ son anillos noetherianos para todo $x\in X$ y
si, además, los $(f_i)_x$ pertenecen al radical de $\mathcal{O}_x$ para todo
$x\in X$ (15.1.11).

Si $X=\operatorname{Spec}(A)$ es un esquema afín, y
$\mathscr{F}=\widetilde{M}$, resulta de (I, 1.3) que, para que la sucesión
$(f_i)$ sea $\mathscr{F}$-regular (resp. $\mathscr{F}$-cuasirregular), es
necesario y suficiente que sea $M$-regular (resp. $M$-cuasirregular).

Cuando $\mathscr{F}=\mathcal{O}_X$, se omitirá la mención de $\mathscr{F}$ en
las definiciones precedentes.
\end{env}

\begin{remark}[15.2.3]
\label{0.15.2.3-es}
Cuando $X=\operatorname{Spec}(A)$ y
$\mathfrak{J}=\sum_{i=1}^n f_iA$, la sucesión $(f_j)$ es cuasirregular para
todo $A$-módulo $M$ en todo abierto $U$ que no interseque
$V(\mathfrak{J})$, pues en todo punto $x\in U$ se tiene
$\mathfrak{J}_x=\mathcal{O}_x$ y los dos miembros de (15.2.1.2) son nulos. La
noción de sucesión $\widetilde{M}$-regular solo tiene interés en un entorno de
los puntos de $V(\mathfrak{J})$.
\end{remark}

\begin{proposition}[15.2.4]
\label{0.15.2.4-es}
Sean $X$ un espacio anillado, $\mathscr{F}$ un $\mathcal{O}_X$-Módulo
coherente, $x$ un punto de $X$, y $(f_i)_{1\leq i\leq n}$ una sucesión de
elementos de $\Gamma(X,\mathcal{O}_X)$. Si la sucesión
$((f_i)_x)_{1\leq i\leq n}$ es $\mathscr{F}_x$-regular, existe un entorno
abierto $U$ de $x$ tal que la sucesión $(f_i|U)_{1\leq i\leq n}$ sea
$(\mathscr{F}|U)$-regular.
\end{proposition}

\begin{proof}
En efecto, para $1\leq i\leq n$, el $\mathcal{O}_X$-Módulo
$\mathscr{F}/(\sum_{j=1}^{i-1}f_j\mathscr{F})$ es coherente
($0_{\mathrm I}$, 5.3.4), por lo que también lo es el núcleo del endomorfismo
de este $\mathcal{O}_X$-Módulo definido por $f_i$
($0_{\mathrm I}$, 5.3.4); como el tallo de este núcleo en el punto $x$ es
nulo, lo mismo ocurre con su restricción a un entorno de $x$
($0_{\mathrm I}$, 5.2.2).
\end{proof}

\begin{proposition}[15.2.5]
\label{0.15.2.5-es}
\oldpage[0\textsubscript{IV-1}]{22}
Sea $u=(\psi,\theta):X\to Y$ un morfismo plano de espacios anillados
($0_{\mathrm I}$, 6.7.1); sean $\mathscr{F}$ un
$\mathcal{O}_Y$-Módulo, $\mathscr{G}=u^*(\mathscr{F})$,
$(f_i)_{1\leq i\leq n}$ una sucesión de elementos de
$\Gamma(Y,\mathcal{O}_Y)$, y $(g_i)_{1\leq i\leq n}$ la sucesión de sus
imágenes
$g_i=\theta(f_i)\in\Gamma(Y,u_*(\mathcal{O}_X))=\Gamma(X,\mathcal{O}_X)$. Si
la sucesión $(f_i)$ es $\mathscr{F}$-regular, la sucesión $(g_i)$ es
$\mathscr{G}$-regular. El recíproco es cierto si $u$ es un morfismo fielmente
plano ($0_{\mathrm I}$, 6.7.8).
\end{proposition}

\begin{proof}
Sean $x$ un punto de $X$ e $y=u(x)=\psi(x)$; se tiene
\[
  \mathscr{G}_x=\psi_x^*(\mathscr{F}_y)
  \otimes_{\psi_x^*(\mathcal{O}_y)}\mathcal{O}_x,
  \qquad\text{y}\qquad
  (g_i)_x=\theta_x^\#(\psi_x^*((f_i)_y))~;
\]
como, por hipótesis,
$\theta_x^\#:\psi_x^*(\mathcal{O}_y)\to\mathcal{O}_x$ hace de
$\mathcal{O}_x$ un $\psi_x^*(\mathcal{O}_y)$-módulo plano, y como el funtor
$\psi^*$ es exacto, la primera afirmación resulta de la afirmación
correspondiente de (15.1.14); lo mismo ocurre con la segunda, teniendo en cuenta
que la hipótesis de que $u$ es fielmente plano significa que $\psi$ es
sobreyectivo y que $\theta_x^\#$ hace de $\mathcal{O}_x$ un
$\psi_x^*(\mathcal{O}_y)$-módulo fielmente plano.
\end{proof}

\section{Dimensión y profundidad en los anillos locales noetherianos}
\label{section:0.16-es}

\subsection{Dimensión de un anillo}
\label{subsection:0.16.1-es}

\begin{env}[16.1.1]
\label{0.16.1.1-es}
Se llama dimensión (o dimensión de Krull) de un anillo $A$, y se denota por
$\dim A$, a la dimensión (combinatoria) de su espectro
$\operatorname{Spec}(A)$ (14.2.1); como los subconjuntos cerrados irreducibles
de $\operatorname{Spec}(A)$ son los
$V(\mathfrak{p})$, donde $\mathfrak{p}$ es un ideal primo de $A$ (Bourbaki,
\emph{Alg.~comm.}, cap.~II, \S~4, n.º~3, cor.~1 de la prop.~14), equivale a
decir que $\dim(A)$ es el supremo de las longitudes de las cadenas (14.1.1)
\[
  \mathfrak{p}_0\subset\mathfrak{p}_1\subset\cdots\subset\mathfrak{p}_n
\]
de ideales primos de $A$. Si $A\neq0$, se tiene $\dim A\geq0$. Un anillo
artiniano $\neq0$ (en particular, un cuerpo) tiene dimensión 0 (Bourbaki,
\emph{Alg.~comm.}, cap.~IV, \S~2, n.º~5, prop.~7); un anillo de Dedekind que
no es un cuerpo (en particular, $\mathbb{Z}$ o un anillo $k[T]$ de polinomios
en una indeterminada sobre un cuerpo $k$) tiene dimensión 1 (Bourbaki,
\emph{Alg.~comm.}, cap.~VII, \S~2). Un anillo noetheriano no tiene
necesariamente dimensión finita [30].

Si $(\mathfrak{p}_\alpha)$ es la familia de los ideales primos minimales de
$A$, se tiene (14.1.2.1)
\[
\label{0.16.1.1.1-es}
  \dim(A)=\sup_\alpha(\dim(A/\mathfrak{p}_\alpha)).
\tag{16.1.1.1}
\]
\end{env}

\begin{env}[16.1.2]
\label{0.16.1.2-es}
Para todo ideal $\mathfrak{a}$ de $A$, los ideales primos de
$A/\mathfrak{a}$ se corresponden biunívocamente con los ideales primos de $A$
que contienen a $\mathfrak{a}$, por lo que
\[
\label{0.16.1.2.1-es}
  \dim(A/\mathfrak{a})\leq\dim(A).
\tag{16.1.2.1}
\]
Si, además, $\mathfrak{a}$ no está contenido en ningún ideal primo minimal
$\mathfrak{p}_\alpha$ de $A$, toda cadena de ideales primos de $A$ que
contenga a $\mathfrak{a}$ puede prolongarse mediante uno de los
$\mathfrak{p}_\alpha$, por lo que se tiene
\[
\label{0.16.1.2.2-es}
  1+\dim(A/\mathfrak{a})\leq\dim(A).
\tag{16.1.2.2}
\]
\end{env}

\begin{env}[16.1.3]
\label{0.16.1.3-es}
\oldpage[0\textsubscript{IV-1}]{23}
Si $S$ es un subconjunto multiplicativo de $A$, los ideales primos de
$S^{-1}A$ se corresponden biunívocamente con los ideales primos de $A$
disjuntos de $S$, por lo que se tiene
\[
\label{0.16.1.3.1-es}
  \dim(S^{-1}A)\leq\dim(A).
\tag{16.1.3.1}
\]
Para todo ideal primo $\mathfrak{p}$ de $A$, se tiene, por definición (14.2.1),
\[
\label{0.16.1.3.2-es}
  \dim(A_{\mathfrak{p}})
  =\operatorname{codim}(V(\mathfrak{p}),\operatorname{Spec}(A)).
\tag{16.1.3.2}
\]
Este número también se llama altura del ideal primo $\mathfrak{p}$ y se denota
por $\operatorname{ht}(\mathfrak{p})$. Más generalmente, para todo ideal
$\mathfrak{a}$ de $A$, se llama altura de $\mathfrak{a}$ al número
\[
\label{0.16.1.3.3-es}
  \operatorname{ht}(\mathfrak{a})
  =\operatorname{codim}(V(\mathfrak{a}),\operatorname{Spec}(A)).
\tag{16.1.3.3}
\]
Si $(\mathfrak{p}_\lambda)$ es la familia de los ideales primos minimales entre
los que contienen a $\mathfrak{a}$, se tiene, pues, (14.2.1)
\[
\label{0.16.1.3.4-es}
  \operatorname{ht}(\mathfrak{a})
  =\inf_\lambda(\operatorname{ht}(\mathfrak{p}_\lambda)).
\tag{16.1.3.4}
\]
Evidentemente, se tiene
\[
\label{0.16.1.3.5-es}
  \dim(A)=\sup_{\mathfrak{m}}(\dim(A_{\mathfrak{m}}))
\tag{16.1.3.5}
\]
donde $\mathfrak{m}$ recorre el conjunto de los ideales máximos de $A$.
\end{env}

\begin{env}[16.1.4]
\label{0.16.1.4-es}
Para todo ideal primo $\mathfrak{p}$ de $A$, se tiene (14.2.2.2)
\[
\label{0.16.1.4.1-es}
  \dim(A_{\mathfrak{p}})+\dim(A/\mathfrak{p})\leq\dim A.
\tag{16.1.4.1}
\]
Se dice que $A$ es catenario (resp. equidimensional, equicodimensional,
biequidimensional) cuando $\operatorname{Spec}(A)$ es catenario (resp.
equidimensional,
equicodimensional, biequidimensional). Cuando $A$ es noetheriano y
biequidimensional, los dos miembros de (16.1.4.1) son iguales para todo ideal
primo de $A$ (14.3.5).

Para todo ideal $\mathfrak{a}$ de $A$, los ideales primos de
$A/\mathfrak{a}$ se corresponden biunívocamente con los ideales primos de $A$
que contienen a $\mathfrak{a}$; por tanto, si $A$ es catenario, también lo es
$A/\mathfrak{a}$. Del mismo modo, para todo subconjunto multiplicativo $S$ de
$A$, los ideales primos de $S^{-1}A$ se corresponden biunívocamente con los
ideales primos de $A$ disjuntos de $S$; por tanto, si $A$ es catenario, también
lo es $S^{-1}A$.

Además, como los ideales primos de $A$ contenidos en un ideal primo
$\mathfrak{p}$ se corresponden biunívocamente con los ideales primos de
$A_{\mathfrak{p}}$, para que $A$ sea catenario es necesario y suficiente que
$A_{\mathfrak{p}}$ lo sea para todo $\mathfrak{p}$. Decir que $A$ es catenario
significa, pues, ((14.3.2) y (16.1.3.2)), que para todo par de ideales primos
$\mathfrak{p}$, $\mathfrak{q}$ de $A$ tales que
$\mathfrak{q}\subset\mathfrak{p}$, se tiene
\[
\label{0.16.1.4.2-es}
  \dim(A_{\mathfrak{q}})
  +\dim(A_{\mathfrak{p}}/\mathfrak{q}A_{\mathfrak{p}})
  =\dim(A_{\mathfrak{p}}).
\tag{16.1.4.2}
\]
Si $A$ es biequidimensional, también lo es $A_{\mathfrak{p}}$ para todo ideal
primo $\mathfrak{p}$ de $A$.
\end{env}

\begin{proposition}[16.1.5]
\label{0.16.1.5-es}
Sea $\rho:A\to B$ un homomorfismo de anillos que hace de $B$ una $A$-álgebra
entera. Entonces se tiene $\dim(B)\leq\dim(A)$; si, además, $\rho$ es
inyectivo, se tiene $\dim(B)=\dim(A)$.
\end{proposition}

\begin{proof}
\oldpage[0\textsubscript{IV-1}]{24}
Si $\mathfrak{n}$ es el núcleo de $\rho$, $\rho(A)$ es isomorfo a
$A/\mathfrak{n}$, luego $\dim\rho(A)\leq\dim(A)$ por (16.1.2.1), y como $B$
es entero sobre $\rho(A)$, basta demostrar la segunda afirmación cuando
$A\subset B$. Si
$\mathfrak{p}'_0\subset\mathfrak{p}'_1\subset\cdots\subset\mathfrak{p}'_n$
es una cadena de ideales primos de $B$, los $\mathfrak{p}'_i\cap A$ son
entonces ideales primos distintos en $A$ (Bourbaki, \emph{Alg.~comm.}, cap.~V,
\S~2, n.º~1, cor.~1 de la prop.~1), por lo que forman una cadena de ideales
primos. Recíprocamente, para toda cadena
$\mathfrak{p}_0\subset\mathfrak{p}_1\subset\cdots\subset\mathfrak{p}_m$
de ideales primos de $A$, se sabe que existe, para cada $i$, un ideal primo
$\mathfrak{p}'_i$ de $B$ tal que $\mathfrak{p}'_i\cap A=\mathfrak{p}_i$ y que
$\mathfrak{p}'_i\subset\mathfrak{p}'_{i+1}$ para $0\leq i\leq m-1$
(\emph{loc.~cit.}, cor.~2 del teo.~1). De aquí se sigue la conclusión.
\end{proof}

\begin{proposition}[16.1.6]
\label{0.16.1.6-es}
Sean $B$ un dominio entero, $A$ un subanillo local de $B$ y $\mathfrak{m}$ su
ideal máximo. Supongamos que $A$ es integralmente cerrado y que $B$ es entero
sobre $A$; entonces, para todo ideal máximo $\mathfrak{n}$ de $B$, se tiene
$\dim(B_{\mathfrak{n}})=\dim(A)$.
\end{proposition}

\begin{proof}
La primera parte del razonamiento de (16.1.5) muestra que
$\dim(B_{\mathfrak{n}})\leq\dim(A)$. Recíprocamente, consideremos una cadena
de ideales primos
$\mathfrak{p}_0\subset\mathfrak{p}_1\subset\cdots\subset
\mathfrak{p}_k=\mathfrak{m}$ de $A$; se sabe que
$\mathfrak{n}\cap A=\mathfrak{m}$ (Bourbaki, \emph{Alg.~comm.}, cap.~V,
\S~2, n.º~1, prop.~1) y, en virtud de las hipótesis formuladas, puede aplicarse
el segundo teorema de Cohen-Seidenberg (Bourbaki, \emph{Alg.~comm.}, cap.~V,
\S~2, n.º~4, teo.~3), que prueba la existencia de una cadena de ideales primos
$\mathfrak{p}'_0\subset\mathfrak{p}'_1\subset\cdots\subset
\mathfrak{p}'_k=\mathfrak{n}$ de $B$ tal que
$\mathfrak{p}'_i\cap A=\mathfrak{p}_i$ para todo $i$; se deduce inmediatamente
$\dim(B_{\mathfrak{n}})\geq\dim(A)$, de donde se sigue la proposición.
\end{proof}

\begin{env}[16.1.7]
\label{0.16.1.7-es}
Sea $M$ un $A$-módulo; se llama dimensión de $M$, y se denota por
$\dim_A(M)$ o $\dim(M)$, a la dimensión del anillo $A/\operatorname{Ann}(M)$
(isomorfo al anillo de homotecias $A_M$ de $M$). Se tiene, pues,
$\dim(M^k)=\dim(M)$ para todo $k>0$. Si $M$ es de tipo finito, los ideales
primos de $A$ que contienen a $\operatorname{Ann}(M)$ son exactamente los que
pertenecen a $\operatorname{Supp}(M)$ ($0_{\mathrm I}$, 1.7.4), y este último
es cerrado en $\operatorname{Spec}(A)$; por tanto, se tiene entonces
\[
\label{0.16.1.7.1-es}
  \dim(M)=\dim(\operatorname{Supp}(M))\leq\dim(A).
\tag{16.1.7.1}
\]
Si $M$, $N$ son dos $A$-módulos tales que $N\subset M$, se tiene
\[
\label{0.16.1.7.2-es}
  \dim(N)\leq\dim(M),\qquad \dim(M/N)\leq\dim(M).
\tag{16.1.7.2}
\]
Si $S$ es un subconjunto multiplicativo de $A$ y $M$ es de tipo finito, se
tiene
\[
\label{0.16.1.7.3-es}
  \dim_{S^{-1}A}(S^{-1}M)\leq\dim_A(M).
\tag{16.1.7.3}
\]
En efecto, por (Bourbaki, \emph{Alg.~comm.}, cap.~II, \S~2, n.º~4), se tiene
$\operatorname{Ann}(S^{-1}M)=S^{-1}\operatorname{Ann}(M)$; además,
$S^{-1}A/S^{-1}\operatorname{Ann}(M)=S^{-1}(A/\operatorname{Ann}(M))$.

Por otra parte, toda cadena maximal de ideales primos que contengan a
$\operatorname{Ann}(M)$ termina en un ideal máximo $\mathfrak{m}$, por lo que
su longitud es igual a la de la cadena de ideales primos correspondientes de
$A_{\mathfrak{m}}$ (que contienen a
$\operatorname{Ann}(M_{\mathfrak{m}})=(\operatorname{Ann}(M))_{\mathfrak{m}}$);
de esta observación y de (16.1.7.3) se obtiene
\[
\label{0.16.1.7.4-es}
  \dim_A(M)=\sup_{\mathfrak{m}}
  (\dim_{A_{\mathfrak{m}}}(M_{\mathfrak{m}}))
\tag{16.1.7.4}
\]
donde $\mathfrak{m}$ recorre el conjunto de los ideales máximos de $A$.

Se dice que $M$ es catenario (resp. equidimensional, equicodimensional,
biequidimensional) cuando $A/\operatorname{Ann}(M)$ lo es.
\end{env}

\begin{proposition}[16.1.8]
\label{0.16.1.8-es}
\oldpage[0\textsubscript{IV-1}]{25}
Sean $A$ un anillo, $M$ un $A$-módulo de tipo finito y $\mathfrak{p}$ un ideal
primo de $A$. Se tiene entonces
\[
\label{0.16.1.8.1-es}
  \dim_{A_{\mathfrak{p}}}(M_{\mathfrak{p}})
  +\dim_A(M/\mathfrak{p}M)\leq\dim_A(M).
\tag{16.1.8.1}
\]
\end{proposition}

\begin{proof}
Se tiene
$\operatorname{Ann}(M_{\mathfrak{p}})=(\operatorname{Ann}(M))_{\mathfrak{p}}$,
por lo que los ideales primos de $A_{\mathfrak{p}}$ que contienen a
$\operatorname{Ann}(M_{\mathfrak{p}})$ se corresponden biunívocamente con los
ideales primos de $A$ que contienen a $\operatorname{Ann}(M)$ y están
contenidos en $\mathfrak{p}$. Por otra parte, si
$\operatorname{Ann}(M)=\mathfrak{n}$,
$V(\operatorname{Ann}(M/\mathfrak{p}M))=V(\mathfrak{p}+\mathfrak{n})$
($0_{\mathrm I}$, 1.7.5), por lo que los ideales primos de $A$ que contienen a
$\operatorname{Ann}(M/\mathfrak{p}M)$ contienen a
$\mathfrak{p}+\mathfrak{n}$; de aquí se sigue la conclusión.

Esta última observación muestra, además, que se tiene
\[
\label{0.16.1.8.2-es}
  \dim_{A/\mathfrak{p}}(M/\mathfrak{p}M)=\dim_A(M/\mathfrak{p}M).
\tag{16.1.8.2}
\]
\end{proof}

\begin{proposition}[16.1.9]
\label{0.16.1.9-es}
Sean $A$ un anillo noetheriano, $\rho:A\to B$ un homomorfismo de anillos y $M$
un $B$-módulo tal que $M_{[\rho]}$ sea un $A$-módulo de tipo finito. Entonces
se tiene
\[
\label{0.16.1.9.1-es}
  \dim_A(M_{[\rho]})=\dim_B(M).
\tag{16.1.9.1}
\]
\end{proposition}

\begin{proof}
En efecto, el anillo $B_M$ de las homotecias del $B$-módulo $M$ se identifica
con un subanillo de $C=\operatorname{End}_A(M_{[\rho]})$, y $C$ es un
$A$-módulo de tipo finito (Bourbaki, \emph{Alg.~comm.}, cap.~III, \S~3,
n.º~1, lema~2), por lo que lo mismo ocurre con $B_M$; por otra parte, el
anillo $A_M$ de las homotecias del $A$-módulo $M_{[\rho]}$ es un subanillo de
$B_M$, imagen canónica de $A$ en $C$, por lo que $B_M$ es finito sobre $A_M$;
la conclusión se deduce entonces de (16.1.5).
\end{proof}

\begin{env}[16.1.10]
\label{0.16.1.10-es}
Sea $A$ un anillo noetheriano; si $M$ es un $A$-módulo de longitud finita, se
sabe (Bourbaki, \emph{Alg.~comm.}, cap.~IV, \S~2, n.º~5, prop.~7 y cor.~1 de
la prop.~7) que $\operatorname{Supp}(M)$ es un espacio finito discreto, por lo
que $\dim(M)=0$ si $M\neq0$; recíprocamente, si $M$ es un $A$-módulo de tipo
finito tal que $\dim(M)=0$, todo punto de $\operatorname{Supp}(M)$ es cerrado,
es decir, es un ideal máximo de $A$, por lo que (\emph{loc.~cit.}, prop.~7)
$M$ es de longitud finita.
\end{env}
\subsection[Dimensión de un anillo semilocal noetheriano]{Dimensión de un
anillo semilocal noetheriano\protect\footnote{Este número reproduce la exposición
impartida por Serre en un curso en el Collège de France en 1958.}}
\label{subsection:0.16.2-es}

\begin{env}[16.2.1]
\label{0.16.2.1-es}
Sean $A$ un anillo semilocal noetheriano y $\mathfrak{r}$ su radical~;
recordemos que un ideal de definición $\mathfrak{q}$ de $A$ es un ideal tal que
$\mathfrak{q}\subset\mathfrak{r}$ y que $\mathfrak{q}$ contiene una potencia de
$\mathfrak{r}$~; esto equivale a decir que los únicos ideales primos que
contienen a $\mathfrak{q}$ son maximales, o que $A/\mathfrak{q}$ es un anillo
artiniano. Sea $M$ un $A$-módulo de tipo finito~; para todo entero $n>0$,
$M/\mathfrak{q}^nM$ es un $A$-módulo de longitud finita, igual a
$\sum_{i=0}^{n-1}\operatorname{long}(\gr_i(M))$, donde $\gr_\bullet(M)$ es el
$(A/\mathfrak{q})$-módulo graduado asociado a $M$ provisto de la filtración
$\mathfrak{q}$-preádica. Se sabe que existe un polinomio $Q(n)$ con
coeficientes racionales tal que $\operatorname{long}(\gr_n(M))=Q(n)$ para $n$
suficientemente grande (III, 2.5.4)~; de ello se deduce inmediatamente que
existe un polinomio $P_{\mathfrak{q}}(M,n)$ en $n$, con coeficientes racionales,
tal que
\[
  \operatorname{long}(M/\mathfrak{q}^nM)=P_{\mathfrak{q}}(M,n)
\]
para $n$ suficientemente grande~; está claro que este polinomio es único y que
su coeficiente dominante es $>0$ si $M\neq0$~; se dice que es el polinomio de
Hilbert--Samuel de $M$ para la filtración $\mathfrak{q}$-preádica.
\end{env}

\begin{lemma}[16.2.2.1]
\label{0.16.2.2.1-es}
\oldpage[0\textsubscript{IV-1}]{26}
Sea $(M_n)$ una filtración $\mathfrak{q}$-buena de $M$ (Bourbaki,
\emph{Alg.~comm.}, cap.~III, \S~3, núm.~1)~; entonces
$\operatorname{long}(M/M_n)$ es igual, para $n$ suficientemente grande, a un
polinomio en $n$ cuyo grado y coeficiente dominante son los mismos que los de
$P_{\mathfrak{q}}(M,n)$.
\end{lemma}

\begin{proof}
En efecto, existe $n_0$ tal que $M_{n+1}=\mathfrak{q}M_n$ para $n\geq n_0$,
de donde las inclusiones
$\mathfrak{q}^{n+n_0}M\subset M_{n+n_0}=\mathfrak{q}^nM_{n_0}
\subset\mathfrak{q}^nM\subset M_n$ para $n$ suficientemente grande, lo que
implica
\[
  P_{\mathfrak{q}}(M,n+n_0)\geq
  \operatorname{long}(M/M_{n+n_0})\geq
  P_{\mathfrak{q}}(M,n)\geq\operatorname{long}(M/M_n)
\]
y, por consiguiente, la conclusión.
\end{proof}

Si $\mathfrak{q}'$ es un segundo ideal de definición, se tiene
$\mathfrak{q}'\supset\mathfrak{q}^m$ para algún entero $m$, luego
$P_{\mathfrak{q}'}(M,n)\leq P_{\mathfrak{q}}(M,mn)$~; por consiguiente, el
grado de $P_{\mathfrak{q}}(M,n)$ no depende del ideal de definición
$\mathfrak{q}$ considerado~; se denota por $d(M)$.

\begin{lemma}[16.2.2.2]
\label{0.16.2.2.2-es}
Sea $0\to M'\to M\to M''\to0$ una sucesión exacta de $A$-módulos de tipo
finito. Entonces el polinomio
\[
  P_{\mathfrak{q}}(M,n)-P_{\mathfrak{q}}(M',n)
  -P_{\mathfrak{q}}(M'',n)
\]
tiene grado $\leq d(M')-1\leq d(M)-1$.
\end{lemma}

\begin{proof}
En efecto, la filtración de los
$M'_n=M'\cap\mathfrak{q}^nM'$ es $\mathfrak{q}$-buena (Bourbaki,
\emph{Alg.~comm.}, cap.~III, \S~3, núm.~1, prop.~1)~; como
\[
  \operatorname{long}(M/\mathfrak{q}^nM)
  =\operatorname{long}(M''/\mathfrak{q}^nM'')
  +\operatorname{long}(M'/M'_n),
\]
el lema resulta inmediatamente de (16.2.2.1) aplicado a $M'$.
\end{proof}

\begin{theorem}[16.2.3]
\label{0.16.2.3-es}
\textup{(Krull-Chevalley-Samuel).}
Sean $A$ un anillo semilocal noetheriano, $\mathfrak{r}$ su radical,
$M\neq0$ un $A$-módulo de tipo finito y $d(M)$ el grado del polinomio de
Hilbert-Samuel de $M$ (para un ideal de definición cualquiera de $A$); sea,
por otra parte, $s(M)$ el ínfimo de los enteros $n$ tales que existen elementos
$x_1,\ldots,x_n$ de $\mathfrak{r}$ para los cuales
$M/(x_1M+\cdots+x_nM)$ es de longitud finita. Entonces se tiene
\[
  d(M)=s(M)=\dim(M).
\]
\end{theorem}

\begin{lemma}[16.2.3.1]
\label{0.16.2.3.1-es}
Para todo $x\in\mathfrak{r}$, sea ${}_xM$ el submódulo de $M$ formado por los
elementos anulados por $x$. Entonces:
\begin{enumerate}
  \item[\rm{(i)}] Se tiene $s(M)\leq s(M/xM)+1$.
  \item[\rm{(ii)}] Sean $\mathfrak{p}_i$ los ideales primos pertenecientes a
  $\operatorname{Supp}(M)$ y tales que
  \[
    \dim(A/\mathfrak{p}_i)=\dim(M)\quad(1\leq i\leq m).
  \]
\end{enumerate}
Si $x\notin\mathfrak{p}_i$ para todo $i$, se tiene
$\dim(M/xM)\leq\dim(M)-1$.
\begin{enumerate}
  \item[\rm{(iii)}] Si $\mathfrak{q}$ es un ideal de definición de $A$, el
  polinomio
  \[
    P_{\mathfrak{q}}({}_xM,n)-P_{\mathfrak{q}}(M/xM,n)
  \]
  tiene grado $\leq d(M)-1$.
\end{enumerate}
\end{lemma}

\begin{proof}
La afirmación (i) es trivial, pues si $(x_i)_{1\leq i\leq m}$ es tal que, para
el módulo $N=M/xM$, $N/(x_1N+\cdots+x_mN)$ es de longitud finita, basta
observar que este módulo es isomorfo a
$M/(xM+x_1M+\cdots+x_mM)$.

Para probar (ii), se observa que, si $\mathfrak{a}$ es el anulador de $M$, los
ideales primos que contienen al anulador de $M/xM$ son los que contienen a
$\mathfrak{a}+Ax$ ($0_{\mathrm I}$, 1.7.5); ninguno de los
$\mathfrak{p}_i$ puede, por tanto, contener a $\mathfrak{a}+Ax$ y, por
definición de $\dim(M)$ (16.1.7), toda cadena de ideales primos que contenga a
$\mathfrak{a}+Ax$ tiene, pues, longitud $\leq\dim(M)-1$.

\oldpage[0\textsubscript{IV-1}]{27}

Por último, para demostrar (iii), se observa que hay dos sucesiones exactas
\[
  0\to{}_xM\to M\xrightarrow{x}xM\to0,\qquad
  0\to xM\to M\to M/xM\to0,
\]
y la afirmación se deduce inmediatamente del lema (16.2.2.2).
\end{proof}

La demostración de (16.2.3) se realiza en tres etapas:

\begin{env}[16.2.3.2]
\label{0.16.2.3.2-es}
Se tiene $\dim(M)\leq d(M)$.
\end{env}

Razonemos por inducción sobre $d(M)$, siendo trivial el caso $d(M)=0$.
Supongamos, pues, que $d(M)\geq1$, y sea
$\mathfrak{p}_0\in\operatorname{Supp}(M)$ tal que
$\dim(A/\mathfrak{p}_0)=\dim(M)$. Esto implica que $\mathfrak{p}_0$ es un
elemento \emph{minimal} de $\operatorname{Supp}(M)$, luego está asociado a $M$
(Bourbaki, \emph{Alg.~comm.}, cap.~IV, \S~1, n.º~4, teo.~2) y, por tanto, $M$
contiene un submódulo $N$ isomorfo a $A/\mathfrak{p}_0$ (\emph{loc.~cit.},
\S~1, n.º~1); como $d(M)\geq d(N)$ por (16.2.2.2), basta probar que
$\dim(N)\leq d(N)$. Sea, pues,
$\mathfrak{p}_0\subset\mathfrak{p}_1\subset\cdots\subset\mathfrak{p}_n$ una
cadena de ideales primos distintos en $A$, y demostremos que $n\leq d(N)$.
Esto es evidente si $n=0$. En caso contrario, existe
$x\in\mathfrak{p}_1\cap\mathfrak{r}$ tal que $x\notin\mathfrak{p}_0$, pues la
relación $\mathfrak{p}_0\supset\mathfrak{p}_1\cap\mathfrak{r}$ implicaría
$\mathfrak{p}_0\supset\mathfrak{r}$, ya que $\mathfrak{p}_0$ no contiene a
$\mathfrak{p}_1$, luego $\mathfrak{p}_0$ sería máximo, lo que es absurdo.
Se tiene $\mathfrak{p}_i\in\operatorname{Supp}(N/xN)$ para $i\geq1$, luego
$n-1\leq\dim(N/xN)$ (16.1.7). Como ${}_xN=0$ en virtud de la elección de $x$,
se tiene $d(N/xN)\leq d(N)-1$ por (16.2.3.1, (iii)). La hipótesis de inducción
implica entonces $d(N/xN)=\dim(N/xN)\geq n-1$, luego
$n-1\leq d(N)-1$.

\begin{env}[16.2.3.3]
\label{0.16.2.3.3-es}
Se tiene $d(M)\leq s(M)$.
\end{env}

Sea, en efecto, $(x_i)_{1\leq i\leq n}$ una familia de elementos de
$\mathfrak{r}$ tal que, si se pone
$\mathfrak{a}=x_1A+\cdots+x_nA$, $M/\mathfrak{a}M$ es de longitud finita.
Los ideales asociados a $M/\mathfrak{a}M$ son entonces máximos (Bourbaki,
\emph{Alg.~comm.}, cap.~IV, \S~2, n.º~5, prop.~7) y, como son los únicos
ideales primos que contienen a
$\mathfrak{q}=\mathfrak{a}+(\mathfrak{r}\cap\operatorname{Ann}(M))$,
$A/\mathfrak{q}$ es artiniano (\emph{loc.~cit.}, prop.~9) y $\mathfrak{q}$ es
un ideal de definición de $A$. Pero se tiene
$\mathfrak{q}^mM/\mathfrak{q}^{m+1}M=
\mathfrak{a}^mM/\mathfrak{a}^{m+1}M$, y si $M$ está generado por $r$ elementos
$z_j$, $\mathfrak{a}^mM/\mathfrak{a}^{m+1}M$ es un
$(A/\mathfrak{q})$-módulo generado por las imágenes canónicas de los
$x_1^{\nu_1}x_2^{\nu_2}\cdots x_n^{\nu_n}z_j$ donde
$\sum_i\nu_i=m$, luego por $r\binom{m+n-1}{n-1}$ elementos; su longitud es,
por tanto, $\leq ar\binom{m+n-1}{n-1}$, donde $a$ es una constante, y se
deduce inmediatamente que el grado de $P_{\mathfrak{q}}(M,n)$ es $\leq n$;
de donde $d(M)\leq s(M)$.

\begin{env}[16.2.3.4]
\label{0.16.2.3.4-es}
Se tiene $s(M)\leq\dim(M)$.
\end{env}

Observemos primero que $n=\dim(M)$ es finito por (16.2.3.2). Razonemos por
inducción sobre $n$; la proposición es evidente para $n=0$, pues entonces $M$
es de longitud finita (16.1.10). Supongamos $n\geq1$, y sean
$\mathfrak{p}_i$ $(1\leq i\leq m)$ los ideales primos de
$\operatorname{Supp}(M)$ tales que $\dim(A/\mathfrak{p}_i)=n$; la definición
de los $\mathfrak{p}_i$ muestra que son elementos minimales de
$\operatorname{Supp}(M)$, luego son finitos en número (Bourbaki,
\emph{Alg.~comm.}, cap.~IV, \S~1, n.º~4, teo.~2). Como $n\geq1$, los
$\mathfrak{p}_i$ no son máximos, y existe, por tanto,
$x\in\mathfrak{r}$ tal que $x\notin\mathfrak{p}_i$ para todo $i$ (Bourbaki,
\emph{Alg.~comm.}, cap.~II, \S~1, n.º~1, prop.~2). En virtud de (16.2.3.1),
se tiene $s(M)\leq s(M/xM)+1$ y $\dim(M)\geq\dim(M/xM)+1$; la hipótesis de
inducción implica que $s(M/xM)\leq\dim(M/xM)$, de donde se sigue la conclusión.

\begin{corollary}[16.2.4]
\label{0.16.2.4-es}
Bajo las hipótesis de (16.2.3), se tiene
\begin{equation}
\label{0.16.2.4.1-es}
\tag{16.2.4.1}
  \dim_{\widehat{A}}(\widehat{M})=\dim_A(M).
\end{equation}
\end{corollary}

\oldpage[0\textsubscript{IV-1}]{28}

\begin{proof}
En efecto, $M/\mathfrak{r}^nM$ y
$\widehat{M}/\mathfrak{r}^n\widehat{M}$ son isomorfos, luego
$d(\widehat{M})=d(M)$.
\end{proof}

\begin{corollary}[16.2.5]
\label{0.16.2.5-es}
La dimensión de un anillo semilocal noetheriano $A$ es finita e igual al número
mínimo de elementos del radical de $A$ que generan un ideal de definición.
\end{corollary}

\begin{proof}
Es la igualdad $s(M)=\dim(M)$ para $M=A$.
\end{proof}

En particular:

\begin{corollary}[16.2.6]
\label{0.16.2.6-es}
Sean $A$ un anillo local noetheriano, $\mathfrak{m}$ su ideal máximo y
$k=A/\mathfrak{m}$ su cuerpo residual; se tiene
\begin{equation}
\label{0.16.2.6.1-es}
\tag{16.2.6.1}
  \dim(A)\leq\operatorname{rg}_k(\mathfrak{m}/\mathfrak{m}^2).
\end{equation}
\end{corollary}

\begin{proof}
En efecto, se sabe que si los $x_i$ $(1\leq i\leq n)$ son elementos de
$\mathfrak{m}$ cuyas clases mód.~$\mathfrak{m}^2$ forman una base del
$k$-espacio vectorial $\mathfrak{m}/\mathfrak{m}^2$, los $x_i$ generan
$\mathfrak{m}$ (Bourbaki, \emph{Alg.~comm.}, cap.~II, \S~3, n.º~2, cor.~2 de
la prop.~4).
\end{proof}

\begin{proposition}[16.2.7]
\label{0.16.2.7-es}
Sean $k$ un campo, $B$ una $k$-álgebra graduada de tipo finito en grados
positivos, generada por sus elementos homogéneos de grado 1 y tal que $B_0=k$;
sea $E$ un $B$-módulo graduado de tipo finito, de modo que (III, 2.5.4), para
$n$ suficientemente grande, $\operatorname{rg}_k(E_n)=r(n)$ es de la forma
$\Phi(n)-\Phi(n-1)$, donde $\Phi$ es un polinomio en $n$ de grado $d$. Entonces
existe en $B$ una familia $(t_i)_{1\leq i\leq d}$ de $d$ elementos homogéneos
de grados $\geq1$ tales que $E/(\sum_i t_iE)$ tiene rango finito sobre $k$.
\end{proposition}

\begin{proof}
Sea $\mathfrak{q}$ el ideal máximo $\bigoplus_{n\geq1}B_n$ de $B$ y, para todo
$n$, sea $E'_n$ el sub-$B$-módulo $\bigoplus_{h\geq n+1}E_h$ de $E$. Todo
elemento de $B-\mathfrak{q}$ determina una homotecia inyectiva en el
$B$-módulo artiniano $E/E'_n$, y esta homotecia es, por tanto, biyectiva
(Bourbaki, \emph{Alg.}, cap.~VIII, \S~2, n.º~2, lema~3), lo que implica que
$(E/E'_n)_{\mathfrak{q}}$ se identifica con $E/E'_n$; como el cuerpo residual
de $B_{\mathfrak{q}}$ es $k$, se deduce de la hipótesis que
$\operatorname{long}_{B_{\mathfrak{q}}}((E/E'_n)_{\mathfrak{q}})$ es un
polinomio en $n$ de grado $d$. Por otra parte, como $E$ es un $B$-módulo de
tipo finito y $B$ está generado por sus elementos homogéneos de grado 1, la
filtración $(E'_n)$ sobre $E$ es $\mathfrak{q}$-buena (Bourbaki,
\emph{Alg.~comm.}, cap.~III, \S~1, n.º~3, lema~1), luego la filtración
$((E'_n)_{\mathfrak{q}})$ sobre $E_{\mathfrak{q}}$ es
$\mathfrak{q}B_{\mathfrak{q}}$-buena. Puesto que $B_{\mathfrak{q}}$ es un
anillo local noetheriano, se concluye de (16.2.3) y (16.2.2.1) que se tiene
$\dim(E_{\mathfrak{q}})=d$. Razonemos entonces por inducción sobre $d$, siendo
evidente el lema si $d=0$. Supongamos $d\geq1$ y observemos que los ideales
primos minimales $\mathfrak{p}_i$ $(1\leq i\leq r)$ en
$\operatorname{Ass}(E)$ son \emph{graduados} (Bourbaki, \emph{Alg.~comm.},
cap.~IV, \S~3, n.º~1, prop.~1), y que $\mathfrak{q}$ es distinto de la unión
de los $\mathfrak{p}_i$ puesto que $d\geq1$; existe, por tanto, un elemento
\emph{homogéneo} $t_1\in\mathfrak{q}$ que no pertenece a ninguno de los
$\mathfrak{p}_i$ (Bourbaki, \emph{Alg.~comm.}, cap.~III, \S~1, n.º~4,
prop.~8). Como los ideales primos $(\mathfrak{p}_i)_{\mathfrak{q}}$ son los
elementos minimales de $\operatorname{Ass}(E_{\mathfrak{q}})$ (Bourbaki,
\emph{Alg.~comm.}, cap.~IV, \S~1, n.º~2, prop.~5), se tiene
$\dim(E_{\mathfrak{q}}/t_1E_{\mathfrak{q}})=d-1$ (16.3.4). Basta entonces
aplicar la hipótesis de inducción al $B$-módulo graduado $E/t_1E$ para obtener
la conclusión.
\end{proof}
\subsection{Sistemas de parámetros en un anillo local noetheriano}
\label{subsection:0.16.3-es}

\begin{proposition}[16.3.1]
\label{0.16.3.1-es}
Sean $A$ un anillo noetheriano, $\mathfrak{p}$ un ideal primo de $A$ y $n$ un
entero. Las condiciones siguientes son equivalentes~:
\begin{enumerate}
  \item[$a)$] $\operatorname{ht}(\mathfrak{p})=
  \dim(A_{\mathfrak{p}})\leq n$~;
  \item[$b)$] Existen $n$ elementos $x_i$ de $A$ $(1\leq i\leq n)$ tales que
  $\mathfrak{p}$ es minimal en el conjunto de los ideales primos que contienen
  el ideal $\mathfrak{a}$ engendrado por los $x_i$ (dicho de otro modo,
  $V(\mathfrak{p})$ es una componente irreducible de $V(\mathfrak{a})$).
\end{enumerate}
\end{proposition}

\begin{proof}
La condición $b)$ implica que $\mathfrak{a}A_{\mathfrak{p}}$ es un ideal de
definición de $A_{\mathfrak{p}}$, de donde se sigue $a)$ en virtud de (16.2.5).
Recíprocamente, si se verifica $a)$, existe un ideal de definición
$\mathfrak{b}$ de $A_{\mathfrak{p}}$, engendrado por $n$ elementos $x_i/s$, con
$s\in A-\mathfrak{p}$. En virtud de la correspondencia biunívoca entre ideales
primos de $A_{\mathfrak{p}}$ e ideales primos de $A$ contenidos en
$\mathfrak{p}$, el ideal $\mathfrak{a}$ de $A$ engendrado por los $x_i$ verifica
$b)$.
\end{proof}

\oldpage[0\textsubscript{IV-1}]{29}

Para $n=1$, se obtiene el caso particular~:

\begin{corollary}[16.3.2]
\label{0.16.3.2-es}
\textup{(``Hauptidealsatz'').}
Para que un ideal primo de un anillo noetheriano sea de altura $\leq1$, es
necesario y suficiente que sea minimal en el conjunto de los ideales primos que
contienen un ideal principal adecuado.
\end{corollary}

\begin{corollary}[16.3.3]
\label{0.16.3.3-es}
\textup{(Artin-Tate).}
Sea $A$ un anillo íntegro noetheriano. Las condiciones siguientes son
equivalentes~:
\begin{enumerate}
  \item[$a)$] $A$ es un anillo semilocal de dimensión $\leq1$.
  \item[$b)$] Existe $f\neq0$ en $A$ tal que $A_f$ es un cuerpo.
\end{enumerate}
\end{corollary}

\begin{proof}
La condición $a)$ equivale a decir que $A$ solo tiene un número finito de
ideales primos $\mathfrak{m}_i\neq0$, y que todos estos ideales son maximales
($0$ es un ideal primo). Como el producto de los $\mathfrak{m}_i$ no es nulo,
puesto que $A$ es íntegro, un elemento $f\neq0$ de este producto pertenece a
todos los $\mathfrak{m}_i$, por lo que $(0)$ es el único ideal primo del anillo
íntegro $A_f$; dicho de otro modo, $A_f$ es un cuerpo. (Obsérvese que esta
parte de la demostración no utiliza el hecho de que $A$ sea noetheriano.)
Demostremos ahora que $b)$ implica $a)$~; la hipótesis $b)$ significa que todo
ideal primo $\mathfrak{p}\neq0$ de $A$ contiene a $f$. Sean $\mathfrak{p}_i$
$(1\leq i\leq n)$ los ideales primos minimales entre los que contienen a $f$
(que son un número finito puesto que $A/fA$ es noetheriano)~; basta demostrar
que los $\mathfrak{p}_i$ son ideales \emph{maximales}, pues entonces todo ideal
primo $\neq0$ contiene necesariamente a uno de los $\mathfrak{p}_i$ y, por
tanto, es igual a él. Supongamos, pues, que existe un ideal maximal
$\mathfrak{m}$ que contiene a uno de los $\mathfrak{p}_i$ y es distinto de
$\mathfrak{p}_i$~; por tanto, $\mathfrak{m}$ no está contenido en la unión de
los $\mathfrak{p}_i$ (Bourbaki, \emph{Alg.~comm.}, cap.~II, \S~1, núm.~1,
prop.~2); dicho de otro modo, existe $g\in\mathfrak{m}$ que no pertenece a
ninguno de los $\mathfrak{p}_i$. Sea $\mathfrak{q}$ uno de los ideales primos
minimales entre los que contienen a $g$~; del ``Hauptidealsatz'' (16.3.2)
resulta que $\mathfrak{q}$ es de altura 1, y, por tanto, $\neq0$~; por
consiguiente, contiene a $f$, luego también a uno de los $\mathfrak{p}_i$, y,
como $g\in\mathfrak{q}$ y $g$ no pertenece a ninguno de los
$\mathfrak{p}_i$, $\mathfrak{q}$ es distinto de los $\mathfrak{p}_i$~; tendría
entonces altura $\geq2$, lo que es contradictorio.
\end{proof}

\begin{proposition}[16.3.4]
\label{0.16.3.4-es}
Sean $A$ un anillo semilocal noetheriano, $\mathfrak{r}$ su radical, $M$ un
$A$-módulo de tipo finito y $\mathfrak{p}_i$ los elementos de
$\operatorname{Supp}(M)$ tales que
$\dim(A/\mathfrak{p}_i)=\dim(M)$. Entonces los $\mathfrak{p}_i$ son elementos
minimales de $\operatorname{Ass}(M)$, y se tiene
$\dim(M/\mathfrak{p}_iM)=\dim(M)$. Para todo $x\in\mathfrak{r}$, se tiene
\begin{equation}
\label{0.16.3.4.1-es}
\tag{16.3.4.1}
  \dim(M/xM)\geq\dim(M)-1.
\end{equation}
Además, para que los dos miembros de (16.3.4.1) sean iguales, es necesario y
suficiente que $x$ no pertenezca a ninguno de los $\mathfrak{p}_i$.
\end{proposition}

\begin{proof}
La primera afirmación es inmediata, pues los $\mathfrak{p}_i$ son, por
definición, elementos minimales de $\operatorname{Supp}(M)$ (16.1.6), y estos
últimos son también elementos minimales de $\operatorname{Ass}(M)$ (Bourbaki,
\emph{Alg.~comm.}, cap.~IV, \S~1, núm.~4, teorema~2). Además,
$\operatorname{Supp}(M/\mathfrak{p}_iM)$ es el conjunto de los ideales primos
que contienen a $\mathfrak{p}_i$ ($0_{\mathrm I}$, 1.7.5), por lo que
$\dim(M/\mathfrak{p}_iM)=\dim(A/\mathfrak{p}_i)=\dim(M)$.

Pongamos $N=M/xM$ y elijamos $x_1,\ldots,x_n$ en $\mathfrak{r}$ tales
que $N/(x_1N+\cdots+x_nN)$ tenga longitud finita y $n=s(N)=\dim(N)$.
Entonces $M/(xM+x_1M+\cdots+x_nM)$ tiene longitud finita, por lo que $s(M)\leq n+1$.
Por (16.2.3), $\dim(M)\leq n+1=\dim(N)+1$, de donde se sigue (16.3.4.1).

El hecho de que los dos miembros de (16.3.4.1) sean iguales cuando $x$ no
pertenece a ninguno de los $\mathfrak{p}_i$ resulta de (16.3.4.1) y de
(16.2.3.1). Recíprocamente, si
\oldpage[0\textsubscript{IV-1}]{30}
$\dim(M/xM)=\dim(M)-1$, ninguno de los $\mathfrak{p}_i$ puede pertenecer a
$\operatorname{Supp}(M/xM)$, y los ideales primos de este soporte son los que
contienen a $\operatorname{Ann}(M)+Ax$~; como los $\mathfrak{p}_i$ contienen a
$\operatorname{Ann}(M)$, no pueden contener a $x$.
\end{proof}

\begin{corollary}[16.3.5]
\label{0.16.3.5-es}
Con las notaciones de (16.3.4), para todo ideal
$\mathfrak{a}\subset\mathfrak{p}_i$, se tiene
$\dim(M/\mathfrak{a}M)=\dim(M)$ y, si se pone
$N=M/\mathfrak{a}M$, la relación $\dim(M/xM)=\dim(M)-1$ implica
$\dim(N/xN)=\dim(N)-1$.
\end{corollary}

\begin{proof}
En efecto, $\operatorname{Supp}(N)$ es el conjunto de los ideales primos que
contienen a $\mathfrak{a}+\operatorname{Ann}(M)$, por lo que $\mathfrak{p}_i$
pertenece a $\operatorname{Supp}(N)$ y, como
$\dim(N)\leq\dim(M)=\dim(A/\mathfrak{p}_i)$, se tiene
$\dim(A/\mathfrak{p}_i)=\dim(N)$~; además, los ideales primos
$\mathfrak{q}\in\operatorname{Supp}(N)$ tales que
$\dim(A/\mathfrak{q})=\dim(N)$ son algunos de los $\mathfrak{p}_j$~; si $x$
no pertenece a ninguno de los $\mathfrak{p}_j$, se tiene, por tanto,
$\dim(N/xN)=\dim(N)-1$ en virtud de (16.3.4).
\end{proof}

\begin{definition}[16.3.6]
\label{0.16.3.6-es}
Sean $A$ un anillo semilocal noetheriano, $\mathfrak{r}$ su radical y $M$ un
$A$-módulo de tipo finito, y pongamos $n=\dim(M)$. Se llama \emph{sistema de
parámetros para $M$} a todo sistema $(x_i)_{1\leq i\leq n}$ de $n$ elementos de
$\mathfrak{r}$ tal que $M/(x_1M+\cdots+x_nM)$ sea de longitud finita.
\end{definition}

Se ha visto en (16.2.3) que siempre existen tales sistemas.

Se ha visto durante la demostración de (16.2.3.3) que, si
$\mathfrak{J}=\sum_{i=1}^n x_iA$ es tal que $M/\mathfrak{J}M$ es de
longitud finita, $\mathfrak{q}=\mathfrak{J}+\operatorname{Ann}(M)$ es un
\emph{ideal de definición} de $A$, y, recíprocamente, si así ocurre,
$M/\mathfrak{J}M$ es un módulo de tipo finito sobre el anillo artiniano
$A/\mathfrak{q}$, y es, por tanto, de longitud finita. Así pues, es equivalente decir
que $(x_i)_{1\leq i\leq n}$ es un sistema de parámetros \emph{para $M$} o
\emph{para $A/\operatorname{Ann}(M)$} (o también que sus imágenes en
$A/\operatorname{Ann}(M)$ forman un sistema de parámetros para este anillo).
Si $(x_i)_{1\leq i\leq n}$ es un sistema de parámetros para $M$, lo mismo vale
para $(x_i^k)$ para todo entero $k\geq1$, pues, en virtud de lo precedente,
podemos limitarnos al caso en que $M=A$, y, si $\mathfrak{J}$ es un ideal de
definición de $A$, el ideal $\sum_{i=1}^n x_i^kA$ contiene a
$\mathfrak{J}^{kn}$, y, por tanto, es también un ideal de definición.

\begin{proposition}[16.3.7]
\label{0.16.3.7-es}
Sean $A$ un anillo semilocal noetheriano, $\mathfrak{r}$ su radical,
$x_1,\ldots,x_k$ elementos de $\mathfrak{r}$ y $M$ un $A$-módulo de tipo
finito. Se tiene entonces
\begin{equation}
\label{0.16.3.7.1-es}
\tag{16.3.7.1}
  \dim(M/(x_1M+\cdots+x_kM))\geq\dim(M)-k.
\end{equation}
Además, para que los dos miembros de (16.3.7.1) sean iguales, es necesario y
suficiente que existan elementos $x_i$
$(k+1\leq i\leq n=\dim(M))$ de $\mathfrak{r}$ tales que
$(x_i)_{1\leq i\leq n}$ sea un sistema de parámetros para $M$.
\end{proposition}

\begin{proof}
La desigualdad (16.3.7.1) resulta de (16.3.4.1) por inducción sobre $k$. Si los
dos miembros de (16.3.7.1) son iguales, y si $x_{k+1},\ldots,x_n$ es un
sistema de parámetros para $M/(x_1M+\cdots+x_kM)$, es claro que
$M/(x_1M+\cdots+x_kM+x_{k+1}M+\cdots+x_nM)$ es de longitud finita, por lo que
$(x_i)_{1\leq i\leq n}$ es un sistema de parámetros para $M$~;
recíprocamente, si existen $x_i$ $(k+1\leq i\leq n)$ con esta propiedad y se
pone $N=M/(x_1M+\cdots+x_kM)$, es claro que
$N/(x_{k+1}N+\cdots+x_nN)$ es de longitud finita, por lo que
$\dim(N)\leq n-k$, y los dos miembros son iguales en virtud de (16.3.7.1).
\end{proof}

\oldpage[0\textsubscript{IV-1}]{31}

\begin{proposition}[16.3.8]
\label{0.16.3.8-es}
Sean $A$ un anillo semilocal noetheriano y $X=\operatorname{Spec}(A)$ su
espectro; sea $\mathfrak{a}$ un ideal de $A$ distinto de $A$ tal que
$\operatorname{codim}(V(\mathfrak{a}),X)=r>0$. Existen entonces $r$ elementos
$x_1,\ldots,x_r$ de $\mathfrak{a}$, que forman parte de un sistema de parámetros
de $A$, tales que
\[
  \operatorname{codim}(V(\mathfrak{a}),
  V(Ax_1+\cdots+Ax_r))=0.
\]
\end{proposition}

\begin{proof}
Sean $\mathfrak{p}_i$ $(1\leq i\leq m)$ los ideales primos minimales del
anillo $A$, y $\mathfrak{q}_j$ $(1\leq j\leq n)$ los ideales primos minimales
entre los que contienen a $\mathfrak{a}$~; se tiene (14.2.1)
\[
  \operatorname{codim}(V(\mathfrak{a}),X)
  =\inf_j(\operatorname{codim}(V(\mathfrak{q}_j),X))
  =\inf_j(\dim(A_{\mathfrak{q}_j}))
\]
(16.1.3.2). La hipótesis $r>0$ implica que $\mathfrak{a}$ no está contenido en
ninguno de los $\mathfrak{p}_i$~; por consiguiente, existe
$x\in\mathfrak{a}$ que no pertenece a ninguno de los $\mathfrak{p}_i$
(Bourbaki, \emph{Alg.~comm.}, cap.~II, \S~1, núm.~1, prop.~2). Se tiene, como
arriba,
\[
  \operatorname{codim}(V(\mathfrak{a}),V(Ax))
  =\inf_j(\dim((A/Ax)_{\mathfrak{q}_j/Ax}))
  =\inf_j(\dim(A_{\mathfrak{q}_j}/A_{\mathfrak{q}_j}x))~;
\]
pero, como $x/1$ no pertenece a ninguno de los ideales primos minimales
de $A_{\mathfrak{q}_j}$, se tiene
$\dim(A_{\mathfrak{q}_j}/A_{\mathfrak{q}_j}x)
=\dim(A_{\mathfrak{q}_j})-1$ (16.3.4), de donde se sigue
$\operatorname{codim}(V(\mathfrak{a}),V(Ax))=r-1$. Por la misma razón, se tiene
$\dim(A/Ax)=\dim(A)-1$. Razonemos ahora por inducción sobre $r$~; si
$r>1$, existen $r-1$ elementos $x'_i$ $(2\leq i\leq r)$ de $A'=A/Ax$,
que forman parte de un sistema de parámetros de este anillo, pertenecen a
$\mathfrak{a}/Ax$ y son tales que
\[
  \operatorname{codim}(V(\mathfrak{a}/Ax),
  V(A'x'_2+\cdots+A'x'_r))=0.
\]
Sea $(x'_i)_{r+1\leq i\leq s}$ un sistema de elementos de $A'$ tal que
$(x'_i)_{2\leq i\leq s}$ sea un sistema de parámetros de $A'$~; para todo
$i$ tal que $2\leq i\leq s$, sea $x_i$ un elemento de la clase $x'_i$ en
$A$, con $x_i\in\mathfrak{a}$ para $2\leq i\leq r$; es inmediato que
$A/(Ax+Ax_2+\cdots+Ax_s)$ es de longitud finita, y como $\dim(A)=s$, se ve
que $x_1=x$ y los $x_i$ de índices $i\geq2$ cumplen las condiciones del
enunciado.
\end{proof}

\begin{proposition}[16.3.9]
\label{0.16.3.9-es}
Sean $A$ y $B$ dos anillos locales noetherianos, $\mathfrak{m}$ el ideal
maximal de $A$, $k$ su cuerpo residual y $\varphi:A\to B$ un homomorfismo
local.
\begin{enumerate}
  \item[\rm{(i)}] Se tiene
  \begin{equation}
  \label{0.16.3.9.1-es}
  \tag{16.3.9.1}
    \dim(B)\leq\dim(A)+\dim(B\otimes_A k).
  \end{equation}
  \item[\rm{(ii)}] Sean $M\neq0$ un $A$-módulo de tipo finito y $N\neq0$ un
  $B$-módulo de tipo finito. Se tiene
  \begin{equation}
  \label{0.16.3.9.2-es}
  \tag{16.3.9.2}
    \dim_B(M\otimes_A N)
    \leq\dim_A(M)+\dim_{B\otimes_A k}(N\otimes_A k).
  \end{equation}
\end{enumerate}
\end{proposition}

\begin{proof}
Sean $m$ la dimensión de $A$ y $(s_i)_{1\leq i\leq m}$ un sistema de
parámetros de $A$ (16.3.6), de modo que, si
$\mathfrak{a}=\sum_i s_iA$, $A/\mathfrak{a}$ es un $A$-módulo de longitud
finita. Entonces el anillo $A/\mathfrak{a}$ es artiniano, por lo que el ideal
maximal $\mathfrak{m}/\mathfrak{a}$ de este anillo es nilpotente; el ideal
$\mathfrak{m}B/\mathfrak{a}B$ de $B/\mathfrak{a}B$, imagen de
$(\mathfrak{m}/\mathfrak{a})\otimes_A B$, es, pues, también nilpotente y, por
consiguiente, se tiene
$\dim(B/\mathfrak{a}B)=\dim(B/\mathfrak{m}B)$ (pues los espectros de estos dos
anillos tienen el mismo espacio subyacente). Por otra parte,
$B/\mathfrak{a}B=B/(s_1B+\cdots+s_mB)$, y las imágenes de los $s_i$ en $B$
pertenecen al ideal maximal de $B$. Por consiguiente (16.3.7), se tiene
\[
  \dim(B)\leq m+\dim(B\otimes_A k),
\]
que no es sino (16.3.9.1).

Para demostrar (16.3.9.2), obsérvese que, si $\mathfrak{r}$
(resp.~$\mathfrak{s}$) es el anulador de $M$ (resp.~$N$)
\oldpage[0\textsubscript{IV-1}]{32}
en $A$ (resp.~$B$), se tiene
$\dim_A(M)=\dim(A/\mathfrak{r})$, $\dim_B(N)=\dim(B/\mathfrak{s})$ por
definición (16.1.7); por otra parte, $\operatorname{Supp}(M\otimes_A N)$ es un
subconjunto cerrado de $\operatorname{Spec}(B)$ que se identifica (I, 9.1.13)
con
\[
  \operatorname{Supp}(M)\times_{\operatorname{Spec}(A)}\operatorname{Supp}(N)
  =\operatorname{Spec}((B/\mathfrak{s})\otimes_A(A/\mathfrak{r}))
  =\operatorname{Spec}(B/(\mathfrak{s}+\mathfrak{r}B)).
\]
Como $N\otimes_A k$ es un $B$-módulo de tipo finito, se tiene
$\dim_{B\otimes_A k}(N\otimes_A k)=\dim_B(N\otimes_A k)$ (16.1.9), y
$\operatorname{Supp}(N\otimes_A k)$ es igual a
$\operatorname{Spec}(B/(\mathfrak{s}+\mathfrak{m}B))$ por el mismo
razonamiento que antes. Aplicando (16.3.9.1) a $A/\mathfrak{r}$ y a
$B/(\mathfrak{s}+\mathfrak{r}B)$, se obtiene
\[
  \dim(B/(\mathfrak{s}+\mathfrak{r}B))
  \leq\dim(A/\mathfrak{r})+
  \dim(B/(\mathfrak{s}+\mathfrak{m}B)),
\]
que no es sino (16.3.9.2).
\end{proof}

\begin{corollary}[16.3.10]
\label{0.16.3.10-es}
Bajo las hipótesis de (16.3.9), si $\mathfrak{m}B$ es un ideal de definición
de $B$, se tiene $\dim(B)\leq\dim(A)$; si, además, $A$ es íntegro y
$\dim(A)=\dim(B)$, $\varphi$ es inyectivo.
\end{corollary}

\begin{proof}
La primera afirmación resulta de (16.3.9), puesto que entonces
$\dim(B/\mathfrak{m}B)=0$. Por lo demás, se puede sustituir $A$ por
$\varphi(A)$, luego $\dim(B)\leq\dim(\varphi(A))$; si
$\mathfrak{a}=\operatorname{Ker}(\varphi)\neq0$, y si $A$ es íntegro, se tiene
$\dim(\varphi(A))=\dim(A/\mathfrak{a})<\dim(A)$ (16.1.2.2), por lo que no se
puede tener entonces $\dim(A)=\dim(B)$ salvo si $\mathfrak{a}=0$.
\end{proof}

\subsection[Profundidad y coprofundidad]{Profundidad y
coprofundidad\protect\footnote{En la tercera parte del cap.~III, desarrollamos
la noción de profundidad desde el punto de vista cohomológico y en un marco
más general.}}
\label{subsection:0.16.4-es}

\begin{proposition}[16.4.1]
\label{0.16.4.1-es}
Sean $A$ un anillo local noetheriano, $\mathfrak{m}$ su ideal maximal y $M$ un
$A$-módulo de tipo finito. Entonces toda sucesión $M$-regular
$(x_i)_{1\leq i\leq r}$ (15.1.7) formada por elementos de $\mathfrak{m}$ forma
parte de un sistema de parámetros para $M$.
\end{proposition}

\begin{proof}
Razonemos por recurrencia sobre $r$; consideremos el $A$-módulo
\[
  N=M/(x_1M+\cdots+x_{r-1}M)~;
\]
por hipótesis, $x_1,\ldots,x_{r-1}$ forman parte de un sistema de parámetros
para $M$, luego $\dim(N)=\dim(M)-(r-1)$ (16.3.7). Como, por hipótesis, la
homotecia de razón $x_r$ en $N$ es inyectiva, $x_r$ no pertenece a ninguno de
los ideales primos asociados a $N$, luego (16.3.4) se tiene
$\dim(N/x_rN)=\dim(N)-1$, lo que también se escribe
\[
  \dim(M/(x_1M+\cdots+x_rM))=\dim(M)-r~;
\]
la conclusión resulta, por tanto, de (16.3.7).
\end{proof}

Para un $A$-módulo $M\neq0$, una sucesión $M$-regular de elementos de
$\mathfrak{m}$ tiene, por tanto, \emph{a lo sumo} $\dim(M)$ elementos.

\begin{lemma}[16.4.2]
\label{0.16.4.2-es}
Sean $A$ un anillo local noetheriano, $\mathfrak{m}$ su ideal maximal,
$k=A/\mathfrak{m}$ su cuerpo residual y $M$ un $A$-módulo. Para que una
sucesión $M$-regular $(x_i)_{1\leq i\leq n}$ de elementos de $\mathfrak{m}$
sea maximal, es necesario y suficiente que se tenga
\[
  \operatorname{Hom}_A(k,M/(x_1M+\cdots+x_nM))\neq0.
\]
\end{lemma}

\begin{proof}
Decir que $(x_i)$ es maximal significa que, para ningún $x\in\mathfrak{m}$,
la homotecia de razón $x$ en $N=M/(x_1M+\cdots+x_nM)$ es inyectiva, o también
que $\mathfrak{m}$ está contenido en la unión de los ideales primos asociados
a $N$ (Bourbaki, \emph{Alg.~comm.}, cap.~IV, \S~1, n.º~1, cor.~2 de la
prop.~2), luego es igual a uno de estos últimos, puesto que es un ideal
\oldpage[0\textsubscript{IV-1}]{33}
maximal (Bourbaki, \emph{Alg.~comm.}, cap.~II, \S~1, n.º~1, prop.~2); dicho de
otro modo, hay un elemento $z\in N$ cuyo anulador es $\mathfrak{m}$ y, por
consiguiente, el submódulo $Az$ de $N$ es isomorfo a
$A/\mathfrak{m}=k$; de donde resulta el lema.
\end{proof}

\begin{lemma}[16.4.3]
\label{0.16.4.3-es}
Sean $A$ un anillo local noetheriano, $\mathfrak{m}$ su ideal maximal,
$k=A/\mathfrak{m}$ su cuerpo residual, $M$ un $A$-módulo y
$(x_i)_{1\leq i\leq n}$ una sucesión $M$-regular de elementos de
$\mathfrak{m}$. Entonces los $A$-módulos
\[
  \operatorname{Hom}_A(k,M/(x_1M+\cdots+x_nM))
  \quad\text{y}\quad \operatorname{Ext}^n_A(k,M)
\]
son isomorfos y, para $n\geq1$, son también isomorfos a
$\operatorname{Ext}^{n-1}_A(k,M/x_1M)$.
\end{lemma}

\begin{proof}
Razonemos por recurrencia sobre $n$, siendo la proposición evidente para
$n=0$. Pongamos $N=M/x_1M$; la sucesión $(x_i)_{2\leq i\leq n}$ es
$N$-regular, luego
\[
  \operatorname{Hom}_A(k,M/(x_1M+\cdots+x_nM))
  =\operatorname{Hom}_A(k,N/(x_2N+\cdots+x_nN))
\]
es isomorfo a $\operatorname{Ext}^{n-1}_A(k,N)$ en virtud de la hipótesis de
recurrencia. Consideremos entonces la sucesión exacta
$0\to M\xrightarrow{x_1}M\to M/x_1M=N\to0$; de ella se deduce la sucesión
exacta de Ext
\begin{equation}
\label{0.16.4.3.1-es}
\tag{16.4.3.1}
  \operatorname{Ext}^{n-1}_A(k,M)
  \to\operatorname{Ext}^{n-1}_A(k,N)
  \to\operatorname{Ext}^n_A(k,M)
  \xrightarrow{x_1}\operatorname{Ext}^n_A(k,M).
\end{equation}
En virtud de la hipótesis de recurrencia, $\operatorname{Ext}^{n-1}_A(k,M)$
es isomorfo a
\[
  \operatorname{Hom}_A(k,M/(x_1M+\cdots+x_{n-1}M)),
\]
que es nulo, puesto que $(x_i)_{1\leq i\leq n-1}$ no es una sucesión
$M$-regular maximal (16.4.2). Por otra parte, como
$x_1\in\mathfrak{m}$, la homotecia de razón $x_1$ en el $A$-módulo
$\operatorname{Hom}_A(k,T)$ es nula para todo $A$-módulo $T$, luego lo mismo
ocurre con la homotecia de razón $x_1$ en todo $A$-módulo
$\operatorname{Ext}^n_A(k,T)$; por tanto, las afirmaciones del lema se deducen
inmediatamente de la sucesión exacta (16.4.3.1).
\end{proof}

\begin{corollary}[16.4.4]
\label{0.16.4.4-es}
Sean $A$ un anillo local noetheriano, $\mathfrak{m}$ su ideal maximal,
$k=A/\mathfrak{m}$ su cuerpo residual y $M$ un $A$-módulo de tipo finito
$\neq0$. Todas las sucesiones $M$-regulares maximales de elementos de
$\mathfrak{m}$ tienen el mismo número de elementos, que es el menor entero $n$
tal que $\operatorname{Ext}^n_A(k,M)\neq0$.
\end{corollary}

\begin{definition}[16.4.5]
\label{0.16.4.5-es}
Sean $A$ un anillo local noetheriano, $\mathfrak{m}$ su ideal maximal y $M$ un
$A$-módulo de tipo finito. Se llama \emph{profundidad} de $M$, y se denota por
$\operatorname{prof}_A M$ o $\operatorname{prof}(M)$, al supremo del número de
elementos de una sucesión $M$-regular formada por elementos de
$\mathfrak{m}$. Se tiene, por tanto, $\operatorname{prof}(M)=+\infty$ si y
solo si $M=0$, y, por (16.4.1),
\begin{equation}
\label{0.16.4.5.1-es}
\tag{16.4.5.1}
  \operatorname{prof}(M)\leq\dim(M)\quad\text{si}\quad M\neq0.
\end{equation}
Evidentemente, se tiene $\operatorname{prof}(M^k)=\operatorname{prof}(M)$ para
todo $k>0$.
\end{definition}

\begin{proposition}[16.4.6]
\label{0.16.4.6-es}
Sean $A$ un anillo local noetheriano, $\mathfrak{m}$ su ideal maximal y $M$ un
$A$-módulo de tipo finito.
\begin{enumerate}
  \item[\rm{(i)}] Para que $\operatorname{prof}(M)=0$, es necesario y
  suficiente que $\mathfrak{m}\in\operatorname{Ass}(M)$ (o también que exista
  un submódulo de $M$ isomorfo al cuerpo residual $k=A/\mathfrak{m}$ de $A$).
  \item[\rm{(ii)}] Para todo elemento $M$-regular $x$ perteneciente a
  $\mathfrak{m}$, se tiene
  \begin{equation}
  \label{0.16.4.6.1-es}
  \tag{16.4.6.1}
    \operatorname{prof}(M/xM)=\operatorname{prof}(M)-1.
  \end{equation}
  \item[\rm{(iii)}] Si $M\neq0$, se tiene
  \begin{equation}
  \label{0.16.4.6.2-es}
  \tag{16.4.6.2}
    \operatorname{prof}(M)
    \leq\inf_{\mathfrak{p}\in\operatorname{Ass}(M)}\dim(A/\mathfrak{p})
    \leq\dim(M).
  \end{equation}
  \oldpage[0\textsubscript{IV-1}]{34}
  \item[\rm{(iv)}] Se tiene
  $\operatorname{prof}_A(M)=
  \operatorname{prof}_{\widehat{A}}(\widehat{M})$.
\end{enumerate}
\end{proposition}

\begin{proof}
(i) Decir que $\operatorname{prof}(M)=0$ significa que no existe ningún
$x\in\mathfrak{m}$ tal que la homotecia de razón $x$ en $M$ sea inyectiva; en
otras palabras, que todo $x\in\mathfrak{m}$ pertenece a un ideal primo
asociado a $M$ (Bourbaki, \emph{Alg.~comm.}, cap.~IV, \S~1, n.º~1, cor.~2 de
la prop.~2); pero $\mathfrak{m}$ solo puede ser unión de los ideales primos
asociados a $M$ si es igual a uno de ellos (Bourbaki,
\emph{Alg.~comm.}, cap.~II, \S~1, n.º~1, prop.~2).

(ii) Si $(x_i)_{1\leq i\leq m}$ es una sucesión $(M/xM)$-regular formada por
elementos de $\mathfrak{m}$, la sucesión formada por $x$ y los $x_i$ es
$M$-regular, y se deduce de la definición que es maximal cuando la sucesión
$(x_i)_{1\leq i\leq m}$ es una sucesión $(M/xM)$-regular maximal, de donde se
sigue la conclusión.

(iii) Razonemos por inducción sobre $n=\operatorname{prof}(M)$, siendo evidente
la afirmación para $n=0$. Utilizaremos el lema siguiente:
\end{proof}

\begin{lemma}[16.4.6.3]
\label{0.16.4.6.3-es}
Sean $A$ un anillo noetheriano, $M$ un $A$-módulo de tipo finito y $t\in A$ un
elemento $M$-regular. Entonces, para todo
$\mathfrak{p}\in\operatorname{Ass}(M)$, todo ideal primo minimal entre los que
contienen a $\mathfrak{p}+At$ está asociado a $M/tM$.
\end{lemma}

\begin{proof}
Se sabe que existe una sucesión exacta
\[
  0\to M'\to M\to M''\to0
\]
tal que $\operatorname{Ass}(M')=\{\mathfrak{p}\}$,
$\operatorname{Ass}(M'')=\operatorname{Ass}(M)-\{\mathfrak{p}\}$ (Bourbaki,
\emph{Alg.~comm.}, cap.~IV, \S~1, n.º~1, prop.~4); como $t$ es
$M$-regular, también es $M'$-regular y $M''$-regular
(\emph{loc.~cit.}, \S~1, n.º~1, cor.~2 de la prop.~2). Se deduce que la
sucesión
\[
  0\to M'/tM'\to M/tM\to M''/tM''\to0
\]
es exacta (15.1.18). Por consiguiente, se tiene
$\operatorname{Ass}(M'/tM')\subset\operatorname{Ass}(M/tM)$. Pero los ideales
primos de $A$ que contienen a $\mathfrak{p}$ son los puntos de
$\operatorname{Supp}(M')$ (\emph{loc.~cit.}, \S~1, n.º~3, prop.~7), luego los
ideales primos que contienen a $\mathfrak{p}+At$ son los puntos del soporte de
$M'/tM'=M'\otimes_A(A/tA)$ ($0_{\mathrm I}$, 1.7.5); un elemento minimal del
conjunto de estos ideales primos pertenece, pues, a
$\operatorname{Ass}(M'/tM')$ (Bourbaki, \emph{loc.~cit.}, \S~1, n.º~3,
cor.~1 de la prop.~7), y \emph{a fortiori} a
$\operatorname{Ass}(M/tM)$.
\end{proof}

\begin{proof}[Continuación de la demostración de (16.4.6)]
Establecido este lema, volvamos a la demostración de (iii). Sea
$x\in\mathfrak{m}$ un elemento $M$-regular, de modo que
$\operatorname{prof}(M/xM)=n-1$ por (ii); para todo
$\mathfrak{p}\in\operatorname{Ass}(M)$, se deduce del lema que existe un ideal
$\mathfrak{p}'$ asociado a $M/xM$ y que contiene a $\mathfrak{p}+xA$; pero,
como $x$ es $M$-regular, se tiene $x\notin\mathfrak{p}$, de donde
$\mathfrak{p}'\neq\mathfrak{p}$ y, por consiguiente,
$\dim(A/\mathfrak{p}')\leq\dim(A/\mathfrak{p})-1$. Ahora bien, la hipótesis de
inducción muestra que $n-1\leq\dim(A/\mathfrak{p}')$; se concluye que
$n\leq\dim(A/\mathfrak{p})$ para todo
$\mathfrak{p}\in\operatorname{Ass}(M)$.

(iv) Basta considerar el caso en que $M\neq0$. Recordemos que
$\widehat{M}=M\otimes_A\widehat{A}$ salvo isomorfismo canónico
($0_{\mathrm I}$, 7.3.3); como $\widehat{A}$ es un $A$-módulo plano, se tienen
isomorfismos canónicos
\[
  \operatorname{Ext}^i_A(k,M)\otimes_A\widehat{A}
  \xrightarrow{\sim}
  \operatorname{Ext}^i_{\widehat{A}}
  (k,M\otimes_A\widehat{A})
\]
puesto que $k=k\otimes_A\widehat{A}$ salvo isomorfismo canónico
($0_{\mathrm{III}}$, 12.3.5). La afirmación (iv) se deduce, pues, de (16.4.4) y
del hecho de que $\widehat{A}$ es un $A$-módulo fielmente plano
($0_{\mathrm I}$, 7.3.5).
\end{proof}

\begin{env}[16.4.6.4]
\label{0.16.4.6.4-es}
Se observará que, si $\mathfrak{p}$ es un ideal primo de $A$, puede ocurrir
tanto que
$\operatorname{prof}_{A_{\mathfrak{p}}}(M_{\mathfrak{p}})
<\operatorname{prof}_A(M)$ como que
$\operatorname{prof}_{A_{\mathfrak{p}}}(M_{\mathfrak{p}})
>\operatorname{prof}_A(M)$; para un ejemplo del primer caso,
\oldpage[0\textsubscript{IV-1}]{35}
basta tomar como $A$ un dominio entero local de dimensión $\geq1$; se tiene
$\operatorname{prof}(A)\geq1$, pero $\operatorname{prof}(K)=0$ para el cuerpo
de fracciones $K$ de $A$. Para un ejemplo del segundo caso, consideremos un
dominio entero local noetheriano $A_0$ de dimensión $\geq2$, con ideal máximo
$\mathfrak{m}_0$, y sea $k=A_0/\mathfrak{m}_0$ su cuerpo residual; sea
$A=A_0\oplus\mathfrak{J}$ la extensión trivial de $A_0$ por el $A_0$-módulo
$k$ (18.2.3), siendo el ideal $\mathfrak{J}$ $A_0$-isomorfo a $k$ (de modo que
la multiplicación en $A$ viene dada por
$(x,z)(x',z')=(xx',xz'+x'z)$). Es claro que $\mathfrak{J}$ es el nilradical de
$A$ y que $A/\mathfrak{J}$ es isomorfo a $A_0$, de modo que todo ideal primo
$\mathfrak{p}$ de $A$ es de la forma
$\mathfrak{p}_0\oplus\mathfrak{J}$, donde $\mathfrak{p}_0$ es un ideal primo
de $A_0$; en particular, $\mathfrak{m}=\mathfrak{m}_0\oplus\mathfrak{J}$ es el
único ideal máximo de $A$. Como todo elemento de $\mathfrak{m}_0$ anula a
$\mathfrak{J}$, se tiene
$\operatorname{prof}(A)=\operatorname{prof}(A_{\mathfrak{m}})=0$; por el
contrario, si $\mathfrak{p}_0\neq\mathfrak{m}_0$, los elementos de
$\mathfrak{J}$ son anulados por los de
$\mathfrak{m}_0-\mathfrak{p}_0$, luego $A_{\mathfrak{p}}$ es canónicamente
isomorfo a $(A_0)_{\mathfrak{p}_0}$ y, como $A_0$ es un dominio entero,
$\operatorname{prof}(A_{\mathfrak{p}})
=\operatorname{prof}((A_0)_{\mathfrak{p}_0})\geq1$ si
$\mathfrak{p}_0\neq0$.
\end{env}

\begin{corollary}[16.4.7]
\label{0.16.4.7-es}
Para que un anillo local noetheriano reducido $A$ tenga profundidad $0$, es
necesario y suficiente que $A$ sea un cuerpo.
\end{corollary}

\begin{proof}
En efecto, como la intersección de los ideales primos minimales de $A$ es $0$,
$A$ no tiene ideales primos asociados inmersos; en virtud de (16.4.6, (i)), si
$\operatorname{prof}(A)=0$, el ideal máximo $\mathfrak{m}$ de $A$ es también
un ideal primo minimal, luego es el único ideal primo de $A$ y, como $A$ es
reducido, $\mathfrak{m}=0$.
\end{proof}

\begin{proposition}[16.4.8]
\label{0.16.4.8-es}
Sean $A$, $B$ dos anillos locales noetherianos,
$\rho:A\to B$ un homomorfismo local y $M$ un $B$-módulo tal que
$M_{[\rho]}$ sea un $A$-módulo de tipo finito. Entonces se tiene
\begin{equation}
\label{0.16.4.8.1-es}
\tag{16.4.8.1}
  \operatorname{prof}_A(M_{[\rho]})=\operatorname{prof}_B(M).
\end{equation}
\end{proposition}

\begin{proof}
Basta considerar el caso $M\neq0$; supongamos que
$\operatorname{prof}_A(M_{[\rho]})=n$. Si $(x_i)_{1\leq i\leq n}$ es una
sucesión $(M_{[\rho]})$-regular maximal formada por elementos del ideal máximo
de $A$, los $\rho(x_i)$ constituyen una sucesión $M$-regular formada por
elementos del ideal máximo de $B$, y se tiene, por tanto, poniendo
$N=M/(\rho(x_1)M+\cdots+\rho(x_n)M)$,
$\operatorname{prof}_B(N)=\operatorname{prof}_B(M)-n$ (16.4.6, (ii)); del
mismo modo,
$N_{[\rho]}=M_{[\rho]}/(x_1M_{[\rho]}+\cdots+x_nM_{[\rho]})$ y
$\operatorname{prof}_A(N_{[\rho]})=0$; nos reducimos así a demostrar la
proposición cuando $n=0$. Sea $k$ el cuerpo residual de $A$; como $k$ es un
cociente de $A$, el $A$-módulo
$P=\operatorname{Hom}_A(k,M_{[\rho]})$ es un submódulo de
$\operatorname{Hom}_A(A,M_{[\rho]})=M_{[\rho]}$, y la hipótesis implica que
$P\neq0$ (16.4.6, (i)); además, $P$ es un $A$-módulo de tipo finito
(pues $M_{[\rho]}$ es un $A$-módulo de tipo finito), luego un $k$-espacio
vectorial de tipo finito, ya que es anulado por el ideal máximo de $A$, y,
finalmente, un $A$-módulo de longitud finita. Por otra parte, $P$ es también un
submódulo del $B$-módulo $M$ y, como tiene longitud finita en tanto que
$A$-módulo, la tiene \emph{a fortiori} en tanto que $B$-módulo. Como es
$\neq0$, contiene un sub-$B$-módulo simple, es decir, isomorfo al cuerpo
residual de $B$; se concluye entonces de (16.4.6, (i)) que, en efecto, se tiene
$\operatorname{prof}_B(M)=0$.
\end{proof}

\begin{definition}[16.4.9]
\label{0.16.4.9-es}
Sean $A$ un anillo local noetheriano y $M$ un $A$-módulo de tipo finito. Si
$M\neq0$, se llama \emph{coprofundidad} de $M$, y se denota por
$\operatorname{coprof}_A(M)$ o $\operatorname{coprof}(M)$, al entero finito
$\dim_A(M)-\operatorname{prof}_A(M)\geq0$ (16.4.5.1). Si $M=0$, se pone
$\operatorname{coprof}(M)=0$.

Se tiene $\operatorname{coprof}(M^k)=\operatorname{coprof}(M)$ para todo
$k>0$.
\end{definition}

\begin{proposition}[16.4.10]
\label{0.16.4.10-es}
Sean $A$ un anillo local noetheriano y $M$ un $A$-módulo de tipo finito.
\oldpage[0\textsubscript{IV-1}]{36}
\begin{enumerate}
  \item[(i)] Para todo elemento $M$-regular $x$ perteneciente al ideal máximo
  de $A$, se tiene
  $\operatorname{coprof}(M/xM)=\operatorname{coprof}(M)$.
  \item[(ii)] Se tiene
  $\operatorname{coprof}_A(M)=\operatorname{coprof}_{\widehat A}(\widehat M)$.
\end{enumerate}
\end{proposition}

\begin{proof}
En efecto, (i) se deduce de (16.3.4) y (16.4.6, (ii)), y (ii) de (16.2.4) y
(16.4.6, (iv)).
\end{proof}

Demostraremos más adelante (IV, 6.11.5) que, para todo ideal primo
$\mathfrak{p}$ de $A$, se tiene
$\operatorname{coprof}_{A_{\mathfrak{p}}}(M_{\mathfrak{p}})
\leq\operatorname{coprof}_A(M)$.

\begin{proposition}[16.4.11]
\label{0.16.4.11-es}
Bajo las hipótesis de (16.4.8), se tiene
\begin{equation}
\label{0.16.4.11.1-es}
\tag{16.4.11.1}
  \operatorname{coprof}_A(M_{[\rho]})=\operatorname{coprof}_B(M).
\end{equation}
\end{proposition}

\begin{proof}
Esto se deduce de (16.1.9) y (16.4.8).
\end{proof}

\subsection{Módulos de Cohen-Macaulay}
\label{subsection:0.16.5-es}

\begin{definition}[16.5.1]
\label{0.16.5.1-es}
Sea $A$ un anillo local noetheriano. Se dice que un $A$-módulo de tipo finito
$M$ es un \emph{módulo de Cohen-Macaulay} si
$\operatorname{coprof}(M)=0$, es decir, si $M=0$ o si $M\neq0$ y
$\dim(M)=\operatorname{prof}(M)$. Se dice que $A$ es un
\emph{anillo de Cohen-Macaulay} si el $A$-módulo $A$ es un módulo de
Cohen-Macaulay.
\end{definition}

Un $A$-módulo de longitud finita es un módulo de Cohen-Macaulay, pues tiene
dimensión $0$ (16.1.10). Un anillo local reducido $A$ de dimensión $1$ es un
anillo de Cohen-Macaulay, pues entonces su ideal máximo $\mathfrak{m}$ no puede
estar asociado a $A$ (Bourbaki, \emph{Alg.~comm.}, cap.~IV, \S~2, n.º~5,
prop.~10), luego $\operatorname{prof}(A)=1$ en virtud de (16.4.6, (i)). Si $A$
es un \emph{dominio entero local e integralmente cerrado de dimensión
$\geq2$}, se tiene $\operatorname{prof}(A)\geq2$: en efecto, si $x\neq0$ es un
elemento de $\mathfrak{m}$, se sabe (Bourbaki, \emph{Alg.~comm.}, cap.~VII,
\S~1, n.º~4, prop.~8) que los ideales primos asociados a $A/xA$ son de altura
$1$, luego distintos de $\mathfrak{m}$, y su unión es, por tanto, distinta de
$\mathfrak{m}$ (Bourbaki, \emph{Alg.~comm.}, cap.~II, \S~1, n.º~1,
prop.~2); existen, pues, sucesiones $(x,y)$ que son $A$-regulares y están
formadas por elementos de $\mathfrak{m}$, de donde se sigue nuestra afirmación.
Se concluye que un anillo local integralmente cerrado de dimensión $2$ es un
anillo de Cohen-Macaulay.

\begin{proposition}[16.5.2]
\label{0.16.5.2-es}
Sean $A$ un anillo local noetheriano y $M$ un $A$-módulo de tipo finito. Para
que $M$ sea un módulo de Cohen-Macaulay, es necesario y suficiente que
$\widehat M$ sea un $\widehat A$-módulo de Cohen-Macaulay.
\end{proposition}
\begin{proof}
Esto resulta inmediatamente de (16.4.10, (ii)).
\end{proof}

\begin{proposition}[16.5.3]
\label{0.16.5.3-es}
Bajo las hipótesis de (16.4.8), para que $M$ sea un $B$-módulo de
Cohen-Macaulay, es necesario y suficiente que $M_{[\rho]}$ sea un $A$-módulo de
Cohen-Macaulay.
\end{proposition}

\begin{proof}
Esto resulta de (16.4.11).
\end{proof}

\begin{proposition}[16.5.4]
\label{0.16.5.4-es}
Sean $A$ un anillo local noetheriano y $M$ un $A$-módulo de tipo finito $\neq0$.
Si $M$ es un $A$-módulo de Cohen-Macaulay, se tiene
$\dim(A/\mathfrak{p})=\dim(M)$ para todo ideal primo
$\mathfrak{p}\in\operatorname{Ass}(M)$; en particular, ningún ideal primo
asociado a $M$ es inmerso.
\end{proposition}

\begin{proof}
Esto resulta inmediatamente de las desigualdades (16.4.6.2).
\end{proof}

También puede decirse que $M$ (o $\operatorname{Supp}(M)$) es
\emph{equidimensional} (16.1.7).

\oldpage[0\textsubscript{IV-1}]{37}

\begin{proposition}[16.5.5]
\label{0.16.5.5-es}
Sean $A$ un anillo local noetheriano, $\mathfrak{m}$ su ideal maximal, $M$
un $A$-módulo de tipo finito $\neq0$ y $x$ un elemento de $\mathfrak{m}$.
Supongamos que $M$ es un módulo de Cohen-Macaulay; entonces las condiciones
siguientes son equivalentes:
\begin{enumerate}
  \item[a)] $x$ es $M$-regular;
  \item[b)] $\dim(M/xM)=\dim(M)-1$;
  \item[c)] $x$ no pertenece a ningún elemento minimal de
  $\operatorname{Supp}(M)$.
\end{enumerate}
Además, $M/xM$ es entonces un módulo de Cohen-Macaulay.
\end{proposition}

\begin{proof}
Se ha visto (16.5.4) que todos los elementos $\mathfrak{p}$ de
$\operatorname{Ass}(M)$ son minimales en $\operatorname{Supp}(M)$ y que
$\dim(A/\mathfrak{p})=\dim(M)$; la equivalencia de a) y c) resulta, pues, de
Bourbaki, \emph{Alg.~comm.}, cap.~IV, \S~1, núm.~1, cor.~2 de la prop.~2;
la equivalencia de b) y c) resulta de (16.3.4). Finalmente, el hecho de que
$M/xM$ sea un módulo de Cohen-Macaulay resulta de (16.4.10, (i)).
\end{proof}

\begin{corollary}[16.5.6]
\label{0.16.5.6-es}
Bajo las hipótesis de (16.5.5), sea $(x_i)_{1\leq i\leq r}$ una sucesión
de elementos de $\mathfrak{m}$. Las condiciones siguientes son equivalentes:
\begin{enumerate}
  \item[a)] La sucesión $(x_i)$ es $M$-regular.
  \item[b)] Se tiene
  $\dim(M/(x_1M+\cdots+x_rM))=\dim(M)-r$.
  \item[c)] Los $x_i$ forman parte de un sistema de parámetros de $M$.
\end{enumerate}
Además, cuando se satisfacen estas condiciones,
$N=M/(x_1M+\cdots+x_rM)$ es un módulo de Cohen-Macaulay, por lo que, para todo
ideal $\mathfrak{p}\in\operatorname{Ass}(N)$, se tiene
$\dim(A/\mathfrak{p})=\dim(M)-r$.
\end{corollary}

\begin{proof}
Ya se sabe que b) y c) son equivalentes (16.3.6) y que a) implica c)
(16.4.1) sin hipótesis sobre $M$. Para ver que c) implica a) y probar la
última afirmación del enunciado, razonemos por inducción sobre $r$, siendo la
afirmación trivial para $r=0$. Puesto que $x_1$ forma parte de un sistema de
parámetros, $\dim(M/x_1M)=\dim(M)-1$ (16.3.6), luego $x_1$ es $M$-regular y
$M/x_1M$ es un módulo de Cohen-Macaulay por (16.5.5). Como la sucesión
$(x_i)_{2\leq i\leq r}$ forma parte de un sistema de parámetros para
$P=M/x_1M$, la hipótesis de inducción muestra que $(x_i)_{2\leq i\leq r}$ es
una sucesión $P$-regular y que
$P/(x_2P+\cdots+x_rP)=N$ es un módulo de Cohen-Macaulay; además,
$(x_i)_{1\leq i\leq r}$ es $M$-regular.
\end{proof}

\begin{remark}[16.5.7]
\label{0.16.5.7-es}
De (16.5.6) resulta que hay identidad entre los \emph{sistemas de
parámetros para $M$} y las \emph{sucesiones $M$-regulares maximales de elementos de
$\mathfrak{m}$} cuando $M$ es un módulo de Cohen-Macaulay. Recíprocamente, de
(16.4.1) resulta que, si un $A$-módulo de tipo finito no nulo es tal que existe
una sucesión $M$-regular que sea un sistema de parámetros para $M$, entonces
$M$ es un módulo de Cohen-Macaulay. La última propiedad enunciada en
(16.5.6) \emph{también caracteriza} los módulos de Cohen-Macaulay:
\end{remark}

\begin{proposition}[16.5.8]
\label{0.16.5.8-es}
Sean $A$ un anillo local noetheriano y $M$ un $A$-módulo de tipo finito. Se
supone que, para toda sucesión $(x_i)_{1\leq i\leq r}$ de elementos de $A$
que forme parte de un sistema de parámetros para $M$, y todo ideal primo
$\mathfrak{p}$ asociado a
$N=M/(x_1M+\cdots+x_rM)$, se tiene
$\dim(A/\mathfrak{p})=\dim(N)$. Entonces $M$ es un módulo de Cohen-Macaulay.
\end{proposition}

\begin{proof}
Sea $(y_i)_{1\leq i\leq n}$ un sistema de parámetros para $M$; basta
mostrar que la sucesión $(y_i)$ es $M$-regular. Razonemos por inducción sobre
$n=\dim(M)$, siendo trivial la proposición para $n=0$. Por hipótesis
(aplicada con $r=0$), los ideales primos $\mathfrak{p}_j$ asociados a $M$
son todos tales que $\dim(A/\mathfrak{p}_j)=\dim(M)$; como $y_1$ forma parte
de un sistema de parámetros
\oldpage[0\textsubscript{IV-1}]{38}
para $M$, de (16.3.7) y (16.3.4) resulta que $y_1$ no pertenece a ninguno
de los $\mathfrak{p}_j$, y por tanto es $M$-regular. Si se pone $P=M/y_1M$,
$(y_i)_{2\leq i\leq n}$ es un sistema de parámetros para $P$, y es
inmediato que $P$ satisface la hipótesis del enunciado; aplicando la hipótesis
de inducción, se ve, pues, que $(y_i)_{2\leq i\leq n}$ es una sucesión
$P$-regular, lo que demuestra la proposición.
\end{proof}

\begin{proposition}[16.5.9]
\label{0.16.5.9-es}
Sean $A$ un anillo local noetheriano y $M$ un $A$-módulo de tipo finito; se
supone que $M$ es un módulo de Cohen-Macaulay. Sea
$\mathfrak{p}\in\operatorname{Supp}(M)$, y pongamos
$r=\dim(M)-\dim(M/\mathfrak{p}M)$. Entonces existe una sucesión $M$-regular
$(x_i)_{1\leq i\leq r}$ formada por elementos de $\mathfrak{p}$; para toda
sucesión que tenga estas propiedades, se tiene
\[
  \dim(M/(x_1M+\cdots+x_rM))
  =\dim(M/\mathfrak{p}M)=\dim(A/\mathfrak{p}),
\]
y $\mathfrak{p}$ es un elemento minimal de
$\operatorname{Ass}(M/(x_1M+\cdots+x_rM))$.
\end{proposition}

\begin{proof}
Probemos primero la existencia de los $x_i$; no hay nada que demostrar para
$r=0$. Para $r>0$, procedamos por inducción sobre $r$. Como para todo ideal
primo $\mathfrak{q}_i$ asociado a $M$ se tiene
$\dim(A/\mathfrak{q}_i)=\dim(M)$ (16.5.4), la hipótesis
$\dim(M/\mathfrak{p}M)<\dim(M)$ implica que $\mathfrak{p}$ no puede estar
contenido en uno de los $\mathfrak{q}_i$, siendo estos últimos los ideales primos
minimales entre los que contienen $\operatorname{Ann}(M)$ (16.1.2.2 y
16.1.7); se concluye que $\mathfrak{p}$ tampoco está contenido en la
unión de los $\mathfrak{q}_i$ (Bourbaki, \emph{Alg.~comm.}, cap.~II, \S~1,
núm.~1, prop.~2), y por consiguiente existe un $x_1\in\mathfrak{p}$ que es
$M$-regular (Bourbaki, \emph{Alg.~comm.}, cap.~IV, \S~1, núm.~1, cor.~2 de
la prop.~2). Se deduce, pues, de (16.5.5) que $N=M/x_1M$ es un módulo de
Cohen-Macaulay, de dimensión $\dim(M)-1$; como $x_1\in\mathfrak{p}$, se tiene
además $N/\mathfrak{p}N=M/\mathfrak{p}M$, de donde
$\dim(N)-\dim(N/\mathfrak{p}N)=r-1$; puede aplicarse, por tanto, a $N$
la hipótesis de inducción, y si $(x_i)_{2\leq i\leq r}$ es una sucesión
$N$-regular formada por elementos de $\mathfrak{p}$, es claro que
$(x_i)_{1\leq i\leq r}$ es una sucesión $M$-regular formada por elementos de
$\mathfrak{p}$.

Si se pone $P=M/(x_1M+\cdots+x_rM)$, se tiene
$P/\mathfrak{p}P=M/\mathfrak{p}M$ puesto que los $x_i$ están en
$\mathfrak{p}$, y $\dim(P)=\dim(P/\mathfrak{p}P)$ por (16.5.6); pero como
$P$ es un módulo de Cohen-Macaulay, se tiene también
$\dim(P)=\dim(A/\mathfrak{p}')$ para todo
$\mathfrak{p}'\in\operatorname{Ass}(P)$ (16.5.4). Como
$\mathfrak{p}\in\operatorname{Supp}(P)$ ($0_{\mathrm I}$, 1.7.5),
$\mathfrak{p}$ contiene uno de los ideales
$\mathfrak{p}'\in\operatorname{Ass}(P)$ (Bourbaki, \emph{Alg.~comm.},
cap.~IV, \S~1, núm.~4, teorema~2), y como
\[
  \dim(P)=\dim(P/\mathfrak{p}P)
  \leq\dim(A/\mathfrak{p})
  \leq\dim(A/\mathfrak{p}')=\dim(P),
\]
se tiene necesariamente $\mathfrak{p}=\mathfrak{p}'$ (16.1.2.2), lo que
termina la demostración.
\end{proof}

\begin{corollary}[16.5.10]
\label{0.16.5.10-es}
Bajo las mismas hipótesis que en (16.5.9):
\begin{enumerate}
  \item[(i)] $M_{\mathfrak{p}}$ es un $A_{\mathfrak{p}}$-módulo de
  Cohen-Macaulay.
  \item[(ii)] Se tiene
  \begin{equation}
  \label{0.16.5.10.1-es}
  \tag{16.5.10.1}
    \dim(M)=\dim_A(M/\mathfrak{p}M)
    +\dim_{A_{\mathfrak{p}}}(M_{\mathfrak{p}}).
  \end{equation}
\end{enumerate}
\end{corollary}

\begin{proof}
Con la notación de (16.5.9), sean $x_i'$ las imágenes canónicas de los $x_i$
en $A_{\mathfrak{p}}$ $(1\leq i\leq r)$; como los $x_i'$ pertenecen al
ideal maximal de $A_{\mathfrak{p}}$ y forman una sucesión
$M_{\mathfrak{p}}$-regular por planitud (15.1.14), se tiene
\[
  \operatorname{prof}_{A_{\mathfrak{p}}}(M_{\mathfrak{p}})
  \geq r=\dim(M)-\dim(M/\mathfrak{p}M)
  \geq\dim_{A_{\mathfrak{p}}}(M_{\mathfrak{p}})
\]
(en virtud de (16.1.8.1)). Pero, teniendo en cuenta (16.4.5.1), los tres
términos de esta desigualdad son necesariamente iguales, de donde resulta el corolario.
\end{proof}

\oldpage[0\textsubscript{IV-1}]{39}

\begin{corollary}[16.5.11]
\label{0.16.5.11-es}
Sea $A$ un anillo local noetheriano; supongamos que existe un $A$-módulo de
tipo finito $M$ de soporte $\operatorname{Spec}(A)$ que es un módulo de
Cohen-Macaulay. Entonces, para todo ideal primo $\mathfrak{p}$ de $A$, se tiene
\begin{equation}
\label{0.16.5.11.1-es}
\tag{16.5.11.1}
  \dim(A)=\dim(A/\mathfrak{p})+\dim(A_{\mathfrak{p}}).
\end{equation}
\end{corollary}

\begin{proof}
En efecto, $\operatorname{Supp}(M/\mathfrak{p}M)
=\operatorname{Spec}(A/\mathfrak{p})$ ($0_{\mathrm I}$, 1.7.5) y
$\operatorname{Supp}(M_{\mathfrak{p}})
=\operatorname{Spec}(A_{\mathfrak{p}})$; por tanto, la relación (16.5.11.1) es un
caso particular de (16.5.10.1).
\end{proof}

\begin{corollary}[16.5.12]
\label{0.16.5.12-es}
Todo anillo cociente de un anillo local noetheriano $A$ que satisfaga las
condiciones de (16.5.11) (en particular, todo anillo cociente de un anillo local
de Cohen-Macaulay) es \emph{catenario}.
\end{corollary}

\begin{proof}
Basta evidentemente probar que $A$ mismo es catenario, es decir, que,
para dos ideales primos $\mathfrak{p}\subset\mathfrak{q}$ de $A$, se tiene
(16.1.4.2)
\[
  \dim(A_{\mathfrak{q}})
  =\dim(A_{\mathfrak{p}})
  +\dim(A_{\mathfrak{q}}/\mathfrak{p}A_{\mathfrak{q}}).
\]
Ahora bien, en virtud de (16.5.10, (i)), $A_{\mathfrak{q}}$ satisface las mismas
hipótesis que $A$, y la relación precedente no es sino
(16.5.11.1) aplicada a $A_{\mathfrak{q}}$ y al ideal primo
$\mathfrak{p}A_{\mathfrak{q}}$ de $A_{\mathfrak{q}}$.
\end{proof}

\begin{env}[16.5.13]
\label{0.16.5.13-es}
Sean $A$ un anillo noetheriano y $M$ un $A$-módulo de tipo finito. Se dice que
$M$ es un \emph{$A$-módulo de Cohen-Macaulay} si, para todo ideal primo
$\mathfrak{p}$ de $A$, $M_{\mathfrak{p}}$ es un $A_{\mathfrak{p}}$-módulo
de Cohen-Macaulay; en virtud de (16.5.10), esta definición coincide con
(16.5.1) cuando $A$ es \emph{local}. Se dice que $A$ es un
\emph{anillo de Cohen-Macaulay} si todos los $A_{\mathfrak{p}}$ lo son. Es
claro que, si $M$ (resp. $A$) es un $A$-módulo de Cohen-Macaulay (resp. un
anillo de Cohen-Macaulay), $S^{-1}M$ es un $S^{-1}A$-módulo de Cohen-Macaulay
(resp. $S^{-1}A$ es un anillo de Cohen-Macaulay) para todo conjunto
multiplicativo $S$ de $A$.
\end{env}

\section{Anillos regulares}
\label{section:0.17-es}

\subsection{Definición de los anillos regulares}
\label{subsection:0.17.1-es}

\begin{proposition}[17.1.1]
\label{0.17.1.1-es}
Sean $A$ un anillo local noetheriano de dimensión $n$, $\mathfrak{m}$ su
ideal maximal, $k=A/\mathfrak{m}$ su cuerpo residual. Las condiciones siguientes
son equivalentes~:
\begin{enumerate}
  \item[\rm a)] El homomorfismo canónico sobreyectivo de $k$-módulos graduados
  (15.1.1.1)
  \[
    \varphi:\mathbf{S}^*_k(\mathfrak{m}/\mathfrak{m}^2)
    \longrightarrow\operatorname{gr}^*_{\mathfrak{m}}(A)
  \]
  (donde el segundo miembro es el $k$-módulo graduado asociado a $A$ provisto de
  la filtración $\mathfrak{m}$-preádica) es biyectivo.
  \item[\rm b)] Se tiene
  $\operatorname{rg}_k(\mathfrak{m}/\mathfrak{m}^2)=n$.
  \item[\rm c)] El ideal $\mathfrak{m}$ admite un sistema de generadores de $n$
  elementos.
  \item[\rm d)] El ideal $\mathfrak{m}$ admite un sistema de generadores que es
  una sucesión $A$-regular.
\end{enumerate}
\end{proposition}


\begin{proof}
En virtud del lema de Nakayama,
$\operatorname{rg}_k(\mathfrak{m}/\mathfrak{m}^2)$ es el menor número de
elementos de un sistema de generadores de $\mathfrak{m}$, por lo que b) y c) son
equivalentes. Por otra parte, si $(x_i)_{1\leq i\leq r}$ es un sistema de
generadores de $\mathfrak{m}$ cuyas clases $\overline{x}_i$ módulo
$\mathfrak{m}^2$ forman una base del $k$-espacio vectorial
$\mathfrak{m}/\mathfrak{m}^2$, $\mathbf{S}^*(\mathfrak{m}/\mathfrak{m}^2)$
es isomorfo al anillo de
\oldpage[0\textsubscript{IV-1}]{40}
polinomios $k[\mathrm{T}_1,\ldots,\mathrm{T}_r]$~; habida cuenta de que todo
$A$-módulo de tipo finito es separado para la filtración $\mathfrak{m}$-preádica
($0_{\mathrm I}$, 7.3.5), se deduce de (15.1.9) que las condiciones a) y d)
son equivalentes~; además, como toda sucesión $A$-regular de elementos de
$\mathfrak{m}$ tiene a lo sumo $n$ elementos (16.4.1), se ve que d) implica c).
Queda demostrar que c) implica a).

Pongamos, para abreviar,
$\mathrm{S}=\mathbf{S}^*(\mathfrak{m}/\mathfrak{m}^2)
=k[\mathrm{T}_1,\ldots,\mathrm{T}_n]$,
$\mathrm{G}=\operatorname{gr}^*_{\mathfrak{m}}(A)$, y consideremos la sucesión
exacta
$0\longrightarrow\mathfrak{J}\longrightarrow\mathrm{S}
\xrightarrow{\varphi}\mathrm{G}\longrightarrow0$, donde el núcleo
$\mathfrak{J}$ de $\varphi$ es, por tanto, un ideal graduado de $\mathrm{S}$. Para
todo entero $s>0$, se tiene
\begin{equation}
\label{0.17.1.1.1-es}
\tag{17.1.1.1}
  \binom{s+n-1}{n-1}
  =\operatorname{long}(\mathrm{S}_s)
  =\operatorname{long}(\mathrm{G}_s)
  +\operatorname{long}(\mathfrak{J}_s).
\end{equation}

Supongamos $\mathfrak{J}\neq0$, de modo que existe un elemento homogéneo
$u\in\mathfrak{J}$ de grado $h>0$~; $\mathfrak{J}_s$ contiene
$u\mathrm{S}_{s-h}$ para $s\geq h$, y como $u\neq0$ y $\mathrm{S}$ es
íntegro, $u\mathrm{S}_{s-h}$ es isomorfo a $\mathrm{S}_{s-h}$ como
$k$-espacio vectorial~; se tendría, por tanto,
$\operatorname{long}(\mathfrak{J}_s)\geq
\operatorname{long}(\mathrm{S}_{s-h})=\binom{s-h+n-1}{n-1}$, de donde, para
$s\geq h$,
\begin{equation}
\label{0.17.1.1.2-es}
\tag{17.1.1.2}
  \operatorname{long}(\mathrm{G}_s)
  \leq\binom{s+n-1}{n-1}-\binom{s-h+n-1}{n-1}.
\end{equation}
Ahora bien, el segundo miembro de (17.1.1.2) es un polinomio en $s$ de grado
$\leq n-2$~; pero se tiene
$\mathrm{G}_s=\mathfrak{m}^s/\mathfrak{m}^{s+1}$, y para $s$ suficientemente grande,
$\operatorname{long}(\mathrm{G}_s)$ es un polinomio en $s$ de grado exactamente
igual a $n-1$ cuyo coeficiente dominante es $>0$ (16.2.1), lo que contradice
la desigualdad (17.1.1.2). C.Q.F.D.
\end{proof}

\begin{definition}[17.1.2]
\label{0.17.1.2-es}
Se dice que un anillo local noetheriano que verifica las condiciones equivalentes de
(17.1.1) es \emph{regular}.
\end{definition}

\begin{corollary}[17.1.3]
\label{0.17.1.3-es}
Un anillo local regular es íntegro, íntegramente cerrado y un anillo de
Cohen-Macaulay.
\end{corollary}

\begin{proof}
La última aserción se deduce de (17.1.1, d))~; por otra parte, si $A$ es
regular, $\operatorname{gr}^*_{\mathfrak{m}}(A)$ es íntegro y completamente
íntegramente cerrado, por ser isomorfo a $k[\mathrm{T}_1,\ldots,\mathrm{T}_n]$~;
se concluye que $A$ posee también estas dos propiedades (Bourbaki,
\emph{Alg.~comm.}, cap.~III, \S~2, no~3, cor. de la prop.~1 y cap.~V,
\S~1, no~5, prop.~15).
\end{proof}

\begin{examples}[17.1.4]
\label{0.17.1.4-es}
\begin{enumerate}
  \item[\rm (i)] Un anillo local regular de \emph{dimensión} $0$, por ser
  íntegro en virtud de (17.1.3), es necesariamente un \emph{cuerpo}, y recíprocamente.
  \item[\rm (ii)] Para que un anillo local noetheriano de \emph{dimensión} $1$
  sea regular, es necesario y suficiente que sea un \emph{anillo de valuación
  discreta}~: en efecto, decir que un anillo local noetheriano de dimensión $1$ es
  regular significa que su ideal maximal es \emph{principal}, y la conclusión
  se deduce de Bourbaki, \emph{Alg.~comm.}, cap.~VI, \S~3, no~6, prop.~9.
  \item[\rm (iii)] Sea $k$ un cuerpo~; el anillo de series formales
  $A=k[[\mathrm{T}_1,\ldots,\mathrm{T}_n]]$ es un \emph{anillo regular de
  dimensión $n$}~; en efecto, es claro que los $\mathrm{T}_i$
  $(1\leq i\leq n)$ generan el ideal maximal $\mathfrak{m}$ de $A$ y
  forman una sucesión $A$-regular, pues
  $A/(\mathrm{T}_1A+\cdots+\mathrm{T}_iA)$ es isomorfo a
  $k[[\mathrm{T}_{i+1},\ldots,\mathrm{T}_n]]$.
\end{enumerate}
\end{examples}

\begin{proposition}[17.1.5]
\label{0.17.1.5-es}
Para que un anillo local noetheriano $A$ sea regular, es necesario y suficiente que
su completado $\widehat{A}$ lo sea.
\end{proposition}

\begin{proof}
\oldpage[0\textsubscript{IV-1}]{41}
En efecto, el ideal maximal de $\widehat{A}$ es $\mathfrak{m}\widehat{A}$ y se
sabe que $\mathfrak{m}^h/\mathfrak{m}^{h+1}$ y
$(\mathfrak{m}\widehat{A})^h/(\mathfrak{m}\widehat{A})^{h+1}$ son
$k$-espacios vectoriales isomorfos~; la condición a) de (17.1.1) es, por tanto, la
misma para $A$ y $\widehat{A}$.
\end{proof}

\begin{definition}[17.1.6]
\label{0.17.1.6-es}
Dado un anillo local noetheriano $A$, de ideal maximal $\mathfrak{m}$, se
llama \emph{sistema regular de parámetros de $A$} a un sistema de parámetros
para $A$ (16.3.5) que genera $\mathfrak{m}$.
\end{definition}

La existencia de tal sistema equivale a que $A$ sea regular (17.1.1)~;
si $(x_i)_{1\leq i\leq n}$ es un sistema regular de parámetros, es un
sistema minimal de generadores de $\mathfrak{m}$, cuyas clases módulo
$\mathfrak{m}^2$ forman una base de $\mathfrak{m}/\mathfrak{m}^2$ sobre $k$~;
en virtud de (17.1.1, a)) y de (15.1.9), $(x_i)$ es una sucesión $A$-regular
maximal (de donde procede la terminología).

Se advertirá, sin embargo, que, en un anillo regular, un sistema de
parámetros para $A$ que sea también una sucesión $A$-regular \emph{no es
necesariamente} un sistema regular de parámetros, como muestra el ejemplo
de las potencias $k$-ésimas de un sistema regular de parámetros (15.1.20).

\begin{proposition}[17.1.7]
\label{0.17.1.7-es}
Sean $A$ un anillo local noetheriano, $\mathfrak{m}$ su ideal maximal,
$k=A/\mathfrak{m}$ su cuerpo residual, $(x_i)_{1\leq i\leq r}$ una sucesión
de elementos de $\mathfrak{m}$, $\mathfrak{J}$ el ideal
$x_1A+\cdots+x_rA$. Las condiciones siguientes son equivalentes~:
\begin{enumerate}
  \item[\rm a)] $A$ es regular y los $x_i$ forman parte de un sistema
  regular de parámetros de $A$.
  \item[\rm a')] $A$ es regular y las clases $\overline{x}_i$ de los $x_i$
  módulo $\mathfrak{m}^2$ son linealmente independientes sobre $k$.
  \item[\rm b)] Los $x_i$ forman parte de un sistema de parámetros para $A$ y
  $A/\mathfrak{J}$ es un anillo regular.
\end{enumerate}
Además, cuando se verifican estas condiciones, $\mathfrak{J}$ es un ideal
primo.
\end{proposition}

\begin{proof}
La última aserción se deduce trivialmente de que $A/\mathfrak{J}$,
por ser regular, es íntegro (17.1.3). Pongamos $n=\dim(A)$. La equivalencia de a)
y a') es inmediata~; en efecto, las clases módulo $\mathfrak{m}^2$ de los elementos
de un sistema regular de parámetros forman una base de
$\mathfrak{m}/\mathfrak{m}^2$ sobre $k$, y recíprocamente, si los
$\overline{x}_i$ $(1\leq i\leq r)$ son linealmente independientes sobre $k$, se
pueden encontrar $n-r$ elementos $x_i\in\mathfrak{m}$ $(r+1\leq i\leq n)$ tales que
los $\overline{x}_i$ $(1\leq i\leq n)$ formen una base de
$\mathfrak{m}/\mathfrak{m}^2$, por lo que $(x_i)_{1\leq i\leq n}$ es un sistema
regular de parámetros. Para ver que a) implica b), observemos que, puesto que
los $x_i$ forman parte de un sistema de parámetros de $A$, se tiene
$\dim(A/\mathfrak{J})=n-r$ (16.3.6)~; sea
$\mathfrak{n}=\mathfrak{m}/\mathfrak{J}$ el ideal maximal de
$A/\mathfrak{J}$. Se tiene una sucesión exacta
\begin{equation}
\label{0.17.1.7.1-es}
\tag{17.1.7.1}
  0\longrightarrow(\mathfrak{m}^2+\mathfrak{J})/\mathfrak{m}^2
  \longrightarrow\mathfrak{m}/\mathfrak{m}^2
  \longrightarrow\mathfrak{n}/\mathfrak{n}^2\longrightarrow0
\end{equation}
puesto que $\mathfrak{n}^2=(\mathfrak{m}^2+\mathfrak{J})/\mathfrak{J}$~;
la hipótesis a) implica que
$\operatorname{rg}_k((\mathfrak{m}^2+\mathfrak{J})/\mathfrak{m}^2)=r$, por lo que
$\operatorname{rg}_k(\mathfrak{n}/\mathfrak{n}^2)=n-r
=\dim(A/\mathfrak{J})$, lo que demuestra b) en virtud de (17.1.1, b)).

Recíprocamente, para demostrar que b) implica a), observemos que el hecho de que los $x_i$
formen parte de un sistema de parámetros implica que
$\dim(A/\mathfrak{J})=n-r$ (16.3.6)~; la hipótesis de que $A/\mathfrak{J}$ es
regular implica, por otra parte,
$\operatorname{rg}_k(\mathfrak{n}/\mathfrak{n}^2)=n-r$ (17.1.1)~; además,
se tiene evidentemente
$\operatorname{rg}_k((\mathfrak{m}^2+\mathfrak{J})/\mathfrak{m}^2)\leq r$,
por lo que se deduce de (17.1.7.1) que
$\operatorname{rg}_k(\mathfrak{m}/\mathfrak{m}^2)\leq r+(n-r)=n$~; así, $A$
es regular en virtud de (16.2.6) y (17.1.1). Además, la relación
$\operatorname{rg}_k(\mathfrak{m}/\mathfrak{m}^2)=n$ implica entonces
$\operatorname{rg}_k((\mathfrak{m}^2+\mathfrak{J})/\mathfrak{m}^2)=r$, y por
consiguiente los $\overline{x}_i$ son linealmente independientes sobre $k$. C.Q.F.D.
\end{proof}

\oldpage[0\textsubscript{IV-1}]{42}

\begin{corollary}[17.1.8]
\label{0.17.1.8-es}
Sean $A$ un anillo local noetheriano, $\mathfrak{m}$ su ideal maximal, $t$ un
elemento de $\mathfrak{m}$. Las condiciones siguientes son equivalentes~:
\begin{enumerate}
  \item[\rm a)] $A/tA$ es regular y $t$ es no divisor de cero en $A$.
  \item[\rm b)] $A$ es regular y $t\notin\mathfrak{m}^2$.
\end{enumerate}
\end{corollary}

\begin{proof}
Es el caso particular $r=1$ de (17.1.7) (habida cuenta de que un elemento de
$\mathfrak{m}$ no divisor de cero forma parte de un sistema de parámetros
(16.4.1)).
\end{proof}

\begin{corollary}[17.1.9]
\label{0.17.1.9-es}
Sean $A$ un anillo local regular, $\mathfrak{m}$ su ideal maximal,
$\mathfrak{J}$ un ideal de $A$ contenido en $\mathfrak{m}$. Las condiciones
siguientes son equivalentes~:
\begin{enumerate}
  \item[\rm a)] El anillo $A/\mathfrak{J}$ es regular.
  \item[\rm b)] $\mathfrak{J}$ está engendrado por una sucesión
  $(x_i)_{1\leq i\leq r}$ que forma parte de un sistema regular de parámetros
  de $A$.
\end{enumerate}
\end{corollary}

\begin{proof}
Ya se ha visto (17.1.7) que b) implica a). Recíprocamente, supongamos que
$A/\mathfrak{J}$ es regular (lo que implica que $\mathfrak{J}$ es primo)
y sean $n=\dim(A)$, $n-r=\dim(A/\mathfrak{J})$; con la notación de
(17.1.7), la sucesión exacta (17.1.7.1) da
$\operatorname{rg}_k((\mathfrak{m}^2+\mathfrak{J})/\mathfrak{m}^2)=r$, luego
existe una sucesión $(x_i)_{1\leq i\leq r}$ de elementos de $\mathfrak{J}$ que
forma parte de un sistema regular de parámetros de $A$. Pongamos
$\mathfrak{J}'=x_1A+\cdots+x_rA$; de (17.1.7) resulta que $\mathfrak{J}'$
es un ideal primo de $A$ y que $A/\mathfrak{J}'$ es regular y de
dimensión $n-r$; pero como
$A/\mathfrak{J}=(A/\mathfrak{J}')/(\mathfrak{J}/\mathfrak{J}')$ y
$\mathfrak{J}/\mathfrak{J}'$ es primo en el anillo íntegro
$A/\mathfrak{J}'$, las dimensiones de $A/\mathfrak{J}$ y de
$A/\mathfrak{J}'$ solo pueden ser iguales si
$\mathfrak{J}=\mathfrak{J}'$ (16.1.2.2).
\end{proof}

\subsection{Recordatorios sobre la dimensión proyectiva y la dimensión inyectiva de los
módulos}
\label{subsection:0.17.2-es}

\begin{env}[17.2.1]
\label{0.17.2.1-es}
Sean $A$ un anillo y $M$ un $A$-módulo. Recordemos (M, VI, 2) que se denomina
\emph{dimensión proyectiva} (resp. \emph{dimensión inyectiva}) de $M$ y que
se denota $\operatorname{dim.proj}(M)$ o
$\operatorname{dim.proj}_A(M)$ (resp. $\operatorname{dim.inj}(M)$ o
$\operatorname{dim.inj}_A(M)$) al menor $n$ (entero o igual a $+\infty$)
tal que existe una resolución proyectiva por la izquierda (resp. una resolución
inyectiva por la derecha) de $M$ que sea \emph{de longitud $n$}. Es equivalente
(\emph{loc. cit.}) decir que $n$ es el menor número tal que para todo
$A$-módulo $N$, se tiene $\operatorname{Ext}^i_A(M,N)=0$ (resp.
$\operatorname{Ext}^i_A(N,M)=0$) para todo $i>n$, o tan solo para $i=n+1$.
Para todo \emph{factor directo} $M'$ de $M$, se tiene, pues,
$\operatorname{dim.proj}(M')\leq\operatorname{dim.proj}(M)$, ya que
$\operatorname{Ext}^i_A(M',N)$ es un factor directo de
$\operatorname{Ext}^i_A(M,N)$; asimismo,
$\operatorname{dim.inj}(M')\leq\operatorname{dim.inj}(M)$. Una condición
equivalente a $\operatorname{dim.proj}(M)=n$ (resp.
$\operatorname{dim.inj}(M)=n$) es que $n$ sea el menor número tal que,
para toda sucesión exacta
\[
  0\longrightarrow R\longrightarrow P_{n-1}\longrightarrow\cdots
  \longrightarrow P_0\longrightarrow M\longrightarrow0
\]
(resp.
$0\longrightarrow M\longrightarrow Q_0\longrightarrow\cdots
\longrightarrow Q_{n-1}\longrightarrow C\longrightarrow0$), donde todos los
$P_i$ para $0\leq i<n$ son proyectivos (resp. todos los $Q_i$ para $0\leq i<n$
son inyectivos), entonces $R$ es proyectivo (resp. $C$ es inyectivo).
\end{env}

\begin{remarks}[17.2.2]
\label{0.17.2.2-es}
\begin{enumerate}
  \item[\rm (i)] Decir que $M$ es un $A$-módulo \emph{proyectivo} (resp.
  \emph{inyectivo}) equivale, pues, a decir que
  $\operatorname{dim.proj}(M)=0$ (resp. $\operatorname{dim.inj}(M)=0$).
  \item[\rm (ii)] Sea $T$ un \emph{funtor covariante} aditivo de la categoría
  de los $A$-módulos en una categoría abeliana. De las definiciones de
  (17.2.1) y de la definición de los
  \oldpage[0\textsubscript{IV-1}]{43}
  funtores derivados resulta inmediatamente que si $\operatorname{dim.proj}(M)\leq n$ (resp.
  $\operatorname{dim.inj}(M)\leq n$), se tiene $L^iT(M)=0$ (resp.
  $R^iT(M)=0$) \emph{para todo $i>n$}.
  \item[\rm (iii)] Si se supone que $A$ es \emph{noetheriano} y $M$ es \emph{de tipo
  finito}, la última interpretación de la dimensión proyectiva dada en
  (17.2.1) muestra que, si $\operatorname{dim.proj}(M)=n$, $M$ admite una
  resolución por la izquierda de longitud $n$ formada por módulos proyectivos \emph{de tipo
  finito}. Si además $A$ es un anillo \emph{local} noetheriano, $M$ admite una
  resolución de longitud $n$ por módulos \emph{libres de tipo finito}, pues un
  $A$-módulo proyectivo de tipo finito es entonces libre (Bourbaki,
  \emph{Alg.~comm.}, cap.~II, \S~3, núm.~2, cor.~2 de la prop.~5).
\end{enumerate}
\end{remarks}

\begin{lemma}[17.2.3]
\label{0.17.2.3-es}
Sean $A$ un anillo y $M$ un $A$-módulo. Para que
$\operatorname{dim.inj}(M)\leq n$, es necesario y suficiente que para todo
$A$-módulo monógeno $N$, se tenga $\operatorname{Ext}^{n+1}_A(N,M)=0$.
\end{lemma}

\begin{proof}
Con la notación de (17.2.1), se trata de probar que $C$ es inyectivo;
para todo $A$-módulo $N$, $\operatorname{Ext}^{n+1}_A(N,M)$ es isomorfo a
$\operatorname{Ext}^1_A(N,C)$ (M, V, 7), luego se tiene
$\operatorname{Ext}^1_A(N,C)=0$ para todo $A$-módulo $N$ de la forma
$A/\mathfrak{J}$, donde $\mathfrak{J}$ es un ideal cualquiera de $A$. La sucesión
exacta de Ext muestra entonces que el homomorfismo canónico
$\operatorname{Hom}(A,C)\longrightarrow\operatorname{Hom}(\mathfrak{J},C)$ es
sobreyectivo para todo ideal $\mathfrak{J}$ de $A$, lo que implica que $C$ es un
$A$-módulo inyectivo (M, I, 3.2).
\end{proof}

\begin{lemma}[17.2.4]
\label{0.17.2.4-es}
Sean $A$ un anillo noetheriano y $M$ un $A$-módulo de tipo finito. Para que
$\operatorname{dim.proj}(M)\leq n$, es necesario y suficiente que para todo
$A$-módulo monógeno $N$, se tenga $\operatorname{Ext}^{n+1}_A(M,N)=0$.
\end{lemma}

\begin{proof}
En efecto, se sabe (M, VI, 2.5) que la condición
$\operatorname{dim.proj}(M)\leq n$ es equivalente a
$\operatorname{Ext}^{n+1}_A(M,N)=0$ para todo $A$-módulo $N$ \emph{de tipo
finito}; para ver que la condición del enunciado también es suficiente, se
razona por inducción sobre el número de generadores $m$ de $N$: existe un
submódulo $N_1$ de $N$ engendrado por $m-1$ elementos y tal que
$N_2=N/N_1$ es monógeno; de la sucesión exacta
$0\longrightarrow N_1\longrightarrow N\longrightarrow N_2\longrightarrow0$,
se deduce entonces la sucesión exacta
$\operatorname{Ext}^{n+1}_A(M,N_1)\longrightarrow
\operatorname{Ext}^{n+1}_A(M,N)\longrightarrow
\operatorname{Ext}^{n+1}_A(M,N_2)$, y la hipótesis de inducción muestra que la
condición del enunciado implica efectivamente $\operatorname{Ext}^{n+1}_A(M,N)=0$.
\end{proof}

\begin{corollary}[17.2.5]
\label{0.17.2.5-es}
Sean $A$ un anillo noetheriano y $M$ un $A$-módulo. Se tiene
\begin{equation}
\label{0.17.2.5.1-es}
\tag{17.2.5.1}
  \operatorname{dim.inj}_A(M)
  =\sup_{\mathfrak{m}}
  \bigl(\operatorname{dim.inj}_{A_{\mathfrak{m}}}(M_{\mathfrak{m}})\bigr)
\end{equation}
y, si $M$ es \emph{de tipo finito},
\begin{equation}
\label{0.17.2.5.2-es}
\tag{17.2.5.2}
  \operatorname{dim.proj}_A(M)
  =\sup_{\mathfrak{m}}
  \bigl(\operatorname{dim.proj}_{A_{\mathfrak{m}}}(M_{\mathfrak{m}})\bigr)
\end{equation}
donde $\mathfrak{m}$ recorre el conjunto de los ideales primos (o el conjunto de los
ideales maximales) de $A$.
\end{corollary}

\begin{proof}
En efecto, si $N$ es un $A$-módulo de tipo finito (por tanto, de presentación finita),
se tiene, para todo conjunto multiplicativo $S$ de $A$ y para todo $i\geq0$,
\[
  S^{-1}\operatorname{Ext}^i_A(N,M)
  =\operatorname{Ext}^i_{S^{-1}A}(S^{-1}N,S^{-1}M)
\]
por planitud, considerando una resolución libre de $M$ y utilizando el
hecho de que la relación precedente es válida para $i=0$ (Bourbaki,
\emph{Alg.~comm.}, cap.~II, \S~2, núm.~7, prop.~19). En particular,
$\operatorname{Ext}^i_{A_{\mathfrak{m}}}
(A_{\mathfrak{m}}/\mathfrak{J}A_{\mathfrak{m}},M_{\mathfrak{m}})
=(\operatorname{Ext}^i_A(A/\mathfrak{J},M))_{\mathfrak{m}}$ para todo ideal
primo $\mathfrak{m}$ y todo ideal $\mathfrak{J}$ de $A$; si se tienen
en cuenta (17.2.3) y el hecho de que todo
\oldpage[0\textsubscript{IV-1}]{44}
ideal de $A_{\mathfrak{m}}$ es de la forma $\mathfrak{J}A_{\mathfrak{m}}$
para un ideal adecuado $\mathfrak{J}$ de $A$, se deduce la fórmula
(17.2.5.1) de lo que precede y de Bourbaki, \emph{Alg.~comm.}, cap.~II,
\S~3, núm.~3, teorema~1. Se procede del mismo modo para (17.2.5.2), utilizando esta vez
(17.2.4) e intercambiando los papeles de $M$ y $N$.
\end{proof}

Para los anillos noetherianos y los módulos de tipo finito, el estudio de la
dimensión proyectiva o de la dimensión inyectiva queda, pues, reducido por
(17.2.5) al caso de los anillos \emph{locales}. Se tiene entonces el

\begin{lemma}[17.2.6]
\label{0.17.2.6-es}
Sean $A$ un anillo local noetheriano, $k$ su cuerpo residual y $M$ un
$A$-módulo de tipo finito. Para que $\operatorname{dim.proj}(M)\leq n$, es
necesario que $\operatorname{Tor}^A_i(M,k)=0$ para $i>n$, y basta que
$\operatorname{Tor}^A_{n+1}(M,k)=0$.
\end{lemma}

\begin{proof}
La necesidad es un caso particular de la observación (17.2.2 (ii)), aplicada al
funtor covariante $M\mapsto k\otimes_A M$. Para probar que la condición es
suficiente, se trata, con la notación de (17.2.1), de establecer que $R$ es
proyectivo cuando se supone que los $P_i$ son de tipo finito; ahora bien,
$\operatorname{Tor}^A_{n+1}(M,k)$ es isomorfo a
$\operatorname{Tor}^A_1(R,k)$ (M, V, 7); y se sabe que, puesto que $R$ es de
tipo finito, la condición $\operatorname{Tor}^A_1(R,k)=0$ implica que $R$ es
libre (Bourbaki, \emph{Alg.~comm.}, cap.~II, \S~3, núm.~2, cor.~2 de la
prop.~5).
\end{proof}

\begin{corollary}[17.2.7]
\label{0.17.2.7-es}
Bajo las hipótesis de (17.2.6), sea $x$ un elemento $M$-regular de $A$
perteneciente al ideal maximal de $A$; entonces se tiene
\begin{equation}
\label{0.17.2.7.1-es}
\tag{17.2.7.1}
  \operatorname{dim.proj}(M/xM)=\operatorname{dim.proj}(M)+1.
\end{equation}
\end{corollary}

\begin{proof}
En efecto, de la sucesión exacta
$0\longrightarrow M\xrightarrow{x}M\longrightarrow M/xM\longrightarrow0$
definida por la multiplicación por $x$ en $M$, se deduce la sucesión exacta de
Tor:
\[
  \operatorname{Tor}^A_i(M,k)\xrightarrow{x}
  \operatorname{Tor}^A_i(M,k)\longrightarrow
  \operatorname{Tor}^A_i(M/xM,k)\longrightarrow
  \operatorname{Tor}^A_{i-1}(M,k)\xrightarrow{x}
  \operatorname{Tor}^A_{i-1}(M,k)
\]
para todo $i\geq1$. Ahora bien, para todo $A$-módulo $N$, la homotecia de razón $x$
en $M\otimes_A N$ procede tanto de la homotecia de razón $x$ en
$M$ como de la homotecia de razón $x$ en $N$; se concluye inmediatamente,
puesto que la homotecia de razón $x$ en $k$ es nula por definición, que,
para todo $i$, el homomorfismo
$\operatorname{Tor}^A_i(M,k)\xrightarrow{x}\operatorname{Tor}^A_i(M,k)$ es
\emph{nulo}; dicho de otro modo, se tiene la sucesión exacta
\[
  0\longrightarrow\operatorname{Tor}^A_i(M,k)\longrightarrow
  \operatorname{Tor}^A_i(M/xM,k)\longrightarrow
  \operatorname{Tor}^A_{i-1}(M,k)\longrightarrow0.
\]
Si $\operatorname{dim.proj}(M)=n$, se tiene
$\operatorname{Tor}^A_n(M,k)\neq0$ y
$\operatorname{Tor}^A_{n+1}(M,k)=\operatorname{Tor}^A_{n+2}(M,k)=0$ en virtud
de (17.2.6). De lo que precede resulta, pues, que se tiene
$\operatorname{Tor}^A_{n+1}(M/xM,k)\neq0$ y
$\operatorname{Tor}^A_{n+2}(M/xM,k)=0$, luego
$\operatorname{dim.proj}(M/xM)=n+1$ por (17.2.6).
\end{proof}

\begin{proposition}[17.2.8]
\label{0.17.2.8-es}
\textnormal{(M.~Auslander)} Sea $A$ un anillo. Las condiciones siguientes son
equivalentes:
\begin{enumerate}
  \item[\rm a)] Todo $A$-módulo es de dimensión proyectiva $\leq n$.
  \item[\rm a')] Todo $A$-módulo de tipo finito es de dimensión proyectiva
  $\leq n$.
  \item[\rm b)] Todo $A$-módulo es de dimensión inyectiva $\leq n$.
  \item[\rm c)] Para todo par de $A$-módulos $M$, $N$, se tiene
  $\operatorname{Ext}^{n+1}_A(M,N)=0$.
  \item[\rm c')] Para todo par de $A$-módulos $M$, $N$ tal que $M$ sea de
  tipo finito (o solamente monógeno), se tiene
  $\operatorname{Ext}^{n+1}_A(M,N)=0$.
\end{enumerate}
\end{proposition}

\begin{proof}
\oldpage[0\textsubscript{IV-1}]{45}
Esto resulta inmediatamente de (17.2.1) y (17.2.3).
\end{proof}

El menor número $n$ (entero o $+\infty$) para el cual se satisfacen las condiciones
equivalentes de (17.2.8) se denomina \emph{dimensión cohomológica
global} (o simplemente \emph{dimensión cohomológica}) de $A$
y se denota por $\operatorname{dim.coh}(A)$.

\begin{proposition}[17.2.9]
\label{0.17.2.9-es}
Sea $A$ un anillo noetheriano. Las condiciones siguientes son equivalentes~:
\begin{enumerate}
  \item[\rm a)] $\operatorname{dim.coh}(A)\leq n$.
  \item[\rm b)] Todo $A$-módulo de tipo finito es de dimensión inyectiva
  $\leq n$.
  \item[\rm c)] Para todo par de $A$-módulos de tipo finito, $M$, $N$, se tiene
  $\operatorname{Ext}^{n+1}_A(M,N)=0$.
\end{enumerate}
\end{proposition}

\begin{proof}
Esto resulta inmediatamente de la definición y de (17.2.4).
\end{proof}

\begin{corollary}[17.2.10]
\label{0.17.2.10-es}
Si $A$ es un anillo noetheriano, se tiene
\begin{equation}
\label{0.17.2.10.1-es}
\tag{17.2.10.1}
  \operatorname{dim.coh}(A)
  =\sup_{\mathfrak{m}}
  \bigl(\operatorname{dim.coh}(A_{\mathfrak{m}})\bigr)
\end{equation}
donde $\mathfrak{m}$ recorre el espectro de $A$ (o el conjunto de los ideales
maximales de $A$).
\end{corollary}

\begin{proof}
Esto se deduce de (17.2.9) y (17.2.5).
\end{proof}

\begin{proposition}[17.2.11]
\label{0.17.2.11-es}
Sean $A$ un anillo local noetheriano y $k$ su cuerpo residual. Para que
$\operatorname{dim.coh}(A)\leq n$, es necesario que
$\operatorname{Tor}^A_i(k,k)=0$, para $i>n$, y basta con que
$\operatorname{Tor}^A_{n+1}(k,k)=0$.
\end{proposition}

\begin{proof}
Habida cuenta de (17.2.6), basta demostrar que las relaciones
$\operatorname{dim.coh}(A)\leq n$ y
$\operatorname{dim.proj}_A(k)\leq n$ son equivalentes. Es claro que la
primera implica la segunda por definición. Recíprocamente, si
$\operatorname{dim.proj}_A(k)\leq n$, se tiene
$\operatorname{Tor}^A_i(M,k)=0$ para $i>n$ y para todo $A$-módulo $M$ en
virtud de (17.2.2, (ii))~; por tanto, $\operatorname{dim.proj}(M)\leq n$, lo que
demuestra la proposición en virtud de (17.2.8).
\end{proof}

\begin{corollary}[17.2.12]
\label{0.17.2.12-es}
Sea $A$ un anillo noetheriano. Para que se tenga
$\operatorname{dim.coh}(A)\leq n$, es necesario que, para todo ideal
maximal $\mathfrak{m}$ de $A$, se tenga
$\operatorname{Tor}^{A_{\mathfrak{m}}}_i(A/\mathfrak{m},A/\mathfrak{m})=0$
para $i>n$, y basta con que estas relaciones se verifiquen para $i=n+1$.
\end{corollary}

\begin{proof}
Esto resulta inmediatamente de (17.2.11) y (17.2.10).
\end{proof}

\begin{proposition}[17.2.13]
\label{0.17.2.13-es}
Sean $A$, $B$ dos anillos locales noetherianos y
$\varphi:A\longrightarrow B$ un homomorfismo local que hace de $B$ un
$A$-módulo plano. Entonces se tiene
\begin{equation}
\label{0.17.2.13.1-es}
\tag{17.2.13.1}
  \operatorname{dim.coh}(A)\leq\operatorname{dim.coh}(B).
\end{equation}
\end{proposition}

\begin{proof}
Supongamos que $\operatorname{dim.coh}(B)=n$ es finita~; basta demostrar
que, para todo par $(M,N)$ de $A$-módulos de tipo finito, se tiene
$\operatorname{Tor}^A_i(M,N)=0$ para $i>n$ (17.2.11). Como $B$ es un
$A$-módulo fielmente plano ($0_{\rm I}$, 6.6.2), demostrar esto equivale
($0_{\rm I}$, 6.4.1) a demostrar que se tiene
\begin{equation}
\label{0.17.2.13.2-es}
\tag{17.2.13.2}
  \operatorname{Tor}^A_i(M,N)\otimes_A B=0.
\end{equation}
Ahora bien, si $L_\bullet=(L_j)$ es una resolución derecha de $M$ por $A$-módulos
libres, del hecho de que $B$ sea un $A$-módulo plano se deduce que
$L_\bullet\otimes_A B=(L_j\otimes_A B)$ es una resolución derecha del
$B$-módulo $M\otimes_A B$ por $B$-módulos libres~; además, se tiene
$(L_j\otimes_A B)\otimes_B(N\otimes_A B)
=(L_j\otimes_A N)\otimes_A B$, de donde se concluye inmediatamente que el primer miembro
\oldpage[0\textsubscript{IV-1}]{46}
de (17.2.13.2) es igual a
$\operatorname{Tor}^B_i(M\otimes_A B,N\otimes_A B)$~; la hipótesis sobre $B$
implica que este $B$-módulo es nulo para $i>n$ (17.2.2, (ii)), de donde resulta la
conclusión.
\end{proof}

\begin{env}[17.2.14]
\label{0.17.2.14-es}
Sean $(X,\mathscr{O}_X)$ un espacio anillado y $\mathscr{F}$ un
$\mathscr{O}_X$-Módulo~; se llama \emph{dimensión proyectiva} (resp.
\emph{inyectiva}) \emph{puntual} de $\mathscr{F}$ y se denota por
$\operatorname{dim.proj}(\mathscr{F})$ (resp.
$\operatorname{dim.inj}(\mathscr{F})$) el número
$\sup_{x\in X}(\operatorname{dim.proj}(\mathscr{F}_x))$ (resp.
$\sup_{x\in X}(\operatorname{dim.inj}(\mathscr{F}_x))$. Se llama
\emph{dimensión cohomológica puntual de $X$} y se denota por
$\operatorname{dim.coh}(X)$ el número
$\sup_{x\in X}(\operatorname{dim.coh}(\mathscr{O}_x))$. Se deduce de
(17.2.5) y (17.2.10) que, cuando $A$ es un anillo noetheriano, $M$ un
$A$-módulo de tipo finito y $X=\operatorname{Spec}(A)$, se tiene
$\operatorname{dim.proj}(\widetilde{M})=\operatorname{dim.proj}(M)$,
$\operatorname{dim.inj}(\widetilde{M})=\operatorname{dim.inj}(M)$ y
$\operatorname{dim.coh}(X)=\operatorname{dim.coh}(A)$. Se llama
\emph{dimensión proyectiva} (resp. \emph{inyectiva}) \emph{de $\mathscr{F}$ en
un punto $x\in X$} a la dimensión proyectiva (resp. inyectiva) de
$\mathscr{F}_x$, y \emph{dimensión cohomológica de $X$ en un punto $x$} a la
dimensión cohomológica de $\mathscr{O}_x$.
\end{env}

\begin{proposition}[17.2.15]
\label{0.17.2.15-es}
Sean $X$, $Y$ dos espacios anillados en anillos locales noetherianos y
$f:X\longrightarrow Y$ un morfismo plano. Si
$\operatorname{dim.coh}(X)\leq n$, entonces $Y$ es de dimensión cohomológica
$\leq n$ en todo punto de $f(X)$.
\end{proposition}

\begin{proof}
Esto resulta inmediatamente de (17.2.13).
\end{proof}

\subsection{Teoría cohomológica de los anillos regulares}
\label{subsection:0.17.3-es}

\begin{theorem}[17.3.1]
\label{0.17.3.1-es}
\textnormal{(Hilbert-Serre)} Sea $A$ un anillo local noetheriano. Para que $A$
sea de dimensión cohomológica finita, es necesario y suficiente que $A$ sea
regular~; se tiene entonces
\begin{equation}
\label{0.17.3.1.1-es}
\tag{17.3.1.1}
  \operatorname{dim.coh}(A)=\dim(A).
\end{equation}
\end{theorem}

\begin{proof}
Supongamos que $A$ es regular~; sean $\mathfrak{m}$ su ideal maximal y
$\mathbf{x}=(x_i)_{1\leq i\leq n}$ un sistema regular de parámetros de $A$
(17.1.6)~; consideremos el complejo del álgebra exterior
$K_\bullet(\mathbf{x})$ (\textbf{III}, 1.1.1), que está formado por $A$-módulos
libres, con $K_i(\mathbf{x})=0$ para $i>n$~; como $(x_i)$ es una sucesión
regular, se tiene $H_i(K_\bullet(\mathbf{x}))=0$ para $i>0$ (\textbf{III},
1.1.4 y 1.1.3.3) y
$H_0(K_\bullet(\mathbf{x}))=A/(x_1A+\cdots+x_nA)=A/\mathfrak{m}=k$~; los
$K_i(\mathbf{x})$ forman, por tanto, una \emph{resolución libre} de $k$ de
longitud $n$. Ahora bien, el hecho de que los $x_i$ pertenezcan a $\mathfrak{m}$
implica inmediatamente que, en el complejo
$K_\bullet(\mathbf{x},k)=K_\bullet(\mathbf{x})\otimes_A k$, el operador de borde
es nulo en todos los grados, de modo que se tiene, por definición,
$\operatorname{Tor}^A_i(k,k)=\bigwedge^i(k^n)$~; la igualdad (17.3.1.1) resulta
entonces inmediatamente de (17.2.11) (este resultado es esencialmente el «~teorema de las
sizigias~» de Hilbert).

Demostremos ahora que, si $A$ es un anillo local noetheriano,
$\mathfrak{m}$ su ideal maximal, y si
$\operatorname{dim.proj}_A(\mathfrak{m})$ es \emph{finita}, entonces $A$ es
regular, lo que terminará la demostración de (17.3.1). Procederemos por
inducción sobre $n=\operatorname{rg}_k(\mathfrak{m}/\mathfrak{m}^2)$. Para
$n=0$, se tiene $\mathfrak{m}=0$ y la aserción es trivial.
\end{proof}

\begin{lemma}[17.3.1.2]
\label{0.17.3.1.2-es}
\textnormal{(Nagata)} Si todo elemento de
$\mathfrak{m}-\mathfrak{m}^2$ es divisor de cero en $A$, existe
$c\neq0$ en $A$ tal que $c\mathfrak{m}=0$ (dicho de otro modo, se tiene
$\mathfrak{m}\in\operatorname{Ass}(A)$).
\end{lemma}

\begin{proof}
Podemos limitarnos al caso en que $\mathfrak{m}\neq0$, y por tanto
$\mathfrak{m}\neq\mathfrak{m}^2$. La hipótesis implica que
$\mathfrak{m}-\mathfrak{m}^2$ está contenido en la unión de los ideales
$\mathfrak{p}_i$ de $\operatorname{Ass}(A)$ (Bourbaki,
\emph{Alg.~comm.}, cap.~IV, \S~1, no~1, cor.~3 de la prop.~2)~; por tanto,
$\mathfrak{m}$ está contenido en la unión de $\mathfrak{m}^2$
\oldpage[0\textsubscript{IV-1}]{47}
y de los $\mathfrak{p}_i$, y por tanto en uno de los $\mathfrak{p}_i$, puesto que
$\mathfrak{m}\not\subset\mathfrak{m}^2$ (Bourbaki, \emph{Alg.~comm.},
cap.~II, \S~1, no~1, prop.~2)~; como $\mathfrak{m}$ es maximal, esto demuestra el
lema.
\end{proof}

\begin{lemma}[17.3.1.3]
\label{0.17.3.1.3-es}
Si $a\in\mathfrak{m}-\mathfrak{m}^2$, $\mathfrak{m}/Aa$ es isomorfo a un
factor directo de $\mathfrak{m}/a\mathfrak{m}$.
\end{lemma}

\begin{proof}
Existe un sistema minimal de generadores de $\mathfrak{m}$ que contiene a $a$
(cuyas imágenes en $\mathfrak{m}/\mathfrak{m}^2$ forman una base de este
$k$-espacio vectorial)~; sea $\mathfrak{b}$ el ideal engendrado por los elementos
de este sistema distintos de $a$. Como la relación $xa\in\mathfrak{b}$ implica
$x\in\mathfrak{m}$ (considerando la imagen de $xa$ en
$\mathfrak{m}/\mathfrak{m}^2$), se tiene
$\mathfrak{b}\cap Aa\subset a\mathfrak{m}$, de donde se obtiene un homomorfismo
$\mathfrak{b}/(\mathfrak{b}\cap Aa)\longrightarrow
\mathfrak{m}/a\mathfrak{m}$ deducido de la inyección
$\mathfrak{b}\longrightarrow\mathfrak{m}$ por paso a los cocientes, y que
es \emph{inyectivo}. Además, $\mathfrak{b}+Aa=\mathfrak{m}$, y
el homomorfismo compuesto
\[
  \mathfrak{m}/Aa=(\mathfrak{b}+Aa)/Aa
  \xrightarrow{\sim}\mathfrak{b}/(\mathfrak{b}\cap Aa)
  \longrightarrow\mathfrak{m}/a\mathfrak{m}
  \longrightarrow\mathfrak{m}/Aa
\]
es la identidad~; de donde resulta el lema.
\end{proof}

\begin{lemma}[17.3.1.4]
\label{0.17.3.1.4-es}
Sean $A$ un anillo local noetheriano, $\mathfrak{m}$ su ideal maximal, $E$
un $A$-módulo de tipo finito y de dimensión proyectiva finita. Si
$a\in\mathfrak{m}$ es $A$-regular y $E$-regular, entonces $E/aE$ es un
$(A/aA)$-módulo de dimensión proyectiva finita, a lo sumo igual a
$\operatorname{dim.proj}_A(E)$.
\end{lemma}
\begin{proof}
Razonemos por inducción sobre $h=\operatorname{dim.proj}_A(E)$, siendo trivial
el caso $h=0$, puesto que entonces $E$ es un $A$-módulo proyectivo, y por tanto
$E/aE$ es un $(A/aA)$-módulo proyectivo. Existe una sucesión exacta
\[
  0\longrightarrow N\longrightarrow L\longrightarrow E\longrightarrow0
\]
donde $L$ es libre y $\operatorname{dim.proj}_A(N)=h-1$ (17.2.2, (iii)), siendo
$N$ de tipo finito. Además, la sucesión
\[
  0\longrightarrow N/aN\longrightarrow L/aL\longrightarrow E/aE
  \longrightarrow0
\]
es exacta (15.1.18). Como $a$ es $A$-regular, también es $N$-regular, puesto
que $L$ es libre; la hipótesis de inducción implica que $N/aN$ es un
$(A/aA)$-módulo de dimensión proyectiva $\leq h-1$, y como $L/aL$ es un
$(A/aA)$-módulo libre, $E/aE$ es un $(A/aA)$-módulo de dimensión proyectiva
$\leq h$.
\end{proof}

Examinemos ahora dos casos:

\medskip\noindent\textbf{I.} --- Supongamos primero que todo elemento de
$\mathfrak{m}-\mathfrak{m}^2$ es divisor de cero en $A$, en cuyo caso
(17.3.1.2) existe $c\neq0$ en $A$ tal que $c\mathfrak{m}=0$. Demostremos
entonces que $\mathfrak{m}=0$. Si no fuera así, observemos primero que
$\mathfrak{m}$ no podría ser un $A$-módulo proyectivo, pues sería libre
(Bourbaki, \emph{Alg.~comm.}, cap.~II, \S~3, n.º~2, cor.~2 de la prop.~5),
lo que contradice la relación $c\mathfrak{m}=0$. Se tendría, por tanto,
$n=\operatorname{dim.coh}(A)\geq1$. Como
$\mathfrak{m}\in\operatorname{Ass}(A)$, existiría una sucesión exacta de
$A$-homomorfismos
\[
  0\longrightarrow k\longrightarrow A\longrightarrow E\longrightarrow0.
\]
Pero esto es absurdo, pues, en virtud de la relación $n\geq1$, la sucesión
exacta de los Tor daría la sucesión exacta
$0\longrightarrow\operatorname{Tor}^A_{n+1}(E,k)\longrightarrow
\operatorname{Tor}^A_n(k,k)\longrightarrow0$; ahora bien, se tiene
$\operatorname{Tor}^A_{n+1}(E,k)=0$ (17.2.2, (ii)) y
$\operatorname{Tor}^A_n(k,k)\neq0$ (17.2.11), y hemos llegado a una
contradicción.

\medskip\noindent\textbf{II.} --- Podemos, pues, limitarnos al caso en que
existe $a\in\mathfrak{m}-\mathfrak{m}^2$ que es un elemento $A$-regular y, por
tanto, también $\mathfrak{m}$-regular. Consideremos el anillo $A'=A/aA$ y su
ideal
\oldpage[0\textsubscript{IV-1}]{48}
máximo $\mathfrak{m}'=\mathfrak{m}/aA$; es claro que
$\operatorname{rg}_k(\mathfrak{m}'/\mathfrak{m}'^2)=n-1$. En virtud de
(17.3.1.4), $\mathfrak{m}/a\mathfrak{m}$ es un $A'$-módulo de dimensión
proyectiva finita, luego lo mismo ocurre con
$\mathfrak{m}'=\mathfrak{m}/aA$, que es un sumando directo suyo (17.3.1.3 y
17.2.1). La hipótesis de inducción implica, por tanto, que $A'=A/aA$ es
regular, lo que demuestra, por (17.1.8), que $A$ es regular.

\begin{corollary}[17.3.2]
\label{0.17.3.2-es}
Si $A$ es un anillo local regular, $A_{\mathfrak{p}}$ es regular para todo
ideal primo $\mathfrak{p}$ de $A$.
\end{corollary}

\begin{proof}
En efecto, se ha visto (17.2.10) que
$\operatorname{dim.coh}(A_{\mathfrak{p}})\leq
\operatorname{dim.coh}(A)$, luego la conclusión resulta inmediatamente de
(17.3.1).
\end{proof}

\begin{proposition}[17.3.3]
\label{0.17.3.3-es}
Sean $A$, $B$ dos anillos locales noetherianos,
$\varphi:A\longrightarrow B$ un homomorfismo local, $\mathfrak{m}$ (resp.
$\mathfrak{n}$) el ideal máximo de $A$ (resp. $B$),
$k=A/\mathfrak{m}$ (resp. $k'=B/\mathfrak{n}$) el cuerpo residual de $A$
(resp. $B$). Se tiene, por tanto, un homomorfismo canónico de $k'$-espacios
vectoriales
\begin{equation}
\label{0.17.3.3.1-es}
\tag{17.3.3.1}
  \psi:(\mathfrak{m}/\mathfrak{m}^2)\otimes_k k'
  \longrightarrow\mathfrak{n}/\mathfrak{n}^2.
\end{equation}
\begin{enumerate}
  \item[\rm (i)] Si $B$ es regular y $B$ es un $A$-módulo plano, $A$ es
  regular.
  \item[\rm (ii)] Las condiciones siguientes son equivalentes:
  \begin{enumerate}
    \item[\rm a)] $B$ es regular y el homomorfismo $\psi$ (17.3.3.1) es
    inyectivo.
    \item[\rm b)] $A$ y $B$ son regulares, y para un sistema regular de
    parámetros $(x_i)_{1\leq i\leq m}$ de $A$, los $\varphi(x_i)$ forman parte
    de un sistema regular de parámetros de $B$ (en cuyo caso esta propiedad se
    verifica para \emph{todo} sistema regular de parámetros de $A$).
    \item[\rm c)] $B$ y $B\otimes_A k$ son regulares, y $B$ es un $A$-módulo
    plano.
    \item[\rm d)] $A$ y $B\otimes_A k$ son regulares, y $B$ es un $A$-módulo
    plano.
    \item[\rm e)] $A$ y $B\otimes_A k$ son regulares, y se tiene
    \begin{equation}
    \label{0.17.3.3.2-es}
    \tag{17.3.3.2}
      \dim(B)=\dim(A)+\dim(B\otimes_A k).
    \end{equation}
  \end{enumerate}
\end{enumerate}
\end{proposition}

\begin{proof}
\textnormal{(i)} Se tiene
$\operatorname{dim.coh}(A)\leq\operatorname{dim.coh}(B)$ por (17.2.13), luego
basta aplicar (17.3.1).

\textnormal{(ii)} Cuando se supone que $A$ es regular y $(x_i)$ es un sistema
regular de parámetros de $A$, decir que $B$ es un $A$-módulo \emph{plano}
equivale, en virtud de (15.1.21), a decir que la sucesión de los
$y_i=\varphi(x_i)$ es $B$-regular (puesto que $B\otimes_A k$ es un
$k$-módulo plano). Por otra parte, como $\dim(A)=m$ y la sucesión $(y_i)$
engendra el ideal $\mathfrak{m}B$, la relación (17.3.3.2), que también se
escribe $\dim(B/\mathfrak{m}B)=\dim(B)-m$, equivale a decir que la sucesión
$(y_i)_{1\leq i\leq m}$ forma parte de un sistema de parámetros de $B$
(16.3.7). Se ve, por tanto, que las condiciones b), d) y e) equivalen,
respectivamente, a las siguientes:
\begin{enumerate}
  \item[\rm b')] $A$ y $B$ son regulares y $(y_i)_{1\leq i\leq m}$ forma
  parte de un sistema regular de parámetros de $B$.
  \item[\rm d')] $A$ es regular,
  $B/(\sum_{i=1}^m y_iB)$ es regular y la sucesión $(y_i)$ es
  $B$-regular.
  \item[\rm e')] $A$ es regular,
  $B/(\sum_{i=1}^m y_iB)$ es regular y la sucesión $(y_i)$ forma parte de un
  sistema de parámetros de $B$.
\end{enumerate}
Ahora bien, b') y e') son equivalentes en virtud de (17.1.7), y como d')
implica e') (16.4.1) y se sigue de b') (17.1.7), es equivalente a ellas.
La conjunción
\oldpage[0\textsubscript{IV-1}]{49}
de b) y d) implica trivialmente c), y, en virtud de (i), c) implica d); se ha
probado, por tanto, la equivalencia de b), c), d) y e). Además, es claro que
b) implica a), pues las clases de los $x_i$ en
$\mathfrak{m}/\mathfrak{m}^2$ forman entonces una base de este $k$-espacio
vectorial y las clases de los $y_i$ en $\mathfrak{n}/\mathfrak{n}^2$ un
sistema libre de este $k'$-espacio vectorial. Queda, pues, probar que a)
implica e). Pongamos $V=\mathfrak{m}/\mathfrak{m}^2$,
$W=\mathfrak{n}/\mathfrak{n}^2$; es inmediato que se tiene
$\psi(V\otimes_k k')=(\mathfrak{n}^2+\mathfrak{m}B)/\mathfrak{n}^2$. Por una
parte, en virtud de (16.3.9),
\begin{equation}
\label{0.17.3.3.3-es}
\tag{17.3.3.3}
  \dim(B)\leq\dim(A)+\dim(B\otimes_A k)~;
\end{equation}
en segundo lugar, según (16.2.6), se tiene
\begin{equation}
\label{0.17.3.3.4-es}
\tag{17.3.3.4}
  \dim(A)\leq\operatorname{rg}_{k'}(V\otimes_k k').
\end{equation}
Por último, en el anillo local $B\otimes_A k=B/\mathfrak{m}B$,
$\mathfrak{n}'=\mathfrak{n}/\mathfrak{m}B$ es el ideal máximo, $k'$ el cuerpo
residual y $\mathfrak{n}'/\mathfrak{n}'^2$ es isomorfo a
$\mathfrak{n}/(\mathfrak{n}^2+\mathfrak{m}B)
=W/\psi(V\otimes_k k')$; se tiene, por tanto, según (16.2.6),
\begin{equation}
\label{0.17.3.3.5-es}
\tag{17.3.3.5}
  \dim(B\otimes_A k)
  \leq\operatorname{rg}_{k'}W
  -\operatorname{rg}_{k'}\psi(V\otimes_k k').
\end{equation}
Finalmente, puesto que se supone que $B$ es regular, se tiene
$\dim(B)=\operatorname{rg}_{k'}W$ (17.1.1); se concluye, por tanto, de
(17.3.3.3), (17.3.3.4) y (17.3.3.5) que se tiene
\[
  \operatorname{rg}_{k'}W
  \leq\dim(A)+\dim(B\otimes_A k)
  \leq\operatorname{rg}_{k'}W
  +\operatorname{rg}_{k'}(V\otimes_k k')
  -\operatorname{rg}_{k'}\psi(V\otimes_k k')
\]
y decir que los dos términos extremos de esta desigualdad son iguales
significa que $\psi$ es \emph{inyectivo}. La condición a) implica, por tanto,
necesariamente que en cada una de las relaciones (17.3.3.3), (17.3.3.4) y
(17.3.3.5), los dos miembros son \emph{iguales}; ahora bien, la igualdad en
(17.3.3.4) (resp. (17.3.3.5)) significa que $A$ (resp. $B\otimes_A k$) es
\emph{regular} (17.1.1); queda así demostrado que a) implica e).
\end{proof}

\begin{proposition}[17.3.4]
\label{0.17.3.4-es}
Sean $A$ un anillo local regular de dimensión $n$, $\mathfrak{m}$ su ideal
máximo. Para todo $A$-módulo no nulo de tipo finito $M$, se tiene
\begin{equation}
\label{0.17.3.4.1-es}
\tag{17.3.4.1}
  \operatorname{prof}(M)+\operatorname{dim.proj}(M)=n.
\end{equation}
\end{proposition}

\begin{proof}
Razonemos por inducción sobre $m=\operatorname{prof}(M)$. Si $m=0$, se sabe
(16.4.6, (i)) que existe un sub-$A$-módulo $N$ de $M$ isomorfo a
$k=A/\mathfrak{m}$; aplicando la sucesión exacta de los Tor a la sucesión
exacta $0\longrightarrow N\longrightarrow M\longrightarrow M/N\longrightarrow0$,
se obtiene una sucesión exacta
\[
  \operatorname{Tor}^A_{n+1}(M/N,k)\longrightarrow
  \operatorname{Tor}^A_n(N,k)\longrightarrow
  \operatorname{Tor}^A_n(M,k),
\]
y la hipótesis de que $A$ es regular implica
$\operatorname{Tor}^A_{n+1}(M/N,k)=0$ ((17.3.1.1) y (17.2.6)), luego
$\operatorname{Tor}^A_n(k,k)$ es isomorfo a un sub-$A$-módulo de
$\operatorname{Tor}^A_n(M,k)$; aplicando de nuevo (17.3.1.1) y (17.2.6), se
ve que $\operatorname{Tor}^A_n(M,k)\neq0$, luego
$\operatorname{dim.proj}(M)\geq n$ (17.2.6); pero como
$\operatorname{dim.proj}(M)\leq n$ por (17.3.1.1), se tiene efectivamente
$\operatorname{dim.proj}(M)=n$.

Supongamos ahora que $m>0$, y sea $x$ un elemento $M$-regular perteneciente a
$\mathfrak{m}$; se sabe entonces (16.4.6, (i)) que se tiene
$\operatorname{prof}(M/xM)=m-1$, y, por otra parte, (17.2.7),
$\operatorname{dim.proj}(M/xM)=\operatorname{dim.proj}(M)+1$; la hipótesis de
inducción prueba inmediatamente la relación (17.3.4.1).
\end{proof}

\oldpage[0\textsubscript{IV-1}]{50}

\begin{corollary}[17.3.5]
\label{0.17.3.5-es}
\begin{enumerate}
  \item[\rm(i)] Sea $A$ un anillo local regular de dimensión $n$; para que un
  $A$-módulo $M\neq0$ de tipo finito sea libre, es necesario y suficiente que
  $M$ sea un $A$-módulo de Cohen-Macaulay de dimensión $n$.
  \item[\rm(ii)] Sean $A$ un anillo local regular, $B$ un anillo local,
  $\rho:A\to B$ un homomorfismo local que hace de $B$ un $A$-módulo de tipo
  finito. Para que $B$ sea un $A$-módulo libre, es necesario y suficiente que
  $B$ sea un anillo de Cohen-Macaulay y que $\rho$ sea inyectivo (o,
  equivalentemente (16.1.5), que $\dim(B)=\dim(A)$).
\end{enumerate}
\end{corollary}

\begin{proof}
(i) Esto resulta de (17.3.4.1) y del hecho de que, para un $A$-módulo de tipo
finito, es lo mismo decir que este módulo es proyectivo que decir que es libre
(Bourbaki, \emph{Alg. comm.}, cap.~II, \S~3, n\textsuperscript{o}~2, cor.~2 de
la prop.~5); los $A$-módulos libres de tipo finito $M$ se caracterizan, por
tanto, por la relación $\operatorname{dim.proj}(M)=0$ (17.2.2, (i)).

(ii) Decir que $B$ es un anillo de Cohen-Macaulay equivale a decir que $B$ es
un $A$-módulo de Cohen-Macaulay (16.5.3), luego basta aplicar (i), puesto que
$\dim(B)=\dim(A)$ (16.1.5).
\end{proof}

\begin{env}[17.3.6]
\label{0.17.3.6-es}
Se dice que un anillo noetheriano $A$ es \emph{regular} si, para todo ideal
primo $\mathfrak{p}$ de $A$, el anillo local $A_{\mathfrak{p}}$ es regular;
cuando el propio $A$ es un anillo local, se sigue de (17.3.2) que esta
definición es equivalente a la de (17.1.2). Para que $A$ sea regular, basta,
en virtud de (17.3.2), que $A_{\mathfrak{m}}$ sea regular para todo ideal
máximo $\mathfrak{m}$ de $A$. Además, se sigue inmediatamente de esta
definición que, para todo conjunto multiplicativamente cerrado $S$ de $A$,
$S^{-1}A$ es regular.
\end{env}

\begin{proposition}[17.3.7]
\label{0.17.3.7-es}
Si $A$ es un anillo noetheriano regular, todo anillo de polinomios
$A[T_1,\ldots,T_n]$ es regular.
\end{proposition}

\begin{proof}
Es evidentemente suficiente probar que el anillo de polinomios $B=A[T]$ es
regular; como $B$ es un $A$-módulo libre, para todo ideal primo
$\mathfrak{q}$ de $B$, $B_{\mathfrak{q}}$ es un
$A_{\mathfrak{p}}$-módulo plano, donde $\mathfrak{p}=\mathfrak{q}\cap A$
($0_{\mathrm I}$, 6.3.1); basta, por tanto, en virtud de (17.1.10), probar que
$B_{\mathfrak{q}}/\mathfrak{p}B_{\mathfrak{q}}$ es regular, y como este anillo
es un anillo local en un ideal primo de
$A_{\mathfrak{p}}[T]/\mathfrak{p}A_{\mathfrak{p}}[T]=k[T]$ (donde $k$ es el
cuerpo residual $A_{\mathfrak{p}}/\mathfrak{p}A_{\mathfrak{p}}$), basta
probar que $C=k[T]$ es regular; ahora bien, como $C$ es un dominio de ideales
principales, los anillos locales en los ideales primos de $C$ son anillos de
valoración discreta o un cuerpo (para el ideal $(0)$), luego son regulares
(17.1.4), lo que termina la demostración.
\end{proof}

\begin{corollary}[17.3.8]
\label{0.17.3.8-es}
Si $A$ es un anillo regular, todo anillo de series formales
$A[[T_1,\ldots,T_n]]$ es regular.
\end{corollary}

\begin{proof}
Sea $\mathfrak{J}$ el ideal engendrado por los $T_i$ en el anillo de polinomios
$A[T_1,\ldots,T_n]$; como este último es regular por (17.3.7), y
$A[[T_1,\ldots,T_n]]$ es el completado de $A[T_1,\ldots,T_n]$ para la
topología $\mathfrak{J}$-preádica (Bourbaki, \emph{Alg. comm.}, cap.~III,
\S~2, n\textsuperscript{o}~12, \emph{Ejemplo 1}), la conclusión resulta del lema siguiente.
\end{proof}

\begin{lemma}[17.3.8.1]
\label{0.17.3.8.1-es}
Sean $A$ un anillo regular, $\mathfrak{J}$ un ideal de $A$,
$\widehat{A}$ el completado separado de $A$ para la topología
$\mathfrak{J}$-preádica. Entonces $\widehat{A}$ es regular.
\end{lemma}

\begin{proof}
Basta, en efecto, (17.3.6) ver que, para todo ideal máximo
$\mathfrak{n}$ de $\widehat{A}$, el anillo local
$\widehat{A}_{\mathfrak{n}}$ es regular; se sabe (Bourbaki,
\emph{Alg. comm.}, cap.~III, \S~3, n\textsuperscript{o}~4, prop.~8) que
$\mathfrak{n}$
\oldpage[0\textsubscript{IV-1}]{51}
es de la forma $\mathfrak{m}\widehat{A}$, donde $\mathfrak{m}$ es un ideal
máximo de $A$ que contiene a $\mathfrak{J}$, y que el homomorfismo canónico
$A\to\widehat{A}$ da un homomorfismo \emph{inyectivo}
$A_{\mathfrak{m}}\to\widehat{A}_{\mathfrak{n}}$, tal que la topología
$\mathfrak{n}\widehat{A}_{\mathfrak{n}}$-preádica de
$\widehat{A}_{\mathfrak{n}}$ induzca sobre $A_{\mathfrak{m}}$ la topología
$\mathfrak{m}A_{\mathfrak{m}}$-preádica, y que $A_{\mathfrak{m}}$ sea denso
en $\widehat{A}_{\mathfrak{n}}$ para la topología
$\mathfrak{n}\widehat{A}_{\mathfrak{n}}$-preádica; por consiguiente, los
\emph{completados} de los anillos locales noetherianos $A_{\mathfrak{m}}$ y
$\widehat{A}_{\mathfrak{n}}$ son canónicamente \emph{isomorfos}. Ahora bien,
por hipótesis $A_{\mathfrak{m}}$ es regular, luego también lo es su completado
(17.1.5), y como el completado de $\widehat{A}_{\mathfrak{n}}$ es regular,
también lo es $\widehat{A}_{\mathfrak{n}}$ por (17.1.5).
\end{proof}

\begin{corollary}[17.3.9]
\label{0.17.3.9-es}
Sea $A$ un anillo noetheriano, cociente de un anillo noetheriano regular $B$.
Si $C$ es una $A$-álgebra de tipo finito, todo anillo de fracciones $S^{-1}C$
de $C$ es cociente de un anillo regular.
\end{corollary}

\begin{proof}
En efecto, $C$ es cociente de un anillo de polinomios $A[T_1,\ldots,T_n]$;
como $A=B/\mathfrak{J}$, donde $\mathfrak{J}$ es un ideal de $B$, se tiene
$A[T_1,\ldots,T_n]=B[T_1,\ldots,T_n]/\mathfrak{J}B[T_1,\ldots,T_n]$; en
virtud de (17.3.7), se puede, por tanto, limitar al caso en que $C$ es cociente
de un anillo regular $B'$; pero si $S'$ es la imagen inversa de $S$ en $B'$,
$S'$ es un conjunto multiplicativamente cerrado de $B'$ y $C$ es un anillo
cociente de $S'^{-1}B'$, de modo que finalmente todo se reduce a ver que,
cuando $C$ es \emph{regular}, también lo es $S^{-1}C$, lo que se ha visto en
(17.3.6).
\end{proof}

La importancia de los cocientes de anillos regulares se debe, entre otras
cosas, a la propiedad precedente y al hecho de que son anillos
\emph{catenarios} (16.5.12).

Todos los anillos importantes en las aplicaciones a la Geometría algebraica
son cocientes de anillos regulares.

\section{Complementos sobre las extensiones de álgebras}
\label{section:0.18-es}

Este parágrafo reúne cierto número de construcciones funtoriales sobre
los anillos, que se utilizarán repetidamente en los \S\S~19 y 20~; no
contiene ningún resultado no trivial.

\subsection{Imágenes inversas de anillos aumentados}
\label{subsection:0.18.1-es}

\begin{env}[18.1.1]
\label{0.18.1.1-es}
Dado un anillo $A$ (no necesariamente conmutativo), la categoría de los
\emph{$A$-anillos} tiene por objetos los pares $(B,\rho)$ formados por un anillo
$B$ y un homomorfismo de anillos $\rho:A\to B$, y por morfismos (llamados también
\emph{$A$-homomorfismos}) los homomorfismos de anillos $u:B\to C$ tales que,
si $\rho:A\to B$ y $\sigma:A\to C$ son los homomorfismos (llamados
\emph{estructurales}) que definen la estructura de $A$-anillo sobre $B$ y $C$
respectivamente, el diagrama
\begin{equation}
\label{0.18.1.1.1-es}
\tag{18.1.1.1}
  \xymatrix{
    B\ar[r]^u & C\\
    & A\ar[lu]^\rho\ar[u]_\sigma
  }
\end{equation}
sea conmutativo. El \emph{núcleo} $\mathfrak{J}$ de $u$ es un ideal bilateral
que es un $A$-\emph{bimódulo} (para $\rho$).

Si $f:A'\to A$ es un homomorfismo de anillos, para todo $A$-anillo
$(B,\rho)$, $(B,\rho\circ f)$ es un $A'$-anillo, y si $u:B\to C$ es un
$A$-homomorfismo del $A$-anillo $(B,\rho)$
\oldpage[0\textsubscript{IV-1}]{52}
en el $A$-anillo $(C,\sigma)$, es también un $A'$-homomorfismo del
$A'$-anillo $(B,\rho\circ f)$ en el $A'$-anillo $(C,\sigma\circ f)$~; se
define así un funtor canónico de la categoría de los $A$-anillos en la de
los $A'$-anillos.
\end{env}

\begin{env}[18.1.2]
\label{0.18.1.2-es}
Sean $B$, $E$, $F$ tres $A$-anillos, $f:E\to B$, $g:F\to B$ dos
$A$-homomorfismos. Recordemos que se llama \emph{producto fibrado} de $E$ y
$F$ sobre $B$ (para los $A$-homomorfismos $f$ y $g$) y se denota por
$E\times_BF$ el subanillo $G$ del anillo producto $E\times F$, formado por los
pares $(x,y)$ tales que $f(x)=g(y)$~; las restricciones $p_1:G\to E$,
$p_2:G\to F$ de las proyecciones $\operatorname{pr}_1$ y
$\operatorname{pr}_2$ a $G$ siguen llamándose las \emph{proyecciones
canónicas}~; la estructura de $A$-anillo de $G$ está definida por
el homomorfismo $\alpha\mapsto(\rho(\alpha),\sigma(\alpha))$ (donde
$\rho:A\to E$ y $\sigma:A\to F$ son los homomorfismos estructurales), que
envía efectivamente $A$ a $G$ en virtud de (18.1.1). La propiedad
característica de $E\times_BF$ es que, para todo par de
$A$-homomorfismos $u:C\to E$, $v:C\to F$ tales que
$f\circ u=g\circ v$, existe un único $A$-homomorfismo $w:C\to G$
tal que $u=p_1\circ w$, $v=p_2\circ w$. También puede decirse que
$E\times_BF$ es el \emph{límite proyectivo} del sistema proyectivo formado por $B$,
$E$, $F$ y los $A$-homomorfismos $f$, $g$, en la categoría de los
$A$-anillos ($0_{\mathrm{III}}$, 8.1.9).
\end{env}

\begin{env}[18.1.3]
\label{0.18.1.3-es}
Sea $\mathfrak{J}$ el ideal bilateral de $E$, núcleo de $f$~; es inmediato
que el núcleo $\mathfrak{J}'$ de $p_2:G\to F$ es el ideal formado por los elementos
$(x,0)$ donde $x\in\mathfrak{J}$~; la restricción
$i_1:\mathfrak{J}'\to\mathfrak{J}$ de $p_1$ es, por tanto, un
\emph{isomorfismo} de anillo \emph{sin elemento unidad}, y también un
isomorfismo de $A$-bimódulos. Del mismo modo, si $\mathfrak{R}$ es el núcleo de
$g$, el núcleo $\mathfrak{R}'$ de $p_1:G\to E$ es el ideal de los elementos
$(0,y)$ donde $y\in\mathfrak{R}$, y la restricción
$i_2:\mathfrak{R}'\to\mathfrak{R}$ de $p_2$ es un isomorfismo (en el mismo
sentido). Finalmente, es claro que el núcleo del $A$-homomorfismo
$q=f\circ p_1=g\circ p_2$ de $E\times_BF$ en $B$ es
$\mathfrak{J}\times\mathfrak{R}=\mathfrak{J}'\oplus\mathfrak{R}'$, de modo
que se tiene el diagrama conmutativo
\begin{equation}
\label{0.18.1.3.1-es}
\tag{18.1.3.1}
  \xymatrix@C=2.2pc@R=1.8pc{
    0\ar[dr] & & 0\ar[d] & 0\ar[d]\\
    & \mathfrak{J}'\oplus\mathfrak{R}'\ar[r]\ar[d]\ar[dr]
      & \mathfrak{R}'\ar[r]^{i_2}_{\sim}\ar[d]^j
      & \mathfrak{R}\ar[d]^{p_2}\\
    0\ar[r] & \mathfrak{J}'\ar[r]\ar[d]^{i_1}_{\sim}
      & G\ar[r]\ar[d]^{p_1}\ar[dr]^q & F\ar[d]^g\\
    0\ar[r] & \mathfrak{J}\ar[r] & E\ar[r]_f & B
  }
\end{equation}

Las definiciones y los resultados de (18.1.2) y (18.1.3) se extienden inmediatamente al
producto fibrado de una \emph{familia} arbitraria $(E_\lambda)_{\lambda\in I}$ de
$A$-anillos definida por una familia de $A$-homomorfismos
$f_\lambda:E_\lambda\to B$. Dejamos al lector la formulación de estos resultados.
\end{env}

\begin{env}[18.1.4]
\label{0.18.1.4-es}
De acuerdo con la terminología de (M, VIII), llamaremos
\emph{$A$-anillo aumentado sobre $B$} a un $A$-anillo $E$ provisto de un
$A$-homomorfismo \emph{sobreyectivo} (llamado \emph{aumentación} de $E$),
$f:E\to B$~; el núcleo $\mathfrak{J}$ de $f$ se llama
\emph{ideal de aumentación}. Se dice que el $A$-anillo aumentado $E$ es
\emph{trivial} si existe un $A$-homomorfismo de $A$-anillos $s:B\to E$
que es \emph{inverso por la derecha} de la aumentación $f:E\to B$ (dicho de otro modo,
$f\circ s=1_B$). La sucesión exacta de $A$-bimódulos
\[
  0\longrightarrow\mathfrak{J}\longrightarrow E
  \xrightarrow{f}B\longrightarrow0
\]
\oldpage[0\textsubscript{IV-1}]{53}
es entonces \emph{escindida}~; dicho de otro modo, se puede identificar el $A$-bimódulo
$E$ con $B\times\mathfrak{J}$, y la multiplicación en $E$ viene dada entonces
por
\[
  (b,z)(b',z')=(bb',bz'+zb'+zz'),
\]
considerándose $\mathfrak{J}$ como $B$-bimódulo por medio de $s:B\to E$.
\end{env}

\begin{env}[18.1.5]
\label{0.18.1.5-es}
Con las notaciones de (18.1.2), supongamos que el $A$-homomorfismo $f$ es
\emph{sobreyectivo}, dicho de otro modo, que hace de $E$ un $A$-anillo \emph{aumentado sobre
$B$}. Entonces es claro que $p_2:G\to F$ es también \emph{sobreyectivo},
dicho de otro modo, define sobre $G$ una estructura de $A$-anillo \emph{aumentado sobre
$F$}, que se llama \emph{imagen inversa} por $g:F\to B$ del anillo
aumentado $E$.
\end{env}

\begin{proposition}[18.1.6]
\label{0.18.1.6-es}
Para que la imagen inversa $G=E\times_BF$ por $g:F\to B$ del $A$-anillo
aumentado $E$ sea un $A$-anillo aumentado trivial, es necesario y suficiente que
exista un $A$-homomorfismo $u:F\to E$ que haga conmutativo el diagrama
\[
  \xymatrix{
    & F\ar[ld]_u\ar[d]^g\\
    E\ar[r]_f & B
  }
\]
\end{proposition}

\begin{proof}
La condición es evidentemente necesaria, tomando $u=p_1\circ s$, donde
$s:F\to G$ es un $A$-homomorfismo inverso por la derecha de $p_2$. Recíprocamente,
si existe un $A$-homomorfismo $u$ que verifica la condición del enunciado,
la existencia del inverso por la derecha $s$ de $p_2$ resulta de la propiedad
universal del producto fibrado (18.1.2) aplicada a los $A$-homomorfismos
$u:F\to E$ y $1_F:F\to F$.

Este resultado implica, en particular, que si $E$ es un anillo aumentado
\emph{trivial}, ocurre lo mismo con todas sus imágenes inversas.
\end{proof}

\begin{env}[18.1.7]
\label{0.18.1.7-es}
Retomemos la situación descrita en (18.1.2) y (18.1.3)~; se tiene evidentemente sobre
$\mathfrak{R}$ una estructura de $G$-bimódulo, procedente de su estructura de
$F$-bimódulo y del homomorfismo de anillos $p_2:G\to F$. Como $i_2$ es
biyectivo, se tiene además una inyección
$\theta=j\circ i_2^{-1}:\mathfrak{R}\to G$, que es un homomorfismo de
$G$-bimódulos. Vamos a ver, recíprocamente, que, cuando se supone además
que $f$ y $g$ son \emph{sobreyectivos}, dicho de otro modo, que $E$ y $F$ son
$A$-anillos \emph{aumentados} sobre $B$, el dato de tal homomorfismo
$\theta$ permite reconstruir el anillo aumentado $E$ (sobre $B$) a partir del
anillo aumentado $G$ (sobre $F$).

Más precisamente, démonos un $A$-anillo $F$ aumentado sobre $B$ y un
$A$-anillo $G$ aumentado sobre $F$, y denotemos por $\mathfrak{R}$ y
$\mathfrak{J}'$, respectivamente, sus ideales de aumentación~:
\begin{equation}
\label{0.18.1.7.1-es}
\tag{18.1.7.1}
  \xymatrix@C=2pc@R=1.5pc{
    & & 0\ar[d] & & \\
    & & \mathfrak{R}\ar[ld]^\theta\ar[d] & & \\
    0\ar[r] & \mathfrak{J}'\ar[r] & G\ar[r]_h & F\ar[r]\ar[d]^g & 0\\
    & & & B\ar[d] & \\
    & & & 0 &
  }
\end{equation}

\oldpage[0\textsubscript{IV-1}]{54}
El homomorfismo $h$ define sobre todos los ideales de $F$ (y, en particular, sobre
$\mathfrak{R}$) una estructura de $G$-bimódulo. Sea
$\theta:\mathfrak{R}\to G$ un homomorfismo de $G$-bimódulos que haga
conmutativo el diagrama (18.1.7.1)~; esto implica que $\theta$ es
\emph{inyectivo} y que $\mathfrak{J}'\cap\theta(\mathfrak{R})=0$~; como
$\theta(\mathfrak{R})$ es un ideal bilateral de $G$, y
$h(\theta(\mathfrak{R}))=\mathfrak{R}$, el homomorfismo compuesto
$g\circ h:G\to B$ se factoriza como
$G\xrightarrow{q}G/\theta(\mathfrak{R})\xrightarrow{f}B$, donde $f$ es
sobreyectivo~; además, la imagen $\mathfrak{J}$ de $\mathfrak{J}'$ en
$E=G/\theta(\mathfrak{R})$ es el núcleo de $f$ (siendo el de $g\circ h$
$\mathfrak{J}'\oplus\theta(\mathfrak{R})$), y la restricción a
$\mathfrak{J}'$ del homomorfismo canónico $q:G\to E$ es inyectiva. Se
puede, por tanto, formar el \emph{producto fibrado} $G'=E\times_BF$, y como los dos
$A$-homomorfismos $q:G\to E$ y $h:G\to F$ son tales que
$f\circ q=g\circ h$, definen un único $A$-homomorfismo $u:G'\to G$
por la propiedad universal de $G'$ (18.1.2)~; vamos a ver que $u$ es
\emph{biyectivo}. Basta demostrarlo cuando $u$ se considera como
homomorfismo de $A$-bimódulos~; pero se observará entonces que $u$ es compatible
con las filtraciones finitas sobre $G'$ y $G$ formadas respectivamente por $G'$
y $\mathfrak{J}''=\operatorname{Ker}(G'\to F)$, y por $G$ y
$\mathfrak{J}'$~; además, se ha visto en (18.1.3) que
$\operatorname{gr}_0u:F\to F$ es la identidad y que
$\operatorname{gr}_1u:\mathfrak{J}''\to\mathfrak{J}'$ es biyectivo, por lo que $u$
mismo es biyectivo (Bourbaki, \emph{Alg. comm.}, cap.~III, \S~2,
n\textsuperscript{o}~8, cor.~3 del teor.~1).
\end{env}

\subsection{Extensiones de un anillo por un bimódulo}
\label{subsection:0.18.2-es}

\begin{env}[18.2.1]
\label{0.18.2.1-es}
Sean $E$ un $A$-anillo aumentado sobre $B$, $f:E\to B$ la aumentación,
$\mathfrak{J}=\operatorname{Ker}(f)$ el ideal de aumentación. Si se tiene
$\mathfrak{J}^2=0$, $\mathfrak{J}$ no es solamente un $E$-bimódulo, sino
también un $B$-\emph{bimódulo}, puesto que $B$ es isomorfo a
$E/\mathfrak{J}$~; más precisamente, todo $b\in B$ es de la forma $f(x)$
con $x\in E$, y si $z,z'$ están en $\mathfrak{J}$ se tiene
$(x+z)z'=xz'$ (resp. $z'(x+z)=z'x$), de modo que el valor de $xz'$ (resp.
$z'x$) no depende del elemento $x$ de $f^{-1}(b)$, y puede escribirse $bz'$
(resp. $z'b$), lo que define la estructura de $B$-bimódulo considerada.
Recíprocamente, si $\mathfrak{J}$ está provisto de una estructura de $B$-bimódulo tal
que $xz'=f(x)z'$ y $z'x=z'f(x)$ para $x\in E$ y
$z'\in\mathfrak{J}$, es claro que $\mathfrak{J}^2=0$.
\end{env}

\begin{env}[18.2.2]
\label{0.18.2.2-es}
Se llama \emph{$A$-extensión de un $A$-anillo $B$ por un $B$-bimódulo $L$}
a una \emph{sucesión exacta} de homomorfismos de $A$-bimódulos
\[
  0\longrightarrow L\xrightarrow{j}E\xrightarrow{f}B\longrightarrow0
\]
donde $E$ es un $A$-anillo, $f$ un $A$-homomorfismo de anillos y se tiene, para
$x\in E$ y $z\in L$,
\[
  j(f(x)z)=xj(z),\qquad j(zf(x))=j(z)x
\]
de donde se deduce, en virtud de (18.2.1), que $j(L)$ es un ideal bilateral de cuadrado nulo de $E$.
Por abuso de lenguaje, se dirá también que $E$ es una extensión de $B$ por $L$.
Se dice que dos $A$-extensiones $E$, $E'$ de $B$ por $L$ son
\emph{$A$-equivalentes} si existe un isomorfismo de $A$-anillos
$u:E\xrightarrow{\sim}E'$ (llamado también \emph{$A$-equivalencia de
$A$-extensiones}) que hace conmutativo el diagrama
\begin{equation}
\label{0.18.2.2.1-es}
\tag{18.2.2.1}
  \xymatrix@C=2.1pc@R=1.2pc{
    & & E\ar[dr]\ar[dd]_u\ar@{}[dd]|{\wr} & & \\
    0\ar[r] & L\ar[ur]\ar[dr] & & B\ar[r] & 0\\
    & & E'\ar[ur] & &
  }
\end{equation}
\end{env}

\oldpage[0\textsubscript{IV-1}]{55}

\begin{env}[18.2.3]
\label{0.18.2.3-es}
Se dice que una $A$-extensión $E$ de $B$ por un $B$-bimódulo $L$ es
\emph{$A$-trivial} si $E$ es un $A$-anillo aumentado (sobre $B$) trivial
(18.1.4). Se define sobre el $A$-bimódulo producto $B\times L$ una estructura de
$A$-anillo poniendo $(x,s)(y,t)=(xy,xt+sy)$, como se comprueba inmediatamente,
y es inmediato que las aplicaciones canónicas $j:L\to B\times L$ y
$f:B\times L\to B$ definen una $A$-extensión de $B$ por $L$ que es
$A$-trivial, siendo la aplicación canónica $g:B\to B\times L$ un
$A$-homomorfismo inverso por la derecha de $f$. Se dice que esta extensión es
\emph{la extensión trivial tipo} de $B$ por $L$ y se denota por $D_B(L)$; es
inmediato que toda $A$-extensión $A$-trivial de $B$ por $L$ es
$A$-equivalente a $D_B(L)$.

Se observará que toda extensión del propio $A$-anillo $A$ por un $A$-bimódulo
es necesariamente $A$-trivial.
\end{env}

\begin{env}[18.2.4]
\label{0.18.2.4-es}
Dadas dos $A$-extensiones
$L\xrightarrow{j}E\xrightarrow{f}B$,
$L'\xrightarrow{j'}E'\xrightarrow{f'}B'$, un \emph{morfismo} de la primera
en la segunda es, por definición, una terna de homomorfismos de
$A$-bimódulos $(u,v,w)$ tal que el diagrama
\begin{equation}
\label{0.18.2.4.1-es}
\tag{18.2.4.1}
  \xymatrix{
    0\ar[r] & L\ar[r]^j\ar[d]_w & E\ar[r]^f\ar[d]_u
      & B\ar[r]\ar[d]_v & 0\\
    0\ar[r] & L'\ar[r]_{j'} & E'\ar[r]_{f'} & B'\ar[r] & 0
  }
\end{equation}
sea conmutativo, siendo $u$ y $v$ $A$-homomorfismos \emph{de anillos} y
$w$ tal que $w(bz)=v(b)w(z)$ y $w(zb)=w(z)v(b)$ para $z\in L$ y
$b\in B$ (en otros términos, el par $(v,w)$ constituye un
\emph{dihomomorfismo} del $B$-bimódulo $L$ en el $B'$-bimódulo $L'$); es
claro que, si $(u',v',w')$ es un morfismo de
$L'\to E'\to B'$ en una $A$-extensión $L''\to E''\to B''$,
$(u'\circ u,v'\circ v,w'\circ w)$ es también un morfismo, lo que justifica la
terminología.

La consideración de los dos cuadrados conmutativos del diagrama (18.2.4.1) va a
conducirnos a dos operaciones sobre las extensiones de $A$-anillos.
\end{env}

\begin{env}[18.2.5]
\label{0.18.2.5-es}
En primer lugar, consideremos una $A$-extensión $E'$ de $B'$ por $L'$
\[
  0\longrightarrow L'\xrightarrow{j'}E'\xrightarrow{f'}B'\longrightarrow0
\]
y un $A$-homomorfismo de anillos $v:B\to B'$, y sea
$F=E'\times_{B'}B$ \emph{la imagen inversa} por $v$ del $A$-anillo aumentado
$E'$ (18.1.5), de modo que se tiene un diagrama conmutativo
\[
  \xymatrix{
    0\ar[r] & L_0\ar[r]\ar[d]^i\ar@{}[d]|{\wr} & F\ar[r]^{p_2}\ar[d]^{p_1}
      & B\ar[r]\ar[d]^v & 0\\
    0\ar[r] & L'\ar[r]_{j'} & E'\ar[r]_{f'} & B'\ar[r] & 0
  }
\]
cuyas filas son exactas, siendo $p_1$ y $p_2$ los homomorfismos
canónicos; se ha visto (18.1.3) que $i$ es biyectivo y de la
definición (18.1.2) resulta también que $L_0^2=0$, de modo que se puede considerar que $F$
es una $A$-extensión de $B$ por $L'$, que se denomina \emph{imagen
inversa por $v$} de la extensión $E'$ de $B'$ por $L'$ (considerándose
naturalmente $L'$ como $B$-bimódulo mediante el homomorfismo
de anillos $v$). El carácter funtorial
\oldpage[0\textsubscript{IV-1}]{56}
del producto fibrado respecto de cada uno de los factores muestra además que, si se
tiene un morfismo de dos extensiones de $B'$
\[
  \xymatrix{
    0\ar[r] & L'_1\ar[r]\ar[d]^h & E'_1\ar[r]\ar[d]^g
      & B'\ar[r]\ar[d]^{1_{B'}} & 0\\
    0\ar[r] & L'_2\ar[r] & E'_2\ar[r] & B'\ar[r] & 0
  }
\]
se deduce de él un morfismo \emph{imagen inversa por $v$}
\[
  \xymatrix{
    0\ar[r] & L'_1\ar[r]\ar[d]^h & E'_1\times_{B'}B\ar[r]
      \ar[d]^{g\times_{B'}1_B} & B\ar[r]\ar[d]^{1_B} & 0\\
    0\ar[r] & L'_2\ar[r] & E'_2\times_{B'}B\ar[r] & B\ar[r] & 0.
  }
\]
En particular, si $E'_1$ y $E'_2$ son $A$-extensiones
$A$-equivalentes de $B'$ por $L'$, sus imágenes inversas por $v$ son
$A$-extensiones $A$-equivalentes de $B$ por $L'$.

La definición del producto fibrado muestra que, cuando se tiene un morfismo
(18.2.4.1) de $A$-extensiones, este se \emph{factoriza a través de} la imagen
inversa de $E'$ por $v$; más precisamente, existe un
$A$-homomorfismo único $u_0:E\to F=E'\times_{B'}B$ que hace conmutativo el
diagrama
\[
  \xymatrix{
    0\ar[r] & L\ar[r]^j\ar[d]^{w_0} & E\ar[r]^f\ar[d]^{u_0}
      & B\ar[r]\ar[d]^{1_B} & 0\\
    0\ar[r] & L_0\ar[r]\ar[d]^i\ar@{}[d]|{\wr} & F\ar[r]^{p_2}\ar[d]^{p_1}
      & B\ar[r]\ar[d]^v & 0\\
    0\ar[r] & L'\ar[r]_{j'} & E'\ar[r]_{f'} & B'\ar[r] & 0
  }
\]
donde $w_0$ es la restricción de $u_0$ a $L$ y
$p_1\circ u_0=u$, $i\circ w_0=w$.
\end{env}

\begin{env}[18.2.6]
\label{0.18.2.6-es}
Estudiemos en particular las imágenes inversas de extensiones por
$A$-homomorfismos \emph{sobreyectivos}. Consideremos un $A$-homomorfismo
sobreyectivo $v:B\to B'$, una $A$-extensión $F$ de $B$ por un $B$-bimódulo $L$,
y el ideal bilateral $\mathfrak{R}$ de $B$, núcleo de $v$ (que puede ser
considerado como $F$-bimódulo mediante la aumentación $F\to B$); se ha visto
(18.1.7) que todo homomorfismo de $F$-bimódulos
$\theta:\mathfrak{R}\to F$ que hace conmutativo el diagrama
\[
  \xymatrix{
    & \mathfrak{R}\ar[ld]^\theta\ar[d]\\
    F\ar[r] & B
  }
\]
determina una extensión $E'$ de $B'=B/\mathfrak{R}$ por $L$ cuya imagen
inversa por $v:B\to B'$ es equivalente a $F$, y que toda $A$-extensión
$E'$ de $B'$ por $L$ que tenga esta última propiedad se obtiene así (salvo
$A$-equivalencia). Además, para que dos homomorfismos
$\theta_1,\theta_2$ de $F$-bimódulos de $\mathfrak{R}$ en $F$ den dos
  $A$-extensiones $A$-equivalentes de $B$ por $L$, es necesario y suficiente que
exista una $A$-equivalencia $u$ de la $A$-extensión $F$ sobre sí misma tal que
\oldpage[0\textsubscript{IV-1}]{57}
$\theta_2=u\circ\theta_1$; esto resulta inmediatamente de lo que se ha visto en
(18.2.5) y de la definición de la biyección canónica de $F$ sobre el producto
fibrado $E'\times_{B'}B$ (18.1.7).
\end{env}

\begin{env}[18.2.7]
\label{0.18.2.7-es}
Consideremos ahora el cuadrado izquierdo de (18.2.4.1), y recordemos primero
la noción de \emph{suma amalgamada} en la categoría de los $A$-bimódulos:
dados tres $A$-bimódulos $X$, $Y$, $Z$ y dos $A$-homomorfismos
$f:X\to Y$, $g:X\to Z$, la suma amalgamada $Y\oplus_XZ$ es el \emph{límite
inductivo} del sistema inductivo formado por $X$, $Y$, $Z$ y los
$A$-homomorfismos $f$, $g$ en la categoría de los $A$-bimódulos
($0_{\mathrm{III}}$, 8.1.11). Se define este $A$-bimódulo como cociente del
producto $Y\times Z$ por el sub-$A$-bimódulo $M$ imagen de $X$ por
el homomorfismo $x\mapsto(f(x),-g(x))$. Su propiedad característica es que,
para todo par de homomorfismos de $A$-bimódulos $u:Y\to T$, $v:Z\to T$
tales que $u\circ f=v\circ g$, existe un único homomorfismo
$w:Y\oplus_XZ\to T$ tal que $u=w\circ j_1$ y $v=w\circ j_2$, donde
$j_1:Y\to Y\oplus_XZ$ y $j_2:Z\to Y\oplus_XZ$ son las aplicaciones
canónicas.
\end{env}

\begin{env}[18.2.8]
\label{0.18.2.8-es}
Consideremos ahora una $A$-extensión $E$ de $B$ por un $B$-bimódulo
$L$:
\[
  0\longrightarrow L\xrightarrow{j}E\xrightarrow{f}B\longrightarrow0
\]
y sea, por otra parte, $w:L\to L'$ un homomorfismo de $B$-bimódulos. Sea $H$
el $A$-bimódulo suma amalgamada $E\oplus_LL'$; mostremos cómo se puede dotar
este $A$-bimódulo de una estructura de $A$-anillo y definir una $A$-extensión
\[
  0\longrightarrow L'\longrightarrow H\longrightarrow B\longrightarrow0.
\]

Observemos para ello que $L'$ está dotado de una estructura de $E$-\emph{bimódulo} por
medio del homomorfismo de aumentación $E\to B$; por tanto, se puede formar la
$A$-extensión trivial tipo $G=D_E(L')$ (18.2.3). Consideremos entonces
la aplicación $\theta:z\mapsto(j(z),-w(z))$ de $L$ en $G$; es un
homomorfismo de $G$-bimódulos (considerándose $L$ como $G$-bimódulo por
medio del homomorfismo canónico $\rho:G\to E$). En efecto, para
$(x,z')\in G$ y $z\in L$, se tiene
$(j(z),-w(z))(x,z')=(j(z)x,j(z)z'-w(z)f(x))$; ahora bien,
$j(z)z'=f(j(z))z'=0$ por definición de la estructura de $E$-bimódulo sobre
$L'$, y $w(z)f(x)=w(zf(x))$ por definición de la estructura de $B$-bimódulo
sobre $L'$; se comprueba del mismo modo que $\theta$ es un homomorfismo de
$G$-módulos izquierdos. Se puede aplicar entonces al diagrama conmutativo
\[
  \xymatrix@C=2pc@R=1.4pc{
    & & 0\ar[d] & & \\
    & & L\ar[ld]^\theta\ar[d]^j & & \\
    0\ar[r] & L'\ar[r] & G\ar[r]_p & E\ar[r]\ar[d]^f & 0\\
    & & & B\ar[d] & \\
    & & & 0 &
  }
\]
el resultado de (18.1.7). Como $H=G/\theta(L)$ por definición, nuestra
afirmación es una consecuencia inmediata de (18.1.7).

\oldpage[0\textsubscript{IV-1}]{58}

Se dice que la $A$-extensión $H$ de $B$ por $L'$ es \emph{deducida de $E$ por
medio del homomorfismo} $w:L\to L'$. El carácter funtorial de la suma
amalgamada en cada uno de los sumandos muestra además que, si se tiene un morfismo
de extensiones
\[
  \xymatrix{
    0\ar[r] & L\ar[r]\ar[d]^{1_L}\ar@{}[d]|{\wr}
      & E_1\ar[r]\ar[d]^g & B_1\ar[r]\ar[d]^h & 0\\
    0\ar[r] & L\ar[r] & E_2\ar[r] & B_2\ar[r] & 0
  }
\]
se deduce canónicamente un morfismo de extensiones
\[
  \xymatrix{
    0\ar[r] & L'\ar[r]\ar[d]^{1_{L'}}\ar@{}[d]|{\wr}
      & E_1\oplus_LL'\ar[r]\ar[d]^{g\oplus_L1_{L'}}
      & B_1\ar[r]\ar[d]^h & 0\\
    0\ar[r] & L'\ar[r] & E_2\oplus_LL'\ar[r] & B_2\ar[r] & 0.
  }
\]
En particular, si $E_1$ y $E_2$ son $A$-extensiones $A$-equivalentes de
$B$ por $L$, las extensiones de $B$ por $L'$ que se deducen de ellas mediante $w$
son $A$-equivalentes.

Cuando se tiene un morfismo (18.2.4.1) de $A$-extensiones, este se \emph{factoriza
a través de} la $A$-extensión $H$ de $B$ por $L'$ deducida de $E$ por medio
del homomorfismo $w:L\to L'$ (considerándose $L'$ como $B$-bimódulo mediante
el homomorfismo $v:B\to B'$): en efecto, la definición de la suma amalgamada
muestra que existe un $A$-homomorfismo único $u_0:H\to E'$ de
$A$-bimódulos que hace conmutativo el diagrama
\[
  \xymatrix{
    0\ar[r] & L\ar[r]^j\ar[d]^{w_0} & E\ar[r]^f\ar[d]^{j_1}
      & B\ar[r]\ar[d]^{1_B}\ar@{}[d]|{\wr} & 0\\
    0\ar[r] & L'_0\ar[r]^{j_2}\ar[d]^{j_0}\ar@{}[d]|{\wr}
      & H\ar[r]^{f_0}\ar[d]^{u_0} & B\ar[r]\ar[d]^v & 0\\
    0\ar[r] & L'\ar[r]_{j'} & E'\ar[r]_{f'} & B'\ar[r] & 0
  }
\]
con $u=u_0\circ j_1$, $w=j_0\circ w_0$, siendo $j_1$ y $j_2$ los
homomorfismos canónicos; se comprueba inmediatamente que $u_0$ es también un
homomorfismo de anillos.

Señalemos, por último, las propiedades funtoriales relativas a las extensiones
triviales:
\end{env}

\begin{proposition}[18.2.9]
\label{0.18.2.9-es}
Sean $B$, $B'$ dos $A$-anillos, $L$ un $B$-bimódulo, $L'$ un
$B'$-bimódulo, $v:B\to B'$ un $A$-homomorfismo de anillos, $w:L\to L'$ un
homomorfismo de $A$-bimódulos tal que $(v,w)$ sea un dihomomorfismo de
bimódulos. Entonces existe un $A$-homomorfismo único de anillos
$u:D_B(L)\to D_{B'}(L')$ que hace conmutativos los diagramas
\[
  \xymatrix@C=2.2pc@R=1.6pc{
    D_B(L)\ar[r]^u & D_{B'}(L') & & D_B(L)\ar[r]^u & D_{B'}(L')\\
    B\ar[u]\ar[r]_v & B'\ar[u] & & L\ar[u]\ar[r]_w & L'\ar[u]
  }
\]
donde las flechas verticales son las inyecciones canónicas.
\end{proposition}

\oldpage[0\textsubscript{IV-1}]{59}

\begin{proof}
En efecto, $u$ no puede ser sino la aplicación
$(x,s)\mapsto(v(x),w(s))$, y queda comprobar que es un
$A$-homomorfismo de anillos, lo que resulta trivialmente de la definición
(18.2.3). Se observará que $u$ hace también conmutativo el diagrama
\[
  \xymatrix@C=2.2pc@R=1.5pc{
    D_B(L)\ar[r]^u\ar[d] & D_{B'}(L')\ar[d]\\
    B\ar[r]_v & B'
  }
\]
donde esta vez las flechas verticales son las aumentaciones.
\end{proof}

\begin{proposition}[18.2.10]
\label{0.18.2.10-es}
Sean $B$ un $A$-anillo, $L$ un $B$-bimódulo y $E$ una $A$-extensión de $B$
por $L$. Se define una aplicación biyectiva del conjunto $G$ de los
homomorfismos de $B$-anillos de $B$ en $E$ (es decir, el conjunto de los
$A$-homomorfismos inversos por la derecha de la aumentación $E\to B$) sobre el
conjunto $G'$ de las $A$-equivalencias de $D_B(L)$ sobre $E$, haciendo
corresponder a todo $g\in G$ la $A$-equivalencia
$g':(x,s)\mapsto g(x)+s$; la aplicación inversa hace corresponder a todo
$g'\in G'$ el $B$-homomorfismo $x\mapsto g'(x,0)$.
\end{proposition}

\begin{proof}
Esto resulta inmediatamente de las definiciones.
\end{proof}

\subsection{El grupo de las clases de $A$-extensiones}
\label{subsection:0.18.3-es}

\begin{env}[18.3.1]
\label{0.18.3.1-es}
Consideremos un $A$-anillo $B$ y un $B$-bimódulo $L$ fijados; entonces la
relación «$E$ y $E'$ son $A$-equivalentes» entre las $A$-extensiones $E$, $E'$
de $B$ por $L$ es una relación de equivalencia, y para esta relación se puede
hablar del \emph{conjunto} de las clases de $A$-extensiones $A$-equivalentes de
$B$ por $L$. Para verlo, basta observar que, si para todo $x\in B$, $c_x$ es
un elemento de $E$ cuya imagen en $B$ es $x$, todo $z\in E$ se escribe de una
manera y solo una en la forma $c_x+t$, donde $t\in L$, y se puede escribir
\[
  c_{x+y}=c_x+c_y+\varphi(x,y),\qquad
  c_{xy}=c_xc_y+\psi(x,y),
\]
donde $\varphi(x,y)$ y $\psi(x,y)$ son elementos de $L$, y las aplicaciones
$\varphi$ y $\psi$ de $B\times B$ en $L$ deben verificar condiciones que
expresan que $E$ es un $A$-anillo y que no es necesario escribir aquí. Toda
$A$-extensión de $B$ por $L$ es, por tanto, $A$-equivalente a una
$A$-extensión cuyo conjunto subyacente es $B\times L$, de donde se obtiene
inmediatamente nuestra conclusión.
\end{env}

\begin{env}[18.3.2]
\label{0.18.3.2-es}
Fijado el $A$-anillo $B$, designemos provisionalmente por $T(L)$ el conjunto
de las clases de $A$-extensiones de $B$ por $L$. Para todo $B$-homomorfismo
$w:L\to L'$ de $B$-bimódulos, se define canónicamente una aplicación
$T(w):T(L)\to T(L')$, haciendo corresponder a la clase de una $A$-extensión
$E$ de $B$ por $L$ la clase de la $A$-extensión $E\oplus_LL'$ deducida de $E$
por $w$, en virtud de (18.2.8). Si $w':L'\to L''$ es un segundo homomorfismo
de $B$-bimódulos, se tiene además
\begin{equation}
\label{0.18.3.2.1-es}
\tag{18.3.2.1}
  T(w'\circ w)=T(w')\circ T(w).
\end{equation}
En efecto, se sabe que existe un isomorfismo canónico de $A$-bimódulos
\[
  E\oplus_LL''\xrightarrow{\sim}(E\oplus_LL')\oplus_{L'}L''
\]
en virtud de las propiedades generales de los límites inductivos (cf., por
ejemplo, ($\mathrm{I}$, 3.3.9)), y es inmediato comprobar que se trata
efectivamente de una $A$-equivalencia de $A$-extensiones, de donde se obtiene

\oldpage[0\textsubscript{IV-1}]{60}

la relación (18.3.2.1). Si $\mathbf C_B$ designa la categoría de los
$B$-bimódulos, se ve que se ha definido así un \emph{funtor covariante}
$T:\mathbf C_B\to\mathbf{Ens}$.
\end{env}

\begin{env}[18.3.3]
\label{0.18.3.3-es}
Consideremos ahora una familia $(L_\alpha)_{\alpha\in I}$ de $B$-bimódulos y
su producto $L=\prod_{\alpha\in I}L_\alpha$; las proyecciones
$\operatorname{pr}_\alpha:L\to L_\alpha$ definen aplicaciones
\[
  T(\operatorname{pr}_\alpha):T(L)\to T(L_\alpha),
\]
de donde se obtiene una aplicación canónica
\begin{equation}
\label{0.18.3.3.1-es}
\tag{18.3.3.1}
  \prod_\alpha T(\operatorname{pr}_\alpha):T(L)
  \longrightarrow\prod_{\alpha\in I}T(L_\alpha).
\end{equation}
Vamos a ver que esta aplicación es \emph{biyectiva}. En efecto, para todo
$\alpha\in I$, sea $E_\alpha$ una $A$-extensión de $B$ por $L_\alpha$, y sea
$F$ el \emph{producto fibrado} de los $E_\alpha$ sobre $B$ (18.1.3); es
inmediato que $F$ es una $A$-extensión de $B$ por
$L=\prod_{\alpha\in I}L_\alpha$ y que, si se reemplaza cada $E_\alpha$ por una
$A$-extensión $A$-equivalente $E'_\alpha$, el producto fibrado $F'$ de los
$E'_\alpha$ sobre $B$ es $A$-equivalente a $F$. Se ha definido así una
aplicación $\prod_{\alpha\in I}T(L_\alpha)\to T(L)$, y es claro que esta
aplicación es la inversa de (18.3.3.1), de donde se sigue nuestra afirmación.
Se identificará canónicamente $T(L)$ con
$\prod_{\alpha\in I}T(L_\alpha)$ mediante la aplicación (18.3.3.1). Se
comprueba además inmediatamente que, si $(L'_\alpha)_{\alpha\in I}$ es una
segunda familia de $B$-bimódulos y, para cada $\alpha$,
$w_\alpha:L_\alpha\to L'_\alpha$ es un homomorfismo de $B$-bimódulos,
entonces, poniendo
$w=\prod_\alpha w_\alpha:L\to L'=\prod_{\alpha\in I}L'_\alpha$, $T(w)$ se
identifica con $\prod_{\alpha\in I}T(w_\alpha)$ cuando se efectúa la
identificación precedente.
\end{env}

\begin{env}[18.3.4]
\label{0.18.3.4-es}
Sentado esto, para un $B$-bimódulo $L$, la adición es un homomorfismo
$s:L\times L\to L$ de $B$-bimódulos, y lo mismo sucede con la simetría
$t:L\to L$ de la ley aditiva de $L$. Se deduce de ello una ley de composición
\[
  T(s):T(L)\times T(L)\longrightarrow T(L)
\]
en $T(L)$, en virtud de (18.3.3), y esta ley es una ley de \emph{grupo
conmutativo} cuya simetría es $T(t)$, como resulta de la definición de un
objeto en grupos mediante diagramas conmutativos
($\mathbf 0_{\mathrm{III}}$, 8.2.5 y 8.2.6). Designaremos por
$\operatorname{Exan}_A(B,L)$ el grupo conmutativo así definido y diremos que
es el \emph{grupo de las clases de $A$-extensiones de $B$ por $L$}.
\end{env}

\begin{env}[18.3.5]
\label{0.18.3.5-es}
Designemos por $\mathbf K$ la categoría cuyos objetos son las ternas
$(A,B,L)$, donde $A$ es un anillo, $B$ un $A$-anillo y $L$ un $B$-bimódulo;
los \emph{morfismos} de esta categoría son las ternas $(u,v,w)$, donde
$u:A'\to A$ y $v:B'\to B$ son dos homomorfismos de anillos que hacen
conmutativo el cuadrado izquierdo del diagrama
\[
  \xymatrix@C=2pc@R=1.5pc{
    A\ar[r] & B & L\ar[d]^w\\
    A'\ar[u]^u\ar[r] & B'\ar[u]^v & L'
  }
\]
donde las flechas horizontales son los homomorfismos estructurales; por último,
$w:L\to L'$ es un homomorfismo de grupos conmutativos tal que se tiene
$w(v(b')z)=b'w(z)$ y

\oldpage[0\textsubscript{IV-1}]{61}

$w(zv(b'))=w(z)b'$ cualesquiera que sean $z\in L$ y $b'\in B'$ (en otros
términos, $w$ es un homomorfismo de $B'$-bimódulos cuando se dota a $L$ de la
estructura de $B'$-bimódulo definida por $v$). La composición de los morfismos
se define por
$(u',v',w')\circ(u,v,w)=(u\circ u',v\circ v',w'\circ w)$, lo que se justifica
inmediatamente. Nos proponemos demostrar que
\begin{equation}
\label{0.18.3.5.1-es}
\tag{18.3.5.1}
  (A,B,L)\longmapsto\operatorname{Exan}_A(B,L)
\end{equation}
es un \emph{funtor covariante} de la categoría $\mathbf K$ en la categoría
$\mathbf{Ab}$ de los grupos conmutativos. Se trata, pues, para toda terna
$(u,v,w)$ como la anterior, de definir un homomorfismo de grupos conmutativos
\[
  (u,v,w)_*:\operatorname{Exan}_A(B,L)
  \longrightarrow\operatorname{Exan}_{A'}(B',L').
\]
En virtud de la definición de los morfismos en $\mathbf K$, se puede escribir
\[
  (u,v,w)=(1_{A'},1_{B'},w)\circ(1_{A'},v,1_L)
  \circ(u,1_B,1_L)
\]
donde, en el primer factor, $L$ está dotado de su estructura de $B'$-bimódulo
definida por $v$; definiremos, por tanto, primero $(u,v,w)_*$ cuando dos de los
homomorfismos $u$, $v$, $w$ sean iguales a la identidad.
\end{env}

\begin{env}[18.3.6]
\label{0.18.3.6-es}
Definiremos primero $(1_A,1_B,w)_*$ como la aplicación
\begin{equation}
\label{0.18.3.6.1-es}
\tag{18.3.6.1}
  w_*:\operatorname{Exan}_A(B,L)\longrightarrow
  \operatorname{Exan}_A(B,L')
\end{equation}
denotada por $T(w)$ en (18.3.2); es inmediato comprobar que se trata de un
homomorfismo de grupos, pues esta propiedad se expresa mediante la
conmutatividad de diagramas que son las imágenes por $T$ de diagramas análogos
para $L$ y $L'$.

La aplicación $(1_A,v,1_L)_*$ es la aplicación
\begin{equation}
\label{0.18.3.6.2-es}
\tag{18.3.6.2}
  v^*:\operatorname{Exan}_A(B,L)\longrightarrow
  \operatorname{Exan}_A(B',L)
\end{equation}
definida del modo siguiente: si $E$ es una $A$-extensión de $B$ por $L$, se ha
visto que $E\times_BB'$ es una $A$-extensión de $B'$ por $L$ (18.2.5) y que,
si se reemplaza $E$ por una $A$-extensión $A$-equivalente $E'$,
$E'\times_BB'$ es $A$-equivalente a $E\times_BB'$; la imagen por $v^*$ de la
clase de $E$ es la clase de $E\times_BB'$. Se comprueba inmediatamente que, si
$w:L\to L'$ es un homomorfismo de $B$-bimódulos, el diagrama
\begin{equation}
\label{0.18.3.6.3-es}
\tag{18.3.6.3}
  \xymatrix@C=3pc@R=2.4pc{
    \operatorname{Exan}_A(B,L)\ar[r]^{v^*}\ar[d]_{w_*}
      & \operatorname{Exan}_A(B',L)\ar[d]^{w_*}\\
    \operatorname{Exan}_A(B,L')\ar[r]_{v^*}
      & \operatorname{Exan}_A(B',L')
  }
\end{equation}
es conmutativo, considerándose $L$ y $L'$ como $B'$-bimódulos por medio de
$v$ en la columna derecha. Reemplazando $L$ y $L'$, respectivamente, por
$L\times L$ y $L$, y $w$ por la adición $s$ en $L$, se concluye que $v^*$ es
efectivamente un \emph{homomorfismo de grupos}.
\end{env}

\oldpage[0\textsubscript{IV-1}]{62}

\begin{env}[18.3.6.4]
\label{0.18.3.6.4-es}
Finalmente, la aplicación $(u,1_B,1_L)_*$ es la aplicación
\[
  u^*:\operatorname{Exan}_A(B,L)\longrightarrow
  \operatorname{Exan}_{A'}(B,L)
\]
que se obtiene haciendo corresponder a una $A$-extensión $E$ de $B$ por $L$
el anillo $E$ considerado como $A'$-\emph{anillo} por medio de $u$ (18.1.1), que
es evidentemente una $A'$-extensión de $B$ por $L$, considerándose también $B$
como $A'$-anillo por medio de $u$~; es claro que una $A$-equivalencia es
también una $A'$-equivalencia, de donde resulta la aplicación (18.3.6.4), que, para todo
homomorfismo $w:L\to L'$ de $B$-bimódulos, hace también conmutativo el diagrama
\begin{equation}
\label{0.18.3.6.5-es}
\tag{18.3.6.5}
  \xymatrix@C=3pc@R=2.4pc{
    \operatorname{Exan}_A(B,L)\ar[r]^{u^*}\ar[d]_{w_*}
      & \operatorname{Exan}_{A'}(B,L)\ar[d]^{w_*}\\
    \operatorname{Exan}_A(B,L')\ar[r]_{u^*}
      & \operatorname{Exan}_{A'}(B,L')
  }
\end{equation}
de donde se concluye, como antes, que $u^*$ es un homomorfismo de grupos.

Dicho esto, se pone
$(u,v,w)_*=(1_{A'},1_{B'},w)_*\circ(1_{A'},v,1_L)_*
\circ(u,1_B,1_L)_*$ y se verifica fácilmente, en razón de la conmutatividad
de los diagramas (18.3.6.3) y (18.3.6.5), que se tiene
\[
  (u\circ u',v\circ v',w'\circ w)_*
  =(u',v',w')_*\circ(u,v,w)_*,
\]
lo que termina de demostrar que (18.3.4.1) es un funtor.

La existencia del homomorfismo de grupos (18.3.6.1) muestra, en particular,
que si $E$ es una $A$-extensión \emph{trivial} de $B$ por $L$, la extensión
$E\oplus_LL'$ de $B$ por $L'$ definida en (18.2.8) es también trivial, lo que
se verifica, por lo demás, sin dificultad de manera directa.

Por otra parte, para todo elemento $z$ del \emph{centro} $Z$ de $B$, la homotecia
$h_z:y\mapsto yz=zy$ es un endomorfismo del $B$-bimódulo $L$, por lo que $(h_z)_*$
es un endomorfismo del grupo conmutativo $\operatorname{Exan}_A(B,L)$, y por
funtorialidad, estos endomorfismos definen sobre
$\operatorname{Exan}_A(B,L)$ una estructura canónica de $Z$-\emph{módulo}.
\end{env}

\begin{env}[18.3.7]
\label{0.18.3.7-es}
Sean $A$, $A'$ dos anillos, $u:A'\to A$ un homomorfismo, $B$ un
$A$-anillo y $L$ un $B$-bimódulo. El \emph{núcleo} del homomorfismo de
grupos
\[
  u^*:\operatorname{Exan}_A(B,L)\longrightarrow
  \operatorname{Exan}_{A'}(B,L)
\]
está formado, por definición, por las clases de $A$-extensiones de $B$ por $L$ que
son $A'$-\emph{triviales} cuando se consideran como $A'$-extensiones por
medio de $u$. Este núcleo se denota por
$\operatorname{Exan}_{A/A'}(B,L)$ cuando ello no da lugar a confusión.

Si $\Lambda$ es un anillo, y si se tiene un diagrama conmutativo de
$\Lambda$-\emph{homomorfismos} de $\Lambda$-anillos
\begin{equation}
\label{0.18.3.7.1-es}
\tag{18.3.7.1}
  \xymatrix@C=2pc@R=1.5pc{
    B'\ar[r] & B\\
    A'\ar[u]\ar[r] & A\ar[u]
  }
\end{equation}

\oldpage[0\textsubscript{IV-1}]{63}

se deducen canónicamente homomorfismos
\begin{equation}
\label{0.18.3.7.2-es}
\tag{18.3.7.2}
  \operatorname{Exan}_{A/\Lambda}(B,L)\longrightarrow
  \operatorname{Exan}_{A'/\Lambda}(B,L)\longrightarrow
  \operatorname{Exan}_{A'/\Lambda}(B',L)
\end{equation}
que proceden de la conmutatividad del diagrama
\[
  \xymatrix@C=2.8pc@R=2.3pc{
    \operatorname{Exan}_A(B,L)\ar[r]\ar[d]
      & \operatorname{Exan}_{A'}(B,L)\ar[r]\ar[d]
      & \operatorname{Exan}_{A'}(B',L)\ar[d]\\
    \operatorname{Exan}_{\Lambda}(B,L)\ar[r]
      & \operatorname{Exan}_{\Lambda}(B,L)\ar[r]
      & \operatorname{Exan}_{\Lambda}(B',L)
  }
\]
donde las flechas se deducen de las de (18.3.7.1) por funtorialidad.
\end{env}

\begin{proposition}[18.3.8]
\label{0.18.3.8-es}
Sean $B$ un anillo, $\mathfrak J$ un ideal bilateral de $B$ y
$C=B/\mathfrak J$ el anillo cociente~; $\mathfrak J/\mathfrak J^2$ está entonces
canónicamente provisto de una estructura de $C$-bimódulo. Para todo $C$-bimódulo
$L$, sea $\operatorname{Hom}_C(\mathfrak J/\mathfrak J^2,L)$ el grupo
aditivo de los homomorfismos de $C$-bimódulos de $\mathfrak J/\mathfrak J^2$
en $L$. Se define entonces un isomorfismo canónico de grupos conmutativos
\begin{equation}
\label{0.18.3.8.1-es}
\tag{18.3.8.1}
  \eta_L:\operatorname{Hom}_C(\mathfrak J/\mathfrak J^2,L)
  \xrightarrow{\sim}\operatorname{Exan}_B(C,L)
\end{equation}
haciendo corresponder a todo $C$-homomorfismo
$w:\mathfrak J/\mathfrak J^2\to L$ (que es \emph{a fortiori} un
$B$-homomorfismo) la clase de la extensión
$(B/\mathfrak J^2)\oplus_{(\mathfrak J/\mathfrak J^2)}L$ deducida de
la extensión $B/\mathfrak J^2$ de $C$ por $\mathfrak J/\mathfrak J^2$ por medio
de $w$~; el isomorfismo inverso hace corresponder a la clase de una
$B$-extensión $E$ de $C$ por $L$ el homomorfismo $w$ tal que el compuesto
$\mathfrak J\to\mathfrak J/\mathfrak J^2\xrightarrow{w}L$ sea la restricción
a $\mathfrak J$ del homomorfismo estructural $B\to E$.
\end{proposition}

\begin{proof}
Sea $F=B/\mathfrak J^2$, que es una $B$-extensión de $C$ por
$\mathfrak J/\mathfrak J^2$. Para toda $B$-extensión
\[
  0\longrightarrow L\xrightarrow{j}E\xrightarrow{p}C\longrightarrow0
\]
de $C$ por $L$, el homomorfismo estructural $f:B\to E$ es tal que el compuesto
$B\xrightarrow{f}E\xrightarrow{p}C$ es el homomorfismo canónico
$B\to B/\mathfrak J=C$. Como la imagen de $\mathfrak J$ por $p\circ f$ es
nula, $f(\mathfrak J)$ está contenida en el núcleo de $p$, es decir,
$j(L)$, y como $j(L)$ es de cuadrado nulo, se tiene $f(\mathfrak J^2)=0$~; por tanto,
$f$ se factoriza como
$B\to B/\mathfrak J^2=F\xrightarrow{u}E$, y si $w$ es la restricción de
$u$ a $\mathfrak J/\mathfrak J^2$, se tiene un diagrama conmutativo
\[
  \xymatrix@C=2pc@R=1.5pc{
    0\ar[r] & \mathfrak J/\mathfrak J^2\ar[r]\ar[d]^w
      & F\ar[r]^g\ar[d]^u
      & C\ar[r]\ar[d]^{1_C}\ar@{}[d]|{\wr} & 0\\
    0\ar[r] & L\ar[r]_j & E\ar[r]_p & C\ar[r] & 0
  }
\]

dicho de otro modo, $(u,1_C,w)$ es un morfismo de extensiones (18.2.4). De él se
deduce un morfismo de extensiones
\[
  \xymatrix@C=2pc@R=1.5pc{
    0\ar[r] & L\ar[r]\ar[d]^{1_L}\ar@{}[d]|{\wr}
      & E'\ar[r]\ar[d]^{u'}
      & C\ar[r]\ar[d]^{1_C}\ar@{}[d]|{\wr} & 0\\
    0\ar[r] & L\ar[r]_j & E\ar[r]_p & C\ar[r] & 0
  }
\]

\oldpage[0\textsubscript{IV-1}]{64}

donde $E'=F\oplus_{(\mathfrak J/\mathfrak J^2)}L$ (18.2.8)~; se trata de
demostrar que $u'$ es una $B$-\emph{equivalencia}, lo que establecerá que $\eta_L$
es biyectivo. En virtud de la construcción realizada en (18.2.8), todo se reduce a
demostrar (habida cuenta de (18.1.7)) que la imagen inversa $G$ de la extensión
$E$ de $C$ por el homomorfismo canónico $g:F\to C$ es $F$-trivial, lo que
es evidente, puesto que $g$ se factoriza como
$F\xrightarrow{u}E\xrightarrow{p}C$ (18.1.6).

Queda ver que $\eta_L$ es un homomorfismo de \emph{grupos}~; ahora bien, para
todo $B$-homomorfismo $h:L\to L'$, es inmediato que el diagrama
\[
  \xymatrix@C=3pc@R=2.4pc{
    \operatorname{Hom}_C(\mathfrak J/\mathfrak J^2,L)
      \ar[r]^{\eta_L}\ar[d]_{\operatorname{Hom}(1,h)}
      & \operatorname{Exan}_B(C,L)\ar[d]^{h_*}\\
    \operatorname{Hom}_C(\mathfrak J/\mathfrak J^2,L')
      \ar[r]_{\eta_{L'}}^{\sim}
      & \operatorname{Exan}_B(C,L')
  }
\]
es conmutativo. Basta aplicar esta observación al homomorfismo
$L\times L\to L$ que define la adición para concluir.
\end{proof}

\subsection{Extensiones de álgebras}
\label{subsection:0.18.4-es}

\begin{env}[18.4.1]
\label{0.18.4.1-es}
Sea $A$ un anillo \emph{conmutativo}. La categoría de las $A$-\emph{álgebras} es
entonces una subcategoría plena de la de los $A$-anillos, caracterizada por el
hecho de que los homomorfismos estructurales $\rho:A\to B$ son \emph{centrales}.

Si $B$ es una $A$-álgebra y $L$ un $B$-bimódulo, es inmediato que la
$A$-extensión trivial tipo $D_B(L)$ es una $A$-álgebra. Del mismo modo, en la
construcción de (18.2.5) (resp. de (18.2.8)), si $B$, $B'$ y $E'$ (resp.
$B$ y $E$) son $A$-álgebras, ocurre lo mismo con
$E'\times_{B'}B$ (resp. con $E\oplus_LL'$). Finalmente, es claro que si $B$ es
una $A$-álgebra y $E$ una $A$-extensión de $B$ por un $B$-bimódulo $L$ que
es una $A$-álgebra, toda $A$-extensión de $B$ por $L$, $A$-\emph{equivalente}
a $E$, es también una $A$-álgebra. Se deduce entonces inmediatamente de la definición
(18.3.4) que las clases de $A$-extensiones equivalentes de $B$ por $L$ que
son $A$-\emph{álgebras} forman un \emph{subgrupo}, denotado por
$\operatorname{Exal}_A(B,L)$, de $\operatorname{Exan}_A(B,L)$. Sea
$\mathbf K'$ la subcategoría plena de la categoría $\mathbf K$ definida en
(18.3.5), cuyos objetos $(A,B,L)$ son tales que $A$ es conmutativo y $B$
una $A$-álgebra. Entonces lo que precede muestra que
\[
  (A,B,L)\longmapsto\operatorname{Exal}_A(B,L)
\]
es un \emph{funtor covariante} de $\mathbf K'$ en $\mathbf{Ab}$. Los
resultados de (18.3.7) no cambian cuando se reemplaza en todas partes
$\operatorname{Exan}$ por $\operatorname{Exal}$.
\end{env}

\begin{env}[18.4.2]
\label{0.18.4.2-es}
Supongamos todavía que el anillo $A$ es conmutativo. Las observaciones de (18.4.1) siguen
siendo válidas cuando se reemplaza «~álgebra~» por «~álgebra conmutativa~» y
«~bimódulo~» por «~módulo~». Si $B$ es una $A$-álgebra conmutativa y $L$ un
$B$-módulo,

\oldpage[0\textsubscript{IV-1}]{65}

las clases de $A$-extensiones equivalentes de $B$ por $L$ que son
$A$-\emph{álgebras conmutativas} (o, lo que equivale a lo mismo,
$A$-\emph{anillos conmutativos}) forman un \emph{subgrupo} de
$\operatorname{Exal}_A(B,L)$, denotado por $\operatorname{Exalcom}_A(B,L)$. Si
$\mathbf K''$ es la subcategoría plena de $\mathbf K'$ formada por los triples
$(A,B,L)$ donde $A$ y $B$ son conmutativos y $L$ un $B$-módulo,
\[
  (A,B,L)\longmapsto\operatorname{Exalcom}_A(B,L)
\]
es también un \emph{funtor covariante} de $\mathbf K''$ en $\mathbf{Ab}$. También se
puede, en (18.3.7), reemplazar en todas partes $\operatorname{Exan}$ por
$\operatorname{Exalcom}$. Finalmente, si en (18.3.8) se supone que $B$ es un
anillo \emph{conmutativo} y que $L$ es un $C$-\emph{módulo}, el mismo
razonamiento da un isomorfismo canónico
\begin{equation}
\label{0.18.4.2.1-es}
\tag{18.4.2.1}
  \operatorname{Hom}_C(\mathfrak J/\mathfrak J^2,L)
  \xrightarrow{\sim}\operatorname{Exalcom}_B(C,L)
\end{equation}
donde el primer miembro es el grupo de los homomorfismos de $C$-\emph{módulos}.
\end{env}

\begin{env}[18.4.3]
\label{0.18.4.3-es}
Sea $A$ un anillo conmutativo. Un caso importante de extensiones de
$A$-álgebras está formado por las $A$-álgebras $E$, extensiones de una $A$-álgebra $B$
por un $B$-bimódulo $L$ tales que $E$, como $A$-\emph{módulo}, es una
extensión \emph{trivial} del $A$-módulo $B$ por el $A$-módulo $L$~; dicho de otro
modo, la sucesión exacta de $A$-\emph{módulos}
$0\to L\to E\to B\to0$ es \emph{escindida}. Se dice entonces que $E$ es una
\emph{extensión de Hochschild} de $B$ por $L$. Este será siempre el caso cuando
$B$ sea un $A$-módulo \emph{proyectivo}, y en particular cuando $A$ sea un
\emph{cuerpo} conmutativo. Como $A$-módulo, se puede identificar $E$ con
$B\times L$, estando dada la multiplicación en $E$ por
$(x,s)(y,t)=(xy,xt+sy+f(x,y))$ con $f(x,y)\in L$. Si se expresa que esta
multiplicación define en $B\times L$ una estructura de $A$-álgebra, se obtiene
(M, XIV, 2) que $f$ debe ser una aplicación $A$-\emph{bilineal} de
$B\times B$ en $L$, tal que
\begin{equation}
\label{0.18.4.3.1-es}
\tag{18.4.3.1}
  f(xy,z)+f(x,y)z=xf(y,z)+f(x,yz),
\end{equation}
es decir, $f$ es un $2$-\emph{cociclo} sobre $B$ con valores en $L$, en el
sentido de Hochschild; para que la extensión $E$ sea $A$-trivial, es necesario
y suficiente que se tenga
\begin{equation}
\label{0.18.4.3.2-es}
\tag{18.4.3.2}
  f(x,y)=xg(y)-g(xy)+g(x)y,
\end{equation}
donde $g$ es una aplicación $A$-\emph{lineal} de $B$ en $L$; es decir, $f$
debe ser un $2$-\emph{coborde} en el sentido de Hochschild. Se deduce
inmediatamente que las clases de extensiones de Hochschild de $B$ por $L$
forman un subgrupo de $\operatorname{Exal}_A(B,L)$, isomorfo al \emph{grupo de
cohomología} de Hochschild $\mathrm H_A^2(B,L)$.

Si $B$ es una $A$-álgebra \emph{conmutativa}, y $L$ un $B$-módulo, las
extensiones de Hochschild \emph{conmutativas} de $B$ por $L$ corresponden a los
$2$-cociclos \emph{simétricos}, es decir, tales que $f(y,x)=f(x,y)$. Las clases
de extensiones de Hochschild conmutativas de $B$ por $L$ forman, por tanto, un
subgrupo de $\operatorname{Exalcom}_A(B,L)$, isomorfo al subgrupo de
$\mathrm H_A^2(B,L)$ que es imagen del grupo de los $2$-cociclos simétricos y
que denotaremos por $\mathrm H_A^2(B,L)^s$.
\end{env}

\begin{env}[18.4.4]
\label{0.18.4.4-es}
Limitémonos al caso en que $A$ y $B$ son \emph{conmutativos} y $L$ es un
$B$-módulo, y recordemos en este caso la definición equivalente de los grupos
de Hochschild $\mathrm H_A^i(B,L)$ (M, IX, 6). Se considera un complejo
$\mathbf P_\bullet=(\mathbf P_n)_{n\geqslant0}$ de $B$-módulos, donde
$\mathbf P_n=B^{\otimes(n+1)}=B\otimes_A B\otimes_A\cdots\otimes_A B$
($n+1$ veces), estando definida la estructura de $B$-módulo

\oldpage[0\textsubscript{IV-1}]{66}

por
\[
  x(y_1\otimes y_2\otimes\cdots\otimes y_{n+1})
  =(xy_1)\otimes y_2\otimes\cdots\otimes y_{n+1}\,;
\]
la derivación, de grado $-1$, $d_n:\mathbf P_n\to\mathbf P_{n-1}$, está
definida por
\begin{align*}
  d_n(x_1\otimes x_2\otimes\cdots\otimes x_{n+1})
    &=(x_1x_2)\otimes x_3\otimes\cdots\otimes x_{n+1}
      -x_1\otimes(x_2x_3)\otimes\cdots\otimes x_{n+1}+\cdots\\
    &\quad+(-1)^{n-1}x_1\otimes x_2\otimes\cdots\otimes x_{n-1}
      \otimes(x_nx_{n+1})
      +(-1)^n(x_{n+1}x_1)\otimes x_2\otimes\cdots\otimes x_n,
\end{align*}
que es efectivamente $B$-lineal puesto que $B$ es conmutativo.

Un $2$-cociclo de este complejo con valores en $L$ es una aplicación
$B$-lineal $h$ de $B\otimes B\otimes B$ en $L$ tal que
$h(d_3(x\otimes y\otimes z\otimes t))=0$; pero como
$h(x\otimes y\otimes z)=xh(1\otimes y\otimes z)$, el cociclo $h$ está
determinado por la aplicación $A$-bilineal
$(y,z)\longmapsto f(y,z)=h(1\otimes y\otimes z)$ de $B\times B$ en $L$, y,
escribiendo para $h$ la condición precedente, se recupera para $f$ la
condición (18.4.3.1). Del mismo modo, un $2$-coborde será una aplicación de la
forma $x\otimes y\otimes z\longmapsto h'(d_2(x\otimes y\otimes z))$, donde
$h':B\otimes_A B\to L$ es lineal, y aquí también $h'$ está determinado por la
aplicación $A$-lineal $g:y\longmapsto h'(1\otimes y)$ de $B$ en $L$; se
obtiene entonces
$h'(d_2(1\otimes x\otimes y))=xg(y)-g(xy)+yg(x)$, lo que vuelve a dar
(18.4.3.2). Se procede del mismo modo para todo $i$, y se ve así que se tiene
\begin{equation}
\label{0.18.4.4.1-es}
\tag{18.4.4.1}
  \mathrm H_A^i(B,L)=\mathrm H^i(\operatorname{Hom}_B(\mathbf P_\bullet,L)).
\end{equation}
\end{env}

\begin{env}[18.4.5]
\label{0.18.4.5-es}
Bajo las condiciones de (18.4.4), se puede interpretar del mismo modo el
grupo $\mathrm H_A^2(B,L)^s$ (18.4.3). Para ello, modifiquemos el complejo
$\mathbf P_\bullet$ en grado $3$, considerando un nuevo complejo
\[
  \mathbf P'_\bullet:\mathbf P'_3\xrightarrow{d'_3}\mathbf P_2
  \xrightarrow{d_2}\mathbf P_1\,;
\]
tomaremos
$\mathbf P'_3=\mathbf P_3\oplus(B\otimes_A B\otimes_A B)$, coincidiendo
$d'_3$ con $d_3$ sobre $\mathbf P_3$, y estando dado en
$B\otimes_A B\otimes_A B$ por
\[
  d'_3(x\otimes y\otimes z)=x\otimes y\otimes z-x\otimes z\otimes y.
\]

La relación $d_2d'_3=0$ resulta de la conmutatividad de $B$. Con las notaciones
introducidas anteriormente, un $2$-cociclo de $\mathbf P'_\bullet$ corresponde
ahora a una aplicación $A$-bilineal $f:B\times B\to L$ que es
\emph{simétrica} y verifica (18.4.3.1); por consiguiente, se tiene
\[
  \mathrm H_A^2(B,L)^s
  =\mathrm H^2(\operatorname{Hom}_B(\mathbf P'_\bullet,L)).
\]
\end{env}

\begin{env}[18.4.6]
\label{0.18.4.6-es}
En el caso particular en que se considera un \emph{cuerpo} conmutativo $k$ y
una \emph{extensión} $K$ de $k$, se tiene
$\operatorname{Ext}_K^1(M,L)=0$ cualesquiera que sean los $K$-espacios
vectoriales $L$, $M$, y, por consiguiente (M, VI, 3.3 $a$)), se tiene un
isomorfismo canónico
\begin{equation}
\label{0.18.4.6.1-es}
\tag{18.4.6.1}
  \mathrm H_k^i(K,L)\xrightarrow{\sim}
  \operatorname{Hom}_K(\mathrm H_i(\mathbf P_\bullet),L)
\end{equation}
y, del mismo modo,
\begin{equation}
\label{0.18.4.6.2-es}
\tag{18.4.6.2}
  \mathrm H_k^2(K,L)^s\xrightarrow{\sim}
  \operatorname{Hom}_K(\mathrm H_2(\mathbf P'_\bullet),L).
\end{equation}
\end{env}

\subsection{Caso de los anillos topológicos}
\label{subsection:0.18.5-es}

\begin{env}[18.5.1]
\label{0.18.5.1-es}
Sean $A$, $B$ dos anillos \emph{topológicos} cuya topología es \emph{lineal},
$\rho:A\to B$ un homomorfismo continuo y $L$ un $B$-bimódulo topológico, y
supongamos que existe un ideal bilateral \emph{abierto} $\mathfrak R_0$ de $B$
tal que $\mathfrak R_0.L=L.\mathfrak R_0=0$, de modo que $L$ puede ser

\oldpage[0\textsubscript{IV-1}]{67}

considerado como un $(B/\mathfrak R)$-bimódulo para todo ideal abierto
$\mathfrak R\subset\mathfrak R_0$. Sean entonces
$\mathfrak R\subset\mathfrak R_0$ un ideal bilateral abierto de $B$,
$\mathfrak J$ un ideal bilateral abierto de $A$ tal que
$\rho(\mathfrak J)\subset\mathfrak R$, de modo que $B/\mathfrak R$ puede ser
considerado como un $(A/\mathfrak J)$-anillo discreto. La observación
precedente prueba que el grupo
$\operatorname{Exan}_{A/\mathfrak J}(B/\mathfrak R,L)$ está definido; además,
si $\mathfrak R'\subset\mathfrak R$, $\mathfrak J'\subset\mathfrak J$ son dos
ideales bilaterales abiertos de $B$ y $A$, respectivamente, tales que
$\rho(\mathfrak J')\subset\mathfrak R'$, se tiene por (18.3.5.1) un
homomorfismo canónico
\begin{equation}
\label{0.18.5.1.1-es}
\tag{18.5.1.1}
  \operatorname{Exan}_{A/\mathfrak J}(B/\mathfrak R,L)
  \longrightarrow
  \operatorname{Exan}_{A/\mathfrak J'}(B/\mathfrak R',L).
\end{equation}

El conjunto de pares de ideales abiertos $(\mathfrak J,\mathfrak R)$ tales que
$\rho(\mathfrak J)\subset\mathfrak R\subset\mathfrak R_0$ está evidentemente
ordenado y es filtrante por la derecha para la relación «~$\mathfrak J\supset\mathfrak J'$
y $\mathfrak R\supset\mathfrak R'$~», y las aplicaciones (18.5.1.1)
definen un \emph{sistema inductivo} de grupos aditivos cuyo conjunto
de índices es este conjunto. Se pone, por abuso de notación (pues ya no se trata
de un grupo en correspondencia biunívoca natural con un conjunto
de extensiones),
\begin{equation}
\label{0.18.5.1.2-es}
\tag{18.5.1.2}
  \operatorname{Exantop}_A(B,L)
  =\varinjlim\operatorname{Exan}_{A/\mathfrak J}(B/\mathfrak R,L).
\end{equation}

Decir que el segundo miembro de (18.5.1.2) es \emph{nulo} significa, pues, que, para
todo par de ideales abiertos $\mathfrak J\subset A$, $\mathfrak R\subset B$
tales que $\rho(\mathfrak J)\subset\mathfrak R\subset\mathfrak R_0$ y toda
$(A/\mathfrak J)$-extensión $E$ de $B/\mathfrak R$ por $L$, existen dos
ideales abiertos $\mathfrak J'\subset\mathfrak J$,
$\mathfrak R'\subset\mathfrak R$ tales que
$\rho(\mathfrak J')\subset\mathfrak R'$ y que la imagen inversa por
el homomorfismo $B/\mathfrak R'\to B/\mathfrak R$ de $E$ sea \emph{trivial}.

Dejamos al lector la definición análoga de los límites inductivos
$\operatorname{Exaltop}_A(B,L)$, $\operatorname{Exalcotop}_A(B,L)$ a partir de
$\operatorname{Exal}_{A/\mathfrak J}(B/\mathfrak R,L)$ y de
$\operatorname{Exalcom}_{A/\mathfrak J}(B/\mathfrak R,L)$ para el caso en que $A$
es conmutativo y $B$ una $A$-álgebra (resp. una $A$-álgebra conmutativa)
topológica.
\end{env}

\begin{env}[18.5.2]
\label{0.18.5.2-es}
Si se tiene un diagrama conmutativo de homomorfismos continuos de anillos
\[
  \xymatrix@C=3pc@R=2pc{
    B'\ar[r] & B\\
    A'\ar[u]\ar[r] & A\ar[u]
  }
\]
se deducen canónicamente dos homomorfismos de grupos aditivos
\[
  \operatorname{Exantop}_A(B,L)
  \longrightarrow\operatorname{Exantop}_{A'}(B,L)
  \longrightarrow\operatorname{Exantop}_{A'}(B',L)
\]
pasando al límite inductivo a partir de (18.3.5.1).

En virtud de la exactitud del funtor $\varinjlim$ en la categoría de los grupos
conmutativos, el núcleo del homomorfismo
\[
  \operatorname{Exantop}_A(B,L)
  \longrightarrow\operatorname{Exantop}_{A'}(B,L)
\]
es el límite inductivo de los núcleos de los homomorfismos
\[
  \operatorname{Exan}_{A/\mathfrak J}(B/\mathfrak R,L)
  \longrightarrow
  \operatorname{Exan}_{A'/\mathfrak J'}(B/\mathfrak R,L)
\]
donde se ha tomado como $\mathfrak J'$ la imagen inversa de $\mathfrak J$; se denota
este núcleo por $\operatorname{Exantop}_{A/A'}(B,L)$. Se definen del mismo modo
$\operatorname{Exaltop}_{A/A'}(B,L)$ y
$\operatorname{Exalcotop}_{A/A'}(B,L)$. Finalmente, si se tiene un homomorfismo

\oldpage[0\textsubscript{IV-1}]{68}

continuo de $B$-bimódulos $L\to L'$, se deduce canónicamente un
homomorfismo de grupos aditivos
\[
  \operatorname{Exantop}_A(B,L)\longrightarrow
  \operatorname{Exantop}_A(B,L')
\]
pasando al límite inductivo a partir de (18.3.6.1).
\end{env}

\begin{env}[18.5.3]
\label{0.18.5.3-es}
Dados un anillo topológico $C$ y dos $C$-bimódulos topológicos
$M$, $N$,\linebreak se denota por $\operatorname{Hom.\ cont}_C(M,N)$ el grupo aditivo de los
$C$-homomorfismos \emph{continuos} de $M$ en $N$.
\end{env}

\begin{lemma}[18.5.3.1]
\label{0.18.5.3.1-es}
Sean $C$ un anillo topológico y $E$, $L$ dos $C$-bimódulos topológicos;
se supone que las topologías son lineales, que $L$ es discreto y está anulado
por un ideal bilateral abierto de $C$. Entonces se tiene un isomorfismo canónico
\begin{equation}
\label{0.18.5.3.2-es}
\tag{18.5.3.2}
  \varinjlim\operatorname{Hom}_{C/\mathfrak R}(E/V,L)
  \xrightarrow{\sim}\operatorname{Hom.\ cont}_C(E,L)
\end{equation}
donde en el primer miembro el límite inductivo se toma con respecto al conjunto
ordenado filtrante de pares $(\mathfrak R,V)$ tales que $\mathfrak R$ es un
ideal bilateral abierto de $C$, $V$ un sub-$C$-bimódulo abierto de $E$, tales que
$\mathfrak R.L=L.\mathfrak R=0$, $\mathfrak R.E\subset V$,
$E.\mathfrak R\subset V$.
\end{lemma}

\begin{proof}
Como $C/\mathfrak R$ y $E/V$ son discretos, se tienen homomorfismos
canónicos
\[
  w_{\mathfrak R,V}:\operatorname{Hom}_{C/\mathfrak R}(E/V,L)
  \longrightarrow\operatorname{Hom.\ cont}_C(E,L)
\]
que forman un sistema inductivo, de donde se obtiene un homomorfismo (18.5.3.2) pasando al
límite inductivo. Como el homomorfismo $E/V'\to E/V$ es sobreyectivo para
$V'\subset V$, de la definición resulta inmediatamente que el homomorfismo
\[
  \operatorname{Hom}_{C/\mathfrak R}(E/V,L)
  \longrightarrow\operatorname{Hom}_{C/\mathfrak R}(E/V',L)
\]
(con $\mathfrak R.L=L.\mathfrak R=0$, $\mathfrak R.E\subset V'$,
$E.\mathfrak R\subset V'$) es inyectivo, y evidentemente sucede lo mismo con
el homomorfismo
\[
  \operatorname{Hom}_{C/\mathfrak R}(E/V,L)
  \longrightarrow\operatorname{Hom}_{C/\mathfrak R'}(E/V,L)
\]
para $\mathfrak R'\subset\mathfrak R$; se concluye que el homomorfismo
(18.5.3.2) es inyectivo. Por otra parte, si $u$ es un $C$-homomorfismo continuo
de $E$ en $L$, su núcleo es un subbimódulo abierto $V_0$ de $E$, y si
$\mathfrak R_0$ es un ideal bilateral abierto de $C$ tal que
$\mathfrak R_0.L=L.\mathfrak R_0=0$ y $\mathfrak R_0.E\subset V_0$,
$E.\mathfrak R_0\subset V_0$, está claro que $u$ es la imagen canónica de un
$(C/\mathfrak R_0)$-homomorfismo de $E/V_0$ en $L$, luego (18.5.3.2) es
sobreyectivo.
\end{proof}

Dicho esto, la proposición (18.3.8) se generaliza como sigue a los anillos
topológicos:

\begin{proposition}[18.5.4]
\label{0.18.5.4-es}
Sean $B$ un anillo topológico linealmente topologizado, $\mathfrak J$ un
ideal bilateral de $B$ y $C=B/\mathfrak J$ el anillo topológico cociente;
$\mathfrak J/\mathfrak J^2$ (donde $\mathfrak J$ está dotado de la topología inducida
por la de $B$ y $\mathfrak J/\mathfrak J^2$ de la topología cociente de
la de $\mathfrak J$) queda entonces canónicamente dotado de una estructura de
$C$-bimódulo topológico. Para todo $C$-bimódulo discreto $L$ anulado por un
ideal abierto de $C$, existe entonces un isomorfismo canónico
\begin{equation}
\label{0.18.5.4.1-es}
\tag{18.5.4.1}
  \operatorname{Hom.\ cont}_C(\mathfrak J/\mathfrak J^2,L)
  \xrightarrow{\sim}\operatorname{Exantop}_B(C,L).
\end{equation}
\end{proposition}

\begin{proof}
En efecto, para todo ideal abierto $\mathfrak R$ de $B$ tal que
$(\mathfrak J+\mathfrak R)/\mathfrak J$ anule $L$, se tiene, según (18.3.9),
un isomorfismo canónico
\[
  \operatorname{Hom}_{B/(\mathfrak J+\mathfrak R)}
    ((\mathfrak J+\mathfrak R)/(\mathfrak J^2+\mathfrak R),L)
  \xrightarrow{\sim}
  \operatorname{Exan}_{B/\mathfrak R}(B/(\mathfrak J+\mathfrak R),L)
\]
y basta pasar al límite inductivo, teniendo en cuenta (18.5.3.1).
\end{proof}

\section{Álgebras formalmente lisas y anillos de Cohen}
\label{section:0.19-es}
\label{0.19-es}

\setcounter{subsection}{-1}
\subsection{Introducción}
\label{subsection:0.19.0-es}
\label{0.19.0-es}

\oldpage[0\textsubscript{IV-1}]{69}

\begin{env}[19.0.1]
\label{0.19.0.1-es}
En el capítulo IV, introduciremos y estudiaremos, entre otras, una clase
importante de morfismos de preesquemas, los morfismos \emph{lisos}\footnote{Llamados
también morfismos \emph{simples} en algunos trabajos recientes (inspirándose
en la terminología clásica de «~punto simple~»); esta terminología se presta,
sin embargo, a confusión, especialmente en teoría de grupos algebraicos.}. Una de
sus propiedades fundamentales (y que, junto con una condición de finitud, los
caracteriza) es una propiedad de \emph{levantamiento de los morfismos}: si
$f:X\to Y$ es un morfismo liso, $g:Y'\to Y$ un morfismo, entonces para todo
morfismo $h:Y''\to Y'$ que haga de $Y''$ un preesquema «~poco diferente~» de
$Y'$, \emph{todo $Y$-morfismo $Y''\to X$ se factoriza} como
$Y''\xrightarrow{h}Y'\to X$. Precisamente, limitémonos al caso en que
$Y=\operatorname{Spec}(A)$, $X=\operatorname{Spec}(B)$ son \emph{afines};
$B$ se llama entonces una $A$-álgebra \emph{formalmente lisa} si, para toda
$A$-álgebra $C$ y todo ideal \emph{nilpotente} $\mathfrak J$ de $C$, todo
$A$-homomorfismo $B\to C/\mathfrak J$ \emph{se levanta} a
$B\to C\to C/\mathfrak J$. En otros términos, la aplicación
\[
  \operatorname{Hom}_A(B,C)\longrightarrow
  \operatorname{Hom}_A(B,C/\mathfrak J)
\]
es \emph{sobreyectiva}. En muchas aplicaciones, $X$ aparecerá como objeto
que representa un \emph{funtor contravariante representable} $Y'\mapsto F(Y')$
de la categoría de los $Y$-preesquemas en la de los conjuntos, de modo
($\mathbf{0}_{\mathrm{III}}$, 8.1.8) que se tendrá
$F(Y')=\operatorname{Hom}_Y(Y',X)$. En el caso afín, si se pone
$F^0(C)=F(\operatorname{Spec}(C))$, la comprobación de que, en las
condiciones anteriores, $F^0(C)\to F^0(C/\mathfrak J)$ es \emph{sobreyectiva}
(que puede hacerse incluso \emph{sin saber} que $F$ es representable) establecerá
que $f$ es liso, siempre que la condición adicional de finitud
se cumpla.
\end{env}

\begin{env}[19.0.2]
\label{0.19.0.2-es}
Para las \emph{álgebras topológicas} sobre un anillo topológico $A$, hay una
noción análoga de álgebra \emph{formalmente lisa} que no precisaremos
aquí (cf.~19.3.1). El estudio de estas nociones se hace primero desde un punto de vista
elemental en el §~19; luego, con ayuda de las propiedades de los diferenciales que
se desarrollarán en los §§~20 y 21, se demostrarán teoremas más delicados
en el §~22. Resumamos aquí los principales resultados sobre las álgebras formalmente
lisas:

\medskip\noindent\textbf{I.} --- Sean $A$, $B$ dos anillos locales
noetherianos, $\varphi:A\to B$ un homomorfismo local, $k$ el cuerpo residual de
$A$, y sea $B_0=B\otimes_A k$; se dota a $A$, $B$ y $B_0$ de sus
topologías preádicas y a $k$ de la topología discreta. Entonces, para que $B$ sea
una $A$-álgebra formalmente lisa, es necesario y suficiente que $B$ sea un
$A$-módulo \emph{plano} y que $B_0$ sea una $k$-álgebra \emph{formalmente
lisa} (19.7.1). Este teorema reduce, pues, la lisura formal, para los anillos
locales noetherianos, a la misma cuestión para los anillos locales noetherianos que
son álgebras sobre un \emph{cuerpo}.

\medskip\noindent\textbf{II.} --- Sean $k$ un cuerpo, $A$ un anillo local
noetheriano que es una $k$-álgebra. Para que $A$ sea formalmente liso, es
necesario y suficiente que $A$ sea \emph{geométricamente regular} sobre $k$,
es decir, que para toda extensión \emph{finita} $k'$ de $k$, el anillo
semilocal $A\otimes_k k'$ sea

\oldpage[0\textsubscript{IV-1}]{70}

\noindent regular (22.5.8); el \emph{completado} $\widehat A$ de $A$ es entonces
un anillo isomorfo a un anillo de series formales
$K[[T_1,\ldots,T_n]]$ (19.6.5). Además, la estructura de $k$-álgebra de $A$,
cuando se supone que $A$ es \emph{completo} y formalmente liso sobre $k$, está
enteramente determinada por el cuerpo residual $K$ de $A$ y la \emph{dimensión}
de $A$; esta puede, por lo demás, ser arbitraria con tal de satisfacer la
desigualdad
\[
  \dim(A)\geq \operatorname{rg}_k(\Gamma_{K/k}),
\]
donde $\Gamma_{K/k}$ es el «~módulo de imperfección~» de $K$ (22.2.6).

En particular, para que una \emph{extensión} $K$ de $k$ sea formalmente
lisa, es necesario y suficiente que $K$ sea una extensión \emph{separable} de $k$
(19.6.1).

\medskip\noindent\textbf{III.} --- Sean $A$ un anillo local noetheriano,
$\mathfrak J$ un ideal de $A$ distinto de $A$, $A_0=A/\mathfrak J$, $B_0$ un
anillo local noetheriano \emph{completo}, $A_0\to B_0$ un homomorfismo local
que hace de $B_0$ una $A_0$-álgebra formalmente lisa. Entonces existe un
anillo local noetheriano \emph{completo} $B$, un homomorfismo local $A\to B$
que hace de $B$ un $A$-módulo plano, y un $A_0$-\emph{isomorfismo}
\[
  B\otimes_A A_0\xrightarrow{\sim}B_0
\]
(de modo que, por I), $B$ es una $A$-álgebra \emph{formalmente lisa}); además,
$B$ está determinado por estas propiedades salvo isomorfismo (19.7.2). Este
teorema contiene en particular los \emph{teoremas de Cohen} sobre la estructura
de los anillos locales noetherianos \emph{completos} (19.8), que desempeñarán un papel
importante en los §§~6 y 7 del capítulo IV.

\medskip\noindent\textbf{IV.} --- El interés del estudio de los anillos locales
noetherianos formalmente lisos sobre otro proviene de la siguiente caracterización
«~puntual~» de los morfismos lisos: si $X$ e $Y$ son
preesquemas localmente noetherianos, $f:X\to Y$ un morfismo localmente de tipo
finito, entonces, para que $f$ sea liso, es necesario y suficiente que para todo
$x\in X$, el anillo $\mathcal O_x$ sea formalmente liso sobre
$\mathcal O_{f(x)}$. En particular, si $Y=\operatorname{Spec}(k)$, donde $k$ es
un cuerpo perfecto, decir que $f$ es liso equivale a decir que $X$ es un
preesquema \emph{regular}.

\medskip\noindent\textbf{V.} --- Finalmente, veremos en los §§~20, 21 y 22 que
la noción de álgebra formalmente lisa se introduce de modo natural en la
teoría de los \emph{diferenciales de Kähler}, esclareciéndose ambas teorías
mutuamente.
\end{env}

\begin{env}[19.0.3]
\label{0.19.0.3-es}
En todo este parágrafo y los siguientes, se supondrá que los anillos y módulos topológicos
están \emph{linealmente topologizados} ($\mathbf{0}_{\mathrm I}$,
7.1.1); se supondrá que los anillos topológicos considerados son \emph{conmutativos},
salvo mención expresa de lo contrario. Recordemos que, si $A$ y $B$ son dos
anillos topológicos, $\rho:A\to B$ un homomorfismo de anillos que define sobre
$B$ una estructura de $A$-álgebra, se dice que $B$ es una \emph{$A$-álgebra
topológica} si $\rho$ es \emph{continuo} para las topologías consideradas.

Para abreviar, en un anillo topológico $A$ (resp. un $A$-módulo topológico
$M$), diremos «~\emph{sistema fundamental de ideales} (resp.
\emph{submódulos}) \emph{abiertos}~» en lugar de «~sistema fundamental de
entornos de 0 formado por ideales (resp. submódulos)~».

Dados un anillo topológico $A$ y un $A$-módulo $M$, los conjuntos
$\mathfrak J M$, donde $\mathfrak J$ recorre un sistema fundamental de ideales
abiertos, forman un sistema fundamental de submódulos abiertos de una topología
sobre $M$ que hace de $M$ un $A$-módulo topológico, y que se llama \emph{inducida}
por la topología de $A$.

Sea $M$ un $A$-módulo topológico cuya topología es \emph{menos fina} que
la topología inducida por la de $A$; entonces, si $N$ es un submódulo
\emph{abierto} de $M$, el $A$-módulo \emph{discreto} $M/N$ está \emph{anulado por
un ideal abierto} de $A$, pues por hipótesis existe tal ideal $\mathfrak R$
tal que $\mathfrak R.M\subset N$.

\oldpage[0\textsubscript{IV-1}]{71}

Si $M$ y $N$ son dos $A$-módulos topológicos cuyas topologías son ambas
inducidas por la de $A$, entonces \emph{todo} $A$-homomorfismo $u:M\to N$
es \emph{continuo}, pues para todo entorno $V$ de 0 en $N$ existe por
hipótesis un ideal abierto $\mathfrak J$ de $A$ tal que $\mathfrak JN\subset V$,
y se tiene, por tanto, $u(\mathfrak JM)\subset\mathfrak JN\subset V$.

Cuando $B$ es una $A$-álgebra topológica, la topología sobre $B$ \emph{inducida}
por la de $A$ es \emph{más fina} que la topología dada, pues para todo
ideal abierto $\mathfrak R$ de $B$, hay por hipótesis un ideal abierto
$\mathfrak J$ de $A$ tal que $\mathfrak JB\subset\mathfrak R$.
\end{env}

\subsection{Epimorfismos y monomorfismos formales}
\label{subsection:0.19.1-es}

\begin{proposition}[19.1.1]
\label{0.19.1.1-es}
Sean $A$ un anillo topológico, $M$, $N$ dos $A$-módulos topológicos,
$(W_\lambda)$ un sistema fundamental de submódulos abiertos en $N$,
$u:M\to N$ un $A$-homomorfismo continuo, $\widehat M$, $\widehat N$ los
completados separados de $M$ y $N$, $\widehat u:\widehat M\to\widehat N$ la
extensión continua de $u$ a los completados separados.
\begin{enumerate}
  \item[(i)] Las condiciones siguientes son equivalentes:
  \begin{enumerate}
    \item[a)] $u(M)$ es \emph{denso} en $N$.
    \item[a$'$)] $\widehat u(\widehat M)$ es \emph{denso} en $\widehat N$.
    \item[a$''$)] Para todo $\lambda$, el homomorfismo compuesto
    $M\xrightarrow{u}N\to N/W_\lambda$ es \emph{sobreyectivo}.
  \end{enumerate}
  \item[(ii)] Las condiciones siguientes son equivalentes:
  \begin{enumerate}
    \item[b)] La imagen inversa por $u$ de la topología de $N$ es igual a la
    topología de $M$.
    \item[b$'$)] $\widehat u$ es un \emph{isomorfismo} del $\widehat A$-módulo
    topológico $\widehat M$ sobre el sub-$\widehat A$-módulo topológico
    $\widehat u(\widehat M)$ de $\widehat N$ (que es necesariamente
    \emph{cerrado}).
    \item[b$''$)] Los $u^{-1}(W_\lambda)$ forman un sistema fundamental de
    entornos de 0 en $M$.
  \end{enumerate}
\end{enumerate}
Esto resulta inmediatamente de la definición de los completados separados y del hecho
de que $N/W_\lambda$ es discreto.
\end{proposition}

\begin{env}[19.1.2]
\label{0.19.1.2-es}
Cuando se cumplen las condiciones equivalentes de (i) (resp. (ii)) de (19.1.1),
se dice que $u$ es un \emph{epimorfismo formal} (resp. un
\emph{monomorfismo formal}). Se dice que $u$ es un \emph{bimorfismo formal}
si es a la vez un monomorfismo formal y un epimorfismo formal; en
virtud de (19.1.1), equivale a decir que $\widehat u$ es un
\emph{isomorfismo} del $\widehat A$-módulo topológico $\widehat M$ sobre el
$\widehat A$-módulo topológico $\widehat N$.
\end{env}

\begin{proposition}[19.1.3]
\label{0.19.1.3-es}
Sean $A$ un anillo topológico, $M$, $N$ dos $A$-módulos topológicos,
$u:M\to N$ un $A$-homomorfismo continuo. Se supone que existen dos sistemas
fundamentales de submódulos abiertos, $(V_\lambda)$, $(W_\lambda)$ en $M$ y
$N$ respectivamente, con el mismo conjunto de índices y tales que
$u(V_\lambda)\subset W_\lambda$ para todo $\lambda$; sea
$u_\lambda:M/V_\lambda\to N/W_\lambda$ el homomorfismo inducido por $u$ al
pasar a los cocientes. Entonces:
\begin{enumerate}
  \item[(i)] Para que $u$ sea un epimorfismo formal, es necesario y suficiente que,
  para todo $\lambda$, $u_\lambda$ sea sobreyectivo.
  \item[(ii)] Para que $u$ sea un monomorfismo formal, es necesario y suficiente
  que, para todo $\lambda$, exista un $\mu$ tal que
  $V_\mu\subset V_\lambda$ y que $\operatorname{Ker}(u_\mu)$ esté contenido
  en el núcleo $V_\lambda/V_\mu$ de la aplicación canónica
  $M/V_\mu\to M/V_\lambda$.
  \item[(iii)] Para que $u$ sea un bimorfismo formal, es necesario y suficiente
  que, para todo $\lambda$, $u_\lambda$ sea sobreyectivo y que exista un
  $\mu$ tal que $V_\mu\subset V_\lambda$ y que el homomorfismo canónico
  $M/V_\mu\to M/V_\lambda$

  \oldpage[0\textsubscript{IV-1}]{72}

  \noindent se factorice como
  $M/V_\mu\xrightarrow{u_\mu}N/W_\mu\to M/V_\lambda$ (donde
  el homomorfismo $N/W_\mu\to M/V_\lambda$ es necesariamente único).
\end{enumerate}
\end{proposition}

\begin{proof}
La afirmación (i) es inmediata; por otra parte,
$\operatorname{Ker}(u_\mu)=u^{-1}(W_\mu)/V_\mu$, y el núcleo de la aplicación
canónica $M/V_\mu\to M/V_\lambda$ es $V_\lambda/V_\mu$; la condición de
(ii) equivale, por tanto, a decir que $u^{-1}(W_\mu)\subset V_\lambda$, y la afirmación
(ii) se deduce de ello inmediatamente; finalmente, cuando $u_\mu$ es sobreyectivo, da lo mismo
decir que $\operatorname{Ker}(u_\mu)$ está contenido en
$V_\lambda/V_\mu$, o a decir que $M/V_\mu\to M/V_\lambda$ se factoriza como
$v\circ u_\mu$, donde $v$ es un homomorfismo
$N/W_\mu\to M/V_\lambda$, ya que $N/W_\mu$ se identifica entonces con
$(M/V_\mu)/\operatorname{Ker}(u_\mu)$.
\end{proof}

\begin{corollary}[19.1.4]
\label{0.19.1.4-es}
Sean $A$ un anillo topológico y $M$, $N$ dos $A$-módulos topológicos cuyas
topologías se deducen de la de $A$, y sea $u:M\to N$ un epimorfismo
formal. Sea $(\mathfrak J_\lambda)$ un sistema fundamental de ideales abiertos
en $A$. Para que $u$ sea un bimorfismo formal, es necesario y suficiente que para
todo $\lambda$, el homomorfismo
$u_\lambda:M/\mathfrak J_\lambda M\to N/\mathfrak J_\lambda N$ deducido de $u$
pasando a los cocientes sea biyectivo.
\end{corollary}

\begin{proof}
En efecto, se tiene $u(\mathfrak J_\lambda M)\subset\mathfrak J_\lambda N$ y se
puede aplicar el criterio (19.1.3, (iii)); pero, si se tiene una factorización
\[
  M/\mathfrak J_\mu M\xrightarrow{u_\mu}
  N/\mathfrak J_\mu N\xrightarrow{v}M/\mathfrak J_\lambda M,
\]
la imagen de $\mathfrak J_\lambda M/\mathfrak J_\mu M$ por $u_\mu$ es
$\mathfrak J_\lambda N/\mathfrak J_\mu N$, luego la imagen de
$\mathfrak J_\lambda N/\mathfrak J_\mu N$ por $v$ es 0 y $v$ se factoriza en
\[
  N/\mathfrak J_\mu N\longrightarrow N/\mathfrak J_\lambda N
  \xrightarrow{v'}M/\mathfrak J_\lambda M~;
\]
pero entonces el automorfismo identidad de $M/\mathfrak J_\lambda M$ se factoriza
en
\[
  M/\mathfrak J_\lambda M\xrightarrow{u_\lambda}
  N/\mathfrak J_\lambda N\xrightarrow{v'}M/\mathfrak J_\lambda M,
\]
lo que muestra que $u_\lambda$ es inyectivo y, como ya se sabe que es
sobreyectivo, esto prueba el corolario.
\end{proof}

\begin{proposition}[19.1.5]
\label{0.19.1.5-es}
Sean $A$ un anillo topológico, $M$, $N$ dos $A$-módulos topológicos y
$u:M\to N$ un $A$-homomorfismo continuo. Las condiciones siguientes son
equivalentes:
\begin{enumerate}
  \item[a)] Para todo $A$-módulo topológico discreto $E$ y todo
  $A$-homomorfismo continuo $v:M\to E$, existe un $A$-homomorfismo continuo
  $w:N\to E$ tal que $v=w\circ u$.
  \item[b)] Para todo submódulo abierto $V'$ de $M$, existen un submódulo
  abierto $W''$ de $N$, un submódulo abierto
  $V''\subset V'\cap u^{-1}(W'')$ y un homomorfismo
  $h:N/W''\to M/V'$ tales que el homomorfismo canónico
  $M/V''\to M/V'$ se factoriza en
  \[
    M/V''\xrightarrow{u''}N/W''\xrightarrow{h}M/V',
  \]
  donde $u''$ se deduce de $u$ pasando a los cocientes.
\end{enumerate}
\end{proposition}

\begin{proof}
Para ver que a) implica b), basta aplicar a) a $E=M/V'$ y al
homomorfismo canónico $v:M\to M/V'$; entonces $W''=w^{-1}(0)$ es un
submódulo abierto de $N$ y $V''=u^{-1}(W'')$ un submódulo abierto de $M$
contenido en $V'$; estos submódulos y el homomorfismo
$h:N/W''\to M/V'$ deducido de $w$ pasando al cociente satisfacen lo
requerido. Recíprocamente, para mostrar que b) implica a), puede uno limitarse al
caso en que $v$ es sobreyectivo, reemplazando $E$ por $v(M)$; entonces
$V'=\operatorname{Ker}(v)$ es un submódulo abierto de $M$, por lo que $E$ es
isomorfo al $A$-módulo discreto $M/V'$, y aplicando b), se obtiene

\oldpage[0\textsubscript{IV-1}]{73}

el $A$-homomorfismo continuo $w$ buscado tomando el compuesto
$N\to N/W''\xrightarrow{h}M/V'$, siendo el diagrama
\[
  \xymatrix{
    M\ar[r]^u\ar[d] & N\ar[d]\\
    M/V''\ar[r]_{u''} & N/W''
  }
\]
conmutativo.
\end{proof}

Cuando se satisfacen las condiciones equivalentes de (19.1.5), diremos que
$u$ es \emph{formalmente invertible por la izquierda}; como la condición b) de
(19.1.5) implica que $u^{-1}(W'')\subset V'$, $u$ es entonces un monomorfismo
formal. La terminología queda además justificada por los corolarios siguientes:

\begin{corollary}[19.1.6]
\label{0.19.1.6-es}
Si existe un $\widehat A$-homomorfismo continuo
$s:\widehat N\to\widehat M$ tal que
$s\circ\widehat u=1_{\widehat M}$, entonces $u$ es formalmente invertible por la
izquierda.
\end{corollary}

Se observará que $\widehat u$ es entonces un isomorfismo topológico de
$\widehat M$ sobre un \emph{factor directo topológico} de $\widehat N$. Para
probar (19.1.6), basta observar que, con la notación de (19.1.5, a)),
$v:M\to E$ se prolonga a un $\widehat A$-homomorfismo continuo
$\widehat v:\widehat M\to E$ puesto que $E$ es discreto; sean
$j:M\to\widehat M$ y $j':N\to\widehat N$ los homomorfismos canónicos;
entonces, poniendo $w=\widehat v\circ s\circ j'$, se tiene
$w\circ u=\widehat v\circ s\circ\widehat u\circ j=\widehat v\circ j=v$.

\begin{corollary}[19.1.7]
\label{0.19.1.7-es}
Supongamos que las topologías de $M$ y de $N$ se deducen de la de $A$,
y sea $(\mathfrak J_\lambda)$ un sistema fundamental de ideales abiertos en
$A$. Para que $u$ sea formalmente invertible por la izquierda, es necesario y suficiente
que, para todo $\lambda$, el homomorfismo
$u_\lambda:M/\mathfrak J_\lambda M\to N/\mathfrak J_\lambda N$, deducido de
$u$ pasando a los cocientes, sea invertible por la izquierda (dicho de otro modo, sea
un isomorfismo de $M/\mathfrak J_\lambda M$ sobre un factor directo de
$N/\mathfrak J_\lambda N$).
\end{corollary}

En efecto, la condición es suficiente, pues, para
$V'=\mathfrak J_\lambda M$ en la condición b) de (19.1.5), se satisface lo
requerido tomando $W''=\mathfrak J_\lambda N$, $V''=V'$ y $h$ tal que
$h\circ u_\lambda$ sea la identidad en $M/\mathfrak J_\lambda M$.
Recíprocamente, si $u$ es formalmente invertible por la izquierda, entonces, para todo
$\lambda$, existen, en virtud de (19.1.5, b)), un $\mu$ tal que
$\mathfrak J_\mu\subset\mathfrak J_\lambda$ y un homomorfismo
$h:N/\mathfrak J_\mu N\to M/\mathfrak J_\lambda M$ tales que el homomorfismo
canónico $M/\mathfrak J_\mu M\to M/\mathfrak J_\lambda M$ se factoriza en
\[
  M/\mathfrak J_\mu M\xrightarrow{u_\mu}N/\mathfrak J_\mu N
  \xrightarrow{h}M/\mathfrak J_\lambda M~;
\]
pero, como
$h(\mathfrak J_\lambda(N/\mathfrak J_\mu N))
\subset\mathfrak J_\lambda(M/\mathfrak J_\lambda M)=0$, $h$ se factoriza
canónicamente en
\[
  N/\mathfrak J_\mu N\longrightarrow N/\mathfrak J_\lambda N
  \xrightarrow{s}M/\mathfrak J_\lambda M,
\]
y es inmediato que $s$ es inverso por la izquierda de $u_\lambda$.

\begin{proposition}[19.1.8]
\label{0.19.1.8-es}
Sea $A$ un anillo topológico que admite un sistema fundamental numerable
decreciente $(\mathfrak J_n)$ de ideales abiertos. Sean $M$, $N$ dos
$A$-módulos topológicos cuyas topologías se deducen de la de $A$,
y sea $u:M\to N$ un $A$-homomorfismo. Pongamos
$M_n=M/\mathfrak J_nM$, $N_n=N/\mathfrak J_nN$ y sea
$u_n:M_n\to N_n$ el $(A/\mathfrak J_n)$-homomorfismo deducido de $u$
pasando a los cocientes. Supongamos que, para todo $n$,
$P_n=\operatorname{Coker}(u_n)$ es un $(A/\mathfrak J_n)$-módulo proyectivo
y que $M$ es separado y completo. Entonces, para que $u$ sea formalmente
invertible por la izquierda, es necesario y suficiente que $u$ sea invertible por la izquierda (y
$u$ es entonces un isomorfismo topológico de $M$ sobre un factor directo
topológico de $N$).
\end{proposition}

\oldpage[0\textsubscript{IV-1}]{74}

\begin{proof}
En virtud de (19.1.7), se tienen diagramas conmutativos
\[
  \xymatrix{
    0\ar[r] & M_n\ar[r]^{u_n} & N_n\ar[r]^{p_n} & P_n\ar[r] & 0\\
    0\ar[r] & M_{n+1}\ar[u]^f\ar[r]_{u_{n+1}} &
      N_{n+1}\ar[u]^g\ar[r]_{p_{n+1}} & P_{n+1}\ar[u]^h\ar[r] & 0
  }
\]
cuyas filas son sucesiones exactas \emph{escindidas}, es decir, existe
para cada $n$ un homomorfismo $s_n:P_n\to N_n$ tal que
$p_n\circ s_n=1_{P_n}$. Vamos a mostrar que se puede, por inducción sobre
$n$, definir un homomorfismo $s'_n:P_n\to N_n$ tal que
$p_n\circ s'_n=1_{P_n}$ y que los diagramas
\[
  \xymatrix{
    P_n\ar[r]^{s'_n} & N_n\\
    P_{n+1}\ar[u]^h\ar[r]_{s'_{n+1}} & N_{n+1}\ar[u]^g
  }
\]
sean conmutativos. En efecto, $g\circ s_{n+1}-s'_n\circ h$ es un
homomorfismo de $P_{n+1}$ en
\[
  u_n(M_n)=u_{n+1}(M_{n+1})/
  \mathfrak J_nu_{n+1}(M_{n+1})~;
\]
como $P_{n+1}$ es un $(A/\mathfrak J_{n+1})$-módulo \emph{proyectivo}, este
homomorfismo se factoriza en
\[
  P_{n+1}\xrightarrow{t_{n+1}}u_{n+1}(M_{n+1})
  \longrightarrow
  u_{n+1}(M_{n+1})/\mathfrak J_nu_{n+1}(M_{n+1}),
\]
de donde se concluye inmediatamente que $s'_{n+1}=s_{n+1}-t_{n+1}$ satisface lo
requerido. Dicho esto, de la descomposición de $N_n$ en suma directa de
$u_n(M_n)$ y de $s'_n(P_n)$, se deduce inmediatamente un homomorfismo
$w_n:N_n\to M_n$ inverso por la izquierda de $u_n$ y tal que los diagramas
\[
  \xymatrix{
    N_n\ar[r]^{w_n} & M_n\\
    N_{n+1}\ar[u]^g\ar[r]_{w_{n+1}} & M_{n+1}\ar[u]^f
  }
\]
sean conmutativos. El sistema proyectivo $(w_n)$ admite entonces un límite
proyectivo $\widehat w:\widehat N\to\widehat M=M$, de donde se obtiene, por composición con
el homomorfismo canónico $N\to\widehat N$, un homomorfismo $w:N\to M$ tal
que, para todo $n$, el endomorfismo $(w\circ u)_n$ de
$M_n=M/\mathfrak J_nM$ deducido de $w\circ u$ pasando a los cocientes sea
la identidad; esto implica que $w\circ u=1_M$ puesto que $M$ es separado y
completo.
\end{proof}

\begin{proposition}[19.1.9]
\label{0.19.1.9-es}
Sean $A$ un anillo topológico preadmisible ($\mathbf{0}_{\mathrm I}$,
7.1.2), $\mathfrak Q$ un ideal de definición de $A$ y
$(\mathfrak J_\lambda)$ un sistema fundamental de ideales abiertos de $A$.
Sean $M$, $N$ dos $A$-módulos topológicos cuyas topologías se deducen
de la de $A$, y tales que para todo $\lambda$,
$N/\mathfrak J_\lambda N$ es un $(A/\mathfrak J_\lambda)$-módulo proyectivo
(cf. (19.2.3)). Sea $u:M\to N$ un $A$-homomorfismo. Entonces las condiciones
siguientes son equivalentes:
\begin{enumerate}
  \item[a)] $u$ es formalmente invertible por la izquierda.

  \oldpage[0\textsubscript{IV-1}]{75}

  \item[b)] El homomorfismo
  $u_0:M/\mathfrak QM\to N/\mathfrak QN$ deducido de $u$ pasando a los
  cocientes es invertible por la izquierda.
\end{enumerate}
\end{proposition}

\begin{proof}
Se ha visto (19.1.7) que la condición a) equivale a decir que $u_\lambda$ es
invertible por la izquierda para todo $\lambda$; como $\mathfrak Q$ es un ideal
abierto y, por tanto, contiene un $\mathfrak J_\lambda$, se deduce inmediatamente que
$u_0$ es invertible por la izquierda. Para mostrar recíprocamente que b) implica a),
observemos que, para todo $\lambda$, $\mathfrak Q/\mathfrak J_\lambda$ es por
hipótesis un ideal nilpotente de $A/\mathfrak J_\lambda$. Nuestra afirmación
resultará de la proposición siguiente.
\end{proof}

\begin{proposition}[19.1.10]
\label{0.19.1.10-es}
Sean $A$ un anillo, $M$, $N$ dos $A$-módulos, siendo $N$ proyectivo, y
$u:M\to N$ un $A$-homomorfismo. Sea $\mathfrak J$ un ideal de $A$; se
supone que se satisface una de las condiciones siguientes:
\begin{enumerate}
  \item[(i)] $\mathfrak J$ es nilpotente.
  \item[(ii)] $\mathfrak J$ está contenido en el radical de $A$ y $M$ es de
  tipo finito.
\end{enumerate}
Entonces, para que $u$ sea invertible por la izquierda, es necesario y suficiente que
el homomorfismo
$u_0:M/\mathfrak JM\to N/\mathfrak JN$ de $(A/\mathfrak J)$-módulos, deducido
de $u$ por paso a los cocientes sea invertible por la izquierda.
\end{proposition}

\begin{proof}
Como la condición es evidentemente necesaria, probemos que es suficiente. Sea
$v_0$ un inverso por la izquierda de $u_0$; el homomorfismo compuesto
$N\to N/\mathfrak JN\xrightarrow{v_0}M/\mathfrak JM$ se factoriza como
$N\xrightarrow{v}M\to M/\mathfrak JM$, puesto que $N$ es proyectivo; entonces
$w=v\circ u$ es un endomorfismo de $M$ tal que el endomorfismo $w_0$ de
$M/\mathfrak JM$ deducido por paso a los cocientes es la identidad, y basta
probar que el propio $w$ es biyectivo (pues entonces $w^{-1}\circ v$ será una
inverso por la izquierda de $u$). Distingamos ahora los dos casos.

(i) Para todo $n$ se tiene el diagrama conmutativo
\[
  \xymatrix{
    (\mathfrak J^n/\mathfrak J^{n+1})\otimes_{A/\mathfrak J}
      (M/\mathfrak JM)\ar[r]\ar[d]_{1\otimes\operatorname{gr}^0(w)} &
      \mathfrak J^nM/\mathfrak J^{n+1}M\ar[d]^{\operatorname{gr}^n(w)}\\
    (\mathfrak J^n/\mathfrak J^{n+1})\otimes_{A/\mathfrak J}
      (M/\mathfrak JM)\ar[r] & \mathfrak J^nM/\mathfrak J^{n+1}M
  }
\]
donde las flechas horizontales son \emph{sobreyectivas}, y como
$\operatorname{gr}^0(w)=w_0$ es la identidad, lo mismo ocurre con
$\operatorname{gr}^n(w)$, que, \emph{a fortiori}, es biyectivo. Como la
filtración $\mathfrak J$-preádica sobre $M$ es finita, puesto que $\mathfrak J$
es nilpotente, se concluye que $w$ es biyectivo (Bourbaki,
\emph{Alg. comm.}, cap.~III, §~2, n\textsuperscript{o}~8, cor.~3 del teorema~1).

(ii) Basta demostrar que, para todo ideal máximo $\mathfrak m$ de $A$, el
endomorfismo $w_{\mathfrak m}$ de $M_{\mathfrak m}$ es biyectivo (Bourbaki,
\emph{Alg. comm.}, cap.~II, §~3, n\textsuperscript{o}~3, teorema~1), y como
$\mathfrak JA_{\mathfrak m}\subset\mathfrak mA_{\mathfrak m}$ y
$A_{\mathfrak m}/\mathfrak JA_{\mathfrak m}=(A/\mathfrak J)_{\mathfrak m}$,
se reduce a probar la proposición cuando $A$ es un anillo local. Además, se
puede suponer que $\mathfrak J$ es el ideal máximo $\mathfrak J_0$ de $A$, pues
si $u_0$ es invertible

\oldpage[0\textsubscript{IV-1}]{76}

por la izquierda, también lo es
$u_{00}:M/\mathfrak J_0M\to N/\mathfrak J_0N$, obtenido tensorizando $u_0$ con
$1_{A/\mathfrak J_0}$, puesto que se tiene
\[
  M/\mathfrak J_0M=(M/\mathfrak JM)\otimes_{A/\mathfrak J}(A/\mathfrak J_0).
\]
Supongamos, pues, que $\mathfrak J$ es máximo, de modo que $A/\mathfrak J$ es
un cuerpo. Es evidentemente suficiente demostrar que $M$ es un $A$-módulo
\emph{libre} bajo las condiciones del enunciado: en efecto, $\det(w_0)$ es
entonces la imagen canónica de $\det(w)$ en $A/\mathfrak J$, luego $\det(w)$
no pertenece al ideal $\mathfrak J$ y es, por tanto, invertible. Ahora bien,
como el $(A/\mathfrak J)$-espacio vectorial $M/\mathfrak JM$ es libre de tipo
finito, existe un $A$-módulo libre de tipo finito $L$ y un $A$-homomorfismo $f:L\to M$
tal que el homomorfismo $f_0:L/\mathfrak JL\to M/\mathfrak JM$ deducido de $f$
por paso a los cocientes es biyectivo. Como $M$ es de tipo finito, se concluye
primero que $f$ es sobreyectivo por el lema de Nakayama (Bourbaki,
\emph{Alg. comm.}, cap.~II, §~3, n\textsuperscript{o}~2, cor.~1 de la
prop.~4); además, si $g=u\circ f$, el homomorfismo
$g_0:L/\mathfrak JL\to N/\mathfrak JN$ deducido por paso a los cocientes es
invertible por la izquierda, y como aquí $L$ es libre, la observación del
comienzo prueba que el propio $g$ es invertible por la izquierda; ahora bien,
esto implica evidentemente que $f$ es \emph{inyectivo}, lo que termina la demostración.
\end{proof}

Mencionemos de paso los siguientes corolarios de (19.1.10):

\begin{corollary}[19.1.11]
\label{0.19.1.11-es}
Sean $A$ un anillo local, $k$ su cuerpo residual, $M$ un $A$-módulo de tipo
finito, $N$ un $A$-módulo proyectivo y $u:M\to N$ un homomorfismo. Para que
$u$ sea invertible por la izquierda, es necesario y suficiente que exista un
sistema de generadores $(x_i)_{1\leq i\leq m}$ de $M$ tal que las imágenes por
$u\otimes 1:M\otimes_Ak\to N\otimes_Ak$ de los $x_i\otimes1$ sean linealmente
independientes en $N\otimes_Ak$; los $x_i$ forman entonces una base de $M$.
\end{corollary}

La condición es evidentemente necesaria, pues si $u$ es invertible por la
izquierda, $M$ es un $A$-módulo proyectivo de tipo finito, y por tanto libre.
Recíprocamente, si se satisface la condición, es claro que los $x_i\otimes1$
forman una base de $M\otimes_Ak$ y que $u\otimes1$ es invertible por la
izquierda; basta, pues, aplicar (19.1.10) al ideal máximo $\mathfrak J$ de $A$.

\begin{corollary}[19.1.12]
\label{0.19.1.12-es}
Sean $A$ un anillo, $M$ un $A$-módulo de tipo finito, $N$ un $A$-módulo
proyectivo y $u:M\to N$ un homomorfismo. Para todo ideal primo
$\mathfrak p\in\operatorname{Spec}(A)$, las condiciones siguientes son
equivalentes:
\begin{enumerate}
  \item[a)] El homomorfismo
  $u_{\mathfrak p}:M_{\mathfrak p}\to N_{\mathfrak p}$ es invertible por la
  izquierda.
  \item[b)] El homomorfismo
  $u\otimes1:M\otimes_Ak(\mathfrak p)\to N\otimes_Ak(\mathfrak p)$ es
  inyectivo (se recuerda que $k(\mathfrak p)$ designa el cuerpo residual de $A$
  en $\mathfrak p$).
  \item[c)] Existe un sistema finito de elementos $x_i\in M$
  $(1\leq i\leq m)$ tal que las imágenes de los $x_i$ en $M_{\mathfrak p}$
  generan $M_{\mathfrak p}$, y un sistema de $m$ formas lineales
  $y_i^*$ sobre $N$ $(1\leq i\leq m)$ tal que
  $\det(\langle y_i^*,u(x_j)\rangle)\notin\mathfrak p$.
  \item[d)] Existe $f\in A-\mathfrak p$ tal que el homomorfismo
  $u_f:M_f\to N_f$ sea invertible por la izquierda.
\end{enumerate}
Además, el conjunto de los $\mathfrak p\in\operatorname{Spec}(A)$ que verifican
estas condiciones es abierto en $\operatorname{Spec}(A)$.
\end{corollary}

La última afirmación es consecuencia trivial de d). Como $N$ es proyectivo, es
sumando directo de un $A$-módulo libre $A^{(I)}$; además, como $M$ es de tipo
finito, $u(M)$ está contenido en un submódulo de $A^{(I)}$ de la forma $A^n$
para $n$ finito; como cada uno de los enunciados a), b), c), d) es equivalente
al enunciado correspondiente en el que se reemplaza $N$ por $N\oplus P$
(suponiéndose siempre que $u$ aplica $M$ en $N$), se reduce al caso en que $N$
es \emph{libre de tipo finito}. Es trivial que d) implica a), y a) y b) son
equivalentes en virtud de (19.1.11); además, a) implica que
$M_{\mathfrak p}$ es libre (19.1.11), luego a) implica c), tomando como $x_i$
una familia cuyas imágenes en $M_{\mathfrak p}$ forman una base de este

\oldpage[0\textsubscript{IV-1}]{77}

$A_{\mathfrak p}$-módulo y observando que (puesto que $N$ es libre de tipo
finito) toda forma lineal sobre el $A_{\mathfrak p}$-módulo
$N_{\mathfrak p}$ se escribe $u=v/f$, donde $f\in A-\mathfrak p$ y $v$ es una
forma lineal sobre el $A$-módulo $N$. Es claro que c) implica b), y queda, por
tanto, ver que a) implica d). Ahora bien, como $N$ es de presentación finita,
$(\operatorname{Hom}_A(N,M))_{\mathfrak p}$ se identifica canónicamente con
$\operatorname{Hom}_{A_{\mathfrak p}}(N_{\mathfrak p},M_{\mathfrak p})$
(Bourbaki, \emph{Alg. comm.}, cap.~II, §~2, n\textsuperscript{o}~7,
prop.~19). Si $w_{\mathfrak p}$ es inverso por la izquierda de $u_{\mathfrak p}$,
existen, por tanto, un homomorfismo $w:N\to M$ y un elemento $f\in A-\mathfrak p$ tales que
$w_{\mathfrak p}=w\otimes(1/f)1_{A_{\mathfrak p}}$. La relación
$w_{\mathfrak p}\circ u_{\mathfrak p}=1_{M_{\mathfrak p}}$ se escribe también
$(w\circ u)\otimes1_{A_{\mathfrak p}}=f.1_{M_{\mathfrak p}}$. Pero como $M$
es un $A$-módulo de tipo finito, existe $g\in A-\mathfrak p$ tal que el
endomorfismo $g((w\circ u)-f.1_M)$ se anula en todos los generadores de $M$, y
es, por tanto, nulo. Si se pone $h=fg$, y
$u_h=u\otimes1_{A_h}$, $w_h=w\otimes1_{A_h}$, se tiene entonces
$w_h(u_h(z))=(f/1)z$ para todo $z\in M_h$; pero como $h/1$ es invertible en
$A_h$, también lo es $f/1=f'$, y $f'^{-1}w_h$ es, por tanto, inverso por la
izquierda de $u_h$, lo que termina la demostración.

\begin{proposition}[19.1.13]
\label{0.19.1.13-es}
Supongamos verificadas las hipótesis de (19.1.9) y supongamos además que
$(\mathfrak J_\lambda)$ es una sucesión decreciente $(\mathfrak J_n)$ y que $M$
es separado y completo. Entonces las condiciones a) y b) de (19.1.9) son
también equivalentes a:
\begin{enumerate}
  \item[c)] $u$ es invertible por la izquierda.
\end{enumerate}
\end{proposition}

Ya se sabe que c) implica a) (19.1.6). Recíprocamente, si se verifica a), se
sabe (con las notaciones de (19.1.8)) que $M_n$ es sumando directo del
$(A/\mathfrak J_n)$-módulo proyectivo $N_n$, luego $P_n$ también es isomorfo a
un sumando directo de $N_n$ y es, por tanto, proyectivo; basta entonces aplicar
(19.1.8).

Señalemos, por último, la proposición siguiente:

\begin{proposition}[19.1.14]
\label{0.19.1.14-es}
Sean $A$ un anillo, $M$ un $A$-módulo de tipo finito, $N$ un $A$-módulo
proyectivo y $u:M\to N$ un $A$-homomorfismo.
\begin{enumerate}
  \item[(i)] Para que $u$ sea invertible por la izquierda, es necesario y
  suficiente que, para todo ideal máximo $\mathfrak m$ de $A$,
  $u_{\mathfrak m}:M_{\mathfrak m}\to N_{\mathfrak m}$ sea invertible por la
  izquierda.
  \item[(ii)] Sea $A'$ una $A$-álgebra que sea un $A$-módulo fielmente plano.
  Para que $u$ sea invertible por la izquierda, es necesario y suficiente que
  $u\otimes1:M\otimes_AA'\to N\otimes_AA'$ sea invertible por la izquierda.
\end{enumerate}
\end{proposition}

Como en (19.1.12), se puede limitar al caso en que $N$ es libre de tipo finito;
decir que $u$ es invertible por la izquierda significa entonces que $u$ es
inyectivo y que el módulo cociente $P=N/u(M)$ es \emph{proyectivo}, pues
$u(M)$ será entonces sumando directo de $N$. Señalemos además que, como $M$ es
de tipo finito, $P$ es de \emph{presentación finita}. Sentado esto:

(i) La condición es evidentemente necesaria. Recíprocamente, si se satisface,
se sabe que $u$ es inyectivo (Bourbaki, \emph{Alg. comm.}, cap.~II, §~3,
n\textsuperscript{o}~3, teorema~1), y como
$P_{\mathfrak m}=N_{\mathfrak m}/u_{\mathfrak m}(M_{\mathfrak m})$ es
proyectivo para todo $\mathfrak m$, se sabe que esto implica que $P$ es
proyectivo (\emph{loc. cit.}, §~5, n\textsuperscript{o}~2, teorema~1).

(ii) También aquí la condición es trivialmente necesaria. Recíprocamente, si
se cumple, se sabe que $u$ es inyectivo ($\mathbf{0}_{\mathrm I}$, 6.4.1), y
como $P\otimes_AA'=\operatorname{Coker}(u\otimes1)$ es proyectivo, y por tanto
plano, se deduce que $P$ es un $A$-módulo plano
($\mathbf{0}_{\mathrm I}$, 6.6.3), y por tanto proyectivo, puesto que es de
presentación finita (Bourbaki, \emph{Alg. comm.}, cap.~II, §~5,
n\textsuperscript{o}~2, cor.~2 del teorema~1).

\oldpage[0\textsubscript{IV-1}]{78}

\begin{remark}[19.1.15]
\label{0.19.1.15-es}
La noción «dual» a la de homomorfismo formalmente invertible por la izquierda
es la de homomorfismo \emph{formalmente invertible por la derecha}; tal
$A$-homomorfismo continuo $u:M\to N$ verifica, por definición, la condición
siguiente: para todo submódulo abierto $W'$ de $N$, existen un submódulo
abierto $V'\subset u^{-1}(W')$ de $M$, un submódulo abierto $W''\subset W'$ de
$N$ y un homomorfismo $h:N/W''\to M/V'$ tales que el homomorfismo canónico
$N/W''\to N/W'$ se factoriza como
\[
  N/W''\xrightarrow{h}M/V'\xrightarrow{u'}N/W',
\]
donde $u'$ se deduce de $u$ por paso a los cocientes. Esto implica que $u$ es
un epimorfismo formal; si existe un $\widehat A$-homomorfismo continuo
$r:\widehat N\to\widehat M$ tal que $\widehat u\circ r=1_{\widehat N}$, se
comprueba inmediatamente que $u$ es formalmente invertible por la derecha. Si
se cumplen las condiciones de (19.1.7), para que $u$ sea formalmente invertible
por la derecha es necesario y suficiente que los $u_\lambda$ sean invertibles
por la derecha (es decir, que el núcleo de $u_\lambda$ sea sumando directo de
$M/\mathfrak J_\lambda M$ y que $u_\lambda$ sea un isomorfismo de un
complementario de $\operatorname{Ker}(u_\lambda)$ sobre
$N/\mathfrak J_\lambda N$). Dejamos al lector el cuidado de enunciar y
demostrar las proposiciones correspondientes a (19.1.8) y (19.1.9) (en el
análogo de (19.1.8), se debe suponer que $M$ es separado y completo y que $M_n$
es un $(A/\mathfrak J_n)$-módulo proyectivo; en el análogo de (19.1.9), se debe
suprimir la hipótesis sobre los $N/\mathfrak J_\lambda N$, pero suponer, en
cambio, que, para todo $\lambda$, $M/\mathfrak J_\lambda M$ es un
$(A/\mathfrak J_\lambda)$-módulo proyectivo).
\end{remark}

\subsection{Módulos formalmente proyectivos}
\label{subsection:0.19.2-es}

\begin{definition}[19.2.1]
\label{0.19.2.1-es}
Sea $A$ un anillo topológico. Se dice que un $A$-módulo topológico $M$ es
\emph{formalmente proyectivo} si verifica la condición siguiente: para todo
ideal abierto $\mathfrak J$ de $A$, todo par de $(A/\mathfrak J)$-módulos
(discretos) $P$, $Q$, todo $A$-homomorfismo sobreyectivo $u:P\to Q$ y todo
$A$-homomorfismo continuo $v:M\to Q$, existe un $A$-homomorfismo continuo
$w:M\to P$ tal que $v=u\circ w$.
\end{definition}

\begin{env}[19.2.2]
\label{0.19.2.2-es}
Para verificar la condición de (19.2.1), es evidente que se puede limitar
(reemplazando $Q$ por $v(M)$ y $P$ por $u^{-1}(v(M))$) al caso en que el propio
$v$ es sobreyectivo; entonces $Q$ es isomorfo a $M/V$, donde $V$ es un
submódulo abierto de $M$ tal que $\mathfrak JM\subset V$; la condición de
(19.2.1) equivale entonces a decir que, para todo $(A/\mathfrak J)$-módulo discreto
$P$ y todo $A$-homomorfismo sobreyectivo $u:P\to M/V$, existe un submódulo
abierto $V'\subset V$ de $M$ y un $A$-homomorfismo $w:M/V'\to P$, tales que
el homomorfismo canónico $M/V'\to M/V$ se factoriza en
\[
  M/V'\xrightarrow{w}P\xrightarrow{u}M/V.
\]
Se observará que basta verificar esta condición cuando $\mathfrak J$
recorre un sistema fundamental de vecindades de 0 en $A$ formado por ideales.
\end{env}

\begin{env}[19.2.3]
\label{0.19.2.3-es}
Supongamos que existe un sistema fundamental $(\mathfrak J_\lambda)$ de ideales
abiertos de $A$ y un sistema fundamental $(V_\lambda)$ de submódulos
abiertos de $M$, con el \emph{mismo} conjunto de índices que
$(\mathfrak J_\lambda)$ y tal que, para todo $\lambda$, $M/V_\lambda$ sea un
$(A/\mathfrak J_\lambda)$-módulo \emph{proyectivo}. Entonces $M$ es formalmente
proyectivo: en efecto, con las notaciones de (19.2.2), basta tomar
$\lambda$ tal que $\mathfrak J_\lambda\subset\mathfrak J$ y
$V_\lambda\subset V$; como $P$ es también un
$(A/\mathfrak J_\lambda)$-módulo, la factorización del homomorfismo canónico
$M/V_\lambda\to M/V$ en $M/V_\lambda\to P\to M/V$ resulta de que se trata
de un $(A/\mathfrak J_\lambda)$-homomorfismo y de que se supone que
$M/V_\lambda$ es un $(A/\mathfrak J_\lambda)$-módulo proyectivo.

Cuando se satisface la condición más estricta de este número, se dice que $M$
es \emph{estrictamente formalmente proyectivo}.
\end{env}

\begin{proposition}[19.2.4]
\label{0.19.2.4-es}
Sean $A$ un anillo topológico y $M$ un $A$-módulo topológico cuya
topología se deduce de la de $A$. Para que $M$ sea formalmente
proyectivo, es necesario y suficiente que sea estrictamente formalmente proyectivo.
\end{proposition}

\oldpage[0\textsubscript{IV-1}]{79}

\begin{proof}
Acabamos de ver que la condición es suficiente. Recíprocamente, supongamos $M$
formalmente proyectivo y sea $(\mathfrak J_\lambda)$ un sistema fundamental
de ideales abiertos de $A$. Para todo $\lambda$, sean $P$ un
$(A/\mathfrak J_\lambda)$-módulo libre y
$u:P\to M/\mathfrak J_\lambda M$ un
$(A/\mathfrak J_\lambda)$-homomorfismo sobreyectivo. Existe, pues, un
$\mathfrak J_\mu\subset\mathfrak J_\lambda$ tal que el homomorfismo canónico
$M/\mathfrak J_\mu M\to M/\mathfrak J_\lambda M$ se factoriza en
\[
  M/\mathfrak J_\mu M\xrightarrow{w}P\xrightarrow{u}M/\mathfrak J_\lambda M;
\]
pero, como
$w(\mathfrak J_\lambda(M/\mathfrak J_\mu M))\subset\mathfrak J_\lambda P=0$,
$w$ se factoriza en $M/\mathfrak J_\lambda M\xrightarrow{v}P$, y está
claro que $u\circ v$ es la identidad en $M/\mathfrak J_\lambda M$, lo que
prueba que este $(A/\mathfrak J_\lambda)$-módulo es proyectivo.
\end{proof}

\begin{proposition}[19.2.5]
\label{0.19.2.5-es}
Sean $A$ un anillo topológico y $M$ un $A$-módulo topológico.
\begin{enumerate}
  \item[(i)] Para que $M$ sea formalmente proyectivo (resp. estrictamente
  formalmente proyectivo), es necesario y suficiente que lo sea el
  $\hat A$-módulo topológico $\hat M$.
  \item[(ii)] Sea $A'$ una $A$-álgebra topológica. Si $M$ es formalmente
  proyectivo (resp. estrictamente formalmente proyectivo), entonces
  $M\otimes_A A'$ (provisto de la topología producto tensorial) es un
  $A'$-módulo topológico formalmente proyectivo (resp. estrictamente formalmente
  proyectivo).
\end{enumerate}
\end{proposition}

\begin{proof}
\begin{enumerate}
  \item[(i)] Basta observar que cuando $\mathfrak J$ (resp. $V$)
  recorre el conjunto de los ideales abiertos de $A$ (resp. el conjunto de los
  submódulos abiertos de $M$), el completado separado $\hat{\mathfrak J}$
  (resp. $\hat V$) recorre el conjunto de los ideales abiertos de $\hat A$
  (resp. el conjunto de los submódulos abiertos de $\hat M$), y
  $\hat A/\hat{\mathfrak J}=A/\mathfrak J$ (resp.
  $\hat M/\hat V=M/V$) salvo isomorfismo canónico; como la noción
  de módulo formalmente proyectivo (resp. estrictamente formalmente proyectivo)
  solo hace intervenir los $A/\mathfrak J$ y los $M/V$, se deduce
  inmediatamente (i).
  \item[(ii)] Supongamos primero que $M$ es formalmente proyectivo y pongamos
  $M'=M\otimes_A A'$; sean $\mathfrak J'$ un ideal abierto de $A'$, $P'$,
  $Q'$ dos $(A'/\mathfrak J')$-módulos discretos, $u':P'\to Q'$ un
  $A'$-homomorfismo sobreyectivo y $v':M'\to Q'$ un $A'$-homomorfismo continuo.
  Hay un ideal abierto $\mathfrak J$ de $A$ tal que
  $\mathfrak JA'\subset\mathfrak J'$, por tanto $P'$ y $Q'$ son también
  $(A/\mathfrak J)$-módulos discretos. Si se considera el
  $A$-homomorfismo compuesto $v:M\xrightarrow{j}M'\xrightarrow{v'}Q'$, que es
  continuo, la hipótesis implica que existe un $A$-homomorfismo continuo
  $w:M\to P'$ tal que $v=u'\circ w$; pero, como $P'$ es un $A'$-módulo
  topológico, $w$ se factoriza en
  $M\xrightarrow{j}M'\xrightarrow{w'}P'$, donde $w'$ es un $A'$-homomorfismo
  continuo, y como $v'\circ j=(u'\circ w')\circ j$, se concluye que
  $v'=u'\circ w'$.

  Supongamos ahora que $M$ es estrictamente formalmente proyectivo; sea
  $(\mathfrak J_\lambda)$ (resp. $(V_\lambda)$) un sistema fundamental
  de ideales abiertos (resp. de submódulos abiertos) en $A$ (resp. $M$) tal
  que $(M/V_\lambda)$ sea un $(A/\mathfrak J_\lambda)$-módulo proyectivo,
  y sea $(\mathfrak J'_\mu)$ un sistema fundamental de ideales abiertos de
  $A'$. Se obtiene un sistema fundamental de vecindades de $0$ en $M'$ tomando
  los submódulos
  \[
    \operatorname{Im}(M\otimes_A\mathfrak J'_\mu)
    +\operatorname{Im}(V_\lambda\otimes_A A')=W_{\lambda\mu},
  \]
  donde $\lambda$ y $\mu$ son tales que
  $\mathfrak J_\lambda A'\subset\mathfrak J'_\mu$. Como
  \[
    (M\otimes_A A')/W_{\lambda\mu}
    =(M/V_\lambda)\otimes_{A/\mathfrak J_\lambda}
      (A'/\mathfrak J'_\mu)
  \]
  y como $(M/V_\lambda)$ es un
  $(A/\mathfrak J_\lambda)$-módulo proyectivo, $M'/W_{\lambda\mu}$
  es un $(A'/\mathfrak J'_\mu)$-módulo proyectivo.
\end{enumerate}
\end{proof}

\subsection{Álgebras formalmente lisas}
\label{subsection:0.19.3-es}

\begin{definition}[19.3.1]
\label{0.19.3.1-es}
Sean $A$ un anillo topológico y $B$ una $A$-álgebra topológica. Se dice que
$B$ es una $A$-álgebra \emph{formalmente lisa} si, para toda
$A$-álgebra topológica discreta $C$ y todo ideal nilpotente $\mathfrak J$ de
$C$, todo $A$-homomorfismo continuo $u:B\to C/\mathfrak J$ se factoriza en
\[
  B\xrightarrow{v}C\xrightarrow{\varphi}C/\mathfrak J,
\]
donde $v$ es un homomorfismo continuo y $\varphi$ el homomorfismo canónico.
\end{definition}

\oldpage[0\textsubscript{IV-1}]{80}

La definición (19.3.1) equivale a decir que se cumple la propiedad siguiente:
\begin{enumerate}
  \item[(P)] Para todo ideal abierto $\mathfrak R$ de la $A$-álgebra $B$ y
  todo $A$-homomorfismo $u':B/\mathfrak R\to C/\mathfrak J$, hay un
  ideal abierto $\mathfrak R'\subset\mathfrak R$ de $B$ tal que
  el homomorfismo
  $B\to B/\mathfrak R\xrightarrow{u'}C/\mathfrak J$ se factoriza en
  \[
    B\to B/\mathfrak R'\xrightarrow{v'}C\to C/\mathfrak J,
  \]
  donde $v'$ es un $A$-homomorfismo.
\end{enumerate}
En efecto, si $u:B\to C/\mathfrak J$ es un $A$-homomorfismo continuo, tiene
por núcleo $\mathfrak R$, un ideal abierto de $B$, por tanto $u$ se factoriza en
$B\to B/\mathfrak R\xrightarrow{u'}C/\mathfrak J$, y si se
verifica (P), basta aplicarla a $u'$ y tomar como $v$ el compuesto
$B\to B/\mathfrak R'\xrightarrow{v'}C$ para satisfacer las condiciones de
(19.3.1). Recíprocamente, supongamos que $B$ es una $A$-álgebra formalmente
lisa; démonos un ideal abierto $\mathfrak R$ de $B$ y un
$A$-homomorfismo $u':B/\mathfrak R\to C/\mathfrak J$, y apliquemos la
definición (19.3.1) a
$u:B\to B/\mathfrak R\xrightarrow{u'}C/\mathfrak J$; si
$v:B\to C$ es un $A$-homomorfismo continuo tal que $u$ se factoriza en
$B\xrightarrow{v}C\to C/\mathfrak J$, el ideal
$\mathfrak R'=\operatorname{Ker}(v)\cap\mathfrak R$ de $B$ es abierto y
está contenido en $\mathfrak R$; por consiguiente, $v$ se factoriza en
$B\to B/\mathfrak R'\xrightarrow{v'}C$, y $v'$ verifica ciertamente la condición
(P).

\begin{proposition}[19.3.2]
\label{0.19.3.2-es}
Sean $A$ un anillo discreto y $V$ un $A$-módulo proyectivo; el álgebra
simétrica $B=\mathbf S_A^\bullet(V)$, provista de la topología discreta, es una
$A$-álgebra formalmente lisa.
\end{proposition}

\begin{proof}
En efecto, teniendo las notaciones $C$ y $\mathfrak J$ el mismo significado que en
(19.3.1), sea $u:B\to C/\mathfrak J$ un homomorfismo de $A$-álgebras, que
por restricción a $V=\mathbf S_A^1(V)$ da un homomorfismo de $A$-módulos
$u_1:V\to C/\mathfrak J$; como $V$ es proyectivo, $u_1$ se factoriza en
$V\xrightarrow{v_1}C\xrightarrow{\varphi}C/\mathfrak J$, y $v_1$ se
prolonga a un homomorfismo de $A$-álgebras
$v:\mathbf S_A^\bullet(V)\to C$, tal que el compuesto
$\varphi\circ v$ coincide con $u$ en $V$, luego
$\varphi\circ v=u$, lo que demuestra la proposición.
\end{proof}

\begin{corollary}[19.3.3]
\label{0.19.3.3-es}
Si $A$ es un anillo discreto, toda álgebra de polinomios
$B=A[T_\alpha]_{\alpha\in I}$, provista de la topología discreta, es una
$A$-álgebra formalmente lisa.
\end{corollary}

\begin{proposition}[19.3.4]
\label{0.19.3.4-es}
Sea $A$ un anillo topológico, y sea
$B=A[[T_\alpha]]_{\alpha\in I}$ un anillo de series formales
$\sum_{\lambda\in\mathbb{N}^{(I)}}c_\lambda T^\lambda$ (álgebra amplia sobre
$A$ del monoide $\mathbb{N}^{(I)}$, identificada como $A$-módulo con el producto
$A^{\mathbb{N}^{(I)}}$). Si se dota $B$ de la topología producto, $B$ es una
$A$-álgebra formalmente lisa.
\end{proposition}

Sea $(\mathfrak L_\mu)_{\mu\in M}$ un sistema fundamental de ideales abiertos
de $A$. Para todo subconjunto finito $H$ de $I$, todo $\mu\in M$ y todo entero
$n$, designemos por $\mathfrak K_{H,\mu,n}$ la vecindad de $0$ en $B$ formada
por los $(c_\lambda)$ tales que, para todo
$\lambda=(\lambda_\alpha)_{\alpha\in I}$ tal que
$\lambda_\alpha=0$ para $\alpha\notin H$ y
$\sum_{\alpha\in H}\lambda_\alpha\leq n$, se tenga
$c_\lambda\in\mathfrak L_\mu$. Se verifica inmediatamente que los
$\mathfrak K_{H,\mu,n}$ forman un sistema fundamental de vecindades de $0$
en $B$ y son ideales de $B$, luego la topología producto es compatible
con la estructura de $A$-álgebra de $B$.

Observemos primero, con las mismas notaciones, el
\begin{lemma}[19.3.4.1]
\label{0.19.3.4.1-es}
Sea $E$ una $A$-álgebra discreta.
\begin{enumerate}
  \item[(i)] Si $f:B\to E$ es un $A$-homomorfismo continuo, existe un
  subconjunto finito $H$ de $I$ tal que $f(T_\alpha)=0$ para $\alpha\notin H$,
  y $f(T_\alpha)$ es nilpotente en $E$ para todo $\alpha\in H$.
  \item[(ii)] Recíprocamente, sean $H$ un subconjunto finito de $I$ y
  $(z_\alpha)_{\alpha\in H}$ una familia de elementos nilpotentes de $E$.
  Existe un único $A$-homomorfismo continuo $g:B\to E$ tal que
  $g(T_\alpha)=z_\alpha$ para $\alpha\in H$ y
  $g(T_\alpha)=0$ para $\alpha\notin H$.
\end{enumerate}
\end{lemma}

\oldpage[0\textsubscript{IV-1}]{81}

\begin{proof}
(i) resulta inmediatamente de que $f^{-1}(0)$ es una vecindad de $0$ en el
$A$-módulo producto $A^{\mathbb{N}^{(I)}}$, de donde $f(T^\lambda)=0$ salvo para un
número finito de valores de $\lambda\in\mathbb{N}^{(I)}$. Para probar (ii),
basta observar que el anillo de polinomios
$B'=A[T_\alpha]_{\alpha\in I}$ es denso en $B$; la existencia y
la unicidad de la restricción $g|B'$ son triviales y su continuidad resulta de
la hipótesis de que los $f(T_\alpha)$ para $\alpha\in H$ son nilpotentes,
pues si $(f(T_\alpha))^n=0$ para todo $\alpha\in H$, se tiene
$g(T^\lambda)=0$ para todo $\lambda=(\lambda_\alpha)_{\alpha\in I}$ de
soporte finito salvo aquellos tales que $\lambda_\alpha=0$ para
$\alpha\notin H$ y $\lambda_\alpha\leq n$ para
$\alpha\in H$, es decir, salvo para un número finito de valores de
$\lambda$.
\end{proof}

\begin{proof}[Demostración de la proposición~{\rm(19.3.4)}]
Establecido este lema, y con las notaciones $C$ y $\mathfrak J$ en el mismo
sentido que en (19.3.1), se tiene $u(T_\alpha)=0$ salvo para
$\alpha\in H$, siendo $H$ una parte finita de $I$, y los
$z_\alpha=u(T_\alpha)$ para $\alpha\in H$ son nilpotentes en
$C/\mathfrak J$; como $\mathfrak J$ es nilpotente, existe una familia
$(x_\alpha)_{\alpha\in H}$ de elementos nilpotentes de $C$ cuyas imágenes
canónicas en $C/\mathfrak J$ son los $z_\alpha$; si $v$ es el
$A$-homomorfismo continuo de $B$ en $C$ tal que
$v(T_\alpha)=0$ para $\alpha\notin H$, $v(T_\alpha)=x_\alpha$ para
$\alpha\in H$, está claro que $u$ se factoriza como
$B\xrightarrow{v}C\to C/\mathfrak J$.
\end{proof}

\begin{proposition}[19.3.5]
\label{0.19.3.5-es}
Sea $A$ un anillo topológico.
\begin{enumerate}
  \item[(i)] $A$ es una $A$-álgebra formalmente lisa.
  \item[(ii)] Si $B$ es una $A$-álgebra formalmente lisa y $C$ una
  $B$-álgebra formalmente lisa, entonces $C$ es una $A$-álgebra formalmente
  lisa.
  \item[(iii)] Sean $B$ una $A$-álgebra formalmente lisa, $A'$ una
  $A$-álgebra topológica; entonces la $A'$-álgebra topológica
  $B\otimes_A A'$ ($\mathbf 0_{\rm I}$, 7.7.5 y 7.7.6) es formalmente
  lisa.
  \item[(iv)] Sean $B$ una $A$-álgebra topológica, $S$ (resp. $T$) un
  subconjunto multiplicativo de $A$ (resp. $B$) tal que la imagen canónica de $S$
  en $B$ esté contenida en $T$. Si $B$ es una $A$-álgebra formalmente
  lisa, entonces $T^{-1}B$ es una $S^{-1}A$-álgebra formalmente lisa.
  \item[(v)] Sean $B_i$ ($1\leq i\leq n$) unas álgebras topológicas sobre $A$.
  Para que $\prod_{i=1}^n B_i$ sea una $A$-álgebra formalmente lisa, es
  necesario y suficiente que cada una de las $B_i$ lo sea.
\end{enumerate}
\end{proposition}

\oldpage[0\textsubscript{IV-1}]{82}

\begin{proof}
\begin{enumerate}
  \item[(i)] Si $C$ es una $A$-álgebra discreta,
  $\varphi:C\to C/\mathfrak J$ el homomorfismo canónico de $C$ sobre un
  $A$-álgebra cociente cualquiera de $C$, el \emph{único} $A$-homomorfismo
  de $A$ en $C/\mathfrak J$ es el homomorfismo compuesto
  $A\xrightarrow{\psi}C\xrightarrow{\varphi}C/\mathfrak J$, donde
  $\psi$ es el homomorfismo que define la estructura de $A$-álgebra de
  $C$; como $\psi$ es continuo, la condición de (19.3.1) se verifica
  trivialmente.
  \item[(ii)] Sean $\alpha:A\to B$, $\beta:B\to C$ los homomorfismos
  continuos que definen respectivamente la estructura de $A$-álgebra sobre $B$ y
  la de $B$-álgebra sobre $C$, de modo que $\beta\circ\alpha$ define la
  estructura de $A$-álgebra sobre $C$. Sean $E$ una $A$-álgebra discreta,
  $\mathfrak L$ un ideal nilpotente de $E$, $u:C\to E/\mathfrak L$ un
  $A$-homomorfismo continuo, de modo que $u\circ\beta\circ\alpha$ es
  el homomorfismo que define la estructura de $A$-álgebra de
  $E/\mathfrak L$. Como $B$ es una $A$-álgebra formalmente lisa, el
  $A$-homomorfismo continuo $u\circ\beta:B\to E/\mathfrak L$ se
  factoriza como $B\xrightarrow{v}E\to E/\mathfrak L$, donde $v$ es un
  $A$-homomorfismo continuo; $v$ y $u\circ\beta$ definen entonces sobre
  $E$ y $E/\mathfrak L$ respectivamente unas estructuras de $B$-álgebra
  topológica, para las cuales $E/\mathfrak L$ sigue siendo el álgebra cociente
  (discreta) de la $B$-álgebra $E$. Además, $u$ es un $B$-homomorfismo
  continuo, por lo que se factoriza como $C\xrightarrow{w}E\to E/\mathfrak L$,
  donde $w$ es un $B$-homomorfismo continuo; como
  $v\circ\alpha$ es el $A$-homomorfismo que define la estructura de
  $A$-álgebra sobre $E$, $w$ es efectivamente un $A$-homomorfismo continuo, de donde nuestra
  afirmación.
  \item[(iii)] Sean $C$ una $A'$-álgebra topológica discreta,
  $\mathfrak J$ un ideal nilpotente de $C$,
  $u:B\otimes_A A'\to C/\mathfrak J$ un $A'$-homomorfismo continuo. Si se
  compone $u$ con el homomorfismo canónico
  $\rho:B\to B\otimes_A A'$, se obtiene (puesto que
  $A\to B\xrightarrow{\rho}B\otimes_A A'$ es igual al compuesto
  $A\to A'\xrightarrow{\sigma}B\otimes_A A'$) un $A$-homomorfismo
  continuo que, por hipótesis, se factoriza, por tanto, como
  $B\xrightarrow{v}C\to C/\mathfrak J$, donde $v$ es un $A$-homomorfismo
  continuo. La igualdad de los compuestos $A\to B\xrightarrow{v}C$ y
  $A\to A'\xrightarrow{w}C$ implica, pues, la existencia de un homomorfismo
  continuo de anillos $f:B\otimes_A A'\to C$ tal que
  $v=f\circ\rho$ y $w=f\circ\sigma$ ($\mathbf 0_{\rm I}$, 7.7.6);
  como los homomorfismos compuestos
  $B\xrightarrow{\rho}B\otimes_A A'\xrightarrow{f}C\to C/\mathfrak J$
  y
  $B\xrightarrow{\rho}B\otimes_A A'\xrightarrow{u}C/\mathfrak J$
  (resp.
  $A'\xrightarrow{\sigma}B\otimes_A A'\xrightarrow{f}C\to C/\mathfrak J$
  y
  $A'\xrightarrow{\sigma}B\otimes_A A'\xrightarrow{u}C/\mathfrak J$)
  son iguales, se tiene efectivamente la factorización
  $u:B\otimes_A A'\xrightarrow{f}C\to C/\mathfrak J$, lo que demuestra
  (iii).
  \item[(iv)] La topología considerada sobre $S^{-1}A$ (resp. $T^{-1}B$) es
  naturalmente aquella para la cual un sistema fundamental de entornos de
  $0$ está formado por los $S^{-1}\mathfrak J$ (resp. $T^{-1}\mathfrak R$), donde
  $\mathfrak J$ (resp. $\mathfrak R$) recorre un sistema fundamental
  de ideales abiertos en $A$ (resp. $B$) ($\mathbf 0_{\rm I}$, 7.6.1). Si
  $\alpha:A\to B$ es el homomorfismo canónico, está claro que
  el homomorfismo canónico $\alpha':S^{-1}A\to T^{-1}B$ inducido por
  $\alpha$ es continuo. Sean entonces $C$ una $S^{-1}A$-álgebra topológica
  discreta, $\mathfrak J$ un ideal nilpotente de esta álgebra,
  $u:T^{-1}B\to C/\mathfrak J$ un $S^{-1}A$-homomorfismo continuo; entonces
  el compuesto $B\to T^{-1}B\xrightarrow{u}C/\mathfrak J$ es un
  $A$-homomorfismo continuo que, por hipótesis, se factoriza como
  $B\xrightarrow{v}C\to C/\mathfrak J$, donde $v$ es un $A$-homomorfismo
  continuo. Como para todo $t\in T$, $t/1$ es invertible en $T^{-1}B$,
  $u(t/1)$ es invertible en $C/\mathfrak J$. Como $\mathfrak J$ es
  nilpotente, todo representante de la clase $u(t/1)$ en $C$, y en particular
  $v(t)$, es invertible en $C$ ($\mathbf 0_{\rm I}$, 7.1.12), y por consiguiente
  $v$ se factoriza como $B\to T^{-1}B\xrightarrow{w}C$; como $v$ es
  continuo, lo mismo ocurre con $w$ ($\mathbf 0_{\rm I}$, 7.6.6), y este es
  un $S^{-1}A$-homomorfismo porque el compuesto
  $A\to B\to T^{-1}B\xrightarrow{w}C$ es igual a
  \[
    A\to S^{-1}A\xrightarrow{\alpha'}T^{-1}B\xrightarrow{w}C,
  \]
  por lo que $S^{-1}A\xrightarrow{\alpha'}T^{-1}B\xrightarrow{w}C$ es
  el homomorfismo canónico que define sobre $C$ la estructura de
  $S^{-1}A$-álgebra. Finalmente, al ser iguales los homomorfismos compuestos
  $B\to T^{-1}B\xrightarrow{w}C\to C/\mathfrak J$ y
  $B\to T^{-1}B\xrightarrow{u}C/\mathfrak J$, el mismo
  razonamiento muestra que $u$ es efectivamente igual al compuesto
  $T^{-1}B\xrightarrow{w}C\to C/\mathfrak J$, de donde (iv).
  \item[(v)] Finalmente, (v) es inmediato, pues el dato de un $A$-homomorfismo
  continuo de $B=\prod_{i=1}^n B_i$ en $C$ (resp. $C/\mathfrak J$)
  equivale al de $n$ $A$-homomorfismos continuos $B_i\to C$
  (resp. $B_i\to C/\mathfrak J$), y todo $A$-homomorfismo continuo
  $B_j\to C$ (resp. $B_j\to C/\mathfrak J$) da por composición un
  $A$-homomorfismo continuo $B\to B_j\to C$
  (resp. $B\to B_j\to C/\mathfrak J$).
\end{enumerate}
\end{proof}

\oldpage[0\textsubscript{IV-1}]{83}

\begin{proposition}[19.3.6]
\label{0.19.3.6-es}
Sean $A$ un anillo topológico, $B$ una $A$-álgebra topológica, $\hat A$
y $\hat B$ los respectivos completados separados de $A$ y $B$, de modo que
$\hat B$ es una $\hat A$-álgebra topológica. Las condiciones siguientes
son equivalentes:
\begin{enumerate}
  \item[a)] $B$ es una $A$-álgebra formalmente lisa.
  \item[b)] $\hat B$ es una $A$-álgebra formalmente lisa.
  \item[c)] $\hat B$ es una $\hat A$-álgebra formalmente lisa.
\end{enumerate}
\end{proposition}

\begin{proof}
Por supuesto, la estructura de $\hat A$-álgebra sobre $\hat B$ está definida
por el homomorfismo $\hat\varphi$, si $\varphi:A\to B$ es
el homomorfismo que define la estructura de $A$-álgebra sobre $B$. Como toda
$A$-álgebra discreta $C$ es separada y completa, también es una
$\hat A$-álgebra (por prolongación canónica del homomorfismo de $A$ en
$C$), y un $A$-homomorfismo continuo de $B$ en $C$ da por prolongación
canónica un $\hat A$-homomorfismo (y \emph{a fortiori} un
$A$-homomorfismo) de $\hat B$ en $C$ (en otros términos, todo
$A$-homomorfismo continuo $B\to C$ se factoriza como $B\to\hat B\to C$ de manera
única). Estas observaciones y la definición (19.3.1) implican inmediatamente la
proposición.
\end{proof}

\begin{corollary}[19.3.7]
\label{0.19.3.7-es}
Bajo las hipótesis de (19.3.5), (iv), la $A\{S^{-1}\}$-álgebra topológica
$B\{T^{-1}\}$ es formalmente lisa.
\end{corollary}

Esto resulta de las definiciones ($\mathbf 0_{\rm I}$, 7.6.1 y 7.6.7) y de
(19.3.5), (iv) y (19.3.6).

\begin{proposition}[19.3.8]
\label{0.19.3.8-es}
Sean $A$ un anillo topológico, $B$ una $A$-álgebra topológica, y supongamos
que para todo ideal abierto $\mathfrak R$ de $B$, $\mathfrak R^2$ sea también
abierto. Sean $A'$, $B'$ los anillos topológicos obtenidos dotando a $A$ y
$B$ de topologías más finas que las topologías iniciales, y supongamos que
el homomorfismo canónico $\varphi:A\to B$ siga siendo un homomorfismo
continuo de $A'$ en $B'$. Entonces, si $B'$ es una $A'$-álgebra formalmente
lisa, $B$ es una $A$-álgebra formalmente lisa.
\end{proposition}

Basta aplicar el lema siguiente:

\begin{lemma}[19.3.8.1]
\label{0.19.3.8.1-es}
Sean $C$ una $A$-álgebra discreta, $\mathfrak J$ un ideal nilpotente de $C$.
Supongamos que el cuadrado de todo ideal abierto de $B$ sea abierto. Entonces, si un
homomorfismo compuesto
$v:B\xrightarrow{u}C\to C/\mathfrak J$ es continuo, el homomorfismo $u$
es continuo.
\end{lemma}

\begin{proof}
En efecto, $\mathfrak R=u^{-1}(\mathfrak J)$ es el núcleo de $v$ y es
abierto por hipótesis; como existe un entero $n$ tal que
$\mathfrak J^n=0$, $\mathfrak R^n$ está contenido en
$\operatorname{Ker}(u)$; pero, por inducción sobre $h$, se deduce de
la hipótesis que $\mathfrak R^{2^h}$ es abierto, por lo que también lo es
$\mathfrak R^n$ para todo $n$, y por consiguiente $\operatorname{Ker}(u)$ es
asimismo abierto, lo que demuestra nuestra afirmación.
\end{proof}

Se observará que la hipótesis sobre $B$ significa que la topología de $B$ es el
supremo de las topologías $\mathfrak R$-preádicas, donde $\mathfrak R$ recorre
el conjunto de los ideales abiertos de $B$. Si $B$ es un anillo preádico
($\mathbf 0_{\rm I}$, 7.1.9), esta condición se verifica siempre.

\begin{env}[19.3.9]
\label{0.19.3.9-es}
Vamos a ver ahora que la propiedad de ser formalmente lisa
implica propiedades de «~levantamiento~» de homomorfismos bajo condiciones
más generales que las de la definición (19.3.1).
\end{env}

\oldpage[0\textsubscript{IV-1}]{84}

\begin{proposition}[19.3.10]
\label{0.19.3.10-es}
Sean $A$ un anillo topológico, $B$ una $A$-álgebra formalmente lisa.
Sean $C$ una $A$-álgebra topológica, $\mathfrak J$ un ideal de $C$,
que verifican las condiciones siguientes:
\begin{enumerate}
  \item[$1^\circ$] $C$ es metrizable y completo.
  \item[$2^\circ$] $\mathfrak J$ es cerrado y la sucesión
  $(\mathfrak J^n)$ tiende a $0$.
\end{enumerate}
Entonces, todo $A$-homomorfismo continuo $u:B\to C/\mathfrak J$ se factoriza como
$B\xrightarrow{v}C\to C/\mathfrak J$, donde $v$ es un
$A$-homomorfismo continuo.
\end{proposition}

\begin{proof}
Sea $(\mathfrak Q_n)$ una sucesión decreciente de ideales de $C$ que forma un
sistema fundamental de entornos de $0$. Para todo $n$, consideremos el
$A$-álgebra discreta $C/\mathfrak Q_n$ y el ideal
$(\mathfrak J+\mathfrak Q_n)/\mathfrak Q_n$ de esta álgebra; como
existe $k$ tal que $\mathfrak J^k\subset\mathfrak Q_n$,
$(\mathfrak J+\mathfrak Q_n)/\mathfrak Q_n$ es nilpotente. Para todo
$n$, sea $u_n$ el homomorfismo continuo
\[
  B\xrightarrow{u}C/\mathfrak J\to C/(\mathfrak J+\mathfrak Q_n)
  =(C/\mathfrak Q_n)/((\mathfrak J+\mathfrak Q_n)/\mathfrak Q_n)~;
\]
por hipótesis $u_n$ se factoriza como
\[
  B\xrightarrow{v_n}C/\mathfrak Q_n
  \xrightarrow{\varphi_n}C/(\mathfrak J+\mathfrak Q_n),
\]
donde $v_n$ es un $A$-homomorfismo continuo. Mostremos que se pueden elegir de
manera sucesiva los $v_n$ de modo que $v_n$ se factorice como
\[
  B\xrightarrow{v_{n+1}}C/\mathfrak Q_{n+1}\to C/\mathfrak Q_n
\]
para todo $n$ (es decir, que los $v_n$ formen un \emph{sistema
proyectivo} de homomorfismos). Esto resultará del lema siguiente.
\end{proof}

\begin{lemma}[19.3.10.1]
\label{0.19.3.10.1-es}
Sean $B$ una $A$-álgebra formalmente lisa, $E$, $E'$ dos $A$-álgebras
topológicas discretas, $\mathfrak R$ (resp. $\mathfrak R'$) un ideal
nilpotente de $E$ (resp. $E'$), $f:E\to E'$ un $A$-homomorfismo sobreyectivo tal
que $f(\mathfrak R)\subset\mathfrak R'$, $f':E/\mathfrak R\to
E'/\mathfrak R'$ el homomorfismo deducido de $f$ por paso a los cocientes.
Sean $u:B\to E/\mathfrak R$ un $A$-homomorfismo continuo, $u'$ el
homomorfismo compuesto
$B\xrightarrow{u}E/\mathfrak R\xrightarrow{f'}E'/\mathfrak R'$, y sea
$v':B\to E'$ un $A$-homomorfismo continuo tal que $u'$ se factorice como
$B\xrightarrow{v'}E'\to E'/\mathfrak R'$. Entonces existe un
$A$-homomorfismo continuo $v:B\to E$ tal que $v'$ se factorice como
$B\xrightarrow{v}E\xrightarrow{f}E'$.
\end{lemma}

En el diagrama conmutativo
\[
\xymatrix{
  E \ar[r]^f \ar[d] & E' \ar[d] \\
  E/\mathfrak R \ar[r]_{f'} & E'/\mathfrak R'
}
\]
todos los homomorfismos son sobreyectivos; si $\mathfrak L$ es el núcleo de
$f$, $E'$ se identifica con $E/\mathfrak L$ y $\mathfrak R'$ con
$(\mathfrak R+\mathfrak L)/\mathfrak L$. Utilicemos ahora el lema
elemental siguiente.

\begin{lemma}[19.3.10.2]
\label{0.19.3.10.2-es}
Sean $F$ un anillo (no necesariamente conmutativo), $\mathfrak a$,
$\mathfrak b$ dos ideales bilaterales de $F$; entonces el producto fibrado
\[
  (F/\mathfrak a)\times_{F/(\mathfrak a+\mathfrak b)}(F/\mathfrak b)
\]
se identifica canónicamente con $F/(\mathfrak a\cap\mathfrak b)$.
\end{lemma}

\begin{proof}
No es sino una aplicación particular de (18.1.7), donde, en el
diagrama (18.1.7.1), se reemplaza $G$ por
$F/(\mathfrak a\cap\mathfrak b)$, $\mathfrak J'$ por
$\mathfrak b/(\mathfrak a\cap\mathfrak b)$, $F$ por
$F/\mathfrak b$, $\mathfrak R$ por
$(\mathfrak a+\mathfrak b)/\mathfrak b$, $B$ por
$F/(\mathfrak a+\mathfrak b)$ y $\theta$ por el $F$-homomorfismo
biyectivo canónico
$(\mathfrak a+\mathfrak b)/\mathfrak b\to
\mathfrak a/(\mathfrak a\cap\mathfrak b)$.
\end{proof}

\oldpage[0\textsubscript{IV-1}]{85}

\begin{proof}[Continuación de la demostración del lema~{\rm(19.3.10.1)}]
Aplicando este lema a la situación de (19.3.10.1), la conmutatividad del
diagrama
\[
\xymatrix{
  B \ar[r]^{v'} \ar[d]_u & E'=E/\mathfrak L \ar[d] \\
  E/\mathfrak R \ar[r]_{f'} & E'/\mathfrak R'
}
\]
muestra la existencia de un homomorfismo
$w:B\to E/(\mathfrak R\cap\mathfrak L)$ tal que $v'$ y $u$ se
factorizan como
$B\xrightarrow{w}E/(\mathfrak R\cap\mathfrak L)\to E/\mathfrak L$
y
$B\xrightarrow{w}E/(\mathfrak R\cap\mathfrak L)\to E/\mathfrak R$
respectivamente; como el núcleo de $w$ contiene la intersección de los de
$v'$ y $u$, es abierto en $B$ y $w$ es continuo. Finalmente, es claro
que $\mathfrak R\cap\mathfrak L$ es nilpotente, luego $w$ se factoriza como
$B\xrightarrow{v}E\to E/(\mathfrak R\cap\mathfrak L)$, donde $v$ es
un $A$-homomorfismo continuo que satisface lo requerido.
\end{proof}

\begin{proof}[Continuación de la demostración de la proposición~{\rm(19.3.10)}]
La existencia de $v_{n+1}$ resulta del lema (19.3.10.1) aplicado reemplazando
$E$, $E'$, $\mathfrak R$, $\mathfrak R'$ por
$C/\mathfrak Q_{n+1}$, $C/\mathfrak Q_n$,
$(\mathfrak J+\mathfrak Q_{n+1})/\mathfrak Q_{n+1}$ y
$(\mathfrak J+\mathfrak Q_n)/\mathfrak Q_n$ respectivamente, $u$, $u'$
y $v'$ por $u_{n+1}$, $u_n$ y $v_n$ respectivamente. Observemos que, como $C$
es separado y completo, es igual a $\varprojlim_n(C/\mathfrak Q_n)$, y
$v=\varprojlim_n v_n$ es un $A$-homomorfismo continuo de $B$ en $C$.
Como $\mathfrak J$ está cerrado en $C$, se tiene
$\mathfrak J=\bigcap_n(\mathfrak J+\mathfrak Q_n)$; como
$C/\mathfrak J$ es metrizable y completo y los
$(\mathfrak J+\mathfrak Q_n)/\mathfrak J$ forman un sistema fundamental
de entornos de $0$ en $C/\mathfrak J$, se tiene
\[
  C/\mathfrak J=\varprojlim_n C/(\mathfrak J+\mathfrak Q_n)
\]
y $\varprojlim_n\varphi_n$ es la aplicación canónica
$\varphi:C\to C/\mathfrak J$. Como $\varprojlim_n u_n=u$, se tiene en efecto
$u=v\circ\varphi$, C.Q.F.D.
\end{proof}

\begin{corollary}[19.3.11]
\label{0.19.3.11-es}
Sean $A$ un anillo topológico, $B$ una $A$-álgebra formalmente lisa,
$C$ un anillo local noetheriano completo, $\mathfrak J$ un ideal propio de
$C$, $\varphi:A\to C$ un homomorfismo continuo que hace de $C$ una
$A$-álgebra topológica. Entonces todo $A$-homomorfismo continuo
$u_0:B\to C_0=C/\mathfrak J$ se factoriza como
$B\xrightarrow{u}C\to C/\mathfrak J$, donde $u$ es un $A$-homomorfismo
continuo.
\end{corollary}

En efecto, se cumplen todas las condiciones de (19.3.10)
($\mathbf 0_{\rm I}$, 7.3.5).

\begin{proposition}[19.3.12]
\label{0.19.3.12-es}
Sean $A$ un anillo topológico, $B$ una $A$-álgebra formalmente lisa,
$C$ una $A$-álgebra topológica, $\mathfrak J$ un ideal de $C$; supongamos que se verifican las
condiciones siguientes:
\begin{enumerate}
  \item[$1^\circ$] Existe un sistema fundamental de ideales abiertos
  $\mathfrak Q_\lambda$ de $C$, tales que los
  $C_\lambda=C/\mathfrak Q_\lambda$ sean anillos artinianos y que
  el homomorfismo canónico $C\to\varprojlim_\lambda C_\lambda$ sea un
  isomorfismo de anillos topológicos.
  \item[$2^\circ$] El ideal $\mathfrak J$ está cerrado en $C$ y es
  topológicamente nilpotente.
  \item[$3^\circ$] El cuadrado de todo ideal abierto de $B$ es abierto.
\end{enumerate}
En estas condiciones, todo $A$-homomorfismo continuo
$u:B\to C/\mathfrak J$ se factoriza como
$B\xrightarrow{v}C\to C/\mathfrak J$, donde $v$ es un $A$-homomorfismo
continuo.
\end{proposition}

\begin{proof}
Sea $\mathfrak J_\lambda=(\mathfrak J+\mathfrak Q_\lambda)/\mathfrak Q_\lambda$
la imagen canónica de $\mathfrak J$ en $C_\lambda$, que es un ideal
nilpotente, y sea $u_\lambda$ el homomorfismo compuesto
$B\xrightarrow{u}C/\mathfrak J\to C_\lambda/\mathfrak J_\lambda
=C/(\mathfrak J+\mathfrak Q_\lambda)$. Mostremos que todo
$A$-homomorfismo $w:B\to C$ tal que $u$ se factorice como
$B\xrightarrow{w}C\to C/\mathfrak J$ es necesariamente continuo; en
efecto, para todo $\lambda$, existe un ideal abierto $\mathfrak R$ de $B$
tal que $u(\mathfrak R)\subset
(\mathfrak J+\mathfrak Q_\lambda)/\mathfrak J$, de donde
$w(\mathfrak R)\subset\mathfrak J+\mathfrak Q_\lambda$, luego existe
una potencia $\mathfrak R^{2^m}$ de $\mathfrak R$ tal que
$w(\mathfrak R^{2^m})\subset\mathfrak Q_\lambda$, lo que demuestra
nuestra afirmación, puesto que $\mathfrak R^{2^m}$ es abierto en $B$. Se puede
por tanto limitar a encontrar un $A$-homomorfismo $w$ con la propiedad anterior
sin preocuparse de sus propiedades de continuidad. Ahora bien, el conjunto
$\operatorname{Hom}(B,C)$ de todos los $A$-homomorfismos de álgebras de $B$
en $C$ es igual a
$\varprojlim_\lambda\operatorname{Hom}(B,C_\lambda)$, y este último es
\emph{cerrado} en el $C_\lambda$-módulo $C_\lambda^B$, provisto de la topología
producto, para la cual es \emph{linealmente compacto}, puesto que
  $C_\lambda$ es artiniano. Para todo $\lambda$, sea $W_\lambda$

\oldpage[0\textsubscript{IV-1}]{86}

el conjunto de los $w_\lambda\in\operatorname{Hom}(B,C_\lambda)$ tales que
$u_\lambda$ se factorice como
$B\xrightarrow{w_\lambda}C_\lambda\to
C_\lambda/\mathfrak J_\lambda$; $W_\lambda$ es una variedad
lineal cerrada en el $C_\lambda$-módulo
  $\operatorname{Hom}(B,C_\lambda)$, no vacía puesto que $B$ es formalmente
  lisa. En el producto
$\prod_{\mu\leq\lambda}\operatorname{Hom}(B,C_\mu)=E_\lambda$,
consideremos la subvariedad lineal $U_\lambda$ formada por los
$(w_\mu)_{\mu\leq\lambda}$ tales que
$w_\mu=\varphi_{\mu\lambda}\circ w_\lambda$ para
$\mu\leq\lambda$ y $w_\lambda\in W_\lambda$, donde
$\varphi_{\mu\lambda}:C_\lambda\to C_\mu$ es
el homomorfismo canónico; esta subvariedad también es cerrada. Finalmente, sea $V_\lambda$
la variedad lineal en el producto
$\prod_\mu\operatorname{Hom}(B,C_\mu)$, producto de $U_\lambda$
y de los $\operatorname{Hom}(B,C_\mu)$ para
$\mu\not\leq\lambda$, que sigue siendo cerrada. Todo se reduce a ver que
la intersección de los $V_\lambda$ es \emph{no vacía}, pues un elemento $w$ de
esta intersección pertenecerá a
$\varprojlim_\lambda\operatorname{Hom}(B,C_\lambda)$ por definición.
Además, como $\mathfrak J$ está cerrado y
$C=\varprojlim_\lambda C_\lambda$ es linealmente compacto,
$C/\mathfrak J$ se identifica con
$\varprojlim_\lambda C_\lambda/\mathfrak J_\lambda$ y
la aplicación canónica $\psi:C\to C/\mathfrak J$ con
$\varprojlim_\lambda\psi_\lambda$, donde $\psi_\lambda$ es
la aplicación canónica
$C_\lambda\to C_\lambda/\mathfrak J_\lambda$. Se concluye, por tanto,
como en (19.3.10), que $\psi\circ w=u$.
\end{proof}

\subsection{Primeros criterios de lisura formal}
\label{subsection:0.19.4-es}

\begin{proposition}[19.4.1]
\label{0.19.4.1-es}
Sean $A$ un anillo topológico, $B$ una $A$-álgebra topológica; supongamos
que existen dos familias filtrantes decrecientes
$(\mathfrak J_\alpha)_{\alpha\in I}$,
$(\mathfrak R_\alpha)_{\alpha\in I}$ de ideales de $A$ y $B$
respectivamente, tales que:
\begin{enumerate}
  \item[$1^\circ$] $(\mathfrak J_\alpha)$ tiende a $0$ en $A$ y
  $(\mathfrak R_\alpha)$ tiende a $0$ en $B$;
  \item[$2^\circ$] para todo $\alpha\in I$ se tiene
  $\mathfrak J_\alpha B\subset\mathfrak R_\alpha$ (de modo que
  $B/\mathfrak R_\alpha$ es una
  $(A/\mathfrak J_\alpha)$-álgebra topológica);
  \item[$3^\circ$] para todo $\alpha\in I$,
  $B/\mathfrak R_\alpha$ es una
  $(A/\mathfrak J_\alpha)$-álgebra formalmente lisa.
\end{enumerate}
Entonces $B$ es una $A$-álgebra formalmente lisa.
\end{proposition}

\begin{proof}
En efecto, sean $C$ una $A$-álgebra discreta, $\mathfrak L$ un ideal
nilpotente de $C$, $\mathfrak R$ un ideal abierto de $B$; por hipótesis existe
un $\alpha\in I$ tal que
$\mathfrak R_\alpha\subset\mathfrak R$, luego $B/\mathfrak R$ es
cociente de $B/\mathfrak R_\alpha$ por un ideal abierto. Todo
$A$-homomorfismo $u':B/\mathfrak R\to C/\mathfrak L$ es también un
$(A/\mathfrak J_\alpha)$-homomorfismo, luego existe un ideal abierto
$\mathfrak R'$ de $B$ tal que
$\mathfrak R_\alpha\subset\mathfrak R'\subset\mathfrak R$ y que
$B/\mathfrak R_\alpha\to B/\mathfrak R\xrightarrow{u'}C/\mathfrak L$
se factorice como
$B/\mathfrak R_\alpha\to B/\mathfrak R'\to C\to C/\mathfrak L$,
donde el segundo homomorfismo es un
$(A/\mathfrak J_\alpha)$-homomorfismo; de donde la conclusión, en virtud de
la observación que sigue a (19.3.1).
\end{proof}

\begin{corollary}[19.4.2]
\label{0.19.4.2-es}
Sean $A$ un anillo topológico, $B$ una $A$-álgebra topológica,
$(\mathfrak J_\alpha)_{\alpha\in I}$ una familia filtrante decreciente
de ideales de $A$ que tiende a $0$. Para que $B$ sea una $A$-álgebra
formalmente lisa, es necesario y suficiente que, para todo $\alpha\in I$, si
se pone $A_\alpha=A/\mathfrak J_\alpha$,
$B_\alpha=B/\mathfrak J_\alpha B$, $B_\alpha$ sea una
$A_\alpha$-álgebra formalmente lisa.
\end{corollary}

La condición es suficiente por (19.4.1), y también es necesaria por
(19.3.5), (iii).

\begin{proposition}[19.4.3]
\label{0.19.4.3-es}
Sean $A$ un anillo topológico y $B$ una $A$-álgebra topológica. Supongamos
que, para toda $A$-álgebra discreta $C$ y todo ideal $\mathfrak J$ de $C$
tal que $\mathfrak J^2=0$, todo $A$-homomorfismo continuo
$u:B\to C/\mathfrak J$ se factoriza en
$B\xrightarrow{v}C\to C/\mathfrak J$, donde $v$ es un $A$-homomorfismo
continuo. Entonces $B$ es una $A$-álgebra formalmente lisa.
\end{proposition}

\begin{proof}
En efecto, con las mismas notaciones, sea $\mathfrak R$ un ideal nilpotente
cualquiera de $C$, y consideremos un $A$-homomorfismo continuo
$u':B\to C/\mathfrak R$. Supongamos que $\mathfrak R^m=0$, y pongamos
$C_i=C/\mathfrak R^i$ para $1\leq i\leq m$, de modo que $C_m=C$, que
$\mathfrak R^{i-1}/\mathfrak R^i$ es un ideal de cuadrado nulo
$\mathfrak J_i$ en $C_i$, y $C_i=C_{i+1}/\mathfrak J_{i+1}$;
la hipótesis implica entonces, por recurrencia sobre $i$, la existencia de
$A$-homomorfismos continuos $u_i:B\to C_i$ tales que $u_1=u$ y que $u_i$
se factoriza en
$B\xrightarrow{u_{i+1}}C_{i+1}\to C_{i+1}/\mathfrak J_{i+1}=C_i$;
de donde se sigue la conclusión.
\end{proof}

\begin{proposition}[19.4.4]
\label{0.19.4.4-es}
Sean $A$ un anillo topológico y $B$ una $A$-álgebra topológica
(conmutativa). Para que $B$ sea una $A$-álgebra formalmente lisa, es necesario y
suficiente que para

\oldpage[0\textsubscript{IV-1}]{87}

todo $B$-módulo topológico discreto $L$, anulado por un ideal abierto de $B$, se
tenga (cf. (18.5.1))
\[
  \operatorname{Exalcotop}_A(B,L)=0.
\]
\end{proposition}

\begin{proof}
Sea $(\mathfrak R_\lambda)$ un sistema fundamental decreciente de ideales
abiertos de $B$ y pongamos $B_\lambda=B/\mathfrak R_\lambda$.
Consideremos una $A$-álgebra topológica discreta $C$ y un ideal
$\mathfrak J$ de $C$ tal que $\mathfrak J^2=0$, de modo que $C$ es una
$A$-extensión de $C/\mathfrak J$ por $\mathfrak J$; supongamos dado un
$A$-homomorfismo $u:B_\lambda\to C/\mathfrak J$ y formemos la
$A$-extensión $E_\lambda$ de $B_\lambda$ por $\mathfrak J$, imagen
inversa de $C$ por el homomorfismo $u$ (18.2.5); es una $A$-álgebra
topológica para la topología discreta. Si
$\operatorname{Exalcotop}_A(B,\mathfrak J)=0$, existe por definición (18.5.1) un
$\mu$ tal que $\mathfrak R_\mu\subset\mathfrak R_\lambda$ y que
la extensión imagen inversa $E_\mu$ sea $A$-trivial; pero esto significa
(18.1.6) que existe un $A$-homomorfismo continuo
$v:B_\mu\to E_\lambda$ tal que el homomorfismo canónico
$B_\mu\to B_\lambda$ se factoriza en
$B_\mu\xrightarrow{v}E_\lambda\to B_\lambda$; se concluye
inmediatamente que $B_\mu\to B_\lambda\xrightarrow{u}C/\mathfrak J$ se
factoriza en
$B_\mu\xrightarrow{v}E_\lambda\to C\to C/\mathfrak J$, y esto
prueba, en virtud de (19.3.1) y (19.4.3), que $B$ es una $A$-álgebra
formalmente lisa. El recíproco es inmediato, aplicando (19.3.1) al caso
en que $C$ es una $A$-álgebra topológica que es una $A$-extensión de
$B_\lambda$ por $L$, y al homomorfismo identidad
$B_\lambda\to B_\lambda=C/L$.
\end{proof}

Cuando $A$ y $B$ son anillos discretos, el criterio de lisura formal
(19.4.4) se reduce a
\begin{equation}
\tag{19.4.4.1}
\label{0.19.4.4.1-es}
\operatorname{Exalcom}_A(B,L)=0\qquad\text{para todo $B$-módulo $L$};
\end{equation}
en otras palabras, \emph{toda $A$-extensión conmutativa de $B$ por un
$B$-módulo debe ser $A$-trivial}.

\begin{corollary}[19.4.5]
\label{0.19.4.5-es}
Sean $A$ un anillo topológico y $B$ una $A$-álgebra topológica
(conmutativa).
\begin{enumerate}
  \item[(i)] Supongamos que $B$ es una $A$-álgebra formalmente lisa.
  Entonces, para todo ideal abierto $\mathfrak R$ de $B$, todo
  $(B/\mathfrak R)$-módulo $L$ y todo $2$-cociclo de Hochschild
  $A$-bilineal y simétrico $f$ de
  $(B/\mathfrak R)\times(B/\mathfrak R)$ en $L$ (18.4.3), existe un
  ideal abierto $\mathfrak R'\subset\mathfrak R$ tal que, si
  $\varphi:B/\mathfrak R'\to B/\mathfrak R$ es el homomorfismo
  canónico, $f\circ(\varphi\times\varphi)$ sea un $2$-coborde de
  Hochschild $A$-bilineal de
  $(B/\mathfrak R')\times(B/\mathfrak R')$ en $L$.
  \item[(ii)] Si $B$ es un $A$-módulo formalmente proyectivo (19.2.1) y si
  se verifica la condición de (i), $B$ es una $A$-álgebra formalmente
  lisa.
\end{enumerate}
\end{corollary}

\begin{proof}
\begin{enumerate}
  \item[(i)] El $2$-cociclo $f$ define una $A$-extensión de Hochschild $C$
  de $B/\mathfrak R$ por $L$ (18.4.3). Aplicando (19.3.1) a $C$, al ideal
  de cuadrado nulo $L$ de $C$ y al homomorfismo identidad
  $B/\mathfrak R\to B/\mathfrak R=C/L$, se deduce la condición (i) en
  virtud de (18.4.3).
  \item[(ii)] Apliquemos el criterio (19.4.3), considerando un ideal abierto
  $\mathfrak R$ de $B$, un ideal abierto $\mathfrak J$ de $A$ tal que
  $\mathfrak JB\subset\mathfrak R$, y una $(A/\mathfrak J)$-extensión $E$
  de $B/\mathfrak R$ por $L$. Como $B$ es un $A$-módulo formalmente
  proyectivo, la aplicación $A$-lineal continua canónica
  $\rho:B\to B/\mathfrak R$ se factoriza en
  $B\xrightarrow{w}E\to B/\mathfrak R$, donde $w$ es una aplicación
  $A$-lineal continua (19.2.1), que se factoriza a su vez en
  $B\to B/\mathfrak R'\xrightarrow{w'}E$, donde $\mathfrak R'$ es un ideal
  abierto de $B$ contenido en $\mathfrak R$; reemplazando $\mathfrak J$ por
  un ideal más pequeño, se puede suponer que
  $\mathfrak JB\subset\mathfrak R'$. Entonces la imagen inversa por
  el homomorfismo canónico $B/\mathfrak R'\to B/\mathfrak R$ de
  la extensión $E$ de $B/\mathfrak R$ por $L$ es una
  $(A/\mathfrak J)$-extensión \emph{de Hochschild} $E'$ de
  $B/\mathfrak R'$ por $L$;

  \oldpage[0\textsubscript{IV-1}]{88}

  aplicando (i) a un cociclo que define esta extensión (18.4.3), se
  concluye que hay un ideal abierto
  $\mathfrak R''\subset\mathfrak R'$ de $B$ tal que la imagen inversa
  $E''$ de $E$ por $B/\mathfrak R''\to B/\mathfrak R$ sea
  $A$-trivial, C.Q.D.
\end{enumerate}
\end{proof}

\begin{corollary}[19.4.6]
\label{0.19.4.6-es}
Sean $A$ un anillo topológico y $B$ una $A$-álgebra topológica que sea un
$A$-módulo formalmente proyectivo. Sea $A'$ una $A$-álgebra \emph{provista de
la topología inducida por la de $A$}. Supongamos además que $A'$ es un
$A$-módulo fielmente plano, y que se cumple una de las condiciones
siguientes:
\begin{enumerate}
  \item[$1^\circ$] Existe un sistema fundamental
  $(\mathfrak J_\lambda)$ de ideales abiertos de $A$ y un sistema fundamental
  $(\mathfrak M_\lambda)$ de ideales abiertos de $B$, con el mismo conjunto
  de índices y tales que, para todo $\lambda$, se tenga
  $\mathfrak J_\lambda B\subset\mathfrak M_\lambda$ y que
  $B/\mathfrak M_\lambda$ sea un
  $(A/\mathfrak J_\lambda)$-módulo proyectivo de tipo finito.
  \item[$2^\circ$] $A'$ es un $A$-módulo proyectivo de tipo finito.
\end{enumerate}
Entonces, para que $B'=B\otimes_A A'$ (provisto de la topología producto tensorial)
sea una $A'$-álgebra formalmente lisa, es necesario y suficiente que $B$ sea una
$A$-álgebra formalmente lisa.
\end{corollary}

\begin{proof}
La suficiencia de la condición está contenida en (19.3.5), (iii), sin ninguna
hipótesis adicional sobre $B$ ni $A'$. Para demostrar el recíproco,
aplicaremos el criterio (19.4.5); bajo la hipótesis $2^\circ$,
seguiremos denotando por $(\mathfrak M_\lambda)$ un sistema fundamental
de ideales abiertos de $B$, y, para todo $\lambda$, por
$\mathfrak J_\lambda$ un ideal abierto de $A$ tal que
$\mathfrak J_\lambda B\subset\mathfrak M_\lambda$; en ambos
casos, pondremos
$A_\lambda=A/\mathfrak J_\lambda$,
$B_\lambda=B/\mathfrak M_\lambda$,
$A'_\lambda=A_\lambda\otimes_A A'$,
$B'_\lambda=B_\lambda\otimes_A A'$. Sea $f$ un $2$-cociclo de
Hochschild $A_\lambda$-bilineal simétrico de
$B_\lambda\times B_\lambda$ en un $B_\lambda$-módulo $L$; por
extensión de escalares, se deduce un $2$-cociclo de Hochschild
$f'=f\otimes1$, $A'_\lambda$-bilineal simétrico, de
$B'_\lambda\times B'_\lambda$ en $L'=L\otimes_A A'$. Como por
hipótesis $B'$ es una $A'$-álgebra formalmente lisa, existe, por (19.4.5),
(i), un índice $\mu$ tal que
$\mathfrak M_\mu\subset\mathfrak M_\lambda$, y tal que, si
$\varphi':B'_\mu\to B'_\lambda$ es la aplicación canónica,
$f'\circ(\varphi'\times\varphi')$ sea un $2$-coborde de Hochschild de
$B'_\mu\times B'_\mu$ en $L'$; en otras palabras, su imagen $c'$ en el
grupo de Hochschild
$\mathrm H^2_{A'_\mu}(B'_\mu,L')^s$ es nula; está claro que si
$\varphi:B_\mu\to B_\lambda$ es el homomorfismo canónico, y $c$ la
clase de $f\circ(\varphi\times\varphi)$ en el grupo de Hochschild
$\mathrm H^2_{A_\mu}(B_\mu,L)^s$, $c'$ es la imagen canónica de $c$.
Ahora bien, si $P_\bullet$ es el complejo relativo a los anillos $A_\mu$ y
$B_\mu$ definido en (18.4.5), utilizado para calcular
$\mathrm H^2_{A_\mu}(B_\mu,L)^s$, el complejo análogo relativo a los
anillos $A'_\mu$ y $B'_\mu$ es evidentemente
$P_\bullet\otimes_A A'$; bajo la hipótesis $1^\circ$, la construcción de
$P_\bullet$ muestra que es un $A_\mu$-módulo proyectivo de tipo finito.
Se concluye, pues, de Bourbaki, \emph{Alg.}, cap.~II, 3\textsuperscript{a} ed., §~5,
n\textsuperscript{o}~3, prop.~7 que, bajo ambas hipótesis, se tiene
\[
  \operatorname{Hom}_{A'_\mu}(P_\bullet\otimes_A A',L\otimes_A A')
  =\operatorname{Hom}_{A_\mu}(P_\bullet,L)\otimes_A A'
\]
salvo isomorfismo canónico; como $A'$ es un $A$-módulo plano, se tiene por
tanto (18.4.5)
\[
  \mathrm H^2_{A'_\mu}(B'_\mu,L')^s
  =\mathrm H^2_{A_\mu}(B_\mu,L)^s\otimes_A A'
\]
y, por tanto, se puede escribir $c'=c\otimes1$; pero, como $A'$ es un $A$-módulo
fielmente plano, la hipótesis $c'=0$ implica $c=0$, lo que termina la
demostración.
\end{proof}

\begin{proposition}[19.4.7]
\label{0.19.4.7-es}
Sean $A$ un anillo topológico preadmisible y $\mathfrak J$ un ideal de
definición de $A$ ($\mathbf 0_{\rm I}$, 7.1.2), y sea $B$ una $A$-álgebra
topológica que sea un $A$-módulo formalmente proyectivo. Consideremos los
anillos topológicos cocientes $A_0=A/\mathfrak J$,
$B_0=B/\mathfrak JB$; entonces, para que $B$ sea una $A$-álgebra formalmente
lisa, es necesario y suficiente que $B_0$ sea una $A_0$-álgebra formalmente
lisa.
\end{proposition}

\oldpage[0\textsubscript{IV-1}]{89}

\begin{proof}
La necesidad de la condición resulta de (19.4.2). Para ver que es
suficiente, observemos primero que, considerando un sistema fundamental de
vecindades abiertas de $0$ en $A$ formadas por ideales
$\mathfrak J_\alpha$ contenidos en $\mathfrak J$, se puede, en virtud de
(19.4.2), reducir al caso en que $A$ es discreto, puesto que
$B/\mathfrak J_\alpha B$ es un
$(A/\mathfrak J_\alpha)$-módulo formalmente proyectivo (19.2.5); por
definición de un anillo preadmisible, $\mathfrak J$ es entonces
\emph{nilpotente}. Basta además demostrar la proposición cuando
$\mathfrak J^2=0$, pues si $\mathfrak J^m=0$, se aplicará sucesivamente
a los anillos $A_k=A/\mathfrak J^k$ y $B_k=B/\mathfrak J^kB$ y al ideal
$\mathfrak J^{k-1}/\mathfrak J^k$ de $A_k$ para $k=2,3,\ldots,m$,
observando por (19.2.5) que $B_k=B\otimes_A A_k$ es un $A_k$-módulo formalmente
proyectivo. Apliquemos el criterio (19.4.5), (ii), considerando un ideal abierto
$\mathfrak R$ de $B$ y un $2$-cociclo de Hochschild $A$-bilineal
simétrico $f$ de $(B/\mathfrak R)\times(B/\mathfrak R)$ en un
$(B/\mathfrak R)$-módulo $L$. Consideremos primero el caso particular en que
$\mathfrak JL=0$, de modo que $L$ también se puede considerar como un
$(B_0/\mathfrak RB_0)$-módulo, y $f$ se factoriza en
\[
  (B/\mathfrak R)\times(B/\mathfrak R)\to
  (B_0/\mathfrak RB_0)\times(B_0/\mathfrak RB_0)
  \xrightarrow{f_0}L,
\]
donde $f_0$ es un $2$-cociclo de Hochschild bilineal simétrico. En virtud de
la hipótesis, existe, pues, un ideal abierto
$\mathfrak R'\subset\mathfrak R$ en $B$ y una aplicación
$A_0$-lineal $g_0:B_0/\mathfrak R'B_0\to L$ tal que
\[
  f_0(\varphi_0(x),\varphi_0(y))
  =xg_0(y)-g_0(xy)+g_0(x)y
\]
para $x,y$ en $B_0/\mathfrak R'B_0$, siendo $\varphi_0$
el homomorfismo canónico
$B_0/\mathfrak R'B_0\to B_0/\mathfrak RB_0$. Se concluye inmediatamente que
la aplicación $A$-lineal compuesta
$g:B/\mathfrak R'\to B_0/\mathfrak R'B_0\xrightarrow{g_0}L$ es tal
que $dg=f\circ(\varphi\times\varphi)$, donde
$\varphi:B/\mathfrak R'\to B/\mathfrak R$ es el homomorfismo
canónico.

Pasemos ahora al caso general, y consideremos primero el
$(B/\mathfrak R)$-módulo $L/\mathfrak JL=L'$, para el cual se tiene
$\mathfrak JL'=0$; si $f'$ es la aplicación $A$-bilineal compuesta
$(B/\mathfrak R)\times(B/\mathfrak R)\xrightarrow{f}L\to L'$, $f'$ sigue
siendo un $2$-cociclo de Hochschild simétrico, y en virtud de lo anterior,
existen un ideal abierto $\mathfrak R'\subset\mathfrak R$ en $B$ y
una aplicación $A$-lineal $g':B/\mathfrak R'\to L'$ que verifica
$dg'=f'\circ(\varphi'\times\varphi')$ para la aplicación canónica
$\varphi':B/\mathfrak R'\to B/\mathfrak R$. Como $B$ es un
$A$-módulo formalmente proyectivo, existen un ideal abierto
$\mathfrak R'_1\subset\mathfrak R'$ y una aplicación $A$-lineal
$g_1:B/\mathfrak R'_1\to L$ tales que el homomorfismo
$B/\mathfrak R'_1\to B/\mathfrak R'\xrightarrow{g'}L'$ se factoriza como
$B/\mathfrak R'_1\xrightarrow{g_1}L\to L'$. Consideremos entonces el
$2$-cociclo de Hochschild
\[
  f_1(x,y)=f(\varphi_1(x),\varphi_1(y))-xg_1(y)+g_1(xy)-g_1(x)y,
\]
aplicación $A$-bilineal simétrica de
$(B/\mathfrak R'_1)\times(B/\mathfrak R'_1)$ en $L$. La elección de
$g_1$ implica que $f_1$ toma sus valores en $\mathfrak JL$. Como
$\mathfrak J(\mathfrak JL)=0$, se puede aplicar de nuevo el primer caso,
y existe, por tanto, un ideal abierto
$\mathfrak R''\subset\mathfrak R'_1$ y una aplicación $A$-lineal
$g_2:B/\mathfrak R''\to\mathfrak JL$ tal que
\[
  f_1(\varphi_2(x),\varphi_2(y))=xg_2(y)-g_2(xy)+g_2(x)y
\]
en $(B/\mathfrak R'')\times(B/\mathfrak R'')$,
siendo $\varphi_2:B/\mathfrak R''\to B/\mathfrak R'_1$ la aplicación
canónica; la aplicación $A$-lineal
$g=g_2+g_1\circ\varphi_2:B/\mathfrak R''\to L$ verifica, por tanto,
$dg=f\circ(\varphi\times\varphi)$ para la aplicación canónica
$\varphi:B/\mathfrak R''\to B/\mathfrak R$. c.q.d.
\end{proof}

\subsection{Lisura formal y anillos graduados asociados}
\label{subsection:0.19.5-es}

\begin{env}[19.5.1]
\label{0.19.5.1-es}
\oldpage[0\textsubscript{IV-1}]{90}
Sean $C$ un anillo topológico (conmutativo), $V$ un $C$-módulo topológico,
y consideremos el álgebra simétrica
$\mathbf S_C^*(V)=\bigoplus_{n\geq0}\mathbf S_C^n(V)$, que dotaremos
canónicamente de una \emph{topología lineal}, compatible con su estructura de
$C$-álgebra. Para ello, consideremos un submódulo abierto $U$ de $V$, y
el \emph{ideal graduado} $U.\mathbf S_C^*(V)$ que genera en
$\mathbf S_C^*(V)$; tomaremos como sistema fundamental de entornos de
$0$ en $\mathbf S_C^*(V)$ el conjunto de las sumas
$\mathfrak a.\mathbf S_C^*(V)+U.\mathbf S_C^*(V)$, donde $\mathfrak a$ (resp.
$U$) recorre un sistema fundamental de ideales abiertos (resp. de submódulos
abiertos) de $C$ (resp. $V$). Se observará que, si la topología de $V$ es menos
fina que la topología inducida por la de $C$, podemos limitarnos a los pares
$(\mathfrak a,U)$ tales que $\mathfrak aV\subset U$, por tanto
$\mathfrak a.\mathbf S_C^*(V)+U.\mathbf S_C^*(V)
=\mathfrak a.1_C+U.\mathbf S_C^*(V)$; en este caso, la topología inducida sobre
cada uno de los $\mathbf S_C^n(V)$ para $n\geq1$ tiene como sistema fundamental de
entornos de $0$ los $U.\mathbf S_C^{n-1}(V)$, donde $U$ recorre los
submódulos abiertos de $V$; en particular, sobre $V=\mathbf S_C^1(V)$, esta
coincide con la topología dada (en general, es menos fina que esta
última). Además, en todos los casos, el \emph{álgebra topológica}
$\mathbf S_C^*(V)$ así definida es, para las categorías de los $C$-módulos
topológicos y de las $C$-álgebras topológicas, la solución del mismo problema
universal que para las categorías de los $C$-módulos y de las $C$-álgebras. En efecto,
sea $E$ una $C$-álgebra topológica, $u:V\to E$ un homomorfismo de
$C$-módulos, $\mathbf S(u)=\widetilde u$ su extensión canónica a
$\mathbf S_C^*(V)$. Supongamos que $u$ es continuo; entonces, si $\mathfrak b$ es un
ideal abierto de $E$, su imagen inversa $U=u^{-1}(\mathfrak b)$ es un
sub-$C$-módulo abierto de $V$, y la imagen por $\widetilde u$ de
$U.\mathbf S_C^*(V)$ está, por tanto, contenida en $\mathfrak b$; como, por otra parte,
existe un ideal abierto $\mathfrak a$ de $C$ tal que
$\mathfrak a.E\subset\mathfrak b$, y por tanto
$\widetilde u(\mathfrak a.\mathbf S_C^*(V))\subset\mathfrak b$, esto demuestra
que $\widetilde u$ es continuo. Recíprocamente, si $\widetilde u$ es continuo,
y si $\mathfrak b$ es un ideal abierto de $E$, existe un submódulo abierto
$U$ de $V$ tal que $\widetilde u(U.\mathbf S_C^*(V))\subset\mathfrak b$, y en
particular $\widetilde u(U.1_C)\subset\mathfrak b$, es decir,
$u(U)\subset\mathfrak b$, por lo que $u$ es continuo. Recordemos además que se tiene
un isomorfismo canónico de $C$-módulos topológicos (discretos)
\begin{equation}
\label{0.19.5.1.1-es}
  \mathbf S_C^n(V)/U.\mathbf S_C^{n-1}(V)
  \xrightarrow{\sim}\mathbf S_C^n(V/U).
  \tag{19.5.1.1}
\end{equation}
\end{env}

\begin{env}[19.5.2]
\label{0.19.5.2-es}
Sean $A$ un anillo topológico, $B$ una $A$-álgebra topológica
(conmutativa), $\mathfrak J$ un ideal de $B$ (no necesariamente abierto ni
cerrado); en todo lo que sigue, se dota a $\mathfrak J$ y a los
$\mathfrak J^n$ ($n\geq2$) de la topología \emph{inducida} por la de $B$,
a los cocientes $C=B/\mathfrak J$,
$\mathfrak J^n/\mathfrak J^{n+1}=\gr_{\mathfrak J}^n(B)$ de la topología
\emph{cociente}, de modo que los $\mathfrak J^n/\mathfrak J^{n+1}$ son
$C$-módulos topológicos; la inyección canónica
$\mathfrak J/\mathfrak J^2\to\gr_{\mathfrak J}^*(B)$ se prolonga a un
homomorfismo (inicialmente no topológico) de $C$-álgebras
\[
  \varphi:\mathbf S_C^*(\mathfrak J/\mathfrak J^2)
  \longrightarrow\gr_{\mathfrak J}^*(B)
\]
que para cada $n$ da un homomorfismo \emph{sobreyectivo} de $C$-módulos
\begin{equation}
\label{0.19.5.2.1-es}
  \varphi_n:\mathbf S_C^n(\mathfrak J/\mathfrak J^2)
  \longrightarrow\mathfrak J^n/\mathfrak J^{n+1}
  =\gr_{\mathfrak J}^n(B).
  \tag{19.5.2.1}
\end{equation}

\oldpage[0\textsubscript{IV-1}]{91}

Cuando $\mathbf S_C^n(\mathfrak J/\mathfrak J^2)$ está dotado de la topología
definida en (19.5.1), los homomorfismos $\varphi_n$ son continuos, en virtud
de la propiedad universal (19.5.1) de
$\mathbf S_C^*(\mathfrak J/\mathfrak J^2)$ aplicada a la $C$-álgebra
topológica $E=\gr_{\mathfrak J/\mathfrak J^n}(B/\mathfrak J^n)$ dotada de la
topología producto de las topologías de los $\mathfrak J^k/\mathfrak J^{k+1}$, y a
la inyección canónica $\mathfrak J/\mathfrak J^2\to E$. Nótese que aquí la
topología de $\mathfrak J/\mathfrak J^2$ es menos fina que la topología
inducida por la de $C$ (siendo también esta última topología sobre
$\mathfrak J/\mathfrak J^2$ la topología inducida por la de $B$).
\end{env}

\begin{theorem}[19.5.3]
\label{0.19.5.3-es}
Sean $A$ un anillo topológico, $B$ una $A$-álgebra topológica,
$\mathfrak J$ un ideal de $B$, $C=B/\mathfrak J$ la $A$-álgebra topológica
cociente. Se supone que la $A$-álgebra discreta $C$ es formalmente lisa.
Entonces:
\begin{enumerate}
  \item[\rm(i)] Si $B$ es una $A$-álgebra formalmente lisa,
  $\mathfrak J/\mathfrak J^2$ es un $C$-módulo topológico formalmente
  proyectivo (19.2.1).
  \item[\rm(ii)] Si $B$ es una $A$-álgebra formalmente lisa y un anillo
  preadmisible ($\mathbf 0_{\rm I}$, 7.1.2), los homomorfismos $\varphi_n$
  (19.5.2) son bimorfismos formales (19.1.2).
  \item[\rm(iii)] Recíprocamente, supongamos que $B$ sea preadmisible, que la
  sucesión $(\mathfrak J^n)$ tienda a $0$ en $B$, que
  $\mathfrak J/\mathfrak J^2$ sea un $C$-módulo formalmente proyectivo y que
  los $\varphi_n$ sean bimorfismos formales. Entonces $B$ es una
  $A$-álgebra formalmente lisa.
\end{enumerate}
\end{theorem}

La demostración de este teorema es larga y está recargada de detalles técnicos;
por ello, comenzaremos por demostrar un corolario más sencillo, en el que
aparecen con mayor claridad las ideas directrices; este corolario es, por lo demás,
el caso particular del teorema (19.5.3) que se empleará con mayor frecuencia
en lo sucesivo.

\begin{corollary}[19.5.4]
\label{0.19.5.4-es}
Sean $A$ un anillo topológico, $B$ una $A$-álgebra topológica,
$\mathfrak J$ un ideal de $B$ tal que la topología de $B$ sea la topología
$\mathfrak J$-preádica. Se supone que la $A$-álgebra discreta
$C=B/\mathfrak J$ es formalmente lisa. Entonces las tres condiciones siguientes
son equivalentes:
\begin{enumerate}
  \item[a)] $B$ es una $A$-álgebra formalmente lisa.
  \item[b)] $\mathfrak J/\mathfrak J^2$ es un $C$-módulo proyectivo y
  el homomorfismo canónico
  \[
    \varphi:\mathbf S_C^*(\mathfrak J/\mathfrak J^2)
    \longrightarrow\gr_{\mathfrak J}^*(B)
  \]
  es biyectivo.
  \item[c)] El completado separado $\widehat B$ de $B$ es una
  $\widehat A$-álgebra topológica isomorfa a una $\widehat A$-álgebra de la
  forma $\widehat{B'}$, donde $B'=\mathbf S_C^*(V)$, siendo $V$ un $C$-módulo
  proyectivo y estando $B'$ dotado de la topología ${B'}^+$-preádica, donde
  ${B'}^+$ es el ideal de aumentación.
\end{enumerate}
\end{corollary}

\begin{env}[19.5.4.1]
\label{0.19.5.4.1-es}
Demostremos primero que $a)$ implica que $\mathfrak J/\mathfrak J^2$ es un
$C$-módulo proyectivo. Sean, pues, $P$ y $Q$ dos $C$-módulos,
$u:P\to Q$ un $C$-homomorfismo sobreyectivo, y
$v:\mathfrak J/\mathfrak J^2\to Q$ un $C$-homomorfismo. Considerando el anillo
$B/\mathfrak J^2$ como una $B$-extensión de $C$ por
$\mathfrak J/\mathfrak J^2$, se deduce de ella mediante $v$ una $B$-extensión
$E=(B/\mathfrak J^2)\oplus_{(\mathfrak J/\mathfrak J^2)}Q$ de $C$ por $Q$
(18.2.8). Como $C$ es una $A$-álgebra discreta formalmente lisa,
la extensión $E$ es $A$-\emph{trivial} (19.4.4) y puede, por tanto, identificarse con
$\mathrm D_C(Q)$ (18.2.3). El homomorfismo sobreyectivo $u:P\to Q$ define entonces
canónicamente un $A$-homomorfismo sobreyectivo
$\psi:\mathrm D_C(P)\to\mathrm D_C(Q)$ (18.2.9) cuyo núcleo es un ideal
$\mathfrak L$ de la extensión $E'=\mathrm D_C(P)$, contenido en $P$, y
\emph{a fortiori} de cuadrado nulo. Sea $f:B\to E=E'/\mathfrak L$
el homomorfismo que define la estructura de $B$-álgebra de $E$: como
$\mathfrak L$ es de cuadrado nulo, y $B$ es una $A$-álgebra formalmente
lisa, $f$ se factoriza como $B\xrightarrow{g}E'\to E'/\mathfrak L$, donde $g$ es
un $A$-homomorfismo continuo.

\oldpage[0\textsubscript{IV-1}]{92}

El diagrama
\[
  \xymatrix{
    E' \ar[r]^{\psi} & E \\
    B \ar[u]^{g} \ar[r]_{\theta} & B/\mathfrak J^2 \ar[u]_{v'}
  }
\]
donde $v'$ se deduce de $v$ (18.2.8), es conmutativo. Además, la imagen de
$\mathfrak J$ por $\psi\circ g$ está contenida en $Q$, por tanto la imagen de
$\mathfrak J$ por $g$ está contenida en $P+\mathfrak L=P$, y la imagen de
$\mathfrak J^2$ por $g$ es nula. Restringiendo $g$ y $\theta$ a
$\mathfrak J$, se obtiene un diagrama conmutativo
\[
  \xymatrix{
    P \ar[r]^{u} & Q \\
    \mathfrak J/\mathfrak J^2 \ar[u]^{w} \ar[ur]_{v}
  }
\]
lo que demuestra que el $C$-módulo $\mathfrak J/\mathfrak J^2$ es proyectivo.
\end{env}

\begin{env}[19.5.4.2]
\label{0.19.5.4.2-es}
Demostremos en segundo lugar que $a)$ implica que $\varphi$ es biyectivo, lo que
acabará de demostrar que $a)$ implica $b)$. Pongamos
$E_n=B/\mathfrak J^{n+1}$, y denotemos por $F_n$ el cociente de
$\mathbf S_C^*(\mathfrak J/\mathfrak J^2)$ por la potencia $(n+1)$-ésima de
su ideal de aumentación. Como $\mathfrak J/\mathfrak J^{n+1}$ es nilpotente
en $E_n$ y $C$ es una $A$-álgebra discreta formalmente lisa,
el automorfismo identidad de $C$ se factoriza como
$C\xrightarrow{f}E_n\to C=E_n/(\mathfrak J/\mathfrak J^{n+1})$, donde $f$ es un
$A$-homomorfismo; por otra parte, como $\mathfrak J/\mathfrak J^2$ es un
$C$-módulo proyectivo según (19.5.4.1), el automorfismo identidad de
$\mathfrak J/\mathfrak J^2$ se factoriza como
$\mathfrak J/\mathfrak J^2\xrightarrow{g}\mathfrak J/\mathfrak J^{n+1}
\to\mathfrak J/\mathfrak J^2$, donde $g$ es una aplicación $C$-lineal; a
partir de $f$ y $g$ se obtiene canónicamente un homomorfismo de $C$-álgebras
$\mathbf S_C^*(\mathfrak J/\mathfrak J^2)\to E_n$, que además (por
definición de $g$) se anula sobre la potencia $(n+1)$-ésima del ideal
de aumentación de $\mathbf S_C^*(\mathfrak J/\mathfrak J^2)$, de donde, al
pasar al cociente, se obtiene un $A$-homomorfismo sobreyectivo de álgebras
$v:F_n\to E_n$ tal que $\gr^0(v)$ y $\gr^1(v)$ sean los automorfismos
identidad de $C$ y de $\mathfrak J/\mathfrak J^2$. Por definición del
homomorfismo canónico $\varphi$, se ve que $\gr^j(v)=\varphi_j$ para todo
$j\leq n$. Observemos ahora que el núcleo $\mathfrak R$ de $v$ es un ideal
nilpotente de $F_n$, de modo que $E_n$ se identifica con $F_n/\mathfrak R$. Como
$B$ es una $A$-álgebra formalmente lisa, el $A$-homomorfismo canónico
$p_n:B\to E_n=B/\mathfrak J^{n+1}$, que es continuo, se factoriza en
$B\xrightarrow{w}F_n\xrightarrow{v}E_n$, donde $w$ es un $A$-homomorfismo
continuo; como $\gr^0(v)$ es la identidad, $w(\mathfrak J)$ está contenido en
el ideal de aumentación de $F_n$, de donde $w(\mathfrak J^{n+1})=0$, por lo que
$w$ se factoriza en
$B\xrightarrow{p_n}B/\mathfrak J^n=E_n\xrightarrow{w'}F_n$, y el compuesto
$E_n\xrightarrow{w'}F_n\xrightarrow{v}E_n$ es la identidad. Además, como
$\gr^0(v)$ y $\gr^1(v)$ son los automorfismos identidad de $C$ y de
$\mathfrak J/\mathfrak J^2$, lo mismo ocurre con $\gr^0(w')$ y con
$\gr^1(w')$. Como $\gr^*(E_n)$ está engendrado por $\gr^1(E_n)$, el
homomorfismo compuesto
\[
  \gr^j(F_n)\xrightarrow{\varphi_j}\gr^j(E_n)
  \xrightarrow{\gr^j(w')}\gr^j(F_n)
\]
es la identidad para todo $j\leq n$, puesto que lo es para $j=0$ y $j=1$;
tomando $j=n$, se prueba así que $\varphi_n$ es inyectivo, lo que termina de
mostrar que $a)$ implica $b)$.
\end{env}

\begin{env}[19.5.4.3]
\label{0.19.5.4.3-es}
\oldpage[0\textsubscript{IV-1}]{93}
Probemos ahora que $b)$ implica $a)$. El mismo razonamiento que al comienzo de
(19.5.4.2) prueba la existencia de un $A$-homomorfismo sobreyectivo de álgebras
$v:F_n\to E_n$ tal que $\gr^j(v)=\varphi_j$ para todo $j$; como $\varphi_j$
es biyectivo para todo $j$ y las filtraciones de $F_n$ y $E_n$ son
finitas, se concluye que $v$ es \emph{biyectivo} (Bourbaki,
\emph{Alg. conm.}, cap.~III, §~2, n\textsuperscript{o}~8, cor.~3 del teo.~1).
Sean ahora $G$ una $A$-álgebra topológica discreta, $\mathfrak R$ un
ideal de cuadrado nulo en $G$ y $f:B\to G/\mathfrak R$ un $A$-homomorfismo
continuo de álgebras. Como $G$ es discreta, existe un entero $m$ tal que
$f$ se anula sobre $\mathfrak J^m$, por tanto $f$ se factoriza en
$B\to E_n\xrightarrow{f_n}G/\mathfrak R$, donde se toma $n=2m$. Se obtiene así por
composición un $A$-homomorfismo continuo de álgebras
$r:C\to F_n\xrightarrow{v}E_n\xrightarrow{f_n}G/\mathfrak R$, y como $C$
es una $A$-álgebra formalmente lisa discreta, $r$ se factoriza en
$C\xrightarrow{r'}G\to G/\mathfrak R$, de modo que $G$ queda dotada por $r'$ de una
estructura de $C$-álgebra. Por otra parte, la restricción a
$\mathfrak J/\mathfrak J^2$ del homomorfismo
$f'_n:F_n\xrightarrow{v}E_n\xrightarrow{f_n}G/\mathfrak R$ es una aplicación
$C$-lineal $g:\mathfrak J/\mathfrak J^2\to G/\mathfrak R$. Puesto que
$\mathfrak J/\mathfrak J^2$ es un $C$-módulo proyectivo, $g$ se factoriza en
$\mathfrak J/\mathfrak J^2\xrightarrow{h}G\to G/\mathfrak R$, donde $h$ es una
aplicación $C$-lineal; por prolongación, se deduce un homomorfismo de
$C$-álgebras $w:\mathbf S_C^*(\mathfrak J/\mathfrak J^2)\to G$, y por
construcción todo elemento de grado $m$ tiene por imagen mediante $w$ un elemento de
$\mathfrak R$, por tanto todo elemento de grado $n=2m$ tiene imagen nula, puesto que
$\mathfrak R^2=0$; en otras palabras, $w$ se factoriza en
$\mathbf S_C^*(\mathfrak J/\mathfrak J^2)\to F_n\xrightarrow{w'}G$. Por
construcción, el compuesto $F_n\xrightarrow{w'}G\xrightarrow{\theta}G/\mathfrak R$
coincide con $f'_n$ en $C$ y en $\mathfrak J/\mathfrak J^2$, luego es
igual a $f'_n$. Finalmente, como el compuesto
$B\to E_n\xrightarrow{v^{-1}}F_n\xrightarrow{w'}G\to G/\mathfrak R$ es
igual a $f$, se ve que $B$ es una $A$-álgebra formalmente lisa.
\end{env}

\begin{env}[19.5.4.4]
\label{0.19.5.4.4-es}
Queda por probar la equivalencia de $a)$ y $c)$. En primer lugar, $c)$ implica
$a)$: en efecto, $B'$ es una $C$-álgebra formalmente lisa para la topología
discreta (19.3.2), luego también para la topología ${B'}^+$-preádica (19.3.8);
como $C$ es una $A$-álgebra formalmente lisa, $B'$ es también una
$A$-álgebra formalmente lisa (19.3.5, (ii)) y finalmente $\widehat{B'}$ es
una $\widehat A$-álgebra formalmente lisa (19.3.6), luego $B$ es una
$A$-álgebra formalmente lisa (19.3.6). Queda por ver que $a)$ implica $c)$.
Observemos primero que, puesto que $C$ es una $A$-álgebra formalmente lisa,
el homomorfismo identidad $C\to B/\mathfrak J$ se factoriza en
$C\to\widehat B\to B/\mathfrak J$, donde $C\to\widehat B$ es un
$A$-homomorfismo (19.3.10); $\widehat B$, y por tanto todos los
$B/\mathfrak J^{n+1}$, quedan así dotados de estructuras de $C$-álgebras. Por otra
parte, como $\mathfrak J/\mathfrak J^2$ es un $C$-módulo proyectivo por $b)$,
la inyección canónica $\mathfrak J/\mathfrak J^2\to B/\mathfrak J^2$ permite
formar un sistema proyectivo de $C$-homomorfismos
$u_n:\mathfrak J/\mathfrak J^2\to B/\mathfrak J^{n+1}$ para $n\geq1$, y por
tanto también (por la propiedad universal del álgebra simétrica) un sistema
proyectivo de homomorfismos de $C$-álgebras
$v_n:\mathbf S_C^*(\mathfrak J/\mathfrak J^2)=B'\to B/\mathfrak J^{n+1}$;
además, está claro que $v_n$ se anula sobre $({B'}^+)^{n+1}$ y, como
$\gr^0(v_n)=\varphi_0$ y $\gr^1(v_n)=\varphi_1$, se tiene
$\gr^j(v_n)=\varphi_j$ para $0\leq j\leq n$. Como las $\varphi_j$ son
isomorfismos por $b)$, y las filtraciones de
$B'/({B'}^+)^{n+1}$ y de $B/\mathfrak J^{n+1}$ son finitas, se concluye que
$v_n$ es \emph{biyectivo} para

\oldpage[0\textsubscript{IV-1}]{94}

todo $n$ (Bourbaki, \emph{Alg. conm.}, cap.~III, §~2,
n\textsuperscript{o}~8, cor.~3 del teo.~1); de donde se sigue $c)$ pasando al límite
proyectivo.
\end{env}

\begin{remark}[19.5.5]
\label{0.19.5.5-es}
Bajo las hipótesis generales de (19.5.4), supongamos además que
$\mathfrak J$ es un ideal \emph{maximal} de $B$, de modo que
$C=B/\mathfrak J$ es un cuerpo, y que $\mathfrak J/\mathfrak J^2$ es un
$C$-espacio vectorial de dimensión \emph{finita}. Entonces las condiciones $a)$,
$b)$ y $c)$ de (19.5.4) siguen siendo equivalentes a la siguiente:
\begin{enumerate}
  \item[d)] \emph{Dada un álgebra de polinomios}
  $E=C[T_1,\ldots,T_n]$ \emph{sobre el cuerpo} $C$, \emph{y denotando por}
  $\mathfrak n$ \emph{el ideal maximal de} $E$ \emph{engendrado por los} $T_i$
  $(1\leq i\leq n)$, \emph{entonces, para todo ideal} $\mathfrak b$ \emph{de}
  $E$ \emph{que contiene una potencia} $\mathfrak n^k$, \emph{todo
  homomorfismo} $v:B\to E/\mathfrak b$ \emph{de} $A$-\emph{álgebras}
  $C$-\emph{aumentadas se factoriza en}
  $B\xrightarrow{w}E/\mathfrak n^k\to E/\mathfrak b$, \emph{donde} $w$ \emph{es
  un} $A$-\emph{homomorfismo}.
\end{enumerate}
En efecto, como $\mathfrak J/\mathfrak J^2$ es aquí un $C$-módulo libre, basta
verificar que la condición $d)$ implica que el homomorfismo canónico
$\varphi:\mathbf S_C^*(\mathfrak J/\mathfrak J^2)\to\gr_{\mathfrak J}^*(B)$
es biyectivo. Ahora bien, aquí $\mathbf S_C^*(\mathfrak J/\mathfrak J^2)$ se identifica con
el álgebra de polinomios $E=C[T_1,\ldots,T_n]$, donde
$n=\operatorname{rg}_C(\mathfrak J/\mathfrak J^2)$, y el ideal de aumentación
de $E$ es el ideal $\mathfrak n$ engendrado por los $T_i$. Para todo entero
$k\geq0$, la hipótesis de que $C$ es una $A$-álgebra discreta formalmente lisa
implica, como en (19.5.4.2), la existencia de un $A$-homomorfismo sobreyectivo
$f:E/\mathfrak n^{k+1}\to B/\mathfrak J^{k+1}$, tal que
$\gr^j(f)=\varphi_j$ para todo $j\leq k$. Si
$\mathfrak b/\mathfrak n^{k+1}=\operatorname{Ker}(f)$,
$B/\mathfrak J^{k+1}$ se identifica, pues, con $E/\mathfrak b$, y la hipótesis $d)$
permite factorizar el homomorfismo canónico $B\to B/\mathfrak J^{k+1}$ en
$B\xrightarrow{w}E/\mathfrak n^{k+1}\xrightarrow{f}B/\mathfrak J^{k+1}$;
como $\gr^0(f)$ es la identidad, se tiene
$w(\mathfrak J)\subset\mathfrak n/\mathfrak n^{k+1}$, por tanto
$w(\mathfrak J^{k+1})=0$; se concluye como en (19.5.4.2) que $\varphi_k$
es inyectivo para todo $k$, lo que termina de probar nuestra afirmación.
\end{remark}

\begin{env}[19.5.6]
\label{0.19.5.6-es}
\emph{Demostración del teorema} (19.5.3). --- Sea
$(\mathfrak b_\alpha)$ un sistema fundamental decreciente de ideales abiertos de
$B$. Pondremos
\[
  B_\alpha=B/\mathfrak b_\alpha,\qquad
  C_\alpha=B/(\mathfrak b_\alpha+\mathfrak J),\qquad
  \mathfrak J_\alpha=(\mathfrak b_\alpha+\mathfrak J)/\mathfrak b_\alpha,
\]
de modo que
\[
  C_\alpha=B_\alpha/\mathfrak J_\alpha\qquad\text{y}\qquad
  \mathfrak J_\alpha^n/\mathfrak J_\alpha^{n+1}
  =(\mathfrak b_\alpha+\mathfrak J^n)/(\mathfrak b_\alpha+\mathfrak J^{n+1}).
\]
Los $C$-módulos
$((\mathfrak b_\alpha\cap\mathfrak J^n)+\mathfrak J^{n+1})/
\mathfrak J^{n+1}$ de $\mathfrak J^n/\mathfrak J^{n+1}$ forman un sistema
fundamental de submódulos abiertos del $C$-módulo topológico
$\mathfrak J^n/\mathfrak J^{n+1}$; como
$(\mathfrak b_\alpha\cap\mathfrak J^n)+\mathfrak J^{n+1}
=\mathfrak J^n\cap(\mathfrak b_\alpha+\mathfrak J^{n+1})$, se tiene
\[
  \mathfrak J^n/((\mathfrak b_\alpha\cap\mathfrak J^n)+\mathfrak J^{n+1})
  =\mathfrak J^n/(\mathfrak J^n\cap(\mathfrak b_\alpha+\mathfrak J^{n+1}))
  =(\mathfrak b_\alpha+\mathfrak J^n)/(\mathfrak b_\alpha+\mathfrak J^{n+1})
  =\mathfrak J_\alpha^n/\mathfrak J_\alpha^{n+1}.
\]
\end{env}

\begin{env}[19.5.6.1]
\label{0.19.5.6.1-es}
Sean $P$, $Q$ dos $C_\alpha$-módulos discretos,
$u:P\to Q$ un $C_\alpha$-homomorfismo sobreyectivo. El anillo discreto
$B/(\mathfrak b_\alpha+\mathfrak J^2)$ es una $B$-extensión de $C_\alpha$ por
el ideal de cuadrado nulo $\mathfrak J_\alpha/\mathfrak J_\alpha^2$. Sea
$v:\mathfrak J/\mathfrak J^2\to Q$ un $C$-homomorfismo continuo; reemplazando
si es necesario $\mathfrak b_\alpha$ por un ideal abierto más pequeño de $B$, se puede
suponer que el núcleo de $v$ contiene el sub-$C$-módulo abierto
$((\mathfrak b_\alpha\cap\mathfrak J)+\mathfrak J^2)/\mathfrak J^2$;
pasando al cociente, se deduce un $C_\alpha$-homomorfismo de módulos
discretos $v':\mathfrak J_\alpha/\mathfrak J_\alpha^2\to Q$; sea
$E_\alpha$ la $B_\alpha$-extensión de $C_\alpha$ por $Q$ deducida de
$B/(\mathfrak b_\alpha+\mathfrak J^2)$ mediante $v'$, y sea
$v'_*:B/(\mathfrak b_\alpha+\mathfrak J^2)\to E_\alpha$ el
$B_\alpha$-homomorfismo correspondiente (18.2.8); $E_\alpha$ es una
$A$-álgebra topológica discreta, y el isomorfismo canónico
$C_\alpha\xrightarrow{\sim}E_\alpha/Q$ da por composición un
$A$-homomorfismo continuo $f:C\to C_\alpha\to E_\alpha/Q$. Pero como $Q$ es
de cuadrado nulo en $E_\alpha$ y $C$ es una $A$-álgebra formalmente
lisa, $f$ se factoriza en $C\xrightarrow{g}E_\alpha\to E_\alpha/Q$, donde $g$
es un $A$-homomorfismo continuo. Como $g$ es continuo y $E_\alpha$ discreta,
existe un $\beta\geq\alpha$ tal que $g$ se factoriza en
$C\to C_\beta\xrightarrow{g'}E_\alpha$. Por otra parte, sea $E_\beta$ la
$A$-extensión de $C_\beta$ por $Q$, imagen inversa de $E_\alpha$ por
el homomorfismo canónico $C_\beta\to C_\alpha$; la existencia de $g'$ (tal que
el compuesto $C_\beta\to E_\alpha\to C_\alpha$ sea el homomorfismo canónico)
equivale a que $E_\beta$ sea una extensión $A$-\emph{trivial}, y puede
por tanto identificarse con $\mathrm D_{C_\beta}(Q)$. Dicho esto, el homomorfismo
sobreyectivo $u:P\to Q$ define canónicamente un homomorfismo sobreyectivo
$\mathrm D_{C_\beta}(P)\to\mathrm D_{C_\beta}(Q)$ (18.2.9), cuyo núcleo es
un ideal $\mathfrak L$ de la extensión $E'_\beta=\mathrm D_{C_\beta}(P)$
contenido en $P$, y \emph{a fortiori} de cuadrado nulo. Sea
$f':B\to E_\beta=E'_\beta/\mathfrak L$ el $A$-homomorfismo continuo que define
la estructura de $B$-álgebra topológica de $E_\beta$; como $\mathfrak L$ es
de cuadrado nulo y $B$ es una $A$-álgebra formalmente lisa, $f'$ se
factoriza en $B\xrightarrow{h'}E'_\beta\to E'_\beta/\mathfrak L$, donde $h'$ es
un $A$-homomorfismo continuo. La construcción de $E'_\beta$ da por
composición un $A$-homomorfismo $\psi:E'_\beta\to E_\beta\to E_\alpha$, y
está claro que el diagrama
\[
  \xymatrix{
    E'_\beta \ar[r]^{\psi} & E_\alpha \\
    B \ar[u]^{h'} \ar[r]_{\theta} &
    B/(\mathfrak b_\alpha+\mathfrak J^2) \ar[u]_{v'_*}
  }
\]
es conmutativo. Además, la imagen por $\psi\circ h'$ de $\mathfrak J$ está
contenida en $Q$, por tanto la imagen por $h'$ de $\mathfrak J$ está contenida en
$P$, y la imagen de $\mathfrak J^2$ por $h'$ es nula. Restringiendo $h'$ y
$\theta$ a $\mathfrak J$, se obtiene un diagrama conmutativo
\[
  \xymatrix{
    P \ar[r]^{u} & Q \\
    \mathfrak J/\mathfrak J^2 \ar[u]^{w} \ar[ur]_{v}
  }
\]
donde $w$ es continuo, lo que prueba que
$\mathfrak J/\mathfrak J^2$ es un $C$-módulo formalmente proyectivo.
\end{env}

\oldpage[0\textsubscript{IV-1}]{95}

\begin{env}[19.5.6.2]
\label{0.19.5.6.2-es}
Para todo entero $n\geq0$, pondremos
\[
  E_n=B/\mathfrak J^{n+1},
  \qquad\text{de modo que}\qquad E_0=C~;
\]
los ideales $(\mathfrak b_\alpha+\mathfrak J^{n+1})/\mathfrak J^{n+1}$ forman
un sistema fundamental de ideales abiertos de $E_n$, y pondremos
\[
  E_{\alpha,n}=B/(\mathfrak b_\alpha+\mathfrak J^{n+1})
  =B_\alpha/\mathfrak J_\alpha^{n+1},
\]
anillo cociente discreto. Del mismo modo, en
$\gr_{\mathfrak J}^n(B)=\mathfrak J^n/\mathfrak J^{n+1}$, hemos visto que los
$((\mathfrak b_\alpha\cap\mathfrak J^n)+\mathfrak J^{n+1})/
\mathfrak J^{n+1}$ forman un sistema fundamental de entornos de $0$, y el
cociente de $\mathfrak J^n/\mathfrak J^{n+1}$ por este submódulo se identifica
canónicamente con $\mathfrak J_\alpha^n/\mathfrak J_\alpha^{n+1}$. Consideremos
el álgebra simétrica
$\mathbf S_{C_\alpha}^*(\mathfrak J_\alpha/\mathfrak J_\alpha^2)$; denotaremos
por $F_{\alpha,n}$ el cociente de
$\mathbf S_{C_\alpha}^*(\mathfrak J_\alpha/\mathfrak J_\alpha^2)$ por la
potencia $(n+1)$-ésima de su ideal de aumentación. Para un $n\geq1$ fijo, se
deduce de (19.5.1.1) que los
$\mathbf S_{C_\alpha}^n(\mathfrak J_\alpha/\mathfrak J_\alpha^2)$ son los
cocientes de $\mathbf S_C^n(\mathfrak J/\mathfrak J^2)$ por un sistema
fundamental de submódulos abiertos en este $C$-módulo topológico.

Para abreviar, diremos que, para $\alpha\leq\beta$, los homomorfismos canónicos
$B_\beta\to B_\alpha$, $C_\beta\to C_\alpha$,
$\mathfrak J_\beta/\mathfrak J_\beta^2\to
\mathfrak J_\alpha/\mathfrak J_\alpha^2$, $E_{\beta,n}\to E_{\alpha,n}$,
$F_{\beta,n}\to F_{\alpha,n}$, etc., son los \emph{homomorfismos de
transición}.
\end{env}

\begin{lemma}[19.5.6.3]
\label{0.19.5.6.3-es}
Supongamos que la $A$-álgebra $C$ es formalmente lisa, el anillo $B$
preadmisible y el $C$-módulo $\mathfrak J/\mathfrak J^2$ formalmente
proyectivo. Entonces:
\begin{enumerate}
  \item[\rm(i)] Para todo $\alpha$, existen $\beta\geq\alpha$ y un
  $A$-homomorfismo sobreyectivo de álgebras
  \[
    v_{\alpha\beta}:F_{\beta,n}\longrightarrow E_{\alpha,n}
  \]
  tales que $\gr^0(v_{\alpha\beta}):C_\beta\to C_\alpha$ y
  $\gr^1(v_{\alpha\beta}):\mathfrak J_\beta/\mathfrak J_\beta^2
  \to\mathfrak J_\alpha/\mathfrak J_\alpha^2$ sean los homomorfismos de
  transición.
  \item[\rm(ii)] Si $\beta$ y $v_{\alpha\beta}$ verifican las condiciones de
  \emph{(i)}, entonces, para todo $\gamma\geq\beta$, existen
  $\delta\geq\gamma$ y un $A$-homomorfismo sobreyectivo de álgebras
  $v_{\gamma\delta}:F_{\delta,n}\to E_{\gamma,n}$ que verifica las condiciones de
  \emph{(i)} (para $\gamma$ y $\delta$) y hace conmutativo el diagrama
  \begin{equation}
  \label{0.19.5.6.4-es}
  \xymatrix@C=3em{
    F_{\beta,n} \ar[r]^{v_{\alpha\beta}} & E_{\alpha,n} \\
    F_{\delta,n} \ar[u] \ar[r]_{v_{\gamma\delta}} & E_{\gamma,n} \ar[u]
  }
  \tag{19.5.6.4}
  \end{equation}
  donde las flechas verticales son los homomorfismos de transición.
\end{enumerate}
\end{lemma}

\begin{env}[19.5.6.4]
\label{0.19.5.6.4-proof-es}
\begin{enumerate}
  \item[\rm(i)] En la $A$-álgebra topológica discreta $E_{\alpha,n}$,
  el ideal $\mathfrak J_\alpha/\mathfrak J_\alpha^{n+1}$ es nilpotente, y
  el isomorfismo identidad
  $C_\alpha\xrightarrow{\sim}
  E_{\alpha,n}/(\mathfrak J_\alpha/\mathfrak J_\alpha^{n+1})$ da por
  composición un $A$-homomorfismo continuo
  $C\to C_\alpha\to
  E_{\alpha,n}/(\mathfrak J_\alpha/\mathfrak J_\alpha^{n+1})$;

\oldpage[0\textsubscript{IV-1}]{96}

  como $C$ es una $A$-álgebra formalmente lisa, este homomorfismo se
  factoriza como
  $C\xrightarrow{f_\alpha}E_{\alpha,n}\to
  E_{\alpha,n}/(\mathfrak J_\alpha/\mathfrak J_\alpha^{n+1})$, donde $f_\alpha$
  es un $A$-homomorfismo continuo; por tanto,
  $\mathfrak J_\alpha/\mathfrak J_\alpha^{n+1}$ se convierte, mediante
  $f_\alpha$, en un $C$-módulo topológico discreto anulado por un ideal abierto de
  $C$. La hipótesis de que $\mathfrak J/\mathfrak J^2$ es un $C$-módulo
  formalmente proyectivo implica entonces la existencia de una aplicación
  $C$-lineal continua
  $g_\alpha:\mathfrak J/\mathfrak J^2\to
  \mathfrak J_\alpha/\mathfrak J_\alpha^{n+1}$ que hace conmutativo el diagrama
  \[
    \xymatrix{
      \mathfrak J_\alpha/\mathfrak J_\alpha^{n+1} \ar[r] &
      \mathfrak J_\alpha/\mathfrak J_\alpha^2 \\
      & \mathfrak J/\mathfrak J^2 \ar[ul]^{g_\alpha} \ar[u]
    }
  \]
  Como $f_\alpha$ y $g_\alpha$ son continuos, existe un
  $\beta\geq\alpha$ tal que estos homomorfismos se factoricen respectivamente como
  \[
    C\to C_\beta\xrightarrow{f'_{\alpha\beta}}E_{\alpha,n},
  \]
  \[
    \mathfrak J/\mathfrak J^2\to
    \mathfrak J_\beta/\mathfrak J_\beta^2
    \xrightarrow{g'_{\alpha\beta}}
    \mathfrak J_\alpha/\mathfrak J_\alpha^{n+1}
  \]
  y, a partir de $f'_{\alpha\beta}$ y $g'_{\alpha\beta}$, se obtiene, pues,
  canónicamente un homomorfismo de $C_\beta$-álgebras
  \[
    \mathbf S_{C_\beta}^*(\mathfrak J_\beta/\mathfrak J_\beta^2)
    \longrightarrow E_{\alpha,n}
  \]
  que además (por definición de $g'_{\alpha\beta}$) se anula sobre la
  potencia $(n+1)$-ésima del ideal de aumentación de
  $\mathbf S_{C_\beta}^*(\mathfrak J_\beta/\mathfrak J_\beta^2)$; pasando al
  cociente por esta potencia $(n+1)$-ésima, se obtiene el homomorfismo
  $v_{\alpha\beta}$ buscado, teniendo en cuenta las definiciones de $f_\alpha$ y
  $g_\alpha$; la sobreyectividad de $v_{\alpha\beta}$ resulta, en efecto, de la
  de los dos homomorfismos $\gr^0(v_{\alpha\beta})$ y
  $\gr^1(v_{\alpha\beta})$, puesto que esto implica que
  $\gr(v_{\alpha\beta})$ es sobreyectivo (el álgebra
  $\gr^*(E_{\alpha,n})$ está engendrada por $\gr^0(E_{\alpha,n})$ y
  $\gr^1(E_{\alpha,n})$), y, como las filtraciones consideradas son finitas, se
  puede aplicar Bourbaki, \emph{Alg. comm.}, cap.~III, §~2,
  n\textsuperscript{o}~8, cor.~3 del teorema~1.

  \item[\rm(ii)] La hipótesis de que $B$ es \emph{preadmisible} significa que
  se pueden suponer todos los $\mathfrak b_\alpha$ contenidos en un mismo
  $\mathfrak b_{\alpha_0}$ cuyas potencias \emph{tienden a} $0$. Esto
  implica, en particular, que el núcleo del homomorfismo de transición
  $E_{\gamma,n}\to E_{\alpha,n}$, igual a
  $(\mathfrak b_\alpha+\mathfrak J^{n+1})/
  (\mathfrak b_\gamma+\mathfrak J^{n+1})$, es \emph{nilpotente}; aplicando el
  lema (19.3.10.1), se ve que se puede suponer elegido $f_\gamma$ de modo que
  el diagrama
  \[
    \xymatrix{
      & E_{\alpha,n} \\
      C \ar[ur]^{f_\alpha} \ar[r]_{f_\gamma} & E_{\gamma,n} \ar[u]
    }
  \]
  (donde la flecha vertical es el homomorfismo de transición) sea conmutativo.
  La hipótesis de que $\mathfrak J/\mathfrak J^2$ es un $C$-módulo formalmente
  proyectivo permite, por otra parte, elegir
  $g_\gamma:\mathfrak J/\mathfrak J^2\to
  \mathfrak J_\gamma/\mathfrak J_\gamma^{n+1}$ de modo que el diagrama
  \[
    \xymatrix@C=5em@R=3em{
      \mathfrak J_\alpha/\mathfrak J_\alpha^{n+1} \ar[r] &
      \mathfrak J_\alpha/\mathfrak J_\alpha^2 \\
      \mathfrak J_\gamma/\mathfrak J_\gamma^{n+1} \ar[u] \ar[r] &
      \mathfrak J_\gamma/\mathfrak J_\gamma^2 \ar[u] \\
      & \mathfrak J/\mathfrak J^2
        \ar[ul]_{g_\gamma} \ar[u]_{p_\gamma} \ar[uul]^{g_\alpha}
    }
  \]
  sea conmutativo (basta observar que la imagen por $(g_\alpha,p_\gamma)$
  de $\mathfrak J/\mathfrak J^2$ en el módulo producto
  $(\mathfrak J_\alpha/\mathfrak J_\alpha^{n+1})
  \times(\mathfrak J_\gamma/\mathfrak J_\gamma^2)$ está contenida (en vista de la
  definición de $g_\alpha$ y de la relación $\beta\leq\gamma$) en la imagen
  canónica $Q$ en este módulo producto del módulo
  $P=\mathfrak J_\gamma/\mathfrak J_\gamma^{n+1}$, y aplicar la definición
  (19.2.1) al homomorfismo \emph{sobreyectivo} $P\to Q$ y al homomorfismo
  $(g_\alpha,p_\gamma)$ de $\mathfrak J/\mathfrak J^2$ en $Q$). Esta elección de
  $f_\gamma$ y $g_\gamma$ permite entonces, construyendo
  $v_{\gamma\delta}$ como en \emph{(i)}, obtener además la conmutatividad
  del diagrama (19.5.6.4).
\end{enumerate}
\end{env}

\oldpage[0\textsubscript{IV-1}]{97}

\begin{lemma}[19.5.6.5]
\label{0.19.5.6.5-es}
Supongamos que las $A$-álgebras $B$ y $C$ son formalmente lisas y que
$B$ es preadmisible (de modo que, en virtud de (19.5.6.1), se satisfacen las
condiciones de (19.5.6.3)). Para todo sistema de dos índices
$\alpha\leq\beta$ y de un homomorfismo $v_{\alpha\beta}$ que verifica las
condiciones de (19.5.6.3, (i)), existen un índice $\lambda\geq\beta$ y un
$A$-homomorfismo de álgebras
\[
  w_{\beta\lambda}:E_{\lambda,n}\longrightarrow F_{\beta,n}
\]
tales que: $1^\circ$ $\gr^0(w_{\beta\lambda}):C_\lambda\to C_\beta$ y
$\gr^1(w_{\beta\lambda}):
\mathfrak J_\lambda/\mathfrak J_\lambda^2\to
\mathfrak J_\beta/\mathfrak J_\beta^2$ sean los homomorfismos de
transición; $2^\circ$ el compuesto
$E_{\lambda,n}\xrightarrow{w_{\beta\lambda}}F_{\beta,n}
\xrightarrow{v_{\alpha\beta}}E_{\alpha,n}$ sea el homomorfismo de transición.
\end{lemma}

\begin{proof}
Apliquemos el lema (19.5.6.3, (ii)) con $\gamma=\beta$, lo que da un
$\delta\geq\beta$ y un
$v_{\beta\delta}:F_{\delta,n}\to E_{\beta,n}$. Recordemos que
$v_{\beta\delta}$ es sobreyectivo; por otra parte, su núcleo $\mathfrak R$ es
\emph{nilpotente}: en efecto, el ideal de aumentación $F_{\delta,n}^+$ de
$F_{\delta,n}$ es nilpotente por definición, y el núcleo de
$\gr^0(v_{\beta\delta}):C_\delta\to C_\beta$ también es nilpotente en virtud de
la hipótesis de que $B$ es preadmisible. Se ve, por tanto, que $E_{\beta,n}$
se identifica con $F_{\delta,n}/\mathfrak R$. Como $B$ es una $A$-álgebra
formalmente lisa, el $A$-homomorfismo canónico
$p_{\beta,n}:B\to E_{\beta,n}$, que es continuo, se factoriza como
$B\xrightarrow{w}F_{\delta,n}\xrightarrow{v_{\beta\delta}}E_{\beta,n}$, donde
$w$ es un $A$-homomorfismo continuo. Como $F_{\delta,n}$ es discreto, existe
$\lambda\geq\delta$ tal que $w$ se anule en $\mathfrak b_\lambda$, luego $w$
se factoriza como
$B\to B/\mathfrak b_\lambda\xrightarrow{w'_\lambda}F_{\delta,n}$. Consideremos
el homomorfismo compuesto
$B/\mathfrak b_\lambda\xrightarrow{w'_\lambda}F_{\delta,n}
\xrightarrow{q_{\beta\delta}}F_{\beta,n}$, donde $q_{\beta\delta}$ es
el homomorfismo de transición; observemos que
$\gr^0(q_{\beta\delta})=\gr^0(v_{\beta\delta}):C_\delta\to C_\beta$ es
el homomorfismo de transición; como el compuesto
$v_{\alpha\beta}\circ q_{\beta\delta}\circ w'_\lambda$ es el homomorfismo
canónico $B/\mathfrak b_\lambda\to
B/(\mathfrak b_\alpha+\mathfrak J^{n+1})$, esto muestra que la imagen de
$\mathfrak J_\lambda=(\mathfrak b_\lambda+\mathfrak J)/\mathfrak b_\lambda$
por $q_{\beta\delta}\circ w'_\lambda$ está contenida en el ideal de aumentación
$F_{\beta,n}^+$, y, por tanto, la imagen de
$\mathfrak J_\lambda^{n+1}=(\mathfrak b_\lambda+\mathfrak J^{n+1})/
\mathfrak b_\lambda$ es nula. En otras palabras,
$q_{\beta\delta}\circ w'_\lambda$ se factoriza como
\[
  B/\mathfrak b_\lambda\to
  B/(\mathfrak b_\lambda+\mathfrak J^{n+1})=E_{\lambda,n}
  \xrightarrow{w_{\beta\lambda}}F_{\beta,n}
\]
de modo que el compuesto
$E_{\lambda,n}\xrightarrow{w_{\beta\lambda}}F_{\beta,n}
\xrightarrow{v_{\alpha\beta}}E_{\alpha,n}$ sea el homomorfismo de transición.
El razonamiento anterior muestra también que $\gr^0(w_{\beta\lambda})$, que es
el compuesto
$B/(\mathfrak b_\lambda+\mathfrak J)\to\gr^0(F_{\delta,n})
\xrightarrow{\gr^0(q_{\beta\delta})}\gr^0(F_{\beta,n})$, es el homomorfismo
de transición $C_\lambda\to C_\beta$, puesto que
$v_{\beta\delta}\circ w$ es el homomorfismo canónico. Además, se tiene también
$\gr^1(q_{\beta\delta})=\gr^1(v_{\beta\delta})$ (19.5.6.3), luego el mismo
razonamiento prueba que $\gr^1(w_{\beta\lambda})$ es el homomorfismo de
transición.
\end{proof}

\begin{env}[19.5.6.6]
\label{0.19.5.6.6-es}
\emph{Demostración de} (19.5.3, (ii)). --- Para demostrar que $\varphi_n$ es un
bimorfismo formal, vamos a aplicar el criterio de (19.1.3, (iii)). Puesto que
las condiciones del lema (19.5.6.5) se satisfacen por hipótesis, determinemos,
para todo índice $\alpha$, $v_{\alpha\beta}$ y $w_{\beta\lambda}$ que verifican
las conclusiones de este lema. El homomorfismo
\[
  \gr^n(v_{\alpha\beta}):\gr^n(F_{\beta,n})
  \longrightarrow\gr^n(E_{\alpha,n})
\]
no es sino el homomorfismo
\[
  \varphi_{\alpha\beta,n}:
  \mathbf S_{C_\beta}^n(\mathfrak J_\beta/\mathfrak J_\beta^2)
  \longrightarrow\mathfrak J_\alpha^n/\mathfrak J_\alpha^{n+1}
\]
deducido del homomorfismo canónico $\varphi_n$ (19.5.2.1) por paso a los
cocientes; en efecto, de (19.5.6.3) resulta que $\gr^0(v_{\alpha\beta})$ y
$\gr^1(v_{\alpha\beta})$ coinciden respectivamente con
$\varphi_{\alpha\beta,0}$ y $\varphi_{\alpha\beta,1}$, y la definición del
homomorfismo canónico $\varphi$ muestra entonces, por recurrencia sobre $j\leq n$,
que $\gr^j(v_{\alpha\beta})$ y $\varphi_{\alpha\beta,j}$ son iguales para todo
$j$. Dicho esto, como $\varphi_n$ es sobreyectivo, es \emph{a fortiori} un
epimorfismo formal; además, para $j\leq n$, el homomorfismo compuesto
\[
  \gr^j(F_{\lambda,n})\xrightarrow{\varphi_{\lambda\lambda,j}}
  \gr^j(E_{\lambda,n})\xrightarrow{\gr^j(w_{\beta\lambda})}
  \gr^j(F_{\beta,n})
\]
es el homomorfismo de transición, pues esto es cierto para $j=0$ y $j=1$ en
virtud de (19.5.6.5), y como $\gr^*(E_{\lambda,n})$ está engendrado por
$\gr^1(E_{\lambda,n})$, esto demuestra la afirmación por recurrencia sobre $j$.
Componiendo con el homomorfismo de transición
$\gr^j(F_{\beta,n})\to\gr^j(F_{\alpha,n})$, se ve, pues, para $j=n$, que
se ha factorizado el homomorfismo de transición
$\gr^n(F_{\lambda,n})\to\gr^n(F_{\alpha,n})$ en
\[
  \gr^n(F_{\lambda,n})\xrightarrow{\varphi_{\lambda\lambda,n}}
  \gr^n(E_{\lambda,n})\longrightarrow\gr^n(F_{\alpha,n}),
\]
lo que es la condición del criterio (19.1.3) para que $\varphi_n$ sea un
bimorfismo formal.
\end{env}

\oldpage[0\textsubscript{IV-1}]{98}

\begin{lemma}[19.5.6.7]
\label{0.19.5.6.7-es}
Supongamos que $B$ es preadmisible, que $\mathfrak J/\mathfrak J^2$ es un
$C$-módulo formalmente proyectivo, que $C$ es una $A$-álgebra formalmente
lisa y que los $\varphi_n$ son bimorfismos formales. Entonces, para todo
par de índices $\alpha$, $\beta$ tal que $\alpha\leq\beta$ y todo
homomorfismo $v_{\alpha\beta}:F_{\beta,n}\to E_{\alpha,n}$ que verifica las
condiciones de (19.5.6.3, (i)), existe un índice $\gamma\geq\beta$ tal que,
para todo índice $\mu\geq\gamma$, haya un diagrama conmutativo de
$A$-homomorfismos
\[
\xymatrix@C=4em@R=3em{
  F_{\beta,n}\ar[r]^{v_{\alpha\beta}} & E_{\alpha,n} \\
  & E_{\mu,n}\ar[ul]_{u_{\beta\mu}}\ar[u]_{p_{\alpha\mu}}
}
\]
donde $p_{\alpha\mu}$ es el homomorfismo de transición.
\end{lemma}

\begin{proof}
Aplicando el criterio (19.1.3) a cada uno de los $\varphi_j$ para $0\leq j\leq n$,
se ve que existe un índice $\gamma\geq\beta$ y homomorfismos
unívocamente determinados (y sobreyectivos) de $C_\gamma$-módulos
\[
  w_{\beta\gamma}^j:\gr^j(E_{\gamma,n})\longrightarrow
  \gr^j(F_{\beta,n})
\]
tales que los compuestos
\[
  \gr^j(F_{\gamma,n})
  \xrightarrow{\varphi_{\gamma\gamma,j}}\gr^j(E_{\gamma,n})
  \xrightarrow{w_{\beta\gamma}^j}\gr^j(F_{\beta,n})
\]
sean los homomorfismos de transición (el hecho de que se pueda elegir el mismo
índice $\gamma$ para todos los $\varphi_j$ resulta de que son en número
finito). Además, la unicidad de los $w_{\beta\gamma}^j$ prueba (puesto que $\varphi$ es
un homomorfismo de álgebras graduadas) que
$w_{\beta\gamma}=(w_{\beta\gamma}^j)_{0\leq j\leq n}$ es un homomorfismo de
$C_\gamma$-álgebras graduadas
\[
  \gr^*(E_{\gamma,n})\longrightarrow
  \gr^*(F_{\beta,n})=F_{\beta,n}.
\]
Además, como $\varphi_{\gamma\gamma,0}$ y $\varphi_{\gamma\gamma,1}$ son
los homomorfismos identidad,
$w_{\beta\gamma}^0:C_\gamma\to C_\beta$ y
$w_{\beta\gamma}^1:\mathfrak J_\gamma/\mathfrak J_\gamma^2
\to\mathfrak J_\beta/\mathfrak J_\beta^2$ son los homomorfismos de
transición, y lo mismo ocurre, por tanto, con
$\gr^0(v_{\alpha\beta}\circ w_{\beta\gamma})$ y
$\gr^1(v_{\alpha\beta}\circ w_{\beta\gamma})$; como
$\gr^*(E_{\gamma,n})$ está engendrado por $\gr^1(E_{\gamma,n})$, se concluye que
$\gr^j(v_{\alpha\beta}\circ w_{\beta\gamma})$ es también el homomorfismo de
transición para $0\leq j\leq n$.

Apliquemos ahora a $\alpha$, $\beta$ y $\gamma$ el lema (19.5.6.3,
(ii)), lo que da el diagrama (19.5.6.4) con $\delta\geq\gamma$; después
repitamos el razonamiento del comienzo reemplazando $\alpha$ y $\beta$ por
$\gamma$ y $\delta$. Se obtiene así un índice $\lambda\geq\delta$ y un
diagrama conmutativo de homomorfismos
\[
\xymatrix@C=4em@R=4em{
  \gr^*(E_{\gamma,n})\ar[r]^{w_{\beta\gamma}} &
  F_{\beta,n}\ar[r]^{v_{\alpha\beta}} & E_{\alpha,n} \\
  \gr^*(E_{\lambda,n})\ar[u]_{\gr^*(p_{\gamma\lambda})}
    \ar[r]_{w_{\delta\lambda}} &
  F_{\delta,n}\ar[u]_{q_{\beta\delta}}\ar[r]_{v_{\gamma\delta}} &
  E_{\gamma,n}\ar[u]_{p_{\alpha\gamma}}
    \ar@{-->}[ul]_{u_{\beta\gamma}}
}
\]
donde las flechas verticales son los homomorfismos de transición. Basta
probar la existencia del homomorfismo $u_{\beta\gamma}$ que hace
conmutativo el diagrama, y para ello basta evidentemente mostrar que se tiene
\[
  \operatorname{Ker}(v_{\gamma\delta})
  \subset\operatorname{Ker}(q_{\beta\delta}),
\]
siendo $v_{\gamma\delta}$ y $q_{\beta\delta}$ sobreyectivos. Como
$w_{\delta\lambda}$ es sobreyectivo, esta última relación equivale a
\[
\begin{split}
  \operatorname{Ker}(v_{\gamma\delta}\circ w_{\delta\lambda})
  &\subset\operatorname{Ker}(q_{\beta\delta}\circ w_{\delta\lambda})\\
  &=\operatorname{Ker}\bigl(w_{\beta\gamma}\circ
  (\gr(p_{\gamma\lambda}))\bigr).
\end{split}
\]
Pero se ha visto más arriba que
$\gr(p_{\gamma\lambda})=\gr(v_{\gamma\delta}\circ w_{\delta\lambda})$ y está
claro que
\[
\begin{split}
  \operatorname{Ker}(v_{\gamma\delta}\circ w_{\delta\lambda})
  &\subset\operatorname{Ker}\bigl(\gr(v_{\gamma\delta}\circ
  w_{\delta\lambda})\bigr)\\
  &=\operatorname{Ker}(\gr(p_{\gamma\lambda}))
  \subset\operatorname{Ker}\bigl(w_{\beta\gamma}\circ
  (\gr(p_{\gamma\lambda}))\bigr),
\end{split}
\]
lo que termina la demostración del lema, pues para todo $\mu\geq\gamma$ bastará
definir $u_{\beta\mu}$ como el compuesto
\[
  E_{\mu,n}\xrightarrow{p_{\gamma\mu}}E_{\gamma,n}
  \xrightarrow{u_{\beta\gamma}}F_{\beta,n}.
\]
\end{proof}

\begin{env}[19.5.6.8]
\label{0.19.5.6.8-es}
\emph{Demostración de} (19.5.3, (iii)). --- Sean $G$ una $A$-álgebra
topológica discreta, $\mathfrak R$ un ideal de cuadrado nulo en $G$ y
$f:B\to G/\mathfrak R$ un $A$-homomorfismo continuo de álgebras. Como
$G/\mathfrak R$ es discreta, hay un índice $\alpha$ tal que $f$ se anula
sobre $\mathfrak b_\alpha$; por hipótesis, existe un entero $m$ tal que
$\mathfrak J^m\subset\mathfrak b_\alpha$, por tanto $f$ se factoriza en
\[
  B\xrightarrow{p_\alpha}E_{\alpha,n}
  \xrightarrow{f_{\alpha,n}}G/\mathfrak R,
\]
donde se toma $n=2m$. Como se satisfacen las hipótesis del lema (19.5.6.3), se
obtiene primero un $\beta\geq\alpha$ y un $A$-homomorfismo compuesto
\begin{equation}
\label{0.19.5.6.9-es}
  f_{\alpha,n}\circ v_{\alpha\beta}:F_{\beta,n}
  \xrightarrow{v_{\alpha\beta}}E_{\alpha,n}
  \xrightarrow{f_{\alpha,n}}G/\mathfrak R.
  \tag{19.5.6.9}
\end{equation}
Como $F_{\beta,n}$ es una $C_\beta$-álgebra, y por tanto \emph{a fortiori} una
$C$-álgebra topológica (discreta), esto da por composición un
$A$-homomorfismo continuo
$r:C\to F_{\beta,n}\to G/\mathfrak R$. Como, por otra parte, $C$ es una
$A$-álgebra formalmente lisa,

\oldpage[0\textsubscript{IV-1}]{99}

este homomorfismo se factoriza en
$r:C\xrightarrow{r'}G\to G/\mathfrak R$, donde $r'$ es un
$A$-homomorfismo continuo, de modo que $G$ queda dotada así de una estructura de
$C$-álgebra topológica (discreta). Por otra parte, por composición con
el homomorfismo canónico
\[
  \mathfrak J/\mathfrak J^2\longrightarrow
  \mathfrak J_\beta/\mathfrak J_\beta^2\longrightarrow F_{\beta,n},
\]
se deduce de (19.5.6.9) un $C$-homomorfismo continuo
$g:\mathfrak J/\mathfrak J^2\to G/\mathfrak R$. Como
$\mathfrak J/\mathfrak J^2$ es un $C$-módulo formalmente proyectivo, $g$ se
factoriza en
\[
  \mathfrak J/\mathfrak J^2\xrightarrow{h}G\longrightarrow G/\mathfrak R,
\]
donde $h$ es un $C$-homomorfismo continuo. Como $G$ es discreta, existe un
índice $\gamma\geq\beta$ tal que la imagen por $h$ de
$((\mathfrak b_\gamma\cap\mathfrak J)+\mathfrak J^2)/\mathfrak J^2$ sea
nula, así como la imagen por $r'$ de
$(\mathfrak b_\gamma+\mathfrak J)/\mathfrak J$, de modo que $h$ se factoriza
en
\[
  \mathfrak J/\mathfrak J^2\longrightarrow
  \mathfrak J_\gamma/\mathfrak J_\gamma^2
  \xrightarrow{h'}G,
\]
donde $h'$ es un $C_\gamma$-homomorfismo. Por prolongación, se deduce, pues, de
$h'$ un homomorfismo de $C_\gamma$-álgebras
\[
  w:\mathbf S_{C_\gamma}^*(\mathfrak J_\gamma/\mathfrak J_\gamma^2)
  \longrightarrow G
\]
y, por construcción, todo elemento de grado $m$ tiene por imagen mediante $w$ un elemento
de $\mathfrak R$, por tanto todo elemento de grado $n=2m$ tiene imagen nula, puesto que
$\mathfrak R^2=0$; en otras palabras, $w$ se factoriza en
\[
  \mathbf S_{C_\gamma}^*(\mathfrak J_\gamma/\mathfrak J_\gamma^2)
  \longrightarrow F_{\gamma,n}\xrightarrow{w'_\gamma}G.
\]
Apliquemos ahora a
\[
  v_{\alpha\gamma}:F_{\gamma,n}
  \xrightarrow{q_{\beta\gamma}}F_{\beta,n}
  \xrightarrow{v_{\alpha\beta}}E_{\alpha,n}
\]
el lema (19.5.6.7), cuyas hipótesis se verifican; existe, pues, un
$\delta\geq\gamma$ y un homomorfismo
$u_{\gamma\delta}:E_{\delta,n}\to F_{\gamma,n}$ tal que el compuesto
$v_{\alpha\gamma}\circ u_{\gamma\delta}$ sea el homomorfismo de transición
$p_{\alpha\delta}:E_{\delta,n}\to E_{\alpha,n}$. Se obtiene así finalmente un
diagrama conmutativo de $A$-homomorfismos continuos
\[
\xymatrix@C=3em@R=2.5em{
  && E_{\alpha,n}\ar[r]^{f_{\alpha,n}} & G/\mathfrak R \\
  && F_{\beta,n}\ar[u]_{v_{\alpha\beta}} & \\
  B\ar[r]_{p_\delta} & E_{\delta,n}\ar[uur]^{p_{\alpha\delta}}
    \ar[r]_{u_{\gamma\delta}} &
  F_{\gamma,n}\ar[u]_{q_{\beta\gamma}}\ar[r]^{w'_\gamma} & G\ar[uu]
}
\]
y el compuesto $f':B\to G$ de los homomorfismos de la línea inferior es, pues,
tal que $f$ se factoriza en $B\xrightarrow{f'}G\to G/\mathfrak R$, lo que
prueba que $B$ es una $A$-álgebra formalmente lisa.

La demostración del teo.~(19.5.3) queda así terminada.
\end{env}

\begin{corollary}[19.5.7]
\label{0.19.5.7-es}
Sean $A$ un anillo topológico, $B$ una $A$-álgebra topológica,
$(\mathfrak b_\lambda)$ un sistema fundamental de ideales abiertos de $B$,
$\mathfrak J$ un ideal de $B$ y $C=B/\mathfrak J$ la $A$-álgebra topológica
cociente. Se pone $C_\lambda=B/(\mathfrak b_\lambda+\mathfrak J)$. Se supone
que: $1^\circ$ para todo $n$, la topología inducida sobre $\mathfrak J^n$ por
la de $B$ es también la topología del $B$-módulo $\mathfrak J^n$ deducida de
la topología de $B$ (19.0.2) (esta condición se satisfará en particular si
$B$ es noetheriano y su topología es preádica ($\mathbf 0_{\mathrm I}$,
7.3.2))~; $2^\circ$ $C$ es una $A$-álgebra formalmente lisa. En estas
condiciones~:
\begin{enumerate}
  \item[\rm(i)] Si $B$ es una $A$-álgebra formalmente lisa, entonces, para todo
  $\lambda$, $(\mathfrak J/\mathfrak J^2)\otimes_C C_\lambda$ es un
  $C_\lambda$-módulo proyectivo.
  \item[\rm(ii)] Si $B$ es una $A$-álgebra formalmente lisa y un anillo
  preadmisible, entonces para todo $\lambda$, el homomorfismo canónico
  \begin{equation}
  \label{0.19.5.7.1-es}
    \varphi_\lambda=\varphi\otimes 1_{C_\lambda}:
    \mathbf S_{C_\lambda}^*((\mathfrak J/\mathfrak J^2)\otimes_C C_\lambda)
    \longrightarrow\gr_{\mathfrak J}^*(B)\otimes_C C_\lambda
    \tag{19.5.7.1}
  \end{equation}
  es biyectivo.

  \oldpage[0\textsubscript{IV-1}]{100}

  \item[\rm(iii)] Recíprocamente, supongamos que $B$ sea preadmisible, que la
  sucesión $(\mathfrak J^n)$ tienda a $0$, y que, para todo $\lambda$,
  $(\mathfrak J/\mathfrak J^2)\otimes_C C_\lambda$ sea un $C_\lambda$-módulo
  proyectivo y el homomorfismo (19.5.4.1) sea biyectivo. Entonces $B$ es una
  $A$-álgebra formalmente lisa.
\end{enumerate}
\end{corollary}

\begin{proof}
En efecto, la topología de $\mathfrak J/\mathfrak J^2$ y las de los
$\mathbf S_C^n(\mathfrak J/\mathfrak J^2)$ son entonces también inducidas por la
de $B$ (19.5.1)~; las conclusiones son entonces consecuencias inmediatas de
(19.5.3) y de los criterios de (19.1.4) y (19.2.4) especiales para este tipo de
topologías.
\end{proof}

\begin{remark}[19.5.8]
\label{0.19.5.8-es}
Supongamos que $\mathfrak J/\mathfrak J^2$ sea un $C$-módulo de tipo finito y
que, para la topología cociente de la de $B$, $C$ sea un anillo de Zariski~;
sea $\mathfrak r$ un ideal de definición de $C$, de modo que los
$\mathfrak r^n$ formen un sistema fundamental de entornos de $0$ en $C$.
Entonces las conclusiones de (i) y (ii) en (19.5.7) pueden sustituirse por
las siguientes~:
\begin{enumerate}
  \item[\rm(i')] $\mathfrak J/\mathfrak J^2$ es un $C$-módulo proyectivo.
  \item[\rm(ii')] El homomorfismo canónico
  $\varphi:\mathbf S_C^*(\mathfrak J/\mathfrak J^2)
  \to\gr_{\mathfrak J}^*(B)$ es biyectivo.
\end{enumerate}
En efecto, está claro que (i') implica la conclusión de (i) en (19.5.7).
Recíprocamente, si $(\mathfrak J/\mathfrak J^2)\otimes_C C_\lambda$ es un
$C_\lambda$-módulo proyectivo para todo $\lambda$, entonces
$(\mathfrak J/\mathfrak J^2)\otimes_C(C/\mathfrak r^n)$ es un
$(C/\mathfrak r^n)$-módulo proyectivo (por tanto plano) para todo $n$~; se concluye
que $\mathfrak J/\mathfrak J^2$ es un $C$-módulo plano
($\mathbf 0_{\mathrm{III}}$, 10.2.2), por tanto proyectivo puesto que es de
presentación finita (Bourbaki, \emph{Alg. comm.}, cap.~II, \textsection5,
n\textsuperscript{o}~2, cor.~2 del teo.~1). Por otra parte, los $C$-módulos
$\mathbf S_C^n(\mathfrak J/\mathfrak J^2)$ y $\gr_{\mathfrak J}^n(B)$ son
de tipo finito, y se sabe que cuando $C$ es un anillo de Zariski, equivale
entonces decir que $\varphi_n$ es biyectivo o que
$\widehat\varphi_n$ es biyectivo (Bourbaki, \emph{Alg. comm.}, cap.~III,
\textsection3, n\textsuperscript{o}~5, prop.~9), por tanto (ii) equivale a
(ii').
\end{remark}

\subsection{Caso de las álgebras sobre un cuerpo}
\label{subsection:0.19.6-es}

\begin{theorem}[19.6.1) (Cohen]
\label{0.19.6.1-es}
Sean $k$ un cuerpo, $K$ una extensión de $k$, estando $k$ y $K$
provistos de las topologías discretas. Para que $K$ sea una $k$-álgebra
formalmente lisa, es necesario y suficiente que $K$ sea una extensión separable
de $k$.
\end{theorem}

\begin{proof}
La necesidad de la condición se establecerá en (19.6.5.1) (y naturalmente no
se utilizará hasta entonces)~; nos limitaremos aquí a probar que la
condición es \emph{suficiente}. Distingamos dos casos~:

\emph{I. --- $K$ es una extensión separable de tipo finito de $k$.} Se sabe
entonces (Bourbaki, \emph{Alg.}, cap.~V, \textsection9,
n\textsuperscript{o}~3, teo.~2) que existe una subextensión puramente trascendente
$K'=k(T_1,\ldots,T_n)$ de $K$ tal que $K$ es una extensión algebraica
finita separable de $K'$. Teniendo en cuenta (19.3.5, (iii)), se puede por tanto
reducir, bien al caso en que $K=K'$, bien al caso en que $K$ es una extensión algebraica finita de $k$.
En el primer caso, se sabe que $A=k[T_1,\ldots,T_n]$ es una $k$-álgebra
formalmente lisa (19.3.3), por tanto ocurre lo mismo con
$k(T_1,\ldots,T_n)=S^{-1}A$, donde $S=A-\{0\}$ (19.3.5, (iv)). En el segundo
caso, \emph{todos} los grupos de cohomología de Hochschild
$\mathrm H_k^n(K,L)$, para un $(K,K)$-bimódulo cualquiera $L$, son nulos~: en
efecto, si se considera la $k$-álgebra producto tensorial $C=K\otimes_k K$,
$L$ es un $C$-módulo izquierdo y el grupo de cohomología
$\mathrm H_k^n(K,L)$ es igual a $\operatorname{Ext}_C^n(K,L)$, donde $K$ es
considerado también como $(K,K)$-bimódulo,

\oldpage[0\textsubscript{IV-1}]{101}

por tanto como $C$-módulo izquierdo (M, IX, 4). Ahora bien, como $K$ es una extensión
finita separable de $k$, se sabe que $K\otimes_k K$ es producto directo
de extensiones de $k$, una de las cuales es el propio $K$ (Bourbaki, \emph{Alg.},
cap.~VIII, \textsection8, prop.~3)~; esto implica por tanto que $K$ es un
$C$-módulo \emph{proyectivo}, de donde nuestra afirmación. Toda $k$-extensión de $K$
de núcleo $L$ es, pues, $k$-trivial, y \emph{a fortiori} también lo son las
$k$-extensiones conmutativas, de donde resulta el teorema en este caso (19.4.4).

\emph{II. --- Caso general.} --- Con las notaciones de (18.4.5), se trata de
probar que $\operatorname{Hom}_K(\mathrm H_2(P'_\bullet),L)=0$ para
\emph{todo} $K$-espacio vectorial $L$, lo que significa evidentemente que
$\mathrm H_2(P'_\bullet)=0$. Si $K$ es unión de una familia filtrante de
subextensiones $K^{(\alpha)}$ de $k$, $P'_\bullet$ es \emph{límite
inductivo} de los complejos correspondientes $P'^{(\alpha)}_\bullet$, puesto que el
funtor $\varinjlim$ conmuta con el producto tensorial en la categoría de los
$k$-módulos~; por la exactitud del funtor $\varinjlim$ en esta categoría,
se tiene por tanto
$\mathrm H_2(P'_\bullet)=\varinjlim\mathrm H_2(P'^{(\alpha)}_\bullet)$.
Como la hipótesis de que $K$ es separable implica que lo mismo ocurre con toda
subextensión de $K$ y que $K$ es unión de las subextensiones de tipo finito,
la primera parte de la demostración implica efectivamente que $K$ es una
$k$-álgebra formalmente lisa. C.Q.F.D.
\end{proof}

\begin{corollary}[19.6.2]
\label{0.19.6.2-es}
Sean $A$ un anillo local separado y completo, $K$ su cuerpo residual, $k$ un
subcuerpo de $K$ tal que $K$ sea una extensión separable de $k$. Entonces
existe un subcuerpo $K'$ de $A$ que contiene a $k$, tal que la restricción a $K'$
del homomorfismo canónico $A\to K$ sea un isomorfismo de $K'$ sobre $K$.
\end{corollary}

\begin{proof}
Sea $\mathfrak m$ el ideal maximal de $A$. En virtud de la hipótesis y de
(19.6.1), $K$ es una $k$-álgebra formalmente lisa~; aplicando (19.3.10) y
sustituyendo $C$ por $A$ y $\mathfrak J$ por $\mathfrak m$, se ve que
el automorfismo identidad de $K=A/\mathfrak m$ se factoriza como
$K\xrightarrow{u}A\to A/\mathfrak m$, donde $u$ es un $k$-homomorfismo, por tanto
necesariamente un $k$-isomorfismo de $K$ sobre un subcuerpo $K'$ que contiene a
$k$.
\end{proof}

\begin{corollary}[19.6.3]
\label{0.19.6.3-es}
Sean $A$ un anillo local noetheriano completo, $\mathfrak m$ su ideal maximal,
$k$ su cuerpo residual. Las condiciones siguientes son equivalentes~:
\begin{enumerate}
  \item[a)] $A$ contiene un subcuerpo.
  \item[b)] Si $p$ es la característica de $k$, se tiene $pA=0$.
  \item[c)] Existe un cuerpo $K$ tal que $A$ es isomorfo a un anillo
  cociente de un anillo de series formales $B=K[[T_1,\ldots,T_n]]$.
\end{enumerate}
Cuando esto ocurre, se puede suponer que $A$ es isomorfo a $B/\mathfrak b$,
donde $\mathfrak b$ está contenido en el cuadrado del ideal maximal de $B$.
\end{corollary}

\begin{proof}
Es inmediato que $c)$ implica $a)$, pues $A$ es $\neq0$ puesto que es un
anillo local~; como $K$ es un cuerpo y el homomorfismo compuesto
$K\xrightarrow{\varphi}B\to A$ (donde $\varphi$ es la inyección canónica) no es
nulo, es inyectivo. Para ver que $a)$ implica $b)$, basta
observar que si $K$ es un subcuerpo de $A$, $K$ es isomorfo a un
subcuerpo de $k$ y tiene por consiguiente la misma característica $p$~; como $p.1=0$ en
$K$, y por tanto en $A$, se tiene $pA=0$. Recíprocamente, $b)$ implica $a)$, pues el
compuesto de los homomorfismos canónicos
$\mathbf Z\xrightarrow{j}A\to k$ tiene por núcleo $p\mathbf Z$, y se tiene
$j(p\mathbf Z)=0$, por tanto $\operatorname{Ker}(j)=p\mathbf Z$ y $A$ contiene un
anillo $A'$ isomorfo a $\mathbf Z/p\mathbf Z$ y que no interseca a
$\mathfrak m$~; si $p>0$, $A'$ es un cuerpo~; en caso contrario, puesto que todo

\oldpage[0\textsubscript{IV-1}]{102}

elemento de $A'$ es invertible en $A$, el cuerpo de fracciones de $A'$
(isomorfo a $\mathbf Q$) está también contenido en $A$.

Probemos finalmente que $a)$ implica $c)$~; la condición $a)$ implica que $A$
contiene un subcuerpo primo $k_0$, isomorfo al subcuerpo primo de $k$,
y del cual $k$ es por consiguiente una extensión separable. Aplicando (19.6.2), se ve
primero que existe un isomorfismo $f$ de $k$ sobre un subcuerpo $K$ de $A$.
Por otra parte, sea $(x_i)_{1\leq i\leq n}$ un sistema de elementos de
$\mathfrak m$ tal que las clases mód.~$\mathfrak m^2$ de los $x_i$ formen una
base (sobre $k$) de $\mathfrak m/\mathfrak m^2$. Puesto que $A$ es completo, existe
entonces un homomorfismo continuo de anillos
$u:k[[T_1,\ldots,T_n]]\to A$ tal que $u$ sea igual a $f$ en $k$ y que
$u(T_i)=x_i$ para todo $i$, y este homomorfismo es sobreyectivo en virtud de la
elección de los $x_i$ (Bourbaki, \emph{Alg. comm.}, cap.~III, \textsection2,
n\textsuperscript{o}~9, prop.~11).
\end{proof}

\begin{theorem}[19.6.4]
\label{0.19.6.4-es}
Sean $k$ un cuerpo, $A$ una $k$-álgebra local noetheriana, $\mathfrak m$ su
ideal maximal, $K=A/\mathfrak m$ su cuerpo residual, estando $A$ provista de la
topología $\mathfrak m$-preádica. Se supone que $K$ sea una extensión
separable de $k$. Entonces las condiciones siguientes son equivalentes~:
\begin{enumerate}
  \item[a)] $A$ es una $k$-álgebra formalmente lisa.
  \item[b)] El completado $\widehat A$ de $A$ es $k$-isomorfo a un anillo de
  series formales $K[[T_1,\ldots,T_n]]$ (cuya estructura de $k$-álgebra está
  definida por el homomorfismo compuesto
  \[
    k\xrightarrow{\varphi}K\xrightarrow{\psi}K[[T_1,\ldots,T_n]],
  \]
  donde $\psi$ es la inyección canónica y $\varphi$ el homomorfismo que define
  la estructura de extensión de $K$).
  \item[b')] $\widehat A$ es un anillo isomorfo a un anillo de series
  formales $K[[T_1,\ldots,T_n]]$.
  \item[c)] $A$ es un anillo local regular.
\end{enumerate}
\end{theorem}

\begin{proof}
El hecho de que $a)$ implica $b)$ es el caso particular de ((19.5.4),
equivalencia de $a)$ y $c)$), aplicado sustituyendo allí $A$ por $k$, $B$ por
$A$ y $\mathfrak J$ por $\mathfrak m$, teniendo en cuenta (19.6.1). Está claro
que $b)$ implica $b')$ y $b')$ implica $c)$ por (17.1.1). Por último, si $c)$ se
verifica, se deduce de (17.1.1, $a)$) y de ((19.5.4), equivalencia de $b)$ y
$a)$), aplicado a $K=A/\mathfrak m$ sustituyendo a $C$, que $A$ es una
$k$-álgebra formalmente lisa.
\end{proof}

\begin{corollary}[19.6.5]
\label{0.19.6.5-es}
Sean $k$ un cuerpo, $A$ una $k$-álgebra local noetheriana, $\mathfrak m$ su
ideal maximal, estando $A$ provista de la topología $\mathfrak m$-preádica.
Supongamos que $A$ sea una $k$-álgebra formalmente lisa~; entonces $A$ es
geométricamente regular sobre $k$, es decir (cf.~$\mathbf{IV}$, 6.7.6),
para toda extensión finita $k'$ de $k$, el anillo semilocal
$A'=A\otimes_k k'$ es regular (17.3.6). Además, si $K$ es el cuerpo
residual de $A$, $\widehat A$ es isomorfo a un anillo de series formales
$K[[T_1,\ldots,T_n]]$.
\end{corollary}

\begin{proof}
Como $\widehat A'=\widehat A\otimes_k k'$, se deduce de (19.3.6) y de
(17.1.5) (aplicado a los anillos locales de $A'$ en sus ideales maximales) que se
puede reducir al caso en que $A$ es completo~; entonces $A'$ es una $k'$-álgebra
formalmente lisa (19.3.5, (iii)) y por otra parte es producto directo de
$k'$-álgebras locales, que también son formalmente lisas (19.3.5, (v)). Se
reduce por tanto a probar que $A$ es \emph{regular}. Sea $k_0$ el subcuerpo
primo de $k$~; como $k$ es una extensión separable de $k_0$, es una
$k_0$-álgebra formalmente lisa (19.6.1) y en virtud de la hipótesis, ocurre lo
mismo con $A$ (19.3.5, (ii)). Como el cuerpo residual $K$ de $A$ es una
extensión separable de $k_0$, se puede entonces aplicar (19.6.4) a $A$ y $k_0$,
de donde resulta la conclusión.
\end{proof}

\oldpage[0\textsubscript{IV-1}]{103}
 
\begin{corollary}[19.6.5.1]
\label{0.19.6.5.1-es}
Sean $A$ un anillo local noetheriano, $k$ un subcuerpo de $A$ tal que $A$
sea una $k$-álgebra formalmente lisa (para su topología preádica). Entonces
todo cuerpo $K$ tal que $k\subset K\subset A$ es una extensión separable de
$k$.
\end{corollary}

\begin{proof}
En efecto, para toda extensión finita $k'$ de $k$, el anillo $K\otimes_k k'$
se identifica con un subanillo de $A'=A\otimes_k k'$; como $A'$ es un anillo
regular, es reducido, luego lo mismo ocurre con $K\otimes_k k'$, lo que prueba que $K$
es una extensión separable de $k$ (Bourbaki, \emph{Alg.}, cap.~VIII,
\textsection7, n\textsuperscript{o}~3, teorema~1).
\end{proof}

Se observará que esto demuestra que la condición del enunciado de (19.6.1) es
\emph{necesaria}.

\begin{remark}[19.6.5.2]
\label{0.19.6.5.2-es}
Si $K$ es una extensión \emph{no separable} de $k$, el anillo de series
formales $K[[T_1,\ldots,T_n]]$ \emph{no es, por tanto}, una $k$-álgebra
formalmente lisa (para la estructura de $k$-álgebra usual recordada en
(19.6.4)). Por otra parte, existen $k$-álgebras formalmente lisas $A$ que
son anillos locales noetherianos completos cuyo cuerpo residual $K$ es una
extensión \emph{no separable} de $k$ (por ejemplo, los completados de las
localizaciones de $B=k[T_1,\ldots,T_n]$ en ideales maximales $\mathfrak n$ tales
que $B/\mathfrak n$ sea una extensión algebraica finita no separable de $k$);
tal anillo no puede, pues, ser $k$-isomorfo a $K[[T_1,\ldots,T_n]]$, aunque
sea isomorfo a él en virtud de (19.6.5).

En ciertos casos se puede precisar el corolario (19.6.5):
\end{remark}

\begin{proposition}[19.6.6]
\label{0.19.6.6-es}
Sean $k$ un cuerpo, $A$ una $k$-álgebra local noetheriana, $\mathfrak m$ su
ideal maximal, estando $A$ provista de la topología $\mathfrak m$-preádica.
Supongamos que el cuerpo residual $K=A/\mathfrak m$ de $A$ sea tal que exista
una extensión radicial finita $k_0$ de $k$ para la cual el cuerpo residual del
anillo local $K\otimes_k k_0$ sea una extensión separable de $k_0$ (más
adelante expresaremos esta propiedad diciendo que $K$ es de
\emph{multiplicidad radicial finita} sobre $k$ ($\mathbf{IV}$, 4.7.4), y
veremos que esta condición se satisface siempre cuando $K$ es una
\emph{extensión de tipo finito} de $k$). Entonces las condiciones siguientes
son equivalentes:
\begin{enumerate}
  \item[a)] $A$ es una $k$-álgebra formalmente lisa.
  \item[b)] Existe una extensión radicial finita $k'$ de $k$ tal que
  $\widehat A\otimes_k k'$ sea $k'$-isomorfa a un álgebra de series
  formales $K'[[T_1,\ldots,T_n]]$, donde $K'$ es una extensión separable de
  $k'$.
  \item[b')] Existe una extensión finita $k'$ de $k$ tal que el anillo
  semilocal $\widehat A\otimes_k k'$ tenga al menos un componente local
  que sea $k'$-isomorfo a un álgebra de series formales
  $K'[[T_1,\ldots,T_n]]$, donde $K'$ es una extensión separable de $k'$.
  \item[c)] $A$ es geométricamente regular sobre $k$ (19.6.5).
\end{enumerate}
\end{proposition}

\begin{proof}
Observemos primero que, si $k'$ es una extensión radicial de $k$, solo hay un
ideal de $A'=A\otimes_k k'$ sobre $\mathfrak m$, formado por los elementos
cuya potencia $p^h$-ésima ($p$ es el exponente característico de $k$) pertenece a
$\mathfrak m$ para un $h$ conveniente (Bourbaki, \emph{Alg. comm.}, cap.~V,
\textsection2, n\textsuperscript{o}~3, lema~4); $A'$ es, por tanto, un anillo
local, y lo mismo ocurre con
$K\otimes_k k'=(A\otimes_k k')/(\mathfrak m\otimes_k k')$; además, los cuerpos
residuales de estos dos anillos son idénticos. Recordemos, por otra parte, que,
si $K$ es una extensión separable de $k$, entonces, para toda extensión finita
$k''$ de $k$, $K\otimes_k k''$ es producto directo de cuerpos (Bourbaki,
\emph{Alg.}, cap.~VIII, \textsection7, n\textsuperscript{o}~3, cor.~1 del
teorema~1), y, por consiguiente, $\mathfrak m\otimes_k k''$ es el radical de
$A\otimes_k k''$, y los cuerpos componentes de $K\otimes_k k''$ son los cuerpos
residuales en los ideales maximales de $A\otimes_k k''$; además,

\oldpage[0\textsubscript{IV-1}]{104}

estos cuerpos son extensiones separables de $k''$, puesto que, para toda
extensión $k_1$ de $k''$,
$K\otimes_k k_1=(K\otimes_k k'')\otimes_{k''}k_1$ es reducido
(\emph{loc. cit.}, teorema~1). Si se aplican estas observaciones al cuerpo $K$ de
multiplicidad radicial finita sobre $k$, se ve que, para \emph{toda} extensión
finita $k''$ de $k_0$, $(K\otimes_k k'')_{\mathrm{red}}$ es separable sobre
$k''$.

Pasemos a la demostración de (19.6.6). Es trivial que $b)$ implica
$b')$. Demostremos que $b')$ implica $a)$. Si se pone
$B'=K'[[T_1,\ldots,T_n]]$, la hipótesis de que $K'$ es separable sobre $k'$
implica, por las observaciones iniciales, que, para toda extensión finita $k''$ de
$k'$, los componentes del anillo semilocal completo $B'\otimes_{k'}k''$
(igual al anillo de series formales
$(K'\otimes_{k'}k'')[[T_1,\ldots,T_n]]$) son los anillos
$K_i''[[T_1,\ldots,T_n]]$, donde los $K_i''$ son los cuerpos componentes de
$K'\otimes_{k'}k''$. Como los componentes locales de $B'\otimes_{k'}k''$ son
también componentes locales de
\[
  \widehat A\otimes_k k''
  =(\widehat A\otimes_k k')\otimes_{k'}k'',
\]
se ve que la hipótesis $b')$ sigue verificándose cuando se reemplaza $k'$ por
cualquiera de sus extensiones finitas. En particular, se puede suponer que
$k'$ es \emph{cuasi-galoisiana}, luego extensión galoisiana de una extensión
radicial $k'_0$ de $k$; $A'=\widehat A\otimes_k k'$ puede escribirse entonces
$A'_0\otimes_{k'_0}k'$, donde $A'_0=\widehat A\otimes_k k'_0$ es un anillo
local completo. Se sabe entonces (Bourbaki, \emph{Alg.}, cap.~VIII,
\textsection8, prop.~4) que $A'$ es producto directo de anillos locales
\emph{isomorfos} a $A'_0$ como $k'_0$-álgebras. Ahora bien, uno de estos anillos
es, por hipótesis, $k'$-isomorfo a un anillo de series formales
$K'[[T_1,\ldots,T_n]]$, donde $K'$ es una extensión separable de $k'$, y, por
tanto, \emph{también} una extensión separable de $k'_0$; teniendo en cuenta
(19.6.4), $A'_0$ es, pues, una $k'_0$-álgebra formalmente lisa. Pero como
$k'_0$ es un $k$-módulo fielmente plano, proyectivo y de tipo finito, se deduce
de (19.4.7) que $\widehat A$ es una $k$-álgebra formalmente lisa.

Ya se ha visto (19.6.5) que $a)$ implica $c)$; queda, pues, verificar que
$c)$ implica $b)$. Ahora bien, si se verifica $c)$, entonces, para toda
extensión radicial finita $k'$ de $k$ que contenga $k_0$, $A\otimes_k k'$ es
un anillo local regular cuyo cuerpo residual $K'$ es una extensión separable de
$k'$; se deduce, por tanto, de (19.6.4) que su completado
$\widehat A\otimes_k k'$ es $k'$-isomorfo a un anillo de series formales
$K'[[T_1,\ldots,T_n]]$.
\end{proof}

\begin{remarks}[19.6.7]
\label{0.19.6.7-es}
\begin{enumerate}
  \item[\rm(i)] Se observará que la hipótesis de que $K$ es de multiplicidad
  radicial finita sobre $k$ solo se ha utilizado en la demostración de la
  implicación $b')\Rightarrow a)$.
  \item[\rm(ii)] Probaremos más adelante (22.5.8) que $a)$ y $c)$ son
  equivalentes, \emph{sin hipótesis} sobre la extensión $K$ de $k$.
\end{enumerate}
\end{remarks}

\subsection{Caso de los homomorfismos locales; teoremas de existencia y
unicidad}
\label{subsection:0.19.7-es}

En este número, cuando un anillo semilocal se considera como anillo
topológico, siempre se sobreentiende que se trata de su topología
$\mathfrak r$-preádica, donde $\mathfrak r$ es su radical. Todo homomorfismo
local de anillos locales es, por tanto, automáticamente \emph{continuo}.

\begin{theorem}[19.7.1]
\label{0.19.7.1-es}
Sean $A$, $B$ dos anillos locales noetherianos, $\mathfrak m$, $\mathfrak n$
sus ideales maximales respectivos, $k=A/\mathfrak m$ el cuerpo residual de $A$;
se suponen $A$ y $B$ provistos respectivamente de las topologías
$\mathfrak m$-preádica y $\mathfrak n$-preádica. Sea
$\varphi:A\to B$ un homomorfismo local, y pongamos $B_0=B\otimes_A k$.
Las propiedades siguientes son equivalentes:
\begin{enumerate}
  \item[a)] $B$ es una $A$-álgebra \emph{formalmente lisa}.

\oldpage[0\textsubscript{IV-1}]{105}

  \item[b)] $B$ es un $A$-módulo \emph{plano} y
  $B_0=B/\mathfrak mB$ (provisto de la topología cociente) es una $k$-álgebra
  \emph{formalmente lisa}.
\end{enumerate}
\end{theorem}

\begin{proof}
La demostración se realiza en varias etapas.

\medskip
\noindent\phantomsection\label{0.19.7.1.1-es}\textbf{(19.7.1.1)}\enspace
Demostremos primero que $b)$ implica $a)$; vamos a aplicar el criterio
(19.4.7), con $\mathfrak J=\mathfrak m$; en virtud de la segunda hipótesis
de $b)$, todo se reduce a demostrar que $B$ es un $A$-módulo
\emph{formalmente proyectivo}. La hipótesis implica que, para todo $h>0$,
$B/\mathfrak m^hB$ es un $(A/\mathfrak m^h)$-módulo plano
($\mathbf 0_{\mathrm{III}}$, 10.2.1); como los $\mathfrak m^h$ forman un
sistema fundamental de entornos de 0 en $A$ y
$(B/\mathfrak m^hB)\otimes_{A/\mathfrak m^h}k=B_0$, se pueden reemplazar $A$ y
$B$ por $A/\mathfrak m^h$ y $B/\mathfrak m^hB$ respectivamente y, por tanto,
suponer $A$ \emph{artiniano} (luego discreto). Como $B_0$ es una $k$-álgebra
formalmente lisa, es un anillo \emph{regular} (19.6.5); sea
$(x_i^0)_{1\leq i\leq n}$ un sistema regular de parámetros para $B_0$
(17.1.6), y, para todo $i$, sea $x_i\in B$ tal que $x_i^0$ sea su imagen
en $B_0=B/\mathfrak mB$; como los $x_i^0$ engendran el ideal maximal
$\mathfrak n_0=\mathfrak n/\mathfrak mB$ de $B_0$, los ideales
\[
  \mathfrak Q_m=\sum_{i=1}^n(x_i^0)^mB_0 \qquad (m>0)
\]
forman un sistema fundamental de entornos de 0 en $B_0$, pues
$\mathfrak Q_m$ está evidentemente contenido en $\mathfrak n_0^m$ y, por otra
parte, contiene $\mathfrak n_0^{mn}$. Pongamos
$\mathfrak J_m=\sum_{i=1}^n x_i^mB$ para todo $m>0$; es claro que
$\mathfrak n=\mathfrak J_1+\mathfrak mB$; como existe un $h>0$ tal que
$\mathfrak m^h=0$, se tiene
$\mathfrak n^m\subset\mathfrak J_1^{m-h}\subset\mathfrak n^{m-h}$ para
$m>h$, y, como se ha visto que $\mathfrak J_m\supset\mathfrak J_1^{mn}$, se ve
que los $\mathfrak J_m$ forman un sistema fundamental de entornos de 0 en
$B$. Todo se reduce, por tanto, a probar que los $B/\mathfrak J_m$ son
$A$-módulos \emph{libres}, y es lo mismo ver que son
$A$-módulos \emph{planos} ($\mathbf 0_{\mathrm{III}}$, 10.1.3). Ahora bien,
la hipótesis de que $(x_i^0)$ es una sucesión $B_0$-regular de elementos del ideal
maximal de $B_0$ implica la misma propiedad para la sucesión de los $(x_i^0)^m$
($1\leq i\leq n$) para todo $m>0$ (15.1.20); la conclusión resulta, por tanto, de
(15.1.16, $b)$ y $c)$).
\end{proof}

\begin{lemma}[19.7.1.2]
\label{0.19.7.1.2-es}
Sean $A$ un anillo topológico, $B$, $C$ dos $A$-álgebras topológicas que
son anillos locales noetherianos. Se supone además que $C$ es completo
y que el cuerpo residual $B/\mathfrak m$ de $B$ es un $A$-módulo de tipo finito.
Sea $E$ el producto tensorial completado $B\widehat\otimes_A C$. Entonces:
\begin{enumerate}
  \item[(i)] $E$ es un anillo semilocal noetheriano completo.
  \item[(ii)] El ideal $\mathfrak mE$ está contenido en el radical de $E$, y
  para todo $h>0$, $E/\mathfrak m^hE$ es isomorfo a
  $(B/\mathfrak m^h)\otimes_A C$.
  \item[(iii)] Si $C$ es un $A$-módulo plano, $E$ es un $B$-módulo plano.
\end{enumerate}
\end{lemma}

\begin{proof}
Por definición, $E$ es el completado separado del producto tensorial $B\otimes_A C$
para la topología definida por los ideales
$\operatorname{Im}(\mathfrak m^i\otimes_A C)+
\operatorname{Im}(B\otimes_A\mathfrak n'{}^j)$
($\mathbf 0_{\mathrm I}$, 7.7.5). Si se pone
\[
  \mathfrak r=\operatorname{Im}(\mathfrak m\otimes_A C)+
  \operatorname{Im}(B\otimes_A\mathfrak n'),
\]
se tiene
$\mathfrak r^{2i}\subset
\operatorname{Im}(\mathfrak m^i\otimes_A C)+
\operatorname{Im}(B\otimes_A\mathfrak n'{}^i)\subset\mathfrak r^i$, luego $E$
es también el completado separado de $B\otimes_A C$ para la topología
$\mathfrak r$-preádica. Por hipótesis,
$(B\otimes_A C)/\mathfrak r=(B/\mathfrak m)\otimes_A(C/\mathfrak n')$ es un
$(C/\mathfrak n')$-módulo de tipo finito, luego un anillo artiniano; además,
$\mathfrak r/\mathfrak r^2$, al ser cociente de
\[
  ((\mathfrak m/\mathfrak m^2)\otimes_A C)
  \oplus(B\otimes_A(\mathfrak n'/\mathfrak n'{}^2)),
\]
es un $(B\otimes_A C)$-módulo de tipo finito; aplicando
($\mathbf 0_{\mathrm I}$, 7.2.11), se ve, por tanto, que $E$ es un anillo
\emph{noetheriano}; además, como $E/\mathfrak rE$, isomorfo a
$(B\otimes_A C)/\mathfrak r$, es artiniano, $E$ es \emph{semilocal}.
Observemos ahora que $E$, que es isomorfo a
\[
  \varprojlim_{i,j}((B/\mathfrak m^i)\otimes_A(C/\mathfrak n'{}^j)),
\]
es también isomorfo, por el teorema del doble límite proyectivo, a
\[
  \varprojlim_i\left(\varprojlim_j
  ((B/\mathfrak m^i)\otimes_A(C/\mathfrak n'{}^j))\right)~;
\]

\oldpage[0\textsubscript{IV-1}]{106}

pero $\varprojlim_j((B/\mathfrak m^i)\otimes_A(C/\mathfrak n'{}^j))$ es el
completado separado de $(B/\mathfrak m^i)\otimes_A C$, y como $C$ es completo,
no es sino el propio $(B/\mathfrak m^i)\otimes_A C$,
siendo $B/\mathfrak m^i$ un $A$-módulo de tipo finito, puesto que
$\mathfrak m/\mathfrak m^i$ es un $(B/\mathfrak m)$-módulo de tipo finito
($\mathbf 0_{\mathrm I}$, 7.3.6). Se tiene, pues,
\[
  E=\varprojlim_i((B/\mathfrak m^i)\otimes_A C).
\]
Para todo entero $h>0$, $\mathfrak m^hE$, al ser un ideal de $E$, es cerrado
en $E$ ($\mathbf 0_{\mathrm I}$, 7.3.5), luego completo, y por otra parte es
evidentemente denso en
\[
  \varprojlim_i
  (\operatorname{Im}(\mathfrak m^h/\mathfrak m^{h+i})\otimes_A C),
\]
luego igual a este último límite proyectivo. Además, todos los sistemas
proyectivos considerados están definidos por homomorfismos \emph{sobreyectivos};
resulta, pues, de ($\mathbf 0_{\mathrm{III}}$, 13.2.2) que
$E/\mathfrak m^hE$ es isomorfo a
\[
  (B/\mathfrak m^h)\otimes_A C
  =(B\otimes_A C)/\operatorname{Im}(\mathfrak m^h\otimes_A C).
\]
En particular, como
$\operatorname{Im}(\mathfrak m\otimes_A C)\subset\mathfrak r$, esto muestra
que $\mathfrak mE$ está contenido en $\mathfrak rE$, y por tanto en el radical de
$E$. Finalmente, la hipótesis de que $C$ es un $A$-módulo plano implica que
$(B/\mathfrak m^h)\otimes_A C=E/\mathfrak m^hE$ es un
$(B/\mathfrak m^h)$-módulo plano para todo $h>0$: como $B$ y $E$ son
noetherianos y $\mathfrak mE$ está contenido en el radical de $E$,
resulta de ($\mathbf 0_{\mathrm{III}}$, 10.2.2) que $E$ es un $B$-módulo
plano.
\end{proof}

\begin{lemma}[19.7.1.3]
\label{0.19.7.1.3-es}
Sean $A$ un anillo local noetheriano, $\mathfrak m$ su ideal maximal, $k$ su
cuerpo residual y $B_0$ una $k$-álgebra; se supone que $B_0$ es un anillo
local noetheriano, completo y regular. Entonces existe una $A$-álgebra
topológica $B$ que es un anillo local noetheriano completo, un $A$-módulo plano,
y tal que $B_0$ sea $k$-isomorfo a
$B\otimes_A k=B/\mathfrak mB$.
\end{lemma}

\begin{proof}
Como $\widehat A$ es plano sobre $A$ y tiene el mismo cuerpo residual, podemos limitarnos
al caso en que $A$ es \emph{completo}.

Sea $K$ el cuerpo residual de $B_0$, y distingamos dos casos:

\medskip
\noindent I) $K$ es una extensión \emph{separable} de $k$. En virtud de
(19.6.4), $B_0$ es $k$-isomorfo a un anillo de series formales
$K[[T_1,\ldots,T_n]]$. Cuando $B_0=K$, el lema ya ha sido demostrado
($\mathbf 0_{\mathrm{III}}$, 10.3.1); sea $C$ un anillo local noetheriano
completo que es un $A$-módulo plano y tal que $C\otimes_A k$ sea isomorfo a
$K$. Para $n\geq1$, basta tomar (con la notación precedente)
$B=C[[T_1,\ldots,T_n]]$; se sabe, en efecto (Bourbaki,
\emph{Alg. conm.}, cap.~III, §~3, n\textsuperscript{o}~4, cor.~3 del teo.~1),
que $B$ es un $C$-módulo plano, luego también un $A$-módulo plano, y por otra parte,
es inmediato que $C[[T_1,\ldots,T_n]]\otimes_A k$ es isomorfo a
$(C/\mathfrak mC)[[T_1,\ldots,T_n]]=B_0$.

\medskip
\noindent II) $K$ es de \emph{característica} $p>0$, y por consiguiente lo mismo
ocurre con $k$. Denotemos por $P$ el cuerpo primo $\mathbf F_p$, y por $W(P)$ el anillo
local completo de los números $p$-ádicos $\mathbf Z_p$, que es un anillo de
valoración discreta (por tanto \emph{regular}), y tiene a $P$ como cuerpo residual.
Mostremos primero que existe un homomorfismo continuo de anillos $W(P)\to A$
que hace así de $A$ una $W(P)$-\emph{álgebra topológica}. En efecto, si
$j:\mathbf Z\to A$ es el homomorfismo canónico, se tiene
$j(p\mathbf Z)\subset\mathfrak m$ por hipótesis, de donde
$j^{-1}(\mathfrak m)=p\mathbf Z$, y por consiguiente $j$ se factoriza en
$\mathbf Z\to\mathbf Z_{p\mathbf Z}\xrightarrow{f}A$, donde $f$ es un
homomorfismo local, por tanto continuo, que (puesto que $A$ es completo) se prolonga
por continuidad al homomorfismo $W(P)\to A$ buscado.

Como $k$ es una extensión separable de $P$, el caso I) muestra que hay un
homomorfismo local $W(P)\to W(k)$, donde $W(k)$ es un anillo local noetheriano
completo y un $W(P)$-módulo plano, tal que $W(k)\otimes_{W(P)}P$ sea isomorfo
a $k$. Además, como el uniformizante $p$ de $W(P)$ es un elemento
$W(k)$-regular por planitud ($\mathbf 0_{\mathrm I}$, 6.3.4) y como

\oldpage[0\textsubscript{IV-1}]{107}

$W(k)/pW(k)=k$, $pW(k)$ es el ideal maximal de $W(k)$, lo que implica que este
último anillo es un \emph{anillo completo de valoración discreta} (Bourbaki,
\emph{Alg. conm.}, cap.~VI, §~3, n\textsuperscript{o}~5, prop.~9). Por
(19.7.1.1) se ve además (puesto que $k$ es separable sobre $P$, y por tanto una
$P$-álgebra formalmente lisa (19.6.1)) que $W(k)$ es una $W(P)$-álgebra
\emph{formalmente lisa}. El $W(P)$-homomorfismo continuo $W(k)\to k$ se
factoriza, pues, en $W(k)\to A\to k$ (19.3.11), lo que permite considerar $A$
como una $W(k)$-álgebra topológica. Aplicando ahora el caso I) a $B_0$
considerado como $P$-álgebra y a $W(P)$, se ve que existe una
$W(P)$-álgebra $B_P$ que es un anillo local noetheriano completo, un
$W(P)$-módulo plano, y tal que $B_P\otimes_{W(P)}P$ sea $P$-isomorfo a
$B_0$. Utilizando de nuevo el hecho de que $W(k)$ es una $W(P)$-álgebra
formalmente lisa, se ve por (19.3.11) que el homomorfismo compuesto
$W(k)\to k\to B_0$ se factoriza en $W(k)\to B_P\to B_0$; además, como
$k=W(k)/pW(k)$, se tiene
$B_P\otimes_{W(k)}k=B_P/pB_P=B_P\otimes_{W(P)}P=B_0$. Mostremos que $B_P$ es
un $W(k)$-módulo \emph{plano}; como $W(k)$ es un anillo de valoración discreta
cuyo uniformizante es $p$, basta verificar que $B_P$ es un
$W(k)$-módulo \emph{sin torsión} ($\mathbf 0_{\mathrm I}$, 6.3.4), o también
que $p$ es un elemento $B_P$-regular, lo que resulta de que $B_P$ es un
$W(P)$-módulo plano ($\mathbf 0_{\mathrm I}$, 6.3.4).

Pongamos ahora
\[
  B=B_P\widehat\otimes_{W(k)}A
\]
y observemos que, al ser el cuerpo residual de $A$ igual al de $W(k)$, es
\emph{a fortiori} un $W(k)$-módulo de tipo finito. Resulta, pues, en primer lugar
de (19.7.1.2) que $B$ es un anillo semilocal noetheriano completo,
estando $\mathfrak mB$ contenido en el radical de $B$; además,
$B/\mathfrak mB$ es $k$-isomorfo a
$B_P\otimes_{W(k)}k=B_0$, luego $B$ es en realidad un anillo local. Como $B_P$
es un $W(k)$-módulo plano, (19.7.1.2) muestra finalmente que $B$ es un $A$-módulo
plano. C.Q.F.D.
\end{proof}

\begin{lemma}[19.7.1.4]
\label{0.19.7.1.4-es}
Sean $A$ un anillo, $\mathfrak J$ un ideal de $A$, y $M$, $N$ dos $A$-módulos
separados para la topología $\mathfrak J$-preádica. Se supone además que $N$
es completo para la topología $\mathfrak J$-preádica y que $M$ es un
$A$-módulo plano. Sea $u:N\to M$ un $A$-homomorfismo; si
$u\otimes1:N\otimes_A(A/\mathfrak J)\to M\otimes_A(A/\mathfrak J)$ es
biyectivo, entonces $u$ es biyectivo.
\end{lemma}

\begin{proof}
Tomando los módulos graduados asociados con respecto a las filtraciones
$\mathfrak J$-preádicas, resulta de las hipótesis sobre $M$ y $N$ relativas
a las topologías $\mathfrak J$-preádicas que basta probar que
$\gr(u):\gr_\bullet(N)\to\gr_\bullet(M)$ es \emph{biyectivo} (Bourbaki,
\emph{Alg. conm.}, cap.~III, §~2, n\textsuperscript{o}~8, cor.~3 del teo.~1).
Ahora bien, se tiene un diagrama conmutativo
\[
\xymatrix{
  \gr_0(N)\otimes_{A_0}\gr_\bullet(A)
    \ar[r]^-{\gr_0(u)\otimes1}\ar[d]_{\varphi_N} &
  \gr_0(M)\otimes_{A_0}\gr_\bullet(A)\ar[d]^{\varphi_M} \\
  \gr_\bullet(N)\ar[r]_-{\gr(u)} & \gr_\bullet(M)
}
\]

\oldpage[0\textsubscript{IV-1}]{108}

donde $A_0=A/\mathfrak J=\gr_0(A)$, y $\varphi_M$ y $\varphi_N$ son las
aplicaciones canónicas ($\mathbf 0_{\mathrm{III}}$, 10.1.1.2). Por
hipótesis, $\gr_0(u)$ es biyectivo, al igual que $\varphi_M$
($\mathbf 0_{\mathrm{III}}$, 10.2.1), y $\varphi_N$ es sobreyectivo; se
deduce primero que $\gr_0(u)\otimes1$ es biyectivo, después que $\varphi_N$ es
inyectivo, luego biyectivo, y finalmente que $\gr(u)$ es biyectivo.
\end{proof}

\begin{lemma}[19.7.1.5]
\label{0.19.7.1.5-es}
Sean $A$ un anillo noetheriano, $\mathfrak J$ un ideal de $A$, y $B$, $B'$ dos
$A$-álgebras que son anillos locales noetherianos, siendo los homomorfismos
$A\to B$, $A\to B'$ continuos para la topología $\mathfrak J$-preádica
sobre $A$. Se supone que:
\begin{enumerate}
  \item[$1^{\mathrm o}$] $B$ y $B'$ son completos para las topologías
  $\mathfrak J$-preádicas;
  \item[$2^{\mathrm o}$] $B$ es una $A$-álgebra formalmente lisa;
  \item[$3^{\mathrm o}$] $B'$ es un $A$-módulo plano.
\end{enumerate}
Pongamos $A_0=A/\mathfrak J$, y sea
$u_0:B\otimes_A A_0\to B'\otimes_A A_0$ un $A_0$-isomorfismo; entonces
existe un $A$-isomorfismo $u:B\to B'$ tal que $u_0=u\otimes1$ (lo que
implica que $B'$ es una $A$-álgebra formalmente lisa y $B$ un $A$-módulo
plano).
\end{lemma}

\begin{proof}
Pongamos $B_0=B\otimes_A A_0$, $B'_0=B'\otimes_A A_0$. Observemos que, si
$\mathfrak m$ y $\mathfrak m'$ son los ideales maximales de $B$ y $B'$, las
topologías $\mathfrak J$-preádicas sobre $B$ y $B'$ son separadas, puesto que
$\mathfrak JB\subset\mathfrak m$ y $\mathfrak JB'\subset\mathfrak m'$; además,
como $\mathfrak JB'$ es cerrado en $B'$ para la topología
$\mathfrak J$-preádica, el homomorfismo compuesto
$B\to B_0\xrightarrow{u_0}B'_0$, que es continuo para las topologías
$\mathfrak J$-preádicas, se factoriza en $B\xrightarrow{u}B'\to B'_0$, donde
$u$ es un $A$-homomorfismo continuo (19.3.10). Se tiene evidentemente
$u_0=u\otimes1$, y la hipótesis de que $u_0$ es biyectivo implica que lo mismo
ocurre con $u$ en virtud de (19.7.1.4).
\end{proof}

\medskip
\noindent\textbf{(19.7.1.6)}\enspace
\emph{Fin de la demostración.} --- Para terminar de probar (19.7.1), hay que
mostrar que $a)$ implica $b)$; ya se sabe que $a)$ implica que $B_0$ es
una $k$-álgebra formalmente lisa (19.3.5, (iii)), por tanto basta probar
que $B$ es un $A$-módulo plano. Equivale a establecer que $\widehat B$
es un $\widehat A$-módulo plano (Bourbaki, \emph{Alg. conm.}, cap.~III, §~5,
n\textsuperscript{o}~4, prop.~4), y se sabe que $\widehat B$ es una
$\widehat A$-álgebra formalmente lisa (19.3.6); podemos, pues, limitarnos al
caso en que $A$ y $B$ son completos. Como $B_0$ es una $k$-álgebra formalmente
lisa, es un anillo regular (19.6.5) y completo
($\mathbf 0_{\mathrm I}$, 6.3.5); aplicando (19.7.1.3), se ve que existe
una $A$-álgebra $B'$ que es un anillo local noetheriano completo y un
$A$-módulo plano, un homomorfismo local $A\to B'$ y un $k$-isomorfismo
$B_0\simeq B'_0=B'\otimes_A k$. Basta entonces aplicar (19.7.1.5),
tomando como $\mathfrak J$ el ideal maximal de $A$, para obtener que $B$ es
$A$-isomorfo a $B'$, y por tanto es un $A$-módulo plano. C.Q.F.D.

\begin{theorem}[19.7.2]
\label{0.19.7.2-es}
Sean $A$ un anillo local noetheriano, $\mathfrak J$ un ideal contenido en
el ideal maximal de $A$, $A_0=A/\mathfrak J$, $B_0$ un anillo local noetheriano
completo y $A_0\to B_0$ un homomorfismo local que hace de $B_0$ una
$A_0$-álgebra formalmente lisa. Entonces existe un anillo local noetheriano
completo $B$, un homomorfismo local $A\to B$ que hace de $B$ un $A$-módulo
plano, y un $A_0$-isomorfismo
$u:B\otimes_A A_0\xrightarrow{\sim}B_0$. Si $(B',u')$ es un par
que satisface las mismas condiciones que $(B,u)$, existe un $A$-isomorfismo
$v:B\xrightarrow{\sim}B'$ que hace conmutativo el diagrama
\[
\xymatrix{
  B\otimes_A A_0\ar[rr]^{v\otimes1}_{\sim}\ar[dr]^{\sim}_u &&
  B'\otimes_A A_0\ar[dl]^{u'}_{\sim} \\
  & B_0
}
\]
\end{theorem}

\oldpage[0\textsubscript{IV-1}]{109}

\begin{proof}
Sea $\mathfrak m$ el ideal maximal de $A$, de modo que
$\mathfrak m_0=\mathfrak m/\mathfrak J$ es el ideal maximal de $A_0$, teniendo $A$ y
$A_0$ el mismo cuerpo residual $k$. Pongamos
$B_{00}=B_0\otimes_{A_0}k$; como $B_{00}$ es una $k$-álgebra formalmente
lisa (19.3.5, (iii)), es un anillo local regular (19.6.5); aplicando
(19.7.1.3), se ve que existe una $A$-álgebra topológica $B$ que es un
anillo local noetheriano completo, un $A$-módulo plano, y para la cual se tiene un
$k$-isomorfismo
$u_{00}:B\otimes_A k\xrightarrow{\sim}B_{00}$. Observemos que, en virtud de (19.7.1),
$B$ es una $A$-álgebra formalmente lisa, luego
$B\otimes_A A_0=B/\mathfrak JB$ es una $A_0$-álgebra formalmente lisa
(19.3.5, (iii)) y un anillo local noetheriano completo; además, $B_0$ es un
$A_0$-módulo plano en virtud de la hipótesis y de (19.7.1); como se tiene un
$k$-isomorfismo
\[
  u_{00}:B\otimes_A k=(B\otimes_A A_0)\otimes_{A_0}k
  \xrightarrow{\sim}B_0\otimes_{A_0}k=B_{00},
\]
se deduce de (19.7.1.5), aplicado sustituyendo allí $A$ por $A_0$ y
$\mathfrak J$ por $\mathfrak m_0$, que existe un $A_0$-isomorfismo
$u:B\otimes_A A_0\xrightarrow{\sim}B_0$ tal que $u_{00}=u\otimes1$. En cuanto a
la afirmación de unicidad, nótese que los ideales $\mathfrak J^hB$ (resp.
$\mathfrak J^hB'$) son cerrados en $B$ (resp. $B'$)
($\mathbf 0_{\mathrm I}$, 7.3.5), luego $B$ y $B'$ son separados y completos
para las topologías $\mathfrak J$-preádicas (Bourbaki, \emph{Top. gén.},
cap.~III, 3\textsuperscript{a} ed., §~3, n\textsuperscript{o}~5, cor.~2 de la
prop.~9); se tiene por hipótesis un $A_0$-isomorfismo
$v_0:B\otimes_A A_0\xrightarrow{\sim}B'\otimes_A A_0$ tal que
$u'\circ v_0=u$; como $B$ es una $A$-álgebra formalmente lisa y $B'$ un
$A$-módulo plano, se puede aplicar (19.7.1.5), de donde resulta la existencia del
$A$-isomorfismo $v$ requerido.
\end{proof}

\begin{remarks}[19.7.3]
\label{0.19.7.3-es}
\begin{enumerate}
  \item[(i)] Se observará que la afirmación de unicidad de (19.7.2) sigue siendo
  válida si se supone solamente que $B$ y $B'$ son completos \emph{para
  las topologías $\mathfrak J$-preádicas}. Ignoramos si se puede
  mejorar del mismo modo la afirmación de existencia, es decir, si se puede
  prescindir de suponer que el anillo local $B_0$ es \emph{completo} (para su topología
  $\mathfrak n_0$-preádica, denotando por $\mathfrak n_0$ su ideal
  maximal) exigiendo solamente que $B$ sea completo \emph{para la topología
  $\mathfrak J$-preádica}. Cuando $A_0$ es completo para la topología
  $\mathfrak m$-preádica, se puede ver que este problema se reduce al siguiente:
  si $B_0$ es un anillo local noetheriano regular (no necesariamente
  completo) que contiene el cuerpo primo $\mathbf F_p=\mathbf Z/p\mathbf Z$,
  ¿existe para todo $n>1$ una $(\mathbf Z/p^n\mathbf Z)$-álgebra plana $B$
  tal que $B/pB$ sea isomorfa a $B_0$?
  \item[(ii)] Se observará que, en general, el isomorfismo $v$ cuya existencia se
  afirma en (19.7.2) \emph{no es único} (cf.~(19.8.7)).
\end{enumerate}
\end{remarks}

\subsection{Álgebras de Cohen y $p$-anillos de Cohen; aplicación a la
estructura de los anillos locales completos}
\label{subsection:0.19.8-es}

Los resultados de esta sección son aplicaciones inmediatas de los teoremas
de (19.7), pero merecen explicitarse en razón de su importancia
práctica.

\begin{definition}[19.8.1]
\label{0.19.8.1-es}
Sean $A$, $B$ dos anillos locales noetherianos, $\mathfrak m$ el ideal maximal
de $A$, $k=A/\mathfrak m$ su cuerpo residual,
$\varphi:A\to B$ un homomorfismo local, que hace de $B$ una $A$-álgebra. Se
dice que $B$ es una \emph{$A$-álgebra de Cohen} si satisface las condiciones
siguientes:
\begin{enumerate}
  \item[(i)] $B$ es un anillo \emph{completo}.
  \item[(ii)] $B$ es un $A$-módulo \emph{plano}.
  \item[(iii)] $B\otimes_A k$ es un \emph{cuerpo} (es decir,
  $\mathfrak mB$ es el ideal maximal de $B$) que es una \emph{extensión
  separable} de $k$.
\end{enumerate}
\end{definition}

\oldpage[0\textsubscript{IV-1}]{110}

\begin{theorem}[19.8.2]
\label{0.19.8.2-es}
Sean $A$ un anillo local noetheriano y $k$ su cuerpo residual.
\begin{enumerate}
  \item[(i)] Si $B$ es una $A$-álgebra de Cohen, $B$ es una $A$-álgebra
  formalmente lisa. Para todo anillo local noetheriano completo $C$, todo
  homomorfismo local $A\to C$ y todo ideal $\mathfrak J\neq C$ de $C$,
  todo $A$-homomorfismo $B\to C/\mathfrak J$ se factoriza, por tanto, en
  $B\xrightarrow{v}C\to C/\mathfrak J$, donde $v$ es un $A$-homomorfismo
  (necesariamente local).
  \item[(ii)] Para todo cuerpo $K$, extensión separable de $k$, existe una
  $A$-álgebra de Cohen $B$ tal que $B\otimes_A k$ sea $k$-isomorfa a $K$,
  y tal $A$-álgebra es única salvo isomorfismo.
\end{enumerate}
\end{theorem}

\begin{proof}
Como $K$ es una $k$-álgebra formalmente lisa (19.6.1), la afirmación (i)
resulta de (19.7.1). Para demostrar (ii), se puede reducir al caso en que $A$ es
\emph{completo}, pues equivale decir que $B$ es un $A$-módulo plano
o un $\widehat A$-módulo plano ($\mathbf 0_{\mathrm{III}}$, 10.2.3), se tiene
$\mathfrak mB=\mathfrak m\widehat AB$ y $k$ es el cuerpo residual de
$\widehat A$. Basta entonces aplicar (19.7.2) tomando
$\mathfrak J=\mathfrak m$ y $B_0=K$ (y utilizando (19.6.1)).
\end{proof}

\begin{definition}[19.8.3]
\label{0.19.8.3-es}
Se llama \emph{anillo local primo} a un anillo local de la forma
$\mathbf Z_{p\mathbf Z}$, donde $p\mathbf Z$ es un ideal primo de $\mathbf Z$.
Se llama \emph{anillo local completo primo} al completado de un anillo local
primo.
\end{definition}

Los anillos locales primos son, por tanto, de dos tipos:
\begin{enumerate}
  \item[$1^{\mathrm o}$] Los que corresponden a los ideales maximales
  $p\mathbf Z$, donde $p\neq0$ es un número primo; $\mathbf Z_{p\mathbf Z}$
  es un anillo de valoración discreta, cuyo completado es \emph{el anillo de los
  enteros $p$-ádicos}, denotado ordinariamente por $\mathbf Z_p$\footnote{Esta notación,
  utilizada actualmente de manera universal, está en este caso en conflicto con la
  notación $A_f$ adoptada en ($\mathbf 0_{\mathrm I}$, 1.2.3): con
  $A=\mathbf Z$ y $f=p$, $A_f$ significa, en efecto, el anillo de los números
  racionales de la forma $k/p^n$ ($k\in\mathbf Z$, $n$ entero $\geq0$);
  evitaremos siempre utilizar la notación $\mathbf Z_p$ para designar este
  último anillo.}.
  \item[$2^{\mathrm o}$] Para el ideal primo $p\mathbf Z=(0)$,
  $\mathbf Z_{p\mathbf Z}$ es el cuerpo de los números racionales $\mathbf Q$,
  idéntico a su completado (siendo la topología, naturalmente, la topología
  de anillo local noetheriano, por tanto aquí la topología discreta).
\end{enumerate}

La terminología de (19.8.3), análoga a la de los «cuerpos primos», se
justifica de la misma manera: para todo anillo local $A$, consideremos
el homomorfismo canónico $\varphi:\mathbf Z\to A$, y sea $p\mathbf Z$
la imagen inversa por este homomorfismo del ideal maximal $\mathfrak m$ de
$A$; $p\mathbf Z$ es un ideal primo de $\mathbf Z$ y el homomorfismo
anterior se factoriza, por tanto, en
$\mathbf Z\to\mathbf Z_{p\mathbf Z}\xrightarrow{\psi}A$; además, como
$\varphi$ es el \emph{único} homomorfismo de $\mathbf Z$ en $A$, $p$ y
$\psi$ están determinados de manera \emph{única}. En otras palabras, para todo anillo
local $A$, hay un \emph{único homomorfismo} $\psi:P\to A$, donde $P$ es un
anillo local primo; si además $A$ es separado y completo, se puede prolongar
por completación este homomorfismo, y hay, por tanto, un \emph{único
homomorfismo} $\widehat\psi:\widehat P\to A$, donde $\widehat P$ es un anillo
local completo primo. Además, por paso a los cocientes, $\psi$ induce un
homomorfismo del cuerpo residual $\mathbf F_p$ si $p>0$ (resp. $\mathbf Q$ si
$p=0$) en el cuerpo residual $k$ de $A$, y $p$ es, por tanto, la
\emph{característica} de $k$.

Si se toma en particular como $A$ un anillo local primo (resp. completo
primo), se ve que en tal anillo no existe más que \emph{un solo
endomorfismo}, a saber, la identidad.

\oldpage[0\textsubscript{IV-1}]{111}

\begin{definition}[19.8.4]
\label{0.19.8.4-es}
Sean $A$ un anillo local, $P\to A$ el único homomorfismo de un anillo local
primo $P$ en $A$, y $p$ la característica de los cuerpos residuales de $P$ y
$A$. Se dice que $A$ es un \emph{anillo de Cohen} si es una $P$-álgebra de
Cohen, es decir (19.8.1), si:
\begin{enumerate}
  \item[$1^{\mathrm o}$] $A$ es noetheriano y completo.
  \item[$2^{\mathrm o}$] $A$ es un $P$-módulo plano (lo que equivale también a
  decir que $A$ es un $\widehat P$-módulo plano (Bourbaki, \emph{Alg. comm.},
  cap.~III, §~5, n\textsuperscript{o}~4, prop.~4)).
  \item[$3^{\mathrm o}$] $A/pA$ es un cuerpo (necesariamente separable sobre el
  cuerpo residual de $P$, pues este cuerpo es primo).
\end{enumerate}
\end{definition}

Si $p=0$, estas condiciones equivalen a decir que $A$ es un \emph{cuerpo de
característica $0$}. Si $p>0$, se tiene necesariamente $pA\neq0$; la condición
$3^{\mathrm o}$ significa que $pA$ es el ideal maximal $\mathfrak m$ de $A$;
la condición $2^{\mathrm o}$ significa que $p$ es $A$-regular, puesto que $P$
es un anillo de valoración discreta ($\mathbf 0_{\mathrm I}$, 6.3.4). Por tanto,
$A$ es un anillo regular (17.1.1, $d)$) de dimensión $1$ y, en consecuencia, un
anillo de valoración discreta, completo en virtud de $1^{\mathrm o}$; en suma:

\begin{proposition}[19.8.5]
\label{0.19.8.5-es}
Los anillos de Cohen son los cuerpos de característica $0$ y los anillos de
valoración discreta completos cuyo cuerpo residual tiene característica
$p>0$ y cuyo ideal maximal está generado por $p\cdot1$ (siendo $1$ la unidad del
anillo).
\end{proposition}

Obsérvese que, en el segundo caso, $p\cdot1\neq0$ porque $p$ es $A$-regular;
por tanto, se puede identificar $p\cdot1$ con el entero $p$; el homomorfismo
canónico $\mathbf Z_p\to A$ es inyectivo, y se identifica $p\cdot1$ con el
elemento $p$ de $\mathbf Z_p$; en este caso se dice que $A$ es un
\emph{$p$-anillo de Cohen}.

\begin{theorem}[19.8.6]
\label{0.19.8.6-es}
\textup{(Cohen).} ---
\begin{enumerate}
  \item[(i)] Sean $W$ un anillo de Cohen, $C$ un anillo local noetheriano
  completo y $\mathfrak J$ un ideal de $C$ distinto de $C$. Entonces todo
  homomorfismo local $u:W\to C/\mathfrak J$ se factoriza como
  $W\xrightarrow{v}C\to C/\mathfrak J$, donde $v$ es un homomorfismo local.
  \item[(ii)] Sea $K$ un cuerpo. Existe un anillo de Cohen $W$ cuyo cuerpo
  residual es isomorfo a $K$. Si $W'$ es un segundo anillo de Cohen y $K'$ su
  cuerpo residual, todo isomorfismo $u:K\xrightarrow{\sim}K'$ proviene, por
  paso a los cocientes, de un isomorfismo $v:W\xrightarrow{\sim}W'$.
\end{enumerate}
\end{theorem}

\begin{proof}
Esto no es sino (19.8.2) aplicado al caso en que $A$ es un anillo local primo.
\end{proof}

\begin{remarks}[19.8.7]
\label{0.19.8.7-es}
\begin{enumerate}
  \item[(i)] Cuando $K$ es de característica $0$, la parte (ii) de (19.8.6)
  se vuelve trivial.
  \item[(ii)] El homomorfismo $v$ de (19.8.6, (i)) \emph{no está
  necesariamente determinado de manera única por $u$}, como muestra ya el
  caso en que $W$ es un cuerpo de característica $0$, $\mathfrak J^2=0$ y $u$
  es un isomorfismo (cf.~(21.5.5)). Del mismo modo, en (19.8.6, (ii)), el
  isomorfismo $v$ \emph{no está necesariamente determinado de manera única
  por $u$} (cf.~(21.5.5)).

  Sin embargo, cuando $K$ es \emph{perfecto} y de característica $p>0$, se
  verá en (21.5.5) que en (19.8.6, (ii)) el isomorfismo $v$ es único. Se verá
  también más adelante que, en este caso, $W$ se identifica con el anillo
  $W_\infty(K)$ de \emph{vectores de Witt de longitud infinita} sobre $K$.
  \item[(iii)] En (19.8.6, (i)), se pueden debilitar las hipótesis sobre $C$
  utilizando (19.3.10) y (19.3.12).
\end{enumerate}
\end{remarks}

\begin{theorem}[19.8.8]
\label{0.19.8.8-es}
\textup{(Cohen).} --- Sean $A$ un anillo local noetheriano completo y $k$ su
cuerpo residual.

\oldpage[0\textsubscript{IV-1}]{112}

\begin{enumerate}
  \item[(i)] Existe un anillo de Cohen $W$ tal que $A$ es isomorfo a un anillo
  cociente de un anillo de series formales $W[[T_1,\ldots,T_n]]$ (y, en
  particular, $A$ es isomorfo a un cociente de un anillo local regular
  completo (17.3.8)). Si $A$ contiene un cuerpo, es isomorfo a un anillo
  cociente de $k[[T_1,\ldots,T_n]]$.
  \item[(ii)] Supongamos además que $A$ es íntegro. Entonces existe un
  subanillo $B$ de $A$ tal que: $1^{\mathrm o}$ $B$ es isomorfo a un anillo de
  series formales sobre un anillo $C$ que es un cuerpo o un anillo de Cohen
  (lo que implica que $B$ es un anillo local regular y completo (17.3.8));
  $2^{\mathrm o}$ $B$ tiene el mismo cuerpo residual que $A$ y la inyección
  $B\to A$ es un homomorfismo local; $3^{\mathrm o}$ $A$ es una $B$-álgebra
  finita.
\end{enumerate}
\end{theorem}

\begin{proof}
Sea $\mathfrak m$ el ideal maximal de $A$. Existe un anillo de Cohen $W$ cuyo
cuerpo residual es isomorfo a $k$ (19.8.6, (ii)); se tiene, pues, un
homomorfismo local $W\to A/\mathfrak m$, que se factoriza entonces como
$W\xrightarrow{u}A\to A/\mathfrak m$, donde $u$ es un homomorfismo local
(19.8.6, (i)). Para toda familia finita $(x_i)_{1\leq i\leq n}$ de elementos
de $\mathfrak m$, existe entonces un homomorfismo local
$v:W[[T_1,\ldots,T_n]]\to A$ que prolonga $u$ y tal que $v(T_i)=x_i$ para todo
$i$ (Bourbaki, \emph{Alg. comm.}, cap.~III, §~4, n\textsuperscript{o}~5,
prop.~6). Cuando $A$ contiene un cuerpo, contiene un cuerpo primo $P$, del cual
$k$ es una extensión (necesariamente separable) y, por tanto, $A$ contiene un
cuerpo isomorfo a $k$ (19.6.2); se puede entonces sustituir $W$ por $k$ en la
definición anterior de $v$.

\medskip
\noindent(i)\enspace
Tomemos primero como $x_i$ un sistema de generadores de $\mathfrak m$. Como
$W$ tiene el mismo cuerpo residual que $A$, y como las clases de los $x_i$ en
el anillo graduado $\gr_\bullet(A)$ generan $\gr_\bullet(A)$ como $k$-álgebra,
$\gr(v):\gr_\bullet(W[[T_1,\ldots,T_n]])\to\gr_\bullet(A)$ es sobreyectivo;
se deduce que $v$ es también sobreyectivo (Bourbaki, \emph{Alg. comm.},
cap.~III, §~2, n\textsuperscript{o}~8, cor.~2 del teorema~1). Recordemos que el
caso en que $A$ contiene un cuerpo ya se ha tratado y solo figura aquí como
recordatorio (19.6.3).

\medskip
\noindent(ii)\enspace
Si $A$ contiene un cuerpo, contiene un cuerpo $k'$ isomorfo a $k$, como se ha
visto; se considera entonces un sistema de parámetros
$(y_j)_{1\leq j\leq m}$ de $A$ (16.3.6), se toma
$B=k[[T_1,\ldots,T_m]]$ y se considera el homomorfismo local $w:B\to A$ cuya
restricción a $k$ es un isomorfismo $k\to k'$ y tal que
$w(T_j)=y_j$ para $1\leq j\leq m$. Si $A$ no contiene ningún cuerpo, el único
homomorfismo $\mathbf Z_{p\mathbf Z}\to A$ (19.8.3) es necesariamente
inyectivo (de lo contrario, como $A$ es íntegro, su núcleo sería el ideal
maximal de $\mathbf Z_{p\mathbf Z}$ y su imagen sería isomorfa a un cuerpo);
además, se tiene entonces $p>0$ por hipótesis, pues
$\mathbf Z_{p\mathbf Z}$ es un cuerpo si $p=0$. El elemento $p\cdot1$ de $A$
(identificado con $p$) no es divisor de cero en $A$ y pertenece a
$\mathfrak m$; por tanto, (16.3.4 y 16.3.7), existe una familia
$(z_i)_{1\leq i\leq m-1}$ que, junto con $p$, forma un sistema de parámetros
de $A$. El anillo de Cohen $W$ considerado al comienzo de la demostración es
entonces un anillo de valoración discreta de cuerpo residual $k$, en el que
$p$ genera el ideal maximal (19.8.5), y el homomorfismo $u:W\to A$
definido al comienzo envía $p$ a sí mismo. Se toma entonces
$B=W[[T_1,\ldots,T_{m-1}]]$ y se considera el homomorfismo local $w:B\to A$
que coincide con $u$ al restringirlo a $W$ y tal que $w(T_i)=z_i$ para
$1\leq i\leq m-1$. En ambos casos, si $\mathfrak n$ es el ideal maximal de
$B$, está claro que $\mathfrak nA$ es un ideal de definición de $A$; como,
además, $B/\mathfrak n$ y $A/\mathfrak m$ son isomorfos, $A$ es un
$B$-módulo cuasifinito ($\mathbf 0_{\mathrm I}$, 7.4.4), y, por tanto, un
$B$-módulo de tipo finito, pues $B$ es completo y $A$ es
separado para las topologías $\mathfrak n$-preádicas
($\mathbf 0_{\mathrm I}$, 7.4.1). Por otra
parte,

\oldpage[0\textsubscript{IV-1}]{113}

en ambos casos se tiene $\dim(B)=\dim(A)=m$; en el primer caso, esto resulta
de (17.1.4, (iii)); en el segundo, se ve directamente que $p$ y los $T_i$
($1\leq i\leq m-1$) forman una sucesión $B$-regular que genera $\mathfrak n$,
o también se puede utilizar el hecho de que estos elementos generan
$\mathfrak n$ y que se tiene $\dim(B)\geq\dim(A)$ por (16.3.10). Como $A$ y
$B$ son íntegros, se deduce finalmente de (16.3.10) que $w$ es inyectivo, lo
que termina la demostración.
\end{proof}

\begin{corollary}[19.8.9]
\label{0.19.8.9-es}
Sea $A$ un anillo local noetheriano íntegro completo que contiene un cuerpo
$k_0$; sea $k$ el cuerpo residual de $A$, y supongamos que $k$ es finito sobre
$k_0$. Entonces, en la conclusión de (19.8.8, (ii)), se pueden sustituir
$1^{\mathrm o}$ y $2^{\mathrm o}$ por la condición de que $B$ sea de la forma
$k_0[[T_1,\ldots,T_n]]$, siendo la inyección canónica $B\to A$ un
$k_0$-homomorfismo local (para la estructura usual de $k_0$-álgebra de $B$).
\end{corollary}

\begin{proof}
En efecto, retomando la demostración de (19.8.8, (ii)), se define esta vez
$w:k_0[[T_1,\ldots,T_n]]\to A$ de modo que su restricción a $k_0$ sea la
identidad y envíe $T_j$ a $y_j$ para $1\leq j\leq n$. La hipótesis de que $k$ es de
grado finito sobre $k_0$ implica de nuevo que $A$ es un $B$-módulo cuasifinito,
y, por tanto, de tipo finito por ($\mathbf 0_{\mathrm I}$, 7.4.1), y se
concluye como en (19.8.8).
\end{proof}

\begin{corollary}[19.8.10]
\label{0.19.8.10-es}
Sea $A$ un anillo local artiniano cuyo ideal maximal $\mathfrak m$ es de
cuadrado nulo; existe entonces un anillo local noetheriano regular $B$, de
ideal maximal $\mathfrak n$, tal que $A$ es isomorfo a
$B/\mathfrak n^2$.
\end{corollary}

\begin{proof}
Sean $K$ el cuerpo residual $A/\mathfrak m$ de $A$ y $n$ el rango de
$\mathfrak m/\mathfrak m^2=\mathfrak m$ sobre $K$. Si $A$ contiene un cuerpo,
se deduce de (19.6.3) que $A$ es isomorfo a $B/\mathfrak b$, donde
$B=K[[T_1,\ldots,T_n]]$ y $\mathfrak b$ está contenido en el cuadrado del
ideal maximal $\mathfrak n$ de $B$; pero, como
$\operatorname{long}(B/\mathfrak n^2)=n+1=\operatorname{long}(A)$, se tiene
necesariamente $\mathfrak b=\mathfrak n^2$.

Supongamos ahora que $A$ no contiene ningún cuerpo; esto implica que $K$ es
de característica $p>0$ y que $p\cdot1\neq0$ en $A$ (19.6.3); por tanto,
$p\cdot1$ es un elemento de $\mathfrak m$, y existen otros $n-1$ elementos
$x_i$ ($2\leq i\leq n$) de $\mathfrak m$ que, junto con $p\cdot1$, forman una
base de $\mathfrak m$ sobre $K$. Sea $W$ un anillo de Cohen cuyo cuerpo
residual es isomorfo a $K$; $W$ es un anillo de valoración discreta cuyo ideal
maximal está generado por $p$; se ha visto en la demostración de (19.8.8) que
existe un homomorfismo $u:W\to A$ que envía $p$ a sí mismo y que, por paso
a los cocientes, da la identidad sobre $K$. Se toma
$B=W[[T_2,\ldots,T_n]]$ y se considera el homomorfismo local $w:B\to A$ que
coincide con $u$ al restringirlo a $W$ y tal que $w(T_i)=x_i$ para $2\leq i\leq n$. Está
claro que $w$ es sobreyectivo y que su núcleo $\mathfrak b$ está contenido en
el cuadrado del ideal maximal $\mathfrak n=pB+BT_2+\cdots+BT_n$ de $B$; como
$\operatorname{long}(B/\mathfrak n^2)=n+1=\operatorname{long}(A)$, se tiene
de nuevo $\mathfrak b=\mathfrak n^2$.
\end{proof}

\begin{proposition}[19.8.11]
\label{0.19.8.11-es}
Sean $A$ un anillo local artiniano, $\mathfrak m$ su ideal maximal y $k$ su
cuerpo residual. Para que $A$ sea isomorfo a un anillo cociente de un anillo
de Cohen, es necesario y suficiente que $\mathfrak m$ esté generado por
$p\cdot1$, donde $p$ es la característica de $k$.
\end{proposition}

\begin{proof}
La condición es evidentemente necesaria (19.8.5). Para ver que es suficiente,
se observa, como al comienzo de la demostración de (19.8.8), que existe un
anillo de Cohen $W$ cuyo cuerpo residual es isomorfo a $k$ y un homomorfismo
local $u:W\to A$. Además, si se considera el homomorfismo compuesto
$\mathbf Z_{p\mathbf Z}\to W\xrightarrow{u}A$ (que es necesariamente el único
homomorfismo de $\mathbf Z_{p\mathbf Z}$ en $A$), se ve que la imagen

\oldpage[0\textsubscript{IV-1}]{114}

por $u$ del elemento $p\cdot1$ de $W$ es el elemento $p\cdot1$ de $A$; como
el elemento $p\cdot1$ de $W$ genera el ideal maximal de este anillo, se deduce
inmediatamente de la hipótesis que
$\gr(u):\gr_\bullet(W)\to\gr_\bullet(A)$ es sobreyectivo y, por tanto, también
lo es $u$ (Bourbaki, \emph{Alg. comm.}, cap.~III, §~2,
n\textsuperscript{o}~8, cor.~2 del teorema~1).
\end{proof}
\subsection{Álgebras relativamente formalmente lisas}
\label{subsection:0.19.9-es}

\begin{definition}[19.9.1]
\label{0.19.9.1-es}
Sean $\Lambda$ un anillo topológico, $A$ una $\Lambda$-álgebra topológica y
$B$ una $A$-álgebra topológica. Se dice que $B$ es una $A$-álgebra
\emph{formalmente lisa relativamente a $\Lambda$} si, para toda $A$-álgebra
topológica discreta $C$ y todo ideal nilpotente $\mathfrak J$ de $C$, todo
$A$-homomorfismo continuo $u_0:B\to C/\mathfrak J$ que se factoriza como
$B\xrightarrow{u}C\to C/\mathfrak J$, donde $u$ es un $\Lambda$-homomorfismo
continuo, se factoriza también como $B\xrightarrow{v}C\to C/\mathfrak J$, donde
$v$ es un $A$-homomorfismo continuo.
\end{definition}

De esta definición resulta que, si $B$ es una $A$-álgebra \emph{formalmente
lisa}, entonces $B$ es también \emph{formalmente lisa relativamente a
$\Lambda$}, para toda estructura de $\Lambda$-álgebra topológica definida sobre
$A$ (es decir, todo homomorfismo continuo de anillos $\Lambda\to A$).

\begin{proposition}[19.9.2]
\label{0.19.9.2-es}
Sean $\Lambda$ un anillo topológico y $A$ una $\Lambda$-álgebra topológica.
\begin{enumerate}
  \item[(i)] $A$ es una $\Lambda$-álgebra formalmente lisa relativamente a
  $\Lambda$.
  \item[(ii)] Si $B$ es una $A$-álgebra formalmente lisa relativamente a
  $\Lambda$ y $C$ una $B$-álgebra formalmente lisa relativamente a $\Lambda$,
  entonces $C$ es una $A$-álgebra formalmente lisa relativamente a $\Lambda$.
  \item[(iii)] Sean $B$ una $A$-álgebra formalmente lisa relativamente a
  $\Lambda$ y $A'$ una $A$-álgebra topológica; entonces la $A'$-álgebra
  topológica $B\otimes_A A'$ es formalmente lisa relativamente a $\Lambda$.
  \item[(iv)] Sean $B$ una $A$-álgebra topológica y $S$ (resp. $T$) un
  subconjunto multiplicativo de $A$ (resp. $B$) tal que la imagen canónica de
  $S$ en $B$ esté contenida en $T$. Si $B$ es una $A$-álgebra formalmente lisa
  relativamente a $\Lambda$, entonces $T^{-1}B$ es una $S^{-1}A$-álgebra
  formalmente lisa relativamente a $\Lambda$.
  \item[(v)] Sean $B_i$ ($1\leq i\leq n$) unas $A$-álgebras topológicas. Para
  que $\prod_{i=1}^n B_i$ sea una $A$-álgebra formalmente lisa relativamente a
  $\Lambda$, es necesario y suficiente que cada una de las $B_i$ lo sea.
\end{enumerate}
\end{proposition}

\begin{proof}
La afirmación (i) es trivial, y la demostración de las demás sigue muy de cerca
las demostraciones de (19.3.5); se deja, pues, al lector.
\end{proof}

\begin{corollary}[19.9.3]
\label{0.19.9.3-es}
Sean $A$ un anillo topológico, y $A'$ y $B$ dos $A$-álgebras topológicas.
Entonces la $A$-álgebra topológica $B'=B\otimes_A A'$ es una $B$-álgebra
formalmente lisa relativamente a $A$.
\end{corollary}

\begin{proof}
Esto resulta de (19.9.2, (i) y (iii)).
\end{proof}

\begin{proposition}[19.9.4]
\label{0.19.9.4-es}
Sean $\Lambda$ un anillo topológico, $A$ una $\Lambda$-álgebra topológica y
$B$ una $A$-álgebra topológica. Las condiciones siguientes son equivalentes:
\begin{enumerate}
  \item[a)] $B$ es una $A$-álgebra formalmente lisa relativamente a $\Lambda$.
  \item[b)] $\widehat B$ es una $A$-álgebra formalmente lisa relativamente a
  $\Lambda$.
  \item[c)] $\widehat B$ es una $\widehat A$-álgebra formalmente lisa
  relativamente a $\Lambda$.
  \item[d)] $\widehat B$ es una $\widehat A$-álgebra formalmente lisa
  relativamente a $\widehat\Lambda$.
\end{enumerate}
\end{proposition}

\begin{proof}
Se deja de nuevo al lector la demostración, calcada de la de (19.3.6).
\end{proof}

\oldpage[0\textsubscript{IV-1}]{115}

\begin{env}[19.9.5]
\label{0.19.9.5-es}
Del mismo modo, el enunciado (19.3.8) sigue siendo válido (con la misma
demostración) cuando se sustituye «formalmente lisa» por «formalmente lisa
relativamente a $\Lambda$». Si en el enunciado de (19.3.10) se sustituye
«formalmente lisa» por «formalmente lisa relativamente a $\Lambda$», la
conclusión se sustituye por la siguiente (permaneciendo la demostración
esencialmente inalterada): \emph{todo $A$-homomorfismo continuo
$u_0:B\to C/\mathfrak J$ que se factoriza como
$B\xrightarrow{u}C\to C/\mathfrak J$, donde $u$ es un $\Lambda$-homomorfismo
continuo, se factoriza también como $B\xrightarrow{v}C\to C/\mathfrak J$, donde
$v$ es un $A$-homomorfismo continuo}.
\end{env}

\begin{env}[19.9.6]
\label{0.19.9.6-es}
Los criterios de lisura formal (19.4.1) y (19.4.2) son válidos cuando se
sustituye en ellos «formalmente lisa» por «formalmente lisa relativamente a
$\Lambda$», permaneciendo las demostraciones prácticamente inalteradas.
\end{env}

\begin{proposition}[19.9.7]
\label{0.19.9.7-es}
Sean $\Lambda$ un anillo topológico, $A$ una $\Lambda$-álgebra topológica y
$B$ una $A$-álgebra topológica. Supongamos que, para toda $A$-álgebra discreta
$C$ y todo ideal $\mathfrak J$ de $C$ tal que $\mathfrak J^2=0$, todo
$A$-homomorfismo continuo $u_0:B\to C/\mathfrak J$ que se factoriza como
$B\xrightarrow{u}C\to C/\mathfrak J$, donde $u$ es un $\Lambda$-homomorfismo
continuo, se factoriza también como $B\xrightarrow{v}C\to C/\mathfrak J$, donde
$v$ es un $A$-homomorfismo continuo. Entonces $B$ es una $A$-álgebra
formalmente lisa relativamente a $\Lambda$.
\end{proposition}

\begin{proof}
La demostración de (19.4.3) se transcribe inmediatamente.
\end{proof}

\begin{proposition}[19.9.8]
\label{0.19.9.8-es}
Sean $\Lambda$ un anillo topológico, $A$ una $\Lambda$-álgebra topológica y
$B$ una $A$-álgebra topológica. Para que $B$ sea una $A$-álgebra formalmente
lisa relativamente a $\Lambda$, es necesario y suficiente que, para todo
$B$-módulo topológico discreto $L$, anulado por un ideal abierto de $B$, se
tenga (cf.~18.4.2)
\[
  \operatorname{Exalcotop}_{A/\Lambda}(B,L)=0.
\]
\end{proposition}

\begin{proof}
Con las notaciones de la demostración de (19.4.4), basta observar aquí que se
puede suponer que la extensión $E_\Lambda$ de $B_\Lambda$ es
$\Lambda$-trivial; el resto de la demostración permanece entonces inalterado.
\end{proof}

Cuando $\Lambda$, $A$ y $B$ son anillos discretos, el criterio (19.9.8) se
reduce a
\begin{equation}
\tag{19.9.8.1}
\label{0.19.9.8.1-es}
  \operatorname{Exalcom}_{A/\Lambda}(B,L)=0
  \quad\text{para todo $B$-módulo $L$}~;
\end{equation}
es decir, \emph{toda $A$-extensión conmutativa de $B$ por un $B$-módulo, que es
$\Lambda$-trivial, es también $A$-trivial}.

\subsection{Álgebras formalmente no ramificadas y álgebras formalmente
étales}
\label{subsection:0.19.10-es}
\label{0.19.10-es}

\begin{env}[19.10.1]
\label{0.19.10.1-es}
También tendremos que utilizar dos nociones próximas a la noción de álgebra
formalmente lisa.
\end{env}

\begin{definition}[19.10.2]
\label{0.19.10.2-es}
Sean $A$ un anillo topológico y $B$ una $A$-álgebra topológica. Se dice que
$B$ es una $A$-álgebra \emph{formalmente no ramificada} si, para todo
$A$-homomorfismo $p:E\to C$ de $A$-álgebras discretas cuyo núcleo es
nilpotente, y todo $A$-homomorfismo continuo $u:B\to C$, existe \emph{a lo sumo
un} $A$-homomorfismo continuo $v:B\to E$ tal que $u$ se factoriza como
$B\xrightarrow{v}E\xrightarrow{p}C$. Se dice que $B$ es una $A$-álgebra
\emph{formalmente étale} si es formalmente lisa y formalmente no ramificada, es
decir, si, para todo $A$-homomorfismo sobreyectivo $p:E\to C$ de $A$-álgebras
discretas cuyo núcleo es nilpotente, y para todo $A$-homomorfismo continuo
$u:B\to C$, existe \emph{uno y solo un} $A$-homomorfismo continuo $v:B\to E$
tal que $u$ se factoriza como $B\xrightarrow{v}E\xrightarrow{p}C$.
\end{definition}

\oldpage[0\textsubscript{IV-1}]{116}

\begin{examples}[19.10.3]
\label{0.19.10.3-es}
Nos limitaremos a los anillos \emph{discretos}.
\begin{enumerate}
  \item[(i)] Si el homomorfismo estructural $\rho:A\to B$ es
  \emph{sobreyectivo}, $B$ es una $A$-álgebra \emph{formalmente no ramificada},
  pues, con las notaciones de (19.10.2), el homomorfismo compuesto
  $A\xrightarrow{\rho}B\xrightarrow{v}E$ debe ser el homomorfismo estructural,
  por lo que $v$ es único (si existe). En cambio, $B$ es formalmente lisa si y
  solo si $\ker(\rho)=\ker(\rho)^2$.
  \item[(ii)] Sean $A$ un anillo y $S$ un subconjunto multiplicativo de $A$;
  entonces $B=S^{-1}A$ es una $A$-álgebra \emph{formalmente étale}. En efecto,
  sean $j:A\to B$ el homomorfismo canónico, $\rho:A\to E$ un homomorfismo de
  anillos, $\mathfrak J$ un ideal nilpotente de $E$ y
  $p:E\to C=E/\mathfrak J$ el homomorfismo canónico. Si $u:B\to C$ es un
  $A$-homomorfismo, $u(s/1)$ es invertible en $C$ para $s\in S$, puesto que
  $s/1$ es invertible en $B$; pero $u(s/1)=p(\rho(s))$, y, como $\mathfrak J$
  es nilpotente, el hecho de que la imagen de $\rho(s)$ en $E/\mathfrak J$ sea
  invertible implica que el propio $\rho(s)$ es invertible en $E$
  ($\mathbf 0_{\mathrm I}$, 7.1.12); se concluye que $\rho$ se factoriza de una
  sola manera como $A\xrightarrow{j}B\xrightarrow{v}E$, y se tiene
  evidentemente $u\circ j=(p\circ v)\circ j$, de donde $u=p\circ v$, lo que
  prueba nuestra afirmación.
  \item[(iii)] Una extensión finita separable $k'$ de un cuerpo $k$ es una
  $k$-álgebra formalmente lisa (19.6.1); veremos más adelante que, para que una
  extensión de \emph{tipo finito} $K$ de un cuerpo $k$ sea una $k$-álgebra
  formalmente étale, es necesario y suficiente que sea formalmente no
  ramificada, y esto equivale también a que $K$ sea una extensión finita y
  separable de $k$ (21.7.4).
  \item[(iv)] Un álgebra de polinomios $A[X_\alpha]_{\alpha\in I}$ sobre un
  anillo $A$ es \emph{formalmente lisa} sobre $A$ (19.3.3), pero \emph{no
  formalmente étale} si $I$ no es vacío, como resulta inmediatamente de las
  definiciones.
\end{enumerate}
\end{examples}

\begin{remark}[19.10.4]
\label{0.19.10.4-es}
El mismo razonamiento que en (19.4.3) muestra que, para verificar que una
$A$-álgebra topológica $B$ es formalmente no ramificada, se puede, en la
definición (19.10.2), limitar al caso en que el núcleo de $p:E\to C$ es de
\emph{cuadrado nulo}; se puede además limitar al caso en que $u(B)\subset p(E)$
(pues de lo contrario no existe ningún $v$ que factorice $u$) y, sustituyendo
$C$ por $u(B)$ y $E$ por $p^{-1}(u(B))$, suponer $u$ y $p$ sobreyectivos.
\end{remark}

\section{Derivaciones y diferenciales}
\label{section:0.20-es}
\label{0.20-es}

Las nociones introducidas en este parágrafo se retomarán en forma
geométrica en el capítulo~IV, §~16, y desempeñarán un papel importante en el
estudio de los preesquemas. Su importancia en este capítulo se debe, en primer
lugar, a sus relaciones con la noción de álgebra formalmente lisa y, en
particular, a los teoremas (20.4.9), (20.5.7) y (20.5.12), que se traducirán al
lenguaje geométrico en el parágrafo del capítulo~IV dedicado a los morfismos
lisos y son de uso constante en las aplicaciones. Por otra parte, las nociones
diferenciales servirán para demostrar en el §~22 importantes criterios de
regularidad que desempeñarán un papel esencial en el estudio profundo de los
anillos locales noetherianos realizado en el §~7 del capítulo~IV.

\subsection{Derivaciones y extensiones de álgebras}
\label{subsection:0.20.1-es}

\oldpage[0\textsubscript{IV-1}]{117}
\begin{proposition}[20.1.1]
\label{0.20.1.1-es}
Sean $A$ un anillo (no necesariamente conmutativo), $B$, $E$, $C$
$A$-anillos, $p:E\to C$ un $A$-homomorfismo cuyo núcleo $\mathfrak{Q}$
es de cuadrado nulo, $u:B\to C$ un $A$-homomorfismo. Supongamos que existe
un $A$-homomorfismo $v_0:B\to E$ tal que $u$ se factoriza como
$B\xrightarrow{v_0}E\xrightarrow{p}C$. Entonces el conjunto de los
$A$-homomorfismos $v:B\to E$ tales que $u=p\circ v$ es idéntico al
conjunto de las aplicaciones $v_0+D$, donde $D:B\to\mathfrak{Q}$ es una
aplicación que verifica las dos condiciones siguientes:
\begin{enumerate}
  \item[\rm(i)] $D$ es un homomorfismo de $A$-bimódulos.
  \item[\rm(ii)] Cualesquiera que sean $f$, $g$ en $B$, se tiene
\end{enumerate}
\begin{equation}
\tag{20.1.1.1}
\label{0.20.1.1.1-es}
  D(fg)=f.D(g)+D(f).g.
\end{equation}
\end{proposition}

\begin{proof}
Decir que $p(v(f))=p(v_0(f))$ para $f\in B$ significa que
$D(f)=v(f)-v_0(f)$ pertenece a $\mathfrak{Q}$; al escribir que
$v(fg)=v(f)v(g)$, se obtiene la relación \sref{0.20.1.1.1-es},
puesto que $\mathfrak{Q}$ es de cuadrado nulo, y la condición (i) resulta de que
$v$ y $v_0$ son $A$-homomorfismos.

Si $\rho:A\to B$ es el homomorfismo estructural, se deduce de
\sref{0.20.1.1.1-es} que se tiene
$D(\rho(a)f)=\rho(a)D(f)+D(\rho(a))f$ para todo $a\in A$; pero se debe
tener $D(\rho(a)f)=\rho(a)D(f)$ por (i), luego, tomando $f=1$, se
obtiene
\[
  D(\rho(a))=0 \qquad\text{para}\quad a\in A\,;
\]
recíprocamente, si $D$ es nula en $\rho(A)$ y verifica
\sref{0.20.1.1.1-es}, también verifica (i).
\end{proof}

\begin{definition}[20.1.2]
\label{0.20.1.2-es}
Dados un $A$-anillo $B$ y un $B$-bimódulo $L$, se llama
\emph{$A$-derivación de $B$ en $L$} a una aplicación $D:B\to L$ que verifica
las condiciones (i) y (ii) de \sref{0.20.1.1-es}.
\end{definition}

Se deduce de \sref{0.20.1.1-es} que el \emph{núcleo} de una $A$-derivación
$D:B\to L$ es un sub-$A$-anillo de $B$.

A veces se llama \emph{derivación} de $B$ en $L$ a una
$\mathbf{Z}$-derivación, es decir, una aplicación \emph{aditiva} de $B$
en $L$ que verifica \sref{0.20.1.1.1-es}; se tiene, por tanto, $D(1)=0$ para tal
aplicación. Si $B$ es un álgebra sobre un \emph{cuerpo primo} $P$, toda
$\mathbf{Z}$-derivación de $B$ es una $P$-derivación: esto es evidente por lo
anterior si $P$ es de característica $>0$ y, en caso contrario,
la relación $n.n^{-1}=1$ para $n\in\mathbf{Z}$ da
\[
  nD(n^{-1})=0
\]
luego $D(n^{-1})=0$, puesto que $P$ es de característica $0$.

Se deduce inmediatamente de la definición \sref{0.20.1.2-es} que, si $D$, $D'$
son dos $A$-derivaciones de $B$ en $L$, lo mismo ocurre con $D-D'$. En otras
palabras, el conjunto de las $A$-derivaciones de $B$ en $L$ está provisto de una
estructura de \emph{grupo aditivo}; este grupo se denota por
$\operatorname{D\acute{e}r}_A(B,L)$. Si $A$ es \emph{conmutativo}, $B$ un
$A$-\emph{álgebra conmutativa}, y $L$ un $B$-módulo, entonces, para toda
$A$-derivación $D$ de $B$ en $L$ y todo $a\in A$, $aD$ sigue siendo una
$A$-derivación de $B$ en $L$; en otras palabras, $\operatorname{D\acute{e}r}_A(B,L)$
está entonces provisto de una estructura de $A$-\emph{módulo}.

\oldpage[0\textsubscript{IV-1}]{118}
La prop. \sref{0.20.1.1-es} se interpreta de la manera siguiente:

\begin{corollary}[20.1.3]
\label{0.20.1.3-es}
Dados dos $A$-homomorfismos de $A$-anillos $p:E\to C$,
$u:B\to C$, el primero de los cuales tiene un núcleo $\mathfrak{Q}$ de cuadrado
nulo, el conjunto de los $A$-homomorfismos $v:B\to E$ tales que $u=p\circ v$
es vacío o es un espacio homogéneo principal bajo el grupo
$\operatorname{D\acute{e}r}_A(B,\mathfrak{Q})$.
\end{corollary}

En particular:

\begin{corollary}[20.1.4]
\label{0.20.1.4-es}
Sean $A$ un anillo, $C$ un $A$-anillo, $L$ un $C$-bimódulo, $E$ una
$A$-extensión de $C$ por $L$, $p:E\to C$ la aumentación. La aplicación que
asocia a toda derivación $D\in\operatorname{D\acute{e}r}_A(C,L)$ la aplicación
$x\mapsto x+D(p(x))$ es un isomorfismo del grupo
$\operatorname{D\acute{e}r}_A(C,L)$ sobre el grupo de las $A$-equivalencias de $E$
sobre sí misma.
\end{corollary}

\begin{proof}
Apliquemos \sref{0.20.1.1-es} para $B=E$, $v_0=\mathrm{I}_E$: el conjunto
de los $A$-homomorfismos $v:E\to E$ tales que $p\circ v=p$ es idéntico al
conjunto de las aplicaciones $v=\mathrm{I}_E+D'$, donde
$D'\in\operatorname{D\acute{e}r}_A(E,L)$. Decir que tal $A$-homomorfismo $v$ es
una $A$-equivalencia equivale a decir que la restricción de $v$ a $L$ es la
inyección canónica, o también que $D'(x)=0$ en $L$; pero esto significa asimismo
que $D'$ se factoriza como $E\xrightarrow{p}C\xrightarrow{D}L$, donde $D$ es una
$A$-derivación. De donde resulta el corolario.
\end{proof}

Para las extensiones triviales, \sref{0.20.1.3-es} da:

\begin{corollary}[20.1.5]
\label{0.20.1.5-es}
Sean $A$ un anillo, $B$, $C$ dos $A$-anillos, $L$ un $C$-bimódulo,
$u:B\to C$ un $A$-homomorfismo; la aplicación que asocia a toda derivación
$D\in\operatorname{D\acute{e}r}_A(B,L)$ la aplicación
$\alpha:x\mapsto(u(x),D(x))$ es una biyección sobre el conjunto de los
$A$-homomorfismos $v:B\to D_C(L)$ (cf. 18.2.3) tales que $u$ se factoriza como
$B\xrightarrow{v}D_C(L)\to C$.
\end{corollary}

Más particularmente:

\begin{corollary}[20.1.6]
\label{0.20.1.6-es}
Sean $A$ un anillo, $B$ un $A$-anillo, $L$ un $B$-bimódulo. Si a toda
derivación $D\in\operatorname{D\acute{e}r}_A(B,L)$ se asocian: $1^\circ$ la
$A$-equivalencia $(x,y)\mapsto(x,y+D(x))$ de la extensión $D_B(L)$ sobre
sí misma; $2^\circ$ el $A$-homomorfismo $x\mapsto(x,D(x))$ de $B$ en
$D_B(L)$, inverso por la derecha del homomorfismo de aumentación $D_B(L)\to B$,
se definen correspondencias biunívocas canónicas entre:
\begin{enumerate}
  \item[\rm(i)] el conjunto $\operatorname{D\acute{e}r}_A(B,L)$;
  \item[\rm(ii)] el conjunto de las $A$-equivalencias de $D_B(L)$ sobre sí misma;
  \item[\rm(iii)] el conjunto de los $A$-homomorfismos inversos por la derecha del
  homomorfismo de aumentación $D_B(L)\to B$.
\end{enumerate}
\end{corollary}

Se observará que la correspondencia biunívoca así establecida entre (i) y (ii)
respeta las estructuras de grupo, y que la que se deduce de ella entre (ii) y
(iii) no es sino la correspondencia biunívoca ya definida en (18.3.8).

\begin{corollary}[20.1.7]
\label{0.20.1.7-es}
Sean $A$ un anillo, $C$ un $A$-anillo, $L$ un $A$-bimódulo, $E$ una
$A$-extensión trivial de $C$ por $L$, $p:E\to C$ la aumentación. El conjunto
de las $A$-derivaciones $D$ de $E$ en $L$ tales que $D(x)=x$ en $L$ es
idéntico al conjunto de las aplicaciones $x\mapsto x-w(p(x))$, donde $w$
recorre el conjunto de los $A$-homomorfismos inversos por la derecha de $p$.
\end{corollary}

\begin{proof}
Basta aplicar de nuevo \sref{0.20.1.1-es} para $B=E$,
$v_0=\mathrm{I}_E$; si $v=\mathrm{I}_E-D$, la condición $D(x)=x$ para
$x\in L$ equivale a $v(x)=0$ para $x\in L$, es decir, a $v=w\circ p$,
donde $w:C\to E$ es un $A$-homomorfismo; además, la condición
$p\circ v=p$ equivale a $p\circ w=\mathrm{I}_C$, en otras palabras, al hecho de
que $w$ sea inverso por la derecha de $p$.
\end{proof}

\subsection{Propiedades funtoriales de las derivaciones}
\label{subsection:0.20.2-es}

\oldpage[0\textsubscript{IV-1}]{119}
\begin{env}[20.2.1]
\label{0.20.2.1-es}
Sean $A$ un anillo, $B$ un $A$-anillo, $L$ un $B$-bimódulo; si $L'$ es
un segundo $B$-bimódulo y $w:L\to L'$ un homomorfismo de $B$-bimódulos, es
claro que la aplicación $D\mapsto w\circ D$ es un homomorfismo de
grupos aditivos
\begin{equation}
\tag{20.2.1.1}
\label{0.20.2.1.1-es}
  w_0:\operatorname{D\acute{e}r}_A(B,L)\to\operatorname{D\acute{e}r}_A(B,L')
\end{equation}
y que, si $w':L'\to L''$ es un segundo homomorfismo de $B$-bimódulos, se tiene
$(w'\circ w)_0=w'_0\circ w_0$. Cuando $A$ es conmutativo, $B$ un
$A$-álgebra conmutativa y $L$ un $B$-módulo, \sref{0.20.2.1.1-es} es un
homomorfismo de $A$-módulos.

En segundo lugar, sean $B'$ un $A$-anillo, $v:B'\to B$ un
$A$-homomorfismo que hace de $L$ un $B'$-bimódulo; entonces la aplicación
$D\mapsto D\circ v$ es un homomorfismo de grupos aditivos,
\begin{equation}
\tag{20.2.1.2}
\label{0.20.2.1.2-es}
  v^0:\operatorname{D\acute{e}r}_A(B,L)\to\operatorname{D\acute{e}r}_A(B',L)
\end{equation}
como se deduce de \sref{0.20.1.1.1-es}; si $v':B''\to B'$ es un segundo
$A$-homomorfismo, se tiene $(v\circ v')^0=v'^0\circ v^0$. Cuando $A$, $B$ y
$B'$ son conmutativos y $L$ un $B$-módulo, \sref{0.20.2.1.2-es} es un
homomorfismo de $A$-módulos.

Finalmente, sea $u:A'\to A$ un homomorfismo de anillos que hace de $B$ un
$A'$-anillo; toda $A$-derivación $D\in\operatorname{D\acute{e}r}_A(B,L)$ es también
una $A'$-derivación, de donde resulta una inyección canónica de grupos conmutativos
\begin{equation}
\tag{20.2.1.3}
\label{0.20.2.1.3-es}
  u^0:\operatorname{D\acute{e}r}_A(B,L)\to\operatorname{D\acute{e}r}_{A'}(B,L)
\end{equation}
y, si $u':A''\to A'$ es un segundo homomorfismo de anillos, se tiene
$(u\circ u')^0=u'^0\circ u^0$; cuando $A$, $A'$ y $B$ son conmutativos y
$L$ un $B$-módulo, \sref{0.20.2.1.3-es} es un dihomomorfismo de módulos
(relativo a $u$).

También se puede decir que
\[
  (A,B,L)\leadsto\operatorname{D\acute{e}r}_A(B,L)
\]
es un \emph{funtor covariante} de la categoría $\boldsymbol{K}$ definida en
(18.3.5) en la categoría $\boldsymbol{Ab}$ de los grupos conmutativos, haciendo
corresponder a toda terna $(u,v,w)$ que constituye un morfismo de
$\boldsymbol{K}$ el homomorfismo $w_0\circ v^0\circ u^0$; la verificación de
la funtorialidad resulta de la conmutatividad de los diagramas
\[
\xymatrix{
  \operatorname{D\acute{e}r}_A(B,L) \ar[r]^{v^0} \ar[d]_{w_0} &
    \operatorname{D\acute{e}r}_A(B',L) \ar[d]^{w_0} \\
  \operatorname{D\acute{e}r}_A(B,L') \ar[r]^{v^0} &
    \operatorname{D\acute{e}r}_A(B',L')
}
\qquad
\xymatrix{
  \operatorname{D\acute{e}r}_A(B,L) \ar[r]^{u^0} \ar[d]_{w_0} &
    \operatorname{D\acute{e}r}_{A'}(B,L) \ar[d]^{w_0} \\
  \operatorname{D\acute{e}r}_A(B,L') \ar[r]^{u^0} &
    \operatorname{D\acute{e}r}_{A'}(B,L')
}
\]
para todo homomorfismo $w:L\to L'$ de $B$-bimódulos.
\end{env}

\oldpage[0\textsubscript{IV-1}]{120}
\begin{theorem}[20.2.2]
\label{0.20.2.2-es}
Sean $u:A\to B$, $v:B\to C$ dos homomorfismos de anillos, $L$ un
$C$-bimódulo. Se tiene una sucesión exacta canónica de grupos conmutativos
\begin{align}
\tag{20.2.2.1}
\label{0.20.2.2.1-es}
  0\to\operatorname{D\acute{e}r}_B(C,L)\xrightarrow{u^0}
  \operatorname{D\acute{e}r}_A(C,L)&\xrightarrow{v^0}
  \operatorname{D\acute{e}r}_A(B,L)\xrightarrow{\partial} \\
  &\to\operatorname{Exan}_B(C,L)\xrightarrow{u^1}
  \operatorname{Exan}_A(C,L)\xrightarrow{v^1}
  \operatorname{Exan}_A(B,L) \notag
\end{align}
donde $u^0$, $v^0$ son los homomorfismos (20.2.1.3) y (20.2.1.2)
respectivamente, $u^1$, $v^1$ los homomorfismos definidos en (18.3.4.1) y
(18.3.3.1) respectivamente, y donde $\partial$ se define como sigue: para
toda $A$-derivación $D$ de $B$ en $L$, $\partial(D)$ es la clase de la
$B$-extensión de $C$ por $L$ definida sobre el anillo $D_C(L)$ por el
$A$-homomorfismo $\alpha:x\mapsto(v(x),D(x))$ (cf.
\sref{0.20.1.5-es}). Además, la sucesión exacta \sref{0.20.2.2.1-es} es
funtorial en $L$ (para los homomorfismos definidos en (20.2.1.1) y
(18.3.1.1) respectivamente).
\end{theorem}

\begin{proof}
Como $D$ es una $A$-derivación (y \emph{a fortiori} una
$\mathbf{Z}$-derivación) de $B$ en $L$, el $A$-homomorfismo
$\alpha:x\mapsto(v(x),D(x))$ define efectivamente sobre $D_C(L)$ una estructura de
$B$-extensión, por tanto $\partial(D)$ está bien definida \sref{0.20.1.5-es}. Hay
que verificar la exactitud en $5$ lugares:

\medskip
\noindent 1)\enspace
La exactitud en $\operatorname{D\acute{e}r}_B(C,L)$ es trivial (cf.
\sref{0.20.2.1-es}).

\medskip
\noindent 2)\enspace
Por definición \sref{0.20.2.1-es} el núcleo de $v^0$ es el conjunto de las
$A$-derivaciones de $C$ en $L$ que se anulan en $v(B)$, es decir,
las $A$-derivaciones que son también $B$-derivaciones
\sref{0.20.1.1-es}; de donde resulta la exactitud en $\operatorname{D\acute{e}r}_A(C,L)$.

\medskip
\noindent 3)\enspace
El núcleo de $\partial$ está formado por las derivaciones
$D\in\operatorname{D\acute{e}r}_A(B,L)$ para las cuales la $B$-extensión definida por
$\alpha:x\mapsto(v(x),D(x))$ es $B$-trivial; esto significa (18.2.3) que
existe un $B$-homomorfismo $z\mapsto(z,w(z))$ de $C$ en $D_C(L)$ (estando la
estructura de $B$-anillo de $D_C(L)$ definida por $\alpha$); pero un tal
homomorfismo, al ser \emph{a fortiori} un $A$-homomorfismo, es de la forma
$z\mapsto(z,D'(z))$ donde $D'\in\operatorname{D\acute{e}r}_A(C,L)$
\sref{0.20.1.6-es}; y al escribir que es un $B$-homomorfismo, se obtiene
\[
  D(x)z+v(x)D'(z)=D'(v(x)z)=v(x)D'(z)+D'(v(x))z
\]
para $x\in B$, $z\in C$, lo que da $D'(v(x))=D(x)$; el núcleo de
$\partial$ es por tanto la imagen de $v^0$.

\medskip
\noindent 4)\enspace
El núcleo de $u^1$ es el conjunto de las clases de $B$-extensiones de $C$ por
$L$ que son $A$-triviales (18.3.7), por tanto (salvo equivalencia) de la forma
$D_C(L)$, donde la estructura de $A$-anillo está definida por el homomorfismo
$t\mapsto(u(t),0)$. Ahora bien, toda estructura de $B$-extensión sobre $D_C(L)$ está
definida por un homomorfismo $\alpha:x\mapsto(v(x),D(x))$, donde $D$ es una
$\mathbf{Z}$-derivación de $B$ en $L$ \sref{0.20.1.5-es}; decir que la
estructura de $A$-anillo de esta $B$-extensión se deduce de su estructura de
$B$-anillo por medio de $u:A\to B$ significa que $D(u(t))=0$ para $t\in A$,
por tanto que $D$ es una $A$-derivación o también que la clase de la $B$-extensión
considerada es de la forma $\partial(D)$; de donde resulta la exactitud en
$\operatorname{Exan}_B(C,L)$.

\medskip
\noindent 5)\enspace
El núcleo de $v^1$ es el conjunto de las clases de las $A$-extensiones $E$ de $C$
por $L$ que se trivializan sobre $v(B)$, es decir, para las cuales existe un
$A$-homomorfismo $w:B\to E$ tal que $v$ se factoriza como
$B\xrightarrow{w}E\to C$; ahora bien, un tal $A$-homomorfismo define sobre $E$ una
estructura de $B$-extensión cuya clase tiene por imagen mediante $u^1$ la clase de
la $A$-extensión dada; siendo trivial el recíproco, se ha probado
la exactitud en $\operatorname{Exan}_A(C,L)$.

Finalmente, la funtorialidad en $L$ se deduce trivialmente de las definiciones.
\end{proof}

\oldpage[0\textsubscript{IV-1}]{121}
\begin{corollary}[20.2.3]
\label{0.20.2.3-es}
Sean $A$, $B$, $C$ tres anillos conmutativos, $u:A\to B$, $v:B\to C$
dos homomorfismos de anillos, $L$ un $C$-módulo. Se tiene una sucesión exacta
canónica de $A$-módulos
\begin{align}
\tag{20.2.3.1}
\label{0.20.2.3.1-es}
  0\to\operatorname{D\acute{e}r}_B(C,L)\xrightarrow{u^0}
  \operatorname{D\acute{e}r}_A(C,L)&\xrightarrow{v^0}
  \operatorname{D\acute{e}r}_A(B,L)\xrightarrow{\partial} \\
  &\to\operatorname{Exalcom}_B(C,L)\xrightarrow{u^1}
  \operatorname{Exalcom}_A(C,L)\xrightarrow{v^1}
  \operatorname{Exalcom}_A(B,L) \notag
\end{align}
funtorial en $L$.
\end{corollary}

\begin{proof}
El razonamiento es el mismo que en \sref{0.20.2.2-es} una vez que se ha
verificado que para toda derivación $D\in\operatorname{D\acute{e}r}_A(B,L)$,
$\partial(D)$ es efectivamente la clase de una $B$-extensión de $C$ por $L$ que es
una $B$-álgebra conmutativa; pero esto resulta inmediatamente de la conmutatividad
de $C$ y del hecho de que $L$ es un $C$-módulo.
\end{proof}

\begin{corollary}[20.2.4]
\label{0.20.2.4-es}
Bajo las hipótesis de \sref{0.20.2.2-es} (resp.
\sref{0.20.2.3-es}), se tiene una sucesión exacta canónica, funtorial en $L$,
\begin{equation}
\tag{20.2.4.1}
\label{0.20.2.4.1-es}
  0\to\operatorname{D\acute{e}r}_B(C,L)\xrightarrow{u^0}
  \operatorname{D\acute{e}r}_A(C,L)\xrightarrow{v^0}
  \operatorname{D\acute{e}r}_A(B,L)\xrightarrow{\partial}
  \operatorname{Exan}_{B/A}(C,L)\to0
\end{equation}
(resp.
\begin{equation}
\tag{20.2.4.2}
\label{0.20.2.4.2-es}
  0\to\operatorname{D\acute{e}r}_B(C,L)\xrightarrow{u^0}
  \operatorname{D\acute{e}r}_A(C,L)\xrightarrow{v^0}
  \operatorname{D\acute{e}r}_A(B,L)\xrightarrow{\partial}
  \operatorname{Exalcom}_{B/A}(C,L)\to0\text{).}
\end{equation}
Esto resulta de la definición de $\operatorname{Exan}_{B/A}(C,L)$ (resp.
$\operatorname{Exalcom}_{B/A}(C,L)$) ((18.3.7) y (18.4.2)).
\end{corollary}

\begin{remark}[20.2.5]
\label{0.20.2.5-es}
Supongamos que se tiene un diagrama conmutativo de homomorfismos de anillos
\[
\xymatrix{
  A \ar[r] & B \ar[r] & C \\
  A' \ar[u] \ar[r] & B' \ar[u] \ar[r] & C' \ar[u]
}
\]
Entonces se tiene un diagrama conmutativo
\[
\xymatrix@C=0.5em{
  0 \ar[r] & \operatorname{D\acute{e}r}_B(C,L) \ar[r] \ar[d] &
  \operatorname{D\acute{e}r}_A(C,L) \ar[r] \ar[d] &
  \operatorname{D\acute{e}r}_A(B,L) \ar[r]^-\partial \ar[d] &
  \operatorname{Exan}_B(C,L) \ar[r] \ar[d] &
  \operatorname{Exan}_A(C,L) \ar[r] \ar[d] &
  \operatorname{Exan}_A(B,L) \ar[d] \\
  0 \ar[r] & \operatorname{D\acute{e}r}_{B'}(C',L) \ar[r] &
  \operatorname{D\acute{e}r}_{A'}(C',L) \ar[r] &
  \operatorname{D\acute{e}r}_{A'}(B',L) \ar[r]^-\partial &
  \operatorname{Exan}_{B'}(C',L) \ar[r] &
  \operatorname{Exan}_{A'}(C',L) \ar[r] &
  \operatorname{Exan}_{A'}(B',L)
}
\]
y del mismo modo para las sucesiones exactas \sref{0.20.2.3.1-es},
\sref{0.20.2.4.1-es} y \sref{0.20.2.4.2-es}.
\end{remark}

\subsection{Derivaciones continuas en los anillos topológicos}
\label{subsection:0.20.3-es}

\begin{env}[20.3.1]
\label{0.20.3.1-es}
Dados dos anillos topológicos $A$, $B$ (linealmente topologizados
como siempre), denotaremos por $\operatorname{Hom.\ cont}(A,B)$
el conjunto de los homomorfismos continuos de $A$ en $B$. Dado un anillo
topológico $A$, la categoría de los $A$-anillos topológicos se define como
la de los $A$-anillos (18.1.4), sustituyendo en todas partes «anillo» por «anillo
topológico» y «homomorfismo» por «homomorfismo continuo»; si $B$

\oldpage[0\textsubscript{IV-1}]{122}
y $C$ son dos $A$-anillos topológicos, se denotará por
$\operatorname{Hom.\ cont}_A(B,C)$ el conjunto de los $A$-homomorfismos continuos
de $B$ en $C$.

Sean $A$ un anillo topológico, $B$ un $A$-anillo topológico, $L$ un
$B$-bimódulo topológico; se denota por $\operatorname{D\acute{e}r.\ cont}_A(B,L)$
el conjunto de las $A$-derivaciones \emph{continuas} de $B$ en $L$; está claro
que es un \emph{subgrupo} de $\operatorname{D\acute{e}r}_A(B,L)$ (y un
sub-$A$-módulo cuando $A$ y $B$ son conmutativos y $L$ un $B$-módulo).
\end{env}

\begin{env}[20.3.2]
\label{0.20.3.2-es}
Es inmediato que la prop. \sref{0.20.1.1-es} subsiste cuando se sustituye en ella
«anillo» por «anillo topológico» y «homomorfismo» por
«homomorfismo continuo». Del mismo modo, \sref{0.20.1.3-es} y
\sref{0.20.1.4-es} siguen siendo válidos sustituyendo $\operatorname{D\acute{e}r}$ por
$\operatorname{D\acute{e}r.\ cont}$, «anillo» (resp. «bimódulo») por «anillo
topológico» (resp. «bimódulo topológico»), «homomorfismo» por
«homomorfismo continuo»; naturalmente hay que suponer en
\sref{0.20.1.4-es} que $p$ es continuo, y sustituir «$A$-equivalencias» por
«$A$-equivalencias continuas». Se tienen los mismos resultados para
\sref{0.20.1.5-es} y \sref{0.20.1.6-es}, a condición de tomar como
topología sobre $D_C(L)$ (resp. $D_B(L)$) la topología \emph{producto} (sobre
$C\times L$, resp. $B\times L$).
\end{env}

\begin{proposition}[20.3.3]
\label{0.20.3.3-es}
Sean $A$ un anillo topológico, $B$ un $A$-anillo topológico, $L$ un
$B$-bimódulo topológico discreto anulado por un ideal bilateral abierto de $B$.
Si en $B$ el cuadrado de todo ideal bilateral abierto es abierto, se tiene
$\operatorname{D\acute{e}r.\ cont}_A(B,L)=\operatorname{D\acute{e}r}_A(B,L)$.
\end{proposition}

\begin{proof}
En efecto, si $\mathfrak{R}$ es un ideal bilateral abierto de $B$ que anula a $L$,
y $D$ una $A$-derivación de $B$ en $L$, se tiene
$D(\mathfrak{R}^2)\subset\mathfrak{R}.D(\mathfrak{R})+
D(\mathfrak{R}).\mathfrak{R}=0$ \sref{0.20.1.1.1-es}, por tanto $D$ es continua.
\end{proof}

\begin{env}[20.3.4]
\label{0.20.3.4-es}
Todos los resultados de \sref{0.20.2.1-es} siguen siendo válidos cuando se sustituye
«anillo» por «anillo topológico», «bimódulo» por «bimódulo
topológico», «homomorfismo» por «homomorfismo continuo» y
$\operatorname{D\acute{e}r}$ por $\operatorname{D\acute{e}r.\ cont}$.
\end{env}

\begin{proposition}[20.3.5]
\label{0.20.3.5-es}
Sean $A$ un anillo topológico, $B$ un $A$-anillo topológico, $L$ un
$B$-bimódulo topológico discreto anulado por un ideal bilateral abierto de $B$.
Se tiene entonces un isomorfismo canónico
\begin{equation}
\tag{20.3.5.1}
\label{0.20.3.5.1-es}
  \varinjlim\operatorname{D\acute{e}r}_{A/\mathfrak{J}}
  (B/\mathfrak{R},L)\xrightarrow{\sim}
  \operatorname{D\acute{e}r.\ cont}_A(B,L)
\end{equation}
donde en el primer miembro el límite inductivo se toma según el conjunto
ordenado filtrante de los pares $(\mathfrak{J},\mathfrak{R})$ de ideales bilaterales
tales que $\mathfrak{R}.L=L.\mathfrak{R}=0$,
$\mathfrak{J}.B\subset\mathfrak{R}$,
$B.\mathfrak{J}\subset\mathfrak{R}$.
\end{proposition}

\begin{proof}
Como $A/\mathfrak{J}$ y $B/\mathfrak{R}$ son discretos, se tienen
homomorfismos canónicos
$w_{\mathfrak{J},\mathfrak{R}}:
\operatorname{D\acute{e}r}_{A/\mathfrak{J}}(B/\mathfrak{R},L)\to
\operatorname{D\acute{e}r.\ cont}_A(B,L)$ que forman un sistema inductivo
\sref{0.20.3.4-es}, de donde se obtiene el homomorfismo \sref{0.20.3.5.1-es} por paso al
límite inductivo. Como el homomorfismo
$B/\mathfrak{R}'\to B/\mathfrak{R}$ es sobreyectivo para
$\mathfrak{R}\supset\mathfrak{R}'$, resulta inmediatamente de la definición que
el homomorfismo
$\operatorname{D\acute{e}r}_{A/\mathfrak{J}}(B/\mathfrak{R},L)\to
\operatorname{D\acute{e}r}_{A/\mathfrak{J}}(B/\mathfrak{R}',L)$ (con
$\mathfrak{R}.L=L.\mathfrak{R}=0$,
$\mathfrak{J}.B\subset\mathfrak{R}'$,
$B.\mathfrak{J}\subset\mathfrak{R}'$) es inyectivo, y evidentemente ocurre
lo mismo con el homomorfismo
$\operatorname{D\acute{e}r}_{A/\mathfrak{J}}(B/\mathfrak{R},L)\to
\operatorname{D\acute{e}r}_{A/\mathfrak{J}'}(B/\mathfrak{R},L)$ para
$\mathfrak{J}'\subset\mathfrak{J}$ \sref{0.20.2.1-es}; se concluye que
el homomorfismo \sref{0.20.3.5.1-es} es inyectivo. Por otra parte, si $D$ es
una $A$-derivación continua de $B$ en $L$, su núcleo contiene un ideal
bilateral abierto $\mathfrak{R}_0$ de $B$, y si $\mathfrak{J}_0$ es un ideal
bilateral abierto de $A$ tal que $\mathfrak{J}_0B\subset\mathfrak{R}_0$ y
$B\mathfrak{J}_0\subset\mathfrak{R}_0$, está claro que $D$ es la imagen
canónica de una $(A/\mathfrak{J}_0)$-derivación de
$B/\mathfrak{R}_0$ en $L$, luego \sref{0.20.3.5.1-es} es sobreyectivo.
\end{proof}

\oldpage[0\textsubscript{IV-1}]{123}
\begin{proposition}[20.3.6]
\label{0.20.3.6-es}
Sean $u:A\to B$, $v:B\to C$ dos homomorfismos continuos de anillos
topológicos, $L$ un $C$-bimódulo discreto anulado por un ideal bilateral abierto
de $C$. Se tiene una sucesión exacta canónica
\begin{align}
\tag{20.3.6.1}
\label{0.20.3.6.1-es}
  0\to\operatorname{D\acute{e}r.\ cont}_B(C,L)\xrightarrow{u^0}
  \operatorname{D\acute{e}r.\ cont}_A(C,L)&\xrightarrow{v^0}
  \operatorname{D\acute{e}r.\ cont}_A(B,L)\xrightarrow{\partial} \\
  &\to\operatorname{Exantop}_B(C,L)\xrightarrow{u^1}
  \operatorname{Exantop}_A(C,L)\xrightarrow{v^1}
  \operatorname{Exantop}_A(B,L) \notag
\end{align}
donde $\partial$ se define por paso al límite inductivo a partir del
homomorfismo $\partial$ de \sref{0.20.2.2.1-es}; esta sucesión exacta es
funtorial en $L$ (en la categoría de los $C$-bimódulos discretos anulados por
ideales bilaterales abiertos).
\end{proposition}

\begin{proof}
Esto resulta de la exactitud del funtor $\varinjlim$, a partir de
\sref{0.20.2.2-es}.
\end{proof}

\begin{corollary}[20.3.7]
\label{0.20.3.7-es}
Sean $A$, $B$, $C$ tres anillos topológicos conmutativos, $u:A\to B$,
$v:B\to C$ dos homomorfismos continuos, $L$ un $C$-módulo discreto anulado
por un ideal abierto de $C$. Se tiene una sucesión exacta canónica de $A$-módulos,
funtorial en $L$,
\begin{align}
\tag{20.3.7.1}
\label{0.20.3.7.1-es}
  &0\to\operatorname{D\acute{e}r.\ cont}_B(C,L)\xrightarrow{u^0}
  \operatorname{D\acute{e}r.\ cont}_A(C,L)\xrightarrow{v^0}
  \operatorname{D\acute{e}r.\ cont}_A(B,L)\xrightarrow{\partial} \\
  &\to\operatorname{Exalcotop}_B(C,L)\xrightarrow{u^1}
  \operatorname{Exalcotop}_A(C,L)\xrightarrow{v^1}
  \operatorname{Exalcotop}_A(B,L). \notag
\end{align}
\end{corollary}

\begin{corollary}[20.3.8]
\label{0.20.3.8-es}
Bajo las hipótesis de \sref{0.20.3.5-es} (resp.
\sref{0.20.3.6-es}) se tiene una sucesión exacta canónica, funtorial en $L$
\begin{equation}
\tag{20.3.8.1}
\label{0.20.3.8.1-es}
  0\to\operatorname{D\acute{e}r.\ cont}_B(C,L)\xrightarrow{u^0}
  \operatorname{D\acute{e}r.\ cont}_A(C,L)\xrightarrow{v^0}
  \operatorname{D\acute{e}r.\ cont}_A(B,L)\to
  \operatorname{Exantop}_{B/A}(C,L)\to0
\end{equation}
(resp.
\begin{equation}
\tag{20.3.8.2}
\label{0.20.3.8.2-es}
  0\to\operatorname{D\acute{e}r.\ cont}_B(C,L)\xrightarrow{u^0}
  \operatorname{D\acute{e}r.\ cont}_A(C,L)\xrightarrow{v^0}
  \operatorname{D\acute{e}r.\ cont}_A(B,L)\to
  \operatorname{Exalcotop}_{B/A}(C,L)\to0\text{).}
\end{equation}
\end{corollary}

Dejamos al lector la tarea de escribir los diagramas análogos a los de
\sref{0.20.2.5-es}.

\subsection{Partes principales y diferenciales}
\label{subsection:0.20.4-es}

En todo lo que sigue de este parágrafo y en los tres siguientes, todos los
anillos se suponen \emph{conmutativos}.

\begin{env}[20.4.1]
\label{0.20.4.1-es}
Sean $A$ un anillo topológico, $B$ una $A$-álgebra topológica; la
$A$-álgebra $B\otimes_A B$ estará provista de la topología \emph{producto
tensorial}, que la convierte en una $A$-álgebra topológica; denotaremos por $p$
(o $p_{B/A}$) el $A$-homomorfismo canónico sobreyectivo
\begin{equation}
\tag{20.4.1.1}
\label{0.20.4.1.1-es}
  p:B\otimes_A B\to B
\end{equation}
tal que $p(b\otimes b')=bb'$; es inmediato que $p$ es continuo. El
\emph{núcleo} de $p$ se denotará por $\mathfrak{J}_{B/A}$ (o solamente
$\mathfrak{J}$ si no hay confusión). Denotaremos por
$j_1:B\to B\otimes_A B$ y $j_2:B\to B\otimes_A B$ los dos
$A$-homomorfismos canónicos, tales que
\[
  j_1(b)=b\otimes1, \qquad j_2(b)=1\otimes b
\]
que son continuos.
\end{env}

\oldpage[0\textsubscript{IV-1}]{124}
\begin{definition}[20.4.2]
\label{0.20.4.2-es}
Se llama \emph{$B$-álgebra aumentada de partes principales de orden $1$ de
$B$ respecto de $A$} y se denota por $\mathrm{P}^1_{B/A}$ la $A$-álgebra
topológica cociente
\begin{equation}
\tag{20.4.2.1}
\label{0.20.4.2.1-es}
  \mathrm{P}^1_{B/A}=(B\otimes_A B)/\mathfrak{J}^2
\end{equation}
provista de la estructura de $B$-álgebra topológica definida por el homomorfismo
$\bar{j}_1:B\to\mathrm{P}^1_{B/A}$ (deducido de $j_1$ por composición con
el homomorfismo canónico $B\otimes_A B\to\mathrm{P}^1_{B/A}$), y de
la aumentación de $B$-álgebra $\varepsilon:\mathrm{P}^1_{B/A}\to B$ (también
denotada por $\varepsilon_{B/A}$) deducida de $p$ por paso al cociente.

Como $p(b\otimes1)=b$ por definición, está claro que $\varepsilon$ es
efectivamente una aumentación de $B$-álgebra.
\end{definition}

\begin{definition}[20.4.3]
\label{0.20.4.3-es}
El núcleo de la aumentación $\varepsilon:\mathrm{P}^1_{B/A}\to B$,
\begin{equation}
\tag{20.4.3.1}
\label{0.20.4.3.1-es}
  \Omega^1_{B/A}=\mathfrak{J}_{B/A}/(\mathfrak{J}_{B/A})^2
\end{equation}
provisto de la topología inducida por la de $\mathrm{P}^1_{B/A}$, que lo convierte
en un $B$-módulo topológico, se llama el $B$-módulo de las
\emph{$1$-diferenciales} (o simplemente de las \emph{diferenciales}) de $B$
respecto de $A$.
\end{definition}

Obsérvese que la topología de $\Omega^1_{B/A}$ es también la topología cociente
de la topología inducida sobre $\mathfrak{J}_{B/A}$ por la de
$B\otimes_A B$ (Bourbaki, \emph{Top. gén.}, cap.~III, 3\textsuperscript{a}
ed., §~2, n\textsuperscript{o}~7, prop.~20). Si $B$ es discreto, lo mismo ocurre
con $\Omega^1_{B/A}$. Se denota por $\widehat{\Omega}^1_{B/A}$ el
\emph{completado separado} del $B$-módulo topológico $\Omega^1_{B/A}$.

Todo anillo topológico $B$ puede considerarse como una
$\mathbf{Z}$-álgebra topológica (donde $\mathbf{Z}$ está provisto de la topología
discreta), por lo que se puede definir el $B$-módulo topológico $\Omega^1_{B/\mathbf{Z}}$,
que también se llama a veces el $B$-módulo de las \emph{diferenciales absolutas}
sobre $B$ y se denota por $\Omega^1_B$. Si $B$ es un álgebra topológica sobre
un cuerpo primo $P$ (discreto), se tiene
$B\otimes_P B=B\otimes_{\mathbf{Z}}B$, luego
$\Omega^1_{B/P}=\Omega^1_B$.

\begin{lemma}[20.4.4]
\label{0.20.4.4-es}
Sean $A$ un anillo, $B$ una $A$-álgebra. El ideal $\mathfrak{J}_{B/A}$ de
$B\otimes_A B$ está generado por los elementos $1\otimes s-s\otimes1$, donde $s$
recorre un conjunto de generadores de la $A$-álgebra $B$.
\end{lemma}

\begin{proof}
Está claro que, para todo $x\in B$, se tiene
$x\otimes1-1\otimes x\in\mathfrak{J}$; por otra parte, cualesquiera que sean $x$,
$y$ en $B$, se tiene
$x\otimes y=xy\otimes1+(x\otimes1)(1\otimes y-y\otimes1)$. Si
$\sum_i(x_i\otimes y_i)\in\mathfrak{J}$, se tiene por definición
$\sum_i x_i y_i=0$, luego
\begin{equation}
\tag{20.4.4.1}
\label{0.20.4.4.1-es}
  \sum_i(x_i\otimes y_i)=
  \sum_i(x_i\otimes1)(1\otimes y_i-y_i\otimes1)
\end{equation}
lo que demuestra que $\mathfrak{J}$ es el ideal generado por los elementos
$1\otimes x-x\otimes1$. Además, si $x=st$, se tiene
\begin{equation}
\tag{20.4.4.2}
\label{0.20.4.4.2-es}
  x\otimes1-1\otimes x=(s\otimes1)(t\otimes1-1\otimes t)
  +(s\otimes1-1\otimes s)(1\otimes t)
\end{equation}
lo que concluye inmediatamente la demostración por recurrencia.
\end{proof}

\begin{proposition}[20.4.5]
\label{0.20.4.5-es}
Sean $A$ un anillo topológico, $B$ una $A$-álgebra topológica. La
topología de $\Omega^1_{B/A}$ es menos fina que la topología deducida de la
de $B$ (19.0.2); si en $B$ el cuadrado de todo ideal abierto es abierto, estas
dos topologías son idénticas.
\end{proposition}

\begin{proof}
La primera aserción es trivial, pues la topología producto tensorial sobre
$B\otimes_A B$ es menos fina que la topología deducida de la de $B$;
\emph{a fortiori}, la topología inducida sobre
$\mathfrak{J}=\mathfrak{J}_{B/A}$

\oldpage[0\textsubscript{IV-1}]{125}
por la de $B\otimes_A B$ es menos fina que la topología sobre $\mathfrak{J}$
deducida de la de $B$. Para demostrar la segunda aserción, escribamos
Denotemos por $M\otimes N$, por abuso de notación, el submódulo
$\operatorname{Im}(M\otimes_A N)$, para dos sub-$A$-módulos $M$, $N$ de $B$.
Utilizando la relación
\[
\begin{aligned}
  (xy)\otimes z-x\otimes(yz)
  &=(x\otimes1)(1\otimes z)(y\otimes1-1\otimes y)\\
  &=xz.(y\otimes1-1\otimes y)
    -x.(z\otimes1-1\otimes z)(y\otimes1-1\otimes y)
\end{aligned}
\]
en el $B$-módulo $B\otimes_A B$ (definido por $j_1$), se ve inmediatamente,
teniendo en cuenta \sref{0.20.4.4-es}, que, si $\mathfrak{R}$ es un ideal de
$B$, se tiene
\[
  ((\mathfrak{R}^2\otimes B)+(B\otimes\mathfrak{R}^2))\cap\mathfrak{J}
  \subset(\mathfrak{R}\otimes\mathfrak{R})\cap\mathfrak{J}
  +\mathfrak{R}\mathfrak{J}+\mathfrak{J}^2
\]
y por otra parte, si $x_i$, $y_i$ son elementos de $\mathfrak{R}$ tales
que $\sum_i(x_i\otimes y_i)\in\mathfrak{J}$, se deduce de
\sref{0.20.4.4.1-es} que se tiene
$\sum_i(x_i\otimes y_i)\in\mathfrak{R}\mathfrak{J}$, de modo que finalmente
\begin{equation}
\tag{20.4.5.1}
\label{0.20.4.5.1-es}
  (\mathfrak{R}^2\otimes B+B\otimes\mathfrak{R}^2)\cap\mathfrak{J}
  \subset\mathfrak{R}\mathfrak{J}+\mathfrak{J}^2.
\end{equation}
Ahora bien, se obtiene un sistema fundamental de entornos de $0$ en $\mathfrak{J}$
(para la topología inducida por la de $B\otimes_A B$) tomando como
entornos de $0$ los conjuntos
$(\mathfrak{R}\otimes B+B\otimes\mathfrak{R})\cap\mathfrak{J}$, donde
$\mathfrak{R}$ recorre el conjunto de los ideales abiertos. Como la topología de
$\Omega^1_{B/A}$ deducida de la de $B$ es también el cociente por
$\mathfrak{J}^2$ de la topología de $\mathfrak{J}$ deducida de la de $B$,
la hipótesis sobre los ideales abiertos de $B$ y la relación
\sref{0.20.4.5.1-es} terminan de demostrar la segunda aserción.
\end{proof}

\begin{definition}[20.4.6]
\label{0.20.4.6-es}
Sean $p_1$ y $p_2$ las aplicaciones compuestas
$B\xrightarrow{j_1}B\otimes_A B\to\mathrm{P}^1_{B/A}$,
$B\xrightarrow{j_2}B\otimes_A B\to\mathrm{P}^1_{B/A}$, que son
$A$-homomorfismos continuos tales que
$\varepsilon\circ p_1=\varepsilon\circ p_2=\mathrm{I}_B$. El
$A$-homomorfismo continuo de $A$-módulos
\begin{equation}
\tag{20.4.6.1}
\label{0.20.4.6.1-es}
  d_{B/A}=p_1-p_2:B\to\Omega^1_{B/A}
\end{equation}
se denomina la \emph{diferencial exterior de $B$ relativa a $A$}; para
todo $x\in B$, $d_{B/A}(x)$ (denotado también $d(x)$ o $dx$) se denomina la
\emph{diferencial de $x$} (relativa a $A$).
\end{definition}

Cuando $A=\mathbf{Z}$, se escribe $d_B$ en lugar de $d_{B/\mathbf{Z}}$; si $B$
es un álgebra sobre un cuerpo primo $P$, se tiene $d_B=d_{B/P}$.

\begin{proposition}[20.4.7]
\label{0.20.4.7-es}
El $B$-módulo $\Omega^1_{B/A}$ está generado por los elementos $d_{B/A}(x)$,
donde $x$ recorre un sistema de generadores del $A$-álgebra $B$.
\end{proposition}

\begin{proof}
Como $d_{B/A}(x)$ es la imagen canónica de $x\otimes1-1\otimes x$ en
$\Omega^1_{B/A}$, la proposición es una consecuencia inmediata de
\sref{0.20.4.4-es}.
\end{proof}

\begin{theorem}[20.4.8]
\label{0.20.4.8-es}
Sean $A$ un anillo topológico y $B$ un $A$-álgebra topológica.
\begin{enumerate}
  \item[\rm(i)] Existe un único isomorfismo de $B$-álgebras topológicas
  aumentadas
\begin{equation}
\tag{20.4.8.1}
\label{0.20.4.8.1-es}
  \varphi:\mathrm{P}^1_{B/A}\xrightarrow{\sim}D_B(\Omega^1_{B/A})
\end{equation}
cuya restricción a $\Omega^1_{B/A}$ es la identidad.
  \item[\rm(ii)] El homomorfismo $d_{B/A}$ es una $A$-derivación de $B$
  en $\Omega^1_{B/A}$, con la siguiente propiedad universal: para todo
  $B$-módulo topológico $L$, la aplicación $u\mapsto u\circ d_{B/A}$ es
  un isomorfismo de $A$-módulos
\begin{equation}
\tag{20.4.8.2}
\label{0.20.4.8.2-es}
  \operatorname{Hom.\ cont}_B(\Omega^1_{B/A},L)
  \xrightarrow{\sim}\operatorname{D\acute{e}r.\ cont}_A(B,L).
\end{equation}
\end{enumerate}
\end{theorem}

\begin{proof}
\oldpage[0\textsubscript{IV-1}]{126}
(i) Es inmediato que $\varphi$ (con las notaciones de
\sref{0.20.4.6-es}) es necesariamente la aplicación
$z\mapsto(\varepsilon(z),z-p_1(\varepsilon(z)))$, y el isomorfismo
inverso la aplicación $(b,x)\mapsto p_1(b)+x$; como estas dos aplicaciones
son continuas, esto demuestra la primera aserción. Obsérvese que esto
implica que la topología de $B$ se identifica (mediante $p_1$) con el \emph{cociente}
de la topología de $B\otimes_A B$ por $\mathfrak{J}_{B/A}$.

(ii) El hecho de que $d_{B/A}$ sea una $A$-derivación de $B$ se deduce de la
definición \sref{0.20.4.6-es} y de \sref{0.20.1.1-es}.

Para demostrar la propiedad universal, recordemos que
$\operatorname{D\acute{e}r.\ cont}_A(B,L)$ se identifica canónicamente con
el conjunto $G$ de los homomorfismos continuos de $A$-álgebras
$u:B\to D_B(L)$ tales que la composición
$B\xrightarrow{u}D_B(L)\xrightarrow{q}B$ sea la identidad (donde $q(b,0)=b$ es
la proyección; cf. \sref{0.20.1.6-es} y \sref{0.20.3.2-es}). Por otra parte,
gracias al isomorfismo $\varphi$,
$\operatorname{Hom.\ cont}_B(\Omega^1_{B/A},L)$ se identifica canónicamente con
el conjunto de los homomorfismos continuos de $B$-álgebras
$v:\mathrm{P}^1_{B/A}\to D_B(L)$ tales que la composición
$\mathrm{P}^1_{B/A}\xrightarrow{v}D_B(L)\xrightarrow{q}B$ sea
el aumento $\varepsilon$. Como por definición se tiene $p_2=p_1-d_{B/A}$,
basta demostrar que todo $u\in G$ se factoriza como
\[
\xymatrix{
  B \ar[r]^u \ar[d]_{p_2} & D_B(L) \\
  \mathrm{P}^1_{B/A} \ar[ur]_v &
}
\]
donde $v$ es un $B$-homomorfismo continuo. Ahora bien, ya se tiene un homomorfismo
continuo de $A$-álgebras $j:b\mapsto(b,0)$ de $B$ en $D_B(L)$, que
pertenece a $G$; por la definición del producto tensorial topológico
de álgebras topológicas ($\mathbf{0_I}$, 7.7.6), existe por tanto un
$A$-homomorfismo continuo de álgebras
$w:B\otimes_A B\to D_B(L)$ que hace conmutativo el diagrama
\[
\xymatrix{
  B\otimes_A B \ar[dr]^w & B \ar[l]_{j_1} \ar[d]^j \\
  B \ar[u]^{j_2} \ar[r]_u & D_B(L)
}
\]
Se tiene, pues, por definición
$w(b\otimes1-1\otimes b)=j(b)-u(b)\in L$, y en virtud de
\sref{0.20.4.4-es}, esto implica $w(\mathfrak{J})\subset L$ y por tanto
$w(\mathfrak{J}^2)=0$; en consecuencia, $w$ se factoriza como
\[
  B\otimes_A B\to\mathrm{P}^1_{B/A}\xrightarrow{v}D_B(L)
\]
donde $v$ es un homomorfismo continuo de $A$-álgebras; además, como
$v\circ p_1=j$ es un homomorfismo de $B$-álgebras, lo mismo ocurre con
$v$ por definición de la estructura de $B$-álgebra de
$\mathrm{P}^1_{B/A}$; como por definición se tiene $u=v\circ p_2$, esto
termina la demostración.
\end{proof}

\begin{theorem}[20.4.9]
\label{0.20.4.9-es}
Supongamos que $B$ es un $A$-álgebra topológica formalmente lisa. Entonces
el $B$-módulo topológico $\Omega^1_{B/A}$ es \emph{formalmente proyectivo}.
\end{theorem}

\begin{proof}
En efecto, la hipótesis implica que $B\otimes_A B$ (dotado de la estructura de
$B$-álgebra topológica definida por $j_1$) es un $B$-álgebra topológica
formalmente lisa \sref{0.19.3.5-es}, (iii), y por consiguiente también un
$A$-álgebra topológica formalmente lisa \sref{0.19.3.5-es}, (ii); como
$B$ es topológicamente isomorfa al $A$-álgebra cociente
$(B\otimes_A B)/\mathfrak{J}$ y es un $A$-álgebra formalmente lisa, la
conclusión se deduce de \sref{0.19.5.3-es}.
\end{proof}

\oldpage[0\textsubscript{IV-1}]{127}
\begin{corollary}[20.4.10]
\label{0.20.4.10-es}
Supongamos además que en $B$ el cuadrado de todo ideal abierto es abierto;
entonces, para todo ideal abierto $\mathfrak{R}$ de $B$,
$\Omega^1_{B/A}\otimes_B(B/\mathfrak{R})
=\Omega^1_{B/A}/\mathfrak{R}.\Omega^1_{B/A}$ es un
$(B/\mathfrak{R})$-módulo proyectivo.
\end{corollary}

\begin{proof}
En efecto, la topología de $\Omega^1_{B/A}$ es entonces la deducida de la de
$B$ \sref{0.20.4.5-es}, y basta aplicar \sref{0.19.2.4-es}.
\end{proof}

\begin{corollary}[20.4.11]
\label{0.20.4.11-es}
Sean $A$, $B$ dos anillos locales noetherianos y $\rho:A\to B$ un
homomorfismo local que hace de $B$ un $A$-álgebra topológica formalmente
lisa (para las topologías preádicas). Entonces, para todo ideal de definición
$\mathfrak{b}$ de $B$,
$\Omega^1_{B/A}\otimes_B(B/\mathfrak{b})$ es un
$(B/\mathfrak{b})$-módulo libre.
\end{corollary}

\begin{proof}
En efecto, de \sref{0.20.4.10-es} se deduce que este módulo es proyectivo, y
como $B/\mathfrak{b}$ es un anillo artiniano, todo
$(B/\mathfrak{b})$-módulo proyectivo es libre
($\mathbf{0_{III}}$, 10.1.3).
\end{proof}

\begin{proposition}[20.4.12]
\label{0.20.4.12-es}
Sean $A$ un anillo topológico y $B$ un $A$-álgebra topológica. Si
el homomorfismo estructural $A\to B$ es sobreyectivo, se tiene
$\Omega^1_{B/A}=0$.
\end{proposition}

\begin{proof}
En efecto, se tiene $\operatorname{D\acute{e}r}_A(B,L)=0$ para todo $B$-módulo
$L$ en virtud de \sref{0.20.1.1-es}, y la proposición se deduce inmediatamente
de \sref{0.20.4.8-es}.
\end{proof}

\begin{examples}[20.4.13]
\label{0.20.4.13-es}
\begin{enumerate}
  \item[\rm(i)] Sean $A$ un anillo y
  $B=A[X_\alpha]_{\alpha\in I}$ un álgebra de polinomios sobre $A$. Entonces
  $\Omega^1_{B/A}$ es un $B$-módulo libre, y los $dX_\alpha$ forman una
  base. En efecto, los $dX_\alpha$ generan este $B$-módulo
  \sref{0.20.4.7-es}. Por otra parte, si $L$ es un $B$-módulo libre, con
  una base $(z_\alpha)_{\alpha\in I}$ cuyo conjunto de índices es $I$, existe
  un $A$-homomorfismo $u$ de $B$ en $D_B(L)$ tal que
  $u(X_\alpha)=(X_\alpha,z_\alpha)$, para todo $\alpha\in I$, y por tanto, por
  \sref{0.20.1.5-es}, una $A$-derivación $D$ de $B$ en $L$ tal que
  $D(X_\alpha)=z_\alpha$ para todo $\alpha$; en virtud de
  \sref{0.20.4.8.1-es}, esto demuestra que los $dX_\alpha$ son linealmente
  independientes.
  \item[\rm(ii)] Sean $A$ un anillo, $L$ un $A$-módulo y $B$ la
  $A$-álgebra $D_A(L)$; entonces el homomorfismo canónico
  $x\mapsto d_{B/A}x$ de $L$ en $\Omega^1_{B/A}$ es \emph{biyectivo}, pues
  es inmediato que las $B$-derivaciones de $B=D_A(L)$ en un $B$-módulo
  $M$ son las aplicaciones de la forma $(b,x)\mapsto u(x)$, donde
  $u\in\operatorname{Hom}_A(L,M)$; se concluye, pues, por
  \sref{0.20.4.8-es}, (ii).
  \item[\rm(iii)] Supongamos que $B$ es el producto de dos $A$-álgebras
  topológicas $B_1$, $B_2$ (identificadas con ideales de $B$). Entonces las
  imágenes de $B_1\otimes_A B_2$ y de $B_2\otimes_A B_1$ por el homomorfismo
  $p$ \sref{0.20.4.1.1-es} son nulas, de donde se deduce inmediatamente que
  $\mathrm{P}^1_{B/A}$ se identifica con el \emph{producto}
  $\mathrm{P}^1_{B_1/A}\times\mathrm{P}^1_{B_2/A}$, y que el $B$-módulo
  $\Omega^1_{B/A}$ es \emph{suma directa} (topológica) de los $B$-módulos
  $\Omega^1_{B_1/A}$ y $\Omega^1_{B_2/A}$ (anulados respectivamente por
  $B_2$ y $B_1$).
  \item[\rm(iv)] Sean $A$, $B$ dos anillos íntegros tales que $A\subset
  B$, que $A$ sea integralmente cerrado, $B$ entero sobre $A$, y que el cuerpo
  de fracciones de $B$ sea una extensión separable del de $A$. Entonces
  el $B$-módulo $\Omega^1_{B/A}$ es un \emph{módulo de torsión}. En efecto,
  para todo $x\in B$, el polinomio minimal de $x$ respecto del cuerpo de
  fracciones $K$ de $A$ es un polinomio $f(T)$ perteneciente a $A[T]$; como
  $x$ es separable sobre $K$, se tiene $f'(x)\neq0$ y, por otra parte, de
  la relación $f(x)=0$ se deduce que $f'(x)d_{B/A}x=0$, de donde se sigue
  la aserción en virtud de \sref{0.20.4.7-es}.
\end{enumerate}
\end{examples}

\begin{remarks}[20.4.14]
\label{0.20.4.14-es}
\begin{enumerate}
  \item[\rm(i)] Obsérvese que el $B$-módulo $\Omega^1_{B/A}$, privado de su
  topología, es independiente de las topologías de $A$ y de $B$.
  \item[\rm(ii)] Introduciremos más adelante el «álgebra de partes
  principales de orden $n$» de $B$ respecto de $A$,
  $\mathrm{P}^n_{B/A}=(B\otimes_A B)/(\mathfrak{J}_{B/A})^{n+1}$, que constituye
  la base del «cálculo diferencial de orden $n$».
\end{enumerate}
\end{remarks}

\subsection{Propiedades funtoriales fundamentales de
\texorpdfstring{$\Omega^1_{B/A}$}{Omega1(B/A)}}
\label{subsection:0.20.5-es}

\oldpage[0\textsubscript{IV-1}]{128}
\begin{env}[20.5.1]
\label{0.20.5.1-es}
En todo este número y el siguiente, salvo mención expresa en contrario, se
supone que los anillos y módulos considerados están provistos de la topología
\emph{discreta}.
\end{env}

\begin{env}[20.5.2]
\label{0.20.5.2-es}
Sean $A$ un anillo, $B$, $C$ dos $A$-álgebras, $u:B\to C$ un
$A$-homomorfismo; se tiene un diagrama conmutativo
\[
\xymatrix{
  B\otimes_A B \ar[r]^{u\otimes u} \ar[d]_{p_{B/A}} &
  C\otimes_A C \ar[d]^{p_{C/A}} \\
  B \ar[r]^u & C
}
\tag{20.5.2.1}
\label{0.20.5.2.1-es}
\]
de donde se obtiene, pasando a los cocientes, un $A$-homomorfismo de álgebras
\begin{equation}
\tag{20.5.2.2}
\label{0.20.5.2.2-es}
  u':\mathrm{P}^1_{B/A}\to\mathrm{P}^1_{C/A}
\end{equation}
tal que el diagrama
\[
\xymatrix{
  \mathrm{P}^1_{B/A} \ar[r]^{u'} & \mathrm{P}^1_{C/A} \\
  B \ar[u]^{p_1} \ar[r]^u & C \ar[u]_{p_1}
}
\]
sea conmutativo; como $u\otimes u$ aplica $\mathfrak{J}_{B/A}$ en
$\mathfrak{J}_{C/A}$, se obtiene, restringiendo $u'$ a
$\Omega^1_{B/A}$, una aplicación
\begin{equation}
\tag{20.5.2.3}
\label{0.20.5.2.3-es}
  u'':\Omega^1_{B/A}\to\Omega^1_{C/A}
\end{equation}
tal que el par $(u'',u)$ sea un \emph{dihomomorfismo} para la
estructura de $B$-módulo de $\Omega^1_{B/A}$ y la estructura de $C$-módulo de
$\Omega^1_{C/A}$; este último hecho permite deducir de él canónicamente un
\emph{homomorfismo de $C$-módulos}
\begin{equation}
\tag{20.5.2.4}
\label{0.20.5.2.4-es}
  u_{C/B/A}:\Omega^1_{B/A}\otimes_B C\to\Omega^1_{C/A}.
\end{equation}

\oldpage[0\textsubscript{IV-1}]{129}
Además, como el diagrama
\[
\xymatrix{
  \mathrm{P}^1_{B/A} \ar[r]^{u'} & \mathrm{P}^1_{C/A} \\
  B \ar[u]^{p_2} \ar[r]^u & C \ar[u]_{p_2}
}
\tag{20.5.2.5}
\label{0.20.5.2.5-es}
\]
es también conmutativo, se deduce que el diagrama
\[
\xymatrix{
  \Omega^1_{B/A} \ar[r]^{u''} & \Omega^1_{C/A} \\
  B \ar[u]^{d_{B/A}} \ar[r]^u & C \ar[u]_{d_{C/A}}
}
\tag{20.5.2.6}
\label{0.20.5.2.6-es}
\]
es conmutativo.

Por último, si $w:C\to D$ es un segundo homomorfismo de $A$-álgebras, se tiene la
propiedad de transitividad
\begin{equation}
\tag{20.5.2.7}
\label{0.20.5.2.7-es}
  (w\circ u)_{D/B/A}
  =w_{D/C/A}\circ(u_{C/B/A}\otimes\mathrm{I}_D)
\end{equation}
como resulta de la definición.
\end{env}

\begin{env}[20.5.3]
\label{0.20.5.3-es}
Sean ahora $A$, $B$ dos anillos, $v:A\to B$ un homomorfismo
de anillos, $C$ una $B$-álgebra que se convierte en una $A$-álgebra mediante $v$;
entonces la aplicación canónica $v_0:C\otimes_A C\to C\otimes_B C$ es un
\emph{dihomomorfismo} sobreyectivo de álgebras (relativo a $v:A\to B$) tal que
el diagrama
\[
\xymatrix{
  C\otimes_A C \ar[r]^{v_0} \ar[d]_{p_{C/A}} &
  C\otimes_B C \ar[d]^{p_{C/B}} \\
  C \ar[r]^{\mathrm{I}_C} & C
}
\tag{20.5.3.1}
\label{0.20.5.3.1-es}
\]
sea conmutativo; pasando a los cocientes, se deduce un
dihomomorfismo de álgebras
\begin{equation}
\tag{20.5.3.2}
\label{0.20.5.3.2-es}
  v':\mathrm{P}^1_{C/A}\to\mathrm{P}^1_{C/B}
\end{equation}
tal que el diagrama
\[
\xymatrix{
  \mathrm{P}^1_{C/A} \ar[r]^{v'} & \mathrm{P}^1_{C/B} \\
  C \ar[u]^{p_1} \ar[r]^{\mathrm{I}_C} & C \ar[u]_{p_1}
}
\]
sea conmutativo. Como $v_0$ aplica $\mathfrak{J}_{C/A}$ en
$\mathfrak{J}_{C/B}$ se obtiene, restringiendo $v'$ a
$\Omega^1_{C/A}$, una aplicación
\begin{equation}
\tag{20.5.3.3}
\label{0.20.5.3.3-es}
  v_{C/B/A}:\Omega^1_{C/A}\to\Omega^1_{C/B}
\end{equation}
que es un homomorfismo de $C$-módulos.

\oldpage[0\textsubscript{IV-1}]{130}
Además, como el diagrama
\[
\xymatrix{
  \mathrm{P}^1_{C/A} \ar[r]^{v'} & \mathrm{P}^1_{C/B} \\
  C \ar[u]^{p_2} \ar[r]^{\mathrm{I}_C} & C \ar[u]_{p_2}
}
\tag{20.5.3.4}
\label{0.20.5.3.4-es}
\]
es también conmutativo, se deduce que el diagrama
\[
\xymatrix{
  \Omega^1_{C/A} \ar[r]^{v_{C/B/A}} & \Omega^1_{C/B} \\
  C \ar[u]^{d_{C/A}} \ar[r]^{\mathrm{I}_C} & C \ar[u]_{d_{C/B}}
}
\tag{20.5.3.5}
\label{0.20.5.3.5-es}
\]
es conmutativo.

Por último, si $s:A'\to A$ es un segundo homomorfismo de anillos, se tiene la
propiedad de transitividad
\begin{equation}
\tag{20.5.3.6}
\label{0.20.5.3.6-es}
  (v\circ s)_{C/B/A'}=v_{C/B/A}\circ s_{C/A/A'}.
\end{equation}
\end{env}

\begin{env}[20.5.4]
\label{0.20.5.4-es}
Si ahora se tiene un diagrama conmutativo de homomorfismos de anillos
\[
\xymatrix{
  B \ar[r]^v & B' \\
  A \ar[u]^f \ar[r]^u & A' \ar[u]_g
}
\]
se deduce de \sref{0.20.5.2.4-es} y \sref{0.20.5.3.3-es}, por composición,
un homomorfismo de $B'$-módulos
\begin{equation}
\tag{20.5.4.1}
\label{0.20.5.4.1-es}
  \Omega^1_{B/A}\otimes_B B'
  \xrightarrow{v_{B'/B/A}}\Omega^1_{B'/A}
  \xrightarrow{u_{B'/A'/A}}\Omega^1_{B'/A'}
\end{equation}
tal que el diagrama de $A'$-homomorfismos
\[
\xymatrix{
  \Omega^1_{B/A}\otimes_B B' \ar[r] & \Omega^1_{B'/A'} \\
  B' \ar[u]^{d_{B/A}\otimes1} \ar[r]^{\mathrm{I}_{B'}} &
  B' \ar[u]_{d_{B'/A'}}
}
\tag{20.5.4.2}
\label{0.20.5.4.2-es}
\]
sea conmutativo.

El homomorfismo \sref{0.20.5.4.1-es} corresponde además a un
dihomomorfismo de $B$-módulos
\begin{equation}
\tag{20.5.4.3}
\label{0.20.5.4.3-es}
  \Omega^1_{B/A}\to\Omega^1_{B'/A'}.
\end{equation}
\end{env}

\begin{proposition}[20.5.5]
\label{0.20.5.5-es}
Si $A'$, $B$ son dos $A$-álgebras y $B'=B\otimes_A A'$,
el homomorfismo canónico \sref{0.20.5.4.1-es}
\begin{equation}
\tag{20.5.5.1}
\label{0.20.5.5.1-es}
  \Omega^1_{B/A}\otimes_B B'\to\Omega^1_{B'/A'}
\end{equation}
es biyectivo.
\end{proposition}

\begin{proof}
\oldpage[0\textsubscript{IV-1}]{131}
El primer miembro de \sref{0.20.5.5.1-es} no es entonces otro que
$\Omega^1_{B/A}\otimes_A A'$. Se puede escribir
$B'\otimes_{A'}B'=(B\otimes_A B)\otimes_A A'$ salvo un isomorfismo canónico
y $p_{B'/A'}$ se identifica entonces con
$p_{B/A}\otimes\mathrm{I}_{A'}$; por consiguiente (como $p_{B/A}$ es sobreyectivo)
$\mathfrak{J}_{B'/A'}=\operatorname{Im}(\mathfrak{J}_{B/A}\otimes_A A')$
y
$\mathfrak{J}^2_{B'/A'}
=\operatorname{Im}(\mathfrak{J}^2_{B/A}\otimes_A A')$, de donde
$\mathrm{P}^1_{B'/A'}=\mathrm{P}^1_{B/A}\otimes_A A'$ salvo un isomorfismo
canónico, que hace corresponder entre sí los ideales de aumentación;
como el $A$-módulo $\mathrm{P}^1_{B/A}$ se identifica canónicamente con la
suma directa de $B$ y de $\Omega^1_{B/A}$, se tiene
\begin{equation}
\tag{20.5.5.2}
\label{0.20.5.5.2-es}
  \Omega^1_{B/A}\otimes_A A'=\Omega^1_{B'/A'}
\end{equation}
por el mismo isomorfismo, y se verifica inmediatamente que el compuesto de este
isomorfismo y del isomorfismo canónico
$\Omega^1_{B/A}\otimes_B B'\xrightarrow{\sim}
\Omega^1_{B/A}\otimes_A A'$ no es otro que \sref{0.20.5.5.1-es}.
\end{proof}

\begin{env}[20.5.6]
\label{0.20.5.6-es}
Los homomorfismos canónicos \sref{0.20.5.2.4-es} y
\sref{0.20.5.3.3-es} dan, por funtorialidad, para todo $C$-módulo $L$,
homomorfismos canónicos
\begin{align}
\tag{20.5.6.1}
\label{0.20.5.6.1-es}
  &\operatorname{Hom}_C(\Omega^1_{C/A},L)
  \to\operatorname{Hom}_C(\Omega^1_{B/A}\otimes_B C,L)
  =\operatorname{Hom}_B(\Omega^1_{B/A},L) \\
\tag{20.5.6.2}
\label{0.20.5.6.2-es}
  &\operatorname{Hom}_C(\Omega^1_{C/B},L)
  \to\operatorname{Hom}_C(\Omega^1_{C/A},L).
\end{align}
Teniendo en cuenta \sref{0.20.4.8.2-es} y los diagramas conmutativos
\sref{0.20.5.2.6-es} y \sref{0.20.5.3.5-es}, estos homomorfismos no son
otros (salvo una identificación canónica) que los homomorfismos
\sref{0.20.2.1.2-es} y \sref{0.20.2.1.3-es} respectivamente.
\end{env}

\begin{theorem}[20.5.7]
\label{0.20.5.7-es}
Sean $u:A\to B$, $v:B\to C$ dos homomorfismos de anillos.
\begin{enumerate}
  \item[\rm(i)] La sucesión de $C$-módulos
\begin{equation}
\tag{20.5.7.1}
\label{0.20.5.7.1-es}
  \Omega^1_{B/A}\otimes_B C
  \xrightarrow{v_{C/B/A}}\Omega^1_{C/A}
  \xrightarrow{u_{C/B/A}}\Omega^1_{C/B}\to0
\end{equation}
es exacta.
  \item[\rm(ii)] Para que $v_{C/B/A}$ sea invertible por la izquierda, es
  necesario y suficiente que $C$ sea una $B$-álgebra formalmente lisa relativamente a
  $A$ (para las topologías discretas (cf. \sref{0.19.9.1-es})); en
  particular, basta para ello que $C$ sea una $B$-álgebra formalmente
  lisa (para las topologías discretas).
\end{enumerate}
\end{theorem}

\begin{proof}
(i) La exactitud de la sucesión \sref{0.20.2.4.2-es} muestra en primer lugar,
teniendo en cuenta \sref{0.20.5.6-es}, que la sucesión
\[
  0\to\operatorname{Hom}_C(\Omega^1_{C/B},L)
  \to\operatorname{Hom}_C(\Omega^1_{C/A},L)
  \to\operatorname{Hom}_C(\Omega^1_{B/A}\otimes_B C,L)
\]
es exacta para todo $C$-módulo $L$. Se sabe que esto implica la exactitud
de la sucesión \sref{0.20.5.7.1-es} (Bourbaki, \emph{Alg.}, cap.~II,
3\textsuperscript{a} ed., §~2, n\textsuperscript{o}~1, teo.~1).

(ii) En virtud de la exactitud de \sref{0.20.5.7.1-es}, decir que
$v_{C/B/A}$ es invertible por la izquierda significa que la sucesión
\begin{equation}
\tag{20.5.7.2}
\label{0.20.5.7.2-es}
  0\to\Omega^1_{B/A}\otimes_B C
  \xrightarrow{v_{C/B/A}}\Omega^1_{C/A}
  \xrightarrow{u_{C/B/A}}\Omega^1_{C/B}\to0
\end{equation}
es \emph{exacta y escindida}; se sabe (Bourbaki, \emph{loc. cit.},
n\textsuperscript{o}~1, prop.~1) que esto equivale a decir que para todo
$C$-módulo $L$, la sucesión
\[
  0\to\operatorname{Hom}_C(\Omega^1_{C/B},L)
  \to\operatorname{Hom}_C(\Omega^1_{C/A},L)
  \to\operatorname{Hom}_C(\Omega^1_{B/A}\otimes_B C,L)\to0
\]
es exacta; teniendo en cuenta \sref{0.20.5.6-es} y
\sref{0.20.2.4.2-es}, esta condición equivale a
\[
  \operatorname{Exalcom}_{B/A}(C,L)=0
\]
para todo $C$-módulo $L$, y la conclusión resulta por tanto de
\sref{0.19.9.8.1-es}.
\end{proof}

\oldpage[0\textsubscript{IV-1}]{132}
Observemos además que si se tiene un diagrama conmutativo de homomorfismos
de anillos
\[
\xymatrix{
  A' \ar[r] & B' \ar[r] & C' \\
  A \ar[u] \ar[r] & B \ar[u] \ar[r] & C \ar[u]
}
\]
se deduce un diagrama conmutativo
\[
\xymatrix{
  \Omega^1_{B/A}\otimes_B C \ar[r] \ar[d] &
  \Omega^1_{C/A} \ar[r] \ar[d] &
  \Omega^1_{C/B} \ar[r] \ar[d] & 0 \\
  \Omega^1_{B'/A'}\otimes_{B'} C' \ar[r] &
  \Omega^1_{C'/A'} \ar[r] &
  \Omega^1_{C'/B'} \ar[r] & 0
}
\tag{20.5.7.3}
\label{0.20.5.7.3-es}
\]
donde las flechas verticales provienen de los dihomomorfismos
\sref{0.20.5.4.3-es}.

\begin{corollary}[20.5.8]
\label{0.20.5.8-es}
Supongamos que el homomorfismo $v:B\to C$ haga de $C$ una $B$-álgebra
formalmente étale (para las topologías discretas
\sref{0.19.10.2-es}); entonces el homomorfismo \sref{0.20.5.3.3-es}
\[
  v_{C/B/A}:\Omega^1_{B/A}\otimes_B C\to\Omega^1_{C/A}
\]
es biyectivo.
\end{corollary}

\begin{proof}
En efecto, si $C$ es una $B$-álgebra formalmente no ramificada para las
topologías discretas, se deduce de \sref{0.19.10.4-es},
\sref{0.20.4.8-es} y \sref{0.20.1.1-es} que se tiene
$\operatorname{Hom}_C(\Omega^1_{C/B},L)=0$ para todo $C$-módulo $L$, por tanto
$\Omega^1_{C/B}=0$ (cf. (20.7.4)); por otra parte, si $C$ es una
$B$-álgebra formalmente lisa para las topologías discretas, la sucesión
\sref{0.20.5.7.2-es} es exacta; de donde resulta el corolario.
\end{proof}

\begin{corollary}[20.5.9]
\label{0.20.5.9-es}
Sean $A$ un anillo, $B$ una $A$-álgebra, $S$ un subconjunto multiplicativo de
$B$; entonces el homomorfismo canónico
\begin{equation}
\tag{20.5.9.1}
\label{0.20.5.9.1-es}
  S^{-1}\Omega^1_{B/A}\to\Omega^1_{S^{-1}B/A}
\end{equation}
es biyectivo.
\end{corollary}

\begin{proof}
Basta aplicar \sref{0.20.5.8-es} a $C=S^{-1}B$, que es una
$B$-álgebra formalmente étale para las topologías discretas
\sref{0.19.10.3-es}, (ii).
\end{proof}

Teniendo en cuenta \sref{0.20.5.5-es}, se puede por tanto escribir
\begin{equation}
\tag{20.5.9.2}
\label{0.20.5.9.2-es}
  \Omega^1_{S^{-1}B/S^{-1}A}
  =S^{-1}\Omega^1_{B/A}
  =\Omega^1_{S^{-1}B/A},
\end{equation}
salvo isomorfismos canónicos.

\begin{corollary}[20.5.10]
\label{0.20.5.10-es}
Si $k$ es un cuerpo y $K=k(X_\alpha)_{\alpha\in I}$ una extensión puramente trascendente de
$k$, los $dX_\alpha$ forman una base del $K$-espacio vectorial
$\Omega^1_{K/k}$.
\end{corollary}

\begin{proof}
Como $K$ es el cuerpo de fracciones del anillo de polinomios
$k[X_\alpha]_{\alpha\in I}$, esto resulta de \sref{0.20.4.13-es}, (i) y de
\sref{0.20.5.9-es}.
\end{proof}

\oldpage[0\textsubscript{IV-1}]{133}
\begin{env}[20.5.11]
\label{0.20.5.11-es}
Sean $A$ un anillo, $B$ una $A$-álgebra, $\mathfrak{R}$ un ideal de $B$,
$C$ la $A$-álgebra cociente $B/\mathfrak{R}$, y consideremos
el homomorfismo compuesto de $A$-módulos
\begin{equation}
\tag{20.5.11.1}
\label{0.20.5.11.1-es}
  \delta:\mathfrak{R}\longrightarrow B\xrightarrow{d_{B/A}}\Omega^1_{B/A}
\end{equation}
donde la primera flecha es la inyección canónica; como
$d(xy)=xdy+ydx$, se ve que se tiene
$\delta(\mathfrak{R}^2)\subset\mathfrak{R}.\Omega^1_{B/A}$, de donde, por
paso a los cocientes, un homomorfismo de $A$-módulos
\begin{equation}
\tag{20.5.11.2}
\label{0.20.5.11.2-es}
  \delta_{C/B/A}:\mathfrak{R}/\mathfrak{R}^2\longrightarrow
  \Omega^1_{B/A}\otimes_B C=
  \Omega^1_{B/A}/(\mathfrak{R}.\Omega^1_{B/A}).
\end{equation}
Pero de hecho, $\delta_{C/B/A}$ es un homomorfismo de $C$-módulos, pues
para $x\in B$ e $y\in\mathfrak{R}$, se tiene
$ydx\in\mathfrak{R}.\Omega^1_{B/A}$, por tanto
$d(xy)\equiv xdy\pmod{\mathfrak{R}.\Omega^1_{B/A}}$, lo que prueba
primero que (20.5.11.2) es un homomorfismo de $B$-módulos, y como
$\mathfrak{R}$ anula los dos miembros, esto establece nuestra afirmación.

Si $B'$ es una segunda $A$-álgebra, $u:B\to B'$ un $A$-homomorfismo,
$\mathfrak{R}'$ un ideal de $B'$ tal que
$u(\mathfrak{R})\subset\mathfrak{R}'$ y $C'=B'/\mathfrak{R}'$ el álgebra
cociente, se tiene un diagrama conmutativo
\begin{equation}
\tag{20.5.11.3}
\label{0.20.5.11.3-es}
\xymatrix{
  \mathfrak{R}/\mathfrak{R}^2 \ar[r]^-{\delta_{C/B/A}} \ar[d] &
  \Omega^1_{B/A}\otimes_B C \ar[d] \\
  \mathfrak{R}'/\mathfrak{R}'^2 \ar[r]_-{\delta_{C'/B'/A}} &
  \Omega^1_{B'/A}\otimes_{B'} C'
}
\end{equation}
donde las flechas verticales provienen de $u$ (20.5.2.4).
\end{env}

\begin{theorem}[20.5.12]
\label{0.20.5.12-es}
Sean $A$ un anillo, $B$ una $A$-álgebra, $C$ la $A$-álgebra cociente
$B/\mathfrak{R}$, $u:B\to C$ el homomorfismo canónico.
\begin{enumerate}
  \item[\rm(i)] Se tiene una sucesión exacta de $C$-módulos
\begin{equation}
\tag{20.5.12.1}
\label{0.20.5.12.1-es}
  \mathfrak{R}/\mathfrak{R}^2
  \xrightarrow{\delta_{C/B/A}}\Omega^1_{B/A}\otimes_B C
  \xrightarrow{u_{C/B/A}}\Omega^1_{C/A}\longrightarrow 0,
\end{equation}
donde $u_{C/B/A}$ y $\delta_{C/B/A}$ están definidos por (20.5.2.4) y
(20.5.11.2) respectivamente.
  \item[\rm(ii)] Si se pone $E=B/\mathfrak{R}^2$, el homomorfismo
canónico (20.5.2.4)
\[
  \Omega^1_{B/A}\otimes_B C\longrightarrow
  \Omega^1_{E/A}\otimes_E C
\]
es biyectivo.
  \item[\rm(iii)] Las tres condiciones siguientes son equivalentes:
  \begin{enumerate}
    \item[\rm a)] $\delta_{C/B/A}$ es invertible por la izquierda.
    \item[\rm b)] Toda $A$-extensión de $C$ por un $C$-módulo $L$, cuya
imagen inversa por $u:B\to C$ es $A$-trivial, es ella misma
$A$-trivial.
    \item[\rm c)] La $A$-álgebra $E=B/\mathfrak{R}^2$ es una extensión
$A$-trivial de $C$ por $\mathfrak{R}/\mathfrak{R}^2$.
  \end{enumerate}
  \item[\rm(iv)] Hay una correspondencia biunívoca canónica entre los
inversos por la izquierda de $\delta_{C/B/A}$ y los inversos por la derecha del
homomorfismo canónico $E\to C$.
\end{enumerate}
\end{theorem}

\begin{proof}
\oldpage[0\textsubscript{IV-1}]{134}
(i) Como $u$ es sobreyectivo, se tiene $\operatorname{D\acute{e}r}_B(C,L)=0$ para todo $C$-módulo
$L$ en virtud de (20.1.1). La sucesión exacta (20.2.3.1) se convierte por tanto en
\begin{equation}
\tag{20.5.12.2}
\label{0.20.5.12.2-es}
  0\longrightarrow\operatorname{D\acute{e}r}_A(C,L)\xrightarrow{v^0}\operatorname{D\acute{e}r}_A(B,L)
  \xrightarrow{\partial}\operatorname{Exalcom}_B(C,L)
  \xrightarrow{v^1}\operatorname{Exalcom}_A(C,L)
  \xrightarrow{u^1}\operatorname{Exalcom}_A(B,L),
\end{equation}
donde $v$ es el homomorfismo $A\to B$. Recordemos, por otra parte, (18.3.8) que
$\operatorname{Exalcom}_B(C,L)$ se identifica canónicamente con
$\operatorname{Hom}_C(\mathfrak{R}/\mathfrak{R}^2,L)$; se deduce, pues, de (20.5.12.2)
y (20.4.8) la sucesión exacta
\begin{equation}
\tag{20.5.12.3}
\label{0.20.5.12.3-es}
\begin{split}
  0\longrightarrow\operatorname{Hom}_C(\Omega^1_{C/A},L)
  \longrightarrow\operatorname{Hom}_C(\Omega^1_{B/A}\otimes_B C,L)
  \xrightarrow{\varphi}\operatorname{Hom}_C(\mathfrak{R}/\mathfrak{R}^2,L)\\
  \xrightarrow{\psi}\operatorname{Ker}(u^1)\longrightarrow 0
\end{split}
\end{equation}
con $\varphi=\partial\circ\eta^{-1}$ y $\psi=v^1\circ\eta$. Volviendo
a las definiciones de $\partial$ (20.2.2) y de $\eta$ (18.3.8), se comprueba
inmediatamente que $\varphi$ es precisamente el homomorfismo
$\operatorname{Hom}(\delta_{C/B/A},1_L)$. La existencia de la sucesión exacta formada por los
4 primeros términos de (20.5.12.3) muestra, pues, que la sucesión (20.5.12.1)
es exacta (Bourbaki, \emph{Álg.}, cap. II, 3\textsuperscript{a} ed.,
\S~2, n\textsuperscript{o}~1, teorema~1).

\medskip
(ii) Apliquemos a $B$ y al ideal $\mathfrak{R}^2$ la sucesión exacta
(20.5.12.1), lo que da
\begin{equation}
\tag{20.5.12.4}
\label{0.20.5.12.4-es}
  \mathfrak{R}^2/\mathfrak{R}^4\longrightarrow
  \Omega^1_{B/A}\otimes_B E\longrightarrow\Omega^1_{E/A}\longrightarrow 0,
\end{equation}
de donde, tensorizando por $C$ (considerado como $E$-álgebra), se obtiene la sucesión
exacta
\[
  \mathfrak{R}^2/\mathfrak{R}^4\otimes_E C\longrightarrow
  \Omega^1_{B/A}\otimes_B C\longrightarrow
  \Omega^1_{E/A}\otimes_E C\longrightarrow 0.
\]
Ahora bien, si $x,y$ son dos elementos de $\mathfrak{R}$, y $\zeta$ la clase
de $xy$ mód. $\mathfrak{R}^4$, la imagen
$\delta'(\zeta\otimes1)$ es, por definición,
$d_{B/A}(xy)\otimes1=(xd_{B/A}(y)+yd_{B/A}(x))\otimes1$; pero, como las
imágenes de $x$ y de $y$ en $C$ son nulas, se tiene también
$d_{B/A}(xy)\otimes1=0$ en $\Omega^1_{B/A}\otimes_B C$, lo que demuestra
la aserción.

\medskip
(iii) Decir que $\delta_{C/B/A}$ es invertible por la izquierda equivale, teniendo
en cuenta la exactitud de (20.5.12.1), a decir que la sucesión
\begin{equation}
\tag{20.5.12.5}
\label{0.20.5.12.5-es}
  0\longrightarrow\mathfrak{R}/\mathfrak{R}^2
  \xrightarrow{\delta_{C/B/A}}\Omega^1_{B/A}\otimes_B C
  \xrightarrow{u_{C/B/A}}\Omega^1_{C/A}\longrightarrow0
\end{equation}
es \emph{exacta y escindida}, y equivale también (Bourbaki,
\emph{loc. cit.}) a decir que $\operatorname{Ker}(u^1)=0$ en la sucesión exacta
(20.5.12.3) para todo $L$, lo que muestra la equivalencia de las condiciones
a) y b) (cf. (18.3.6.2)).

El hecho de que b) implique c) proviene de que la imagen inversa por
$u:B\to C$ de la $A$-extensión $E$ de $C$ por
$\mathfrak{R}/\mathfrak{R}^2$ es $A$-trivial, pues $u$ es la composición de los
homomorfismos canónicos $B\to E\to C$ (18.1.6). Recíprocamente, c)
implica b), pues toda $B$-extensión de $C$ por $L$ es $B$-equivalente
a $E\oplus_{\mathfrak{R}/\mathfrak{R}^2}L$ (18.3.8), es decir, su
clase es la imagen de la clase de $E$ (considerada como $B$-extensión)
por el homomorfismo

\oldpage[0\textsubscript{IV-1}]{135}
$w_*:\operatorname{Exan}_B(C,\mathfrak{R}/\mathfrak{R}^2)
\to\operatorname{Exan}_B(C,L)$ correspondiente a un $B$-homomorfismo
$w:\mathfrak{R}/\mathfrak{R}^2\to L$. El hecho de que c) implique b) se deduce
entonces de la conmutatividad del diagrama (18.3.6.5)
\[
\xymatrix{
  \operatorname{Exan}_B(C,\mathfrak{R}/\mathfrak{R}^2)
    \ar[r]^{w_*} \ar[d]_{u^*} &
  \operatorname{Exan}_B(C,L) \ar[d]^{u^*} \\
  \operatorname{Exan}_A(C,\mathfrak{R}/\mathfrak{R}^2)
    \ar[r]_{w_*} & \operatorname{Exan}_A(C,L).
}
\]

\medskip
(iv) Se ha visto (20.1.7) que los inversos por la derecha de $E\to C$
corresponden canónica y biyectivamente al conjunto de los elementos
$D\in\operatorname{D\acute{e}r}_A(E,\mathfrak{R}/\mathfrak{R}^2)$ tales que $D(x)=x$ en
$\mathfrak{R}/\mathfrak{R}^2$, y por tanto también, por (20.4.8), al conjunto de los
$E$-homomorfismos
$h:\Omega^1_{E/A}\to\mathfrak{R}/\mathfrak{R}^2$ tales que la composición
\[
  \mathfrak{R}/\mathfrak{R}^2\xrightarrow{d_{E/A}}
  \Omega^1_{E/A}\xrightarrow{h}\mathfrak{R}/\mathfrak{R}^2
\]
sea la identidad. Tensorizando con $C$, se deduce de ello (puesto que
$\mathfrak{R}/\mathfrak{R}^2$ es un $C$-módulo) que la composición
\[
  \mathfrak{R}/\mathfrak{R}^2\xrightarrow{\delta_{C/E/A}}
  \Omega^1_{E/A}\otimes_E C\xrightarrow{h\otimes1}
  \mathfrak{R}/\mathfrak{R}^2
\]
es la identidad; ahora bien, puesto que $\mathfrak{R}/\mathfrak{R}^2$ es un
$C$-módulo, $h\mapsto h\otimes1$ es un isomorfismo del conjunto
$\operatorname{Hom}_E(\Omega^1_{E/A},\mathfrak{R}/\mathfrak{R}^2)$ sobre
$\operatorname{Hom}_C(\Omega^1_{E/A}\otimes_E C,\mathfrak{R}/\mathfrak{R}^2)$; por
otra parte, (ii) demuestra que se pueden identificar canónicamente
$\Omega^1_{E/A}\otimes_E C$ y $\Omega^1_{B/A}\otimes_B C$,
identificándose entonces $\delta_{C/E/A}$ con $\delta_{C/B/A}$. C.Q.D.
\end{proof}

\begin{example}[20.5.13]
\label{0.20.5.13-es}
Sean $B=A[X_\alpha]_{\alpha\in I}$ un álgebra de polinomios sobre $A$,
$\mathfrak{R}$ un ideal de $B$, $(P_\lambda)$ un sistema de generadores
de $\mathfrak{R}$ y $C=B/\mathfrak{R}$; se sabe que
$\Omega^1_{B/A}$ es un $B$-módulo libre y que los $dX_\alpha$ forman una
base (20.4.13, (i)); por tanto, los $dX_\alpha$ forman también una base del
$C$-módulo libre $\Omega^1_{B/A}\otimes_B C$. Por otra parte, se deduce
inmediatamente de la definición que la imagen de
$\mathfrak{R}/\mathfrak{R}^2$ por $\delta_{C/B/A}$ es el sub-$C$-módulo
generado por los
\[
  dP_\lambda=\sum_\alpha
  \frac{\partial P_\lambda}{\partial X_\alpha}\,dX_\alpha.
\]
Se concluye que $\Omega^1_{C/A}$ es isomorfo al cociente del
$C$-módulo libre que tiene los $dX_\alpha$ por base, por el sub-$C$-módulo
generado por los $dP_\lambda$, lo que da una descripción de un módulo
de diferenciales de un álgebra \emph{cualquiera}, pues toda $A$-álgebra
$C$ puede obtenerse de la manera precedente.
\end{example}

\begin{corollary}[20.5.14]
\label{0.20.5.14-es}
Si $C$ es un $A$-álgebra formalmente lisa (para las topologías
discretas), la sucesión
\begin{equation}
\tag{20.5.14.1}
\label{0.20.5.14.1-es}
  0\longrightarrow\mathfrak{R}/\mathfrak{R}^2
  \xrightarrow{\delta_{C/B/A}}\Omega^1_{B/A}\otimes_B C
  \xrightarrow{u_{C/B/A}}\Omega^1_{C/A}\longrightarrow0
\end{equation}
es exacta y escindida.
\end{corollary}

\begin{proof}
En efecto, toda $A$-extensión de $C$ por un $C$-módulo es entonces trivial
(19.4.4.1).
\end{proof}

\oldpage[0\textsubscript{IV-1}]{136}
\begin{remark}[20.5.15]
\label{0.20.5.15-es}
Sea $u:A\to B$ un homomorfismo \emph{sobreyectivo} de anillos; entonces, para
todo homomorfismo de anillos $v:B\to C$, el homomorfismo canónico
\begin{equation}
\tag{20.5.15.1}
\label{0.20.5.15.1-es}
  u_{C/B/A}:\Omega^1_{C/A}\longrightarrow\Omega^1_{C/B}
\end{equation}
es \emph{biyectivo}; esto se deduce, en efecto, de la sucesión exacta (20.5.7.1),
puesto que $\Omega^1_{B/A}=0$ (20.4.12).
\end{remark}

\subsection{Módulos de imperfección y homomorfismos característicos}
\label{subsection:0.20.6-es}
\label{0.20.6-es}

\begin{definition}[20.6.1]
\label{0.20.6.1-es}
Dados dos homomorfismos de anillos $u:A\to B$, $v:B\to C$, se
denomina \emph{módulo de imperfección del $B$-álgebra $C$ relativo a
$A$}, y se denota $\Upsilon_{C/B/A}$, al $C$-módulo núcleo del
homomorfismo
$v_{C/B/A}:\Omega^1_{B/A}\otimes_B C\to\Omega^1_{C/A}$.
Se tiene, pues, por definición (cf. (20.5.7)) la sucesión exacta
\begin{equation}
\tag{20.6.1.1}
\label{0.20.6.1.1-es}
  0\longrightarrow\Upsilon_{C/B/A}\longrightarrow
  \Omega^1_{B/A}\otimes_B C\xrightarrow{v_{C/B/A}}
  \Omega^1_{C/A}\xrightarrow{u_{C/B/A}}\Omega^1_{C/B}
  \longrightarrow0.
\end{equation}
\end{definition}

Cuando $A=\mathbb{Z}$ (siendo entonces los módulos $\Omega^1_{B/A}$ y $\Omega^1_{C/A}$
los módulos de diferenciales «absolutas» $\Omega^1_B$ y
$\Omega^1_C$), se escribe $\Upsilon_{C/B}$ en lugar de
$\Upsilon_{C/B/\mathbb{Z}}$; cuando $B$ y $C$ son álgebras sobre un
cuerpo primo $P$, se tiene $\Upsilon_{C/B/P}=\Upsilon_{C/B}$.

Sean $R$, $S$ partes multiplicativas de $B$ y $C$, respectivamente,
tales que la imagen de $R$ esté contenida en $S$. Se deduce entonces de la
sucesión exacta (20.6.1.1) y de (20.5.9) que se tiene
\begin{equation}
\tag{20.6.1.2}
\label{0.20.6.1.2-es}
  \Upsilon_{S^{-1}C/R^{-1}B/A}=S^{-1}\Upsilon_{C/B/A}.
\end{equation}

\begin{proposition}[20.6.2]
\label{0.20.6.2-es}
Si $C$ es un $B$-álgebra formalmente lisa relativamente a $A$ (y, en
particular, si $C$ es un $B$-álgebra formalmente lisa), se tiene
$\Upsilon_{C/B/A}=0$.
\end{proposition}

\begin{proof}
Esto se deduce de (20.5.7, (ii)).
\end{proof}

\begin{proposition}[20.6.3]
\label{0.20.6.3-es}
Sean $k$ un cuerpo y $K$ una extensión de $k$. Para que $K$ sea una
extensión separable de $k$, es necesario y suficiente que $\Upsilon_{K/k}=0$
(es decir, que el homomorfismo canónico
$\Omega^1_k\otimes_k K\to\Omega^1_K$ sea inyectivo o también (20.4.8), que
toda derivación de $k$ en $K$ se prolongue a una derivación de $K$
en sí mismo).
\end{proposition}

\begin{proof}
En efecto, equivale decir que $K$ es separable sobre $k$ o que es un
$k$-álgebra formalmente lisa (19.6.1); por otra parte, si $P$ es el
cuerpo primo de $k$, $K$ es separable sobre $P$, y por tanto un $P$-álgebra
formalmente lisa, y equivale, pues, decir que $K$ es un
$k$-álgebra formalmente lisa o un $k$-álgebra formalmente lisa
relativamente a $P$. Finalmente, decir que $K$ es un $k$-álgebra formalmente
lisa relativamente a $P$ equivale, según (20.5.7, (ii)), a decir que
el homomorfismo
$v_{K/k/P}:\Omega^1_k\otimes_k K\to\Omega^1_K$ es invertible por la izquierda;
pero, como $K$ es un cuerpo, esta última condición equivale a decir
que el núcleo de $v_{K/k/P}$, es decir, $\Upsilon_{K/k}$, es nulo.
\end{proof}

\begin{env}[20.6.4]
\label{0.20.6.4-es}
Consideremos un diagrama conmutativo
\begin{equation}
\tag{20.6.4.1}
\label{0.20.6.4.1-es}
\xymatrix{
  A' \ar[r]^{u'} & B' \ar[r]^{v'} & C' \\
  A \ar[u]^f \ar[r]_u & B \ar[u]^g \ar[r]_v & C \ar[u]^h
}
\end{equation}

\oldpage[0\textsubscript{IV-1}]{137}
de homomorfismos de anillos conmutativos. La conmutatividad del diagrama
correspondiente (20.5.7.3) implica la existencia de un único
$C$-homomorfismo
\begin{equation}
\tag{20.6.4.2}
\label{0.20.6.4.2-es}
  \Upsilon_{C/B/A}\longrightarrow\Upsilon_{C'/B'/A'}
\end{equation}
deducido canónicamente de (20.6.4.1) y que hace conmutativo el diagrama
\begin{equation}
\tag{20.6.4.3}
\label{0.20.6.4.3-es}
\xymatrix{
  0 \ar[r] & \Upsilon_{C/B/A} \ar[r] \ar[d] &
  \Omega^1_{B/A}\otimes_B C \ar[r]^{v_{C/B/A}} \ar[d] &
  \Omega^1_{C/A} \ar[r]^{u_{C/B/A}} \ar[d] &
  \Omega^1_{C/B} \ar[r] \ar[d] & 0 \\
  0 \ar[r] & \Upsilon_{C'/B'/A'} \ar[r] &
  \Omega^1_{B'/A'}\otimes_{B'} C' \ar[r]_{v_{C'/B'/A'}} &
  \Omega^1_{C'/A'} \ar[r]_{u_{C'/B'/A'}} &
  \Omega^1_{C'/B'} \ar[r] & 0
}
\end{equation}

El dato del homomorfismo (20.6.4.2) equivale además al de un
$C'$-homomorfismo
\begin{equation}
\tag{20.6.4.4}
\label{0.20.6.4.4-es}
  \Upsilon_{C/B/A}\otimes_C C'\longrightarrow\Upsilon_{C'/B'/A'}
\end{equation}
que, compuesto con el homomorfismo canónico
$\Upsilon_{C/B/A}\to\Upsilon_{C/B/A}\otimes_C C'$, vuelve a dar (20.6.4.2).
Está claro que (20.6.4.2) goza de una evidente propiedad de
transitividad, lo que permite decir que $\Upsilon_{C/B/A}$ es un
\emph{funtor} de la terna $(A,B,C)$.
\end{env}

\begin{env}[20.6.5]
\label{0.20.6.5-es}
Nos resultará conveniente para lo que sigue, bajo las condiciones de (20.6.1),
introducir un \emph{complejo} (de cadenas) \emph{de $C$-módulos}
$\mathrm{K}_\bullet(C/B/A)$ cuyos términos son nulos salvo en grados 0
y 1, donde se toma
\begin{equation}
\tag{20.6.5.1}
\label{0.20.6.5.1-es}
\left\{
\begin{aligned}
  \mathrm{K}_0(C/B/A)&=\Omega^1_{C/A},\\
  \mathrm{K}_1(C/B/A)&=\Omega^1_{B/A}\otimes_B C,
\end{aligned}
\right.
\end{equation}
siendo el operador diferencial $\mathrm{K}_1\to\mathrm{K}_0$
$v_{C/B/A}$. Esto permite escribir (salvo isomorfismos canónicos)
$\Omega^1_{C/B}$ y $\Upsilon_{C/B/A}$ como los \emph{módulos
de homología} de este complejo:
\begin{equation}
\tag{20.6.5.2}
\label{0.20.6.5.2-es}
  \mathrm{H}_0(\mathrm{K}_\bullet(C/B/A))=\Omega^1_{C/B},\qquad
  \mathrm{H}_1(\mathrm{K}_\bullet(C/B/A))=\Upsilon_{C/B/A}.
\end{equation}
\end{env}

Del mismo modo:
\begin{proposition}[20.6.6]
\label{0.20.6.6-es}
Bajo las hipótesis de (20.6.1), para todo $C$-módulo $L$, se tienen
$C$-isomorfismos canónicos
\begin{align}
\tag{20.6.6.1}
\label{0.20.6.6.1-es}
  \mathrm{H}^0(\mathrm{K}_\bullet(C/B/A),L)&\simeq\operatorname{D\acute{e}r}_B(C,L),\\
\tag{20.6.6.2}
\label{0.20.6.6.2-es}
  \mathrm{H}^1(\mathrm{K}_\bullet(C/B/A),L)&\simeq
  \operatorname{Exalcom}_{B/A}(C,L).
\end{align}
\end{proposition}

\begin{proof}
En efecto, el complejo de cocadenas
$\operatorname{Hom}_C(\mathrm{K}_\bullet(C/B/A),L)$ no es otro, en virtud de
(20.4.8)
y (20.5.6), que el complejo
\[
  \cdots\longrightarrow0\longrightarrow\operatorname{D\acute{e}r}_A(C,L)
  \longrightarrow\operatorname{D\acute{e}r}_A(B,L)\longrightarrow0\longrightarrow\cdots
\]
\oldpage[0\textsubscript{IV-1}]{138}
donde el operador diferencial es $v^0$ (con las notaciones de (20.2.1)).
La proposición resulta entonces de la sucesión exacta (20.2.4.2) y de la
definición de los módulos de cohomología
\begin{equation}
\tag{20.6.6.3}
\label{0.20.6.6.3-es}
  \mathrm{H}^i(\mathrm{K}_\bullet,L)=
  \mathrm{H}^i(\operatorname{Hom}_C(\mathrm{K}_\bullet,L)).
\end{equation}
\end{proof}

Si se tiene un diagrama conmutativo de homomorfismos de anillos
\[
\xymatrix{
  A' \ar[r] & B' \ar[r] & C' \\
  A \ar[u] \ar[r] & B \ar[u] \ar[r] & C \ar[u]
}
\]
los dihomomorfismos (20.5.4.3) definen un dihomomorfismo de
complejos de módulos
\begin{equation}
\tag{20.6.6.4}
\label{0.20.6.6.4-es}
  \mathrm{K}_\bullet(C/B/A)\longrightarrow
  \mathrm{K}_\bullet(C'/B'/A')
\end{equation}
y los dihomomorfismos que de él se deducen para la homología o la
cohomología se identifican, por las fórmulas (20.6.5.2), (20.6.6.1) y
(20.6.6.2), con los dihomomorfismos ya definidos en (20.5.4.3),
(20.6.4.2), (20.2.1) y (18.3.7.2) (teniendo en cuenta, para el último, la
definición del operador $\partial$ en la sucesión exacta (20.2.4.2)).

\begin{env}[20.6.7]
\label{0.20.6.7-es}
Se sabe que, para un complejo $\mathrm{K}_\bullet$ de $C$-módulos y un
$C$-módulo $L$, se tienen homomorfismos canónicos
\[
  \alpha_i:\mathrm{H}^i(\mathrm{K}_\bullet,L)\longrightarrow
  \operatorname{Hom}_C(\mathrm{H}_i(\mathrm{K}_\bullet),L)
\]
(M, IV, 6). Aquí, el homomorfismo canónico
\[
  \alpha_1:\mathrm{H}^1(\mathrm{K}_\bullet(C/B/A),L)
  \longrightarrow
  \operatorname{Hom}_C(\mathrm{H}_1(\mathrm{K}_\bullet(C/B/A)),L)
\]
se define inmediatamente haciendo pasar al cociente por la imagen
de $\operatorname{Hom}_C(\mathrm{K}_0,L)$ el homomorfismo de restricción
\[
  \operatorname{Hom}_C(\mathrm{K}_1(C/B/A),L)\longrightarrow
  \operatorname{Hom}_C(\mathrm{H}_1(\mathrm{K}_\bullet(C/B/A)),L)
\]
puesto que $\mathrm{H}_1(\mathrm{K}_\bullet(C/B/A))$ no es otro que el
núcleo de $\mathrm{K}_1\to\mathrm{K}_0$; teniendo en cuenta (20.6.6.2) y
(20.6.5.2), se obtiene, pues, un $C$-homomorfismo canónico
\begin{equation}
\tag{20.6.7.1}
\label{0.20.6.7.1-es}
  \operatorname{Exalcom}_{B/A}(C,L)\longrightarrow
  \operatorname{Hom}_C(\Upsilon_{C/B/A},L)
\end{equation}
que se explicita de la manera siguiente: en virtud de (20.2.4.2), toda
$B$-extensión de $C$ por $L$ que es $A$-trivial procede del dato
de una $A$-derivación $D$ de $B$ en $L$, luego, por (20.4.8), de un
$C$-homomorfismo $f$ de $\Omega^1_{B/A}\otimes_B C$ en $L$; se hace
corresponder a la clase de esta extensión la \emph{restricción} de $f$
a $\Upsilon_{C/B/A}$, que solo depende, en efecto, de esta clase y
no de la elección de $D$.
\end{env}

\begin{definition}[20.6.8]
\label{0.20.6.8-es}
Sean $u:A\to B$, $v:B\to C$ dos homomorfismos de anillos. Para toda
$B$-extensión $E$ de $C$ por un $C$-módulo $L$, que sea $A$-trivial
(18.3.7), se llama \emph{homomorfismo característico de $E$}, y se
denota por $\chi_E$, al $C$-homomorfismo $\Upsilon_{C/B/A}\to L$, imagen de la
clase de $E$ por el homomorfismo canónico (20.6.7.1).
\end{definition}

\oldpage[0\textsubscript{IV-1}]{139}
Se puede definir el homomorfismo $\chi_E$ de otra manera:
\begin{proposition}[20.6.9]
\label{0.20.6.9-es}
Sea $E$ una $B$-extensión de $C$ por un $C$-módulo $L$, que sea
$A$-trivial (18.3.7); entonces el diagrama
\begin{equation}
\tag{20.6.9.1}
\label{0.20.6.9.1-es}
\xymatrix{
  0 \ar[r] & \Upsilon_{C/B/A} \ar[r] \ar[d]^{\chi_E} &
  \Omega^1_{B/A}\otimes_B C \ar[r]^{v_{C/B/A}}
    \ar[d]^{q_{E/B/A}\otimes1_C} &
  \Omega^1_{C/A} \ar[d]^{1} \\
  0 \ar[r] & L \ar[r]_{\delta_{C/E/A}} &
  \Omega^1_{E/A}\otimes_E C \ar[r]_{p_{C/E/A}} &
  \Omega^1_{C/A} \ar[r] & 0
}
\end{equation}
en el que $q:B\to E$ define la estructura de $B$-extensión de $E$ y
$p:E\to C$ es el homomorfismo de aumentación, es conmutativo; sus
filas son exactas.
\end{proposition}

\begin{proof}
La fila inferior del diagrama es la sucesión (20.5.12.5) relativa a los
dos homomorfismos $A\xrightarrow{q\circ u}E$ y
$p:E\to C=E/L$; como $p$ es sobreyectivo y $E$ es una extensión
$A$-trivial, esta sucesión es exacta y escindida en virtud de (20.5.12,
(iii)). La conmutatividad del cuadrado derecho de (20.6.9.1) resulta de la
relación $v=p\circ q$ (20.5.2.7); la imagen por
$q_{E/B/A}\otimes1_C$ del núcleo $\Upsilon_{C/B/A}$ de $v_{C/B/A}$ está
por tanto contenida en el núcleo $L$ de $p_{C/E/A}$. Por otra parte, sea
$h:C\to E$ un $A$-homomorfismo inverso por la derecha de $p$, y sea
$j:L\to E$ la inyección canónica, de modo que se tiene
$q(b)=h(v(b))+j(D(b))$ para $b\in B$, donde $D$ es la $A$-derivación de
$B$ en $L$ que define la $B$-extensión $E$; se puede escribir
$D=f\circ d_{B/A}$, donde $f:\Omega^1_{B/A}\to L$ es un
$B$-homomorfismo. En virtud de (20.5.2.6), se tiene
\[
  q_{E/B/A}(d_{B/A}(b))=d_{E/A}(q(b))
  =d_{E/A}(h(v(b)))+d_{E/A}(j(f(d_{B/A}(b))))
\]
para $b\in B$. Sea entonces
$z=\sum_i d_{B/A}(b_i)\otimes c_i$, donde $b_i\in B$ y $c_i\in C$, un
elemento de $\Omega^1_{B/A}\otimes_B C$; se tiene
\[
  (q_{E/B/A}\otimes1_C)(z)=
  \sum_i d_{E/A}(h(v(b_i)))\otimes c_i+
  \sum_i d_{E/A}(j(f(d_{B/A}(b_i))))\otimes c_i.
\]
En la primera suma, se tiene
$d_{E/A}(h(v(b_i)))=h_{E/C/A}(d_{C/A}(v(b_i)))$, luego esta suma es
$(h_{E/C/A}\otimes1_C)(v_{C/B/A}(z))$ en virtud de (20.5.2.6). Si se
toma $z\in\Upsilon_{C/B/A}$, esta suma es, por tanto, nula, y queda,
por definición de $\delta_{C/E/A}$,
\[
  (q_{E/B/A}\otimes1_C)(z)=
  \sum_i\delta_{C/E/A}(c_i f(d_{B/A}(b_i)))
  =\delta_{C/E/A}(f(z)),
\]
lo que demuestra la conmutatividad del cuadrado izquierdo de (20.6.9.1).
\end{proof}

Obsérvese que, cuando solo se supone que $\delta_{C/E/A}$ es
\emph{inyectivo} (y no necesariamente invertible por la izquierda), esta
interpretación permitiría aún definir $\chi_E$ como la
restricción de $q_{E/B/A}\otimes1_C$ a $\Upsilon_{C/B/A}$.

\begin{corollary}[20.6.10]
\label{0.20.6.10-es}
Si $\mathfrak{R}$ es un ideal de $B$, $C=B/\mathfrak{R}$,
$E=B/\mathfrak{R}^2$, y si $E$ es una $B$-extensión $A$-trivial
(18.3.7) de $C$ por $\mathfrak{R}/\mathfrak{R}^2$, entonces
el homomorfismo característico
$\chi_E:\Upsilon_{C/B/A}\to\mathfrak{R}/\mathfrak{R}^2$ es biyectivo.
\end{corollary}
\begin{proof}
\oldpage[0\textsubscript{IV-1}]{140}
En efecto, en el diagrama (20.6.9.1), las dos flechas verticales de la
derecha son homomorfismos biyectivos (20.5.12, (ii)).
\end{proof}

\begin{theorem}[20.6.11]
\label{0.20.6.11-es}
Sean $u:A\to B$, $v:B\to C$ dos homomorfismos de anillos, y $L$ un
$C$-módulo. Supongamos verificada una de las condiciones siguientes:
\begin{enumerate}
  \item[\rm(i)] $L$ es un $C$-módulo inyectivo.
  \item[\rm(ii)] $\Upsilon_{C/B/A}$ es factor directo del $C$-módulo
$\Omega^1_{B/A}\otimes_B C$ y
$u_{C/B/A}:\Omega^1_{C/A}\to\Omega^1_{C/B}$ es invertible por la derecha.
\end{enumerate}
Entonces el homomorfismo canónico (20.6.7.1)
\[
  \operatorname{Exalcom}_{B/A}(C,L)\longrightarrow
  \operatorname{Hom}_C(\Upsilon_{C/B/A},L)
\]
es biyectivo.

En particular, si $\Omega^1_{C/B}$ y $\Omega^1_{C/A}$ son
$C$-módulos proyectivos, el homomorfismo canónico (20.6.7.1) es
biyectivo.
\end{theorem}

\begin{proof}
El hecho de que cada una de las condiciones (i), (ii) implique que (20.6.7.1) es
biyectivo se sigue, en ambos casos, de la definición de $\alpha_1$. Señalemos
además que la condición (ii) es \emph{necesaria y
suficiente} para que el homomorfismo (20.6.7.1) sea biyectivo para
\emph{todo} $C$-módulo $L$ (Bourbaki, \emph{Alg.}, cap. II,
3\textsuperscript{a} ed., \S~2, n\textsuperscript{o}~1, prop.~1). Si se
supone que $\Omega^1_{C/B}$ y $\Omega^1_{C/A}$ son $C$-módulos
proyectivos, entonces, en la sucesión exacta (20.6.1.1),
$\operatorname{Ker}(u_{C/B/A})$ es un $C$-módulo proyectivo, pues la sucesión exacta
\[
  0\longrightarrow\operatorname{Ker}(u_{C/B/A})\longrightarrow\Omega^1_{C/A}
  \longrightarrow\Omega^1_{C/B}\longrightarrow0
\]
es escindida, ya que $\Omega^1_{C/B}$ es proyectivo; como
$\operatorname{Ker}(u_{C/B/A})=\operatorname{Im}(v_{C/B/A})$, la sucesión exacta
\[
  0\longrightarrow\Upsilon_{C/B/A}\longrightarrow
  \Omega^1_{B/A}\otimes_B C\longrightarrow
  \operatorname{Im}(v_{C/B/A})\longrightarrow0
\]
es escindida.
\end{proof}

\begin{corollary}[20.6.12]
\label{0.20.6.12-es}
Supongamos que $C$ es una $A$-álgebra formalmente lisa. Existe entonces
un homomorfismo canónico
\begin{equation}
\tag{20.6.12.1}
\label{0.20.6.12.1-es}
  \operatorname{Exalcom}_B(C,L)\longrightarrow
  \operatorname{Hom}_C(\Upsilon_{C/B/A},L).
\end{equation}
Además, este homomorfismo es biyectivo si se verifica una de las condiciones (i),
(ii) de (20.6.11).
\end{corollary}

\begin{proof}
En efecto, se sigue de la hipótesis sobre $C$ que
$\operatorname{Exalcom}_B(C,L)=\operatorname{Exalcom}_{B/A}(C,L)$, y
el homomorfismo (20.6.12.1) no es otro que (20.6.7.1).
\end{proof}

\begin{corollary}[20.6.13]
\label{0.20.6.13-es}
Si $C$ es una $A$-álgebra formalmente lisa y si $\Omega^1_{C/B}$ es
un $C$-módulo proyectivo, el homomorfismo (20.6.12.1) es biyectivo.
\end{corollary}

\begin{proof}
En efecto, se sabe entonces que $\Omega^1_{C/A}$ es un $C$-módulo proyectivo
(20.4.9), siendo discretas las topologías.
\end{proof}

\begin{env}[20.6.14]
\label{0.20.6.14-es}
Manteniendo las mismas notaciones, supongamos ahora además que $A$,
$B$, $C$ son $\Lambda$-álgebras y que $u$, $v$ son
$\Lambda$-homomorfismos, lo que equivale a darse tres homomorfismos
de anillos
\[
  \Lambda\xrightarrow{s}A\xrightarrow{u}B\xrightarrow{v}C.
\]

\oldpage[0\textsubscript{IV-1}]{141}
Se tienen, pues, además del módulo de imperfección $\Upsilon_{C/B/A}$, los
módulos de imperfección $\Upsilon_{B/A/\Lambda}$,
$\Upsilon_{C/A/\Lambda}$ y $\Upsilon_{C/B/\Lambda}$, y ya se han definido
homomorfismos canónicos de $C$-módulos (20.6.4.2)
\begin{align}
\tag{20.6.14.1}
\label{0.20.6.14.1-es}
  u':\Upsilon_{C/A/\Lambda}&\longrightarrow\Upsilon_{C/B/\Lambda},\\
\tag{20.6.14.2}
\label{0.20.6.14.2-es}
  s':\Upsilon_{C/B/\Lambda}&\longrightarrow\Upsilon_{C/B/A}.
\end{align}

Como en el diagrama conmutativo (20.5.7.3)
\begin{equation}
\tag{20.6.14.3}
\label{0.20.6.14.3-es}
\xymatrix{
  \Omega^1_{A/\Lambda}\otimes_A B \ar[r] \ar[d] &
  \Omega^1_{B/\Lambda} \ar[r] \ar[d] &
  \Omega^1_{B/A} \ar[r] \ar[d] & 0 \\
  \Omega^1_{A/\Lambda}\otimes_A C \ar[r] &
  \Omega^1_{C/\Lambda} \ar[r] & \Omega^1_{C/A} \ar[r] & 0
}
\end{equation}
la fila inferior está formada por $C$-módulos, se deduce por
tensorización un diagrama conmutativo
\begin{equation}
\tag{20.6.14.4}
\label{0.20.6.14.4-es}
\xymatrix{
  \Omega^1_{A/\Lambda}\otimes_A C \ar[r] \ar@{=}[d] &
  \Omega^1_{B/\Lambda}\otimes_B C \ar[r] \ar[d] &
  \Omega^1_{B/A}\otimes_B C \ar[r] \ar[d] & 0 \\
  \Omega^1_{A/\Lambda}\otimes_A C \ar[r] &
  \Omega^1_{C/\Lambda} \ar[r] & \Omega^1_{C/A} \ar[r] & 0
}
\end{equation}
donde la primera fila sigue siendo exacta y la flecha vertical de la izquierda
es la identidad; si se pone
\begin{equation}
\tag{20.6.14.5}
\label{0.20.6.14.5-es}
  \Upsilon^C_{B/A/\Lambda}=
  \operatorname{Ker}(u_{B/A/\Lambda}\otimes1_C)=
  \operatorname{Ker}(\Omega^1_{A/\Lambda}\otimes_A C\longrightarrow
       \Omega^1_{B/\Lambda}\otimes_B C)
\end{equation}
se ve, teniendo en cuenta la definición de $\Upsilon_{C/A/\Lambda}$, que
se tiene un único $C$-homomorfismo
\begin{equation}
\tag{20.6.14.6}
\label{0.20.6.14.6-es}
  v':\Upsilon^C_{B/A/\Lambda}\longrightarrow\Upsilon_{C/A/\Lambda}
\end{equation}
que hace conmutativo el diagrama
\begin{equation}
\tag{20.6.14.7}
\label{0.20.6.14.7-es}
\xymatrix{
  0 \ar[r] & \Upsilon^C_{B/A/\Lambda} \ar[r] \ar[d]^{v'} &
  \Omega^1_{A/\Lambda}\otimes_A C \ar[r] \ar@{=}[d] &
  \Omega^1_{B/\Lambda}\otimes_B C \ar[r] \ar[d] &
  \Omega^1_{B/A}\otimes_B C \ar[r] \ar[d] & 0 \\
  0 \ar[r] & \Upsilon_{C/A/\Lambda} \ar[r] &
  \Omega^1_{A/\Lambda}\otimes_A C \ar[r] &
  \Omega^1_{C/\Lambda} \ar[r] & \Omega^1_{C/A} \ar[r] & 0
}
\end{equation}
cuyas filas son exactas.

Cuando $\Lambda=\mathbb{Z}$, escribiremos $\Upsilon^C_{B/A}$ en lugar de
$\Upsilon^C_{B/A/\mathbb{Z}}$. Si $\Lambda$ es un cuerpo primo, se tiene
$\Upsilon^C_{B/A/\Lambda}=\Upsilon^C_{B/A}$.
\end{env}

\begin{env}[20.6.15]
\label{0.20.6.15-es}
Para estudiar las relaciones entre los módulos anteriores, introduciremos
por una parte el complejo de $B$-módulos
$\mathrm{K}_\bullet(B/A/\Lambda)$, y por otra, los complejos
\oldpage[0\textsubscript{IV-1}]{142}
de $C$-módulos $\mathrm{K}_\bullet(C/A/\Lambda)$ y
$\mathrm{K}_\bullet(C/B/\Lambda)$ (20.6.5), y además los complejos de
$C$-módulos siguientes. Pondremos en primer lugar
\begin{equation}
\tag{20.6.15.1}
\label{0.20.6.15.1-es}
  \mathrm{K}^C_\bullet(B/A/\Lambda)=
  \mathrm{K}_\bullet(B/A/\Lambda)\otimes_B C.
\end{equation}
Por otra parte, designaremos por $\mathrm{T}_\bullet(C/B/\Lambda)$ el
complejo de $C$-módulos cuyos términos son nulos salvo en grados 0 y 1,
y en el que
\begin{equation}
\tag{20.6.15.2}
\label{0.20.6.15.2-es}
  \mathrm{T}_0(C/B/\Lambda)=\mathrm{T}_1(C/B/\Lambda)=
  \Omega^1_{B/\Lambda}\otimes_B C
\end{equation}
siendo el operador diferencial la identidad, de modo que este complejo es
homotópico a 0; pongamos finalmente
\begin{equation}
\tag{20.6.15.3}
\label{0.20.6.15.3-es}
  \mathrm{K}'_\bullet(C/A/\Lambda)=
  \mathrm{K}_\bullet(C/A/\Lambda)\oplus
  \mathrm{T}_\bullet(C/B/\Lambda).
\end{equation}
En virtud del carácter trivial de $\mathrm{T}_\bullet$, está claro que
se tiene
\begin{equation}
\tag{20.6.15.4}
\label{0.20.6.15.4-es}
  \mathrm{H}^i(\mathrm{K}'_\bullet,L)=
  \mathrm{H}^i(\mathrm{K}_\bullet,L)
  \quad\text{y}\quad
  \mathrm{H}_i(\mathrm{K}'_\bullet,L)=
  \mathrm{H}_i(\mathrm{K}_\bullet,L)
\end{equation}
para todo $C$-módulo $L$ y todo $i$.
\end{env}
\begin{env}[20.6.16]
\label{0.20.6.16-es}
Definamos ahora una sucesión exacta de complejos, \emph{escindida} en cada
grado
\begin{equation}
\tag{20.6.16.1}
\label{0.20.6.16.1-es}
  0\to\mathrm{K}^C_\bullet(B/A/\Lambda)
  \xrightarrow{j}\mathrm{K}'_\bullet(C/A/\Lambda)
  \xrightarrow{p}\mathrm{K}_\bullet(C/B/\Lambda)\to0
\end{equation}
del modo siguiente: denotemos por un momento con
\[
  f:\Omega^1_{A/\Lambda}\otimes_A C\to
    \Omega^1_{B/\Lambda}\otimes_B C,
  \qquad
  g:\Omega^1_{B/\Lambda}\otimes_B C\to\Omega^1_{C/\Lambda}
\]
los homomorfismos canónicos $u_{B/A/\Lambda}\otimes1_C$ y
$v_{C/B/\Lambda}$ respectivamente, cuya composición es
$g\circ f=(v\circ u)_{C/A/\Lambda}$ (cf.\ \sref{0.20.6.14.4-es}).
Se tomarán $j_1(x)=(x,f(x))$, $p_1(y,z)=z-f(y)$,
$j_0(x)=(g(x),x)$, $p_0(y,z)=g(z)-y$, de modo que
$\operatorname{Im}(j_1)=\operatorname{Ker}(p_1)$ es el grafo de $f$,
complementario de $\{0\}\oplus\mathrm{T}_1$, y
$\operatorname{Im}(j_0)=\operatorname{Ker}(p_0)$ es el grafo de $g$,
complementario de $\mathrm{K}_1(C/A/\Lambda)\oplus\{0\}$; la verificación
de la conmutatividad del diagrama
\[
\xymatrix{
  0 \ar[r] & \mathrm{K}^C_1(B/A/\Lambda) \ar[r]^{j_1} \ar[d] &
  \mathrm{K}'_1(C/A/\Lambda) \ar[r]^{p_1} \ar[d] &
  \mathrm{K}_1(C/B/\Lambda) \ar[r] \ar[d] & 0 \\
  0 \ar[r] & \mathrm{K}^C_0(B/A/\Lambda) \ar[r]_{j_0} &
  \mathrm{K}'_0(C/A/\Lambda) \ar[r]_{p_0} &
  \mathrm{K}_0(C/B/\Lambda) \ar[r] & 0
}
\]
donde las flechas verticales son los operadores diferenciales, es inmediata.
\end{env}

\begin{theorem}[20.6.17]
\label{0.20.6.17-es}
Se tiene una sucesión exacta de $C$-módulos
\begin{equation}
\tag{20.6.17.1}
\label{0.20.6.17.1-es}
  0\to\Upsilon^C_{B/A/\Lambda}\xrightarrow{v'}
  \Upsilon_{C/A/\Lambda}\xrightarrow{u'}
  \Upsilon_{C/B/\Lambda}\xrightarrow{\partial}
  \Omega^1_{B/A}\otimes_B C\xrightarrow{v_{C/B/A}}
  \Omega^1_{C/A}\xrightarrow{u_{C/B/A}}\Omega^1_{C/B}\to0
\end{equation}
\oldpage[0\textsubscript{IV-1}]{143}
donde el operador de borde $\partial$ es la composición
\begin{equation}
\tag{20.6.17.2}
\label{0.20.6.17.2-es}
  \Upsilon_{C/B/\Lambda}\xrightarrow{s'}\Upsilon_{C/B/A}
  \longrightarrow\Omega^1_{B/A}\otimes_B C
\end{equation}
siendo la segunda flecha la inyección canónica.
\end{theorem}

\begin{proof}
Escribiendo la sucesión exacta de homología para la sucesión exacta de complejos
\sref{0.20.6.16.1-es}, se obtiene \sref{0.20.6.17.1-es}, pues la homología es
nula en grados distintos de 0 y de 1; el hecho de que los homomorfismos de
esta sucesión exacta que provienen por funtorialidad de $j$ y de $p$ sean
precisamente los del enunciado es inmediato. Queda por comprobar que $\partial$ es igual
a \sref{0.20.6.17.2-es}; ahora bien, un elemento $z\in\Upsilon_{C/B/\Lambda}$ es
la imagen por $p_1$ de $(0,z)$, de donde se deduce inmediatamente que $\partial(z)$ es
la imagen de $z$ por el homomorfismo canónico
\[
  s_{B/A/\Lambda}\otimes1_C:
  \Omega^1_{B/\Lambda}\otimes_B C\to\Omega^1_{B/A}\otimes_B C.
\]
La aserción se deduce de la conmutatividad del diagrama
\begin{equation}
\tag{20.6.17.3}
\label{0.20.6.17.3-es}
\xymatrix{
  \Upsilon_{C/B/\Lambda} \ar[r]^{s'} \ar[dd] & \Upsilon_{C/B/A} \ar[dd] \\
  & \\
  \Omega^1_{B/\Lambda}\otimes_B C \ar[r] &
  \Omega^1_{B/A}\otimes_B C
}
\end{equation}
(cf.\ \sref{0.20.6.4.3-es}).
\end{proof}

\begin{corollary}[20.6.18]
\label{0.20.6.18-es}
\begin{enumerate}
  \item[\rm(i)] La sucesión de $C$-módulos
\begin{equation}
\tag{20.6.18.1}
\label{0.20.6.18.1-es}
  0\to\Upsilon^C_{B/A/\Lambda}\xrightarrow{v'}
  \Upsilon_{C/A/\Lambda}\xrightarrow{u'}
  \Upsilon_{C/B/\Lambda}\xrightarrow{s'}\Upsilon_{C/B/A}\to0
\end{equation}
es exacta.
  \item[\rm(ii)] Si $B$ es un $A$-álgebra formalmente lisa relativamente
  a $\Lambda$, se tiene $\Upsilon^C_{B/A/\Lambda}=0$.
\end{enumerate}
\end{corollary}

\begin{proof}
(i) En \sref{0.20.6.17.1-es}, la imagen de $\Upsilon_{C/B/\Lambda}$ por
$\partial$ es el núcleo de $v_{C/B/A}$, es decir, $\Upsilon_{C/B/A}$ por
definición.

(ii) La hipótesis implica que la sucesión
\[
  0\to\Omega^1_{A/\Lambda}\otimes_A B\to
  \Omega^1_{B/\Lambda}\to\Omega^1_{B/A}\to0
\]
es exacta y escindida \sref{0.20.5.7-es}; por tensorización con $C$, la
sucesión
\[
  0\to\Omega^1_{A/\Lambda}\otimes_A C\to
  \Omega^1_{B/\Lambda}\otimes_B C\to
  \Omega^1_{B/A}\otimes_B C\to0
\]
permanece, pues, exacta, de donde se deduce la aserción.
\end{proof}

\begin{corollary}[20.6.19]
\label{0.20.6.19-es}
Sean $K$ un cuerpo y $E$, $F$ dos extensiones de $K$ tales que
$K\subset E\subset F$.
\begin{enumerate}
  \item[\rm(i)] Si $F$ es una extensión separable de $E$, se tiene
  $\Upsilon_{F/E/K}=0$ (es decir, el homomorfismo canónico
  $\Omega^1_{E/K}\otimes_E F\to\Omega^1_{F/K}$ es inyectivo).
  \item[\rm(ii)] Recíprocamente, si $F$ es una extensión separable de $K$ y
  si se tiene $\Upsilon_{F/E/K}=0$, entonces $F$ es una extensión separable de
  $E$.
\end{enumerate}
\end{corollary}

\begin{proof}
Si $P$ es el cuerpo primo de $K$, se tiene, en efecto, la sucesión exacta
\sref{0.20.6.18-es}
\[
  \Upsilon_{F/K/P}\to\Upsilon_{F/E/P}\to\Upsilon_{F/E/K}\to0.
\]
Si $\Upsilon_{F/E/P}=\Upsilon_{F/E}=0$, se tiene, pues,
$\Upsilon_{F/E/K}=0$, de donde (i) en virtud de \sref{0.20.6.3-es};
recíprocamente, si $\Upsilon_{F/E/K}=0$ y
$\Upsilon_{F/K/P}=\Upsilon_{F/K}=0$, se tiene
$\Upsilon_{F/E/P}=0$, de donde (ii) en virtud de \sref{0.20.6.3-es}.
\end{proof}

\oldpage[0\textsubscript{IV-1}]{144}
\begin{corollary}[20.6.20]
\label{0.20.6.20-es}
\begin{enumerate}
  \item[\rm(i)] Si $K$ es una extensión algebraica separable de $k$, se tiene
  $\Omega^1_{K/k}=0$.
  \item[\rm(ii)] Sean $k$ un cuerpo de característica 0 y $K$ una extensión
  de $k$. Para que $\Omega^1_{K/k}=0$, es necesario y suficiente que $K$ sea una
  extensión algebraica de $k$. En particular, para que un cuerpo $K$ de
  característica 0 sea tal que $\Omega^1_K=0$, es necesario y suficiente que
  $K$ sea una extensión algebraica de $\mathbb{Q}$ (cf.\ (21.4.4) y
  (21.7.5)).
\end{enumerate}
\end{corollary}

\begin{proof}
(i) Para todo $x\in K$, sea $f$ el polinomio minimal de $x$ sobre $k$; como
$f'(x)\ne0$ y $d_{K/k}(f(x))=f'(x)d_{K/k}(x)=0$, se tiene
$d_{K/k}(x)=0$ y la aserción se deduce de \sref{0.20.4.7-es}.

(ii) Existe una extensión puramente trascendente $L$ de $k$ tal que
$k\subset L\subset K$ y que $K$ es una extensión algebraica de $L$.
Como $K$ es separable sobre $L$, se deduce de \sref{0.20.6.19-es}, (i) que
la sucesión \sref{0.20.5.7.2-es}
\[
  0\to\Omega^1_{L/k}\otimes_L K\to\Omega^1_{K/k}
  \to\Omega^1_{K/L}\to0
\]
es exacta, y de (i) se deduce que $\Omega^1_{K/L}=0$. Por tanto, la relación
$\Omega^1_{K/k}=0$ equivale a $\Omega^1_{L/k}=0$, y como
$L$ es una extensión puramente trascendente de $k$, se deduce de \sref{0.20.5.10-es} que
la relación $\Omega^1_{L/k}=0$ equivale a $L=k$.
\end{proof}

\begin{remarks}[20.6.21]
\label{0.20.6.21-es}
\begin{enumerate}
  \item[\rm(i)] Como $\Upsilon_{B/A/\Lambda}$ es el núcleo de
  \[
    u_{B/A/\Lambda}:\Omega^1_{A/\Lambda}\otimes_A B
    \to\Omega^1_{B/\Lambda}
  \]
  y $\Upsilon^C_{B/A/\Lambda}$ el núcleo de
  $u_{B/A/\Lambda}\otimes1_C$, se tiene un homomorfismo canónico
\begin{equation}
\tag{20.6.21.1}
\label{0.20.6.21.1-es}
  \Upsilon_{B/A/\Lambda}\otimes_B C\to\Upsilon^C_{B/A/\Lambda}.
\end{equation}
Este homomorfismo es biyectivo cuando la sucesión
\begin{equation}
\tag{20.6.21.2}
\label{0.20.6.21.2-es}
  0\to\Upsilon_{B/A/\Lambda}\otimes_B C\to
  \Omega^1_{A/\Lambda}\otimes_A C\to
  \Omega^1_{B/\Lambda}\otimes_B C\to
  \Omega^1_{B/A}\otimes_B C\to0
\end{equation}
es exacta, lo que ocurre en los casos siguientes:
\begin{enumerate}
  \item[$1^\circ$] $C$ es un $B$-módulo \emph{plano}.
  \item[$2^\circ$] Los $B$-módulos $\Omega^1_{B/\Lambda}$ y
  $\Omega^1_{B/A}$ son \emph{planos}: en efecto, entonces ocurre lo mismo con
  $\operatorname{Ker}(\Omega^1_{B/\Lambda}\to\Omega^1_{B/A})$
  ($\mathbf{0_I}$, 6.1.2), y la sucesión \sref{0.20.6.21.2-es} es entonces
  exacta en virtud de ($\mathbf{0_I}$, 6.1.2).
\end{enumerate}
  \item[\rm(ii)] Consideremos un diagrama conmutativo de homomorfismos
  de anillos
\[
\xymatrix{
  \Lambda' \ar[r] & A' \ar[r] & B' \ar[r] & C' \\
  \Lambda \ar[u] \ar[r] & A \ar[u] \ar[r] & B \ar[u] \ar[r] & C \ar[u]
}
\]
Entonces las definiciones de \sref{0.20.6.16-es} muestran que se tiene un
diagrama conmutativo de complejos (donde las flechas verticales provienen de
\sref{0.20.6.6.4-es})
\[
\xymatrix{
  0 \ar[r] & \mathrm{K}^C_\bullet(B/A/\Lambda) \ar[r] \ar[d] &
  \mathrm{K}'_\bullet(C/A/\Lambda) \ar[r] \ar[d] &
  \mathrm{K}_\bullet(C/B/\Lambda) \ar[r] \ar[d] & 0 \\
  0 \ar[r] & \mathrm{K}^{C'}_\bullet(B'/A'/\Lambda') \ar[r] &
  \mathrm{K}'_\bullet(C'/A'/\Lambda') \ar[r] &
  \mathrm{K}_\bullet(C'/B'/\Lambda') \ar[r] & 0
}
\]
\oldpage[0\textsubscript{IV-1}]{145}
de donde se obtiene, pasando a la homología, un \emph{diagrama conmutativo}
\begin{equation}
\tag{20.6.21.3}
\label{0.20.6.21.3-es}
\xymatrix@C=0.7em{
  0 \ar[r] & \Upsilon^C_{B/A/\Lambda} \ar[r] \ar[d] &
  \Upsilon_{C/A/\Lambda} \ar[r] \ar[d] &
  \Upsilon_{C/B/\Lambda} \ar[r] \ar[d] &
  \Omega^1_{B/A}\otimes_B C \ar[r] \ar[d] &
  \Omega^1_{C/A} \ar[r] \ar[d] &
  \Omega^1_{C/B} \ar[r] \ar[d] & 0 \\
  0 \ar[r] & \Upsilon^{C'}_{B'/A'/\Lambda'} \ar[r] &
  \Upsilon_{C'/A'/\Lambda'} \ar[r] &
  \Upsilon_{C'/B'/\Lambda'} \ar[r] &
  \Omega^1_{B'/A'}\otimes_{B'} C' \ar[r] &
  \Omega^1_{C'/A'} \ar[r] & \Omega^1_{C'/B'} \ar[r] & 0
}
\end{equation}
Se tiene un diagrama conmutativo análogo para \sref{0.20.6.18.1-es}.
\end{enumerate}
\end{remarks}

\begin{proposition}[20.6.22]
\label{0.20.6.22-es}
Sean $s:\Lambda\to A$, $u:A\to B$ dos homomorfismos de anillos,
$\mathfrak{R}$ un ideal de $B$ y $C$ el anillo cociente $B/\mathfrak{R}$.
Supongamos que $E=B/\mathfrak{R}^2$ es una $B$-extensión
$\Lambda$-trivial de $C$ por $\mathfrak{R}/\mathfrak{R}^2$. Se tiene entonces una
sucesión exacta
\begin{equation}
\tag{20.6.22.1}
\label{0.20.6.22.1-es}
  0\to\Upsilon^C_{B/A/\Lambda}\xrightarrow{v'}
  \Upsilon_{C/A/\Lambda}\xrightarrow{\chi_E}
  \mathfrak{R}/\mathfrak{R}^2\xrightarrow{\delta_{C/B/A}}
  \Omega^1_{B/A}\otimes_B C\xrightarrow{v_{C/B/A}}
  \Omega^1_{C/A}\to0.
\end{equation}
\end{proposition}
 
\begin{proof}
En efecto, como $v:B\to C$ es sobreyectivo, se tiene $\Omega^1_{C/B}=0$
\sref{0.20.4.12-es}. Además, resulta de \sref{0.20.6.10-es} que
$\Upsilon_{C/B/\Lambda}$ se identifica canónicamente con
$\mathfrak{R}/\mathfrak{R}^2$. Basta entonces aplicar las sucesiones exactas
\sref{0.20.6.17.1-es} y \sref{0.20.6.18.1-es}, observando que se tiene un
diagrama conmutativo
\[
\xymatrix{
  \Upsilon_{C/A/\Lambda} \ar[r] \ar[d]_{\chi_E} &
  \Omega^1_{A/\Lambda}\otimes_A C \ar[d] \\
  \Upsilon_{C/E/\Lambda}=\Upsilon_{C/B/\Lambda}=\mathfrak{R}/\mathfrak{R}^2
    \ar[r]^-{\delta_{C/B/A}} \ar[d]_{\ell} &
  \Omega^1_{B/A}\otimes_B C=\Omega^1_{E/A}\otimes_E C \ar[d] \\
  \mathfrak{R}/\mathfrak{R}^2 \ar[r]^{\delta_{C/B/A}} &
  \Omega^1_{B/A}\otimes_B C
}
\]
y utilizando la conmutatividad del diagrama \sref{0.20.6.17.3-es}.
\end{proof}

Así se ha precisado el \emph{núcleo} de $\delta_{C/B/A}$ en el caso en que
existe un anillo $\Lambda$ tal que $A$ sea una $\Lambda$-álgebra y que
$B/\mathfrak{R}^2$ sea $\Lambda$-trivial; obsérvese que este es el caso
en particular cuando $C$ es una \emph{$\Lambda$-álgebra formalmente
lisa} (para la topología discreta).

Supongamos además que se tiene un diagrama conmutativo de homomorfismos
de anillos
\[
\xymatrix{
  \Lambda' \ar[r] & A' \ar[r] & B' \\
  \Lambda \ar[u] \ar[r] & A \ar[u] \ar[r] & B \ar[u]^f
}
\]
\oldpage[0\textsubscript{IV-1}]{146}
que $\mathfrak{R}'$ es un ideal de $B'$ tal que
$f(\mathfrak{R})\subset\mathfrak{R}'$, y que
$E'=B'/\mathfrak{R}'^2$ sea una $B'$-extensión $\Lambda'$-trivial de
$C'=B'/\mathfrak{R}'$ por $\mathfrak{R}'/\mathfrak{R}'^2$. Se tiene entonces un
diagrama conmutativo
\begin{equation}
\tag{20.6.22.2}
\label{0.20.6.22.2-es}
\xymatrix@C=0.7em{
  0 \ar[r] & \Upsilon^{C'}_{B'/A'/\Lambda'} \ar[r] &
  \Upsilon_{C'/A'/\Lambda'} \ar[r] &
  \mathfrak{R}'/\mathfrak{R}'^2 \ar[r] &
  \Omega^1_{B'/A'}\otimes_{B'} C' \ar[r] &
  \Omega^1_{C'/A'} \ar[r] & 0 \\
  0 \ar[r] & \Upsilon^C_{B/A/\Lambda} \ar[r] \ar[u] &
  \Upsilon_{C/A/\Lambda} \ar[r] \ar[u] &
  \mathfrak{R}/\mathfrak{R}^2 \ar[r] \ar[u] &
  \Omega^1_{B/A}\otimes_B C \ar[r] \ar[u] &
  \Omega^1_{C/A} \ar[r] \ar[u] & 0
}
\end{equation}
como resulta de \sref{0.20.6.21.3-es} y \sref{0.20.5.11.3-es}.

\begin{corollary}[20.6.23]
\label{0.20.6.23-es}
Bajo las hipótesis de \sref{0.20.6.22-es}, supongamos además que $B$ sea
una $A$-álgebra formalmente lisa. Se tiene entonces una sucesión exacta
\begin{equation}
\tag{20.6.23.1}
\label{0.20.6.23.1-es}
  0\to\Upsilon_{C/A/\Lambda}\xrightarrow{\chi_E}
  \mathfrak{R}/\mathfrak{R}^2\xrightarrow{\delta_{C/B/A}}
  \Omega^1_{B/A}\otimes_B C\to\Omega^1_{C/A}\to0.
\end{equation}
\end{corollary}

\begin{proof}
Esto resulta, en efecto, de \sref{0.20.6.18-es}, (ii).
\end{proof}

\begin{env}[20.6.24]
\label{0.20.6.24-es}
Cuando se satisfacen las hipótesis de \sref{0.20.6.22-es}, se dice también
que el homomorfismo característico $\chi_E$ es el
\emph{homomorfismo característico de la $A$-álgebra $B$, respecto de
$\Lambda$ y del ideal $\mathfrak{R}$} (omitiéndose estas últimas
precisiones cuando no haya riesgo de confusión); se denotará
a veces por $\chi_B$ o $\chi_{B/A}$.
\end{env}

\begin{proposition}[20.6.25]
\label{0.20.6.25-es}
Sean $s:\Lambda\to A$, $u:A\to B$, $v:B\to C$ tres homomorfismos
de anillos, $L$ un $C$-módulo. Se tiene entonces una sucesión exacta
\begin{equation}
\tag{20.6.25.1}
\label{0.20.6.25.1-es}
\begin{split}
  0\to\operatorname{D\acute{e}r}_B(C,L)\xrightarrow{u^1}
  \operatorname{D\acute{e}r}_A(C,L)\xrightarrow{v^1}
  \operatorname{D\acute{e}r}_A(B,L)\xrightarrow{\partial} \\
  \to\operatorname{Exalcom}_{B/\Lambda}(C,L)\xrightarrow{u^1}
  \operatorname{Exalcom}_{A/\Lambda}(C,L)\xrightarrow{v^1}
  \operatorname{Exalcom}_{A/\Lambda}(B,L)\to0
\end{split}
\end{equation}
donde $u^1$, $v^1$ son los homomorfismos definidos en (18.3.6.4) y
(18.3.6.2), y $\partial$ se define como en \sref{0.20.2.2-es}.
\end{proposition}

\begin{proof}
En efecto, como la sucesión exacta \sref{0.20.6.16.1-es} es escindida, se
deduce de ella una sucesión exacta
\[
  0\to\operatorname{Hom}_C(\mathrm{K}_\bullet(C/B/\Lambda),L)\to
  \operatorname{Hom}_C(\mathrm{K}'_\bullet(C/A/\Lambda),L)\to
  \operatorname{Hom}_C(\mathrm{K}^C_\bullet(B/A/\Lambda),L)\to0.
\]
Si se aplica a este complejo la sucesión exacta de cohomología, teniendo en cuenta
\sref{0.20.6.15.4-es} y \sref{0.20.6.6-es}, se obtiene
\sref{0.20.6.25.1-es}, pues se tiene
\[
  \operatorname{Hom}_C(\mathrm{K}^C_\bullet(B/A/\Lambda),L)
  =\operatorname{Hom}_B(\mathrm{K}_\bullet(B/A/\Lambda),L)
\]
por definición; la identificación de $u^1$ y $v^1$ con los homomorfismos
de (18.3.6.4) y (18.3.6.2), respectivamente, se deduce de
\sref{0.20.6.6.4-es}.
\end{proof}

\begin{remark}[20.6.26]
\label{0.20.6.26-es}
En este número, los complejos $\mathrm{K}_\bullet(C/B/A)$ han aparecido como
un artificio técnico destinado a simplificar la exposición de ciertos
comportamientos funtoriales. En realidad, estos complejos, considerados como objetos
de la categoría de los
\oldpage[0\textsubscript{IV-1}]{147}
complejos de $C$-módulos «salvo homotopía» (es decir, donde los morfismos
son las clases de homomorfismos homotópicos), son invariantes notables,
más finos que el par formado por $\Omega^1_{C/B}$ y
$\Upsilon_{C/B/A}$. Cuando $k$ es un cuerpo primo, y $C$ una $k$-álgebra
formalmente lisa (por ejemplo, un anillo regular de tipo finito sobre una
extensión de $k$ (cf.\ ($\mathbf{IV}$, 6.8.6))), se puede demostrar que el
complejo $\mathrm{K}_\bullet(C/A/k)$ puede describirse haciendo intervenir
únicamente $C$ y $A$ (con exclusión de $k$): se expresa $C$ como cociente
de un álgebra de polinomios $B$ sobre $A$ por un ideal $\mathfrak{Q}$, y se
considera el complejo $\mathrm{F}_\bullet(C/A)$ de 2 términos no nulos
\[
  \cdots\to0\to\mathfrak{Q}/\mathfrak{Q}^2\to
  \Omega^1_{B/A}\otimes_B C\to0\to\cdots
\]
(cuya homología coincide efectivamente con la de
$\mathrm{K}_\bullet(C/A/k)$ en virtud de \sref{0.20.6.23.1-es}). Estos
complejos $\mathrm{F}_\bullet(C/A)$, que desde el punto de vista del álgebra
homológica desempeñan el papel de un fibrado conormal para $\operatorname{Spec}(C)$
sobre $\operatorname{Spec}(A)$, ocuparán un lugar importante en
los capítulos de esta obra dedicados a la dualidad de haces coherentes y
al teorema de Riemann--Roch.
\end{remark}

\subsection{Generalizaciones a los anillos topológicos}
\label{subsection:0.20.7-es}

\begin{env}[20.7.1]
\label{0.20.7.1-es}
Resulta inmediatamente de las definiciones que, si en \sref{0.20.5.2-es} y
\sref{0.20.5.3-es} se supone que los anillos $A$, $B$, $C$ son anillos
topológicos y que los homomorfismos de anillos $u$, $v$ son continuos, entonces los
homomorfismos $u_{C/B/A}$ y $v_{C/B/A}$ son \emph{continuos} (naturalmente,
hay que tomar sobre $\Omega^1_{B/A}\otimes_B C$ la topología producto
tensorial).

Además, $u_{C/B/A}:\Omega^1_{C/A}\to\Omega^1_{C/B}$ es un
\emph{morfismo estricto sobreyectivo} de $C$-módulos topológicos; en efecto, lo
mismo ocurre con el homomorfismo canónico $C\otimes_A C\to C\otimes_B C$, teniendo
en cuenta la definición de la topología producto tensorial; pasando a los
cocientes se deduce que el homomorfismo correspondiente
$\mathrm{P}^1_{C/A}\to\mathrm{P}^1_{C/B}$ es un morfismo estricto sobreyectivo,
y $u_{C/B/A}$ es la restricción de este último a $\Omega^1_{C/A}$.

En \sref{0.20.5.5-es}, si $A'$ y $B$ son $A$-álgebras topológicas,
y si $B'$ y $\Omega^1_{B/A}\otimes_B B'$ están provistos de las topologías producto
tensorial, el homomorfismo canónico \sref{0.20.5.5.1-es} es un isomorfismo
topológico, teniendo en cuenta que $\mathrm{P}^1_{B/A}$ es suma directa
topológica de $B$ y de $\Omega^1_{B/A}$ \sref{0.20.4.8-es}.
\end{env}

\begin{proposition}[20.7.2]
\label{0.20.7.2-es}
Sean $u:A\to B$, $v:B\to C$ dos homomorfismos continuos de anillos
topológicos. Para que el homomorfismo continuo
$v_{C/B/A}:\Omega^1_{B/A}\otimes_B C\to\Omega^1_{C/A}$ sea formalmente
invertible por la izquierda (cf.\ \sref{0.19.1.5-es}), es necesario y suficiente que $C$
sea una $B$-álgebra formalmente lisa respecto de $A$
\sref{0.19.9.1-es} (y \emph{a fortiori} basta que $C$ sea una
$B$-álgebra formalmente lisa).
\end{proposition}

\begin{proof}
Decir que $v_{C/B/A}$ es formalmente invertible por la izquierda significa, en efecto,
puesto que las topologías de $\Omega^1_{C/A}$ y $\Omega^1_{B/A}\otimes_B C$
son menos finas que las deducidas de la topología de $C$
\sref{0.20.4.5-es}, que para todo $C$-módulo \emph{discreto} $L$, anulado por un
ideal abierto de $C$, el homomorfismo canónico
\[
  \operatorname{Hom.\ cont}_C(\Omega^1_{C/A},L)\to
  \operatorname{Hom.\ cont}_C(\Omega^1_{B/A}\otimes_B C,L)
\]
\oldpage[0\textsubscript{IV-1}]{148}
es \emph{sobreyectivo} \sref{0.19.1.5-es}; como
$\operatorname{Hom.\ cont}_C(\Omega^1_{B/A}\otimes_B C,L)=
\operatorname{Hom.\ cont}_B(\Omega^1_{B/A},L)$ por definición de la topología
producto tensorial, equivale, en virtud de \sref{0.20.4.8-es}, a decir
que el homomorfismo canónico
\[
  \operatorname{D\acute{e}r.\ cont}_A(C,L)\to
  \operatorname{D\acute{e}r.\ cont}_A(B,L)
\]
es sobreyectivo. Pero la sucesión exacta \sref{0.20.3.8.2-es} muestra que esta
condición equivale a $\operatorname{Exalcotop}_{B/A}(C,L)=0$, es decir,
precisamente a que $C$ sea formalmente lisa respecto de $A$
\sref{0.19.9.8-es}.
\end{proof}

\begin{corollary}[20.7.3]
\label{0.20.7.3-es}
Supongamos que, en $B$ y en $C$, el cuadrado de todo ideal abierto es abierto.
Para que $C$ sea una $B$-álgebra formalmente lisa respecto de $A$, es
necesario y suficiente que, si se denota por $(\mathfrak{R}_\lambda)$ un sistema
fundamental de entornos de $0$ formado por ideales de $C$, entonces, para todo
$\lambda$, el homomorfismo
\begin{equation}
\tag{20.7.3.1}
\label{0.20.7.3.1-es}
  v_{C/B/A}\otimes 1_{C/\mathfrak{R}_\lambda}:
  \Omega^1_{B/A}\otimes_B(C/\mathfrak{R}_\lambda)\to
  \Omega^1_{C/A}\otimes_C(C/\mathfrak{R}_\lambda)
\end{equation}
sea invertible por la izquierda.
\end{corollary}

\begin{proof}
Se sabe, en efecto, que la topología de $\Omega^1_{B/A}$ (resp.\
$\Omega^1_{C/A}$) se deduce de la de $B$ (resp.\ de $C$)
\sref{0.20.4.5-es}; se concluye inmediatamente que la topología de
$\Omega^1_{B/A}\otimes_B C$ se deduce también de la de $C$; el corolario
resulta entonces de \sref{0.20.7.2-es} y \sref{0.19.1.7-es}.
\end{proof}

\begin{proposition}[20.7.4]
\label{0.20.7.4-es}
Sean $A$ un anillo topológico, $B$ una $A$-álgebra topológica. Para que
$B$ sea formalmente no ramificada \sref{0.19.10.2-es}, es necesario y suficiente
que el completado separado $\widehat{\Omega}^1_{B/A}$ sea nulo.
\end{proposition}

\begin{proof}
En efecto, se deduce inmediatamente de \sref{0.19.10.4-es} y
\sref{0.20.1.1-es} que, para que $B$ sea formalmente no ramificada, es necesario y
suficiente que, para todo ideal abierto $\mathfrak{R}$ de $B$, todo ideal abierto
$\mathfrak{H}$ de $A$ tal que $\mathfrak{H}B\subset\mathfrak{R}$ y todo
$(B/\mathfrak{R})$-módulo $L$, se tenga
$\operatorname{D\acute{e}r}_{A/\mathfrak{H}}(B/\mathfrak{R},L)=0$,
es decir, $\operatorname{D\acute{e}r.\ cont}_A(B,L)=0$ para todo
$B$-módulo discreto $L$ anulado por un ideal abierto de $B$; en virtud de
\sref{0.20.4.8-es}, esto equivale a
$\operatorname{Hom.\ cont}_B(\Omega^1_{B/A},L)=0$ para tal $B$-módulo,
de donde se sigue inmediatamente la proposición.

Cuando $B$ es \emph{discreto}, la condición del enunciado de
\sref{0.20.7.4-es} equivale por tanto a $\Omega^1_{B/A}=0$.
\end{proof}

\begin{corollary}[20.7.5]
\label{0.20.7.5-es}
Sean $A$ un anillo, $\mathfrak{m}$ un ideal de $A$, $B$ una $A$-álgebra;
dotemos $A$ de la topología $\mathfrak{m}$-preádica, $B$ de la topología
deducida de la de $A$, y pongamos $A_0=A/\mathfrak{m}$,
$B_0=B/\mathfrak{m}B=B\otimes_A A_0$. Entonces, para que $B$ sea una
$A$-álgebra formalmente no ramificada, es necesario y suficiente que
$\Omega^1_{B_0/A_0}=0$ (o también que $B_0$ sea una $A_0$-álgebra
formalmente no ramificada).
\end{corollary}

\begin{proof}
En efecto, se tiene
$\Omega^1_{B/A}\otimes_B B_0=\Omega^1_{B/A}/\mathfrak{m}.\Omega^1_{B/A}
=\Omega^1_{B_0/A_0}$ en virtud de \sref{0.20.5.5-es}; escribir que todo
submódulo abierto de $\Omega^1_{B/A}$ es igual a $\Omega^1_{B/A}$ equivale
por \sref{0.20.4.5-es} a escribir que
$\mathfrak{m}.\Omega^1_{B/A}=\Omega^1_{B/A}$, de donde se obtiene la conclusión.
\end{proof}

\begin{proposition}[20.7.6]
\label{0.20.7.6-es}
Sean $u:A\to B$, $v:B\to C$ dos homomorfismos continuos de anillos
topológicos, y supongamos que $v$ haga de $C$ una $B$-álgebra formalmente
étale; entonces $v_{C/B/A}:\Omega^1_{B/A}\otimes_B C\to\Omega^1_{C/A}$ es un
bimorfismo formal \sref{0.19.1.2-es}.
\end{proposition}

\begin{proof}
En efecto, para todo $C$-módulo discreto $L$ anulado por un ideal abierto de $C$,
se tiene entonces por \sref{0.20.7.4-es}
$\operatorname{D\acute{e}r.\ cont}_B(C,L)=0$ y, por consiguiente,
\sref{0.20.3.6.1-es}, el homomorfismo canónico
\oldpage[0\textsubscript{IV-1}]{149}
$\operatorname{D\acute{e}r.\ cont}_A(C,L)\to
\operatorname{D\acute{e}r.\ cont}_A(B,L)$ es inyectivo; además, es
sobreyectivo en virtud de \sref{0.20.7.2-es}, por lo que es biyectivo; se deduce
que la imagen de $v_{C/B/A}$ es necesariamente densa en $\Omega^1_{C/A}$,
pues, de lo contrario, el cociente de $\Omega^1_{C/A}$ por la adherencia de
$\operatorname{Im}(v_{C/B/A})$ sería separado y $\neq0$ y tendría por tanto un
cociente discreto $L\neq0$, en contradicción con lo que acabamos de ver
(teniendo en cuenta \sref{0.20.4.8-es}). Por consiguiente, $v_{C/B/A}$, que es un
monomorfismo formal en virtud de \sref{0.20.7.2-es}, es también un epimorfismo formal
\sref{0.19.1.2-es}, y por tanto un bimorfismo formal.
\end{proof}

\begin{corollary}[20.7.7]
\label{0.20.7.7-es}
Supongamos que en $B$ y en $C$ el cuadrado de todo ideal abierto sea abierto.
Si $C$ es una $B$-álgebra formalmente étale, entonces, para todo ideal abierto
$\mathfrak{R}_\lambda$ de $C$, el homomorfismo \sref{0.20.7.3.1-es} es
biyectivo.
\end{corollary}

\begin{proof}
En efecto, en virtud de \sref{0.19.1.1-es}, este homomorfismo es sobreyectivo, y
es inyectivo por \sref{0.20.7.3-es}.
\end{proof}

\begin{proposition}[20.7.8]
\label{0.20.7.8-es}
Sean $u:A\to B$ un homomorfismo continuo de anillos topológicos,
$\mathfrak{R}$ un ideal de $B$, $C$ el anillo topológico cociente
$B/\mathfrak{R}$, $v:B\to C$ el homomorfismo canónico. Entonces:
\begin{enumerate}
  \item[\rm(i)] En la sucesión exacta \sref{0.20.5.12.1-es}
\[
  \mathfrak{R}/\mathfrak{R}^2\xrightarrow{\delta_{C/B/A}}
  \Omega^1_{B/A}\otimes_B C\xrightarrow{v_{C/B/A}}\Omega^1_{C/A}\to0
\]
el homomorfismo $\delta_{C/B/A}$ es continuo y el homomorfismo $v_{C/B/A}$
es un morfismo estricto de $C$-módulos topológicos.
  \item[\rm(ii)] Para que $\delta_{C/B/A}$ sea formalmente invertible por la
izquierda \sref{0.19.1.5-es}, es necesario y suficiente que, para todo $C$-módulo
discreto $L$ anulado por un ideal abierto de $C$, el homomorfismo canónico
\begin{equation}
\tag{20.7.8.1}
\label{0.20.7.8.1-es}
  \operatorname{Exalcotop}_A(C,L)\to\operatorname{Exalcotop}_A(B,L)
\end{equation}
sea inyectivo.
\end{enumerate}
\end{proposition}

\begin{proof}
(i) La primera afirmación es evidente. Por otra parte, el homomorfismo
canónico $v\otimes v:B\otimes_A B\to C\otimes_A C$ es un morfismo estricto
por definición de la topología producto tensorial, y se deduce inmediatamente
(cf.\ \sref{0.20.5.2-es}) que lo mismo ocurre con $v_{C/B/A}$.

(ii) Decir que $\delta_{C/B/A}$ es formalmente invertible por la izquierda significa
que, para todo $C$-módulo discreto $L$ anulado por un ideal abierto de $C$,
el homomorfismo canónico
\[
  \operatorname{Hom.\ cont}_C(\Omega^1_{B/A}\otimes_B C,L)\to
  \operatorname{Hom.\ cont}_C(\mathfrak{R}/\mathfrak{R}^2,L)
\]
es sobreyectivo. Ahora bien, teniendo en cuenta \sref{0.18.4.3-es} y
\sref{0.20.4.8-es}, esto equivale a decir que el homomorfismo canónico
\[
  \operatorname{D\acute{e}r.\ cont}_A(B,L)\to
  \operatorname{Exalcotop}_B(C,L)
\]
es sobreyectivo, y la conclusión se deduce entonces de la sucesión exacta
\sref{0.20.3.6.1-es}.

El hecho de que \sref{0.20.7.8.1-es} sea inyectivo se expresa aún de la manera
siguiente, teniendo en cuenta la definición de ambos miembros
\sref{0.18.4.1-es}: para todo ideal abierto $\mathfrak{M}$ de $A$, todo ideal
abierto $\mathfrak{N}$ de $B$ tal que $\mathfrak{M}B\subset\mathfrak{N}$, y
toda $(A/\mathfrak{M})$-extensión $E$ de
$B/(\mathfrak{R}+\mathfrak{N})$ por un
$(B/(\mathfrak{R}+\mathfrak{N}))$-módulo $L$, tal que la imagen recíproca de
$E$ por el homomorfismo canónico
$B/\mathfrak{N}\to B/(\mathfrak{R}+\mathfrak{N})$ sea
$(A/\mathfrak{M})$-trivial, existe un ideal abierto
$\mathfrak{M}'\subset\mathfrak{M}$ de $A$,
\oldpage[0\textsubscript{IV-1}]{150}
un ideal abierto $\mathfrak{N}'\subset\mathfrak{N}$ de $B$ tales que
$\mathfrak{M}'B\subset\mathfrak{N}'$ y que la imagen recíproca de $E$ por
el homomorfismo canónico
$B/(\mathfrak{R}+\mathfrak{N}')\to B/(\mathfrak{R}+\mathfrak{N})$ sea
$(A/\mathfrak{M}')$-trivial.
\end{proof}

\begin{corollary}[20.7.9]
\label{0.20.7.9-es}
Si la $A$-álgebra topológica $C=B/\mathfrak{R}$ es formalmente lisa,
el homomorfismo canónico $\delta_{C/B/A}$ es formalmente invertible por la
izquierda.
\end{corollary}

\begin{env}[20.7.10]
\label{0.20.7.10-es}
En \sref{0.20.6.1-es}, cuando $A$, $B$, $C$ son anillos topológicos,
$u$ y $v$ homomorfismos continuos, se dota a $\Upsilon_{C/B/A}$ de la
topología inducida por la de $\Omega^1_{B/A}\otimes_B C$; los
homomorfismos \sref{0.20.6.4.2-es} y \sref{0.20.6.4.4-es} son entonces
continuos siempre que también lo sean los del diagrama
\sref{0.20.6.4.1-es}. Además, si, en \sref{0.20.6.7-es}, se supone que
$L$ es un $C$-módulo discreto anulado por un ideal abierto de $C$, se deduce,
por paso al límite inductivo, un $C$-homomorfismo canónico
\begin{equation}
\tag{20.7.10.1}
\label{0.20.7.10.1-es}
  \operatorname{Exalcotop}_{B/A}(C,L)\to
  \operatorname{Hom.\ cont}_C(\Upsilon_{C/B/A},L)
\end{equation}
teniendo en cuenta \sref{0.18.5.3.1-es}: para todo ideal abierto
$\mathfrak{M}$ de $A$, todo ideal abierto $\mathfrak{N}$ de $B$ tal que
$\mathfrak{M}B\subset\mathfrak{N}$, todo ideal abierto $\mathfrak{P}$ de $C$
tal que $\mathfrak{N}C\subset\mathfrak{P}$ y $\mathfrak{P}.L=0$, toda
$(B/\mathfrak{N})$-extensión $E$ de $C/\mathfrak{P}$ por $L$ que sea
$(A/\mathfrak{M})$-trivial procede de los datos de un $C$-homomorfismo
continuo $f$ de $\Omega^1_{B/A}\otimes_B C$ en $L$, y el homomorfismo
\sref{0.20.7.10.1-es} hace corresponder a la imagen en
$\operatorname{Exalcotop}_{B/A}(C,L)$ de la clase de $E$, la restricción de
$f$ a $\Upsilon_{C/B/A}$, llamado \emph{homomorfismo característico de $E$}
y que también se denota $\chi_E$.
\end{env}

\begin{proposition}[20.7.11]
\label{0.20.7.11-es}
Supongamos que la topología de $C$ sea tal que el cuadrado de un ideal abierto
sea abierto. Si $C$ es una $A$-álgebra topológica formalmente lisa y si
$\Omega^1_{C/B}$ es un $C$-módulo formalmente proyectivo, se tiene un
isomorfismo canónico
\begin{equation}
\tag{20.7.11.1}
\label{0.20.7.11.1-es}
  \operatorname{Exalcotop}_B(C,L)\xrightarrow{\sim}
  \operatorname{Hom.\ cont}_C(\Upsilon_{C/B/A},L)
\end{equation}
para todo $C$-módulo discreto $L$ anulado por un ideal abierto de $C$.
\end{proposition}

\begin{proof}
En efecto, en la sucesión exacta \sref{0.20.3.7.1-es},
$\operatorname{Exalcotop}_A(C,L)=0$ \sref{0.19.4.4-es}, por lo que
$\operatorname{Exalcotop}_B(C,L)=\operatorname{Exalcotop}_{B/A}(C,L)$ y
el homomorfismo \sref{0.20.7.8.1-es} no es otro que
\sref{0.20.7.10.1-es}; el hecho de que sea biyectivo se deduce de
\sref{0.20.6.13-es} por paso al límite inductivo, teniendo en cuenta que
la topología de $\Omega^1_{C/B}$ se deduce entonces de la de $C$
\sref{0.20.4.5-es} y de \sref{0.19.2.4-es}.
\end{proof}

\begin{env}[20.7.12]
\label{0.20.7.12-es}
En \sref{0.20.6.14-es} se dota de nuevo a $\Upsilon^C_{B/A/\Lambda}$ de la
topología inducida por la de $\Omega^1_{A/\Lambda}\otimes_A C$ y entonces
el homomorfismo \sref{0.20.6.14.6-es} es continuo, cuando los anillos
considerados son topológicos y los homomorfismos de anillos son continuos.
\end{env}

\begin{env}[20.7.13]
\label{0.20.7.13-es}
Si, en \sref{0.20.6.23-es}, se supone que los anillos $\Lambda$, $A$, $B$,
$C$ son topológicos, los homomorfismos $s$, $u$, $v$ continuos y $L$ un
$C$-módulo discreto anulado por un ideal abierto de $C$, entonces se puede pasar
al límite inductivo como en \sref{0.20.3.6-es}, y se tiene una sucesión exacta
\begin{equation}
\tag{20.7.13.1}
\label{0.20.7.13.1-es}
\begin{split}
  0\to\operatorname{D\acute{e}r.\ cont}_B(C,L)\to
  \operatorname{D\acute{e}r.\ cont}_A(C,L)\to
  \operatorname{D\acute{e}r.\ cont}_A(B,L)\to \\
  \to\operatorname{Exalcotop}_{B/\Lambda}(C,L)\to
  \operatorname{Exalcotop}_{A/\Lambda}(C,L)\to
  \operatorname{Exalcotop}_{A/\Lambda}(B,L)\to0.
\end{split}
\end{equation}
\end{env}

\begin{env}[20.7.14]
\label{0.20.7.14-es}
Sean $A$ un anillo topológico, $B$ una $A$-álgebra topológica,
$\mathfrak{M}'$ un ideal abierto de $A$, $\mathfrak{N}'$ un ideal abierto de
$B$ tales que $\mathfrak{M}'B\subset\mathfrak{N}'$; pongamos
$A'=A/\mathfrak{M}'$,
\oldpage[0\textsubscript{IV-1}]{151}
$B'=B/\mathfrak{N}'$; el núcleo del homomorfismo
$B\otimes_A B\to B'\otimes_{A'}B'$ es
$\mathfrak{U}'=\operatorname{Im}(\mathfrak{N}'\otimes B+B\otimes\mathfrak{N}')$,
de donde se deduce inmediatamente que el núcleo del homomorfismo
\begin{equation}
\tag{20.7.14.1}
\label{0.20.7.14.1-es}
  \varphi_{(\mathfrak{M}',\mathfrak{N}')}:
  \Omega^1_{B/A}\to\Omega^1_{B'/A'}
\end{equation}
es $((\mathfrak{J}\cap\mathfrak{U}')+\mathfrak{J}^2)/\mathfrak{J}^2$ ; por
otra parte, el homomorfismo \sref{0.20.7.14.1-es} es sobreyectivo, como se
desprende de \sref{0.20.4.7-es}. Si $\mathfrak{M}''$ (resp.\
$\mathfrak{N}''$) es un segundo ideal abierto de $A$ (resp.\ $B$) tal que
$\mathfrak{M}''\subset\mathfrak{M}'$, $\mathfrak{N}''\subset\mathfrak{N}'$
y $\mathfrak{M}''B\subset\mathfrak{N}''$, y si se pone
$A''=A/\mathfrak{M}''$, $B''=B/\mathfrak{N}''$, se tiene de igual modo un
homomorfismo sobreyectivo
\[
  \varphi_{(\mathfrak{M}'',\mathfrak{N}''),(\mathfrak{M}',\mathfrak{N}')}:
  \Omega^1_{B''/A''}\to\Omega^1_{B'/A'}
\]
y estos homomorfismos forman evidentemente un \emph{sistema proyectivo}. Si
se observa que los
$((\mathfrak{J}\cap\mathfrak{U}')+\mathfrak{J}^2)/\mathfrak{J}^2$ forman
un sistema fundamental de entornos de $0$ en $\Omega^1_{B/A}$, se ve
entonces que el \emph{completado separado} $\widehat{\Omega}^1_{B/A}$ del
módulo $B$-topológico $\Omega^1_{B/A}$ viene dado, salvo un isomorfismo
canónico, por
\begin{equation}
\tag{20.7.14.2}
\label{0.20.7.14.2-es}
  \widehat{\Omega}^1_{B/A}=\varprojlim(\Omega^1_{B'/A'}).
\end{equation}

Además, el homomorfismo canónico
$j:\Omega^1_{B/A}\to\widehat{\Omega}^1_{B/A}$ es el límite proyectivo del
sistema proyectivo de los $\varphi_{(\mathfrak{M}',\mathfrak{N}')}$, por lo que
$\varphi_{(\mathfrak{M}',\mathfrak{N}')}$ se factoriza como
$\Omega^1_{B/A}\xrightarrow{j}\widehat{\Omega}^1_{B/A}\to
\Omega^1_{B'/A'}$, y como es sobreyectivo se concluye que el homomorfismo
canónico
\begin{equation}
\tag{20.7.14.3}
\label{0.20.7.14.3-es}
  \widehat{\Omega}^1_{B/A}\to\Omega^1_{B'/A'}
\end{equation}
es sobreyectivo para todo par $(\mathfrak{M}',\mathfrak{N}')$.

Además, en lo anterior se puede reemplazar en todas partes $A'$ por $A$.

Finalmente, si $L$ es un $B$-módulo topológico \emph{separado y completo}, todo
$B$-homomorfismo continuo de $\Omega^1_{B/A}$ en $L$ se prolonga de forma
única a un $\widehat B$-homomorfismo continuo de
$\widehat{\Omega}^1_{B/A}$ en $L$, y recíprocamente un homomorfismo de este
tipo vuelve a dar por composición con $\Omega^1_{B/A}\to\widehat{\Omega}^1_{B/A}$
un $B$-homomorfismo continuo, por lo que se tiene un isomorfismo canónico
\[
  \operatorname{Hom.\ cont}_B(\Omega^1_{B/A},L)\xrightarrow{\sim}
  \operatorname{Hom.\ cont}_{\widehat B}(\widehat{\Omega}^1_{B/A},L).
\]
En particular, si $L$ es un $B$-módulo discreto anulado por un ideal
abierto de $B$, se ve que el isomorfismo canónico
\sref{0.20.4.8.2-es} también puede escribirse
\begin{equation}
\tag{20.7.14.4}
\label{0.20.7.14.4-es}
  \operatorname{Hom.\ cont}_{\widehat B}(\widehat{\Omega}^1_{B/A},L)
  \xrightarrow{\sim}\operatorname{D\acute{e}r.\ cont}_A(B,L).
\end{equation}
\end{env}

\begin{proposition}[20.7.15]
\label{0.20.7.15-es}
Sean $A$ un anillo discreto, $B$ un $A$-álgebra topológica adica
($\mathbf{0}_{\mathrm I}$, 7.1.9), $\mathfrak{n}$ un ideal de definición de
$B$, $B_0=B/\mathfrak{n}$. Se supone que $\Omega^1_{B_0/A}$ y
$\mathfrak{n}/\mathfrak{n}^2$ son módulos de tipo finito sobre $B_0$. Entonces
$\widehat{\Omega}^1_{B/A}$ es un $\widehat B$-módulo de tipo finito.
\end{proposition}

\begin{proof}
Como el cuadrado de todo ideal abierto de $B$ es abierto, la topología de
$\Omega^1_{B/A}$ es la topología $\mathfrak{n}$-preádica
\sref{0.20.4.5-es}. Teniendo en cuenta la hipótesis de que $B$ es un anillo adico,
basta, por tanto, en virtud de ($\mathbf{0}_{\mathrm I}$, 7.2.7 y 7.2.9), demostrar
\oldpage[0\textsubscript{IV-1}]{152}
que $\Omega^1_{B/A}/\mathfrak{n}.\Omega^1_{B/A}=
\Omega^1_{B/A}\otimes_B B_0$ es un $B_0$-módulo de tipo finito. Pero esto
se deduce de la hipótesis y de la sucesión exacta \sref{0.20.5.12.1-es}
\[
  \mathfrak{n}/\mathfrak{n}^2\to
  \Omega^1_{B/A}\otimes_B B_0\to\Omega^1_{B_0/A}\to0.
\]
\end{proof}

\begin{env}[20.7.16]
\label{0.20.7.16-es}
La proposición \sref{0.20.7.15-es} se aplica, por ejemplo, cuando $A$ es un
cuerpo $k$, $B=k'[[T_1,\ldots,T_n]]$ es un álgebra de series formales provista de
su topología usual, y $k'$ es una extensión \emph{finita} de $k$ (cf.\
\sref{0.21.9.2-es}). Se observará que el cuerpo de fracciones $K$ de $B$ tiene
grado de trascendencia infinito sobre $k$; cuando $k$ tiene característica
$0$, se deduce inmediatamente (con ayuda, en particular, de
\sref{0.20.4.13-es}, (i), y de \sref{0.20.5.9-es}, usando también el hecho de
que toda derivación de un cuerpo de característica $0$ se prolonga a toda
extensión) que $\Omega^1_{B/k}$ \emph{no es un $B$-módulo de tipo finito}.
\end{env}

\begin{env}[20.7.17]
\label{0.20.7.17-es}
Sean $A$, $B$, $C$ tres anillos topológicos, y sean $u:A\to B$, $v:B\to C$
dos homomorfismos continuos; reemplazando $A$, $B$, $C$ por cocientes por
ideales abiertos $A'=A/\mathfrak{M}'$, $B'=B/\mathfrak{N}'$,
$C'=C/\mathfrak{P}'$, con $u(\mathfrak{M}')\subset\mathfrak{N}'$,
$v(\mathfrak{N}')\subset\mathfrak{P}'$, de modo que se tienen homomorfismos
$u':A'\to B'$, $v':B'\to C'$, se obtienen homomorfismos canónicos
$u'_{C'/B'/A'}$, $v'_{C'/B'/A'}$ que, en virtud de
\sref{0.20.5.4-es}, forman sistemas proyectivos y dan por consiguiente,
pasando al límite, homomorfismos canónicos, prolongaciones a los completados
separados de los homomorfismos de la sucesión exacta
\sref{0.20.5.7.1-es}
\[
  \Omega^1_{B/A}\otimes_B C\xrightarrow{v_{C/B/A}}\Omega^1_{C/A}
  \xrightarrow{u_{C/B/A}}\Omega^1_{C/B}\to0:
\]
\begin{align}
  \widehat v_{C/B/A}:\widehat{\Omega}^1_{B/A}
  \widehat\otimes_{\widehat B}\widehat C&\to\widehat{\Omega}^1_{C/A},
  \tag{20.7.17.1}\label{0.20.7.17.1-es}\\
  \widehat u_{C/B/A}:\widehat{\Omega}^1_{C/A}&\to
  \widehat{\Omega}^1_{C/B}. \tag{20.7.17.2}\label{0.20.7.17.2-es}
\end{align}
y en la sucesión
\begin{equation}
\tag{20.7.17.3}
\label{0.20.7.17.3-es}
  \widehat{\Omega}^1_{B/A}\widehat\otimes_{\widehat B}\widehat C
  \xrightarrow{\widehat v_{C/B/A}}\widehat{\Omega}^1_{C/A}
  \xrightarrow{\widehat u_{C/B/A}}\widehat{\Omega}^1_{C/B}\to0
\end{equation}
la composición de dos homomorfismos consecutivos es $0$, pero la sucesión
\emph{no es necesariamente exacta}. Sin embargo, si $B$ y $C$ son
\emph{metrizable}, el homomorfismo $\widehat u_{C/B/A}$ es \emph{sobreyectivo},
y $\operatorname{Im}(\widehat v_{C/B/A})$ es \emph{densa} en
$\operatorname{Ker}(\widehat u_{C/B/A})$: esto se deduce inmediatamente
(cf.\ ($\mathbf{0}_{\mathrm I}$, 7.3.1)) del hecho de que, si
$(\mathfrak{M}_k)$ (resp.\ $(\mathfrak{N}_k)$) es una sucesión decreciente
de ideales de $B$ (resp.\ $C$), que forma un sistema fundamental de entornos
de $0$, y si se ponen $B_k=B/\mathfrak{M}_k$, $C_k=C/\mathfrak{N}_k$, los
homomorfismos de transición
$\Omega^1_{C_{k+1}/B_{k+1}}\to\Omega^1_{C_k/B_k}$,
$\Omega^1_{C_{k+1}/A}\to\Omega^1_{C_k/A}$ y
$\Omega^1_{B_{k+1}/A}\to\Omega^1_{B_k/A}$ son sobreyectivos
\sref{0.20.7.14-es}.
\end{env}

\begin{proposition}[20.7.18]
\label{0.20.7.18-es}
Sean $A$, $B$, $C$ tres anillos topológicos, y sean $u:A\to B$, $v:B\to C$
dos homomorfismos continuos. Se supone que $B$ y $C$ son admisibles
($\mathbf{0}_{\mathrm I}$, 7.1.2) y metrizables. Para que el homomorfismo
canónico
\[
  \widehat v_{C/B/A}:\widehat{\Omega}^1_{B/A}
  \widehat\otimes_{\widehat B}\widehat C\to\widehat{\Omega}^1_{C/A}
\]
admita una inversa por la izquierda que sea un $C$-homomorfismo continuo, es
necesario y suficiente que $C$ sea una $B$-álgebra formalmente lisa respecto
de $A$.
\end{proposition}

\begin{proof}
La condición es necesaria en virtud de \sref{0.20.7.2-es} y
\sref{0.19.1.6-es}. Para ver que es suficiente, obsérvese que el $C$-módulo
topológico
$L=\widehat{\Omega}^1_{B/A}\widehat\otimes_{\widehat B}\widehat C$ es
metrizable y
\oldpage[0\textsubscript{IV-1}]{153}
completo en virtud de la hipótesis, por lo que
$E=D_C(L)$, provisto de la topología producto, es metrizable y completo; es
además admisible, pues si $\mathfrak{R}$ es un ideal de definición en $C$, la
sucesión $(\mathfrak{R}\oplus L)^n=\mathfrak{R}^n\oplus\mathfrak{R}^{n-1}L$
tiende a $0$. Como la aplicación compuesta
\[
  D:B\xrightarrow{d_{B/A}}\Omega^1_{B/A}\to
  \widehat{\Omega}^1_{B/A}\widehat\otimes_{\widehat B}\widehat C=L
\]
es una $A$-derivación continua de $B$ en $L$, el $A$-homomorfismo continuo
\[
  f:x\longmapsto(v(x),D(x))
\]
de $B$ en $E$ define en $E$ una estructura de extensión de $B$. Como $L$ es
un ideal cerrado en $E$, se deduce de \sref{0.19.9.5-es} y de la hipótesis que
la aplicación identidad $C\to E/L$ (que es un $B$-homomorfismo) se factoriza
como $C\xrightarrow{g}E\to E/L$, donde $g$ es un homomorfismo continuo tal
que $g\circ v=f$; por consiguiente, según \sref{0.20.1.3-es}, $g$ es de la
forma $y\mapsto(y,D'(y))$, donde $D'$ es una $B$-derivación continua de $C$ en
$L$, es decir, $D'\circ v=D$. Teniendo en cuenta \sref{0.20.7.14.4-es}, se
tiene $D'=h\circ\widehat d_{C/A}$, donde
$h:\widehat{\Omega}^1_{C/A}\to L$ es un $C$-homomorfismo continuo; pero por
definición, $\widehat d_{C/A}\circ v=\widehat v_{C/B/A}\circ D$, y como la
imagen de $B$ por $D$ genera (topológicamente) el $C$-módulo $L$
\sref{0.20.4.7-es}, la relación $D'\circ v=D$ da precisamente
$h\circ\widehat v_{C/B/A}=1_L$.
\end{proof}

\begin{corollary}[20.7.19]
\label{0.20.7.19-es}
Bajo las hipótesis de \sref{0.20.7.18-es}, supongamos además que en $B$ y
$C$ el cuadrado de todo ideal abierto sea abierto; entonces, para que $C$ sea
una $B$-álgebra formalmente lisa respecto de $A$, es necesario y suficiente
que $\widehat v_{C/B/A}$ sea invertible por la izquierda.
\end{corollary}

\begin{proof}
En efecto, las topologías de
$\widehat{\Omega}^1_{B/A}\widehat\otimes_{\widehat B}\widehat C$ y de
$\widehat{\Omega}^1_{C/A}$ se deducen entonces de la de $C$
\sref{0.20.4.5-es}, y todo $C$-homomorfismo de uno en el otro es
necesariamente continuo.
\end{proof}

\begin{env}[20.7.20]
\label{0.20.7.20-es}
Sean $A$ un anillo topológico, $B$ una $A$-álgebra topológica
\emph{metrizable y completa}, $\mathfrak{R}$ un ideal \emph{cerrado} de $B$,
$C=B/\mathfrak{R}$ el anillo topológico cociente, que es metrizable y
completo. Sea $(\mathfrak{M}_k)$ un sistema fundamental decreciente de
entornos de $0$ en $B$ formado por ideales, y pongamos
$B_k=B/\mathfrak{M}_k$,
$\mathfrak{R}_k=(\mathfrak{R}+\mathfrak{M}_k)/\mathfrak{M}_k$,
$C_k=B_k/\mathfrak{R}_k$. Se tiene un sistema proyectivo de homomorfismos
\[
  \delta_{C_k/B_k/A}:\mathfrak{R}_k/\mathfrak{R}_k^2\to
  \Omega^1_{B_k/A}\otimes_{B_k}C_k
\]
\sref{0.20.5.11.3-es}, de donde se deduce, pasando al límite, un
homomorfismo canónico
\[
  \widehat\delta_{C/B/A}:\mathfrak{R}/\mathfrak{R}^2\to
  \widehat{\Omega}^1_{B/A}\widehat\otimes_{\widehat B}\widehat C
\]
y, razonando como en \sref{0.20.7.17-es}, se ve que el homomorfismo canónico
$\widehat u_{C/B/A}:\widehat{\Omega}^1_{B/A}
\widehat\otimes_{\widehat B}\widehat C\to\widehat{\Omega}^1_{C/A}$ es
\emph{sobreyectivo} y que $\operatorname{Im}(\widehat\delta_{C/B/A})$ es
\emph{densa} en $\operatorname{Ker}(\widehat u_{C/B/A})$.
\end{env}

\section{Diferenciales en los anillos de característica p}
Los resultados del presente párrafo, de naturaleza más especial y técnica que
los de los §§19, 20 y 22, solo se utilizarán excepcionalmente en el curso del
cap. IV. Su función principal aquí reside en la demostración de tres teoremas del §22
(22.3.3, 22.5.8 y 22.7.3), de los cuales el primero y el último intervienen
esencialmente en la teoría «fina» de los anillos locales noetherianos del
cap. IV, §7.
\oldpage[0\textsubscript{IV-1}]{154}
\subsection{Sistemas de $p$-generadores y $p$-bases}
\label{0.21.1-es}

\begin{env}[21.1.1]
\label{0.21.1.1-es}
Dado un número $p$ que sea o bien $0$ o bien un número primo, diremos que un
anillo $A$ es de característica $p$ si existe un homomorfismo de anillos
$P\to A$, donde $P$ es el cuerpo primo de característica $p$; observemos que
este homomorfismo es entonces único, pues la composición
$\mathbf Z\to P\to A$ es el único homomorfismo de $\mathbf Z$ en $A$.
Si $A\neq0$, equivale a decir que $A$ contiene un cuerpo de característica
$p$, ya que la imagen de $P$ es necesariamente un cuerpo isomorfo a $P$ (y,
por lo demás, el único subcuerpo de $A$ isomorfo a $P$).
\end{env}

\begin{env}[21.1.2]
\label{0.21.1.2-es}
Si $p>0$, decir que $A$ es de característica $p$ equivale a decir que en
$A$ se tiene $p\cdot1=0$, o también $pA=0$. Si $p=0$, decir que $A$ es de
característica $p$ equivale a decir que para todo entero $n\neq0$, $n\cdot1$
es invertible en $A$. Si $A\neq0$, no puede existir más que un único $p$
(primo o $0$) tal que $A$ sea de característica $p$; esto se deduce de la
identidad de Bézout $ap+bq=1$ para dos números primos distintos $p,q$.
En cambio, el anillo reducido a $0$ es de característica $p$ para todo $p$.
\end{env}

\begin{env}[21.1.3]
\label{0.21.1.3-es}
Si $A$ es de característica $p$, lo mismo sucede con toda álgebra sobre $A$.
En particular, para todo ideal primo $\mathfrak p$ de $A$, el cuerpo residual
de $A$ en $\mathfrak p$ es de característica $p$. Recíprocamente, si $p=0$
y, para todo ideal maximal $\mathfrak m$ de $A$, el cuerpo residual de $A$ en
$\mathfrak m$ es de característica $0$, lo mismo sucede con $A$, pues para
todo entero $n\neq0$, $n\cdot1$ es entonces invertible en todos los
$A_{\mathfrak m}$ y, por tanto, también en $A$. En cambio, si $p>0$, un
anillo local puede tener un cuerpo residual de característica $p$ sin ser él
mismo de característica $p$, como muestra el ejemplo del anillo primo
$\mathbf Z_{p\mathbf Z}$ (19.8.3) o del anillo local artiniano
$\mathbf Z/p^n\mathbf Z$ para $n\geq2$, que no contienen cuerpos.

Observemos finalmente que, para un anillo (incluso reducido), los cuerpos
residuales en sus ideales primos pueden tener características distintas, como
muestra el ejemplo de $\mathbf Z$.
\end{env}

\begin{env}[21.1.4]
\label{0.21.1.4-es}
En todo lo que sigue suponemos fijado un número primo $p$ y todos los anillos
se supondrán de característica $p$, salvo mención expresa en contrario. Para
un anillo $A$ de este tipo, la aplicación $x\mapsto x^p$ es un endomorfismo de
$A$, que denotamos por $F_A$; si $A$ es reducido, $F_A$ es inyectivo. Ponemos
$A^p=F_A(A)$ (subanillo de $A$ formado por los $x^p$, $x\in A$); podemos
considerar naturalmente $A$ como una $A^p$-álgebra.

También podemos considerar $A$ como una $A$-álgebra mediante el homomorfismo
$F_A:x\mapsto x^p$ de $A$ en $A$; en otras palabras, se trata del álgebra
$A^{(p)}$ para la cual el producto $\lambda\cdot x$ de $x\in A$ por un
escalar $\lambda\in A$ es el producto $\lambda^p x$ en el anillo $A$. Es
claro que, para todo homomorfismo de anillos $u:A\to B$, el par $(u,u)$ es un
dihomomorfismo de álgebras $A^{(p)}\to B^{(p)}$.
Para todo $A$-módulo $E$, ponemos $E^{(p)}=E\otimes_A A^{(p)}$, donde
$A^{(p)}$ se considera como un $(A,A)$-bimódulo, siendo la estructura de
$A$-módulo por la izquierda la que acabamos de definir y la estructura de
$A$-módulo por la derecha la definida por la multiplicación en $A$; así,
para $x\in E$, $\alpha,\beta\in A$, se tiene
\[
  \alpha(x\otimes\beta)=x\otimes(\alpha\beta),\qquad
  (\alpha x)\otimes\beta=x\otimes\alpha^p\beta.
\]
\end{env}

\oldpage[0\textsubscript{IV-1}]{155}
Poniendo $x^{(p)}=x\otimes1$, se tiene, por tanto,
$(\alpha x)^{(p)}=\alpha^p x^{(p)}$. Cuando $F_A$ es inyectivo, podemos
identificar $A^{(p)}$ con el anillo $A$ considerado como álgebra sobre su
subanillo $A^p$.

\begin{proposition}[21.1.5]
\label{0.21.1.5-es}
Sean $A$, $B$ dos anillos y $u:A\to B$ un homomorfismo.
\begin{enumerate}
  \item[(i)] Toda $A$-derivación de $B$ en un $B$-módulo $L$ se anula en
  $A[B^p]$ (por tanto es una $A[B^p]$-derivación).
  \item[(ii)] Para toda sub-$A$-álgebra $A'$ de $A[B^p]$, si
  $j:A'\to B$ es la inyección canónica, el homomorfismo canónico
  \[
  j_{B/A'/A}:\Omega^1_{B/A}\longrightarrow\Omega^1_{B/A'}
  \]
  es biyectivo.
  \item[(iii)] Supongamos que existe un entero $s\geq0$ tal que
  $u(A)\supset B^{p^s}$. Entonces, en el anillo $B\otimes_A B$,
  $\mathfrak J_{B/A}$ es un ideal nil.
\end{enumerate}
\end{proposition}

\begin{proof}
(i) Por inducción a partir de (20.1.1.1), se deduce, para toda
$A$-derivación $D:B\to L$, que $D(x^n)=nx^{n-1}D(x)$, y en particular
$D(x^p)=0$.

(ii) En vista de (20.4.8), la afirmación (ii) no es más que una traducción
de (i), que se escribe
$\operatorname{Der}_{A[B^p]}(B,L)=\operatorname{Der}_A(B,L)$ para todo
$B$-módulo $L$.

(iii) Para todo $x\in B$, se tiene
$(x\otimes1-1\otimes x)^{p^s}=x^{p^s}\otimes1-1\otimes x^{p^s}=0$,
puesto que $x^{p^s}\in u(A)$. La conclusión resulta de que los elementos
$1\otimes x-x\otimes1$ generan $\mathfrak J_{B/A}$ (20.4.4).
\end{proof}

Se deduce inmediatamente de (21.1.5) que, para todo par de homomorfismos de
anillos $A\to B\to C$, se tiene
\begin{equation}
\tag{21.1.5.1}
\label{0.21.1.5.1-es}
\Upsilon_{C/B/A}=\Upsilon_{C/B/A'}
\end{equation}
para toda sub-$A$-álgebra $A'\subset A[B^p]$ de $B$.

Por otra parte, (21.1.5) muestra en particular que, para los módulos de
diferenciales «absolutas»,
\begin{equation}
\tag{21.1.5.2}
\label{0.21.1.5.2-es}
\Omega^1_A=\Omega^1_{A/A^p}.
\end{equation}

\begin{corollary}[21.1.6]
\label{0.21.1.6-es}
Supongamos que $B$ es una $A$-álgebra de tipo finito y que existe un entero
$s\geq0$ tal que $u(A)\supset B^{p^s}$. Si $\Omega^1_{B/A}=0$, $u$ es
sobreyectivo.
\end{corollary}

\begin{proof}
\oldpage[0\textsubscript{IV-1}]{156}
Según (21.1.5, (iii)), $\mathfrak J_{B/A}$ es un ideal nil; además, en
virtud de (20.4.4) y de la hipótesis, $\mathfrak J_{B/A}$ es un ideal de tipo
finito y, por tanto, es nilpotente; como la relación $\Omega^1_{B/A}=0$
significa que $\mathfrak J_{B/A}=\mathfrak J_{B/A}^2$, se concluye que
$\mathfrak J_{B/A}=0$, o, equivalentemente, que
$\pi:B\otimes_A B\to B$ es biyectivo.
Por otra parte, todo elemento de $B$ cuya potencia $p^s$-ésima pertenece a
$u(A)$ es entero; por tanto, $B$ es entero sobre $A$ y, al ser de tipo finito,
es una $A$-álgebra finita. Se reduce así la cuestión a demostrar el lema
siguiente (en el que no
se supone que los anillos considerados sean de característica $p$):
\end{proof}

\begin{lemma}[21.1.6.1]
\label{0.21.1.6.1-es}
Sean $R$ un anillo y $S$ una $R$-álgebra finita; si el homomorfismo canónico
$\pi:S\otimes_R S\to S$ es biyectivo, entonces el homomorfismo estructural
$u:R\to S$ es sobreyectivo.
\end{lemma}

\begin{proof}
Basta demostrar que, para todo ideal maximal $\mathfrak m$ de $R$, si se
pone $T=R-\mathfrak m$, el homomorfismo
$u_{\mathfrak m}:R_{\mathfrak m}\to T^{-1}S$ es sobreyectivo (Bourbaki,
\emph{Alg. comm.}, cap.~II, §~3, n$^\circ$~3, teorema~1); ahora bien, la
hipótesis implica que el homomorfismo
$(T^{-1}S)\otimes_{R_{\mathfrak m}}(S/\mathfrak mS)\to T^{-1}S$
es biyectivo (0$_I$, 1.3.4), y como $T^{-1}S$ es una
$R_{\mathfrak m}$-álgebra finita, se ve que podemos limitarnos al caso en
que $R$ es un anillo local.

Denotando todavía por $\mathfrak m$ su ideal maximal, basta demostrar que
$u\otimes1:R/\mathfrak m\to S/\mathfrak mS$ es sobreyectivo, en virtud del
lema de Nakayama ($S$ es un $R$-módulo de tipo finito); como el homomorfismo
canónico
$(S/\mathfrak mS)\otimes_{R/\mathfrak m}(S/\mathfrak mS)\to S/\mathfrak mS$
es biyectivo y $S/\mathfrak mS$ es una $(R/\mathfrak m)$-álgebra finita,
finalmente nos reducimos al caso en que $R$ es un cuerpo; pero entonces los
rangos de $S\otimes_R S$ y de $S$ sobre $R$ solo pueden ser iguales si $S$
es de rango $0$ o $1$, es decir, si $u:R\to S$ es sobreyectivo.
\end{proof}

\begin{proposition}[21.1.7]
\label{0.21.1.7-es}
Sean $A$ un anillo, $B$ una $A$-álgebra, $(x_\alpha)_{\alpha\in I}$ una
familia de elementos de $B$. Consideremos las propiedades siguientes:
\begin{enumerate}
  \item[a)] $B=A[B^p,(x_\alpha)]$, es decir, la $A[B^p]$-álgebra $B$ está
  generada por la familia $(x_\alpha)_{\alpha\in I}$.
  \item[b)] El $A[B^p]$-módulo $B$ está generado por los monomios
  $\prod_\alpha x_\alpha^{n(\alpha)}$, donde $(n(\alpha))_{\alpha\in I}$ es
  una familia de enteros de soporte finito tal que $0\leq n(\alpha)<p$ para
  todo $\alpha\in I$.
  \item[c)] El $B$-módulo $\Omega^1_{B/A}$ está generado por los
  $d_{B/A}(x_\alpha)$ $(\alpha\in I)$.
\end{enumerate}
Entonces las propiedades a) y b) son equivalentes y entrañan c); además, si
$B$ es una $A[B^p]$-álgebra de tipo finito, c) es equivalente a a) y b).
\end{proposition}

\begin{proof}
Está claro que b) implica a), y, recíprocamente, a) implica b), pues todo
monomio $\prod_\alpha x_\alpha^{m(\alpha)}$, donde los $m(\alpha)$ son
enteros $\geq0$ (formando una familia de soporte finito), puede escribirse
$(\prod_\alpha x_\alpha^{q(\alpha)})^p\prod_\alpha x_\alpha^{r(\alpha)}$,
dividiendo cada $m(\alpha)$ por $p$, lo que da
$m(\alpha)=p q(\alpha)+r(\alpha)$ con $0\leq r(\alpha)<p$.
El hecho de que a) entrañe c) se deduce de (21.1.5, (ii)) y de (20.4.7).

Supongamos recíprocamente que c) se verifica y que $B$ es una
$A[B^p]$-álgebra de tipo finito; sea $B'$ la sub-$A[B^p]$-álgebra de $B$
generada por los $x_\alpha$; en la sucesión exacta (20.5.7.1)
\[
\Omega^1_{B'/A[B^p]}\otimes_{B'}B\longrightarrow
\Omega^1_{B/A[B^p]}\longrightarrow\Omega^1_{B/B'}\longrightarrow0
\]
la hipótesis entraña que la flecha de la izquierda es sobreyectiva (teniendo
en cuenta (21.1.5, (ii))); por tanto, $\Omega^1_{B/B'}=0$, y como
$B^p\subset B'\subset B$ y $B$ es una $B'$-álgebra de tipo finito, se tiene
necesariamente $B'=B$ por (21.1.6).
\end{proof}

\begin{remark}[21.1.8]
\label{0.21.1.8-es}
Cuando $B$ es un cuerpo, demostraremos la equivalencia de las propiedades a),
b) y c) sin hipótesis de finitud (21.4.5).
\end{remark}

\begin{definition}[21.1.9]
\label{0.21.1.9-es}
Sean $A$ un anillo (de característica $p$), $B$ una $A$-álgebra. Diremos
que una familia $(x_\alpha)_{\alpha\in I}$ de elementos de $B$ es
$p$-libre sobre $A$ (resp. un sistema de $p$-generadores de $B$ sobre $A$,
resp. una $p$-base de $B$ sobre $A$) si la familia de monomios
$\prod_\alpha x_\alpha^{n(\alpha)}$ ($0\leq n(\alpha)<p$,
$(n(\alpha))_{\alpha\in I}$ de soporte finito) es una familia libre (resp.
un sistema de generadores, resp. una base) en el $A[B^p]$-módulo $B$.
\end{definition}

Se tienen definiciones correspondientes para un conjunto $M$ de elementos de
$B$, considerando la familia definida por la inyección canónica $M\to B$.
Cuando se toma $A$ igual al cuerpo primo $\mathbf F_p$ (en cuyo caso
$A[B^p]=B^p$), se omite la mención de $A$ en la definición anterior (o se
dice también que una familia es «absolutamente» $p$-libre, resp. un sistema
«absoluto» de $p$-generadores, resp. una $p$-base «absoluta»).

\oldpage[0\textsubscript{IV-1}]{157}
\begin{lemma}[21.1.10]
\label{0.21.1.10-es}
Sean $C$ una sub-$A$-álgebra de $B$ tal que $B^p\subset C$, $M$ una
parte de $C$ y $N$ una parte de $B$.
\begin{enumerate}
  \item[(i)] Si $M$ es un sistema de $p$-generadores de $C$ sobre $A$ y $N$
  es un sistema de $p$-generadores de $B$ sobre $C$, entonces $M\cup N$ es un
  sistema de $p$-generadores de $B$ sobre $A$.
  \item[(ii)] Supongamos que $M$ es una $p$-base de $C$ sobre $A$; entonces,
  para que $N$ sea $p$-libre sobre $C$, es necesario y suficiente que $M\cup N$
  sea $p$-libre sobre $A$.
\end{enumerate}
\end{lemma}

\begin{proof}
Podemos limitarnos al caso en que $B^p\subset A\subset C\subset B$. Entonces
(i) es un caso particular del hecho de que, si $P$ (resp. $Q$) es un sistema
de generadores del $A$-módulo $C$ (resp. del $C$-módulo $B$), el conjunto de
los productos $xy$, donde $x\in P$ y $y\in Q$, es un sistema de generadores
del $A$-módulo $B$.

Conservando las mismas notaciones, si $P$ es una base del $A$-módulo $C$,
decir que $Q$ es una familia libre sobre $C$ significa que la relación
$\sum_{\lambda,\mu}a_{\lambda\mu}x_\lambda y_\mu=0$, donde
$a_{\lambda\mu}\in A$, $x_\lambda\in P$, $y_\mu\in Q$ (siendo los
$x_\lambda$ (resp. los $y_\mu$) distintos dos a dos), equivale a
$a_{\lambda\mu}=0$ para todo par $(\lambda,\mu)$, o también a
$\sum_\lambda a_{\lambda\mu}x_\lambda=0$ para todo $\mu$; de ahí la
afirmación (ii).
\end{proof}

\subsection{$p$-bases y lisitud formal}
\label{0.21.2-es}

\begin{theorem}[21.2.1]
\label{0.21.2.1-es}
Sean $B$ un anillo, $A$ un subanillo de $B$ tal que $B^p\subset A\subset B$,
$(x_\alpha)_{\alpha\in I}$ una $p$-base de $B$ sobre $A$.
Sean $E$ una $A$-álgebra, $u:A\to E$ el homomorfismo estructural,
$q:E\to B$ un $A$-homomorfismo sobreyectivo, $\mathfrak R$ su núcleo, y
supongamos que $\mathfrak R^p=0$. Entonces:
\begin{enumerate}
  \item[(i)] Para que exista un $A$-homomorfismo $v:B\to E$ inverso por la
  derecha del homomorfismo $q:E\to B$, es necesario y suficiente que
  $u(q(z)^p)=z^p$ para todo $z\in E$.
  \item[(ii)] Cuando se satisface la condición de (i), para toda familia
  $(z_\alpha)_{\alpha\in I}$ de elementos de $E$ tal que
  $q(z_\alpha)=x_\alpha$ para todo $\alpha\in I$, existe un único
  $A$-homomorfismo $v:B\to E$ tal que $v(x_\alpha)=z_\alpha$ para todo
  $\alpha\in I$, y $v$ es inverso por la derecha de $q$.
\end{enumerate}
\end{theorem}

\begin{proof}
Si existe un $A$-homomorfismo $v:B\to E$, se debe tener $v(a)=u(a)$ para
todo $a\in A\subset B$, y por tanto $v(q(z)^p)=u(q(z)^p)$ para todo $z\in E$,
ya que $B^p\subset A$. Pero se tiene $v(q(z)^p)=(v(q(z)))^p$ y, por definición,
$v(q(z))\equiv z\pmod{\mathfrak R}$ si $q\circ v=1_B$, de modo que
$(v(q(z)))^p=z^p$ puesto que $\mathfrak R^p=0$; de ahí la necesidad de (i).
La suficiencia de la condición (i) resultará de (ii).

Ahora bien, bajo las hipótesis de (ii), la unicidad de $v$ es evidente, pues
los monomios $\prod_\alpha x_\alpha^{n(\alpha)}$ generan el $A$-módulo $B$;
como estos monomios forman además una \emph{base} del $A$-módulo $B$, existe
una aplicación $A$-\emph{lineal} $v$ de $B$ en $E$ tal que
$v(\prod_\alpha x_\alpha^{n(\alpha)})=\prod_\alpha z_\alpha^{n(\alpha)}$
para toda familia $(n(\alpha))_{\alpha\in I}$ de soporte finito con
$0\leq n(\alpha)<p$ para todo $\alpha$. Queda por ver que $v$ es un
homomorfismo de anillos. Ahora bien, se puede escribir
$(\prod_\alpha x_\alpha^{m(\alpha)})(\prod_\alpha x_\alpha^{n(\alpha)})
=a\prod_\alpha x_\alpha^{r(\alpha)}$, donde
$r(\alpha)=m(\alpha)+n(\alpha)$ si $m(\alpha)+n(\alpha)<p$,
$r(\alpha)=m(\alpha)+n(\alpha)-p$ en caso contrario, y $a\in A$
\oldpage[0\textsubscript{IV-1}]{158}
es el producto de los $x_\alpha^p$ para los $\alpha\in I$ tales que
$m(\alpha)+n(\alpha)\geq p$; se trata de comprobar que
$u(a)=\prod_\alpha z_\alpha^p$; pero como $x_\alpha=q(z_\alpha)$, esto se
deduce de la hipótesis.
\end{proof}

\begin{corollary}[21.2.2]
\label{0.21.2.2-es}
Sean $B$ un anillo, $A$ un subanillo de $B$ tal que $B^p\subset A\subset B$.
Si existe una $p$-base de $B$ sobre $A$, $B$ es una $A$-álgebra formalmente
lisa respecto de $B^p$ (para las topologías discretas).
\end{corollary}

\begin{proof}
Sean $E$ una $A$-extensión de $B$ por un $B$-módulo $L$, $q:E\to B$ la
aumentación, $u:A\to E$ el homomorfismo estructural. Decir que $E$ es
$B^p$-trivial significa que existe un homomorfismo de anillos $v:B\to E$
tal que $q(v(b))=b$ y $v(b^p)=u(b^p)$ para todo $b\in B$. Se deduce, para
$z\in E$, que $u(q(z)^p)=v(q(z)^p)=(v(q(z)))^p$; pero, en virtud de la
relación $L^2=0$, también se tiene $L^p=0$, y como $v(q(z))-z\in L$, se
tiene $(v(q(z)))^p=z^p$. Por tanto se satisface la condición de (21.2.1),
(i), y $E$ es $B^p$-trivial.
\end{proof}

\begin{corollary}[21.2.3]
\label{0.21.2.3-es}
Sean $B$ un anillo, $A$ un subanillo de $B$ tal que $B^p\subset A\subset B$,
y $\left(x_\alpha\right)_{\alpha\in I}$ una $p$-base de $B$ sobre $A$. Sea
$L$ un $B$-módulo. Entonces:
\begin{enumerate}
  \item[(i)] Para que una derivación $D$ de $A$ en $L$ se prolongue a una
  derivación de $B$ en $L$, es necesario y suficiente que $D$ se anule en
  $B^p$.
  \item[(ii)] Si $D$ se anula en $B^p$, entonces, para toda familia
  $(y_\alpha)_{\alpha\in I}$ de elementos de $L$, existe una única derivación
  $D'$ de $B$ en $L$ que prolonga $D$ y tal que
  $D'(x_\alpha)=y_\alpha$ para todo $\alpha\in I$.
\end{enumerate}
\end{corollary}

\begin{proof}
Dada una derivación $D$ de $A$ en $L$, consideremos el anillo $E=D_B(L)$ y
el homomorfismo $u:A\to E$ definido por $u(a)=(a,D(a))$; un
$A$-homomorfismo $v:B\to E$ inverso por la derecha del homomorfismo canónico
$q:E\to B$ es entonces de la forma $x\mapsto(x,D'(x))$, donde $D'$ es una
derivación de $B$ en $L$ que prolonga $D$ (20.1.5); como
$u(q(z)^p)=(q(z)^p,D(q(z)^p))$ para $z\in E$, el corolario se deduce
inmediatamente de (21.2.1) aplicado a $E$.
\end{proof}

Se deduce de (21.2.3) que la sucesión
\begin{equation}
\tag{21.2.3.1}
\label{0.21.2.3.1-es}
0\to\operatorname{Der}_A(B,L)\to\operatorname{Der}_{B^p}(B,L)
\to\operatorname{Der}_{B^p}(A,L)\to0
\end{equation}
(cf. (20.2.2.1)) es exacta, y que la aplicación
\begin{equation}
\tag{21.2.3.2}
\label{0.21.2.3.2-es}
D'\longmapsto(D'(x_\alpha))_{\alpha\in I}
\end{equation}
es un isomorfismo del $B$-módulo $\operatorname{Der}_A(B,L)$ sobre el
$B$-módulo producto $L^I$ (tomando $D=0$ en (21.2.3)).

\begin{corollary}[21.2.4]
\label{0.21.2.4-es}
Bajo las hipótesis de (21.2.3), la sucesión
\begin{equation}
\tag{21.2.4.1}
\label{0.21.2.4.1-es}
0\to\Omega^1_{A/B^p}\otimes_A B\to\Omega^1_{B/B^p}
\to\Omega^1_{B/A}\to0
\end{equation}
es exacta y escindida, y la familia $(d_{B/A}(x_\alpha))_{\alpha\in I}$ es
una base del $B$-módulo $\Omega^1_{B/A}$.
\end{corollary}

\begin{proof}
Esto se deduce inmediatamente de (21.2.3) y de la fórmula
$\operatorname{Der}_A(B,L)=\operatorname{Hom}_B(\Omega^1_{B/A},L)$ (20.4.8).
\end{proof}

\begin{corollary}[21.2.5]
\label{0.21.2.5-es}
Sean $A$ un anillo, $B$ una $A$-álgebra y $(x_\alpha)_{\alpha\in I}$ una
$p$-base de $B$ sobre $A$. Entonces $(d_{B/A}(x_\alpha))_{\alpha\in I}$ es
una base del $B$-módulo $\Omega^1_{B/A}$.
\end{corollary}

\begin{proof}
Usando (21.1.5), (ii), nos reducimos al caso en que $A=A[B^p]$ y basta
entonces aplicar (21.2.4).
\end{proof}

\oldpage[0\textsubscript{IV-1}]{159}
\begin{env}[21.2.6]
\label{0.21.2.6-es}
Sean $A$, $B$ dos anillos, $u:A\to B$ un homomorfismo de anillos; se tiene
(con las notaciones de (21.1.4)) el diagrama conmutativo
\begin{equation}
\tag{21.2.6.1}
\label{0.21.2.6.1-es}
\xymatrix{
 A^{(p)} \ar[r]^{u} & B^{(p)}\\
 A \ar[u]^{F_A} \ar[r]^{u} & B \ar[u]_{F_B}
}
\end{equation}
De él se deduce canónicamente un homomorfismo de $B$-álgebras
\begin{equation}
\tag{21.2.6.2}
\label{0.21.2.6.2-es}
u\otimes1:A^{(p)}\otimes_A B\longrightarrow B^{(p)}
\end{equation}
cuya imagen es el anillo $A[B^p]$ (que es una $B$-álgebra para el
homomorfismo $F_B:B\to B^p$).
\end{env}

\begin{theorem}[21.2.7]
\label{0.21.2.7-es}
Sean $A$ un anillo, $B$ una $A$-álgebra, $u:A\to B$ el homomorfismo
estructural. Supongamos verificadas las condiciones siguientes:
\begin{enumerate}
  \item[(i)] El homomorfismo (21.2.6.2) deducido de $u$ es inyectivo.
  \item[(ii)] $B$ admite una $p$-base $(x_\alpha)_{\alpha\in I}$ respecto de
  $A$.
\end{enumerate}
Entonces $B$ es una $A$-álgebra formalmente lisa (para las topologías
discretas). Más precisamente, dotemos $A$ y $B$ de las topologías discretas;
sean $E$ una $A$-álgebra topológica admisible (0$_I$, 7.1.2),
$\mathfrak R$ un ideal de definición de $E$, $C=E/\mathfrak R$,
$v:B\to C$ un $A$-homomorfismo, $q:E\to C$ la aumentación. Entonces, para
toda familia $(z_\alpha)_{\alpha\in I}$ de elementos de $E$ tal que
$q(z_\alpha)=v(x_\alpha)$ para todo $\alpha\in I$, existe un único
$A$-homomorfismo $w:B\to E$ tal que $v=q\circ w$ y
$w(x_\alpha)=z_\alpha$ para todo $\alpha\in I$.
\end{theorem}

\begin{proof}
Consideremos primero el caso en que $E$ es discreto, de modo que $\mathfrak R$
es nilpotente. Podemos limitarnos al caso en que $v$ es sobreyectivo; además,
el razonamiento de (19.4.3) permite suponer $\mathfrak R^2=0$ y, por tanto,
$\mathfrak R^p=0$. Finalmente, considerando la imagen inversa por $v$ de la
extensión $E$ de $C$ por $\mathfrak R$, podemos limitarnos al caso en que $C=B$
y $v$ es la identidad (19.4.4). Como el homomorfismo de anillos
$F_E:z\rightsquigarrow z^p$ se anula en $\mathfrak R$, se factoriza como
$E\xrightarrow q B\xrightarrow{F'}E$, y $F'$, considerado como homomorfismo de
$B$ en $E^{(p)}$, es un \emph{$A$-homomorfismo} por definición de la estructura de
$A$-álgebra de $E^{(p)}$ mediante el homomorfismo compuesto
$A\xrightarrow r E\xrightarrow{F_E}E^{(p)}$, donde $r:A\to E$ es el homomorfismo
estructural de la $A$-álgebra $E$; por otra parte,
$r:A^{(p)}\to E^{(p)}$ es también un $A$-homomorfismo. Existe, pues, un único
$A$-homomorfismo $f:A^{(p)}\otimes_A B\to E^{(p)}$ tal que las composiciones de
$f$ con los homomorfismos canónicos $A^{(p)}\to A^{(p)}\otimes_A B$ y
$B\to A^{(p)}\otimes_A B$ sean, respectivamente, $r$ y $F'$. Ahora bien, por
hipótesis, podemos identificar $A^{(p)}\otimes_A B$ con $A[B^p]$, identificándose
los homomorfismos canónicos $A^{(p)}\to A^{(p)}\otimes_A B$ y
$B\to A^{(p)}\otimes_A B$ respectivamente con $u$ y con $F_B$. Podemos, pues,
considerar ahora $E$ como una $A[B^p]$-álgebra mediante el homomorfismo $f$, y,
por construcción, tenemos $f(q(z)^p)=z^p$ para todo $z\in E$; se cumplen, por
tanto, las condiciones de aplicación de (21.2.1), de donde se deduce el
teorema en este caso.

Para pasar al caso general, consideremos un sistema fundamental
$(\mathfrak J_\lambda)$ de ideales abiertos de $E$, y pongamos
$E_\lambda=E/\mathfrak J_\lambda$,
$\mathfrak R_\lambda=(\mathfrak R+\mathfrak J_\lambda)/\mathfrak J_\lambda$,
$C_\lambda=E_\lambda/\mathfrak R_\lambda=E/(\mathfrak R+\mathfrak J_\lambda)$;
como $E$ es admisible, tenemos $E=\varprojlim E_\lambda$ (0$_I$, 7.2.2).
Para cada par $(\alpha,\lambda)$, denotemos por $z_{\alpha\lambda}$ la imagen
canónica de $z_\alpha$ en $E_\lambda$, por $p_\lambda:C\to C_\lambda$ el
homomorfismo canónico,
\oldpage[0\textsubscript{IV-1}]{160}
y por $q_\lambda$ el homomorfismo canónico $E_\lambda\to C_\lambda$.
Como, por hipótesis, $\mathfrak R_\lambda$ es nilpotente en $E_\lambda$, la primera
parte de la demostración prueba la existencia y unicidad de un
$A$-homomorfismo $w_\lambda:B\to E_\lambda$ tal que
$p_\lambda\circ v=q_\lambda\circ w_\lambda$ y
$w_\lambda(x_\alpha)=z_{\alpha\lambda}$ para todo $\alpha$.
La unicidad de los $w_\lambda$ muestra entonces inmediatamente que $(w_\lambda)$
es un sistema proyectivo de homomorfismos, y $w=\varprojlim w_\lambda$ responde
a la cuestión.
\end{proof}

\begin{remark}[21.2.8]
\label{0.21.2.8-es}
La verificación de la existencia y unicidad del homomorfismo $w:B\to E$ tal
que $w(x_\alpha)=z_\alpha$ para todo $\alpha$ se deduce directamente del
hecho de que $B$ es una $A$-álgebra formalmente lisa y de que los
$d_{B/A}(x_\alpha)$ forman una base del $B$-módulo $\Omega^1_{B/A}$ (21.2.4),
sin hacer intervenir el hecho de que se trata de una $p$-base (de modo que el
resultado es válido sin suponer que los anillos sean de característica $p$).
En efecto, podemos limitarnos al caso en que $\mathfrak R^2=0$; como
$\operatorname{Der}_A(B,\mathfrak R)=\operatorname{Hom}_B(\Omega^1_{B/A},\mathfrak R)$
(20.4.8), existe una única $A$-derivación $D$ de $B$ en $\mathfrak R$ tal que
$D(x_\alpha)=z'_\alpha$ para toda familia $(z'_\alpha)$ de elementos de
$\mathfrak R$; la conclusión se deduce de (20.1.1).
\end{remark}

\subsection{$p$-bases y módulos de imperfección}
\label{0.21.3-es}

\begin{env}[21.3.1]
\label{0.21.3.1-es}
Sean $A$, $B$ dos anillos (de característica $p$), $i:A\to B$,
$j:B\to A$ dos homomorfismos de anillos tales que
\begin{align}
\tag{21.3.1.1}\label{0.21.3.1.1-es}
j(i(a))&=a^p\quad\text{para todo }a\in A,\\
\tag{21.3.1.2}\label{0.21.3.1.2-es}
i(j(b))&=b^p\quad\text{para todo }b\in B.
\end{align}
Con frecuencia $i$ será inyectivo, de modo que $A$ se identificará mediante
$i:A\to B$ con un subanillo de $B$; hecha esta identificación, la existencia
de $j$ implica que $B^p\subset A$, y $j$ se identifica entonces con $F_B$.
\end{env}

\begin{env}[21.3.2]
\label{0.21.3.2-es}
Si $h:A^p\to A$ es la inyección canónica, se tiene, según (21.3.1.1) y
(21.3.1.2), un diagrama conmutativo
\begin{equation}
\tag{21.3.2.1}\label{0.21.3.2.1-es}
\xymatrix{
 B \ar[r]^{j} & A\\
 A \ar[u]^{i} \ar[r]^{F_A} & A^p \ar[u]_{h}
}
\end{equation}
por lo que el par $(j,F_A)$ puede considerarse un dihomomorfismo del
$A$-álgebra $B$ (para $i$) en el $A^p$-álgebra $A$ (para $h$).
De ello se deduce un homomorfismo canónico de $A$-módulos
\begin{equation}
\tag{21.3.2.2}\label{0.21.3.2.2-es}
\pi_{B/A}:\Omega^1_{B/A}\otimes_B A_{[j]}\longrightarrow
\Omega^1_{A/A^p}=\Omega^1_A
\end{equation}
(cf. (20.5.4); la identificación de $\Omega^1_{A/A^p}$ con el módulo de
diferenciales absolutas $\Omega^1_A$ proviene de (21.1.5), (ii)).
\begin{align}
\tag{21.3.2.3}\label{0.21.3.2.3-es}
\Theta_{B/A}&=\Omega^1_{B/A}\otimes_B A_{[j]},\\
\tag{21.3.2.4}\label{0.21.3.2.4-es}
\Xi_{B/A}&=\operatorname{Ker}(\pi_{B/A}).
\end{align}
\oldpage[0\textsubscript{IV-1}]{161}
Por tanto se tiene la sucesión exacta
\begin{equation}
\tag{21.3.2.5}\label{0.21.3.2.5-es}
0\to\Xi_{B/A}\to\Theta_{B/A}\xrightarrow{\pi_{B/A}}\Omega^1_A
\end{equation}
(se observará que estas notaciones pueden inducir a confusión, pues
$\Theta_{B/A}$ y $\Xi_{B/A}$ dependen no solo de $A$ y $B$, sino también de
$i$ y $j$).
\end{env}

\begin{env}[21.3.3]
\label{0.21.3.3-es}
Como, en virtud de (21.3.1.2), se tiene $F_B=i\circ j$, se puede escribir,
para todo $B$-módulo $M$ (cf. (21.1.4)),
\begin{equation}
\tag{21.3.3.1}\label{0.21.3.3.1-es}
M^{(p)}=M\otimes_B B^{(p)}=(M\otimes_B A_{[j]})\otimes_A B_{[i]},
\end{equation}
de donde, en particular,
\begin{equation}
\tag{21.3.3.2}\label{0.21.3.3.2-es}
(\Omega^1_{B/A})^{(p)}=\Theta_{B/A}\otimes_A B_{[i]},
\end{equation}
y, por tanto, de (21.3.2.2) se deduce un homomorfismo canónico de
$B$-módulos
\begin{equation}
\tag{21.3.3.3}\label{0.21.3.3.3-es}
\pi_{B/A}\otimes1_B:(\Omega^1_{B/A})^{(p)}\longrightarrow
\Omega^1_A\otimes_A B_{[i]}.
\end{equation}
Teniendo en cuenta (20.5.4.2), la imagen de este homomorfismo es el $B$-módulo
generado por los elementos $d_A(j(b))\otimes1$ para $b\in B$. Su imagen por
el homomorfismo canónico $i_{B/A}:\Omega^1_A\otimes_A B\to\Omega^1_B$
deducido de $i$ (20.5.2) está, por tanto, en el sub-$B$-módulo de
$\Omega^1_B$ generado por los $d_B(i(j(b)))$ para $b\in B$; en virtud de
(21.3.1.2), esta imagen es nula. En otras palabras, se tiene una sucesión de
homomorfismos
\begin{equation}
\tag{21.3.3.4}\label{0.21.3.3.4-es}
0\to\Xi_{B/A}\otimes_A B_{[i]}\to(\Omega^1_{B/A})^{(p)}
\to\Upsilon_{B/A}
\end{equation}
que no es necesariamente exacta, pero en la que la composición de dos
homomorfismos consecutivos es nula.
\end{env}

\begin{proposition}[21.3.4]
\label{0.21.3.4-es}
Si $B$ es un $A$-módulo plano (para $i$), la sucesión (21.3.3.4) es exacta.
\end{proposition}
\begin{proof}
Esto se deduce de la definición de la planitud. La proposición se aplica en
particular cuando $B^p\subset A\subset B$, siendo $i$ y $j$ respectivamente
la inyección canónica y $F_B$, y $B$ admite una $p$-base sobre $A$ (de modo
que $B$ es entonces un $A$-módulo libre). Pero incluso en este caso el núcleo
$\Xi_{B/A}$ no es necesariamente nulo. Sin embargo:
\end{proof}

\begin{proposition}[21.3.5]
\label{0.21.3.5-es}
Sean $B$ un anillo reducido, $A$ un subanillo de $B$ tal que
$B^p\subset A\subset B$. Supongamos que existe una $p$-base
$(x_\lambda)_{\lambda\in L}$ de $B$ sobre $A$ y una $p$-base
$(y_\mu)_{\mu\in M}$ de $A$ sobre $B^p$; entonces el homomorfismo canónico
(21.3.3.4)
\begin{equation}
\tag{21.3.5.1}\label{0.21.3.5.1-es}
(\Omega^1_{B/A})^{(p)}\longrightarrow\Upsilon_{B/A}
\end{equation}
es biyectivo, y los elementos $d_A(x_\lambda^p)\otimes1$ forman una base del
$B$-módulo $\Upsilon_{B/A}$.
\end{proposition}

\begin{proof}
\oldpage[0\textsubscript{IV-1}]{162}
Como $B$ es reducido, $x\mapsto x^p$ es un isomorfismo estructural de $B$
sobre $B^p$; transportando la estructura mediante este isomorfismo, se ve
que $(x_\lambda^p)_{\lambda\in L}$ y los $y_\mu$ forman una $p$-base de
$B^p$ sobre $A^p$ (21.1.10); se deduce entonces de (21.1.5), (ii), y (21.2.5)
que los $d_A(x_\lambda^p)$ y los $d_A(y_\mu)$ forman una base del $A$-módulo
$\Omega^1_A$. Ahora bien, la imagen de $d_A(x_\lambda^p)\otimes1$ por
$i_{B/A}$ es $d_B(x_\lambda^p)=0$; por otra parte, la imagen de $d_A(y_\mu)$
por $i_{B/A}$ es $d_B(i(y_\mu))$, y como los $y_\mu$ y los $x_\lambda$
forman una $p$-base de $B$ sobre $B^p$ (21.1.10), los $d_B(y_\mu)$ son una
parte de una base de $\Omega^1_B$ (21.2.5), por lo que son linealmente
independientes sobre $B$. Se deduce que el núcleo de $i_{B/A}$ tiene por base
los $d_A(x_\lambda^p)\otimes1$; como estos son las imágenes por (21.3.5.1) de
los elementos correspondientes de $(\Omega^1_{B/A})^{(p)}$, y estos últimos
forman una base del $B$-módulo $(\Omega^1_{B/A})^{(p)}$ (21.2.5), el
homomorfismo (21.3.5.1) de $B$-módulos es biyectivo.
\end{proof}

\begin{remarks}[21.3.6]
\label{0.21.3.6-es}
\begin{enumerate}
  \item[(i)] Sean $A'$, $B'$ dos anillos de característica $p$, y supongamos
  dados dos homomorfismos $i':A'\to B'$, $j':B'\to A'$ que satisfacen
  (21.3.1.1) y (21.3.1.2), así como homomorfismos de anillos
  $f:A\to A'$, $g:B\to B'$ que hacen conmutativo el diagrama
  \[
  \xymatrix{
  A'\ar[r]^{i'}&B'\ar[r]^{j'}&A'\\
  A\ar[u]^f\ar[r]^i&B\ar[u]^g\ar[r]^j&A\ar[u]^f
  }
  \]
  Entonces el homomorfismo canónico $\Omega^1_{B/A}\to\Omega^1_{B'/A'}$
  (20.5.4.3) da un dihomomorfismo canónico
  $\Theta_{B/A}\to\Theta_{B'/A'}$ que hace conmutativo el diagrama
  \[
  \xymatrix{
  0\ar[r]&\Xi_{B/A}\ar[r]\ar[d]&\Theta_{B/A}\ar[r]^{\pi_{B/A}}\ar[d]&\Omega^1_A\ar[d]\\
  0\ar[r]&\Xi_{B'/A'}\ar[r]&\Theta_{B'/A'}\ar[r]^{\pi_{B'/A'}}&\Omega^1_{A'}
  }
  \]
  \item[(ii)] Sea $A'$ una $A$-álgebra cualquiera y pongamos
  $B'=B\otimes_A A'$; entonces $i'=i\otimes1:A'\to B'$ y
  $j'=j\otimes1:B'\to A'$ satisfacen (21.3.1.1) y (21.3.1.2), y se tiene
  $\Omega^1_{B'/A'}=\Omega^1_{B/A}\otimes_A A'=\Omega^1_{B/A}\otimes_B B'$
  (20.5.5); se deduce un $A'$-isomorfismo canónico
  \begin{equation}
  \tag{21.3.6.1}\label{0.21.3.6.1-es}
  \Theta_{B/A}\otimes_A A'\xrightarrow{\sim}\Theta_{B'/A'}
  \end{equation}
  y también un $B'$-homomorfismo canónico
  \begin{equation}
  \tag{21.3.6.2}\label{0.21.3.6.2-es}
  (\Omega^1_{B/A})^{(p)}\otimes_B B'\xrightarrow{\sim}
  (\Omega^1_{B'/A'})^{(p)}.
  \end{equation}
\end{enumerate}
\end{remarks}

\subsection{Caso de las extensiones de cuerpos}
\label{0.21.4-es}

\begin{env}[21.4.0]
\label{0.21.4.0-es}
Sean $K$ un cuerpo de característica $p>0$, $k$ un subcuerpo de $K$; entonces
el anillo $k[K^p]$ es igual al cuerpo $k(K^p)$, puesto que $k$ es algebraico
sobre $K^p$. Por tanto se pueden aplicar los resultados de los números
anteriores reemplazando en todas partes $A$, $B$ y $A[B^p]$ por $k$, $K$ y
$k(K^p)$.
\end{env}

\begin{lemma}[21.4.1]
\label{0.21.4.1-es}
Sean $k$ un cuerpo de característica $p>0$, $K$ una extensión de $k$. Para
que un elemento $x\in K$ sea $p$-libre sobre $k$, es necesario y suficiente
que $x\notin k(K^p)$.
\end{lemma}

\begin{proof}
En efecto, $x$ es raíz del polinomio $X^p-x^p$ de $k(K^p)[X]$, y se sabe
(Bourbaki, \emph{Alg.}, cap.~V, §~8, n$^\circ$~1, prop.~1) que, si
$x\notin k(K^p)$, este polinomio es irreducible, de modo que los elementos
$1,x,\ldots,x^{p-1}$ forman una base del $k(K^p)$-módulo $k(K^p)(x)$.
\end{proof}

\begin{theorem}[21.4.2]
\label{0.21.4.2-es}
Sean $k$ un cuerpo de característica $p>0$, $K$ una extensión de $k$,
$S$ un sistema de $p$-generadores de $K$ sobre $k$, $L\subset S$ una parte
$p$-libre sobre $k$. Existe entonces una $p$-base $B$ de $K$ sobre $k$ tal que
$L\subset B\subset S$. En particular, toda extensión de $k$ admite una
$p$-base sobre $k$.
\end{theorem}

\begin{proof}
Podemos limitarnos al caso en que $K^p\subset k$. En virtud del teorema de
Zorn, existe en $K$ una parte $B$ tal que $L\subset B\subset S$, $p$-libre
sobre $k$ y maximal entre las partes de $S$ con estas propiedades. Basta ver
que el subcuerpo $K'$ de $K$ generado por $k$ y $B$ es igual a $K$.
En caso contrario existiría $x\in S$ que no pertenece a $K'$; como
$K^p\subset k$, se tiene $K'=K'(K^p)$; por tanto $x$ sería $p$-libre sobre
$K'$ (21.4.1), y, en consecuencia, $B\cup\{x\}$ sería $p$-libre sobre $k$
(21.1.10), en contradicción con la definición de $B$. La última afirmación de
(21.4.2) se obtiene tomando $L=\varnothing$, $S=K$.
\end{proof}

\begin{corollary}[21.4.3]
\label{0.21.4.3-es}
Sean $k$ un cuerpo de característica $p>0$, $K$ una extensión de $k$.
Para que una familia $(x_\alpha)$ de elementos de $K$ sea $p$-independiente sobre $k$,
es necesario y suficiente que, para todo $\alpha$, $x_\alpha$ no pertenezca al
cuerpo $K_\alpha$ generado por $k(K^p)$ y por los $x_\beta$ de índice
$\beta\neq\alpha$.
\end{corollary}

\begin{proof}
La condición es necesaria en virtud de (21.1.10). Recíprocamente, supongamos que se
cumple; podemos limitarnos al caso en que $(x_\alpha)$ sea un sistema de
$p$-generadores de $K$. Existe entonces una subfamilia de $(x_\alpha)$ que es una
$p$-base de $K$ (21.4.2), pero esta familia no puede ser distinta de $(x_\alpha)$,
pues, en virtud de la hipótesis, no sería un sistema de $p$-generadores de $K$.
\end{proof}

\oldpage[0\textsubscript{IV-1}]{163}
\begin{corollary}[21.4.4]
\label{0.21.4.4-es}
Sean $k$ un cuerpo de característica $p>0$, $K$ una extensión de $k$.
Para que $\Omega^1_{K/k}=0$, es necesario y suficiente que $K=k(K^p)$.
En particular, para que $\Omega^1_K=0$, es necesario y suficiente que $K$ sea
un cuerpo perfecto.
\end{corollary}

\begin{proof}
En efecto, si $(x_\alpha)$ es una $p$-base de $K$ sobre $k$, los
$d_{K/k}(x_\alpha)$ forman una base del espacio vectorial $\Omega^1_{K/k}$
(21.2.5).
\end{proof}

\begin{theorem}[21.4.5]
\label{0.21.4.5-es}
Sean $k$ un cuerpo de característica $p>0$, $K$ una extensión de $k$,
$(x_\alpha)$ una familia de elementos de $K$. Para que $(x_\alpha)$ sea
$p$-independiente sobre $k$ (resp. una familia de $p$-generadores de $K$ sobre $k$,
resp. una $p$-base de $K$ sobre $k$), es necesario y suficiente que la familia
$(d_{K/k}(x_\alpha))$ sea una familia libre (resp. un sistema de
generadores, resp. una base) en el espacio vectorial $\Omega^1_{K/k}$.
\end{theorem}

\begin{proof}
\oldpage[0\textsubscript{IV-1}]{164}
Sea $K'$ el subcuerpo (igual al subanillo) de $K$ generado por
$k(K^p)$ y los $x_\alpha$; teniendo en cuenta (21.4.3), basta ver que,
para cada $\alpha$, $x_\alpha$ no pertenece a $K_\alpha$. Ahora bien, en la
sucesión exacta
\[
\Omega^1_{K_\alpha/k(K^p)}\otimes_{K_\alpha}K\longrightarrow
\Omega^1_{K/k(K^p)}\longrightarrow\Omega^1_{K/K_\alpha}\longrightarrow0,
\]
las imágenes por la flecha de la izquierda de los $d_{K_\alpha/k}(x_\beta)$ para
$\beta\neq\alpha$ son los $d_{K/k}(x_\beta)$, por lo que generan un subespacio
vectorial de $\Omega^1_{K/k(K^p)}$ que no contiene a $d_{K/k}(x_\alpha)$; como
este subespacio es el núcleo de la flecha de la derecha, se ve que
$d_{K/K_\alpha}(x_\alpha)\neq0$, y por tanto $x_\alpha\notin K_\alpha$.

Si ahora $(x_\alpha)$ es una familia $p$-independiente, es una subfamilia
de una $p$-base de $K$ sobre $k$ (21.4.2); los $d_{K/k}(x_\alpha)$ forman entonces
parte de una base del espacio vectorial $\Omega^1_{K/k}$ (21.2.5), y son
por consiguiente linealmente independientes sobre $K$. Recíprocamente, si los
$d_{K/k}(x_\alpha)$ son linealmente independientes sobre $K$, los dos
párrafos precedentes muestran que $(x_\alpha)$ es $p$-independiente sobre $k$.
\end{proof}

\begin{corollary}[21.4.6]
\label{0.21.4.6-es}
Sean $k$ un cuerpo de característica $p>0$, $K$ una extensión de $k$,
$x$ un elemento de $K$. Las tres condiciones siguientes son equivalentes:
\begin{enumerate}
  \item[a)] $x\notin k(K^p)$.
  \item[b)] $d_{K/k}(x)\neq0$.
  \item[c)] El elemento $x$ es $p$-independiente sobre $k$.
\end{enumerate}
\end{corollary}

\begin{proposition}[21.4.7]
\label{0.21.4.7-es}
Sean $L$ un cuerpo de característica $p>0$, $K$ un subcuerpo de $L$ tal
que $L^p\subset K\subset L$. Entonces la sucesión de espacios vectoriales sobre $L$
\begin{equation}
\tag{21.4.7.1}\label{0.21.4.7.1-es}
0\to\Omega^1_{K/L^p}\otimes_KL\to\Omega^1_{L/L^p}
\to\Omega^1_{L/K}\to0
\end{equation}
es exacta; en otros términos, se tiene $\Upsilon_{L/K/L^p}=0$.
\end{proposition}
\begin{proof}
Es un caso particular de (21.2.4), teniendo en cuenta (21.4.2).
\end{proof}

\begin{proposition}[21.4.8]
\label{0.21.4.8-es}
Bajo las hipótesis de (21.4.7), el homomorfismo canónico
\begin{equation}
\tag{21.4.8.1}\label{0.21.4.8.1-es}
(\Omega^1_{L/K})^{(p)}\longrightarrow\Upsilon_{L/K}
\end{equation}
es biyectivo; si $(x_\lambda)$ es una $p$-base de $L$ sobre $K$, los elementos
$d_K(x_\lambda^p)\otimes1$ forman una base del espacio vectorial $\Upsilon_{L/K}$.
\end{proposition}
\begin{proof}
Es un caso particular de (21.3.5).
\end{proof}

\subsection{Aplicación: criterios de separabilidad}
\label{0.21.5-es}

En este número y en los dos siguientes, ya no suponemos que los anillos
considerados sean de característica $p>0$.

\begin{env}[21.5.1]
\label{0.21.5.1-es}
Notemos primero que el criterio (21.2.7) permite demostrar una parte del
teorema de Cohen sobre las extensiones separables (19.6.1), a saber, que si $k$
es de característica $p>0$ y si $K$ es una extensión separable de $k$,
entonces $K$ es una $k$-álgebra formalmente lisa. En efecto, $K$ admite una
$p$-base sobre $k$ (21.4.2), y por otra parte se deduce del criterio de MacLane
(Bourbaki, \emph{Alg.}, cap.~V, §~8, n$^\circ$~2, prop.~3) que, en una
clausura algebraica de $K$, $k^{p^{-1}}$ y $K$ son linealmente disjuntos sobre
$k$, y por consiguiente que el homomorfismo canónico
$k^{p^{-1}}\otimes_k K\to k^{p^{-1}}(K)$ es biyectivo, lo que es precisamente
la condición (i) de (21.2.7), después de transportar la estructura mediante el
isomorfismo $x\mapsto x^p$.
\end{env}
\begin{proposition}[21.5.2]
\label{0.21.5.2-es}
\textnormal{(MacLane).}
Sean $A$, $B$ dos anillos de valoración discreta completos, $\mathfrak m$,
$\mathfrak n$ sus ideales máximos respectivos, $K=A/\mathfrak m$ y
$L=B/\mathfrak n$ sus cuerpos residuales, $u:A\to B$ un homomorfismo tal que
$Bu(\mathfrak m)=\mathfrak n$, $u_0=u\otimes 1:K\to L$ el homomorfismo
correspondiente para los cuerpos residuales. Consideremos las condiciones siguientes:
\begin{enumerate}
  \item[\rm a)] $L$ es una extensión separable de $K$ (para $u_0$).
  \item[\rm b)] Para todo anillo de valoración discreta completo $B'$ de ideal
  máximo $\mathfrak n'$ y de cuerpo residual $L'$,
\oldpage[0\textsubscript{IV-1}]{165}
  todo homomorfismo
  $u':A\to B'$ tal que $B'u'(\mathfrak m)=\mathfrak n'$ y todo
  $K$-isomorfismo $\sigma:L\to L'$ (relativo a $u_0$ y
  $u'_0=u'\otimes1:K\to L'$), existe un isomorfismo $w:B\to B'$ tal que
  $u'=w\circ u$ y $w_0=w\otimes1:L\to L'$ sea igual a $\sigma$.
  \item[\rm b$'$)] Para todo homomorfismo $u':A\to B$ tal que
  $Bu'(\mathfrak m)=\mathfrak n$ y tal que $u'_0=u'\otimes1:K\to L$ sea igual
  a $u_0$, existe un automorfismo $w$ de $B$ tal que $u'=w\circ u$ y
  $w_0=w\otimes1:L\to L$ sea la identidad.
  \item[\rm c)] Designando por $u_1:A/\mathfrak m^2\to B/\mathfrak n^2$
  el homomorfismo deducido de $u$ por paso a los cocientes, entonces, para todo
  homomorfismo local $u'_1:A/\mathfrak m^2\to B/\mathfrak n^2$ tal que
  $\operatorname{gr}_0(u'_1)=\operatorname{gr}_0(u_1)(=u_0)$, existe un
  automorfismo $w_1$ de $B/\mathfrak n^2$ tal que
  $\operatorname{gr}_0(w_1)$ sea la identidad y $u'_1=w_1\circ u_1$.
  \item[\rm c$'$)] Para todo homomorfismo local $u'_1:A/\mathfrak m^2\to
  B/\mathfrak n^2$ tal que $\operatorname{gr}_0(u'_1)=\operatorname{gr}_0(u_1)$
  y $\operatorname{gr}_1(u'_1)=\operatorname{gr}_1(u_1)$
  (homomorfismo de $\mathfrak m/\mathfrak m^2$ en
  $\mathfrak n/\mathfrak n^2$), existe un automorfismo $w_1$ de $B/\mathfrak n^2$
  tal que $\operatorname{gr}_0(w_1)$ sea la identidad y $u'_1=w_1\circ u_1$.
\end{enumerate}
Entonces se tienen las implicaciones $c)\Leftrightarrow c')\Leftrightarrow a)\Rightarrow b)
 \Rightarrow b')$. Si además $A$ es un anillo de Cohen, las cinco condiciones
anteriores son equivalentes.

\end{proposition}
\begin{proof}
Las implicaciones $b)\Rightarrow b')$ y $c)\Rightarrow c')$ son triviales.
Demostremos que $a)$ implica $b)$.
El homomorfismo $u$ hace de $B$ un $A$-módulo sin torsión, pues transforma por
hipótesis un uniformizante de $A$ en un uniformizante de $B$; se deduce que
$B$ es un $A$-módulo plano (0$_{\rm I}$, 6.3.4), luego un $A$-álgebra de
Cohen (19.8.1); el hecho de que $a)$ implique $b)$ es entonces consecuencia de
(19.8.2, (i)) aplicada a $G=B'$, $\mathfrak J=\mathfrak n'$,
considerando $B'$ como $A$-álgebra mediante $u'$ y siendo el homomorfismo
$B\to B'/\mathfrak n'$ el compuesto
$B\to B/\mathfrak n\to B'/\mathfrak n'$, que es un $A$-homomorfismo en
virtud de la hipótesis $u'_0=\sigma\circ u_0$. El mismo razonamiento prueba que
$a)$ implica $c)$, tomando esta vez $G=B/\mathfrak n^2$,
$\mathfrak J=\mathfrak n/\mathfrak n^2$, y considerando $G$ como
$A$-álgebra mediante $u'_1$.
Probemos en segundo lugar que $c')$ implica $a)$. Los dos homomorfismos
$u_1$ y $u'_1$ son tales que $u'_1=u_1+D$, donde $D$ es una derivación de
$A/\mathfrak m^2$ en $\mathfrak n/\mathfrak n^2$ (20.1.1); además,
la hipótesis $\operatorname{gr}_1(u'_1)=\operatorname{gr}_1(u_1)$ significa
que $D$ es nula en $\mathfrak m/\mathfrak m^2$, y por tanto puede
considerarse como una derivación de $K=A/\mathfrak m$ en
$\mathfrak n/\mathfrak n^2$, y $\mathfrak n/\mathfrak n^2$ se identifica
(eligiendo un uniformizante de $B$) con el $K$-módulo $L$ (para $u_0$).
Por otra parte, las condiciones impuestas a $w_1$ implican que
$\operatorname{gr}_1(w_1)$ es también la identidad (pues
$u_1(\mathfrak m/\mathfrak m^2)$ genera
$\mathfrak n/\mathfrak n^2$); por tanto $w_1$ es de la forma
$x\mapsto x+D'(x)$, donde $D'$ es ahora una derivación de
$B/\mathfrak n^2$ en el $(B/\mathfrak n^2)$-módulo
$\mathfrak n/\mathfrak n^2$; como $w_1$ es la identidad en
$\mathfrak n/\mathfrak n^2$, $D'$ puede también (mediante la identificación
anterior de $\mathfrak n/\mathfrak n^2$ con $L$) considerarse como una
derivación de $L$ en $L$; finalmente, la relación $u'_1=w_1\circ u_1$
significa que $D'$ prolonga $D$. Por otra parte, toda derivación $D$
de $K$ en $L$ corresponde a un homomorfismo $u'_1$ que verifica las condiciones
de $c')$ (20.1.1); por tanto, la condición $c')$ significa que toda derivación de
$K$ en $L$ se prolonga a una derivación de $L$ en sí mismo, es decir,
(20.6.5), que $L$ es separable sobre $K$.
Probemos finalmente que, cuando $A$ es un anillo de Cohen, $b')$ implica $c')$.
Probemos primero que bajo las hipótesis de $c')$, existe un homomorfismo
$u':A\to B$ que verifica las hipótesis de $b')$ y que, por paso a los cocientes,
da $u'_1:A/\mathfrak m^2\to B/\mathfrak n^2$. En efecto, esto se deduce de
(19.8.6, (i)) aplicado al homomorfismo compuesto
$v':A\to A/\mathfrak m^2\xrightarrow{u'_1}B/\mathfrak n^2$;
\oldpage[0\textsubscript{IV-1}]{166}
este último se factoriza entonces como $A\xrightarrow{u'}B\to B/\mathfrak n^2$
y las hipótesis sobre $\operatorname{gr}_0(u'_1)$ y
$\operatorname{gr}_1(u'_1)$ implican
$\operatorname{gr}_0(u')=\operatorname{gr}_0(u)=u_0$ y
$\operatorname{gr}_1(u')=\operatorname{gr}_1(u)$, de modo que la imagen por
$u'$ de un uniformizante de $A$ es un uniformizante de $B$, y por tanto
$Bu'(\mathfrak m)=\mathfrak n$. Basta entonces, para obtener un automorfismo
$w_1$ que responda a la cuestión, tomar el automorfismo inducido por paso a los
cocientes por el automorfismo $w$ de $B$ proporcionado por la aplicación de
$b')$ al homomorfismo $u'$.
\end{proof}
\begin{remarks}[21.5.3]
\label{0.21.5.3-es}
\begin{enumerate}
  \item[(i)] Las propiedades diferenciales de los cuerpos permiten resolver la
  cuestión de la unicidad del cuerpo de representantes en un anillo local
  noetheriano completo $C$ (19.8.7, (ii)). Sean $\mathfrak J$ el ideal máximo
  de $C$, $K=C/\mathfrak J$ el cuerpo residual de $C$; podemos limitarnos al
  caso $\mathfrak J\ne0$. Supongamos que existe un homomorfismo $u:K\to C$
  que, compuesto con la augmentación $C\to K$, da la identidad; entonces, para
  que este homomorfismo sea único, es necesario y suficiente que
  $\Omega^1_K=0$. En efecto, la condición $\Omega^1_K=0$ implica que $K$ es
  formalmente no ramificado sobre su cuerpo primo (20.7.4); si existiera un
  segundo homomorfismo $v\ne u$ que respondiera a la cuestión, habría un
  entero $n$ máximo tal que $\mathfrak J^n$ contuviese el conjunto de
  $u(x)-v(x)$ para $x\in K$; por paso al cociente, $u$ y $v$ darían dos
  homomorfismos distintos $u_1,v_1$ de $K$ en
  $C/\mathfrak J^{n+1}$, cuyos compuestos con
  $C/\mathfrak J^{n+1}\to C/\mathfrak J^n$ serían iguales, lo que contradice
  la definición (19.10.2). Recíprocamente, supongamos que $\Omega^1_K\ne0$;
  existe entonces una derivación $D\ne0$ de $K$ en
  $\mathfrak J/\mathfrak J^2$ (20.4.8), y por tanto un homomorfismo
  $v_1:K\to C/\mathfrak J^2$ tal que, si $u_1:K\to C/\mathfrak J^2$ se
  obtiene de $u$ por paso al cociente, se tiene $v_1=u_1+D$
  (20.1.1). Si $k$ es el cuerpo primo de $K$, $C$ es una $k$-álgebra y
  $K$ una $k$-álgebra formalmente lisa, y las restricciones de
  $u_1$ y $v_1$ a $k$ coinciden, por lo que $v_1$ se factoriza como
  $K\xrightarrow{v}C\to C/\mathfrak J^2$ y se tiene $v\ne u$.
  Recordemos ((20.6.20) y (21.4.4)) que la condición $\Omega^1_K=0$ significa
  que $K$ es perfecto si es de característica $\ne0$, y una extensión
  algebraica de $\mathbf Q$ si es de característica $0$.
  \item[(ii)] De la misma manera, sean $W$ un anillo de Cohen, $C$ un anillo
  local noetheriano completo, y sea $\mathfrak J\ne0$ un ideal de $C$ contenido en
  el ideal máximo; para que la factorización $W\to C\to C/\mathfrak J$
  en (19.8.6, (i)) sea única, es necesario y suficiente que el cuerpo residual
  $K$ del anillo de Cohen $W$ satisfaga $\Omega^1_K=0$. En efecto, si
  $\Omega^1_K\ne0$ y si $\mathfrak m$ designa el ideal máximo de $W$, basta
  componer una derivación $D\ne0$ de $K$ en
  $\mathfrak J/\mathfrak m\mathfrak J$ con la augmentación $W\to K$ para
  obtener una derivación $D'\ne0$ de $W$ en
  $\mathfrak J/\mathfrak m\mathfrak J$ y se termina el razonamiento como
  en (i), formando mediante $D'$ un homomorfismo
  $v_1:W\to C/\mathfrak m\mathfrak J$ distinto del homomorfismo
  $u_1:W\to C/\mathfrak m\mathfrak J$ obtenido de $u$ por paso al cociente;
  se concluye invocando esta vez (19.8.6, (i)). Si, por el contrario,
  $\Omega^1_K=0$, la unicidad de $v$ ya resulta de (i) cuando $W=K$ es un
  cuerpo de característica $0$. En el caso contrario, se tiene
  $\widehat{\Omega}^1_W=0$; en efecto, entonces $\mathfrak m=pW$, siendo $p$
  la característica de $K$ (19.8.5), y el homomorfismo canónico
  $\mathfrak m/\mathfrak m^2\to\Omega^1_W/\mathfrak m\Omega^1_W$ (20.5.11.2)
  es por tanto nulo. La sucesión exacta (20.5.12.1) aplicada a $W$ y a
  $K=W/\mathfrak m$ implica entonces que $\Omega^1_W=\mathfrak m\Omega^1_W$,
  de donde se sigue nuestra afirmación. Pero entonces (20.7.4) $W$ es
  formalmente no ramificado (para su topología adítica) sobre $\mathbf Z$, y la
  unicidad de $v$ se demuestra como en (i).
\end{enumerate}
\end{remarks}

\oldpage[0\textsubscript{IV-1}]{167}
\subsection{Cuerpos admisibles para una extensión}
\label{section:0.21.6-es}

\begin{env}[21.6.1]
\label{0.21.6.1-es}
Dados cuatro cuerpos $k_0\subset k\subset K\subset L$, se deduce de
(20.6.16) y (20.6.17) que se tiene una sucesión exacta
\begin{equation}
\tag{21.6.1.1}
\label{0.21.6.1.1-es}
0\to\Upsilon_{K/k/k_0}\otimes_KL\xrightarrow{v}\Upsilon_{L/k/k_0}
\xrightarrow{u}\Upsilon_{L/K/k_0}\xrightarrow{s}\Upsilon_{L/K/k}\to0.
\end{equation}

Cuando fijamos $k_0$, $K$ y $L$ y hacemos «variar» el cuerpo intermedio
$k$ entre $k_0$ y $K$, evidentemente se tiene
$\Upsilon_{L/K/k_0}=\Upsilon_{L/K/k}$ cuando $k=k_0$. Cuando el homomorfismo
canónico $s$ de (21.6.1.1) es además biyectivo, decimos que $k$ es un cuerpo
$k_0$-admisible para la extensión $L$ de $K$. El interés de la existencia, bajo
ciertas condiciones, de tales cuerpos $k$, que sean sin embargo
«suficientemente cercanos» a $K$ (por ejemplo, tales que $[K:k]$ sea finito), es
que permiten sustituir en ciertas cuestiones los módulos de diferenciales
$\Omega^1_{K/k_0}$ y $\Omega^1_{L/k_0}$ (que pueden ser «demasiado grandes»,
por ejemplo cuando $k_0$ es el cuerpo primo) por
$\Omega^1_{K/k}$ y $\Omega^1_{L/k}$, más fáciles de manejar.

Cuando $k_0$ es el \emph{cuerpo primo}, diremos «cuerpo admisible» en lugar
de «cuerpo $k_0$-admisible».
 
Se pondrá
\end{env}

\begin{env}[21.6.1.2]
\label{0.21.6.1.2-es}
\[
\Delta(L/K,k/k_0)=\operatorname{Coker}
(\Upsilon_{K/k/k_0}\otimes_KL\to\Upsilon_{L/k/k_0})
\simeq\operatorname{Ker}(\Upsilon_{L/K/k_0}\to\Upsilon_{L/K/k})
\]
(espacio vectorial sobre $L$); su rango se denotará $d(L/K,k/k_0)$ y se llamará
\emph{defecto de $k_0$-admisibilidad de $k$ para la extensión $L$ de $K$}
(es evidentemente nulo si y solo si $k$ es $k_0$-admisible para esta
extensión). Cuando $k_0$ es el cuerpo primo, escribiremos
$\Delta(L/K,k)$ y $d(L/K,k)$ en lugar de $\Delta(L/K,k/k_0)$ y
$d(L/K,k/k_0)$.
\end{env}

\begin{proposition}[21.6.2]
\label{0.21.6.2-es}
Sean $k_0\subset k\subset K\subset L$ cuatro cuerpos.
\begin{enumerate}
  \item[\rm (i)] Las condiciones siguientes son equivalentes:
  \begin{enumerate}
    \item[\rm a)] El cuerpo $k$ es $k_0$-admisible para la extensión $L$ de $K$
    (en otras palabras, el homomorfismo $\Upsilon_{L/K/k_0}\to
    \Upsilon_{L/K/k}$ es inyectivo y, por tanto, biyectivo).
    \item[\rm b)] El homomorfismo canónico
    $u:\Upsilon_{K/k/k_0}\to\Upsilon_{L/k/k_0}$ es nulo.
    \item[\rm c)] El homomorfismo canónico
    $v:\Upsilon_{K/k/k_0}\otimes_KL\to\Upsilon_{L/k/k_0}$ es sobreyectivo
    (y, por tanto, biyectivo).
    \item[\rm d)] Se tiene $d(L/K,k/k_0)=0$ (o
    $\Delta(L/K,k/k_0)=0$).
  \end{enumerate}
  \item[\rm (ii)] Las condiciones equivalentes de (i) se verifican cuando
  se está en uno de los casos siguientes: $\alpha)$ $L$ es separable sobre $k$;
  $\beta)$ $L$ es separable sobre $K$; $\gamma)$ se tiene
  $k\subset k_0(K^p)$, designando por $p$ el exponente de característica de $k_0$.
\end{enumerate}
\end{proposition}

\begin{proof}
\begin{enumerate}
  \item[(i)] Las afirmaciones se deducen trivialmente de la exactitud de la sucesión
  (21.6.1.1).
  \item[(ii)] Si $L$ es separable sobre $k$, se tiene
  $\Upsilon_{L/k/k_0}=0$ (20.6.19), por lo que se cumple la condición b) de (i);
  si $L$ es separable sobre $K$, se tiene
  $\Upsilon_{L/K/k_0}=0$ (20.6.19), por lo que se cumple la condición a) de (i);
  finalmente, si se tiene $k\subset k_0(K^p)$, se deduce que
  $\Upsilon_{L/K/k_0}=\Upsilon_{L/K/k}$ en virtud de (21.1.5.1), y se cumple
  la condición a) de (i).
\end{enumerate}
\end{proof}
\oldpage[0\textsubscript{IV-1}]{168}
\begin{env}[21.6.3]
\label{0.21.6.3-es}
Supongamos que tenemos un diagrama conmutativo de monomorfismos de cuerpos
\[
\xymatrix{
 k'_0\ar[r]&k'\ar[r]&K'\ar[r]&L'\\
 k_0\ar[u]\ar[r]&k\ar[u]\ar[r]&K\ar[u]\ar[r]&L\ar[u]
}
\]
Se deduce entonces de (20.6.17, (ii)) que se tiene un homomorfismo canónico
\begin{equation}
\tag{21.6.3.1}
\label{0.21.6.3.1-es}
\Delta(L/K,k/k_0)\longrightarrow\Delta(L'/K',k'/k'_0)
\end{equation}
con una propiedad de transitividad evidente, de modo que podemos decir que
$\Delta(L/K,k/k_0)$ es un \emph{functor} en el cuádruple
$(k_0,k,K,L)$.
\end{env}

\begin{proposition}[21.6.4]
\label{0.21.6.4-es}
\begin{enumerate}
  \item[\rm (i)] Sean $k\subset k'\subset k''\subset K\subset L$ cinco cuerpos.
  Se tiene una sucesión exacta de homomorfismos canónicos
\begin{equation}
\tag{21.6.4.1}
\label{0.21.6.4.1-es}
0\to\Delta(L/K,k'/k)\to\Delta(L/K,k''/k)\to\Delta(L/K,k''/k')\to0
\end{equation}
y, por consiguiente, la igualdad
\begin{equation}
\tag{21.6.4.2}
\label{0.21.6.4.2-es}
d(L/K,k''/k)=d(L/K,k''/k')+d(L/K,k'/k).
\end{equation}
  \item[\rm (ii)] Sean $k_0\subset k\subset K\subset L\subset M$ cinco cuerpos.
  Se tiene una sucesión exacta de homomorfismos canónicos
\begin{equation}
\tag{21.6.4.3}
\label{0.21.6.4.3-es}
0\to\Delta(L/K,k/k_0)\otimes_LM\to\Delta(M/K,k/k_0)
\to\Delta(M/L,k/k_0)\to0
\end{equation}
y, por consiguiente, la igualdad
\begin{equation}
\tag{21.6.4.4}
\label{0.21.6.4.4-es}
d(M/K,k/k_0)=d(M/L,k/k_0)+d(L/K,k/k_0).
\end{equation}
\end{enumerate}
\end{proposition}

\begin{proof}
\begin{enumerate}
  \item[(i)] Consideremos el diagrama conmutativo
\[
\xymatrix{
0\ar[r]&\Upsilon_{K/k'/k}\otimes_KL\ar[r]\ar[d]&
\Upsilon_{K/k''/k}\otimes_KL\ar[r]\ar[d]&
\Upsilon_{K/k''/k'}\otimes_KL\ar[r]\ar[d]&0\\
0\ar[r]&\Upsilon_{L/k'/k}\ar[r]\ar[d]&
\Upsilon_{L/k''/k}\ar[r]\ar[d]&
\Upsilon_{L/k''/k'}\ar[r]\ar[d]&0\\
0\ar[r]&\Delta(L/K,k'/k)\ar[r]&\Delta(L/K,k''/k)\ar[r]&
\Delta(L/K,k''/k')\ar[r]&0
}
\]
donde podemos considerar las tres filas como \emph{complejos}
$T_1,T_2,T_3$, respectivamente; la sucesión exacta (21.6.1.1) y la definición
de $\Delta$ (21.6.1.2) muestran que se tiene una sucesión exacta de complejos
$0\to T_1\to T_2\to T_3\to0$; apliquémosle la sucesión exacta de
\oldpage[0\textsubscript{IV-1}]{169}
cohomología, y observemos que, en virtud de la exactitud de (21.6.1.1), la
cohomología de $T_1$ y la de $T_2$ son nulas salvo en un mismo grado, para el
cual los módulos de cohomología son ambos iguales a
$\Upsilon_{k''/k'/k}\otimes_{k''}L$; como $T_1$ y $T_2$ tienen así la misma
cohomología, la de $T_3$ es necesariamente nula, lo que prueba (i).
  \item[(ii)] Consideremos del mismo modo el diagrama conmutativo
\[
\xymatrix{
0\ar[r]&\Delta(L/K,k/k_0)\otimes_LM\ar[r]\ar[d]&
\Delta(M/K,k/k_0)\ar[r]\ar[d]&\Delta(M/L,k/k_0)\ar[r]&0\\
0\ar[r]&\Upsilon_{L/K/k_0}\otimes_LM\ar[r]\ar[d]&
\Upsilon_{M/K/k_0}\ar[r]\ar[d]&\Upsilon_{M/L/k_0}\ar[r]&0\\
0\ar[r]&\Upsilon_{L/K/k}\otimes_LM\ar[r]&
\Upsilon_{M/K/k}\ar[r]&\Upsilon_{M/L/k}\ar[r]&0
}
\]
donde de nuevo consideramos las tres filas como \emph{complejos}
$T'_1,T'_2,T'_3$; la sucesión exacta (21.6.1.1) y la definición de $\Delta$
(21.6.1.2) dan aquí una sucesión exacta de complejos
$0\to T'_1\to T'_2\to T'_3\to0$, a la que aplicamos nuevamente la sucesión
exacta de cohomología; esta vez, en virtud de la exactitud de (21.6.1.1), la
cohomología de $T'_2$ y la de $T'_3$ son nulas salvo en un mismo grado, para el
cual los módulos de cohomología son ambos iguales a
$\Upsilon_{M/L/k_0}$; concluimos por tanto aquí que la cohomología de $T'_1$ es
nula, lo que establece (ii).
\end{enumerate}
\end{proof}

\begin{corollary}[21.6.5]
\label{0.21.6.5-es}
\begin{enumerate}
  \item[(i)] Dados cinco cuerpos $k\subset k'\subset k''\subset K\subset L$,
  para que $k''$ sea $k$-admisible para la extensión $L$ de $K$, es necesario y
  suficiente que $k'$ sea $k$-admisible y que $k''$ sea $k'$-admisible
  para la extensión $L$ de $K$.
  \item[(ii)] Dados cinco cuerpos
  $k_0\subset k\subset K\subset L\subset M$, para que $k$ sea
  $k_0$-admisible para la extensión $M$ de $K$, es necesario y suficiente que lo
  sea para la extensión $L$ de $K$ y para la extensión $M$ de $L$.
\end{enumerate}

Esto se deduce inmediatamente de las relaciones (21.6.4.2) y (21.6.4.4), siendo
los valores de $d$ $\geqslant0$.
\end{corollary}

\begin{corollary}[21.6.6]
\label{0.21.6.6-es}
Sean $k_0\subset k\subset K\subset L$ cuatro cuerpos, y supongamos que $k$
sea $k_0$-admisible para la extensión $L$ de $K$. Entonces, si $k'_0$, $k'$,
$K'$, $L'$ son cuatro cuerpos tales que
$k_0\subset k'_0\subset k'\subset k\subset K\subset K'\subset L'\subset L$,
$k'$ es $k'_0$-admisible para la extensión $L'$ de $K'$.
\end{corollary}

\subsection{La igualdad de Cartier}
\label{subsection:0.21.7-es}

El resultado siguiente traduce en términos de diferenciales un teorema de
MacLane sobre las derivaciones:

\oldpage[0\textsubscript{IV-1}]{170}
\begin{theorem}[21.7.1]
\label{0.21.7.1-es}
\textnormal{(Cartier).}
Sean $K$ un cuerpo, $L$ una extensión de tipo finito de $K$. Entonces
$\Omega^1_{L/K}$ y $\Upsilon_{L/K}$ son espacios vectoriales de rango
finito sobre $L$, y se tiene
\begin{equation}
\tag{21.7.1.1}
\label{0.21.7.1.1-es}
\operatorname{rg}\Omega^1_{L/K}-\operatorname{rg}\Upsilon_{L/K}
=\deg.\operatorname{tr}_K L.
\end{equation}
\end{theorem}

\begin{proof}
Si $L'$ es un cuerpo tal que $K\subset L'\subset L$, se tiene la sucesión exacta
(20.6.15.1) (aplicada a $\Lambda=P$, cuerpo primo de $K$, $A=K$,
$B=L'$, $C=L$)
\[
0\to\Upsilon_{L'/K}\otimes_{L'}L\to\Upsilon_{L/K}
\to\Omega^1_{L'/K}\otimes_{L'}L\to\Omega^1_{L/K}
\to\Omega^1_{L/L'}\to0
\]
de donde (0$_{\rm III}$, 11.10.2)
\[
\operatorname{rg}_L\Omega^1_{L/K}-\operatorname{rg}_L\Upsilon_{L/K}
=(\operatorname{rg}_{L'}\Omega^1_{L/L'}-\operatorname{rg}_{L'}\Upsilon_{L/L'})
+(\operatorname{rg}_{L'}\Omega^1_{L'/K}-\operatorname{rg}_{L'}\Upsilon_{L'/K}).
\]
Como, por otra parte, $\deg.\operatorname{tr}_K L=\deg.\operatorname{tr}_K L'
+\deg.\operatorname{tr}_{L'}L$, se ve, por inducción sobre el número de
generadores de la extensión $L$, que basta demostrar (21.7.1.1) cuando
$L=K(x)$. Distingamos entonces tres casos:
\begin{enumerate}
  \item[\rm a)] $x$ es trascendente sobre $K$; como $L$ es separable sobre
  $K$, se tiene $\Upsilon_{L/K}=0$ (20.6.3); por otra parte, (20.5.9) y
  (20.4.13) muestran que $\Omega^1_{L/K}$ tiene rango $1$, de donde se
  obtiene (21.7.1.1) en este caso.
  \item[\rm b)] $L$ es una extensión algebraica separable de $K$, de modo que
  se tiene de nuevo $\Upsilon_{L/K}=0$ (20.6.3). Por otra parte,
  $\Omega^1_{L/K}=0$ por (20.6.20), y se obtiene de nuevo (21.7.1.1).
  \item[\rm c)] $L$ es una extensión algebraica inseparable de $K$; el
  razonamiento inicial muestra que podemos limitarnos al caso $L^p\subset K$;
  se deduce entonces de (21.4.8) que
  $\operatorname{rg}_L\Upsilon_{L/K}=\operatorname{rg}_L\Omega^1_{L/K}$,
  de donde se obtiene de nuevo (21.7.1.1).
\end{enumerate}
\end{proof}

\begin{corollary}[21.7.2]
\label{0.21.7.2-es}
Sean $K$ un cuerpo, $L$ una extensión de tipo finito de $K$, $k$ un subcuerpo
de $K$. Entonces $\Omega^1_{L/K}$ y $\Upsilon_{L/K/k}$ son espacios
vectoriales de rango finito sobre $L$, y se tiene
\begin{equation}
\tag{21.7.2.1}
\label{0.21.7.2.1-es}
\operatorname{rg}\Omega^1_{L/K}-\operatorname{rg}\Upsilon_{L/K/k}
=\deg.\operatorname{tr}_K L+d(L/K,k).
\end{equation}
Por consiguiente, se tiene la desigualdad
\begin{equation}
\tag{21.7.2.2}
\label{0.21.7.2.2-es}
\operatorname{rg}\Omega^1_{L/K}-\operatorname{rg}\Upsilon_{L/K/k}
\geqslant\deg.\operatorname{tr}_K L.
\end{equation}
Además, para que los dos miembros de (21.7.2.2) sean iguales, es necesario y
suficiente que $k$ sea un cuerpo admisible para la extensión $L$ de $K$.
\end{corollary}

\begin{proof}
En efecto, como el homomorfismo
$\Upsilon_{L/K}\to\Upsilon_{L/K/k}$ es sobreyectivo, se tiene por definición
(21.6.1.2) $\operatorname{rg}\Upsilon_{L/K}-\operatorname{rg}\Upsilon_{L/K/k}
=d(L/K,k)$, y el corolario se deduce inmediatamente de (21.7.1) y de
(21.6.2).
\end{proof}

\begin{corollary}[21.7.3]
\label{0.21.7.3-es}
Sean $k\subset K\subset L$ tres cuerpos tales que $L$ sea una extensión de
tipo finito de $k$. Entonces se tiene
\begin{equation}
\tag{21.7.3.1}
\label{0.21.7.3.1-es}
\operatorname{rg}_L\Omega^1_{L/k}-\operatorname{rg}_K\Omega^1_{K/k}
=\deg.\operatorname{tr}_K L+d(L/K,k).
\end{equation}
\oldpage[0\textsubscript{IV-1}]{171}
Por consiguiente, se tiene la desigualdad
\begin{equation}
\tag{21.7.3.2}
\label{0.21.7.3.2-es}
\operatorname{rg}_L\Omega^1_{L/k}-\operatorname{rg}_K\Omega^1_{K/k}
\geqslant\deg.\operatorname{tr}_K L.
\end{equation}
La igualdad se alcanza si y solo si $k$ es un cuerpo admisible para
la extensión $L$ de $K$.
\end{corollary}

\begin{proof}
En efecto, se tiene la sucesión exacta (20.6.1.1)
\[
0\to\Upsilon_{L/K/k}\to\Omega^1_{K/k}\otimes_KL\to\Omega^1_{L/k}
\to\Omega^1_{L/K}\to0
\]
por lo que el primer miembro de (21.7.3.1) es igual a
\[
\operatorname{rg}_L\Omega^1_{L/k}-\operatorname{rg}_L\Upsilon_{L/K/k}
\]
y basta aplicar (21.7.2).
\end{proof}

El interés de este último corolario es que solo intervienen en él módulos de
diferenciales, con exclusión de los módulos de imperfección. Veremos además más
abajo (21.8.6) que para toda extensión de tipo finito $L$ de $K$ existe un
subcuerpo $k$ de $K$ tal que $[K:k]$ sea finito y que sea admisible para la
extensión $L$ de $K$, de modo que los dos miembros de (21.7.3.2) sean entonces
iguales.

\begin{corollary}[21.7.4]
\label{0.21.7.4-es}
Sea $K$ una extensión de tipo finito de un cuerpo $k$. Las condiciones siguientes
son equivalentes:
\begin{enumerate}
  \item[\rm a)] $K$ es una extensión finita separable de $k$.
  \item[\rm b)] $K$ es una $k$-álgebra formalmente no ramificada (19.10.2).
  \item[\rm c)] $K$ es una $k$-álgebra formalmente étale (19.10.2).
  \item[\rm d)] Se tiene $\Omega^1_{K/k}=0$.
\end{enumerate}
\end{corollary}

\begin{proof}
En efecto, se sabe (siendo discretas las topologías) que las condiciones b) y
d) son equivalentes (20.7.4) y que a) implica que $K$ es una $k$-álgebra
formalmente lisa (19.6.1), por lo que la conjunción de a) y b) equivale a c).
Todo se reduce a ver que d) implica a). Ahora bien, en virtud de (21.7.1.1),
la relación $\Omega^1_{K/k}=0$ implica que $K$ es algebraico (y por tanto finito)
sobre $k$ y que $\Upsilon_{K/k}=0$, es decir (20.6.3), que $K$ es separable
sobre $k$.
\end{proof}

\begin{remark}[21.7.5]
\label{0.21.7.5-es}
En virtud de (20.6.3.2) y de (0$_{\rm III}$, 11.10.2), el primer miembro de
(21.7.1.1) no es otro que la \emph{característica de Euler-Poincaré} del
complejo $K_\bullet(C/B/A)$ introducido en (20.6.3), para $C=L$, $B=K$ y
$A=P$ (cuerpo primo de $K$). En el capítulo dedicado a la teoría de
intersecciones y al teorema de Riemann--Roch, también desempeñarán un papel
importante las características de Euler-Poincaré generalizadas (con valores
en grupos de clases de $\mathcal O_X$-Módulos) de complejos que generalizan
los complejos $F_\bullet(C/A)$ considerados en (20.6.22).
\end{remark}

\subsection{Criterios de admisibilidad}
\label{subsection:0.21.8-es}

Volvemos a nuestras convenciones anteriores y suponemos, por tanto, que todos
los cuerpos considerados en este número son de característica $p>0$.

\begin{lemma}[21.8.1]
\label{0.21.8.1-es}
Sean $K$ un cuerpo, $k$ un subcuerpo de $K$, $(k_\alpha)_{\alpha\in I}$ una
familia filtrante decreciente de subcuerpos de $K$ tal que
$k=\bigcap_{\alpha\in I}k_\alpha$. Sean $V$ un espacio vectorial sobre $K$,
$(a_i)_{1\leqslant i\leqslant n}$ una familia finita de vectores de $V$; si la
familia $(a_i)$ es libre sobre $k$, existe un índice $\gamma$ tal que también
sea libre sobre $k_\gamma$.
\end{lemma}

\oldpage[0\textsubscript{IV-1}]{172}
\begin{proof}
Sea $r$ el rango de la familia $(a_i)_{1\leqslant i\leqslant n}$ sobre $K$, y
razonemos por inducción sobre $n-r$; la proposición es evidente para $n=r$,
pues entonces la familia $(a_i)_{1\leqslant i\leqslant r}$ es libre sobre $K$, y
por tanto sobre todo subcuerpo de $K$. Supongamos, por ejemplo, que
$(a_i)_{1\leqslant i\leqslant r}$ sea libre sobre $K$, y escribamos
$a_{r+1}=\sum_{i=1}^r\lambda_i a_i$ con $\lambda_i\in K$; como la familia
$(a_i)_{1\leqslant i\leqslant r+1}$ es libre sobre $k$, los $\lambda_i$ no
pueden pertenecer todos a $k$; supongamos, por ejemplo, que
$\lambda_{i_0}\notin k$. Entonces existe un índice $\beta$ tal que
$\lambda_{i_0}\notin k_\beta$; concluimos que la familia
$(a_i)_{1\leqslant i\leqslant r+1}$ es libre sobre $k_\beta$; en efecto, como la
familia $(a_i)_{1\leqslant i\leqslant r}$ es libre sobre todo subcuerpo de
$K$, si la familia $(a_i)_{1\leqslant i\leqslant r+1}$ no fuera libre sobre
$k_\beta$, $a_{r+1}$ sería igual a una combinación lineal de los $a_i$ con
coeficientes en $k_\beta$, y como estos coeficientes son necesariamente los
$\lambda_i$, se llega a una contradicción. Basta ahora aplicar la hipótesis de
inducción sustituyendo $K$ por $k_\beta$ y la familia $(k_\alpha)_{\alpha\in I}$
por la subfamilia de los $k_\alpha$ contenidos en $k_\beta$.
\end{proof}

\begin{lemma}[21.8.2]
\label{0.21.8.2-es}
Sean $K$ un cuerpo, $p$ su exponente de característica, $k_0$ un subcuerpo de
$K$, $(k_\alpha)_{\alpha\in I}$ una familia filtrante decreciente de
subcuerpos de $K$ tal que $\bigcap_{\alpha\in I}k_\alpha(K^p)=k_0(K^p)$.
Si $(x_i)_{1\leqslant i\leqslant n}$ es una familia finita de elementos de $K$
que es $p$-independiente sobre $k_0$, existe un índice $\alpha$ tal que $(x_i)$ sea
$p$-independiente sobre $k_\alpha$.

En efecto, decir que la familia $(x_i)$ es $p$-independiente sobre un subcuerpo $k$
de $K$ significa que la familia finita de monomios
$\prod_{i=1}^n x_i^{m(i)}$ con $0\leqslant m(i)<p$ es libre sobre $k(K^p)$;
basta, pues, aplicar el lema (21.8.1) a esta familia de monomios en el
espacio vectorial $V=K$, y a los subcuerpos $k_\alpha(K^p)$ y $k_0(K^p)$
de $K$.
\end{lemma}

\begin{theorem}[21.8.3]
\label{0.21.8.3-es}
Sean $K$ un cuerpo de característica $p>0$, $k_0$ un subcuerpo de $K$,
$(k_\alpha)_{\alpha\in I}$ una familia filtrante decreciente de subcuerpos
de $K$ que contienen a $k_0$. Las condiciones siguientes son equivalentes:
\begin{enumerate}
  \item[\rm a)] $\bigcap_{\alpha\in I}k_\alpha(K^p)=k_0(K^p)$.
  \item[\rm b)] Para toda extensión $L$ de $K$ tal que
  $\Upsilon_{L/K/k_\alpha}$ sea un espacio vectorial sobre $L$ de rango finito
  (lo que ocurre en particular si $L$ es una extensión de tipo finito en
  virtud de (21.7.2)), existe $\alpha\in I$ tal que $k_\alpha$ sea
  $k_0$-admisible para la extensión $L$ de $K$.
  \item[\rm b$'$)] Para toda extensión $L=K(x)$ de $K$, con $x^p\in K$,
  existe $\alpha\in I$ tal que $k_\alpha$ sea $k_0$-admisible para
  la extensión $L$ de $K$.
  \item[\rm c)] La aplicación canónica
\begin{equation}
\tag{21.8.3.1}
\label{0.21.8.3.1-es}
\Omega^1_{K/k_0}\longrightarrow\varprojlim_{\alpha}\Omega^1_{K/k_\alpha}
\end{equation}
es inyectiva.
\end{enumerate}
\end{theorem}

\begin{proof}
La aplicación canónica (21.8.3.1) se obtiene, naturalmente, pasando al límite
proyectivo en el sistema proyectivo de homomorfismos
$\Omega^1_{K/k_0}\to\Omega^1_{K/k_\alpha}$ (20.5.3.3). Probaremos el
teorema según el esquema lógico $c)\Rightarrow b)\Rightarrow b')
\Rightarrow a)\Rightarrow c)$.

Decir que $k_\alpha$ es $k_0$-admisible para la extensión $L$ de $K$ significa
que el homomorfismo
\oldpage[0\textsubscript{IV-1}]{173}
canónico $\Upsilon_{L/K/k_0}\to\Upsilon_{L/K/k_\alpha}$ es inyectivo; ahora
bien, si $N_\alpha$ es el núcleo de este homomorfismo, los $N_\alpha$ forman un
sistema filtrante decreciente de subespacios vectoriales de
$\Upsilon_{L/K/k_0}$, y como $\Upsilon_{L/K/k_0}$ es por hipótesis de rango
finito, equivale a decir que uno de los $N_\alpha$ es nulo o que su
intersección es nula. Pero esta intersección no es sino el núcleo del
homomorfismo límite proyectivo
$\Upsilon_{L/K/k_0}\to\varprojlim_\alpha\Upsilon_{L/K/k_\alpha}$. Ahora bien,
se tiene el diagrama conmutativo
\[
\xymatrix{
\Upsilon_{L/K/k_0}\ar[r]\ar[d]&\varprojlim\Upsilon_{L/K/k_\alpha}\ar[d]\\
\Omega^1_{K/k_0}\otimes_KL\ar[r]&
\varprojlim(\Omega^1_{K/k_\alpha}\otimes_KL)
}
\]
donde la flecha vertical de la izquierda es inyectiva por definición. Para
demostrar que $c)$ implica $b)$, basta, por tanto, mostrar que la aplicación
canónica
\begin{equation}
\tag{21.8.3.2}
\label{0.21.8.3.2-es}
\Omega^1_{K/k_0}\otimes_KV\longrightarrow
\varprojlim_{\alpha}(\Omega^1_{K/k_\alpha}\otimes_KV)
\end{equation}
es inyectiva para todo espacio vectorial $V$ sobre $K$ (y en particular para
$V=L$). Ahora bien, esto es evidente si $V=K^n$, pues entonces
$W\otimes_KV=W^n$ para todo espacio vectorial $W$ sobre $K$ y los productos y
los límites proyectivos conmutan. Por otra parte, para todo elemento $z$ del
primer miembro de (21.8.3.2), existe un subespacio $V'$ de $V$ de rango finito
tal que $z\in\Omega^1_{K/k_0}\otimes_KV'$ (identificado canónicamente con un
subespacio de $\Omega^1_{K/k_0}\otimes_KV$); si $z\ne0$, su imagen en
$\varprojlim_\alpha(\Omega^1_{K/k_\alpha}\otimes_KV')$ no es, por tanto, nula;
como el funtor $\varprojlim$ es exacto por la izquierda, la imagen de $z$ en el
segundo miembro de (21.8.3.2) también es $\ne0$, lo que termina de demostrar
que $c)$ implica $b)$.

Es trivial que $b)$ implica $b')$; mostremos que $b')$ implica $a)$. Con las
notaciones de $b')$, podemos suponer que $L\ne K$. Resulta entonces de (21.4.7)
que se tiene $\operatorname{rg}\Upsilon_{L/K}\leqslant1$. Si
$\Upsilon_{L/K/k_\alpha}=0$, no hay nada que demostrar en virtud de
(21.6.1.1). En caso contrario, $\Upsilon_{L/K/k_\alpha}$ se identifica
canónicamente con $\Upsilon_{L/K}$ y, si se pone $a=x^p$, $\Upsilon_{L/K}$
tiene una base formada por el único elemento $d_K(a)\otimes1$ (21.4.7);
$\Upsilon_{L/K/k_\alpha}$ tiene, por tanto, una base formada por el único
elemento $d_{K/k_\alpha}(a)\otimes1$, y decir que un subcuerpo $k_\alpha$ es
$k_0$-admisible para la extensión $L$ de $K$ significa que se tiene
$d_{K/k_\alpha}(a)\ne0$, o también (21.4.5) que
$a\notin k_\alpha(K^p)$. Ahora bien, para todo $a\notin k_0(K^p)$, se tiene
$d_{K/k_0}(a)\ne0$ y $K(a^{1/p})\ne K$, por lo que se puede aplicar $b')$, que
demuestra la existencia de un $\alpha$ tal que $a\notin k_\alpha(K^p)$; en
otras palabras, $b')$ implica $a)$.

Queda por mostrar que $a)$ implica $c)$. Sea $(x_\mu)_{\mu\in M}$ una
$p$-base de $K$ sobre $k_0$; entonces, los $d_{K/k_0}(x_\mu)$ forman una base
del $K$-espacio vectorial $\Omega^1_{K/k_0}$ (21.4.5), y la condición $c)$
significa que para toda parte finita $J$ del conjunto de índices $M$, existe
un $\alpha\in I$ tal que los $d_{K/k_\alpha}(x_\mu)$, para $\mu\in J$, sean
linealmente independientes en $\Omega^1_{K/k_\alpha}$; pero esto significa
también (21.4.5) que los $x_\mu$, para $\mu\in J$, forman una familia
$p$-libre sobre $k_\alpha$, y la existencia de un $\alpha$ con esta propiedad
resulta de (21.8.2) y de la hipótesis $a)$.
\end{proof}
\oldpage[0\textsubscript{IV-1}]{174}

\begin{corollary}[21.8.4]
\label{0.21.8.4-es}
Sean $K$ un cuerpo de característica $p>0$, $k_0$ un subcuerpo de $K$ y
$(k_\alpha)_{\alpha\in I}$ una familia filtrante decreciente de subcuerpos de
$K$, que contienen a $k_0$, tal que
\[
  \bigcap_{\alpha\in I}k_\alpha(K^p)=k_0(K^p).
\]
Sea $L$ una extensión de $K$ tal que $\Upsilon_{L/K/k_0}$ sea un
$L$-espacio vectorial de rango finito; entonces, para todo cuerpo $K'$ tal que
$K\subset K'\subset L$, existe $\alpha\in I$ tal que $k_\alpha$ sea
$k_0$-admisible para la extensión $L$ de $K'$.
\end{corollary}

\begin{proof}
Basta aplicar (21.8.3) y (21.6.6).
\end{proof}

\begin{corollary}[21.8.5]
\label{0.21.8.5-es}
Sean $K$ un cuerpo de característica $p>0$, $k_0$ un subcuerpo de $K$ y
$(k_\alpha)_{\alpha\in I}$ una familia filtrante decreciente de subcuerpos de
$K$ que contienen a $k_0$, tal que
\[
  \bigcap_{\alpha\in I}k_\alpha(K^p)=k_0(K^p).
\]
Entonces, para toda extensión $L$ de $K$ tal que $\Upsilon_{L/K/k_0}$ sea un
$L$-espacio vectorial de rango finito, se tiene
\[
  \bigcap_{\alpha\in I}k_\alpha(L^p)=k_0(L^p).
\]
\end{corollary}

\begin{proof}
Supongamos, en efecto, que $M$ sea una extensión de $L$ tal que
$\Upsilon_{M/L/k_0}$ sea un $M$-espacio vectorial de rango finito. La sucesión
exacta (21.6.1.1)
\[
  0\longrightarrow \Upsilon_{L/K/k_0}\otimes_L M
  \longrightarrow \Upsilon_{M/K/k_0}
  \longrightarrow \Upsilon_{M/L/k_0}
\]
y la hipótesis muestran entonces que $\Upsilon_{M/K/k_0}$ también es un
$M$-espacio vectorial de rango finito. Existe, por tanto, un índice $\alpha$
tal que $k_\alpha$ sea $k_0$-admisible para la extensión $M$ de $K$ (21.8.3),
y, por consiguiente, también para la extensión $M$ de $L$ (21.6.6); como esto
tiene lugar para toda extensión $M$ de $L$ tal que $\Upsilon_{M/L/k_0}$ sea de
rango finito, el corolario resulta de la equivalencia de $a)$ y $b)$ en
(21.8.3).
\end{proof}

\begin{corollary}[21.8.6]
\label{0.21.8.6-es}
Sean $K$ un cuerpo, $p$ su exponente característico y $k_0$ un subcuerpo de
$K$. Si $L$ es una extensión de $K$ tal que $\Upsilon_{L/K/k_0}$ sea un
$L$-espacio vectorial de rango finito, existe un subcuerpo $k$ de $K$, que
contiene a $k_0(K^p)$, tal que $[K:k]$ sea finito y que sea $k_0$-admisible
para la extensión $L$ de $K$.
\end{corollary}

\begin{proof}
Basta, en efecto, en virtud de (21.8.3), construir una familia filtrante
decreciente $(k_\alpha)_{\alpha\in I}$ de subcuerpos de $K$, que contengan a
$K^p$ y a $k_0$, para los cuales $[K:k_\alpha]<+\infty$ y
$\bigcap_\alpha k_\alpha=k_0(K^p)$. Para ello se considera una $p$-base
$(x_\lambda)_{\lambda\in J}$ de $K$ sobre $k_0$ y, para toda parte finita
$H$ de $J$, se considera el subcuerpo $k_H$ de $K$ engendrado por
$k_0(K^p)$ y los $x_\lambda$ de índice $\lambda\in J-H$; resulta de esta
definición que $(x_\lambda)_{\lambda\in H}$ es una $p$-base de $K$ sobre
$k_H$, y se concluye inmediatamente que los $k_H$ verifican las condiciones
requeridas.
\end{proof}

\begin{remarks}[21.8.7]
\label{0.21.8.7-es}
\begin{enumerate}
  \item[(i)] Ya se ha visto (21.7.2) que si $L$ es una extensión de
  \emph{tipo finito} de $K$, $\Upsilon_{L/K/k_0}$ es de rango finito para
  todo subcuerpo $k_0$ de $K$. Lo mismo sucede si $L$ es una extensión
  \emph{separable} de $K$, pues, en virtud de (20.6.19), se tiene
  $\Upsilon_{L/K/k_0}=0$. Finalmente, si $L$ es una \emph{extensión de tipo
  finito de una extensión separable $L_0$ de $K$}, el mismo razonamiento que
  en (21.8.5) muestra que $\Upsilon_{L/K/k_0}$ sigue siendo de rango finito
  (y de hecho es isomorfo a un subespacio de $\Upsilon_{L/L_0/k_0}$).

  \oldpage[0\textsubscript{IV-1}]{175}
  \item[(ii)] En el enunciado de (21.8.5), y por tanto también en el de
  (21.8.3, $b)$), no se puede omitir la hipótesis de que
  $\Upsilon_{L/K/k_0}$ es de rango finito sobre $L$. Tomemos como $k_0$ el
  cuerpo primo $\mathbf F_p$, y como $K$ un cuerpo tal que $[K:K^p]$ sea
  infinito numerable (por ejemplo, el cuerpo de fracciones racionales
  $\mathbf F_p(X_1,\ldots,X_n,\ldots)$ en una infinidad de indeterminadas); el
  procedimiento de construcción de (21.8.6) muestra inmediatamente que existe
  una sucesión infinita estrictamente decreciente $(K_n)$ de subcuerpos de $K$
  tal que $K_0=K$ y $\bigcap_nK_n=K^p$. Vamos a construir una sucesión
  creciente de extensiones finitas $M_n$ de $K$, tal que, si $M$ es la unión
  de los $M_n$, la extensión $L=M^{1/p}$ de $K$ contradiga (21.8.5). Para
  ello, sea $x$ un elemento de $K-K^p$, pongamos $M_0=K^p$, y
  $M_n=K^p(a_1x,a_2x,\ldots,a_nx)$ para $n\geqslant1$, donde los $a_n$ se
  construyen por recurrencia de modo que $a_n\in K_n$ y
  $a_{n+1}\notin M_n(x)$ para todo $n$: esto es posible, pues $M_n(x)$ es de
  grado finito sobre $K^p$, mientras que $K_n$ no lo es. Se concluye
  inmediatamente que $x\notin M=L^p$, pero como $x=(a_nx)a_n^{-1}$, se tiene
  $x\in K_n(M)=K_n(L^p)$ para todo $n$.
\end{enumerate}
\end{remarks}

\begin{proposition}[21.8.8]
\label{0.21.8.8-es}
Sea $k$ un cuerpo de característica $p>0$, y sean
$A=k[[T_1,\ldots,T_r]]$ el anillo de series formales en $r$ indeterminadas
sobre $k$ y $K=k((T_1,\ldots,T_r))$ su cuerpo de fracciones. Entonces existe
una familia filtrante decreciente $(A_\alpha)$ de subanillos noetherianos de
$A$ tal que $A$ sea un $A_\alpha$-módulo libre de tipo finito para todo
$\alpha$ y que, si $K_\alpha$ es el cuerpo de fracciones de $A_\alpha$, se
tenga $\bigcap_\alpha K_\alpha=K^p$.
\end{proposition}

\begin{proof}
Se puede escribir $K^p=k^p((T_1^p,\ldots,T_r^p))$; se ha visto en la
demostración de (21.8.6) que existe una familia decreciente $(k_\alpha)$ de
subcuerpos de $k$ tal que $[k:k_\alpha]$ sea finito para todo $\alpha$ y
$\bigcap_\alpha k_\alpha=k^p$; está claro que, si se pone
$A_\alpha=k_\alpha[[T_1^p,\ldots,T_r^p]]$, $A$ es un $A_\alpha$-módulo libre
de tipo finito; todo se reduce, por tanto, a demostrar la relación
\begin{equation}
\tag{21.8.8.1}
\label{0.21.8.8.1-es}
  \bigcap_\alpha k_\alpha((T_1^p,\ldots,T_r^p))
  =k^p((T_1^p,\ldots,T_r^p)).
\end{equation}
Como $\bigcap_\alpha k_\alpha=k^p$, esto resultará de los dos lemas siguientes.
\end{proof}

\begin{lemma}[21.8.8.2]
\label{0.21.8.8.2-es}
Si $k'$ es una extensión de un cuerpo $k$, se tiene
\[
  k'[[T_1,\ldots,T_r]]\cap k((T_1,\ldots,T_r))=k[[T_1,\ldots,T_r]].
\]
\end{lemma}

\begin{proof}
En efecto, pongamos $C=k[[T_1,\ldots,T_r]]$ y
$D=k'[[T_1,\ldots,T_r]]$; como $k((T_1,\ldots,T_r))$ es el cuerpo de
fracciones de $C$, bastará demostrar que $D$ es un $C$-módulo
\emph{fielmente plano} (Bourbaki, \emph{Alg. comm.}, cap.~I, \S~3,
n\textsuperscript{o}~5, prop.~10). Ahora bien, $C$ y $D$ son anillos locales
noetherianos, y si $\mathfrak m$ es el ideal maximal de $C$, se tiene
$D/\mathfrak m^jD=(C/\mathfrak m^j)\otimes_k k'$, por lo que
$D/\mathfrak m^jD$ es un $(C/\mathfrak m^j)$-módulo plano; basta, por tanto,
aplicar ($\mathbf 0_{\mathrm{III}}$, 10.2.1) y
($\mathbf 0_{\mathrm I}$, 6.6.2).
\end{proof}

\begin{lemma}[21.8.8.3]
\label{0.21.8.8.3-es}
Sean $k$ un cuerpo, $(k_\alpha)$ una familia filtrante decreciente de
subcuerpos de $k$, y pongamos $k_0=\bigcap_\alpha k_\alpha$. Se supone que
existe una potencia $q$ del exponente característico de $k$ tal que
$k^q\subset k_0$. Entonces se tiene
\begin{equation}
\tag{21.8.8.4}
\label{0.21.8.8.4-es}
  \bigcap_\alpha k_\alpha((T_1,\ldots,T_r))=k_0((T_1,\ldots,T_r)).
\end{equation}
\end{lemma}

\begin{proof}
\oldpage[0\textsubscript{IV-1}]{176}
Basta demostrar que un elemento $f\ne0$ del primer miembro de (21.8.8.4)
pertenece al segundo miembro. Sean $\gamma$ un índice y
$g\in k_\gamma[[T_1,\ldots,T_r]]$ tales que
$gf\in k_\gamma[[T_1,\ldots,T_r]]$; reemplazando, si es necesario, $g$ por
$g^q$, se puede suponer que $g\in k_0[[T_1,\ldots,T_r]]$. Entonces, para todo
$\alpha\geqslant\gamma$, se tiene
\[
  gf\in k_\gamma[[T_1,\ldots,T_r]]\cap k_\alpha((T_1,\ldots,T_r))
  =k_\alpha[[T_1,\ldots,T_r]]
\]
en virtud del lema (21.8.8.2). Pero está claro que la intersección de los
anillos $k_\alpha[[T_1,\ldots,T_r]]$ no es sino
$k_0[[T_1,\ldots,T_r]]$, y se tiene, por tanto,
$f\in k_0((T_1,\ldots,T_r))$.
\end{proof}

\subsection{Módulos de diferenciales completados en los anillos de series formales}
\label{subsection:0.21.9-es}

En este número, ya no se supone necesariamente que los cuerpos sean de
característica $>0$.

\begin{lemma}[21.9.1]
\label{0.21.9.1-es}
Sean $k$ un cuerpo, $A$ un anillo local noetheriano completo que es una
$k$-álgebra y $K$ el cuerpo residual de $A$. Si $K$ es una extensión de tipo
finito de $k$, el $A$-módulo $\widehat\Omega^1_{A/k}$ es de tipo finito.
\end{lemma}

\begin{proof}
En virtud de (20.7.15), basta demostrar que $\Omega^1_{K/k}$ es un
$K$-espacio vectorial de rango finito. Ahora bien, por hipótesis, $K$ es el
cuerpo de fracciones de una $k$-álgebra de tipo finito $B$; como
$\Omega^1_{B/k}$ es un $B$-módulo de tipo finito en virtud de (20.4.7), la
conclusión resulta de (20.5.9).
\end{proof}

\begin{proposition}[21.9.2]
\label{0.21.9.2-es}
Sean $k_0$ un cuerpo, $A$ un anillo local noetheriano completo que es una
$k_0$-álgebra formalmente lisa, $\mathfrak m$ su ideal maximal y
$K=A/\mathfrak m$ su cuerpo residual (de modo que, como anillo, $A$ es
isomorfo a un anillo de series formales $K[[T_1,\ldots,T_r]]$ (19.6.5)).
Si $K$ es una extensión de tipo finito de $k_0$, entonces
$\widehat\Omega^1_{A/k_0}$ es un $A$-módulo libre de rango igual a
$\dim(A)+\deg.\operatorname{tr}_{k_0}K$. Además, para todo subcuerpo $k$ de
$k_0$ tal que $\Omega^1_{k_0/k}$ sea de rango finito,
$\widehat\Omega^1_{A/k}$ es un $A$-módulo libre de rango igual a
\[
  \dim(A)+\deg.\operatorname{tr}_{k_0}K
  +\operatorname{rg}_{k_0}(\Omega^1_{k_0/k}).
\]
\end{proposition}

\begin{proof}
Está claro que, si $P$ es el subcuerpo primo de $k_0$, $K$ es una
$P$-álgebra formalmente lisa (19.6.1). Observemos, por otra parte, que se tiene
$\Upsilon^K_{A/k_0/P}=0$: en efecto, esto equivale a ver que el homomorfismo
\[
  u_{A/k_0/P}\otimes1_K:
  \Omega^1_{k_0/P}\otimes_{k_0}K
  \longrightarrow \Omega^1_{A/P}\otimes_A K
\]
es inyectivo (20.6.14.5), donde $u$ designa el homomorfismo estructural
$k_0\to A$. Pero como $A$ es una $k_0$-álgebra formalmente lisa por hipótesis,
resulta de (20.7.2) que $u_{A/k_0/P}$ es formalmente invertible por la
izquierda, por lo que (19.1.7) $u_{A/k_0/P}\otimes1_K$ es invertible por la
izquierda y, \emph{a fortiori}, inyectivo, lo que demuestra nuestra afirmación.
Aplicando la sucesión exacta (20.6.22.1), donde $\Lambda,A,B,C$ se sustituyen
por $P,k_0,A,K$, se obtiene la sucesión exacta
\[
  0\longrightarrow \Upsilon^K_{A/k_0}
  \longrightarrow \mathfrak m/\mathfrak m^2
  \longrightarrow \Omega^1_{A/k_0}\otimes_A K
  \longrightarrow \Omega^1_{K/k_0}\longrightarrow0,
\]
de modo que se tiene
\[
\begin{split}
  \operatorname{rg}_K(\Omega^1_{A/k_0}\otimes_AK)
  &={\operatorname{rg}}_K(\mathfrak m/\mathfrak m^2)
    +\bigl(\operatorname{rg}_K(\Omega^1_{K/k_0})
    -\operatorname{rg}_K(\Upsilon_{K/k_0})\bigr)\\
  &={\operatorname{rg}}_K(\mathfrak m/\mathfrak m^2)
    +\deg.\operatorname{tr}_{k_0}K
\end{split}
\]
en virtud de (21.7.1). Por otra parte, $\widehat\Omega^1_{A/k_0}$ es un
$A$-módulo formalmente proyectivo (20.4.9), y como su topología es la topología
$\mathfrak m$-preádica (20.4.5),
$\widehat\Omega^1_{A/k_0}/\mathfrak m^{j+1}\widehat\Omega^1_{A/k_0}$ es un
$(A/\mathfrak m^{j+1})$-módulo proyectivo para todo $j$ (19.2.4), luego libre
($\mathbf 0_{\mathrm{III}}$, 10.1.2) y de rango

\oldpage[0\textsubscript{IV-1}]{177}
$n=\operatorname{rg}_K(\Omega^1_{A/k_0}\otimes_AK)$. En virtud de (21.9.1),
$\widehat\Omega^1_{A/k_0}$ es un $A$-módulo de tipo finito y, además, este
$A$-módulo es plano en virtud de ($\mathbf 0_{\mathrm{III}}$, 10.2.1); es,
por tanto, libre ($\mathbf 0_{\mathrm{III}}$, 10.1.3) de rango $n$, de donde
resulta la primera afirmación.

Como por hipótesis $\Omega^1_{k_0/k}$ es de rango finito sobre $k_0$, el
producto tensorial completado
$\Omega^1_{k_0/k}\widehat\otimes_{k_0}A$ se identifica con
$\Omega^1_{k_0/k}\otimes_{k_0}A$; como $A$ es una $k_0$-álgebra formalmente
lisa, resulta de (20.7.18) que el homomorfismo
\[
  \widehat u_{A/k_0/k}:
  \Omega^1_{k_0/k}\otimes_{k_0}A
  \longrightarrow\widehat\Omega^1_{A/k}
\]
admite un inverso por la izquierda que es un $A$-homomorfismo continuo, lo que
implica en particular que su imagen es cerrada en
$\widehat\Omega^1_{A/k}$; la sucesión exacta
\[
  0\longrightarrow \Omega^1_{k_0/k}\otimes_{k_0}A
  \xrightarrow{\widehat u_{A/k_0/k}}\widehat\Omega^1_{A/k}
  \longrightarrow\widehat\Omega^1_{A/k_0}\longrightarrow0
\]
es, por tanto, exacta y escindida en virtud de (20.7.17) y de lo que precede;
esto termina de demostrar la proposición.
\end{proof}

\begin{corollary}[21.9.3]
\label{0.21.9.3-es}
Sean $k$ un cuerpo y $A$ el anillo de series formales
$k[[T_1,\ldots,T_r]]$, provisto de su estructura usual de $k$-álgebra.
Entonces $\widehat\Omega^1_{A/k}$ es un $A$-módulo libre de rango
$r=\dim(A)$.
\end{corollary}

\begin{proof}
Basta observar que $A$ es una $k$-álgebra formalmente lisa (19.3.4), y aplicar
(21.9.2) con $k_0=k=K$.
\end{proof}

\begin{lemma}[21.9.4]
\label{0.21.9.4-es}
Sean $k$ un cuerpo de característica $p>0$, $A$ una $k$-álgebra que es un
anillo local noetheriano completo y $B$ una sub-$k$-álgebra de $A$ isomorfa a
una $k$-álgebra topológica de la forma $k[[T_1,\ldots,T_r]]$ y tal que $A$ sea
una $B$-álgebra de tipo finito. Si $B_1$ es la subálgebra
$k[[T_1^p,\ldots,T_r^p]]$ de $B$, $\widehat\Omega^1_{A/k}$ se identifica
canónicamente con $\Omega^1_{A/B_1}$.
\end{lemma}

\begin{proof}
Toda derivación continua de $B_1$ en un $A$-módulo topológico $P$ que sea la
restricción de una $k$-derivación continua de $A$ en $P$ es \emph{nula}, pues
lo es en el subanillo $k[T_1^p,\ldots,T_r^p]$ de los polinomios, y este último
es denso en $B_1$. La sucesión exacta (20.3.7.1) muestra, por tanto, que el
homomorfismo canónico
\[
  \operatorname{D\acute{e}r.\,cont}_{B_1}(A,P)
  \longrightarrow \operatorname{D\acute{e}r.\,cont}_k(A,P)
\]
es biyectivo; la conclusión resulta de (20.7.14.4), teniendo en cuenta que
$\Omega^1_{A/B_1}$ es un $A$-módulo de tipo finito (20.4.7), luego separado y
completo, puesto que $A$ es un anillo local noetheriano completo.
\end{proof}

\begin{proposition}[21.9.5]
\label{0.21.9.5-es}
Sean $A$ un anillo local noetheriano completo íntegro y $A_0$ un subanillo de
$A$ isomorfo a un anillo de series formales $k[[T_1,\ldots,T_r]]$ sobre un
cuerpo $k$, tal que $A$ sea una $A_0$-álgebra finita y que el cuerpo de
fracciones $E$ de $A$ sea una extensión separable del cuerpo de fracciones
$L_0$ de $A_0$. Entonces se tiene
\begin{equation}
\tag{21.9.5.1}
\label{0.21.9.5.1-es}
  \operatorname{rg}_A\widehat\Omega^1_{A/k}
  =\operatorname{rg}_E(\widehat\Omega^1_{A/k}\otimes_AE)=\dim(A).
\end{equation}
\end{proposition}

\begin{proof}
Observemos que si $\mathfrak m_0$ es el ideal maximal de $A_0$, la topología
de $A$ es la topología $\mathfrak m_0$-ádica, puesto que $A$ es una
$A_0$-álgebra finita (Bourbaki, \emph{Alg. comm.}, cap.~IV, \S~2,
n\textsuperscript{o}~5, cor.~3 de la prop.~9), e induce en $A_0$ la topología
$\mathfrak m_0$-ádica (Bourbaki, \emph{Alg. comm.}, cap.~III, \S~3,
n\textsuperscript{o}~4, teo.~3). Se sabe (21.9.1) que
$\widehat\Omega^1_{A/k}$ es un $A$-módulo de tipo finito y,

\oldpage[0\textsubscript{IV-1}]{178}
en virtud de (21.9.3), $\widehat\Omega^1_{A_0/k}$ es un $A_0$-módulo libre de
rango $r=\dim(A_0)=\dim(A)$ (16.1.5); el producto tensorial
$\widehat\Omega^1_{A_0/k}\otimes_{A_0}A$ es, por tanto, un $A$-módulo libre de
rango $r$ idéntico a su completado separado; finalmente,
$\widehat\Omega^1_{A/A_0}$ es un $A$-módulo de tipo finito (20.4.7), luego
idéntico a su completado separado ($\mathbf 0_{\mathrm I}$, 7.3.6). Sentado
esto, en la sucesión de homomorfismos
\begin{equation}
\tag{21.9.5.2}
\label{0.21.9.5.2-es}
  \widehat\Omega^1_{A_0/k}\otimes_{A_0}A
  \xrightarrow{u}\widehat\Omega^1_{A/k}
  \xrightarrow{v}\widehat\Omega^1_{A/A_0}\longrightarrow0
\end{equation}
se sabe que $v$ es sobreyectivo y que $\operatorname{Im}(u)$ es densa en
$\operatorname{Ker}(v)$ (20.7.17); pero todo sub-$A_0$-módulo de
$\widehat\Omega^1_{A/k}$ es cerrado para la topología $\mathfrak m_0$-ádica
($\mathbf 0_{\mathrm I}$, 7.3.5), por lo que la sucesión (21.9.5.2) es exacta.
Observemos, por otra parte, que como $A$ es entero sobre $A_0$, $A_0$ es
íntegramente cerrado y $E$ es separable sobre $L_0$,
$\widehat\Omega^1_{A/A_0}$ es un $A$-módulo de torsión (20.4.13, (iv)); al
tensorizar la sucesión exacta (21.9.5.2) por $E$, se obtiene una sucesión
exacta
\[
  \widehat\Omega^1_{A_0/k}\otimes_{A_0}E
  \longrightarrow \widehat\Omega^1_{A/k}\otimes_AE\longrightarrow0
\]
de donde $\operatorname{rg}_E(\widehat\Omega^1_{A/k}\otimes_AE)\leqslant r$.
Queda, por tanto, mostrar que existen, en
$\widehat\Omega^1_{A/k}\otimes_AE$, $r$ elementos linealmente independientes
sobre $E$. Ahora bien, sean $D_i=\partial/\partial T_i$ las derivaciones
canónicas de $A_0$ en sí mismo $(1\leqslant i\leqslant r)$; se prolongan de
manera única a derivaciones (también denotadas $D_i$) de $E$ en sí mismo,
puesto que $E$ es una extensión separable finita de $L_0$; por restricción a
$A$, estas derivaciones dan derivaciones de $A$ en $E$, y es inmediato que
estas derivaciones toman sus valores en un mismo sub-$A$-módulo de tipo finito
de $E$; como son continuas en $A_0$ y la topología de $A$ es la topología
$\mathfrak m_0$-ádica, se han formado así $r$ derivaciones continuas de $A$
en un $A$-módulo topológico, que son evidentemente linealmente independientes,
puesto que $D_i(T_j)=\delta_{ij}$; esto termina de demostrar la proposición.
\end{proof}

\begin{proposition}[21.9.6]
\label{0.21.9.6-es}
Sean $k_0\subset k\subset K_0$ tres cuerpos de característica $p>0$, tales
que $[K_0:k]<+\infty$; se pone
\[
  L_0=K_0((T_1,\ldots,T_r)),\qquad
  M_0=K_0((T_1^p,\ldots,T_r^p)),
\]
\[
  L=k((T_1,\ldots,T_r)),\qquad
  M=k((T_1^p,\ldots,T_r^p)),\qquad
  N=k((T_1^{p^2},\ldots,T_r^{p^2})).
\]
Sea $E$ una extensión de tipo finito de $L_0$.
\begin{enumerate}
  \item[(i)] Si $\Upsilon_{K_0/k_0}$ es de rango finito, se tiene
  \begin{equation}
  \tag{21.9.6.1}
  \label{0.21.9.6.1-es}
    \operatorname{rg}_E\Omega^1_{E/M}
    -\operatorname{rg}_{K_0}\Omega^1_{K_0/k}
    =(r+\deg.\operatorname{tr}_{L_0}E)
      +d(E/K_0,k_0)+d(E/M_0,N/k_0)
  \end{equation}
  (notaciones de (21.6.1)), y en particular se tiene
  \begin{equation}
  \tag{21.9.6.2}
  \label{0.21.9.6.2-es}
    \operatorname{rg}_E\Omega^1_{E/M}
    -\operatorname{rg}_{K_0}\Omega^1_{K_0/k}
    \geqslant r+\deg.\operatorname{tr}_{L_0}E+d(E/K_0,k_0).
  \end{equation}
  Además, si los dos miembros de (21.9.6.2) son iguales, siguen siéndolo al
  sustituir $k$ por un cuerpo $k'$ tal que $k_0\subset k'\subset k$.

  \item[(ii)] Se supone además que $[k_0:k_0^p]<+\infty$. Sea
  $(k_\alpha)$ una familia filtrante decreciente de subcuerpos de $K_0$ que
  contienen a $k_0$, tales que $[K_0:k_\alpha]<+\infty$ para todo $\alpha$ y
  que $\bigcap_\alpha k_\alpha(K_0^p)=k_0(K_0^p)$; entonces existe un
  $\alpha$ tal que, para $k=k_\alpha$, los dos miembros de (21.9.6.2) sean
  iguales.
\end{enumerate}
\end{proposition}

\begin{proof}
\begin{enumerate}
  \oldpage[0\textsubscript{IV-1}]{179}
  \item[(i)] Se sabe (21.1.5, (ii)) que
  $\operatorname{rg}_E\Omega^1_{E/M}$ no cambia cuando se sustituye $M$ por
  $M(L_0^p)$; como
  $L_0^p=K_0^p((T_1^p,\ldots,T_r^p))$ y $[k(K_0^p):k]<+\infty$, se tiene
  \[
    M(L_0^p)=k(K_0^p)((T_1^p,\ldots,T_r^p))~;
  \]
  asimismo, $\Omega^1_{K_0/k}$ no cambia cuando en él se sustituye $k$ por
  $k(K_0^p)$, de modo que se puede \emph{suponer que $K_0^p\subset k$}, como
  haremos en todo lo que sigue. Introduciremos también los cuerpos
  \[
    k'_0=k_0(K_0^p),\qquad
    N'=k'_0((T_1^{p^2},\ldots,T_r^{p^2}))
  \]
  de modo que se tiene el diagrama de cuerpos
  \[
    \xymatrix{
      K_0\ar[r] & M_0\ar[r] & L_0\ar[r] & E\\
      k\ar[u]\ar[r] & N\ar[u]\ar[r] & M\ar[u]\ar[r] & L\ar[u]\\
      k_0\ar[u]\ar[r] & k'_0\ar[u]\ar[r] & N'\ar[u] &
    }
  \]

  Nos proponemos evaluar la diferencia $\delta$ entre el primer y el segundo
  miembro de (21.9.6.2).

  Resulta de (21.7.3) que se tiene
  \[
    \operatorname{rg}_E\Omega^1_{E/M}
    -\operatorname{rg}_{M_0}\Omega^1_{M_0/M}
    =\deg.\operatorname{tr}_{M_0}E+d(E/M_0,M)
  \]
  y como $L_0$ es una extensión finita de $M_0$, se tiene
  $\deg.\operatorname{tr}_{L_0}E=\deg.\operatorname{tr}_{M_0}E$. Por otra
  parte, una $p$-base de $K_0$ sobre $k=k(K_0^p)$ es también una $p$-base de
  $M_0$ sobre $M$, por lo que (21.4.5) se tiene
  \[
    \operatorname{rg}_{M_0}\Omega^1_{M_0/M}
    =\operatorname{rg}_{K_0}\Omega^1_{K_0/k}
  \]
  y, por consiguiente,
  \begin{equation}
  \tag{21.9.6.3}
  \label{0.21.9.6.3-es}
    \delta=d(E/M_0,M)-d(E/K_0,k_0)-r.
  \end{equation}

  Observemos ahora el lema clásico:

  \begin{lemma}[21.9.6.4]
  \label{0.21.9.6.4-es}
  Para todo cuerpo $k$, el cuerpo de series formales
  $K=k((T_1,\ldots,T_r))$ es una extensión separable de $k$.
  \end{lemma}

  \begin{proof}
  Recordemos brevemente la demostración de este lema, para ser completos.
  Basta (Bourbaki, \emph{Alg.}, cap.~VIII, \S~7,
  n\textsuperscript{o}~3, demostración del teo.~1) demostrar que, para toda
  extensión finita $k'$ de $k$, $K\otimes_k k'$ no tiene elementos nilpotentes;
  pero si se pone $A=k[[T_1,\ldots,T_r]]$ y $S=A-\{0\}$,
  $K\otimes_k k'$ es igual a $S^{-1}(A\otimes_k k')$, y
  $A\otimes_k k'$ se identifica canónicamente con el anillo íntegro
  $A'=k'[[T_1,\ldots,T_r]]$, y $A$ con un subanillo de $A'$, de donde resulta
  la conclusión.
  \end{proof}

  Utilizando este lema y (21.6.2, (ii)), se tiene $d(L_0/K_0,k_0)=0$, y la
  fórmula (21.6.4.4) da, por tanto,
  \[
    d(E/K_0,k_0)=d(E/L_0,k_0).
  \]

  \oldpage[0\textsubscript{IV-1}]{180}
  Las fórmulas (21.6.4.4) y (21.6.4.2) dan, por otra parte,
  \[
    d(E/M_0,M)=d(E/L_0,M)+d(L_0/M_0,M),
  \]
  \[
    d(L_0/M_0,M)=d(L_0/M_0,M/N)+d(L_0/M_0,N),
  \]
  de donde
  \[
    \delta=d(E/L_0,M)+d(L_0/M_0,N)-d(E/L_0,k_0)
    +d(L_0/M_0,M/N)-r.
  \]

  Vamos a demostrar que se tiene la relación
  \begin{equation}
  \tag{21.9.6.5}
  \label{0.21.9.6.5-es}
    d(L_0/M_0,M/N)=r.
  \end{equation}
  Para ello, observemos que, por definición, se tiene
  \[
    d(L_0/M_0,M/N)
    =\operatorname{rg}_{L_0}\Upsilon_{L_0/M/N}
    -\operatorname{rg}_{M_0}\Upsilon_{M_0/M/N}.
  \]
  Teniendo en cuenta las dos sucesiones exactas
  \[
  \begin{aligned}
    0&\longrightarrow\Upsilon_{L_0/M/N}
      \longrightarrow\Omega^1_{M/N}\otimes_ML_0
      \longrightarrow\Omega^1_{L_0/N}
      \longrightarrow\Omega^1_{L_0/M}\longrightarrow0,\\
    0&\longrightarrow\Upsilon_{M_0/M/N}
      \longrightarrow\Omega^1_{M/N}\otimes_MM_0
      \longrightarrow\Omega^1_{M_0/N}
      \longrightarrow\Omega^1_{M_0/M}\longrightarrow0,
  \end{aligned}
  \]
  se obtiene
  \[
  \begin{split}
    d(L_0/M_0,M/N)={}&\operatorname{rg}_{L_0}\Omega^1_{L_0/M}
      -\operatorname{rg}_{L_0}\Omega^1_{L_0/N}\\
      &-\operatorname{rg}_{M_0}\Omega^1_{M_0/M}
      +\operatorname{rg}_{M_0}\Omega^1_{M_0/N}.
  \end{split}
  \]
  Ahora bien, se tiene $N(L_0^p)\supset M$, luego
  $\Omega^1_{L_0/N}=\Omega^1_{L_0/M}$ (21.1.5, (ii)). Por otra parte, los
  $T_i^p$ $(1\leqslant i\leqslant r)$ forman una $p$-base de $M$ sobre $N$,
  y una $p$-base de $K_0$ sobre $k$ es también una $p$-base de $M_0$ sobre
  $M$, por lo que, como $M_0^p\subset N$,
  \[
    \operatorname{rg}_{M_0}\Omega^1_{M_0/N}
    =r+\operatorname{rg}_{K_0}\Omega^1_{K_0/k}
    =r+\operatorname{rg}_{M_0}\Omega^1_{M_0/M}
  \]
  en virtud de (21.4.5) y (21.1.10), lo que demuestra (21.9.6.5).

  Se tiene además, por (21.6.4.2),
  \[
    d(E/L_0,M)=d(E/L_0,k_0)+d(E/L_0,M/k_0)
  \]
  y
  \[
    d(L_0/M_0,N)=d(L_0/M_0,N/k)+d(L_0/M_0,k).
  \]
  Pero como $L_0$ es separable sobre $K_0$ (21.9.6.4), se tiene
  $d(L_0/K_0,k)=0$, y por consiguiente también $d(L_0/M_0,k)=0$ ((21.6.2) y
  (21.6.4)). Se obtiene, por tanto,
  \[
    \delta=d(E/L_0,M/k_0)+d(L_0/M_0,N/k).
  \]
  Apliquemos de nuevo dos veces (21.6.4.2) al primer término del segundo
  miembro de esta fórmula; se obtiene
  \[
    d(E/L_0,M/k_0)=d(E/L_0,M/N)+d(E/L_0,N/k)+d(E/L_0,k/k_0).
  \]
  Como $M\subset N(L_0^p)$, se tiene $d(E/L_0,M/N)=0$ (21.6.2), luego (por
  (21.6.4.4))
  \[
    \delta=d(E/L_0,k/k_0)+d(E/M_0,N/k).
  \]
  Pero $M_0$ y $L_0$ son extensiones separables de $K_0$ (21.9.6.4), por lo
  que
  \[
    d(L_0/K_0,k/k_0)=d(M_0/K_0,k/k_0)=0
  \]
  (21.6.2), y aplicando de nuevo dos veces (21.6.4.4),

  \oldpage[0\textsubscript{IV-1}]{181}
  se tiene
  \[
    d(E/L_0,k/k_0)=d(E/K_0,k/k_0)=d(E/M_0,k/k_0),
  \]
  de donde, por una última aplicación de (21.6.4.2),
  \[
    \delta=d(E/M_0,N/k_0),
  \]
  lo que demuestra (21.9.6.1). Además, cuando $k$ se sustituye por una
  subextensión $k'$ de $k_0$, $N$ se sustituye por una subextensión $N'$, por
  lo que $d(E/M_0,N'/k_0)\leqslant d(E/M_0,N/k_0)$ (21.6.4.2), lo que
  demuestra la última afirmación de (i).

  \item[(ii)] Para todo $\alpha$, pongamos
  $N_\alpha=k_\alpha((T_1^{p^2},\ldots,T_r^{p^2}))$. En virtud de (21.9.6.1)
  y de (21.8.3), se trata de demostrar (teniendo en cuenta que
  $K_0^p\subset k_\alpha$ y que $E$ es una extensión de tipo finito de $M_0$)
  que se tiene
  \[
    \bigcap_\alpha N_\alpha=k_0(M_0^p).
  \]
  Ahora bien, se tiene
  \begin{equation}
  \tag{21.9.6.6}
  \label{0.21.9.6.6-es}
    \bigcap_\alpha N_\alpha=N'
    =(k_0(K_0^p))((T_1^{p^2},\ldots,T_r^{p^2}))
  \end{equation}
  en virtud de (21.8.8.3) y de la relación
  $\bigcap_\alpha k_\alpha=k_0(K_0^p)$ (21.8.6). Por otra parte, se tiene
  $k_0(M_0^p)=k_0(K_0^p((T_1^{p^2},\ldots,T_r^{p^2})))$. Para terminar la
  demostración de (21.9.6, (ii)), queda, por tanto, demostrar el

  \begin{lemma}[21.9.6.7]
  \label{0.21.9.6.7-es}
  Sean $k_0$ un cuerpo de característica $p>0$ tal que
  $[k_0:k_0^p]<+\infty$ y $K_0$ una extensión de $k_0$. Entonces se tiene
  \begin{equation}
  \tag{21.9.6.8}
  \label{0.21.9.6.8-es}
    (k_0(K_0^p))((T_1,\ldots,T_r))
    =k_0\bigl(K_0^p((T_1,\ldots,T_r))\bigr).
  \end{equation}
  \end{lemma}

  \begin{proof}
  Como $k_0^p\subset K_0^p$, se tiene
  \[
    k_0^p\bigl(K_0^p((T_1,\ldots,T_r))\bigr)
    =K_0^p((T_1,\ldots,T_r)).
  \]
  Si $(c_i)_{1\leqslant i\leqslant n}$ es una base de $k_0$ sobre $k_0^p$,
  es inmediato que $(c_i)$ es también un sistema de generadores de cada uno
  de los dos miembros de (21.9.6.8) sobre
  $K_0^p((T_1,\ldots,T_r))$.
  \end{proof}
\end{enumerate}
\end{proof}

\begin{remark}[21.9.7]
\label{0.21.9.7-es}
La desigualdad (21.9.6.2) puede ser estricta si se omite la hipótesis de que
$[k_0:k_0^p]$ es finito. Tomemos, en efecto, como $k_0$ el cuerpo
$\mathbf F_p(X_n)_{n\geqslant0}$ de las fracciones racionales en una infinidad
de indeterminadas, de modo que los $X_n$ formen una $p$-base de $k_0$ sobre
$\mathbf F_p$, y por consiguiente $[k_0:k_0^p]=+\infty$. En el enunciado de
(21.9.6), tomemos $K_0=k_0$ (por tanto, necesariamente $k=k_0$),
$L_0=k_0((T))$, y como $E$ la extensión $L_0(y)$, donde $y$ es raíz del
polinomio
\[
  Y^p-\sum_{n=0}^{\infty}X_nT^{np}
\]
de $L_0[Y]$. Entonces la diferencia $\delta$ entre los dos miembros de
(21.9.6.2) es \emph{no nula}. En efecto, $E$ es \emph{separable} sobre $k_0$,
pues de lo contrario $y^p$ sería elemento de $k_0(L_0^p)$, y se ve
inmediatamente que esto no es posible (debido a que hay una infinidad de
$X_n$ algebraicamente independientes); se tiene, por tanto,
$d(E/K_0,k_0)=0$, la fórmula (21.9.6.3) se reduce a
$\delta=d(E/M_0,M_0)-1$, y está claro que
$d(E/M_0,M_0)=\operatorname{rg}_E\Upsilon_{E/M_0}$. Todo se reduce a comprobar
que este rango es $2$. Ahora bien, se tiene $E^p\subset M_0$, una $p$-base de
$E$ sobre $L_0$ está formada por el único elemento $y$, y $L_0$ tiene
evidentemente una $p$-base sobre $M_0$ formada por el único elemento $T$; por
tanto, como $E$ es una extensión algebraica de $M_0$, nuestra afirmación
resulta de (21.7.1) y (21.4.5).
\end{remark}

\begin{proposition}[21.9.8]
\label{0.21.9.8-es}
Sean $k_0$ un cuerpo (de característica cualquiera), $K_0$ una extensión
separable de $k_0$ y $A$ un anillo local noetheriano completo, que es una
$k_0$-álgebra y cuyo cuerpo residual

\oldpage[0\textsubscript{IV-1}]{182}
es una extensión finita de $K_0$. Para todo ideal primo $\mathfrak p$ de $A$,
sea $k(\mathfrak p)$ el cuerpo residual de $A$ en $\mathfrak p$. Entonces,
para todo cuerpo $k$ tal que $k_0\subset k\subset K_0$ y tal que
$[K_0:k]<+\infty$, se tiene
\begin{equation}
\tag{21.9.8.1}
\label{0.21.9.8.1-es}
\begin{split}
  &\operatorname{rg}_{k(\mathfrak p)}
    \bigl(\widehat\Omega^1_{(A/\mathfrak p)/k}
      \otimes_{A/\mathfrak p}k(\mathfrak p)\bigr)
    -\operatorname{rg}_{K_0}\Omega^1_{K_0/k}\\
  &\hspace{35mm}\geqslant
    \dim(A/\mathfrak p)
    +\operatorname{rg}_{k(\mathfrak p)}\Upsilon_{k(\mathfrak p)/k_0}.
\end{split}
\end{equation}
Supongamos además que $[k_0:k_0^p]<+\infty$ (donde $p$ es el exponente
característico de $k_0$). Entonces se puede encontrar un cuerpo $k$ tal que
$k_0(K_0^p)\subset k\subset K_0$ y $[K_0:k]<+\infty$, para el cual los dos
miembros de (21.9.8.1) son iguales.
\end{proposition}

\begin{proof}
Como $A/\mathfrak p$ tiene el mismo cuerpo residual $K$ que $A$ y es completo,
se puede reducir al caso en que $A$ es íntegro y $\mathfrak p=0$; se pondrá
entonces $E=k(\mathfrak p)$, cuerpo de fracciones de $A$. Como $K_0$ es
formalmente liso sobre $k_0$ (19.6.1), existe un $k_0$-monomorfismo
$K_0\to A$ que hace de $A$ una $K_0$-álgebra; se sabe además que existe una
sub-$K_0$-álgebra $A_0$ de $A$, $K_0$-isomorfa a un álgebra de series formales
$K_0[[T_1,\ldots,T_r]]$ y tal que $A$ sea una $A_0$-álgebra \emph{finita}
(19.8.9), lo que implica que $E$ es una extensión finita del cuerpo de
fracciones $L_0=K_0((T_1,\ldots,T_r))$ de $A_0$. Se tiene entonces
$\Upsilon_{K_0/k_0}=0$, puesto que $K_0$ es separable sobre $k_0$, y
$\deg.\operatorname{tr}_{L_0}E=0$; además, como $A$ es una $A_0$-álgebra
finita, se tiene $\dim(A)=\dim(A_0)$ (16.1.5).

Si $p=1$, se tiene
$\operatorname{rg}_E(\widehat\Omega^1_{A/k}\otimes_AE)=r$ por (21.9.5) y
(17.1.4, (iii)), y $\Omega^1_{K_0/k}=0$, puesto que $K_0$ es una extensión
finita separable de $k$ (21.7.4); por otra parte,
$\Upsilon_{E/k_0}=0$, ya que $E$ es una extensión separable de $k_0$; los dos
miembros de (21.9.8.1) son, por tanto, iguales en este caso.

Supongamos ahora que $p>1$; si $B=k[[T_1,\ldots,T_r]]$, $A$ es una
$B$-álgebra finita, por lo que, poniendo
$B_1=k[[T_1^p,\ldots,T_r^p]]$, $\widehat\Omega^1_{A/k}$ se identifica con
$\Omega^1_{A/B_1}$ (21.9.4), y, designando por $M$ el cuerpo de fracciones
$k((T_1^p,\ldots,T_r^p))$, resulta de (20.5.9) que
$\widehat\Omega^1_{A/k}\otimes_AE$ se identifica con $\Omega^1_{E/M}$; la
desigualdad (21.9.8.1) no es entonces sino (21.9.6.2), y la última afirmación
de la proposición resulta de (21.9.6, (ii)).\footnote{La fuente francesa
imprime «del corolario»; el resultado que se está demostrando es la
proposición (21.9.8).}
\end{proof}

\section{Criterios diferenciales de lisitud formal y de regularidad}
\label{section:0.22-es}

Los principales resultados de este parágrafo son los siguientes:
\begin{enumerate}
  \item[a)] Los criterios (22.2.2) y (22.5.8), que dan condiciones
  necesarias y suficientes para que un anillo local noetheriano completo $A$
  que contiene un cuerpo $k$ sea formalmente liso sobre $k$; el criterio
  (22.5.8) se enuncia sin hacer intervenir nociones diferenciales y completa
  (19.6.6); el enunciado (22.2.2), de naturaleza algo más técnica, permite dar
  una \emph{clasificación}, para $k$ fijado, de los anillos locales $A$
  considerados (22.2.5): vienen dados por su dimensión $n$ y por su cuerpo
  residual $K$, con la única condición
  $n\geqslant\operatorname{rg}_K\Upsilon_{K/k}$.

  \item[b)] El teorema (22.3.3) y el corolario (22.7.6) del criterio
  jacobiano de Nagata (22.7.3), que dan condiciones que permiten afirmar que
  una \emph{localización} de un anillo local noetheriano completo $A$ que
  contiene un cuerpo $k$ es geométricamente regular sobre $k$; desempeñarán
  un papel importante en el estudio de las propiedades «~finas~» de los
  anillos locales realizado en el cap.~IV, §~7. Por el momento no se conoce
  ninguna demostración de (22.7.6) independiente del criterio jacobiano de
  Nagata.

  \item[c)] El criterio jacobiano de Zariski (22.6.7) y sus variantes, que
  son consecuencias fáciles de (20.7.8).
\end{enumerate}

\oldpage[0\textsubscript{IV-1}]{183}

\subsection{Levantamiento de la lisitud formal}
\label{subsection:0.22.1-es}

\begin{theorem}[22.1.1]
\label{0.22.1.1-es}
Sean $s:\Lambda\to A$, $u:A\to B$ dos homomorfismos de anillos,
$\mathfrak J$ un ideal de $B$ y $C$ el anillo cociente $B/\mathfrak J$. Se
supone que $C$ es una $\Lambda$-álgebra formalmente lisa para la topología
discreta.
\begin{enumerate}
  \item[(i)] Supongamos que se cumplen las condiciones siguientes:
  \begin{enumerate}
    \item[$\alpha$)] $\mathfrak J/\mathfrak J^2$ es un $C$-módulo proyectivo.
    \item[$\beta$)] El homomorfismo canónico
    $\mathbf S_C^*(\mathfrak J/\mathfrak J^2)\to
    \operatorname{gr}_{\mathfrak J}^*(B)$ (19.5.2) es biyectivo.
    \item[$\gamma$)] El homomorfismo característico
    $\chi:\Upsilon_{C/B/\Lambda}\to\mathfrak J/\mathfrak J^2$ (20.6.20)
    es inyectivo.
    \item[$\delta$)] El $C$-módulo
    $\Omega^1_{B/A}\otimes_BC$ es proyectivo.
  \end{enumerate}
  Entonces $B$, provisto de la topología $\mathfrak J$-preádica, es una
  $A$-álgebra formalmente lisa.

  \item[(ii)] Recíprocamente, supongamos que $B$ sea una $A$-álgebra
  formalmente lisa para la topología $\mathfrak J$-preádica y que $A$ sea una
  $\Lambda$-álgebra formalmente lisa para la topología discreta. Entonces se
  cumplen las condiciones $\alpha$), $\beta$), $\gamma$), $\delta$) de (i).
\end{enumerate}
\end{theorem}

\begin{proof}
En virtud de la sucesión exacta (20.6.22.1) (que es aplicable, puesto que
$B/\mathfrak J^2$ es una extensión $\Lambda$-trivial de $C$ por
$\mathfrak J/\mathfrak J^2$, al suponerse que $C$ es una $\Lambda$-álgebra
formalmente lisa (19.4.4)), la condición $\gamma$) equivale a
$\Upsilon^C_{B/A/\Lambda}=0$, o también a la siguiente:
\begin{enumerate}
  \item[$\gamma'$)] El homomorfismo
  \[
    u_{B/A/\Lambda}\otimes 1_C:
    \Omega^1_{A/\Lambda}\otimes_A C
    \longrightarrow \Omega^1_{B/\Lambda}\otimes_B C
  \]
  es inyectivo.
\end{enumerate}

(i) Las condiciones $\alpha$) y $\beta$) implican que $B$ es una
$\Lambda$-álgebra formalmente lisa para la topología
$\mathfrak J$-preádica (19.5.4). Por tanto, equivale decir que $B$ es (para la
topología $\mathfrak J$-preádica) una $A$-álgebra formalmente lisa o una
$A$-álgebra formalmente lisa \emph{relativamente a $\Lambda$}. Para ver que
$B$ es una $A$-álgebra formalmente lisa basta, en virtud de (20.7.2), probar
que el homomorfismo continuo de $B$-módulos topológicos
\[
  u_{B/A/\Lambda}:\Omega^1_{A/\Lambda}\otimes_A B
  \longrightarrow\Omega^1_{B/\Lambda}
\]
es \emph{formalmente invertible por la izquierda}. Ahora bien, puesto que $B$
es una $\Lambda$-álgebra formalmente lisa, el $B$-módulo topológico
$\Omega^1_{B/\Lambda}$ es formalmente proyectivo (20.4.9) y su topología se
deduce de la de $B$ (20.4.5). En virtud de (19.1.9), para que
$u_{B/A/\Lambda}$ sea formalmente invertible por la izquierda basta que
$u_{B/A/\Lambda}\otimes 1_C$ sea invertible por la izquierda. Pero, en virtud
de $\gamma'$), esta última aplicación es inyectiva, por lo que se tiene la
sucesión exacta (20.6.14.7)
\[
  0\longrightarrow\Omega^1_{A/\Lambda}\otimes_A C
  \longrightarrow\Omega^1_{B/\Lambda}\otimes_B C
  \longrightarrow\Omega^1_{B/A}\otimes_B C\longrightarrow0.
\]
Finalmente, en virtud de $\delta$), el $C$-módulo
$\Omega^1_{B/A}\otimes_BC$ es proyectivo, luego la sucesión exacta anterior
se escinde, lo que termina de probar (i).

\oldpage[0\textsubscript{IV-1}]{184}

(ii) Si $B$ es una $A$-álgebra formalmente lisa para la topología
$\mathfrak J$-preádica y $A$ una $\Lambda$-álgebra formalmente lisa para la
topología discreta, $B$ es una $\Lambda$-álgebra formalmente lisa para la topología
$\mathfrak J$-preádica (19.3.5, (ii)); por tanto, las condiciones $\alpha$) y
$\beta$) resultan de (19.5.6). Además, (20.7.2) muestra que
$u_{B/A/\Lambda}$ es formalmente invertible por la izquierda, luego
$u_{B/A/\Lambda}\otimes 1_C$ es invertible por la izquierda (19.1.7) y, \emph{a
fortiori}, inyectivo. Finalmente, $\Omega^1_{B/A}$ es un $B$-módulo
formalmente proyectivo (20.4.9) y, por consiguiente,
$\Omega^1_{B/A}\otimes_BC$ es un $C$-módulo proyectivo (19.2.4).
\end{proof}

\begin{corollary}[22.1.2]
\label{0.22.1.2-es}
Sean $s:\Lambda\to A$, $u:A\to B$ dos homomorfismos de anillos,
$\mathfrak m$ un ideal maximal de $B$ y $K=B/\mathfrak m$ el cuerpo cociente.
Se supone que las $\Lambda$-álgebras $A$ y $K$ son formalmente lisas para la
topología discreta. Entonces las condiciones siguientes son equivalentes:
\begin{enumerate}
  \item[a)] $B$ es una $A$-álgebra formalmente lisa para la topología
  $\mathfrak m$-preádica.
  \item[b)] El homomorfismo canónico
  \begin{equation}
  \tag{22.1.2.1}
  \label{0.22.1.2.1-es}
    \mathbf S_K^*(\mathfrak m/\mathfrak m^2)
    \longrightarrow\operatorname{gr}_{\mathfrak m}^*(B)
  \end{equation}
  es biyectivo, y el homomorfismo característico
  \begin{equation}
  \tag{22.1.2.2}
  \label{0.22.1.2.2-es}
    \chi:\Upsilon_{K/B/\Lambda}\longrightarrow\mathfrak m/\mathfrak m^2
  \end{equation}
  es inyectivo.
\end{enumerate}
\end{corollary}

\begin{proof}
Esto resulta inmediatamente de (22.1.1), pues las condiciones $\alpha$) y
$\delta$) se cumplen aquí trivialmente, ya que $K$ es un cuerpo.
\end{proof}

\begin{remark}[22.1.3]
\label{0.22.1.3-es}
Supongamos que se cumplen las hipótesis de (22.1.2) y, además, que
$\mathfrak m/\mathfrak m^2$ es un $K$-espacio vectorial de rango \emph{finito}.
Entonces, para que $B$ sea una $A$-álgebra formalmente lisa para la topología
$\mathfrak m$-preádica, basta que se satisfagan las condiciones siguientes:

\begin{env}[22.1.3.1]
\label{0.22.1.3.1-es}
Dados un álgebra de polinomios $E=K[T_1,\ldots,T_n]$, el ideal maximal
$\mathfrak n$ de $E$ engendrado por los $T_i$ $(1\leqslant i\leqslant n)$ y
un ideal $\mathfrak b\supset\mathfrak n^k$ de $E$, todo
$\Lambda$-homomorfismo $v:B\to E/\mathfrak b$ que, al pasar a los cocientes,
da la identidad $K\to K$, se factoriza como
\[
  B\xrightarrow{w}E/\mathfrak n^k\longrightarrow E/\mathfrak b,
\]
donde $w$ es un $\Lambda$-homomorfismo.
\end{env}

\begin{env}[22.1.3.2]
\label{0.22.1.3.2-es}
Si $F$ es la extensión trivial $D_K(K)$ («~álgebra de los números duales
sobre $K$~»), $\rho:A\to F$ un homomorfismo de anillos tal que el compuesto
$A\xrightarrow{\rho}F\to K$ sea el homomorfismo canónico, entonces todo
$\Lambda$-homomorfismo $v:B\to K$ que, al pasar a los cocientes, da la
identidad $K\to K$, se factoriza como
\[
  B\xrightarrow{w}F\longrightarrow K,
\]
donde $w$ es un $\Lambda$-homomorfismo (para la estructura de $A$-álgebra
sobre $F$ definida por $\rho$).
\end{env}

En efecto, se ha visto (19.5.5) que la condición (22.1.3.1) implica que el
homomorfismo canónico (22.1.2.1) es biyectivo. En virtud de (22.1.2), basta
entonces ver que el homomorfismo canónico
\[
  u_{B/A/\Lambda}\otimes1_K:
  \Omega^1_{A/\Lambda}\otimes_AK
  \longrightarrow\Omega^1_{B/\Lambda}\otimes_BK
\]
es inyectivo, o también (puesto que se trata de espacios vectoriales sobre
$K$) que, para todo $K$-espacio vectorial $L$, el homomorfismo canónico
$\operatorname{D\acute{e}r}_\Lambda(B,L)\to\operatorname{D\acute{e}r}_\Lambda(A,L)$ es
\emph{sobreyectivo} (en virtud de (20.4.8)); es claro que, para ello, es
necesario y suficiente que el homomorfismo canónico
$\operatorname{D\acute{e}r}_\Lambda(B,K)\to\operatorname{D\acute{e}r}_\Lambda(A,K)$ sea
sobreyectivo, de donde resulta nuestra afirmación en virtud de (20.1.5).

Se observará que las dos condiciones (22.1.3.1) y (22.1.3.2) son consecuencia
de la siguiente:

\begin{env}[22.1.3.3]
\label{0.22.1.3.3-es}
Para toda $\Lambda$-álgebra local artiniana $E$ de cuerpo residual igual a
$K$, y todo ideal nilpotente $\mathfrak r$ de $E$, todo
$\Lambda$-homomorfismo $B\to E/\mathfrak r$ que, al pasar a los cocientes,
da la identidad $K\to K$, se factoriza como
$B\xrightarrow{u}E\to E/\mathfrak r$, donde $u$ es un
$\Lambda$-homomorfismo.
\end{env}
\end{remark}

\oldpage[0\textsubscript{IV-1}]{185}

Este resultado se generaliza como sigue:

\begin{proposition}[22.1.4]
\label{0.22.1.4-es}
Sea $\varphi:A\to B$ un homomorfismo local de anillos locales noetherianos,
$k$ (resp. $K$) el cuerpo residual de $A$ (resp. de $B$). Para que $B$ sea una
$A$-álgebra formalmente lisa (para las topologías preádicas), es necesario y
suficiente que, para todo anillo local artiniano $C$, de cuerpo residual igual
a $K$, todo homomorfismo local $A\to C$ (que hace de $C$ una $A$-álgebra
topológica), y todo ideal nilpotente $\mathfrak J$ de $C$, todo
$A$-homomorfismo local $B\to C/\mathfrak J$ tal que el homomorfismo
$K\to K$ obtenido al pasar a los cocientes sea la identidad, se factorice como
\[
  B\xrightarrow{f}C\longrightarrow C/\mathfrak J,
\]
donde $f$ es un $A$-homomorfismo local.
\end{proposition}

\begin{proof}
Como la condición es necesaria por definición (19.3.1), mostremos que es
suficiente. Puesto que los $\widehat A$-homomorfismos locales de $\widehat B$
en una $\widehat A$-álgebra local discreta proceden, por extensión, de los
$A$-homomorfismos locales de $B$ en esta álgebra, podemos limitarnos al caso
en que $A$ y $B$ sean \emph{completos}, en virtud de (19.3.6). Cuando $A$ es
un \emph{cuerpo} $k$, la proposición resulta de (22.1.3.3), tomando para
$\Lambda$ el subcuerpo primo de $k$, y de (19.6.1).

En el caso general, pongamos $B_0=B\otimes_Ak=B/\mathfrak mB$
($\mathfrak m$ ideal maximal de $A$), que es completo. Si $C_0$ es una
$k$-álgebra local artiniana de cuerpo residual igual a $K$, $\mathfrak J_0$
un ideal nilpotente de $C_0$, todo $k$-homomorfismo local
$g_0:B_0\to C_0/\mathfrak J_0$ que da la identidad sobre $K$ al pasar a los
cocientes da, por composición, un $A$-homomorfismo local
\[
  g:B\longrightarrow B_0\xrightarrow{g_0}C_0/\mathfrak J_0
\]
que, por hipótesis, se factoriza por tanto como
$B\xrightarrow{f}C_0\to C_0/\mathfrak J_0$, donde $f$ es un
$A$-homomorfismo local; pero, como $C_0$ es una $k$-álgebra, $f$ se factoriza
como $B\to B_0\xrightarrow{f_0}C_0$, donde $f_0$ es un
$k$-homomorfismo local. Se ve así que las hipótesis del enunciado siguen
cumpliéndose cuando se sustituyen $B_0$ y $k$ por $B$ y $A$, respectivamente;
por tanto, $B_0$ es una $k$-álgebra \emph{formalmente lisa} por lo que se acaba
de ver.

Aplicando (19.7.2) y (19.7.1), se ve que existen un anillo local noetheriano
completo $B'$, un homomorfismo local $A\to B'$ que hace de $B'$ un $A$-módulo
\emph{plano} y una $A$-álgebra \emph{formalmente lisa}, y un
$k$-isomorfismo
\[
  u'_0:B'_0=B'\otimes_Ak\xrightarrow{\sim}B_0.
\]
Sea $v'_0$ el isomorfismo recíproco de $u'_0$; nos proponemos mostrar que
existe un $A$-homomorfismo local $v:B\to B'$ tal que $v'_0$ proceda de $v$ al
pasar a los cocientes. Para ello, designemos por $\mathfrak n$ y
$\mathfrak n'$ los ideales maximales respectivos de $B$ y $B'$ y, para todo
$j$, sea
\[
  w_j:B/(\mathfrak n^{j+1}+\mathfrak mB)
  \longrightarrow B'/(\mathfrak n'^{j+1}+\mathfrak mB')
\]
el homomorfismo deducido de $v'_0$ al pasar a los cocientes; bastará construir,
para todo $j$, un $A$-homomorfismo local
$v_j:B/\mathfrak n^{j+1}\to B'/\mathfrak n'^{j+1}$ tal que los diagramas
\[
\xymatrix@C=4em@R=3em{
  B/\mathfrak n^{j+1} \ar[r]^{v_j} \ar[d] &
  B'/\mathfrak n'^{j+1} \ar[d] \\
  B/\mathfrak n^j \ar[r]_{v_{j-1}} & B'/\mathfrak n'^j
}
\qquad
\xymatrix@C=4em@R=3em{
  B/\mathfrak n^{j+1} \ar[r]^{v_j} \ar[d] &
  B'/\mathfrak n'^{j+1} \ar[d] \\
  B/(\mathfrak n^{j+1}+\mathfrak mB) \ar[r]_{w_j} &
  B'/(\mathfrak n'^{j+1}+\mathfrak mB')
}
\]
sean conmutativos; como $B$ y $B'$ son completos, el homomorfismo
$v=\varprojlim v_j$ responderá a la cuestión.

\oldpage[0\textsubscript{IV-1}]{186}

Ahora bien, la construcción por recurrencia de $v_j$ resulta de la hipótesis
sobre $B$ y del mismo razonamiento que en (19.3.10.1), utilizando el lema
(19.3.10.2) y el hecho de que
\[
  \mathfrak n'^j+(\mathfrak n'^{j+1}+\mathfrak mB')
  =\mathfrak n'^j+\mathfrak mB'.
\]
Resulta entonces de (19.7.1.4) que $v$ es un $A$-\emph{isomorfismo}, pues
$B'$ es un $A$-módulo plano y $B$ es completo para la topología
$\mathfrak m$-preádica, pues los ideales $\mathfrak m^iB$ son cerrados en $B$
para la topología $\mathfrak n$-ádica ($\mathbf 0_{\mathrm I}$, 7.3.5).
Q.E.D.
\end{proof}

\subsection{Caracterización diferencial de las álgebras locales formalmente
lisas sobre un cuerpo}
\label{subsection:0.22.2-es}

\begin{env}[22.2.1]
\label{0.22.2.1-es}
Sean $k$ un cuerpo, $P$ su subcuerpo primo, $A$ una $k$-álgebra que es un
anillo \emph{local}, $\mathfrak m$ su ideal maximal y $K=A/\mathfrak m$ su
cuerpo residual. Como $K$ es separable sobre $P$, $K$ es una $P$-álgebra
formalmente lisa para la topología discreta (19.6.1), y por consiguiente la
$P$-extensión $A/\mathfrak m^2$ de $K$ por $\mathfrak m/\mathfrak m^2$ es
$P$-\emph{trivial}, y el homomorfismo característico
\begin{equation}
\tag{22.2.1.1}
\label{0.22.2.1.1-es}
  \chi_{A/k}:\Upsilon_{K/k}\longrightarrow\mathfrak m/\mathfrak m^2
\end{equation}
queda así definido (20.6.23).
\end{env}

\begin{theorem}[22.2.2]
\label{0.22.2.2-es}
Bajo las condiciones de (22.2.1), para que $A$ sea una $k$-álgebra
formalmente lisa para la topología $\mathfrak m$-preádica, es necesario y
suficiente que se satisfagan las condiciones siguientes:
\begin{enumerate}
  \item[(i)] El homomorfismo canónico
  $\mathbf S_K^*(\mathfrak m/\mathfrak m^2)\to
  \operatorname{gr}_{\mathfrak m}^*(A)$ es biyectivo.
  \item[(ii)] El homomorfismo característico
  $\chi_{A/k}:\Upsilon_{K/k}\to\mathfrak m/\mathfrak m^2$ es inyectivo.
\end{enumerate}
\end{theorem}

\begin{proof}
Es un caso particular de (22.1.2), donde se sustituye $\Lambda$ por $P$, $A$
por $k$ y $B$ por $A$; en efecto, como $K$ y $k$ son separables sobre $P$,
son $P$-álgebras formalmente lisas (19.6.1).
\end{proof}

\begin{remark}[22.2.3]
\label{0.22.2.3-es}
La condición (ii) de (22.2.2) puede sustituirse por cualquiera de las
siguientes:
\begin{enumerate}
  \item[a)] El homomorfismo canónico
  $\Omega^1_k\otimes_kK\to\Omega^1_A\otimes_AK$ es inyectivo.
  \item[b)] Para todo $n$, el homomorfismo canónico
  \[
    \Omega^1_k\otimes_k(A/\mathfrak m^{n+1})
    \longrightarrow
    \Omega^1_A\otimes_A(A/\mathfrak m^{n+1})
  \]
  es invertible por la izquierda.
  \item[c)] El homomorfismo canónico
  $\Omega^1_k\widehat\otimes_kA\to\widehat\Omega^1_A$ es invertible por la
  izquierda.
  \item[d)] $A$ es una $k$-álgebra formalmente lisa (para la topología
  $\mathfrak m$-preádica) \emph{relativamente al cuerpo primo}.
\end{enumerate}
En efecto, ya se ha observado en (22.1.1) que la inyectividad de $\chi_{A/k}$
equivale a $a)$, y que la conjunción de $a)$ y de la condición (i) de (22.2.2)
equivale a la de $d)$ y (i). Finalmente, también se ha visto en (22.1.1) que
$b)$ y $d)$ son equivalentes, y la equivalencia de $b)$ y $c)$ resulta de
(19.1.8).
\end{remark}

\begin{env}[22.2.4]
\label{0.22.2.4-es}
La principal aplicación de (22.2.2) se refiere a las $k$-álgebras locales
noetherianas y \emph{completas} (caso en que $\mathfrak m/\mathfrak m^2$ es un
$K$-espacio vectorial de rango \emph{finito}). Precisamente, sean $k$ un
cuerpo, $K$ una extensión de $k$ y $V$ un $K$-espacio vectorial de

\oldpage[0\textsubscript{IV-1}]{187}

rango finito; consideramos las ternas $(A,\alpha,\beta)$, donde $A$ es una
$k$-álgebra local noetheriana completa \emph{formalmente lisa} de ideal
maximal $\mathfrak m$, $\alpha:A/\mathfrak m\xrightarrow{\sim}K$ un
isomorfismo de $k$-álgebras y
$\beta:\mathfrak m/\mathfrak m^2\xrightarrow{\sim}V$ un diisomorfismo de
espacios vectoriales (para el isomorfismo $\alpha$). Dadas dos de tales ternas
$(A,\alpha,\beta)$, $(A',\alpha',\beta')$, una \emph{equivalencia} de la
primera sobre la segunda es un $k$-isomorfismo $u:A\xrightarrow{\sim}A'$ tal
que (si $\mathfrak m'$ es el ideal maximal de $A'$) los diagramas
\[
\xymatrix@C=4em@R=3em{
  A/\mathfrak m \ar[r]^{\operatorname{gr}^0(u)} \ar[d]_{\alpha} &
  A'/\mathfrak m' \ar[d]^{\alpha'} \\
  K \ar[r]_{1_K} & K
}
\qquad
\xymatrix@C=4em@R=3em{
  \mathfrak m/\mathfrak m^2 \ar[r]^{\operatorname{gr}^1(u)}
    \ar[d]_{\beta} &
  \mathfrak m'/\mathfrak m'^2 \ar[d]^{\beta'} \\
  V \ar[r]_{1_V} & V
}
\]
sean conmutativos. Es inmediato que las \emph{clases de equivalencia} de las
ternas $(A,\alpha,\beta)$ forman un \emph{conjunto} que denotaremos por
$\operatorname{FL}(K/k,V)$. Observemos ahora que a toda terna
$(A,\alpha,\beta)$ se asocia el $K$-homomorfismo compuesto de espacios
vectoriales
\[
  \Upsilon_{K/k}\xrightarrow{\chi_{A/k}}
  \mathfrak m/\mathfrak m^2\xrightarrow{\beta}V
\]
(donde $\chi_{A/k}$ se considera como un dihomomorfismo relativo a
$\alpha^{-1}$); la definición anterior demuestra (por (20.6.22.2)) que este
$K$-homomorfismo \emph{solo depende de la clase de equivalencia} de
$(A,\alpha,\beta)$ en $\operatorname{FL}(K/k,V)$; en otras palabras, se ha
definido una aplicación canónica
\begin{equation}
\tag{22.2.4.1}
\label{0.22.2.4.1-es}
  \operatorname{FL}(K/k,V)\longrightarrow
  \operatorname{Hom}_K(\Upsilon_{K/k},V).
\end{equation}
\end{env}

\begin{proposition}[22.2.5]
\label{0.22.2.5-es}
Sean $k$ un cuerpo, $K$ una extensión de $k$ y $V$ un $K$-espacio vectorial
de rango finito. La aplicación canónica (22.2.4.1) es una biyección de
$\operatorname{FL}(K/k,V)$ sobre el conjunto de los $K$-homomorfismos
inyectivos de $\Upsilon_{K/k}$ en $V$.
\end{proposition}

\begin{proof}
(i) Para mostrar que (22.2.4.1) es inyectiva, hay que probar lo siguiente:
sean $A$, $A'$ dos $k$-álgebras locales noetherianas completas formalmente
lisas, $\mathfrak m$, $\mathfrak m'$ sus ideales maximales respectivos,
$\alpha:A/\mathfrak m\xrightarrow{\sim}K$,
$\alpha':A'/\mathfrak m'\xrightarrow{\sim}K$ dos $k$-isomorfismos;
supongamos dados dos $k$-isomorfismos
\[
  v^0:A/\mathfrak m\xrightarrow{\sim}A'/\mathfrak m',\qquad
  v^1:\mathfrak m/\mathfrak m^2\xrightarrow{\sim}
      \mathfrak m'/\mathfrak m'^2
\]
tales que los diagramas
\begin{equation}
\tag{22.2.5.1}
\label{0.22.2.5.1-es}
\xymatrix@C=4em@R=3em{
  A/\mathfrak m \ar[r]^{v^0} \ar[dr]_{\alpha} &
  A'/\mathfrak m' \ar[dl]^{\alpha'} \\
  & K
}
\qquad
\xymatrix@C=4em@R=3em{
  \mathfrak m/\mathfrak m^2 \ar[r]^{v^1} &
  \mathfrak m'/\mathfrak m'^2 \\
  \Upsilon_{K/k} \ar[u]^{\chi_A} \ar[ur]_{\chi_{A'}} &
}
\end{equation}
sean conmutativos. Se trata de probar que existe un $k$-isomorfismo
$u:A\xrightarrow{\sim}A'$ tal que $\operatorname{gr}^0(u)=v^0$ y
$\operatorname{gr}^1(u)=v^1$. Observemos primero que, puesto que $K$ es un
cuerpo, el homomorfismo canónico (20.6.7.1) es \emph{biyectivo} (20.6.11),
luego ya existe un $k$-isomorfismo
$v:A/\mathfrak m^2\xrightarrow{\sim}A'/\mathfrak m'^2$ tal que
$\operatorname{gr}^0(v)=v^0$ y $\operatorname{gr}^1(v)=v^1$. Observemos ahora
que, puesto que $A'$ es metrizable y completo y $A$ una $k$-álgebra
formalmente lisa, el homomorfismo compuesto
$A\to A/\mathfrak m^2\xrightarrow{v}A'/\mathfrak m'^2$ se factoriza como
$A\xrightarrow{u}A'\to A'/\mathfrak m'^2$, donde $u$ es

\oldpage[0\textsubscript{IV-1}]{188}

un $k$-homomorfismo (19.3.11), y se tiene ya, por tanto,
$\operatorname{gr}^0(u)=v^0$ y $\operatorname{gr}^1(u)=v^1$. Ahora bien, se
tiene el diagrama conmutativo
\[
\xymatrix@C=5em@R=3em{
  \mathbf S_K^*(\mathfrak m/\mathfrak m^2) \ar[r] \ar[d]_{w} &
  \operatorname{gr}_{\mathfrak m}^*(A) \ar[d]^{\operatorname{gr}^*(u)} \\
  \mathbf S_K^*(\mathfrak m'/\mathfrak m'^2) \ar[r] &
  \operatorname{gr}_{\mathfrak m'}^*(A')
}
\]
donde las flechas horizontales son los homomorfismos canónicos y $w$ procede
de $v^0$ y $v^1$; como $v^0$ y $v^1$ son biyectivos, también lo es $w$, y,
puesto que $A$ y $A'$ son $k$-álgebras formalmente lisas, se deduce de
(22.2.2, (i)) que $\operatorname{gr}(u)$ es biyectivo; pero, como $A$ y $A'$
son separados y completos, esto implica que el propio $u$ es
\emph{biyectivo} (Bourbaki, \emph{Alg. comm.}, cap.~III, §~2,
n\textsuperscript{o}~8, cor.~3 del th.~1).

(ii) Queda por mostrar que la imagen de (22.2.4.1) es el conjunto de los
homomorfismos \emph{inyectivos} de $\Upsilon_{K/k}$ en $V$; ya se sabe por
(22.2.2, (ii)) que está contenida en este conjunto; bastará entonces probar el
lema siguiente.
\end{proof}

\begin{lemma}[22.2.5.2]
\label{0.22.2.5.2-es}
Para todo $K$-homomorfismo $h:\Upsilon_{K/k}\to V$, existe una $k$-álgebra
$A$ que es un anillo local noetheriano regular y completo, de ideal maximal
$\mathfrak m$, de cuerpo residual $K$, y un $K$-isomorfismo
$\beta:\mathfrak m/\mathfrak m^2\xrightarrow{\sim}V$ tales que
$h=\beta\circ\chi_{A/k}$.
\end{lemma}

\begin{proof}
En efecto, si se prueba este lema, el hecho de que $A$ sea regular implica que
se satisface la condición (i) de (22.2.2) (17.1.1); si además $h$ es
inyectivo, (22.2.2) mostrará que $A$ es una $k$-álgebra formalmente lisa cuya
clase tendrá por imagen $h$ (en cambio, si $h$ no es inyectivo, la $k$-álgebra
$A$ cuya existencia prueba (22.2.5.2) no es formalmente lisa).

Para probar (22.2.5.2), sean $B$ la $K$-álgebra $\mathbf S_K^*(V)$,
$\mathfrak n$ el ideal de aumentación de $B$ y $A=\widehat B$ la $K$-álgebra
completada de $B$ para la topología $\mathfrak n$-preádica (de modo que $A$ es
isomorfa al álgebra de series formales $K[[T_1,\ldots,T_n]]$ si
$n=\operatorname{rg}_KV$). En virtud de (20.6.11), existe sobre
$A/\mathfrak m^2$ (donde $\mathfrak m=\mathfrak nA$) una estructura de
$k$-extensión de $K$ por $V$ (en general \emph{distinta} de la estructura
$k$-trivial definida por la inyección canónica $k\to K\to A$) tal que $h$ sea
el homomorfismo característico de esta extensión. Si
$f:k\to A/\mathfrak m^2$ es el homomorfismo estructural, el hecho de que $k$
sea separable sobre el cuerpo primo permite (19.6.1) factorizar $f$ como
$k\xrightarrow{g}A\to A/\mathfrak m^2$, y la $k$-álgebra $A$ así definida
responde a la cuestión (20.6.22.2).
\end{proof}

\begin{theorem}[22.2.6]
\label{0.22.2.6-es}
Sean $k$ un cuerpo y $K$ una extensión de $k$. Para que exista una $k$-álgebra
local noetheriana $A$, de ideal maximal $\mathfrak m$, tal que
$A/\mathfrak m=K$ y que $A$ sea formalmente lisa (para su topología
$\mathfrak m$-preádica), es necesario y suficiente que $\Upsilon_{K/k}$ sea
un $K$-espacio vectorial de rango finito. La $k$-álgebra $A$ queda entonces
determinada (salvo un $k$-isomorfismo que, al pasar a los cocientes, da la
identidad $K\to K$) por su dimensión, que está sometida a la única condición
de ser $\geqslant\operatorname{rg}_K\Upsilon_{K/k}$.
\end{theorem}

\begin{proof}
Como $\mathfrak m/\mathfrak m^2$ es de rango finito sobre $K$, la necesidad
de la condición resulta de (22.2.2, (ii)); recíprocamente, (22.2.5) muestra que
la condición es suficiente y que se puede tomar como
$\mathfrak m/\mathfrak m^2$ un $K$-espacio vectorial de rango cualquiera
$\geqslant\operatorname{rg}_K(\Upsilon_{K/k})$.
\end{proof}

\oldpage[0\textsubscript{IV-1}]{189}

En particular, el homomorfismo identidad de $\Upsilon_{K/k}$ asocia
\emph{canónicamente} a $K$ una $k$-álgebra noetheriana completa, formalmente
lisa, para la cual $\Upsilon_{K/k}$ es isomorfo a
$\mathfrak m/\mathfrak m^2$.

\begin{remarks}[22.2.7]
\label{0.22.2.7-es}
\begin{enumerate}
  \item[(i)] Sea $\Lambda$ un anillo local noetheriano, de cuerpo residual
  $k$; resulta de (19.7.1) y (19.7.2) que la determinación de las
  $\Lambda$-álgebras $A$ que son anillos locales noetherianos completos y que
  son formalmente lisas es \emph{equivalente} al mismo problema en el que
  $\Lambda$ se sustituye por $k$, y queda por tanto, en principio, enteramente
  resuelta por (22.2.5), que reduce la cuestión a la estructura de los
  \emph{cuerpos residuales} de las $k$-álgebras buscadas, como extensiones de
  $k$.

  \item[(ii)] El hecho de que $\Upsilon_{K/k}$ sea de rango finito sobre $K$
  \emph{no implica} que $K$ sea «~de multiplicidad radicial finita~» sobre
  $k$ (cf. (19.6.6)). Por ejemplo, sean $k_0$ un cuerpo perfecto de
  característica $p>0$, $k=k_0(T)$ el cuerpo de las fracciones racionales en
  una indeterminada sobre $k_0$ y $K=k^{p^{-\infty}}$ la menor extensión
  perfecta de $k$. Se tiene entonces $\Omega^1_K=\Omega^1_{K/k}=0$ (21.4.4),
  luego, por definición (20.6.1.1),
  $\Upsilon_{K/k}=\Omega^1_k\otimes_kK$; y como $k^p=k_0(T^p)$, $T$ es una
  $p$-base de $k$ sobre $k^p$, luego (21.4.5) $\Omega^1_k$ es de rango $1$
  sobre $k$ y, por consiguiente, $\Upsilon_{K/k}$ de rango $1$ sobre $K$.
  Pero, para todo $n$, el cuerpo residual del anillo local
  $K\otimes_kk^{p^{-n}}$ es isomorfo a $K$ y, por tanto, no es separable sobre
  $k^{p^{-n}}$.
\end{enumerate}
\end{remarks}

El teorema (22.2.2) se generaliza como sigue:

\begin{proposition}[22.2.8]
\label{0.22.2.8-es}
Bajo las condiciones de (22.2.1), supongamos que la característica $p$ de $k$
sea $>0$; sea $(k_\alpha)$ una familia filtrante decreciente de subcuerpos de
$k$ tal que $\bigcap_\alpha k_\alpha(k^p)=k^p$. Entonces la condición (ii) de
(22.2.2) puede sustituirse por cualquiera de las siguientes:
\begin{enumerate}
  \item[a)] Existe $\alpha_0$ tal que el homomorfismo canónico
  \[
    \Omega^1_{k/k_\alpha}\otimes_kK
    \longrightarrow\Omega^1_{A/k_\alpha}\otimes_AK
  \]
  sea inyectivo para todo $\alpha\geqslant\alpha_0$.
  \item[b)] Existe $\alpha_0$ tal que $A$ sea una $k$-álgebra formalmente
  lisa (para la topología $\mathfrak m$-preádica) relativamente a $k_\alpha$,
  para todo $\alpha\geqslant\alpha_0$.
  \item[c)] Para todo $\alpha$ y todo entero $n\geqslant0$, el homomorfismo
  canónico
  \[
    \Omega^1_{k/k_\alpha}\otimes_k(A/\mathfrak m^{n+1})
    \longrightarrow
    \Omega^1_{A/k_\alpha}\otimes_A(A/\mathfrak m^{n+1})
  \]
  es invertible por la izquierda.
\end{enumerate}
\end{proposition}

\begin{proof}
En efecto, si $A$ es una $k$-álgebra formalmente lisa, lo es \emph{a
fortiori} relativamente a todo subcuerpo de $k$, lo que implica $c)$ en virtud
de (20.7.2) y (19.1.7); es claro que $c)$ implica $b)$, en virtud de (20.7.2),
y que $b)$ implica $a)$. Finalmente, teniendo en cuenta (22.2.3), queda por
probar que $a)$ implica que el homomorfismo canónico
\[
  u:\Omega^1_k\otimes_kK\longrightarrow\Omega^1_A\otimes_AK
\]
es inyectivo. Ahora bien, se tiene para todo $\alpha$ un diagrama conmutativo
\[
\xymatrix@C=5em@R=3em{
  \Omega^1_k\otimes_kK \ar[r]^{u} \ar[d]_{v_\alpha} &
  \Omega^1_A\otimes_AK \ar[d] \\
  \Omega^1_{k/k_\alpha}\otimes_kK \ar[r]_{u_\alpha} &
  \Omega^1_{A/k_\alpha}\otimes_AK
}
\]

\oldpage[0\textsubscript{IV-1}]{190}

y la hipótesis de que $u_\alpha$ es inyectivo para
$\alpha\geqslant\alpha_0$ implica que $\operatorname{Ker}(u)$ está contenido
en la \emph{intersección} de los $\operatorname{Ker}(v_\alpha)$. Pero se ha
visto (21.8.3) que esta intersección es nula en virtud de la hipótesis
$\bigcap_\alpha k_\alpha(k^p)=k^p$, lo que termina la demostración.
\end{proof}

\begin{proposition}[22.2.9]
\label{0.22.2.9-es}
Bajo las hipótesis de (22.2.1), las condiciones siguientes son equivalentes:
\begin{enumerate}
  \item[a)] $A/\mathfrak m^2$ es una $k$-extensión $k$-trivial de $K$ por
  $\mathfrak m/\mathfrak m^2$.
  \item[b)] $\chi_{A/k}=0$.
  \item[c)] El homomorfismo canónico (20.5.11.2)
  \[
    \delta_{K/A/k}:\mathfrak m/\mathfrak m^2
    \longrightarrow\Omega^1_{A/k}\otimes_AK
  \]
  es inyectivo (en otras palabras, la sucesión
  \begin{equation}
  \tag{22.2.9.1}
  \label{0.22.2.9.1-es}
    0\longrightarrow\mathfrak m/\mathfrak m^2
    \longrightarrow\Omega^1_{A/k}\otimes_AK
    \longrightarrow\Omega^1_{K/k}\longrightarrow0
  \end{equation}
  es exacta).
\end{enumerate}
\end{proposition}

\begin{proof}
La equivalencia de $b)$ y $c)$ resulta de la sucesión exacta (20.6.22.1); la
equivalencia de $c)$ y $a)$ resulta de (20.5.12), pues la sucesión (22.2.9.1) se
escinde si es exacta, ya que se trata de espacios vectoriales.

Se observará que, si $K$ es \emph{separable sobre $k$}, se satisfacen las
condiciones equivalentes de (22.2.9), en virtud de la definición de
$\chi_{A/k}$ (22.2.1) y de (20.6.3).
\end{proof}

\begin{corollary}[22.2.10]
\label{0.22.2.10-es}
Bajo las hipótesis de (22.2.1), consideremos un subcuerpo $k_0$ de $k$. Las
condiciones siguientes son equivalentes:
\begin{enumerate}
  \item[a)] $A/\mathfrak m^2$ es una $k_0$-extensión $k_0$-trivial de $K$
  por $\mathfrak m/\mathfrak m^2$.
  \item[b)] El homomorfismo compuesto
  \[
    \Upsilon_{K/k_0}\longrightarrow\Upsilon_{K/k}
    \xrightarrow{\chi_{A/k}}\mathfrak m/\mathfrak m^2
  \]
  es nulo.
  \item[c)] El homomorfismo característico
  $\chi_{A/k}:\Upsilon_{K/k}\to\mathfrak m/\mathfrak m^2$ se factoriza como
  \[
    \Upsilon_{K/k}\longrightarrow\Upsilon_{K/k/k_0}
    \xrightarrow{\chi'}\mathfrak m/\mathfrak m^2.
  \]
\end{enumerate}
Además, cuando se cumplen estas condiciones, el homomorfismo $\chi'$ queda
determinado de manera única y coincide con el homomorfismo característico de
la $k_0$-extensión $k_0$-trivial $A/\mathfrak m^2$ de $K$ por
$\mathfrak m/\mathfrak m^2$.
\end{corollary}

\begin{proof}
Como el compuesto de $b)$ no es otro que $\chi_{A/k_0}$, la equivalencia de
$a)$ y $b)$ resulta de (22.2.9); por otra parte, $b)$ y $c)$ son equivalentes
en virtud de
la sucesión exacta (20.6.18.1)
\[
  \Upsilon_{K/k_0}\longrightarrow\Upsilon_{K/k}
  \longrightarrow\Upsilon_{K/k/k_0}\longrightarrow0,
\]
y, como el homomorfismo canónico
$\Upsilon_{K/k}\to\Upsilon_{K/k/k_0}$ es sobreyectivo, esto implica la
unicidad de $\chi'$; el hecho de que $\chi'$ sea el homomorfismo característico
de la $k_0$-extensión $A/\mathfrak m^2$ resulta de (20.6.22.2).
\end{proof}

\begin{corollary}[22.2.11]
\label{0.22.2.11-es}
Bajo las hipótesis de (22.2.1), sea $p$ el exponente característico de $k$ y
sea $(k_\alpha)$ una familia filtrante decreciente de subcuerpos de $k$ tal
que $\bigcap_\alpha k_\alpha(k^p)=k^p$. Si $\Upsilon_{K/k}$ es de rango finito
sobre $K$, existe un índice $\alpha$ tal que $A/\mathfrak m^2$ sea una
$k_\alpha$-extensión $k_\alpha$-trivial de $K$ por
$\mathfrak m/\mathfrak m^2$.
\end{corollary}

\begin{proof}
En efecto, se sabe (21.8.3) que existe entonces un índice $\alpha$ tal que el
homomorfismo canónico $\Upsilon_{K/k}\to\Upsilon_{K/k/k_\alpha}$ sea
biyectivo, luego se satisface la condición $b)$ de (22.2.10) (donde $k_\alpha$
sustituye a $k_0$).
\end{proof}

\oldpage[0\textsubscript{IV-1}]{191}

\subsection{Aplicación a las relaciones entre ciertos anillos locales y sus completados}
\label{subsection:0.22.3-es}

\begin{lemma}[22.3.1]
\label{0.22.3.1-es}
Sean $A_0$ un anillo semilocal noetheriano y $A$ una $A_0$-álgebra finita
(que es, por tanto, un anillo semilocal noetheriano, cf. Bourbaki,
\emph{Alg. comm.}, cap.~IV, §~2, n\textsuperscript{o}~5, cor.~3 de la
prop.~9). Si $\widehat A_0$ y $\widehat A$ son los completados de $A_0$ y $A$
para sus topologías preádicas respectivas, entonces, cuando se provee a
$A_0$, $A$ y $\widehat A$ de las topologías discretas, $\widehat A$ es una
$A$-álgebra formalmente lisa relativamente a $A_0$ (19.9.1).
\end{lemma}

\begin{proof}
En efecto, se sabe que
$\widehat A=A\otimes_{A_0}\widehat A_0$ (Bourbaki,
\emph{Alg. comm.}, cap.~IV, §~2, n\textsuperscript{o}~5, cor.~3 de la
prop.~9 y cap.~III, §~3, n\textsuperscript{o}~4, th.~1), luego la proposición
resulta de (19.9.3).
\end{proof}

\begin{proposition}[22.3.2]
\label{0.22.3.2-es}
Sean $A$ un anillo local noetheriano regular, $K$ su cuerpo de fracciones y
$p$ la característica de $K$. Se supone que se cumple una de las hipótesis
siguientes:
\begin{enumerate}
  \item[(i)] $p=0$.
  \item[(ii)] $p>0$ y existe una familia filtrante decreciente $(A_\alpha)$
  de subanillos noetherianos de $A$, tales que $A$ sea una
  $A_\alpha$-álgebra finita para todo $\alpha$, que
  $A^{p^{h(\alpha)}}\subset A_\alpha$ para un entero $h(\alpha)$ conveniente,
  y que, si $K_\alpha$ denota el cuerpo de fracciones de $A_\alpha$, se tenga
  \[
    \bigcap_\alpha K_\alpha(K^p)=K^p.
  \]
\end{enumerate}
Sean entonces $B$ una $A$-álgebra finita íntegra que contiene a $A$,
$\mathfrak n$ un ideal primo de $B$, $C$ el anillo local $B_\mathfrak n$ y
$\mathfrak q$ un ideal primo del completado $\widehat C$ tal que
$\mathfrak q\cap C=0$, de modo que el anillo local
$\widehat C_\mathfrak q$ es un álgebra sobre el cuerpo de fracciones $L$ de
$B$. Entonces $\widehat C_\mathfrak q$ es una $L$-álgebra formalmente lisa
para su topología $\mathfrak q$-ádica y, por consiguiente, un anillo
geométricamente regular sobre $L$ (19.6.5).
\end{proposition}

\begin{proof}
Distinguamos dos casos.

I) Supongamos que $\mathfrak n$ contiene el ideal maximal $\mathfrak m$ de
$A$, y por consiguiente $\mathfrak n\cap A=\mathfrak m$. Entonces
$\mathfrak n$ es maximal (Bourbaki, \emph{Alg. comm.}, cap.~V, §~2,
n\textsuperscript{o}~1, prop.~1) y, si $\mathfrak r$ es el radical del anillo
semilocal $B$, $\widehat C$ es el completado separado de $B$ para la topología
$\mathfrak n$-preádica, y uno de los componentes del anillo semilocal
$\widehat B$, completado de $B$ para la topología $\mathfrak r$-preádica
(Bourbaki, \emph{Alg. comm.}, cap.~III, §~3, n\textsuperscript{o}~4,
prop.~8); se puede escribir, por tanto,
$\widehat C_\mathfrak q=\widehat B_{\mathfrak q'}$, donde $\mathfrak q'$ es
el ideal primo de $\widehat B$ imagen recíproca de $\mathfrak q$. Observemos
ahora que $B$ es un subanillo de $\widehat B$ y la hipótesis
$\mathfrak q\cap C=0$ implica $\mathfrak q'\cap B=0$, luego
$\widehat B_{\mathfrak q'}$ es también una localización del anillo
$\widehat B\otimes_B L$ en uno de sus ideales primos $\mathfrak q''$, y
$\mathfrak q'$ es su imagen recíproca en $\widehat B$. Además, se tiene
$\widehat B=\widehat A\otimes_A B$ (Bourbaki, \emph{Alg. comm.}, cap.~IV,
§~2, n\textsuperscript{o}~5, cor. de la prop.~9 y cap.~III, §~3,
n\textsuperscript{o}~4, th.~1), luego
$\widehat B\otimes_B L=\widehat A\otimes_A L$; finalmente, si $\mathfrak p$
es el ideal primo de $\widehat A$, imagen recíproca de $\mathfrak q''$, el
hecho de que $A\to B$ sea inyectivo y la conmutatividad del diagrama
\[
  \xymatrix{
    B \ar[r] & \widehat B\\
    A \ar[u] \ar[r] & \widehat A \ar[u]
  }
\]
implican que $\mathfrak p\cap A=0$, luego
$(\widehat A\otimes_A L)_{\mathfrak q''}$ es también una localización del
anillo
\[
  \widehat A_\mathfrak p\otimes_A L
  =(\widehat A_\mathfrak p\otimes_A K)\otimes_K L
\]
en uno de sus ideales primos.

\oldpage[0\textsubscript{IV-1}]{192}

Para mostrar que este anillo local es una $L$-álgebra formalmente lisa, basta
probar que $\widehat A_\mathfrak p\otimes_A K$ es una $K$-álgebra formalmente
lisa (19.3.5, (iii) y (iv)); por otra parte,
$\widehat A_\mathfrak p\otimes_A K=\widehat A_\mathfrak p$, puesto que
$\mathfrak p\cap A=0$. Vamos a aplicar los criterios (22.2.2) y (22.2.8).
En primer lugar, como $A$ es regular, también lo es $\widehat A$ (17.1.5),
luego $\widehat A_\mathfrak p$ (17.3.2), y se cumple la condición (i) de
(22.2.2) (17.1.1). Basta por consiguiente (cuando $p>0$) verificar la
condición b) de (22.2.8). Ahora bien, (22.3.1) muestra que, para todo
$\alpha$, $\widehat A$ es una $A$-álgebra formalmente lisa relativamente a
$A_\alpha$ (para las topologías discretas); por consiguiente,
\[
  \widehat A\otimes_{A_\alpha}K_\alpha
  =\widehat A\otimes_A(A\otimes_{A_\alpha}K_\alpha)
\]
es una $(A\otimes_{A_\alpha}K_\alpha)$-álgebra formalmente lisa relativamente
a $K_\alpha$, para las topologías discretas (19.9.2, (iii)). Pero, como $A$
es una $A_\alpha$-álgebra finita, se tiene
$A\otimes_{A_\alpha}K_\alpha=K$; como $\widehat A_\mathfrak p$ es un anillo
de fracciones de $\widehat A\otimes_A K$, se concluye (19.9.2, (iv)) que
$\widehat A_\mathfrak p$ es una $K$-álgebra formalmente lisa relativamente a
$K_\alpha$, para la topología discreta, y \emph{a fortiori} para su topología
$\mathfrak p$-ádica ((19.3.8) y (19.9.5)). Pero esta es precisamente la
condición b) de (22.2.8), luego la proposición queda demostrada en este caso
para $p>0$. Cuando $p=0$, el cuerpo residual de
$\widehat A_\mathfrak p$ es separable sobre $K$, y como
$\widehat A_\mathfrak p$ es un anillo regular, es una $K$-álgebra formalmente
lisa en virtud de (19.6.4).

II) Caso general. Sea $\mathfrak s$ el ideal primo $\mathfrak n\cap A$ de
$A$, y pongamos $S=A-\mathfrak s$, de modo que $S^{-1}A=A_\mathfrak s$ y se
tiene $C=(S^{-1}B)_{S^{-1}\mathfrak n}$, donde ahora
$S^{-1}\mathfrak n$ contiene el ideal maximal
$\mathfrak sA_\mathfrak s$ de $A_\mathfrak s$. Como $A_\mathfrak s$ es
regular (17.3.2) y $S^{-1}B$ una $A_\mathfrak s$-álgebra finita íntegra que
contiene a $A_\mathfrak s$, todo se reduce a verificar la condición (ii) para
$A_\mathfrak s$ cuando $p>0$: ahora bien, pongamos
$\mathfrak s_\alpha=\mathfrak s\cap A_\alpha$ y consideremos el anillo local
$(A_\alpha)_{\mathfrak s_\alpha}$. Para todo $t\notin\mathfrak s$, por
hipótesis, poniendo $q=p^{h(\alpha)}$, se tiene $t^q\in A_\alpha$ y
$t^q\notin\mathfrak s$, luego $t^q\notin\mathfrak s_\alpha$; si se pone
$S_\alpha=A_\alpha-\mathfrak s_\alpha$, se ve por tanto que
$A_\mathfrak s=S_\alpha^{-1}A$, y la hipótesis implica que
$A_\mathfrak s$ es una $(A_\alpha)_{\mathfrak s_\alpha}$-álgebra finita;
además, $K_\alpha$ es también el cuerpo de fracciones de
$(A_\alpha)_{\mathfrak s_\alpha}$, lo que termina la demostración.
\end{proof}

\begin{theorem}[22.3.3]
\label{0.22.3.3-es}
Sean $A$ un anillo local noetheriano completo, $\mathfrak p$ un ideal primo
de $A$, $B$ el anillo localizado $A_\mathfrak p$, $\mathfrak q$ un ideal
primo del completado $\widehat B$ y $\mathfrak r=\mathfrak q\cap B$.
Entonces el anillo localizado $\widehat B_\mathfrak q$ de $\widehat B$ es una
$B_\mathfrak r$-álgebra formalmente lisa para las topologías preádicas.
\end{theorem}

\begin{proof}
El ideal primo $\mathfrak r$ de $B$ es de la forma
$\mathfrak n_\mathfrak p$, donde $\mathfrak n\subset\mathfrak p$ es un ideal
primo de $A$, y se tiene $B_\mathfrak r=A_\mathfrak n$; sea $k$ el cuerpo
residual de $B_\mathfrak r$, que es también el cuerpo de fracciones del anillo
local íntegro y completo $A'=A/\mathfrak n$. Como $\widehat B$ es un
$B$-módulo plano $(\mathbf 0_\mathrm I,~7.3.5)$,
$\widehat B_\mathfrak q$ es un $B_\mathfrak r$-módulo plano
$(\mathbf 0_\mathrm I,~6.3.2)$, y para probar que
$\widehat B_\mathfrak q$ es una $B_\mathfrak r$-álgebra formalmente lisa
basta probar que
\[
  \widehat B_\mathfrak q\otimes_{B_\mathfrak r}k
  =\widehat B_\mathfrak q/\mathfrak r\widehat B_\mathfrak q
\]
es una $k$-álgebra formalmente lisa (19.7.1). Ahora bien,
$\widehat B_\mathfrak q/\mathfrak r\widehat B_\mathfrak q
=(\widehat B/\mathfrak r\widehat B)_{\mathfrak q'}$, donde $\mathfrak q'$ es
la imagen canónica de $\mathfrak q$ en
$\widehat B/\mathfrak r\widehat B$; por definición se tiene
$\mathfrak q'\cap(B/\mathfrak r)=0$ y
$B/\mathfrak r=A_{\mathfrak p'}$, donde
$\mathfrak p'=\mathfrak p/\mathfrak n$; además,
$\widehat B/\mathfrak r\widehat B$ es el completado de $B/\mathfrak r$.

Quedamos así reducidos a demostrar el corolario siguiente:
\end{proof}

\begin{corollary}[22.3.4]
\label{0.22.3.4-es}
Sean $A$ un anillo local noetheriano completo e íntegro, $K$ su cuerpo de
fracciones, $\mathfrak p$ un ideal primo de $A$ y $B$ el anillo localizado
$A_\mathfrak p$. Entonces, para todo ideal primo $\mathfrak q$

\oldpage[0\textsubscript{IV-1}]{193}

del completado $\widehat B$, tal que $\mathfrak q\cap B=0$, el anillo local
$\widehat B_\mathfrak q$ es una $K$-álgebra formalmente lisa para su topología
$\mathfrak q$-ádica y, por tanto, un anillo geométricamente regular sobre
$K$.
\end{corollary}

\begin{proof}
En efecto, resulta de (19.8.9) que existe un subanillo $A_0$ de $A$ que es un
anillo de series formales sobre un cuerpo de característica $p>0$ o sobre un
anillo de Cohen (se recuerda que los anillos de Cohen tienen un cuerpo de
fracciones de característica $0$), y tal que $A$ es una $A_0$-álgebra finita.
Ahora bien, se sabe que el anillo $A_0$ es regular (17.3.8) y cumple una de las
hipótesis (i), (ii) de (22.3.2), en virtud de (21.8.8); basta entonces aplicar
(22.3.2), sustituyendo $B$ por $A$ y $A$ por $A_0$.
\end{proof}
\subsection{Resultados preliminares sobre las extensiones finitas de anillos locales cuyo ideal maximal es de cuadrado nulo}
\label{subsection:0.22.4-es}
\label{0.22.4-es}

\begin{proposition}[22.4.1]
\label{0.22.4.1-es}
Sean $A$ un anillo local cuyo ideal maximal $\mathfrak m$ es de cuadrado
nulo, $K=A/\mathfrak m$ su cuerpo residual; denotemos por $V$ el ideal
$\mathfrak m$ considerado como $K$-espacio vectorial. Sean $T_i$
$(1\leqslant i\leqslant r)$ indeterminadas y, para cada $i$, sea
$F_i$ un polinomio mónico de $A[T_i]$ de grado $d_i$; denotemos por $A'$ el
cociente del anillo de polinomios $A[T_1,\ldots,T_r]$ por el ideal
$\mathfrak b$ engendrado por los $r$ polinomios $F_i$. Supongamos que $A'$ es
un anillo local, y sean $\mathfrak m'$ su ideal maximal y
$K'=A'/\mathfrak m'$ su cuerpo residual. Entonces:
\begin{enumerate}
  \item[(i)] Si se pone $d=\prod_{i=1}^r d_i$, $A'$ es un $A$-módulo
  libre de rango $d$, y se tiene $[K':K]\leqslant d$.
  \item[(ii)] Para que $[K':K]=d$, es necesario y suficiente que
  $A'\otimes_AK=A'/\mathfrak mA'$ sea un cuerpo (necesariamente isomorfo a
  $K'$), es decir, que $\mathfrak m'=\mathfrak mA'$. En este caso se tiene
  $\mathfrak m'^2=0$, y si $V'$ denota el ideal $\mathfrak m'$ considerado
  como $K'$-espacio vectorial, el homomorfismo canónico
  $V\otimes_KK'\to V'$ es biyectivo (y, por consiguiente,
  $\operatorname{rg}_{K'}V'=\operatorname{rg}_K V$).
\end{enumerate}
\end{proposition}

\begin{proof}
Sin suponer que $A'$ sea local, es evidente, por división euclídea e
inducción sobre $r$, que si $t_i$ es la clase de $T_i$ mod.~$\mathfrak b$,
los monomios
$M_\nu=t_1^{\nu_1}t_2^{\nu_2}\cdots t_r^{\nu_r}$, donde
$0\leqslant\nu_i\leqslant d_i-1$, forman una base del $A$-módulo $A'$. Si
$A'$ se supone local, $K'$ es un cociente de
$A'/\mathfrak mA'=A'\otimes_AK$, por tanto
$[K':K]\leqslant[A'\otimes_AK:K]=d$, y no puede haber igualdad salvo si
$\mathfrak m'=\mathfrak mA'$, lo que demuestra (i) y la primera afirmación
de (ii). Por otra parte, está claro que cuando
$\mathfrak m'=\mathfrak mA'$, se tiene $\mathfrak m'^2=0$, y
$V\otimes_KK'=\mathfrak m\otimes_AA'=\mathfrak mA'=V'$ puesto que $A'$ es
un $A$-módulo libre.
\end{proof}

\begin{remark}[22.4.2]
\label{0.22.4.2-es}
Cuando $A$ es artiniano (es decir, $\operatorname{rg}_K(V)$ es finito),
la hipótesis de que $A'$ sea local implica siempre que
$\operatorname{rg}_{K'}(\mathfrak m'/\mathfrak m'^2)
\geqslant\operatorname{rg}_K(V)$, como veremos más adelante (22.5.2).
Obsérvese que estos dos rangos pueden ser iguales sin que se tenga
$\mathfrak m'=\mathfrak mA'$.
\end{remark}

\begin{proposition}[22.4.3]
\label{0.22.4.3-es}
Siendo las hipótesis sobre $A$ y las notaciones las de (22.4.1), supongamos
que los $F_i$ son de la forma
\begin{equation}
\label{0.22.4.3.1-es}
  F_i(T_i)=T_i^{d_i}+\sum_{1\leqslant k\leqslant d_i}
  c_{ik}T_i^{d_i-k},
  \tag{22.4.3.1}
\end{equation}
donde $d_i\geqslant2$ para todo $i$, y $c_{ik}\in\mathfrak m$
(«polinomios de Eisenstein»). Entonces:
\begin{enumerate}
  \item[(i)] $A'$ es un anillo local, y su cuerpo residual
  $K'=A'/\mathfrak m'$ es isomorfo a $K$.

  \oldpage[0\textsubscript{IV-1}]{194}

  \item[(ii)] El núcleo del homomorfismo canónico
  $V\to V'=\mathfrak m'/\mathfrak m'^2$ es el ideal $\mathfrak n$ de $A$
  engendrado por los $\xi_i=c_{i,d_i}$ («términos constantes» de los $F_i$), y
  el conúcleo de este homomorfismo tiene por base sobre $K'$ las imágenes de los $T_i$;
  en particular, se tiene
  \begin{equation}
  \label{0.22.4.3.2-es}
    \operatorname{rg}_{K'}(V')=\operatorname{rg}_K(V)+r-r',
    \tag{22.4.3.2}
  \end{equation}
  donde $r'=\operatorname{rg}_K(\mathfrak n)$. Cuando
  $\operatorname{rg}_K(V)$ es finito (es decir, cuando $A$ es
  artiniano), para que
  $\operatorname{rg}_{K'}(V')=\operatorname{rg}_K(V)$, es necesario y suficiente
  que los $\xi_i$ sean linealmente independientes sobre $K$.
\end{enumerate}
\end{proposition}

\begin{proof}
(i) Está claro que $A'\otimes_AK=A'/\mathfrak mA'$ es isomorfo a
$K[T_1,\ldots,T_r]/\mathfrak b'$, donde $\mathfrak b'$ es el ideal engendrado
por los $r$ polinomios $T_i^{d_i}$; es, pues, un anillo local de cuerpo
residual $K$, cuyo ideal maximal está engendrado por las clases de los $T_i$
mod.~$\mathfrak b'$, de donde (i).

(ii) Sea $t_i$ $(1\leqslant i\leqslant r)$ la clase de $T_i$
mod.~$\mathfrak b$; se deduce de (i) que el ideal maximal $\mathfrak m'$ de
$A'$ viene dado por
\[
  \mathfrak m'=\mathfrak mA'+\sum_{i=1}^r t_iA',
\]
de donde, teniendo en cuenta que $\mathfrak m^2=0$,
\[
  \mathfrak m'^2=\sum_{i,j}t_it_jA'+\sum_i t_i\mathfrak mA'.
\]
Se concluye que $A'/\mathfrak m'^2$ es isomorfo al anillo
$A[T_1,\ldots,T_r]/\mathfrak c$, donde $\mathfrak c$ es el ideal engendrado por
los elementos
\[
  F_i\ (1\leqslant i\leqslant r),\qquad
  T_iT_j\ (1\leqslant i\leqslant r,\ 1\leqslant j\leqslant r),\qquad
  xT_i\ (1\leqslant i\leqslant r,\ x\in\mathfrak m).
\]
Como $d_i\geqslant2$ y los $T_i^2$ pertenecen a
$\mathfrak c$, las hipótesis hechas implican que $\mathfrak c$ está también
engendrado por los elementos
\[
  \xi_i\ (1\leqslant i\leqslant r),\qquad
  T_iT_j\ (1\leqslant i\leqslant r,\ 1\leqslant j\leqslant r),\qquad
  T_ix\ (1\leqslant i\leqslant r,\ x\in\mathfrak m).
\]
Sean
\[
  \mathfrak c_1=\sum_{i,j}T_iT_jA[T_1,\ldots,T_r],\qquad
  \mathfrak c_2=\mathfrak c_1+\sum_i\mathfrak mT_iA[T_1,\ldots,T_r],
\]
de modo que
$\mathfrak c=\mathfrak c_2+\sum_i\xi_iA[T_1,\ldots,T_r]$. Está claro que
$A[T_1,\ldots,T_r]/\mathfrak c_1$ es el $A$-módulo suma directa de $A$ y
de los $At_i^{(1)}$, donde $t_i^{(1)}$ es la clase de $T_i$
mod.~$\mathfrak c_1$. Se deduce que
$A[T_1,\ldots,T_r]/\mathfrak c_2$ es suma directa de $A$ y de los
$Kt_i^{(2)}$, donde $t_i^{(2)}$ es la clase de $T_i$
mod.~$\mathfrak c_2$. Finalmente, $A[T_1,\ldots,T_r]/\mathfrak c$ es suma
directa de $A/\mathfrak n$ y de los $Kt_i'$, donde $t_i'$ es la clase de $T_i$
mod.~$\mathfrak c$, y $\mathfrak n$ es el ideal de $A$ engendrado por los
$\xi_i$; en la identificación de $A'/\mathfrak m'^2$ con
$A[T_1,\ldots,T_r]/\mathfrak c$, la imagen de $V$ se identifica con
$\mathfrak m/\mathfrak n$, y lo anterior demuestra, por tanto, (ii).
\end{proof}

\begin{proposition}[22.4.4]
\label{0.22.4.4-es}
Siendo las hipótesis sobre $A$ y las notaciones las de (22.4.1), supongamos
que $K$ es de característica $p>0$ (por tanto $p\cdot1\in\mathfrak m$ en
$A$) y que los $F_i$ son de la forma
\begin{equation}
\label{0.22.4.4.1-es}
  F_i(T_i)=T_i^p+p\sum_{k=1}^{p-1}c_{ik}T_i^{p-k}-f_i
  \tag{22.4.4.1}
\end{equation}
con $c_{ik}\in A$ y $f_i\in A$ para todo $i$ y todo $k$ tal que
$1\leqslant k\leqslant p-1$; además, si $a_i$ es la clase de $f_i$
mod.~$\mathfrak m$, se supone que la familia $(a_i)_{1\leqslant i\leqslant r}$
es $p$-libre sobre $K^p$ («absolutamente $p$-libre»). Entonces:
\begin{enumerate}
  \item[(i)] $A'$ es un anillo local, de ideal maximal $\mathfrak mA'$,
  cuyo cuerpo residual $K'$ es isomorfo a la extensión de $K$ obtenida por
  adjunción de los $a_i^{1/p}$ $(1\leqslant i\leqslant r)$.

  \oldpage[0\textsubscript{IV-1}]{195}

  \item[(ii)] Los elementos $d_Af_i\otimes1_{K'}$ forman una base del núcleo
  $N=\Upsilon^{K'}_{A'/A}$ del homomorfismo canónico
  \[
    \Omega^1_A\otimes_AK'\longrightarrow
    \Omega^1_{A'}\otimes_{A'}K'.
  \]
\end{enumerate}
\end{proposition}

\begin{proof}
(i) Esta vez $A'\otimes_AK=A'/\mathfrak mA'$ es isomorfo a
$K[T_1,\ldots,T_r]/\mathfrak b'$, donde $\mathfrak b'$ es el ideal engendrado
por los polinomios $T_i^p-a_i$; la hipótesis hecha sobre los $a_i$ implica que
este anillo cociente es un cuerpo (21.1.8), de donde las afirmaciones de (i).

Antes de demostrar (ii), estableceremos el

\begin{lemma}[22.4.4.2]
\label{0.22.4.4.2-es}
Sean $\Lambda$ un anillo, $A$ una $\Lambda$-álgebra y $\mathfrak J$ un
ideal de $A$ tal que el anillo $C=A/\mathfrak J$ sea de característica
$p>0$ (lo que equivale a decir que $p\cdot1\in\mathfrak J$ en $A$). Sea
$\pi$ la imagen canónica de $p\cdot1$ en
$\mathfrak J/\mathfrak J^2$. Si $C$ es una $(\Lambda/p\Lambda)$-álgebra
formalmente lisa para las topologías discretas, se tiene una sucesión exacta
canónica
\begin{equation}
\label{0.22.4.4.3-es}
  0\longrightarrow(\mathfrak J/\mathfrak J^2)/C\cdot\pi
  \longrightarrow\Omega^1_{A/\Lambda}\otimes_AC
  \longrightarrow\Omega^1_{C/\Lambda}\longrightarrow0.
  \tag{22.4.4.3}
\end{equation}
\end{lemma}

\begin{proof}[Demostración del lema]
Pongamos, en efecto, $A'=A/pA$, $\mathfrak J'=\mathfrak J/pA$, de modo que
$C=A'/\mathfrak J'$. Como
$A'=A\otimes_\Lambda(\Lambda/p\Lambda)$, se tiene
$\Omega^1_{A/\Lambda}\otimes_AA'=\Omega^1_{A'/\Lambda'}$ salvo
isomorfismo canónico, donde $\Lambda'$ denota el anillo
$\Lambda/p\Lambda$ (20.5.5). Aplicando entonces la sucesión exacta (20.5.14.1) a
la $\Lambda'$-álgebra formalmente lisa $C=A'/\mathfrak J'$, se obtiene la
sucesión exacta
\[
  0\longrightarrow\mathfrak J'/\mathfrak J'^2
  \longrightarrow\Omega^1_{A'/\Lambda'}\otimes_{A'}C
  \longrightarrow\Omega^1_{C/\Lambda'}\longrightarrow0
\]
y como el homomorfismo $\Lambda\to\Lambda'$ es sobreyectivo, se tiene
$\Omega^1_{C/\Lambda'}=\Omega^1_{C/\Lambda}$ salvo isomorfismo canónico
(20.5.15); de donde la sucesión exacta (22.4.4.3) en virtud de lo anterior.
\end{proof}

(ii) Como se tiene $\mathfrak m'^2=0$, denotaremos también por $V'$ el ideal
$\mathfrak m'$ considerado como $K'$-espacio vectorial, y pondremos
\begin{equation}
\label{0.22.4.4.4-es}
  V_0=V/K\cdot p=\mathfrak m/(\mathfrak m^2+pA),\qquad
  V_0'=V'/K'\cdot p=\mathfrak m'/(\mathfrak m'^2+pA').
  \tag{22.4.4.4}
\end{equation}
Aplicando la sucesión exacta (22.4.4.3) al caso en que $\Lambda=\mathbf Z$, a la
$\mathbf Z$-álgebra $A$ (resp. $A'$) y a su ideal $\mathfrak m$
(resp. $\mathfrak m'$), se obtienen, puesto que $K$ (resp. $K'$) es separable sobre
el cuerpo primo $\mathbf Z/p\mathbf Z$, y por tanto una
$(\mathbf Z/p\mathbf Z)$-álgebra formalmente lisa (19.6.1), las sucesiones
exactas
\begin{equation}
\label{0.22.4.4.5-es}
  0\longrightarrow V_0\longrightarrow\Omega^1_A\otimes_AK
  \longrightarrow\Omega^1_K\longrightarrow0,\qquad
  0\longrightarrow V_0'\longrightarrow\Omega^1_{A'}\otimes_{A'}K'
  \longrightarrow\Omega^1_{K'}\longrightarrow0.
  \tag{22.4.4.5}
\end{equation}
Consideremos entonces el diagrama (22.4.4.6), donde las dos columnas centrales son
la segunda sucesión (22.4.4.5) y la primera tensorizada por $K'$; como la
primera sucesión exacta (22.4.4.5) está formada por espacios vectoriales sobre $K$,
las dos columnas centrales de (22.4.4.6) son exactas; además, el
diagrama formado por estas columnas y las flechas horizontales que las
conectan es conmutativo, en virtud de la demostración de (22.4.4.2) y de
la conmutatividad del diagrama (20.5.11.3). La última fila de (22.4.4.6)
es la sucesión exacta obtenida a partir de (20.6.1.1) sustituyendo $A$, $B$,
$C$ por $\mathbf Z$, $K$ y $K'$; la fila central se obtiene por
tensorización con $K'$ de la sucesión exacta de $K$-módulos deducida de
(20.6.1.1) sustituyendo $A$, $B$, $C$ por $\mathbf Z$, $A$, $K$
respectivamente. El diagrama formado por estas dos filas y las flechas
verticales que las conectan es conmutativo en virtud de la conmutatividad de
(20.6.4.3). Observemos, por otra parte, que se está en las condiciones de aplicación
de (22.4.1, (ii)), por lo que la fila superior del diagrama (22.4.4.6) es
exacta, y las columnas extremas del diagrama

\oldpage[0\textsubscript{IV-1}]{196}

\begin{equation}
\label{0.22.4.4.6-es}
  \xymatrix{
    && 0 \ar[d] & 0 \ar[d] && \\
    & 0 \ar[r] & V_0\otimes_KK' \ar[r] \ar[d] & V_0' \ar[r] \ar[d] & 0 & \\
    0 \ar[r] & N \ar[r] \ar[d] & \Omega^1_A\otimes_AK' \ar[r] \ar[d] &
      \Omega^1_{A'}\otimes_{A'}K' \ar[r] \ar[d] &
      \Omega^1_{A'/A}\otimes_{A'}K' \ar[r] \ar[d] & 0 \\
    0 \ar[r] & \Upsilon_{K'/K} \ar[r] \ar[d] &
      \Omega^1_K\otimes_KK' \ar[r] \ar[d] & \Omega^1_{K'} \ar[r] \ar[d] &
      \Omega^1_{K'/K} \ar[r] \ar[d] & 0 \\
    & 0 & 0 & 0 & 0 &
  }
  \tag{22.4.4.6}
\end{equation}
están formadas, respectivamente, por los núcleos y los conúcleos de las flechas
horizontales centrales, lo que termina de probar que el diagrama es
conmutativo. Además:

\begin{lemma}[22.4.4.7]
\label{0.22.4.4.7-es}
Los homomorfismos
\[
  N\longrightarrow\Upsilon_{K'/K}\qquad\text{y}\qquad
  \Omega^1_{A'/A}\otimes_{A'}K'\longrightarrow\Omega^1_{K'/K}
\]
del diagrama (22.4.4.6) son biyectivos.
\end{lemma}

\begin{proof}
En efecto, el diagrama de la serpiente (Bourbaki, \emph{Alg. conmut.}, cap.~I,
§~1, n\textsuperscript{o}~4, prop.~2), aplicado a las dos columnas centrales de
(22.4.4.6), da una sucesión exacta
\[
  0\longrightarrow N\longrightarrow\Upsilon_{K'/K}\longrightarrow0
  \longrightarrow\Omega^1_{A'/A}\otimes_{A'}K'
  \longrightarrow\Omega^1_{K'/K}\longrightarrow0.
\]
Ahora bien, si $t_i$ es la imagen canónica de $T_i$ en $A'$, la relación
$F_i(t_i)=0$, junto con el hecho de que en $A'$ se tiene
$p\cdot1\in\mathfrak mA'\subset\mathfrak m'$, muestra inmediatamente que
$d_{A'}f_i\in\mathfrak m'\Omega^1_{A'}$, por tanto
$d_{A'}f_i\otimes1_{K'}=0$; esto demuestra que los
$d_Af_i\otimes1_{K'}$ pertenecen a $N$. Además, sus imágenes en
$\Upsilon_{K'/K}$ son los $d_Ka_i\otimes1_{K'}$, que forman una base de
$\Upsilon_{K'/K}$ sobre $K'$ (21.4.7); teniendo en cuenta (22.4.4.7), esto termina
de demostrar (22.4.4).
\end{proof}
\end{proof}
\begin{env}[22.4.5]
\label{0.22.4.5-es}
Consideremos ahora dos anillos $A$, $B$ \emph{de característica} $p>0$, y
dos homomorfismos
\[
  A\xrightarrow{i}B\xrightarrow{j}A
\]

\oldpage[0\textsubscript{IV-1}]{197}

que verifican las condiciones de (21.3.1). Sea además $\mathfrak m$ un ideal
\emph{de cuadrado nulo} en $A$; pongamos
\[
  K=A/\mathfrak m,\qquad
  B_{(K)}=B\otimes_AK=B/Bi(\mathfrak m)
\]
y sean
\[
  \varphi:A\longrightarrow K,\qquad
  \psi:B\longrightarrow B_{(K)}
\]
los homomorfismos canónicos.

Como $\varphi$ es sobreyectivo, $\Omega^1_{B_{(K)}/A}$ se identifica
canónicamente con $\Omega^1_{B_{(K)}/K}$ (20.5.15). Observemos que, puesto que
$p\geqslant2$, para todo $x\in\mathfrak m$, se tiene $j(i(x))=x^p=0$, es decir,
$j(i(\mathfrak m))=0$; pasando a los cocientes, se deduce, pues, de $j$ un
homomorfismo
\[
  j':B_{(K)}\longrightarrow A
\]
tal que, si $i'=\psi\circ i:A\longrightarrow B_{(K)}$, se tiene
\[
  j'(i'(x))=x^p\qquad\text{y}\qquad i'(j'(y))=y^p
  \qquad\text{para }x\in A\text{ e }y\in B_{(K)}.
\]
Poniendo, según las notaciones de (21.3.2),
\[
  \Theta_{B_{(K)}/A}
  =\Omega^1_{B_{(K)}/A}\otimes_{B_{(K)}}A_{[j']}
  =\Omega^1_{B_{(K)}/K}\otimes_{B_{(K)}}A_{[j']},
\]
se tiene, pues, (21.3.2.2) un homomorfismo canónico
\begin{equation}
\label{0.22.4.5.1-es}
  \pi_{B_{(K)}/A}:\Theta_{B_{(K)}/A}\longrightarrow\Omega^1_A.
  \tag{22.4.5.1}
\end{equation}

Por otra parte, se deducen de $i$ y $j$, pasando a los cocientes (teniendo en
cuenta que $j(i(\mathfrak m))=0$), homomorfismos
\[
  K\xrightarrow{i_0}B_{(K)}\xrightarrow{j_0}K
\]
de modo que se tiene el diagrama conmutativo
\[
  \xymatrix{
    A \ar[r]^i \ar[d]_\varphi & B \ar[r]^j \ar[d]^\psi & A \ar[d]^\varphi \\
    K \ar[r]_{i_0} & B_{(K)} \ar[r]_{j_0} \ar[ur]_{j'} & K
  }
\]
Está claro, además, que se tiene
\[
  j_0(i_0(s))=s^p\qquad\text{y}\qquad i_0(j_0(t))=t^p
  \qquad\text{para }s\in K\text{ y }t\in B_{(K)}.
\]
Se define, pues, del mismo modo
\[
  \Theta_{B_{(K)}/K}
  =\Omega^1_{B_{(K)}/K}\otimes_{B_{(K)}}K_{[j_0]}
  =\Theta_{B_{(K)}/A}\otimes_AK,
\]
por lo que, tensorizando (22.4.5.1) por $K$, se obtiene un
$K$-\emph{homomorfismo canónico}
\begin{equation}
\label{0.22.4.5.2-es}
  \pi'_{B_{(K)}/K}=\pi_{B_{(K)}/A}\otimes1_K:
  \Theta_{B_{(K)}/K}\longrightarrow\Omega^1_A\otimes_AK.
  \tag{22.4.5.2}
\end{equation}

\oldpage[0\textsubscript{IV-1}]{198}

Se deduce inmediatamente de las definiciones que el diagrama
\[
  \xymatrix{
    \Theta_{B_{(K)}/K} \ar[r]^{\pi'_{B_{(K)}/K}} \ar[dr]_{\pi_{B_{(K)}/K}}
      & \Omega^1_A\otimes_AK \ar[d]^{\varphi_{K/A}} \\
    & \Omega^1_K
  }
\]
es conmutativo, siendo $\pi_{B_{(K)}/K}$ el homomorfismo canónico
correspondiente a $i_0$ y $j_0$ (21.3.2.2). Por \emph{restricción} a
$\Xi_{B_{(K)}/K}=\operatorname{Ker}(\pi_{B_{(K)}/K})$, el homomorfismo
$\pi'_{B_{(K)}/K}$ define, pues, un \emph{homomorfismo canónico}
\begin{equation}
\label{0.22.4.5.3-es}
  \Xi_{B_{(K)}/K}\longrightarrow
  \operatorname{Ker}(\Omega^1_A\otimes_AK\longrightarrow\Omega^1_K)=\Upsilon_{K/A}.
  \tag{22.4.5.3}
\end{equation}
\end{env}

\begin{env}[22.4.6]
\label{0.22.4.6-es}
Siendo las hipótesis y notaciones las de (22.4.5), supongamos además que
$\mathfrak m$ es un ideal \emph{maximal de cuadrado nulo}, de modo que $K$ es un
\emph{cuerpo}, y denotemos por $V$ el ideal $\mathfrak m$ considerado como
$K$-espacio vectorial. Como $K$ es un álgebra formalmente lisa sobre su cuerpo primo
$\mathbf F_p$ (21.5.2) y $A$ es un álgebra sobre $\mathbf F_p$ (de modo
que $\Omega^1_A=\Omega^1_{A/\mathbf F_p}$), se deduce de (20.5.14) que se tiene
una \emph{sucesión exacta}
\begin{equation}
\label{0.22.4.6.1-es}
  0\longrightarrow V\longrightarrow\Omega^1_A\otimes_AK
  \longrightarrow\Omega^1_K\longrightarrow0.
  \tag{22.4.6.1}
\end{equation}
Se tiene, pues, por (22.4.5.3) un homomorfismo denotado por
\begin{equation}
\label{0.22.4.6.2-es}
  \chi'_{B/A}:\Xi_{B_{(K)}/K}\longrightarrow V
  \tag{22.4.6.2}
\end{equation}
que llamaremos \emph{homomorfismo característico} de la $A$-álgebra
$B$ (relativo a $i$, $j$ y $\mathfrak m$). Se deduce de las definiciones que para
que $\chi'_{B/A}$ sea \emph{inyectivo}, es necesario y suficiente que
$\pi'_{B_{(K)}/K}$ sea inyectivo, pues el núcleo de este último está evidentemente
contenido en $\Xi_{B_{(K)}/K}$.
\end{env}

\begin{proposition}[22.4.7]
\label{0.22.4.7-es}
Sean $A$ un anillo local artiniano cuyo ideal maximal $\mathfrak m$ es de
cuadrado nulo, $K=A/\mathfrak m$ su cuerpo residual y $V$ el ideal $\mathfrak m$
considerado como $K$-espacio vectorial; se supone que $A$ es de
característica $p>0$; sean
\begin{equation}
\label{0.22.4.7.1-es}
  F_i(T_i)=T_i^p-f_i\qquad\text{con }f_i\in A\quad(1\leqslant i\leqslant r)
  \tag{22.4.7.1}
\end{equation}
y denotemos por $B$ el anillo cociente de $A[T_1,\ldots,T_r]$ por el ideal
engendrado por los $r$ polinomios $F_i$. Entonces:
\begin{enumerate}
  \item[\rm(i)] $B$ es un anillo local.
  \item[\rm(ii)] Si $\mathfrak m'$ es el ideal maximal de $B$,
  $K'=B/\mathfrak m'$ su cuerpo residual y $V'$ el $B$-módulo
  $\mathfrak m'/\mathfrak m'^2$ considerado como $K'$-espacio vectorial, entonces,
  para que $\operatorname{rg}_K(V)=\operatorname{rg}_{K'}(V')$, es necesario y suficiente que el homomorfismo
  característico de $K$-espacios vectoriales
  \begin{equation}
  \label{0.22.4.7.2-es}
    \chi'_{B/A}:\Xi_{B_{(K)}/K}\longrightarrow V
    \tag{22.4.7.2}
  \end{equation}
  (cf. (22.4.6.2)) sea inyectivo.
\end{enumerate}
\end{proposition}

\begin{proof}
Denotemos también por $a_i$ $(1\leqslant i\leqslant r)$ la clase de $f_i$
mod.~$\mathfrak m$; los $a_i$ ya no son necesariamente $p$-independientes
sobre $K^p$, pero siempre se puede suponer que la subfamilia

\oldpage[0\textsubscript{IV-1}]{199}

$(a_i)_{s+1\leqslant i\leqslant r}$ es $p$-libre sobre $K^p$ y forma una
$p$-base sobre $K^p$ del cuerpo $K^p(a_1,\ldots,a_r)$, de modo que se tiene
\[
  a_i\in K^p(a_{s+1},\ldots,a_r)\qquad\text{para }1\leqslant i\leqslant s.
\]
Denotemos por $A''$ el anillo cociente de $A[T_{s+1},\ldots,T_r]$ por el ideal
engendrado por los $F_i$ de índice $i\geqslant s+1$; entonces, por (22.4.4), $A''$
es un anillo local cuyo ideal maximal es $\mathfrak m''=\mathfrak mA''$
(por tanto de cuadrado nulo) y cuyo cuerpo residual es
$K''=K(a_{s+1}^{1/p},\ldots,a_r^{1/p})$. Se tiene, pues,
$a_i\in K''^p$ para $1\leqslant i\leqslant s$ y, por consiguiente, en $A''$,
el elemento $f_i$, para $1\leqslant i\leqslant s$, puede escribirse
\begin{equation}
\label{0.22.4.7.3-es}
  f_i=g_i^p+h_i
  \tag{22.4.7.3}
\end{equation}
donde $g_i\in A''$ y $h_i\in\mathfrak m''$ (como $\mathfrak m''^2=0$ y
$p\mathfrak m''=0$, es inmediato que $h_i$ está determinado de manera
\emph{única} por estas condiciones). Sustituyendo $T_i$ por $T_i+g_i$, se ve, pues,
que si se pone
\[
  G_i(T_i)=T_i^p+\sum_{k=1}^{p-1}\binom pk g_i^{p-k}T_i^k-h_i
  \qquad(1\leqslant i\leqslant s),
\]
$B$ es el anillo cociente de $A''[T_1,\ldots,T_s]$ por el ideal engendrado por los
$G_i$; ahora bien, todos los coeficientes de $G_i$, salvo el coeficiente dominante,
pertenecen a $\mathfrak m''$, por lo que se está en la situación de (22.4.3),
$B$ es un anillo local, y su cuerpo residual $K'$ es \emph{isomorfo} a
$K''$, lo que demuestra (i) (que, por lo demás, resulta directamente de que
$B^p\subset A$, y por tanto solo hay un ideal de $B$ por encima de
$\mathfrak m$).

Observemos ahora que para que (22.4.7.2) sea inyectivo, es necesario y suficiente
que lo sea
\[
  \pi'_{B_{(K)}/K}:\Omega^1_{B_{(K)}/K}\otimes_{B_{(K)}}K
  \longrightarrow\Omega^1_A\otimes_AK
\]
(por (22.4.6)). Ahora bien, $B_{(K)}=B/\mathfrak mB$ es el álgebra cociente de
$C=K[T_1,\ldots,T_r]$ por el ideal engendrado por los polinomios
$\overline F_i=T_i^p-a_i$, y como $d_{C/K}\overline F_i=0$, se deduce de
(20.5.13) que $\Omega^1_{B_{(K)}/K}$ es un $B_{(K)}$-módulo libre que tiene por
\emph{base} los $d_{B_{(K)}/K}t_i$, donde $t_i$ es la imagen canónica de $T_i$
$(1\leqslant i\leqslant r)$. Pero, por definición (22.4.5), la imagen canónica por
$j':B_{(K)}\longrightarrow A$ de una clase mod.~$\mathfrak mB$ de un elemento
$x\in B$ es el elemento $x^p\in A$; se deduce inmediatamente, puesto que $t_i^p=f_i$,
que la imagen de $d_{B_{(K)}/K}t_i\otimes1$ por $\pi'_{B_{(K)}/K}$ es
$d_Af_i\otimes1$ en $\Omega^1_A\otimes_AK$; la condición de que
$\pi'_{B_{(K)}/K}$ sea inyectivo equivale, pues, a que los
$d_Af_i\otimes1_K$ sean \emph{linealmente independientes sobre $K$}, o también
a que sus imágenes $d_Af_i\otimes1_{K'}$ sean \emph{linealmente independientes
sobre $K'$} en $\Omega^1_A\otimes_AK'$.

Ahora bien, si se aplica (22.4.4) a la $A$-álgebra $A''$, se ve que el núcleo del
homomorfismo canónico
$\Omega^1_A\otimes_AK'\longrightarrow\Omega^1_{A''}\otimes_{A''}K'$ tiene por
base los $d_Af_i\otimes1_{K'}$ para $s+1\leqslant i\leqslant r$. La condición
precedente equivale, pues, también a decir que los $d_{A''}f_i\otimes1_{K'}$ son
linealmente independientes en $\Omega^1_{A''}\otimes_{A''}K'$ para
$1\leqslant i\leqslant s$. Ahora bien, se tiene $d_{A''}f_i=d_{A''}h_i$ por (22.4.7.3);
por otra parte, si $V''$ es el $K'$-espacio vectorial $\mathfrak m''$, se tiene
$\operatorname{rg}_{K'}(V'')=\operatorname{rg}_K(V)$ (22.4.1, (ii)), y la condición
$\operatorname{rg}_{K'}(V')=\operatorname{rg}_K(V)$ equivale, pues, a
$\operatorname{rg}_{K'}(V'')=\operatorname{rg}_{K'}(V')$. Pero se deduce de (22.4.3) que esta última
relación equivale a que las clases $h_i$

\oldpage[0\textsubscript{IV-1}]{200}

en $V''$ sean \emph{linealmente independientes sobre $K'$}. Ahora bien, en
$\Omega^1_{A''}\otimes_{A''}K'$, los $d_{A''}h_i\otimes1_{K'}$ son las imágenes
canónicas (20.5.11.2) de los $h_i$ por la inyección
$V''\longrightarrow\Omega^1_{A''}\otimes_{A''}K'$ (cf. (22.4.6.1)), y esto
termina de demostrar (22.4.7).
\end{proof}

Además, se ha demostrado:

\begin{corollary}[22.4.7.4]
\label{0.22.4.7.4-es}
Para que $\operatorname{rg}_K(V)=\operatorname{rg}_{K'}(V')$, es necesario y suficiente que los elementos
$d_Af_i\otimes1_K$ $(1\leqslant i\leqslant r)$ sean linealmente
independientes en el $K$-espacio vectorial $\Omega^1_A\otimes_AK$.
\end{corollary}

\begin{remark}[22.4.8]
\label{0.22.4.8-es}
Se pueden definir los homomorfismos (22.4.5.2) y (22.4.6.2) en condiciones
algo distintas. Sean $A$ un anillo, $\mathfrak m$ un ideal \emph{de cuadrado
nulo} en $A$, $K=A/\mathfrak m$ y $p$ un número primo tal que
$p\cdot1\in\mathfrak m$ (es decir, $K$ es de característica $p$, pero no
necesariamente $A$). Sea, por otra parte, $B$ una $A$-álgebra que sea un $A$-módulo
\emph{fielmente plano} (de modo que $A\subset B$), y supongamos que se tiene
\begin{equation}
\label{0.22.4.8.1-es}
  y^p\in A+pB\subset A+\mathfrak mB
  \tag{22.4.8.1}
\end{equation}
para todo $y\in B$. Si $y^p=x+pz$ con $x\in A$, $z\in B$, la \emph{clase}
mod.~$(A\cap pB)$ de $x$ está bien determinada, y como $A\cap pB=pA$ puesto que
$B$ es un $A$-módulo fielmente plano, se ha definido así una aplicación
canónica $j:B\longrightarrow A/pA$, tal que, para $x\in A$, $j(x)$ sea la
clase de $x^p$ mod.~$pA$. Por otra parte, si $y'\in B$ y
$y'^p=x'+pz'$ con $x'\in A$, $z'\in B$, es inmediato que los elementos
$(y+y')^p-x-x'$ y $(yy')^p-xx'$ pertenecen a $pB$, en virtud de que los
coeficientes binomiales $\binom pi$ son múltiplos de $p$ para
$1\leqslant i\leqslant p-1$. La aplicación $j$ es, pues, un
\emph{homomorfismo de anillos}. Por composición, se deduce de $j$ un
homomorfismo $j':B\xrightarrow{j}A/pA\longrightarrow A/\mathfrak m=K$ y, por
consiguiente, un homomorfismo $j_0:B_{(K)}\longrightarrow K$ tal que $j'$ sea el
compuesto $B\longrightarrow B_{(K)}=B\otimes_AK\xrightarrow{j_0}K$; además,
si $i_0:K\longrightarrow B_{(K)}$ es la aplicación canónica, $i_0$ y $j_0$
verifican las condiciones de (21.3.1) para los anillos \emph{de característica
$p$}, $K$ y $B_{(K)}$. Dicho esto, para definir un $K$-homomorfismo
$\pi'_{B_{(K)}/K}:\Omega^1_{B_{(K)}/K}\otimes_{B_{(K)}}K_{[j_0]}
\longrightarrow\Omega^1_A\otimes_AK$, basta (20.4.8) definir una
$K$-derivación
\begin{equation}
\label{0.22.4.8.2-es}
  D_0:B_{(K)}\longrightarrow\Omega^1_A\otimes_AK
  \tag{22.4.8.2}
\end{equation}
donde el segundo miembro se considera como un $B_{(K)}$-módulo por medio de $j_0$.
Para ello basta definir una $A$-derivación
\begin{equation}
\label{0.22.4.8.3-es}
  D:B\longrightarrow\Omega^1_A\otimes_AK
  \tag{22.4.8.3}
\end{equation}
que \emph{se anule en} $\mathfrak mB$ (considerándose el segundo miembro como
$B$-módulo por medio de $j'$). Ahora bien, si se tiene la relación $y^p=x+pz$ con
$y\in B$, $x\in A$, $z\in B$, el elemento $d_Ax\otimes1_K$ de
$\Omega^1_A\otimes_AK$ solo depende de la clase $j(y)$ mod.~$pA$ de $x$
(por tanto solo de $y$), pues para todo $t\in A$ se tiene
$d_A(pt)=p\,d_A(t)$, por tanto $d_A(pt)\otimes1=0$ en
$\Omega^1_A\otimes_AK=\Omega^1_A/\mathfrak m\Omega^1_A$, puesto que
$p\cdot1\in\mathfrak m$. Se puede tomar, pues,
\[
  D(y)=d_Ax\otimes1_K.
\]
El hecho de que $D$ sea una derivación se deduce fácilmente de que $j$ es un
homomorfismo de anillos; además, para $t\in A$, se tiene
$d_A(t^px)=t^p d_A(x)+pt^{p-1}x d_A(t)$, y de ello se deduce que
$D(ty)=tD(y)$, por tanto $D$ es una $A$-derivación; finalmente, si se tiene además
$t\in\mathfrak m$, se tiene

\oldpage[0\textsubscript{IV-1}]{201}

$D(ty)=tD(y)=0$, puesto que $\Omega^1_A\otimes_AK$ está anulado por $\mathfrak m$.
Se ha definido así el homomorfismo (22.4.5.2) buscado, y es inmediato
verificar que el homomorfismo compuesto
\[
  \Theta_{B_{(K)}/K}\xrightarrow{\pi'_{B_{(K)}/K}}
  \Omega^1_A\otimes_AK\longrightarrow\Omega^1_K
\]
es todavía el homomorfismo canónico $\pi_{B_{(K)}/K}$ (relativo a $i_0$ y
$j_0$). Se tiene, pues, también un homomorfismo canónico análogo a (22.4.5.3).

Si ahora $K$ es un \emph{cuerpo}, $A/pA$ es una $\mathbf F_p$-álgebra,
y se tiene la sucesión exacta (20.5.14)
\[
  0\longrightarrow V/pA\longrightarrow
  \Omega^1_{A/pA}\otimes_{A/pA}K\longrightarrow\Omega^1_K\longrightarrow0
\]
y, por otra parte, se tiene la sucesión exacta (20.5.12.1)
\[
  pA\longrightarrow\Omega^1_A\otimes_A(A/pA)
  \longrightarrow\Omega^1_{A/pA}\longrightarrow0
\]
que muestra, tensorizando con $K$, que los $K$-espacios vectoriales
$\Omega^1_A\otimes_AK$ y
$\Omega^1_{A/pA}\otimes_{A/pA}K$ son isomorfos. El análogo de (22.4.6.2) es,
por tanto, aquí un homomorfismo
\begin{equation}
\label{0.22.4.8.4-es}
  \Xi_{B_{(K)}/K}\longrightarrow V/pA=\mathfrak m/(\mathfrak m^2+pA).
  \tag{22.4.8.4}
\end{equation}

Dicho esto, \emph{el criterio (22.4.7), en el que se sustituyen $V$ y $V'$ por $V/pA$
y $V'/pB$ respectivamente, sigue siendo válido suponiendo solamente que $K$ sea
de característica $p>0$} (es decir, que $p\cdot1\in\mathfrak m$ en
$A$), \emph{y que los $F_i(T_i)$ sean de la forma más general
(22.4.4.1)}: en efecto, $B$ es un $A$-módulo libre (por tanto fielmente plano), y
basta retomar el razonamiento de (22.4.7), sustituyendo (22.4.6.2)
por (22.4.8.4) y $V''$ por $V''/pA''$.
\end{remark}

\subsection{Álgebras geométricamente regulares y álgebras formalmente lisas}
\label{subsection:0.22.5-es}

\begin{proposition}[22.5.1]
\label{0.22.5.1-es}
Sean $A$ un anillo local regular y $A'$ un anillo local que contiene a $A$ y que
es una $A$-álgebra finita, $\mathfrak m$ (resp. $\mathfrak m'$) el ideal maximal
de $A$ (resp. $A'$), y $K=A/\mathfrak m$ (resp. $K'=A'/\mathfrak m'$) su cuerpo
residual. Para que $A'$ sea un anillo regular, es necesario y suficiente que se
tenga
\begin{equation}
\label{0.22.5.1.1-es}
  \operatorname{rg}_K(\mathfrak m/\mathfrak m^2)
  =\operatorname{rg}_{K'}(\mathfrak m'/\mathfrak m'^2).
  \tag{22.5.1.1}
\end{equation}
\end{proposition}

\begin{proof}
En efecto, el primer miembro de (22.5.1.1) es $\dim(A)$ (17.1.1), y como
$A\subset A'$ y $A'$ es una $A$-álgebra finita,
$\dim(A')=\dim(A)$ (16.1.5); la igualdad (22.5.1.1) es, pues, necesaria y
suficiente para que $A'$ sea regular (17.1.1).
\end{proof}

\begin{remark}[22.5.2]
\label{0.22.5.2-es}
\begin{enumerate}
  \item[\rm(i)] Sean $A_1=A/\mathfrak m^2$,
  $A'_1=A'\otimes_AA_1=A'/\mathfrak m^2A'$; como
  $\mathfrak m^2A'\subset\mathfrak m'^2$,
  $\mathfrak m'_1=\mathfrak m'/\mathfrak m^2A'$ es el ideal maximal de $A'_1$
  y $\mathfrak m'_1/\mathfrak m_1'^2$ es un $K'$-espacio vectorial isomorfo
  a $\mathfrak m'/\mathfrak m'^2$; la condición de regularidad de $A'$ solo
  depende, pues, de la estructura de la $A_1$-álgebra $A'_1$, lo que nos
  permitirá a continuación aplicar los resultados preliminares de (22.4) sobre
  las álgebras finitas sobre anillos artinianos.
  \item[\rm(ii)] Se deduce de la demostración de (22.5.1) y de (16.2.6.1)
  que se tiene en todo caso
  \begin{equation}
  \label{0.22.5.2.1-es}
    \operatorname{rg}_K(\mathfrak m/\mathfrak m^2)
    \leqslant\operatorname{rg}_{K'}(\mathfrak m'/\mathfrak m'^2).
    \tag{22.5.2.1}
  \end{equation}

  \oldpage[0\textsubscript{IV-1}]{202}

  \item[\rm(iii)] Supongamos que $A$ y $A'$ verifican las hipótesis generales
  de (22.4.1). Se sabe (19.8.10) que existe un anillo local \emph{regular}
  $B$, de ideal maximal $\mathfrak n$, tal que $A$ es isomorfo a
  $B/\mathfrak n^2$; denotemos por $G_i\in B[T_i]$ un polinomio mónico cuya
  imagen canónica en $A[T_i]$ sea $F_i$ $(1\leqslant i\leqslant r)$;
  entonces, si $B'$ es el anillo cociente de $B[T_1,\ldots,T_r]$ por el ideal
  engendrado por los $G_i$, está claro que $B'$ es un $B$-módulo libre de rango
  $d$ y que $A'=B'\otimes_BA=B'/\mathfrak n^2B'$; como $A'$ se supone un
  anillo local, lo mismo sucede con $B'$, y si $\mathfrak n'$ es el ideal
  maximal de $B'$, $B'/\mathfrak n'$ es isomorfo a $K'$ y
  $\mathfrak n'/\mathfrak n'^2$ es isomorfo a $\mathfrak m'/\mathfrak m'^2$ como
  espacio vectorial sobre $K'$; la desigualdad (22.5.2.1) muestra, pues, que
  se tiene, bajo las hipótesis de (22.4.1),
  \begin{equation}
  \label{0.22.5.2.2-es}
    \operatorname{rg}_K(V)\leqslant\operatorname{rg}_{K'}(\mathfrak m'/\mathfrak m'^2).
    \tag{22.5.2.2}
  \end{equation}
\end{enumerate}
\end{remark}

\begin{corollary}[22.5.3]
\label{0.22.5.3-es}
Sean $A$ un anillo local regular, $\mathfrak m$ su ideal maximal,
$T_i$ $(1\leqslant i\leqslant r)$ indeterminadas y, para cada $i$, sea
\begin{equation}
\label{0.22.5.3.1-es}
  F_i(T_i)=T_i^{d_i}+\sum_{1\leqslant k\leqslant d_i}c_{ik}T_i^{d_i-k}
  \tag{22.5.3.1}
\end{equation}
un polinomio mónico de $A[T_i]$, tal que $d_i\geqslant2$ para todo $i$ y
$c_{ik}\in\mathfrak m$. Sea $A'$ el cociente del anillo de polinomios
$A[T_1,\ldots,T_r]$ por el ideal engendrado por los $r$ polinomios $F_i$. Entonces:
\begin{enumerate}
  \item[\rm(i)] $A'$ es un anillo local, y su cuerpo residual $K'$ es
  isomorfo al cuerpo residual $K$ de $A$.
  \item[\rm(ii)] Para que $A'$ sea un anillo regular, es necesario y suficiente
  que las clases $\xi_i$ mod.~$\mathfrak m^2$ de los elementos $c_{i,d_i}$
  $(1\leqslant i\leqslant r)$ sean linealmente independientes sobre $K$.
\end{enumerate}
\end{corollary}

\begin{proof}
En efecto, con las notaciones de (22.5.2, (i)), si $\overline F_i$ denota el
polinomio de $A_1[T_i]$ obtenido aplicando a los coeficientes de $F_i$
el homomorfismo $A\longrightarrow A_1$, $A'_1$ se define a partir de $A_1$ y
de los polinomios $\overline F_i$ como en (22.4.1), y la conclusión se deduce,
por tanto, de (22.5.1) y de (22.4.3).
\end{proof}

\begin{theorem}[22.5.4]
\label{0.22.5.4-es}
Sean $A$ un anillo local regular y $K$ su cuerpo residual; se supone que
$K$ es de característica $p>0$; sean
\begin{equation}
\label{0.22.5.4.1-es}
  F_i(T_i)=T_i^p+p\sum_{k=1}^{p-1}c_{ik}T_i^{p-k}-f_i
  \tag{22.5.4.1}
\end{equation}
con $c_{ik}\in A$ y $f_i\in A$ $(1\leqslant i\leqslant r)$; sea $B$
el anillo cociente de $A[T_1,\ldots,T_r]$ por el ideal engendrado por los $r$
polinomios $F_i$. Entonces $B$ es un anillo local, y las condiciones siguientes
son equivalentes:
\begin{enumerate}
  \item[a)] $B$ es regular;
  \item[b)] los elementos $d_Af_i\otimes1_K$ $(1\leqslant i\leqslant r)$ son
  linealmente independientes en el $K$-espacio vectorial
  $\Omega^1_A\otimes_AK$;
  \item[c)] si $B'=B/\mathfrak m^2B$ y si se pone
  $\Xi_{B_{(K)}/K}=\Xi_{B'_{(K)}/K}$ (22.4.5), el homomorfismo
  característico $\Xi_{B_{(K)}/K}\longrightarrow\mathfrak m/\mathfrak m^2$
  (22.4.6.2) es inyectivo.
\end{enumerate}
\end{theorem}

\begin{proof}
En efecto, $B'$ se define también a partir de $A_1=A/\mathfrak m^2$ y de los
polinomios $\overline F_i$ obtenidos aplicando a los $F_i$ el homomorfismo
$A\longrightarrow A_1$. Teniendo en cuenta que $\Omega^1_A\otimes_AK$ es
canónicamente isomorfo a $\Omega^1_{A_1}\otimes_{A_1}K$ (20.5.12, (ii)), el
teorema se deduce de (22.5.1) y de (22.4.7.4), utilizando la observación
(22.4.8) cuando $A_1$ no es de característica $p$.
\end{proof}

\oldpage[0\textsubscript{IV-1}]{203}

\begin{env}[22.5.5]
\label{0.22.5.5-es}
Sean $k$ un cuerpo de característica $p>0$, $A$ un anillo local y $k'$ una
extensión \emph{radicial finita} de $k$ tal que $k'^p\subset k$,
$\mathfrak m$ el ideal maximal de $A$ y $K=A/\mathfrak m$ su cuerpo residual;
recordemos que $A'=A\otimes_kk'$ es un anillo local cuyo cuerpo residual es
el mismo que el del anillo local $K\otimes_kk'$ (Bourbaki,
\emph{Alg. conmut.}, cap.~V, §~2, n\textsuperscript{o}~3, lema~4). Observemos que
se tiene $A'_{(K)}=A'\otimes_AK=K\otimes_kk'$, y por consiguiente (21.3.6)
\begin{equation}
\label{0.22.5.5.1-es}
  \Theta_{A'_{(K)}/K}=\Theta_{k'/k}\otimes_kK
  \tag{22.5.5.1}
\end{equation}
salvo isomorfismo canónico. Además, se sabe (21.4.7) que se tiene
$\Xi_{k'/k}=0$, es decir, el homomorfismo canónico (21.3.2.2)
\[
  \pi_{k'/k}:\Theta_{k'/k}\longrightarrow\Omega^1_k
\]
es inyectivo, de modo que se tiene una \emph{inyección canónica}
\begin{equation}
\label{0.22.5.5.2-es}
  \Theta_{A'_{(K)}/K}\longrightarrow\Omega^1_k\otimes_kK
  \tag{22.5.5.2}
\end{equation}
que se explicita de la manera siguiente (teniendo en cuenta (22.5.5.1)): para todo
$x'\in k'$, al elemento
$d_{A'_{(K)}/K}(x'\otimes1_K)\otimes1_K$ de
$\Omega^1_{A'_{(K)}/K}\otimes_{A'_{(K)}}K$ se hace corresponder el elemento
$d_k(x'^p)\otimes1_K$ de $\Omega^1_k\otimes_kK$.
\end{env}

\begin{lemma}[22.5.6]
\label{0.22.5.6-es}
Con las notaciones de (22.5.5):
\begin{enumerate}
  \item[\rm(i)] El núcleo del homomorfismo canónico
  $\pi_{A'_{(K)}/K}$ (21.3.2.2) se identifica por (22.5.5.2) con
  $(\Theta_{k'/k}\otimes_kK)\cap\Upsilon_{K/k}$.
  \item[\rm(ii)] Sean $A_1=A/\mathfrak m^2$,
  $A'_1=A_1\otimes_kk'=A'/\mathfrak m^2A'$; entonces el homomorfismo
  característico $\chi'_{A'_1/A_1}$ (22.4.6.2) se identifica con la restricción del
  homomorfismo característico
  $\chi_{A/k}:\Upsilon_{K/k}\longrightarrow V=\mathfrak m/\mathfrak m^2$
  (20.6.24) a $(\Theta_{k'/k}\otimes_kK)\cap\Upsilon_{K/k}$.
\end{enumerate}
\end{lemma}

\begin{proof}
(i) En virtud de (20.6.9), aplicado sustituyendo $A$ por el cuerpo primo, $B$
por $k$, $C$ por $K$ y $E$ por $A_1=A/\mathfrak m^2$ (extensión de $K$ por
$V=\mathfrak m/\mathfrak m^2$), basta (cf. (22.5.5)) verificar que si
$x'\in k'$ y $z\in A_1$ es un elemento cuya clase mod.
$\mathfrak m/\mathfrak m^2$ es $x'^p\in k$, entonces la imagen de
$d_k(x'^p)\otimes1_K$ por el homomorfismo canónico
$\Omega^1_k\otimes_kK\longrightarrow\Omega^1_{A_1}\otimes_{A_1}K$ es
$d_{A_1}z\otimes1_K$, lo que se deduce de las definiciones.

(ii) Como $A'_{(K)}=K\otimes_kk'=(A'_1)_{(K)}$ (en virtud de la relación
$K=A_1/(\mathfrak m/\mathfrak m^2)$), $\Theta_{A'_{(K)}/K}$ se identifica con
$\Theta_{(A'_1)_{(K)}/K}$ y la verificación resulta del mismo cálculo que en
(i) y de la definición (22.4.5.2) de $\pi'_{(A'_1)_{(K)}/K}$.
\end{proof}

\begin{proposition}[22.5.7]
\label{0.22.5.7-es}
Sean $k$ un cuerpo de característica $p>0$ y $A$ una $k$-álgebra que es un
anillo local regular, $\mathfrak m$ su ideal maximal y $K=A/\mathfrak m$ su
cuerpo residual; sea $V=\mathfrak m/\mathfrak m^2$, considerado como
$K$-espacio vectorial. Sean $k'$ una extensión finita de $k$ tal que
$k'^p\subset k$, y $A'=A\otimes_kk'$; para que el anillo local $A'$ sea
regular, es necesario y suficiente que la restricción del homomorfismo
característico $\chi_{A/k}:\Upsilon_{K/k}\longrightarrow V$ (20.6.24) a
$(\Theta_{k'/k}\otimes_kK)\cap\Upsilon_{K/k}$ sea inyectiva.
\end{proposition}

\begin{proof}
Se trata de demostrar que la condición del enunciado equivale a
(22.5.1.1), denotando por $\mathfrak m'$ el ideal maximal de $A'$ y por
$K'=A'/\mathfrak m'$ su cuerpo residual. Con las notaciones de la observación
(22.5.2, (i)), se tiene $A'_1=A_1\otimes_kk'$. Ahora bien, sea
$(a_i)_{1\leqslant i\leqslant r}$ una $p$-base de $k'$ sobre $k$, de modo que
$a_i^p=b_i\in k$, y $k'$ es isomorfo al cociente del anillo de

\oldpage[0\textsubscript{IV-1}]{204}

polinomios $k[T_1,\ldots,T_r]$ por el ideal engendrado por los polinomios
$F_i=T_i^p-b_i$ $(1\leqslant i\leqslant r)$; se deduce que $A'_1$ es
isomorfo al cociente del anillo de polinomios $A_1[T_1,\ldots,T_r]$ por
el ideal engendrado por los $F_i$. Se puede aplicar entonces el criterio (22.4.7), y
la condición del enunciado equivale, por tanto, a que $\chi'_{A'_1/A_1}$ sea
inyectivo; se concluye con ayuda de (22.5.6).
\end{proof}

Podemos ahora demostrar el recíproco de (19.6.5):

\begin{theorem}[22.5.8]
\label{0.22.5.8-es}
Sean $k$ un cuerpo, $p$ su exponente característico y $A$ una $k$-álgebra
local noetheriana. Las condiciones siguientes son equivalentes:
\begin{enumerate}
  \item[a)] $A$ es una $k$-álgebra formalmente lisa (para su topología
  preádica).
  \item[b)] $A$ es geométricamente regular sobre $k$.
  \item[b')] Para toda extensión finita $k'$ de $k$ tal que $k'^p\subset k$,
  $A'=A\otimes_kk'$ es un anillo regular.
\end{enumerate}
\end{theorem}

\begin{proof}
Teniendo en cuenta (19.6.6), se puede reducir al caso $p>1$.

El hecho de que a) implique b) no es otro que (19.6.5), y b) implica
trivialmente b'). Demostremos que b') implica que se satisfacen las condiciones de
(22.2.2); esto es evidente para la primera (17.1.1), puesto que b') implica
en primer lugar que $A$ es regular. Por otra parte, sea $(x_\alpha)_{\alpha\in I}$ una
$p$-base de $k^{1/p}$ sobre $k$; se sabe que los elementos $d_k(x_\alpha^p)$
forman una base del $k$-espacio vectorial $\Omega^1_k$ (21.4.5), por lo que los
elementos $d_k(x_\alpha^p)\otimes1_K$ forman una base del $K$-espacio vectorial
$\Omega^1_k\otimes_kK$. Se concluye (cf. (22.5.5)) que cuando $k'$ recorre
el conjunto de las subextensiones de $k^{1/p}$, finitas sobre $k$, la familia de
subespacios $\Theta_{k'/k}\otimes_kK$ de $\Omega^1_k\otimes_kK$ es filtrante
creciente y tiene por unión $\Omega^1_k\otimes_kK$. Se deduce entonces de
(22.5.7) que la condición b') implica que $\chi_{A/k}$ es inyectivo, lo que
no es sino la condición (ii) de (22.2.2).
\end{proof}

\begin{corollary}[22.5.9]
\label{0.22.5.9-es}
Sean $k$ un cuerpo y $A$ una $k$-álgebra local noetheriana formalmente
lisa. Entonces, para todo ideal primo $\mathfrak p$ de $A$, el anillo local
$A_{\mathfrak p}$ es una $k$-álgebra formalmente lisa (para la topología
$\mathfrak p$-preádica).
\end{corollary}

\begin{proof}
En efecto, con las notaciones de (22.5.8, b')), basta demostrar que
el anillo local $A_{\mathfrak p}\otimes_kk'$ es regular; pero como es un
anillo de fracciones del anillo local $A\otimes_kk'$ y este último es
regular por hipótesis, lo mismo sucede con $A_{\mathfrak p}\otimes_kk'$
(17.3.6).
\end{proof}

\begin{remark}[22.5.10]
\label{0.22.5.10-es}
\begin{enumerate}
  \item[\rm(i)] La conclusión de (22.5.8) falla cuando, en la
  condición b'), uno se limita a las extensiones \emph{monógenas} $k'=k(x)$ de $k$
  (con $x^p\in k$). Puede suceder, en efecto, que para todas estas extensiones $k'$,
  se tenga
  $\Upsilon_{K/k}\cap(\Theta_{k'/k}\otimes_kK)=0$, aunque
  $\Upsilon_{K/k}\neq0$; esto significa que para todo $x\in k^{1/p}$ tal que
  $x\notin k$, debe tenerse
  $d_k(x^p)\otimes1_K\notin\Upsilon_{K/k}$, o también $d_K(x^p)\neq0$,
  es decir, $x^p\notin K^p$; en otras palabras, debe tenerse
  $K^p\cap k=k^p$, aunque $K$ sea una extensión \emph{inseparable} de $k$.
  Ahora bien, se construyen fácilmente ejemplos de tales extensiones: partamos de un
  cuerpo perfecto $k_0$ de característica $p>0$, sean $X$, $Y$, $Z$ tres
  indeterminadas y pongamos $k=k_0(X,Y)$ y
  $K=k(X^{1/p}+ZY^{1/p})$; se verifica fácilmente que $K$ satisface las condiciones
  precedentes, por lo que $\Upsilon_{K/k}\neq0$. Apliquemos ahora el lema
  (22.2.5.2) tomando $V=K$ y como $h$ el homomorfismo nulo; $A$ es entonces
  un anillo de valoración discreta que tiene a $K$ por cuerpo residual, con
  $\chi_{A/k}=0$, y no es, por tanto, una $k$-álgebra formalmente lisa por
  (22.2.2); sin embargo, $A\otimes_kk'$ es un anillo regular para toda
  extensión \emph{monógena} $k'\subset k^{1/p}$ de $k$.

  \oldpage[0\textsubscript{IV-1}]{205}

  \item[\rm(ii)] El corolario (22.5.9) lleva a considerar la cuestión
  siguiente: si $A$ es una $k$-álgebra noetheriana, ¿bajo qué condiciones
  el conjunto de los ideales primos $\mathfrak p\in\operatorname{Spec}(A)$ tales que
  $A_{\mathfrak p}$ sea una $k$-álgebra formalmente lisa es
  \emph{abierto}? Abordaremos más adelante algunos casos particulares.
\end{enumerate}
\end{remark}
\subsection{Criterio jacobiano de Zariski}
\label{subsection:0.22.6-es}

\begin{theorem}[22.6.1]
\label{0.22.6.1-es}
\textnormal{(criterio jacobiano de lisitud formal).}
Sean $A$, $B$ dos anillos topológicos,
$u:A\longrightarrow B$ un homomorfismo continuo que hace de $B$ una
$A$-álgebra formalmente lisa, $\mathfrak J$ un ideal de $B$ (no
necesariamente cerrado), $C=B/\mathfrak J$ el $A$-álgebra topológica cociente.
Para que $C$ sea una $A$-álgebra formalmente lisa, es necesario y suficiente que
el homomorfismo canónico (cf. (20.5.11.2))
\begin{equation}
\label{0.22.6.1.1-es}
  \delta_{C/B/A}:\mathfrak J/\mathfrak J^2
  \longrightarrow\Omega^1_{B/A}\otimes_BC
  \tag{22.6.1.1}
\end{equation}
sea \emph{formalmente invertible por la izquierda} (cf. (19.1.5)).
\end{theorem}

\begin{proof}
En efecto, decir que $B$ (resp. $C$) es una $A$-álgebra formalmente lisa
significa (19.4.4) que se tiene
$\operatorname{Exalcotop}_A(B,L)=0$ (resp.
$\operatorname{Exalcotop}_A(C,L)=0$) para todo $B$-módulo (resp. $C$-módulo)
discreto $L$ anulado por un ideal abierto de $B$ (resp. $C$). Como, por hipótesis,
$\operatorname{Exalcotop}_A(B,L)=0$ para todo $C$-módulo discreto $L$ anulado
por un ideal abierto de $C$, decir que $C$ es una $A$-álgebra formalmente lisa
significa, pues, que el homomorfismo canónico
$\operatorname{Exalcotop}_A(C,L)\longrightarrow
\operatorname{Exalcotop}_A(B,L)$ es inyectivo, y por tanto el teorema resulta de
(20.7.8).
\end{proof}

Recordemos (19.5.3) que, cuando $B$ y $C$ son $A$-álgebras formalmente
lisas, $\mathfrak J/\mathfrak J^2$ es un $C$-módulo topológico
\emph{formalmente proyectivo}; además, los homomorfismos canónicos
$\varphi_n:\mathbf S_C^n(\mathfrak J/\mathfrak J^2)
\longrightarrow\mathfrak J^n/\mathfrak J^{n+1}$ son
\emph{bimorfismos formales} cuando se supone que $B$ es un anillo
\emph{preadmisible}.
\begin{corollary}[22.6.2]
\label{0.22.6.2-es}
Siendo las hipótesis y notaciones las de (22.6.1), supongamos además que
en $B$ el cuadrado de todo ideal abierto es abierto, y que sobre $\mathfrak J$ la
topología inducida por la de $B$ es idéntica a la topología deducida de
la de $B$ (19.0) (se observará que estas dos condiciones se verifican cuando
$B$ es noetheriano y su topología es preádica ($0_{\mathrm I}$, 7.3.2), o si la
topología de $B$ es la topología $\mathfrak J$-preádica). Sea
$(\mathfrak b_\lambda)$ un sistema fundamental de ideales abiertos de $B$ y
pongamos $B_\lambda=B/\mathfrak b_\lambda$ para todo $\lambda$. Entonces:
\begin{enumerate}
  \item[\rm(i)] Las condiciones siguientes son equivalentes:
  \begin{enumerate}
    \item[a)] $C$ es una $A$-álgebra formalmente lisa.
    \item[b)] Para todo $\lambda$, el homomorfismo de $B_\lambda$-módulos
    \begin{equation}
    \label{0.22.6.2.1-es}
      (\mathfrak J/\mathfrak J^2)\otimes_BB_\lambda
      \longrightarrow\Omega^1_{B/A}\otimes_BB_\lambda
      \tag{22.6.2.1}
    \end{equation}
    deducido de $\delta_{C/B/A}$ por tensorización, es invertible por la izquierda
    (en otras palabras, un isomorfismo sobre un factor directo de
    $\Omega^1_{B/A}\otimes_BB_\lambda$).
  \end{enumerate}
  \item[\rm(ii)] Supongamos además que el anillo topológico $C$ es
  preadmisible ($0_{\mathrm I}$, 7.1.2) y

  \oldpage[0\textsubscript{IV-1}]{206}

  sea $\mathfrak L$ un ideal de definición de $C$; entonces las condiciones a)
  y b) de (i) también equivalen a:
  \begin{enumerate}
    \item[c)] El homomorfismo de $(C/\mathfrak L)$-módulos
    \[
      (\mathfrak J/\mathfrak J^2)\otimes_C(C/\mathfrak L)
      \longrightarrow\Omega^1_{B/A}\otimes_B(C/\mathfrak L)
    \]
    es invertible por la izquierda.
  \end{enumerate}
\end{enumerate}
\end{corollary}

\begin{proof}
Las hipótesis implican que sobre $\Omega^1_{B/A}$ la topología es deducida de
la de $B$ (20.4.5), por lo que (i) resulta de (22.6.1) y de (19.1.7). Por otra
parte, recordemos (20.4.9) que, puesto que $B$ es una $A$-álgebra formalmente
lisa, $\Omega^1_{B/A}$ es un $B$-módulo formalmente proyectivo, luego
$\Omega^1_{B/A}\otimes_BB_\lambda$ es un $B_\lambda$-módulo proyectivo para
todo $\lambda$ (19.2.4); la equivalencia de c) con a) y b) cuando $C$ es
preadmisible resulta entonces de (19.1.9).
\end{proof}

En particular:

\begin{corollary}[22.6.3]
\label{0.22.6.3-es}
Sean $A$ un anillo, $B$ una $A$-álgebra, $\mathfrak J$ un ideal de $B$,
$C=B/\mathfrak J$; se supone que $A$, $B$, $C$ están provistos de las topologías discretas
y que $B$ es una $A$-álgebra formalmente lisa. Para que $C$ sea una
$A$-álgebra formalmente lisa, es necesario y suficiente que el homomorfismo
canónico (22.6.1.1) sea invertible por la izquierda.
\end{corollary}

Se observará que, puesto que toda $A$-álgebra $C$ es cociente de un álgebra de
polinomios $A[T_\alpha]_{\alpha\in I}$, y que esta última es formalmente
lisa (19.3.3), el criterio (22.6.3) permite en principio reconocer si $C$
es una $A$-álgebra formalmente lisa.

\begin{proposition}[22.6.4]
\label{0.22.6.4-es}
Sean $A$ un anillo, $B$ una $A$-álgebra formalmente lisa (para las
topologías discretas), $\mathfrak J$ un ideal de $B$, $C=B/\mathfrak J$; se
supone que $\mathfrak J/\mathfrak J^2$ es un $C$-módulo de tipo finito. Sean
$\mathfrak p$ un ideal primo de $C$, $k(\mathfrak p)$ el cuerpo residual de
$C_{\mathfrak p}$, $\mathfrak p'$ la imagen inversa de $\mathfrak p$ en
$B$, $\mathfrak q$ el ideal primo de $A$ imagen inversa de $\mathfrak p'$.
Las condiciones siguientes son equivalentes:
\begin{enumerate}
  \item[a)] $C_{\mathfrak p}$ es una $A_{\mathfrak q}$-álgebra formalmente
  lisa (para las topologías discretas).
  \item[a')] $C_{\mathfrak p}$ es una $A$-álgebra formalmente lisa (para
  las topologías discretas).
  \item[b)] El homomorfismo canónico
  \begin{equation}
  \label{0.22.6.4.1-es}
    (\mathfrak J/\mathfrak J^2)\otimes_Ck(\mathfrak p)
    \longrightarrow\Omega^1_{B/A}\otimes_Bk(\mathfrak p)
    \tag{22.6.4.1}
  \end{equation}
  es inyectivo.
  \item[c)] Existe $f\in C-\mathfrak p$ tal que $C_f$ es una
  $A$-álgebra formalmente lisa.
\end{enumerate}
Si además $B$ es noetheriano, las condiciones anteriores también son
equivalentes a:
\begin{enumerate}
  \item[d)] $C_{\mathfrak p}$ es una $A_{\mathfrak q}$-álgebra formalmente
  lisa para la topología $\mathfrak p$-preádica sobre $C_{\mathfrak p}$ y la
  topología discreta (o la topología $\mathfrak q$-preádica) sobre
  $A_{\mathfrak q}$.
\end{enumerate}
\end{proposition}

\begin{proof}
Como $A_{\mathfrak q}$ es una $A$-álgebra formalmente lisa (19.3.5, (iv)
y (ii)), ya se sabe que a) y a') son equivalentes (19.3.5, (ii)). Se tiene
$C_{\mathfrak p}=B_{\mathfrak p'}/\mathfrak J B_{\mathfrak p'}$, y
$B_{\mathfrak p'}$ es una $A$-álgebra formalmente lisa para la topología
discreta (19.3.5, (iv)); además,
$\Omega^1_{B_{\mathfrak p'}/A}=(\Omega^1_{B/A})_{\mathfrak p'}$ (20.5.9).
La equivalencia de a') y b) resulta entonces de aplicar (22.6.3) a
$B_{\mathfrak p'}$ y $C_{\mathfrak p}$, teniendo en cuenta que
$\Omega^1_{B/A}$ es un $B$-módulo proyectivo (20.4.9) y (19.2.11), y que
$\mathfrak J/\mathfrak J^2$ es un $B$-módulo de tipo finito, y utilizando
(19.1.12). La aplicación de (19.1.12) prueba también la equivalencia de b) y c).
Por último, observemos que, puesto que $B_{\mathfrak p'}$ es una
$A_{\mathfrak q}$-álgebra formalmente lisa para las topologías

\oldpage[0\textsubscript{IV-1}]{207}

discretas (19.3.5, (iv)), $B_{\mathfrak p'}$, provisto de la topología
$\mathfrak p'$-preádica, sigue siendo una $A_{\mathfrak q}$-álgebra
formalmente lisa cuando se toma sobre $A_{\mathfrak q}$ la topología
discreta o la topología $\mathfrak q$-preádica (19.3.8). Para demostrar
la equivalencia de b) y d) cuando $B$ es noetheriano, apliquemos (22.6.2) al
anillo $A_{\mathfrak q}$ discreto (o $\mathfrak q$-preádico) y al anillo
$B_{\mathfrak p'}$, provisto de la topología $\mathfrak p'$-preádica; como
$C_{\mathfrak p}$ es preadmisible y $\mathfrak pC_{\mathfrak p}$ es
un ideal de definición para su topología, se puede invocar la equivalencia de
c) y a) en (22.6.2).
\end{proof}

\begin{corollary}[22.6.5]
\label{0.22.6.5-es}
Bajo las hipótesis generales de (22.6.4), el conjunto de los
$\mathfrak p\in\operatorname{Spec}(C)$ tales que $C_{\mathfrak p}$ sea una $A$-álgebra
formalmente lisa (o una $A_{\mathfrak q}$-álgebra formalmente lisa, designando
por $\mathfrak q$ la imagen inversa de $\mathfrak p$ en $A$)
para las topologías discretas, es abierto en $\operatorname{Spec}(C)$.
\end{corollary}

Esto resulta de la forma c) de (22.6.4). Se observará que, cuando $B$ es
noetheriano, se pueden sustituir las topologías discretas por las topologías
preádicas.

\begin{corollary}[22.6.6]
\label{0.22.6.6-es}
Bajo las hipótesis generales de (22.6.4), las condiciones siguientes son
equivalentes:
\begin{enumerate}
  \item[a)] $C$ es una $A$-álgebra formalmente lisa (para las topologías
  discretas).
  \item[b)] Para todo $\mathfrak p\in\operatorname{Spec}(C)$ (o solamente para todo ideal
  máximo $\mathfrak p$ de $C$), $C_{\mathfrak p}$ es una $A$-álgebra
  formalmente lisa (o una $A_{\mathfrak q}$-álgebra formalmente lisa, designando
  por $\mathfrak q$ la imagen inversa de $\mathfrak p$ en $A$)
  para las topologías discretas.
\end{enumerate}
Además, cuando $B$ es noetheriano, se pueden sustituir en b) las topologías
discretas por las topologías preádicas sobre $A_{\mathfrak q}$ y
$C_{\mathfrak p}$.
\end{corollary}

\begin{proof}
La última afirmación resulta de (22.6.4). En virtud de (22.6.3), la condición
a) equivale a que el homomorfismo (22.6.1.1) sea invertible por la izquierda;
la equivalencia de a) y b) resulta, pues, de (19.1.14), teniendo en cuenta que
$\Omega^1_{B/A}$ es un $B$-módulo proyectivo y $\mathfrak J/\mathfrak J^2$
un $C$-módulo de tipo finito.
\end{proof}

\begin{proposition}[22.6.7]
\label{0.22.6.7-es}
Sean $k_0$ un cuerpo, $k$ una extensión separable de $k_0$, $C$ una
$k$-álgebra de tipo finito.
\begin{enumerate}
  \item[\rm(i)] Sea $\mathfrak p$ un ideal primo de $C$. Las condiciones
  siguientes son equivalentes:
  \begin{enumerate}
    \item[a)] $C_{\mathfrak p}$ es una $k_0$-álgebra formalmente lisa para
    la topología discreta.
    \item[b)] $C_{\mathfrak p}$ es una $k_0$-álgebra formalmente lisa para
    la topología $\mathfrak p$-preádica.
    \item[c)] Existe $f\in C-\mathfrak p$ tal que $C_f$ es una
    $k_0$-álgebra formalmente lisa para la topología discreta.
    \item[d)] $C_{\mathfrak p}$ es un anillo geométricamente regular
    (19.6.5) sobre $k_0$.
  \end{enumerate}
  Además, el conjunto de los $\mathfrak p\in\operatorname{Spec}(C)$ que poseen cualquiera de
  estas propiedades es abierto en $\operatorname{Spec}(C)$.
  \item[\rm(ii)] Las condiciones siguientes son equivalentes:
  \begin{enumerate}
    \item[a)] $C$ es una $k_0$-álgebra formalmente lisa para la topología
    discreta.
    \item[b)] Todo $\mathfrak p\in\operatorname{Spec}(C)$ verifica las condiciones
    equivalentes de (i).
    \item[c)] Todo ideal máximo $\mathfrak m$ de $C$ verifica las condiciones
    equivalentes de (i).
    \item[d)] Para toda extensión $k'$ de $k_0$, $C\otimes_{k_0}k'$ es un
    anillo regular (17.3.6); se dice también que $C$ es un anillo
    geométricamente regular sobre $k_0$.
  \end{enumerate}

  \oldpage[0\textsubscript{IV-1}]{208}

  \item[\rm(iii)] Sean $B=k[T_1,\ldots,T_r]$ un álgebra de polinomios,
  $\mathfrak J$ un ideal de $B$ tales que $C$ sea isomorfa a
  $B/\mathfrak J$. Sea $\mathfrak q$ un ideal primo de $B$ que contiene
  $\mathfrak J$ y sea $\mathfrak p=\mathfrak q/\mathfrak J$. Las condiciones
  de (i) son entonces también equivalentes a la siguiente («criterio jacobiano de
  Zariski»):
  \begin{enumerate}
    \item[e)] Existe un sistema de elementos
    $(f_i)_{1\leq i\leq m}$ de $\mathfrak J$, que engendra el ideal
    $\mathfrak J_{\mathfrak q}$ en $B_{\mathfrak q}$, y de
    $k_0$-derivaciones $D_j$ de $B$ en sí mismo $(1\leq j\leq m)$ tales que
    $\det(D_j(f_i))\notin\mathfrak q$.
  \end{enumerate}
\end{enumerate}
\end{proposition}

\begin{proof}
Se sabe que $C$ es siempre isomorfa a una $k$-álgebra de la forma
$B/\mathfrak J$. Como $B$ es una $k$-álgebra formalmente lisa (19.3.3) y
$k$, al ser separable sobre $k_0$, es una $k_0$-álgebra formalmente lisa
(19.6.1), $B$ es una $k_0$-álgebra formalmente lisa (19.3.5, (ii)).
La equivalencia de las condiciones a), b), c) de (i) resulta, pues, de (22.6.4), y
la de a), b), c) en (ii) resulta de (22.6.6); la equivalencia de a) y d)
en (i) resulta de (22.5.8); como todo localizado de $C\otimes_{k_0}k'$ es
también un localizado de $C_{\mathfrak p}\otimes_{k_0}k'$ para un ideal primo
$\mathfrak p\in\operatorname{Spec}(C)$ conveniente, la equivalencia de d) y b) en (ii)
resulta de la equivalencia de d) y a) en (i). Finalmente, como
$\Omega^1_{B/k_0}$ es un $B$-módulo proyectivo, la equivalencia del criterio de
Zariski e) y de las demás condiciones de (i) se deduce de (19.1.12) y de
(22.6.3), teniendo en cuenta (20.4.8).
\end{proof}

Zariski se interesaba, de hecho, por un criterio diferencial de regularidad para los
anillos locales $C_{\mathfrak p}$. Como es equivalente decir que un anillo
local noetheriano que contiene un cuerpo es regular o es formalmente liso en
tanto que álgebra sobre su subcuerpo primo (19.6.4), se obtiene inmediatamente tal
criterio sustituyendo $k_0$ por el subcuerpo primo de $k$ en (22.6.7); en
particular, se obtiene el resultado siguiente, que volveremos a encontrar más adelante
(IV, 6.12.5) como caso particular de resultados más generales de Nagata:

\begin{corollary}[22.6.8]
\label{0.22.6.8-es}
\textnormal{(Zariski).}
Sea $C$ un álgebra de tipo finito sobre un cuerpo. Entonces el conjunto de los
$\mathfrak p\in\operatorname{Spec}(C)$ tales que $C_{\mathfrak p}$ sea un anillo local
regular es abierto en $\operatorname{Spec}(C)$.
\end{corollary}

\begin{remark}[22.6.9]
\label{0.22.6.9-es}
\begin{enumerate}
  \item[\rm(i)] Los enunciados de (22.6.7) siguen siendo válidos si, en lugar de
  suponer $k$ separable sobre $k_0$, se supone solamente que es de
  \emph{multiplicidad radicial finita}, es decir (19.6.6), que existe
  una extensión radicial finita $k'_0$ de $k_0$ tal que el cuerpo residual del
  anillo artiniano local $k\otimes_{k_0}k'_0$ sea una extensión separable
  $k'$ de $k'_0$. Existe entonces un $k'_0$-monomorfismo
  $k'\longrightarrow k\otimes_{k_0}k'_0$ que, compuesto con el homomorfismo
  canónico $k\otimes_{k_0}k'_0\longrightarrow k'$, da la identidad, puesto que
  $k'$ es una $k'_0$-álgebra formalmente lisa (19.6.1); se concluye que
  $C\otimes_{k_0}k'_0$ está provista de una estructura de $k'$-álgebra de tipo finito,
  y, como $k'$ es separable sobre $k'_0$, se puede aplicar (22.6.7) a esta
  $k'$-álgebra; se concluye nuestra afirmación aplicando (19.4.7). Como
  no habremos de utilizar esta generalización, dejamos los detalles al
  lector. Ignoramos, en cambio, si los resultados de (22.6.7) son válidos sin
  \emph{ninguna} hipótesis sobre la extensión $k$ de $k_0$.
  \item[\rm(ii)] Bajo las hipótesis generales de (22.6.7), sea $(k_\alpha)$
  una familia filtrante decreciente de subcuerpos de $k$ que contienen $k_0$ y
  tal que
  $\bigcap_\alpha k_\alpha(k^p)=k_0(k^p)$ (donde $p$ es el exponente
  característico de $k_0$). Utilizando un método de cómputo de dimensiones
  debido a Nagata, se puede demostrar que, en el criterio jacobiano (22.6.7, (iii)),
  se puede uno limitar a derivaciones $D_j$ de $B$ que sean
  $k_\alpha$-derivaciones para un $\alpha$ conveniente.

  \oldpage[0\textsubscript{IV-1}]{209}

  El interés de este resultado es que siempre existen tales familias
  $(k_\alpha)$ para las cuales $[k:k_\alpha]$ es finito (21.8.6). No
  demostraremos este refinamiento del criterio de Zariski, que no habremos de
  utilizar. En (22.7), vamos a dar, para los anillos locales completos,
  una variante (debida también a Nagata) del criterio de Zariski, que se demuestra
  esencialmente por el mismo método (con dificultades técnicas algo
  mayores).
\end{enumerate}
\end{remark}

\subsection{El criterio jacobiano de Nagata}
\label{subsection:0.22.7-es}

\begin{env}[22.7.1]
\label{0.22.7.1-es}
El criterio jacobiano de Nagata es el análogo del criterio jacobiano de Zariski,
pero para anillos cocientes de anillos de \emph{series formales} sobre un
cuerpo. Daremos, como Nagata [31], dos versiones, presentadas aquí
como criterios de lisitud formal.
\end{env}

\begin{proposition}[22.7.2]
\label{0.22.7.2-es}
Sean $k$ un cuerpo, $A=k[[T_1,\ldots,T_r]]$ un anillo de series formales sobre
$k$, provisto de su estructura usual de $k$-álgebra, $\mathfrak q$ un ideal de
$A$, $B=A/\mathfrak q$. Sea $\mathfrak p$ un ideal primo de $A$ que contiene
$\mathfrak q$. Se supone que existe una sub-$k$-álgebra $C'$ de
$C=A/\mathfrak p$, isomorfa a un álgebra de series formales
$k[[X_1,\ldots,X_s]]$, tal que $C$ sea finita sobre $C'$ y que el cuerpo de
fracciones $L$ de $C$ sea una extensión separable de
$L'=k((X_1,\ldots,X_s))$ (hipótesis siempre satisfecha si $k$ es de
característica 0). Entonces las condiciones siguientes son equivalentes:
\begin{enumerate}
  \item[a)] $B_{\mathfrak p}=A_{\mathfrak p}/\mathfrak qA_{\mathfrak p}$ es
  una $k$-álgebra formalmente lisa para la topología
  $\mathfrak p$-preádica.
  \item[b)] Existen $k$-derivaciones $D_i$ $(1\leq i\leq m)$ de $A$ en
  sí mismo, y elementos $f_i$ $(1\leq i\leq m)$ de $\mathfrak q$, tales
  que las imágenes de los $f_i$ en $\mathfrak qA_{\mathfrak p}$ engendran este
  ideal de $A_{\mathfrak p}$, y se tiene
  $\det(D_i(f_j))\notin\mathfrak p$.
  \item[c)] $B_{\mathfrak p}$ es un anillo regular.
\end{enumerate}
\end{proposition}

\begin{proof}
El cuerpo residual de $B_{\mathfrak p}$ es $k(\mathfrak p)$; como
$k((X_1,\ldots,X_s))$ es separable sobre $k$ (21.9.6.4), $k(\mathfrak p)$ es
separable sobre $k$ por hipótesis, y ya se sabe que en estas condiciones las
propiedades a) y c) son equivalentes (19.6.4). Por otra parte,
$A_{\mathfrak p}$ es un anillo regular ((17.1.4) y (17.3.2)), y como su
cuerpo residual $k(\mathfrak p)$ es separable sobre $k$, $A_{\mathfrak p}$ es
formalmente lisa sobre $k$ para la topología $\mathfrak p$-preádica
(19.6.4). Resulta, pues, de (22.6.2, (ii)) que la condición a) equivale a que
el homomorfismo de $L$-espacios vectoriales
\begin{equation}
\label{0.22.7.2.1-es}
  j_0:(\mathfrak q/\mathfrak q^2)\otimes_AL
  \longrightarrow\Omega^1_{A_{\mathfrak p}/k}\otimes_{A_{\mathfrak p}}L
  =\widehat\Omega^1_{A/k}\otimes_AL
  \tag{22.7.2.1}
\end{equation}
sea inyectivo.

Consideremos, por otra parte, el homomorfismo compuesto
\begin{equation}
\label{0.22.7.2.2-es}
  j:(\mathfrak q/\mathfrak q^2)\otimes_AL
  \xrightarrow{j_0}\Omega^1_{A/k}\otimes_AL
  \longrightarrow\widehat\Omega^1_{A/k}\otimes_AL
  \tag{22.7.2.2}
\end{equation}
y demostremos que la condición b) equivale a decir que $j$ es inyectivo. Observemos, en
efecto, que
\[
  (\mathfrak q/\mathfrak q^2)\otimes_AL
  =(\mathfrak qA_{\mathfrak p}/\mathfrak q^2A_{\mathfrak p})
  \otimes_{A_{\mathfrak p}}L~;
\]
la condición b) significa que, si $F_i=j(f'_i\otimes1)$, donde $f'_i$ es la imagen
de $f_i$ en
$\mathfrak qA_{\mathfrak p}/\mathfrak q^2A_{\mathfrak p}$, la matriz
$(\langle F_i,D_j\otimes1\rangle)$ es invertible; por tanto, esto implica que los
$F_i$ son linealmente independientes, y \emph{a fortiori} que ocurre lo
mismo con los $f'_i\otimes1$; pero como estos últimos engendran
$(\mathfrak qA_{\mathfrak p}/\mathfrak q^2A_{\mathfrak p})
\otimes_{A_{\mathfrak p}}L$, se concluye que $j$ es inyectivo. Recíprocamente,
supongamos $j$ inyectivo, y tomemos los $f_i$ de modo que los $f'_i\otimes1$
formen

\oldpage[0\textsubscript{IV-1}]{210}

una base de
$(\mathfrak qA_{\mathfrak p}/\mathfrak q^2A_{\mathfrak p})
\otimes_{A_{\mathfrak p}}L$; entonces los $F_i$ forman una base de la imagen de
$j$; pero se sabe (21.9.3) que $\widehat\Omega^1_{A/k}$ es un $A$-módulo
libre de rango $r$ y que las $k$-derivaciones de $A$ en sí mismo engendran su
dual; el hecho de que los $F_i$ sean linealmente independientes implica, pues,
la existencia de $k$-derivaciones $D_j$ de $A$ en sí mismo tales que la matriz
$(\langle F_i,D_j\otimes1\rangle)$ sea invertible, es decir, la condición
b).

Estas observaciones muestran ya que b) implica a). Para probar recíprocamente que
la condición a) implica que $j$ es inyectivo, consideremos el diagrama
conmutativo
\begin{equation}
\label{0.22.7.2.3-es}
\xymatrix{
  (\mathfrak q/\mathfrak q^2)\otimes_AL
    \ar[r]^j \ar[d]_\alpha &
  \widehat\Omega^1_{A/k}\otimes_AL \ar@{=}[d] \\
  (\mathfrak p/\mathfrak p^2)\otimes_AL
    \ar[r]_i &
  \widehat\Omega^1_{A/k}\otimes_AL
}
\tag{22.7.2.3}
\end{equation}
y observemos el lema siguiente:
\end{proof}

\begin{lemma}[22.7.2.4]
\label{0.22.7.2.4-es}
Sean $R$ un anillo local regular, $\mathfrak m$ su ideal máximo,
$K=R/\mathfrak m$ su cuerpo residual. Si $\mathfrak n$ es un ideal de $R$ tal
que $R/\mathfrak n$ sea regular, entonces el homomorfismo canónico
\begin{equation}
\label{0.22.7.2.5-es}
  \alpha:(\mathfrak n/\mathfrak n^2)\otimes_RK
  \longrightarrow(\mathfrak m/\mathfrak m^2)\otimes_RK
  \tag{22.7.2.5}
\end{equation}
es inyectivo.
\end{lemma}

\begin{proof}
En efecto, se sabe (17.1.6) que $\mathfrak n$ está engendrado por una sucesión
$(x_i)_{1\leq i\leq t}$ que forma parte de un sistema regular de parámetros
$(x_i)_{1\leq i\leq n}$ de $R$; las imágenes de los $x_i$ $(1\leq i\leq t)$
forman una base de $(\mathfrak n/\mathfrak n^2)\otimes_RK$, y sus imágenes
por $\alpha$ forman parte de la base de
$(\mathfrak m/\mathfrak m^2)\otimes_RK$ constituida por las imágenes de los $x_i$ para
$1\leq i\leq n$, de donde resulta el lema.
\end{proof}

\begin{proof}[Continuación de la demostración de la proposición~{\rm(22.7.2)}]
Este lema y el diagrama (22.7.2.3) reducen, pues, (en virtud de la hipótesis a))
la cuestión a probar que $i$ es inyectivo. Ahora bien, en la sucesión
\begin{equation}
\label{0.22.7.2.6-es}
  \mathfrak p/\mathfrak p^2\xrightarrow{h}
  \widehat\Omega^1_{A/k}\widehat\otimes_A(A/\mathfrak p)
  \xrightarrow{g}\widehat\Omega^1_{(A/\mathfrak p)/k}\longrightarrow0
  \tag{22.7.2.6}
\end{equation}
se sabe (20.7.20) que $g$ es sobreyectivo y que la imagen de $h$ es densa en
el núcleo de $g$ para la topología $\mathfrak m$-ádica ($\mathfrak m$ es
el ideal máximo de $A$); como $\widehat\Omega^1_{A/k}$ es un $A$-módulo de
tipo finito, todos los submódulos del $(A/\mathfrak p)$-módulo
\[
  \widehat\Omega^1_{A/k}\otimes_A(A/\mathfrak p)
  =\widehat\Omega^1_{A/k}\widehat\otimes_A(A/\mathfrak p)
\]
son cerrados para la topología $\mathfrak m$-ádica ($0_{\mathrm I}$, 7.3.5),
por lo que la sucesión (22.7.2.6) es exacta. Tensorizando por $L$, se obtiene la sucesión
exacta
\[
  (\mathfrak p/\mathfrak p^2)\otimes_AL\xrightarrow{i}
  \widehat\Omega^1_{A/k}\otimes_AL\longrightarrow
  \widehat\Omega^1_{(A/\mathfrak p)/k}\otimes_{A/\mathfrak p}L
  \longrightarrow0.
\]

Ahora bien, la hipótesis de separabilidad hecha sobre $L$ implica, en virtud de (21.9.5),
que $\widehat\Omega^1_{(A/\mathfrak p)/k}\otimes_{A/\mathfrak p}L$ tiene
rango $s$ sobre $L$; como $\widehat\Omega^1_{A/k}\otimes_AL$ tiene rango $r$
sobre $L$ (21.9.3), se tiene $\operatorname{rg}_L(\operatorname{Im}i)=r-s$. Pero como
$A/\mathfrak p$ es finita sobre una subálgebra isomorfa a
$k[[T_1,\ldots,T_s]]$, se tiene $\dim(A/\mathfrak p)=s$ ((17.1.4) y (16.1.5));
por tanto,
\[
  r-s=\dim(A)-\dim(A/\mathfrak p)=\dim(A_{\mathfrak p})
\]

\oldpage[0\textsubscript{IV-1}]{211}

por (16.5.11). Ahora bien, como $A_{\mathfrak p}$ es regular (17.3.2),
$\dim(A_{\mathfrak p})$ es igual al rango sobre $L$ de
$\mathfrak pA_{\mathfrak p}/\mathfrak p^2A_{\mathfrak p}$ (17.1.1), y por tanto al
rango sobre $L$ de $(\mathfrak p/\mathfrak p^2)\otimes_AL$, lo que acaba de
probar que $i$ es inyectivo.
\end{proof}

\begin{theorem}[22.7.3]
\label{0.22.7.3-es}
Sean $k$ un cuerpo, $A$ un anillo local noetheriano completo de cuerpo residual
$K$. Se supone que:
\begin{enumerate}
  \item[$1^\circ$] $[k:k^p]<+\infty$ (donde $p$ es el exponente característico
  de $k$);
  \item[$2^\circ$] $K$ es una extensión finita de una extensión separable
  $K_0$ de $k$;
  \item[$3^\circ$] $A$ está provisto de una estructura de $K_0$-álgebra formalmente
  lisa (para la topología preádica).
\end{enumerate}
Sean $\mathfrak q$ un ideal de $A$, $B=A/\mathfrak q$, $\mathfrak p$ un
ideal primo de $A$ que contiene $\mathfrak q$. Las condiciones siguientes son
equivalentes:
\begin{enumerate}
  \item[a)] $B_{\mathfrak p}$ es una $k$-álgebra formalmente lisa (para la
  topología $\mathfrak p$-preádica).
  \item[b)] Existen $k$-derivaciones $D_i$ de $A$ en sí mismo
  $(1\leq i\leq m)$ y elementos $f_i$ $(1\leq i\leq m)$ de
  $\mathfrak q$, tales que las imágenes de los $f_i$ en
  $\mathfrak qA_{\mathfrak p}$ engendran este ideal de $A_{\mathfrak p}$, y
  se tiene $\det(D_i(f_j))\notin\mathfrak p$.
  \item[b')] Existe una subextensión $k'$ de $K_0$, que contiene $K_0^p$,
  tal que $[K_0:k']<+\infty$, $k'$-derivaciones $D_i$ de $A$ en
  sí mismo $(1\leq i\leq m)$ y elementos $f_i$ de $\mathfrak q$
  $(1\leq i\leq m)$, tales que las imágenes de los $f_i$ en
  $\mathfrak qA_{\mathfrak p}$ engendran este ideal de $A_{\mathfrak p}$, y
  se tiene $\det(D_i(f_j))\notin\mathfrak p$.
\end{enumerate}
\end{theorem}

\begin{proof}
Distingamos dos casos, según que $p=1$ o $p>1$.

\emph{A)} $k$ \emph{es de característica} $0$. Como $K$ es entonces una
extensión separable de $K_0$, resulta de (19.6.4) que $A$ es
$K_0$-isomorfo a un anillo de series formales $K[[T_1,\ldots,T_r]]$ provisto de
su estructura usual de $K$-álgebra. Pero entonces, teniendo en cuenta (19.8.8, (ii))
aplicado a $A/\mathfrak p$, se ve que las condiciones generales de (22.7.2)
se verifican sustituyendo allí $k$ por $K$. Además, en virtud de (19.6.4), es
equivalente decir que $B_{\mathfrak p}$ es una $k$-álgebra
formalmente lisa o una $K$-álgebra formalmente lisa, pues ambas condiciones
equivalen a que $B_{\mathfrak p}$ sea un anillo regular. Se pueden
aplicar, pues, las conclusiones de (22.7.2), y es inmediato que esto
prueba la equivalencia de a), b) y b') (con $k'=K_0$ en b')).

\emph{B)} $k$ \emph{es de característica} $p>0$. Como $A$ es una
$K_0$-álgebra formalmente lisa y $K_0$ es separable sobre $k$,
resulta de (19.3.5, (ii)) y (19.6.1) que $A$ es una $k$-álgebra
formalmente lisa; en virtud de (22.5.9), $A_{\mathfrak p}$ es también una
$k$-álgebra formalmente lisa para la topología $\mathfrak p$-preádica. Resulta,
pues, de nuevo de (22.6.2, (ii)) que la condición a) equivale a que
el homomorfismo $j_0$ definido en (22.7.2.1) (donde $L=k(\mathfrak p)$) sea
inyectivo. Por tanto, (19.1.12, c)) muestra ya que b) implica a); como, por
otra parte, b') implica trivialmente b), todo se reduce a demostrar que la hipótesis
de que $j_0$ es inyectivo implica b').

Designemos primero por $k'$ una subextensión cualquiera de $K_0$ tal que
$[K_0:k']<+\infty$. Observemos que $\widehat\Omega^1_{A/k'}$ es un $A$-módulo
libre de tipo finito. Por el razonamiento de (22.7.2), basta demostrar que
(para una elección conveniente de $k'$), el homomorfismo compuesto
\begin{equation}
\label{0.22.7.3.1-es}
  j':(\mathfrak q/\mathfrak q^2)\otimes_AL
  \xrightarrow{j_0}\Omega^1_{A/k}\otimes_AL
  \longrightarrow\widehat\Omega^1_{A/k'}\otimes_AL
  \tag{22.7.3.1}
\end{equation}

\oldpage[0\textsubscript{IV-1}]{212}

es inyectivo. Consideremos de nuevo el diagrama conmutativo
\begin{equation}
\label{0.22.7.3.2-es}
\xymatrix{
  (\mathfrak q/\mathfrak q^2)\otimes_AL
    \ar[r]^{j_0} \ar[d]_\alpha &
  \Omega^1_{A/k}\otimes_AL \ar[r] \ar@{=}[d] &
  \widehat\Omega^1_{A/k'}\otimes_AL \ar@{=}[d] \\
  (\mathfrak p/\mathfrak p^2)\otimes_AL
    \ar[r]_{i_0} &
  \Omega^1_{A/k}\otimes_AL \ar[r] &
  \widehat\Omega^1_{A/k'}\otimes_AL
}
\tag{22.7.3.2}
\end{equation}

Como $j_0$ es inyectivo, se tiene $\operatorname{Ker}(i_0\circ\alpha)=0$. Pero como
$A_{\mathfrak p}$ es un anillo regular y la hipótesis a) implica que
$B_{\mathfrak p}=A_{\mathfrak p}/\mathfrak qA_{\mathfrak p}$ es también un
anillo regular (19.6.5), el lema (22.7.2.4) muestra que $\alpha$ es inyectivo,
de donde $\alpha^{-1}(\operatorname{Ker} i_0)=0$. Si $i'$ es el compuesto de los homomorfismos de la
segunda línea de (22.7.3.2), se tiene asimismo
$\operatorname{Ker}(j')=\alpha^{-1}(\operatorname{Ker}(i'))$, y todo se reduce, pues, a demostrar que, para una
elección conveniente de $k'$, se tiene
\begin{equation}
\label{0.22.7.3.3-es}
  \operatorname{Ker}(i')=\operatorname{Ker}(i).
  \tag{22.7.3.3}
\end{equation}

Pero el mismo razonamiento que en (22.7.2) muestra que se tiene una sucesión exacta
\[
  (\mathfrak p/\mathfrak p^2)\otimes_AL\xrightarrow{i'}
  \widehat\Omega^1_{A/k'}\otimes_AL\longrightarrow
  \widehat\Omega^1_{(A/\mathfrak p)/k'}\otimes_{A/\mathfrak p}L
  \longrightarrow0.
\]
Ahora bien, las hipótesis hechas permiten aplicar (21.9.8), por lo que existe una
subextensión $k'$ de $K_0$ que contiene $K_0^p$, tal que
$[K_0:k']<+\infty$, y para la cual se tiene
\[
  \operatorname{rg}_L\bigl(\widehat\Omega^1_{(A/\mathfrak p)/k'}
  \otimes_{A/\mathfrak p}L\bigr)
  =\dim(A/\mathfrak p)+\operatorname{rg}_L\Upsilon_{L/k}
  +\operatorname{rg}_{K_0}\Omega^1_{K_0/k'}.
\]
Pero se tiene
\[
  \dim(A/\mathfrak p)=\dim(A)-\dim(A_{\mathfrak p})
\]
por (16.5.11),
\[
  \dim(A_{\mathfrak p})
  =\operatorname{rg}_L\bigl((\mathfrak p/\mathfrak p^2)\otimes_AL\bigr)
\]
puesto que $A_{\mathfrak p}$ es regular (17.1.1), y finalmente
\[
  \dim(A)+\operatorname{rg}_{K_0}\Omega^1_{K_0/k'}
  =\operatorname{rg}_L\bigl(\widehat\Omega^1_{A/k'}\otimes_AL\bigr)
\]
en virtud de (21.9.2); se obtiene, pues,
\begin{equation}
\label{0.22.7.3.4-es}
  \operatorname{rg}_L(\operatorname{Ker}(i'))=\operatorname{rg}_L\Upsilon_{L/k}.
  \tag{22.7.3.4}
\end{equation}

Por otra parte, siendo el cuerpo residual de $A_{\mathfrak p}$ formalmente liso
sobre el cuerpo primo de $k$, se tiene una sucesión exacta (20.6.22.1)
\[
  \Upsilon_{L/k}\xrightarrow{\chi_{A_{\mathfrak p}/k}}
  (\mathfrak pA_{\mathfrak p})/(\mathfrak pA_{\mathfrak p})^2
  \xrightarrow{i}\Omega^1_{A_{\mathfrak p}/k}
  \otimes_{A_{\mathfrak p}}L,
\]
y como $A_{\mathfrak p}$ es una $k$-álgebra formalmente lisa para la
topología $\mathfrak p$-preádica, resulta de (22.2.2) que
$\chi_{A_{\mathfrak p}/k}$ es un homomorfismo inyectivo, por lo que $\operatorname{Ker}(i)$ es
isomorfo a $\Upsilon_{L/k}$ y, teniendo en cuenta (22.7.3.4) y que
$\operatorname{Ker}(i)\subset\operatorname{Ker}(i')$, esto acaba de probar (22.7.3.3) y, por consiguiente, el
teorema.
\end{proof}

\oldpage[0\textsubscript{IV-1}]{213}

\begin{remark}[22.7.4]
\label{0.22.7.4-es}
De hecho se ha demostrado (utilizando (21.9.6)) que la condición b') sigue siendo
equivalente a las demás condiciones de (22.7.3) cuando se obliga a $k'$ a
ser uno de los cuerpos de una familia filtrante decreciente $(k_\alpha)$ de
subcuerpos de $K_0$ que contienen $K_0^p$, tales que
$[K_0:k_\alpha]<+\infty$ para todo $\alpha$ y que
$\bigcap_\alpha k_\alpha=K_0^p$.
\end{remark}

\begin{corollary}[22.7.5]
\label{0.22.7.5-es}
Bajo las hipótesis generales de (22.7.3), el conjunto de los ideales primos
$\mathfrak n\in\operatorname{Spec}(B)$ tales que $B_{\mathfrak n}$ sea una
$k$-álgebra formalmente lisa (para la topología $\mathfrak n$-preádica)
es abierto en $\operatorname{Spec}(B)$.
\end{corollary}

\begin{proof}
Con las notaciones de (22.7.3), basta ver que, si $B_{\mathfrak p}$
es formalmente lisa sobre $k$, ocurre lo mismo con $B_{\mathfrak p'}$, para
todo $\mathfrak p'\in V(\mathfrak q)$ suficientemente próximo a $\mathfrak p$ en
$\operatorname{Spec}(A)$. Ahora bien, si las $k$-derivaciones $D_i$ y los elementos
$f_i$ de $\mathfrak q$ verifican el criterio b) de (22.7.3), se tiene también
$f=\det(D_if_j)\notin\mathfrak p'$ para $\mathfrak p'\in D(f)$; por otra
parte, existe un $g\notin\mathfrak p$ tal que las imágenes de los $f_i$ en
$\mathfrak qA_{\mathfrak p'}$ engendran $\mathfrak qA_{\mathfrak p'}$ para
$\mathfrak p'\in D(g)$ (Bourbaki, \emph{Alg. comm.}, chap.~II, §~5,
n\textsuperscript{o}~1, prop.~2), lo que acaba de probar el corolario.
\end{proof}

\begin{corollary}[22.7.6]
\label{0.22.7.6-es}
\textnormal{(Nagata).} Sea $B$ un anillo local noetheriano completo que contiene
un cuerpo. Entonces el conjunto de los $\mathfrak n\in\operatorname{Spec}(B)$ tales
que $B_{\mathfrak n}$ sea regular es abierto en $\operatorname{Spec}(B)$.
\end{corollary}

\begin{proof}
En efecto, $B$ es un álgebra sobre un cuerpo primo $k$, por tanto (19.8.8, (ii))
de la forma $A/\mathfrak q$, donde
$A=K_0[[T_1,\ldots,T_n]]$ es un anillo de series formales y $K_0$ una
extensión de $k$; por otra parte, decir que $B_{\mathfrak n}$ es regular
equivale a decir que es una $k$-álgebra formalmente lisa para la topología
$\mathfrak n$-preádica (19.6.4). Se verifican, pues, todas las condiciones de aplicación de
(22.7.5).
\end{proof}

\begin{remarks}[22.7.7]
\label{0.22.7.7-es}
\begin{enumerate}
  \item[(i)] Veremos más adelante, con Nagata, utilizando (22.7.6), que
  se puede extender la conclusión de (22.7.6) al caso de un anillo local
  noetheriano completo cualquiera (IV, 6.12.7).

  \item[(ii)] La conclusión del teorema (22.7.3) ya no es necesariamente
  exacta cuando no se supone que $[k:k^p]$ sea finito. Tomemos, por
  ejemplo, $k=\mathbf F_p(X_n)_{n\geqslant0}$,
  $A=k[[T,U]]$, $\mathfrak q$ igual al ideal principal $Af$, donde
  \[
    f=U^p-\sum_{n=0}^{\infty}X_nT^{np}~;
  \]
  $A$ es un anillo factorial en el cual $f$ es un elemento irreducible, pues es
  irreducible en el anillo de polinomios $k((T))[U]$, y por tanto también en
  el anillo de series formales $k((T))[[U]]$ (Bourbaki,
  \emph{Alg. comm.}, chap.~VII, §~3), y \emph{a fortiori} en $A$.
  Tomemos $\mathfrak p=\mathfrak q$, de modo que $B_{\mathfrak p}$ es el
  cuerpo designado por $E$ en (21.9.7), que es separable sobre $k$
  (\emph{loc. cit.}), luego una $k$-álgebra formalmente lisa. Pero se tiene aquí,
  con las notaciones de (22.7.3), $k=K_0=K$, y las $k$-derivaciones de $A$
  en sí mismo son las combinaciones lineales de $\partial/\partial T$ y
  $\partial/\partial U$; como se tiene
  $\partial f/\partial T=\partial f/\partial U=0$, los criterios b) y b') de
  (22.7.3) no se verifican.
\end{enumerate}
\end{remarks}

\section{Anillos japoneses}
\label{section:0.23-es}

Los resultados de esta sección se completarán en IV, (7.6) y (7.7).

\subsection{Anillos japoneses}
\label{subsection:0.23.1-es}

\begin{definition}[23.1.1]
\label{0.23.1.1-es}
Se dice que un anillo íntegro $A$ es un anillo japonés si, para toda
extensión finita $K'$ de su cuerpo de fracciones $K$, la clausura integral
$A'$ de $A$ en $K'$ es

\oldpage[0\textsubscript{IV-1}]{214}

un $A$-módulo de tipo finito (en otras palabras, una $A$-álgebra finita). Se dice
que un anillo $A$ es universalmente japonés si toda $A$-álgebra íntegra de
tipo finito es un anillo japonés.
\end{definition}

Está claro que si $A$ es un anillo japonés, todo anillo de fracciones
$S^{-1}A$ es entonces un anillo japonés, pues (con las notaciones precedentes)
$S^{-1}A'$ es la clausura integral de $S^{-1}A$ en $K'$, y es
evidentemente un $S^{-1}A$-módulo de tipo finito.

Un anillo íntegro noetheriano, incluso un \emph{anillo de valoración discreta},
no es necesariamente un anillo japonés [30].

\begin{proposition}[23.1.2]
\label{0.23.1.2-es}
Sean $A$ un anillo noetheriano íntegro y $K$ su cuerpo de fracciones. Si, para
toda extensión radicial finita $K'$ de $K$, la clausura integral de $A$
en $K'$ es un $A$-módulo de tipo finito, entonces $A$ es un anillo japonés.
\end{proposition}

\begin{proof}
Como $A$ es noetheriano, para comprobar que $A$ es un anillo japonés basta
ver que, para toda extensión cuasi-Galois finita $L$ de $K$, la clausura
integral $B$ de $A$ en $L$ es un $A$-módulo de tipo
finito. Ahora bien, $L$ es
una extensión galoisiana de la mayor extensión radicial $K'$ de $K$ contenida
en $L$; y si $A'$ es la clausura integral de $A$ en $K'$, $B$ es la clausura
integral de $A'$ en $L$. Pero se sabe que en una extensión separable de $K'$,
la clausura integral de $A'$ es un $A'$-módulo de tipo finito (Bourbaki,
\emph{Alg. comm.}, cap.~V, §~1, n\textsuperscript{o}~6, cor.~1 de la
prop.~20), de donde se sigue la proposición.
\end{proof}

De (23.1.2) se deduce que, cuando $K$ es de característica $0$, decir que
$A$ es un anillo japonés significa que su \emph{clausura integral} $A'$ es
un $A$-módulo de tipo finito.

\begin{theorem}[23.1.3]
\label{0.23.1.3-es}
\textnormal{(Tate).} Sean $A$ un anillo íntegro noetheriano y $x\neq0$ un
elemento de $A$. Supongamos que se cumplen las condiciones siguientes~:
\begin{enumerate}
  \item[(i)] $A$ es íntegramente cerrado.
  \item[(ii)] $\mathfrak p=xA$ es primo y $A$ es separado y completo para
  la topología $\mathfrak p$-preádica.
  \item[(iii)] $A/xA$ es un anillo japonés.
\end{enumerate}
Entonces $A$ es un anillo japonés.
\end{theorem}

Utilizaremos el lema siguiente~:

\begin{lemma}[23.1.3.1]
\label{0.23.1.3.1-es}
Sean $A$ un anillo y $x$ un elemento de $A$ no divisor de $0$ en $A$, tal
que $xA=\mathfrak p$ sea primo; entonces, para todo entero $n>0$, la imagen
inversa de $x^nA_{\mathfrak p}$ en $A$ es $x^nA$.
\end{lemma}

\begin{proof}
En efecto, supongamos que $b\in A$ es un elemento tal que
$b/1=x^na/s$ en $A_{\mathfrak p}$, donde $a\in A$ y
$s\notin\mathfrak p$; existe entonces $s'\notin\mathfrak p$ tal que
$s'sb=x^nas'$, de donde $s'sb\in\mathfrak p$, y como
$s's\notin\mathfrak p$, esto implica $b\in\mathfrak p$, es decir,
$b=xb'$ con $b'\in A$; puesto que $x$ no es divisor de $0$, se concluye
$s'sb'=x^{n-1}as'$\footnote{Lectura diplomática: el factor final impreso es s, no s prima; la cancelación exige conservar s prima.}, y basta razonar por inducción sobre $n$.
\end{proof}

\begin{proof}
Para demostrar el teorema, podemos limitarnos al caso en que el cuerpo de fracciones
$K$ de $A$ es de característica $p>0$, puesto que $A$ es íntegramente cerrado.
Sea $K'$ una extensión radicial finita de $K$, de modo que existe una
potencia $q=p^e$ tal que $K'^q\subset K$; sustituyendo $K'$ por una
extensión radicial mayor, podemos incluso suponer que existe $y\in K'$ tal
que $y^q=x$. Sea $A'$ la clausura integral de $A$ en $K'$; como
$A$ es íntegramente cerrado, $A'$ es el conjunto de los $x'\in K'$ tales que
$x'^q\in A'\cap K=A$. Pongamos $V=A_{\mathfrak p}$; como el ideal maximal
$\mathfrak m=xV$ de $V$ es principal, se sabe que el anillo local
noetheriano $V$ es un \emph{anillo de valoración discreta} (17.1.4); como
$V$ es íntegramente cerrado, el mismo

\oldpage[0\textsubscript{IV-1}]{215}

razonamiento anterior muestra que la clausura integral $V'$ de $V$ en
$K'$ es el conjunto de los $x'\in K'$ tales que $x'^q\in V$; se sabe
(Bourbaki, \emph{Alg. comm.}, cap.~VI, §~8, n\textsuperscript{o}~6,
prop.~6, y cap.~V, §~2, n\textsuperscript{o}~3, lema~4) que $V'$ es un
\emph{anillo de valoración discreta}, cuyo ideal maximal $\mathfrak m'$ es
el conjunto de los $x'\in K'$ tales que $x'^q\in\mathfrak m$ y, además, el cuerpo
residual $V'/\mathfrak m'$ es una extensión \emph{finita} de
$V/\mathfrak m$ (\emph{loc. cit.}, cap.~VI, §~8,
n\textsuperscript{o}~1, lema~2). Demostremos que, para todo entero $n>0$, se tiene
\begin{equation}
\label{0.23.1.3.2-es}
  \mathfrak m'^n\cap A'=y^nA'.
  \tag{23.1.3.2}
\end{equation}

En efecto, como $y^q=x$, se tiene $y\in\mathfrak m'$, por tanto
$y^nA'\subset\mathfrak m'^n\cap A'$. Recíprocamente, sea
$x'\in\mathfrak m'^n\cap A'$, y pongamos $x'=z'y^n$ con $z'\in K'$. Se tiene
$x'^q\in A'$; por otra parte, $x'$ es suma de productos
$t'_1\cdots t'_n$ con $t'_i\in\mathfrak m'$, por tanto
$t_i'^q\in\mathfrak m$, y de ello se concluye que
$x'^q\in\mathfrak m^n\cap A$. Ahora bien, el lema (23.1.3.1) prueba que
$\mathfrak m^n\cap A=x^nA$, por lo que se tiene $z'^qx^n\in x^nA$, o también
$z'^q\in A$, puesto que $x\neq0$, lo que establece (23.1.3.2).

Demostremos, en segundo lugar, que sobre $A'$ la topología $xA'$-preádica es
\emph{separada}; esta topología es también la topología $yA'$-preádica de
$A'$, y de (23.1.3.2) se deduce que esta topología es inducida por la
topología $\mathfrak m'$-preádica de $V'$, que es separada porque $V'$ es
un anillo de valoración discreta.

Probemos ahora que $A'/xA'$ es un $A$-módulo de tipo finito. Como $x=y^q$
y $y^kA'/y^{k+1}A'$ es isomorfo a $A'/yA'$, podemos limitarnos a
demostrar que $A'/yA'$ es un $A$-módulo de tipo finito; pero la fórmula
(23.1.3.2), aplicada para $n=1$, muestra que $A'/yA'$ es un subanillo de
$V'/\mathfrak m'$, es decir, del cuerpo residual de $V'$, que es una
extensión finita del cuerpo residual $V/\mathfrak m$ de $V$. Ahora bien,
$V/\mathfrak m$ es el cuerpo de fracciones de $A/\mathfrak p$, y como
$A'$ es entero sobre $A$, $A'/yA'$ es entero sobre $A/\mathfrak p$, por tanto
está contenido en la clausura integral de $A/\mathfrak p$ en
$V'/\mathfrak m'$; como por hipótesis $A/\mathfrak p$ es un anillo
japonés noetheriano, $A'/yA'$ es un $(A/\mathfrak p)$-módulo de tipo finito,
y \emph{a fortiori} un $A$-módulo de tipo finito.

Así pues, puesto que la topología $xA'$-preádica de $A'$ es separada, $A'$ es un
subanillo de su completado $\widehat{A'}$ para esta topología; pero, puesto que
$A'/xA'$ es un $A$-módulo de tipo finito y $A$ es separado y completo para
la topología $xA$-preádica, $\widehat{A'}$ es un $A$-módulo de tipo finito en
virtud de (0\textsubscript{I}, 7.2.9), por tanto lo mismo vale para $A'$, C.Q.D.
\end{proof}

\begin{corollary}[23.1.4]
\label{0.23.1.4-es}
Sea $A$ un anillo japonés noetheriano e íntegramente cerrado. Entonces todo
anillo de series formales $A[[T_1,\ldots,T_r]]$ es un anillo japonés.
\end{corollary}

\begin{proof}
Se sabe que, para todo anillo noetheriano íntegramente cerrado $B$, el anillo de
series formales $B[[T]]$ es noetheriano e íntegramente cerrado (Bourbaki,
\emph{Alg. comm.}, cap.~V, §~1, n\textsuperscript{o}~4, prop.~14); por
tanto, por inducción sobre $r$\footnote{Lectura diplomática: la fuente imprime n; se corrige a r, que es el número de variables ligado en el enunciado.}, podemos limitarnos a probar que $A[[T]]$ es un anillo
japonés. Ahora bien, el elemento $x=T$ cumple todas las condiciones de (23.1.3), de donde
se sigue la conclusión.
\end{proof}

\begin{theorem}[23.1.5]
\label{0.23.1.5-es}
\textnormal{(Nagata).} Todo anillo local noetheriano íntegro completo $A$ es
un anillo japonés.
\end{theorem}

\begin{proof}
Se sabe (19.8.8, (ii)) que existe un subanillo $B$ de $A$ que es un
anillo local completo y regular, tal que $A$ es una $B$-álgebra finita; basta
evidentemente probar

\oldpage[0\textsubscript{IV-1}]{216}

que $B$ es un anillo japonés; en otras palabras, podemos limitarnos al caso en que
$A$ es además \emph{regular}. Razonemos por inducción sobre
$n=\dim(A)$; el teorema es evidente para $n=0$. Supongamos, pues, $n>0$ y,
denotando por $\mathfrak m$ el ideal maximal de $A$, sea $x$ un elemento de
$\mathfrak m-\mathfrak m^2$; se sabe (17.1.8) que $A/xA$ es un anillo
regular, y además ((17.1.7) y (16.3.4)) que
$\dim(A/xA)=n-1$; además, el ideal $xA$ es cerrado en $A$
(0\textsubscript{I}, 7.3.5), por lo que $A/xA$ es completo. En virtud de la hipótesis
de inducción, $A/xA$ es, por tanto, un anillo japonés, y de
(23.1.3) se deduce que lo mismo vale para $A$.
\end{proof}

\begin{corollary}[23.1.6]
\label{0.23.1.6-es}
Sean $A$ un anillo local noetheriano íntegro completo y $K$ su cuerpo de
fracciones; sea $K'$ una extensión finita de $K$; entonces la clausura integral
$A'$ de $A$ en $K'$ es un anillo local completo.
\end{corollary}

\begin{proof}
Ya se sabe por (23.1.5) que $A'$ es una $A$-álgebra finita, por tanto un anillo
semilocal noetheriano completo y, en consecuencia, un \emph{producto directo}
de anillos locales; pero como $A'$ es íntegro, es necesariamente un
anillo local.
\end{proof}

\begin{proposition}[23.1.7]
\label{0.23.1.7-es}
Sean $A$ un anillo local noetheriano íntegro y $K$ su cuerpo de fracciones;
sean $K'$ una extensión finita de $K$ y $A'$ la clausura integral de $A$ en
$K'$. Se supone que el completado $\widehat A$ es reducido, y se denota por
$R$ su anillo total de fracciones.
\begin{enumerate}
  \item[(i)] Si el anillo $R\otimes_KK'$ es reducido (lo que ocurrirá en
  particular cuando $K'$ sea una extensión separable de $K$, siendo $R$ un
  producto directo de cuerpos que son extensiones de $K$ (Bourbaki, \emph{Alg.},
  cap.~VIII, §~7, n\textsuperscript{o}~3, cor.~1 del teo.~1)), entonces $A'$ es
  un $A$-módulo de tipo finito.
  \item[(ii)] Si, en particular, $R$ es una $K$-álgebra separable, entonces $A$
  es un anillo japonés.
\end{enumerate}
\end{proposition}

\begin{proof}
\emph{(i)} Pongamos, para simplificar, $A_1=\widehat A$,
$A'_1=A'\otimes_AA_1$, $K'_1=K'\otimes_AA_1$. Como $A_1$ es un
$A$-módulo plano (0\textsubscript{I}, 7.3.5), $A'_1$ se identifica con un
subanillo de $K'_1$ y es evidentemente entero sobre $A_1$. Por otra parte, si
se pone $K_1=K\otimes_AA_1$, se puede escribir
$K'_1=K'\otimes_KK_1$; como $A$ es íntegro y $A_1$ es un $A$-módulo
plano, todo elemento $\neq0$ de $A$ es $A_1$-regular
(0\textsubscript{I}, 6.3.4); como $K=S^{-1}A$, donde $S=A-\{0\}$, se tiene
$K_1=S^{-1}A_1$, y la observación anterior prueba que $K_1$ se identifica con un
subanillo del anillo total de fracciones $R$ de $A_1$. Puesto que todo
$K$-módulo es plano, $K'_1$ se identifica con un subanillo de $K'\otimes_KR$,
que por hipótesis es reducido y es una $R$-álgebra finita. Designemos por
$\mathfrak q_i$ ($1\leqslant i\leqslant n$) los ideales primos minimales de
$A_1$, por $B_i$ el anillo íntegro $A_1/\mathfrak q_i$, por $L_i$ el cuerpo de
fracciones de $B_i$, de modo que $R$ es producto directo de los $L_i$, y
$K'\otimes_KR$ producto directo de los $K'\otimes_KL_i$; estas últimas
$K$-álgebras, al ser reducidas, son productos directos de cuerpos,
extensiones finitas de los $L_i$. Como $B_i$ es un anillo local íntegro
\emph{completo}, el teorema de Nagata (23.1.5) prueba que la clausura
integral de $B_i$ en $K'\otimes_KL_i$ es un $B_i$-módulo de tipo finito, y
\emph{a fortiori} un $A_1$-módulo de tipo finito; como todo elemento de
$K'\otimes_KL_i$ que es entero sobre $A_1$ es también entero sobre $B_i$ (al ser
anulado por $\mathfrak q_i$), se concluye finalmente que la clausura
integral de $A_1$ en $K'\otimes_KR$ es un $A_1$-módulo de tipo finito. Pero
como $A'_1$ está contenido en esta clausura integral y $A_1$ es
noetheriano, $A'_1$ es también un $A_1$-módulo de tipo finito. Por último, como $A_1$
es un $A$-módulo fielmente plano (0\textsubscript{I}, 7.3.5), se concluye
que $A'$ es un $A$-módulo de tipo finito (Bourbaki, \emph{Alg. comm.},
cap.~I, §~3, n\textsuperscript{o}~6, prop.~11).

\oldpage[0\textsubscript{IV-1}]{217}

\emph{(ii)} La hipótesis implica que los $L_i$ son extensiones separables
de $K$, por tanto $K'\otimes_KL_i$ es reducido (Bourbaki, \emph{Alg.}, cap.~VIII,
§~7, n\textsuperscript{o}~3, teo.~1) y se puede aplicar (i) a toda
extensión finita $K'$ de $K$, lo que demuestra nuestra afirmación.
\end{proof}

\subsection{Clausura integral de un anillo local noetheriano íntegro}
\label{subsection:0.23.2-es}

\begin{env}[23.2.1]
\label{0.23.2.1-es}
En lo que sigue, para todo anillo $A$, denotaremos por $A_{\mathrm{red}}$, para
abreviar, el cociente de $A$ por su nilradical (de modo que si
$X=\operatorname{Spec}(A)$, se tiene
$X_{\mathrm{red}}=\operatorname{Spec}(A_{\mathrm{red}})$).

Diremos que un anillo \emph{local} $A$ es \emph{unirrama} si el anillo
$A_{\mathrm{red}}$ es íntegro y si la clausura integral de
$A_{\mathrm{red}}$ es un anillo \emph{local}, lo que generaliza la definición
dada en (III, 4.3.6). Diremos que $A$ es \emph{geométricamente
unirrama} si es unirrama y si el cuerpo residual del anillo local,
clausura integral de $A_{\mathrm{red}}$, es una extensión \emph{radicial}
del de $A$. Es claro que un anillo local íntegramente cerrado es
geométricamente unirrama~; de (23.1.6) resulta que un anillo local
noetheriano íntegro \emph{completo} es unirrama.
\end{env}

\begin{lemma}[23.2.2]
\label{0.23.2.2-es}
\textnormal{(*)}\footnote{En lo que sigue de este número, utilizamos
nociones y resultados expuestos en el cap.~IV, §§~2 y 5~; como los resultados de
este número no se utilizan antes del cap.~IV, §~6, esto no da lugar a un
círculo vicioso.} Sean $A$ un anillo íntegro, $A'$ su clausura integral~; para
que el morfismo canónico
$f:\operatorname{Spec}(A')\to\operatorname{Spec}(A)$ sea radicial, es necesario y
suficiente que para todo $\mathfrak p\in\operatorname{Spec}(A)$,
$A_{\mathfrak p}$ sea geométricamente unirrama~; cuando así ocurre,
$f$ es un homeomorfismo.
\end{lemma}

\begin{proof}
En efecto, para todo $\mathfrak p\in\operatorname{Spec}(A)$, la clausura
integral de $A_{\mathfrak p}$ es $A'_{\mathfrak p}$ (Bourbaki,
\emph{Alg. comm.}, cap.~V, §~1, n\textsuperscript{o}~5, prop.~16), y todos
los ideales primos de $A'_{\mathfrak p}$ situados sobre
$\mathfrak pA_{\mathfrak p}$ son maximales (\emph{loc. cit.}, §~2,
n\textsuperscript{o}~1, prop.~1). Decir que $f$ es inyectivo significa, por tanto, que
para todo $\mathfrak p\in\operatorname{Spec}(A)$, $A'_{\mathfrak p}$ es un
anillo \emph{local}, es decir, que cada uno de los $A_{\mathfrak p}$ es
unirrama. Decir que todo $A_{\mathfrak p}$ es geométricamente unirrama
significa entonces que $f$ es radicial en virtud de (I, 3.5.8). Cuando así
ocurre, $f$ es sobreyectivo y cerrado (II, 6.1.10), luego es un homeomorfismo.
\end{proof}

\begin{lemma}[23.2.3]
\label{0.23.2.3-es}
Sean $A$ un anillo íntegro, $K$ su cuerpo de fracciones, $B$ un anillo
noetheriano, $\varphi:A\to B$ un homomorfismo de anillos que hace de $B$ un
$A$-módulo plano. Sea $K'$ una extensión de $K$~; pongamos
$B_1=B_{\mathrm{red}}$, y sea $R$ el anillo total de fracciones de $B_1$.
Entonces, para toda sub-$A$-álgebra $A'$ de $K'$, el homomorfismo canónico
\begin{equation}
\label{0.23.2.3.1-es}
  (A'\otimes_AB)_{\mathrm{red}}
  \longrightarrow(K'\otimes_AR)_{\mathrm{red}}
  \tag{23.2.3.1}
\end{equation}
es inyectivo.\footnote{Lectura diplomática: la fuente numera esta fórmula 23.1.8.1; el contexto del lema impone 23.2.3.1.}
\end{lemma}

\begin{proof}
En efecto, se puede considerar el homomorfismo canónico
$h:A'\otimes_AB\to K'\otimes_AR$ como compuesto de los homomorfismos siguientes
\[
  A'\otimes_AB\xrightarrow{u}K'\otimes_AB
  \xrightarrow{v}K'\otimes_AB_1\xrightarrow{w}K'\otimes_AR.
\]
Como $B$ es un $A$-módulo plano, $u$ es inyectivo~; asimismo, como $K'$ es un
$K$-módulo plano y $K$ un $A$-módulo plano, $K'$ es un $A$-módulo plano, y
$w$ es, por tanto, inyectivo, puesto que $B_1$ es reducido y, por consiguiente, el homomorfismo
$B_1\to R$ es inyectivo. Por último, por la misma razón, si $\mathfrak N$ es

\oldpage[0\textsubscript{IV-1}]{218}

el nilradical de $B$, el núcleo de $v$ es $K'\otimes_A\mathfrak N$, por lo que está
contenido en el nilradical de $K'\otimes_AB$~; se concluye inmediatamente que si
la imagen por $h=w\circ v\circ u$ de un elemento $x\in A'\otimes_AB$ es
nilpotente, $x$ es nilpotente, lo que demuestra el lema.
\end{proof}

\begin{proposition}[23.2.4]
\label{0.23.2.4-es}
Sean $A$ un anillo íntegro, $K$ su cuerpo de fracciones, $B$ un anillo
noetheriano, $\mathfrak q_i$ ($1\leqslant i\leqslant n$) sus ideales primos
minimales, $\varphi:A\to B$ un homomorfismo de anillos. Se supone que los
$B/\mathfrak q_i$ son anillos japoneses y que $B$ es un $A$-módulo
plano. Sea $K'$ una extensión finita de $K$, y sea $(C_\lambda)$ la familia
filtrante creciente de subanillos de $K'$ que son $A$-álgebras
finitas y admiten $K'$ como cuerpo de fracciones~; la unión $A'$ de los
$C_\lambda$ es, por tanto, la clausura integral de $A$ en $K'$. Entonces~:
\begin{enumerate}
  \item[(i)] Existe un índice $\alpha$ tal que para $\lambda\geqslant
  \alpha$, el homomorfismo canónico
  \begin{equation}
  \label{0.23.2.4.1-es}
    (C_\alpha\otimes_AB)_{\mathrm{red}}
    \longrightarrow(C_\lambda\otimes_AB)_{\mathrm{red}}
    \tag{23.2.4.1}
  \end{equation}
  es biyectivo.
  \item[(ii)] Si además $B$ es un $A$-módulo fielmente plano, el morfismo
  $\operatorname{Spec}(A')\to\operatorname{Spec}(C_\alpha)$ es radicial.
\end{enumerate}
\end{proposition}

\begin{proof}
\emph{(i)} Sea $B_1=B_{\mathrm{red}}$, y sea $R$ el anillo total de
fracciones de $B_1$~; como $A$ es íntegro y $B$ un $A$-módulo plano, todo
elemento $\neq0$ de $A$ es $B$-regular (0\textsubscript{I}, 6.3.4), por lo que
el homomorfismo compuesto $A\xrightarrow{\varphi}B\to B_1$ se prolonga a un
homomorfismo $K\to R$~; se puede escribir entonces
$K'\otimes_AR=(K'\otimes_AK)\otimes_KR$, y como $K'\otimes_AK$
se identifica con $K'$, se tiene $K'\otimes_AR=K'\otimes_KR$. Como $R$ es producto
directo de los cuerpos de fracciones $L_i$ de los $B/\mathfrak q_i$,
$K'\otimes_AR$ es producto directo de los $K'\otimes_KL_i$, y por consiguiente
$(K'\otimes_AR)_{\mathrm{red}}$ es producto directo de un número finito
de extensiones finitas de los $L_i$. En virtud de la hipótesis, la clausura
integral de $B/\mathfrak q_i$ en una extensión finita de $L_i$ es un
$(B/\mathfrak q_i)$-módulo de tipo finito, luego un $B$-módulo de tipo finito~;
se concluye que la clausura integral $B'$ de $B$ en
$(K'\otimes_AR)_{\mathrm{red}}$ es un $B$-módulo de tipo finito. En virtud de
(23.2.3), los $(C_\lambda\otimes_AB)_{\mathrm{red}}$ se identifican
canónicamente con subanillos de $(K'\otimes_AR)_{\mathrm{red}}$ que son
$B$-álgebras finitas, y están, por tanto, contenidos en $B'$. Como $B$ es noetheriano
y $B'$ un $B$-módulo de tipo finito, la familia filtrante de los
$(C_\lambda\otimes_AB)_{\mathrm{red}}$ admite un \emph{elemento mayor}
$(C_\alpha\otimes_AB)_{\mathrm{red}}$, lo que demuestra (i).

\emph{(ii)} Como $B$ es un $A$-módulo fielmente plano, basta
(IV, 2.6.1, (v)) demostrar que el morfismo
$\operatorname{Spec}(A'\otimes_AB)\to
\operatorname{Spec}(C_\alpha\otimes_AB)$ es radicial, o, lo que es equivalente
(I, 5.1.6), que el morfismo
$\operatorname{Spec}(A'\otimes_AB)_{\mathrm{red}}\to
\operatorname{Spec}(C_\alpha\otimes_AB)_{\mathrm{red}}$ es radicial~; pero
se tiene $A'\otimes_AB=\varinjlim_\lambda(C_\lambda\otimes_AB)$, luego
(IV, 5.13.2)
$(A'\otimes_AB)_{\mathrm{red}}=
\varinjlim_\lambda(C_\lambda\otimes_AB)_{\mathrm{red}}$, y la conclusión
resulta de (i).
\end{proof}

\begin{corollary}[23.2.5]
\label{0.23.2.5-es}
Sean $A$ un anillo local noetheriano íntegro, $K$ su cuerpo de fracciones,
$K'$ una extensión finita de $K$. Sea $(C_\lambda)$ la familia filtrante
creciente de subanillos de $K'$ que son $A$-álgebras finitas (por tanto,
anillos noetherianos semilocales (Bourbaki, \emph{Alg. comm.}, cap.~IV,
§~2, n\textsuperscript{o}~5, cor.~3 de la prop.~9)) y admiten $K'$ como
cuerpo de fracciones. Entonces existe $\alpha$ tal que el homomorfismo
$(\widehat C_\alpha)_{\mathrm{red}}\to
(\widehat C_\lambda)_{\mathrm{red}}$ sea un isomorfismo para
$\lambda\geqslant\alpha$, y si $A'$ es la clausura integral de $A$ en
$K'$, el morfismo
$\operatorname{Spec}(A')\to\operatorname{Spec}(C_\alpha)$ es radicial
(cf. (23.2.2)).
\end{corollary}

\oldpage[0\textsubscript{IV-1}]{219}

\begin{proof}
Se aplica (23.2.4) tomando $B=\widehat A$~; como los
$B/\mathfrak q_i$ son anillos locales noetherianos completos, son anillos
japoneses (23.1.5) y $B$ es un $A$-módulo fielmente plano
(0\textsubscript{I}, 7.3.5)~; además, se tiene
$C_\lambda\otimes_AB=\widehat C_\lambda$ (Bourbaki, \emph{Alg. comm.},
cap.~III, §~3, n\textsuperscript{o}~4, th.~3 y cap.~IV, §~2,
n\textsuperscript{o}~5, cor.~3 de la prop.~9).
\end{proof}

\begin{corollary}[23.2.6]
\label{0.23.2.6-es}
Bajo las hipótesis de (23.2.5), la clausura integral $A'$ de $A$ en
$K'$ es un anillo semilocal~; si $\mathfrak m$ es el ideal maximal de $A$,
entonces, para todo ideal maximal $\mathfrak m'$ de $A'$,
$A'/\mathfrak m'$ es una extensión finita de $A/\mathfrak m$.
\end{corollary}

\begin{proof}
El hecho de que el homomorfismo
$(\widehat C_\alpha)_{\mathrm{red}}\to
(\widehat C_\lambda)_{\mathrm{red}}$ sea biyectivo para
$\lambda\geqslant\alpha$ implica que el número de ideales maximales de
$C_\lambda$ es constante para $\lambda\geqslant\alpha$ y que si
$\mathfrak m_\alpha$ es un ideal maximal de $C_\alpha$ y
$\mathfrak m_\lambda$ el único ideal maximal de $C_\lambda$ situado sobre
$\mathfrak m_\alpha$, los cuerpos $C_\alpha/\mathfrak m_\alpha$ y
$C_\lambda/\mathfrak m_\lambda$ son canónicamente isomorfos. La conclusión
resulta de que $\mathfrak m'=\varinjlim\mathfrak m_\lambda$ y
$A'/\mathfrak m'=\varinjlim(C_\lambda/\mathfrak m_\lambda)$
(0\textsubscript{III}, 10.3.1.1.3).
\end{proof}

\begin{proposition}[23.2.7]
\label{0.23.2.7-es}
\textnormal{(Y. Mori).} Sean $A$ un anillo local noetheriano íntegro, $K$ su
cuerpo de fracciones, $A'$ la clausura integral de $A$. Entonces $A'$ es un
anillo de Krull (Bourbaki, \emph{Alg. comm.}, cap.~VII, §~1) semilocal~;
en otras palabras, existe una familia $(V_\lambda)$ de anillos de valoración
discreta que tienen $K$ como cuerpo de fracciones, tales que
$A'=\bigcap_\lambda V_\lambda$ y que todo $x\in K$ pertenece a todos los
$V_\lambda$ salvo a un número finito de ellos. Además, existe un
subanillo $C$ de $K$ que es un $A$-álgebra finita, tal que los
$V_\lambda$ son las clausuras integrales de los anillos
$C_{\mathfrak p_\lambda}$, donde $(\mathfrak p_\lambda)$ es la familia de los
ideales primos de altura $1$ del anillo local $C$.
\end{proposition}

Demostremos primero el lema siguiente~:

\begin{lemma}[23.2.7.1]
\label{0.23.2.7.1-es}
Sean $A$, $B$ dos anillos, $\varphi:A\to B$ un homomorfismo que hace de $B$
un $A$-módulo fielmente plano. Se supone que $A$ es íntegro~; entonces, si
$C=B_{\mathrm{red}}$, el homomorfismo compuesto
$\psi:A\to B\to B_{\mathrm{red}}=C$ es inyectivo y se prolonga a un
homomorfismo inyectivo del cuerpo de fracciones $K$ de $A$ en el anillo total
de fracciones $R$ de $C$~; además, si $C'$ es la clausura integral de
$C$ en $R$, $C'\cap K$ es la clausura integral de $A$.
\end{lemma}

\begin{proof}
Como $A$ es íntegro y $\varphi$ es inyectivo, la intersección de $\varphi(A)$
y del nilradical $\mathfrak N$ de $B$ se reduce a $0$, y como
$C=B/\mathfrak N$, $\psi$ es inyectivo. Todo $a\neq0$ de $A$, al no ser
divisor de cero en $A$, tampoco lo es en $B$ por planitud
(0\textsubscript{I}, 6.3.4)~; se deduce que $a$ tampoco es divisor
de $0$ en $C$, pues si se tuviera $ax\in\mathfrak N$ para un
$x\notin\mathfrak N$ de $B$, se obtendría $a^nx^n=0$ para un entero $n$,
lo que contradice lo anterior, puesto que $x^n\neq0$. Se puede, por tanto, prolongar
$\psi$ a un homomorfismo inyectivo de $K$ en $R$. Para demostrar la última
afirmación, observemos que es claro que la clausura integral $A'$ de $A$ está
contenida en $C'\cap K$. Recíprocamente, sea $x\in C'\cap K$~; $x$ es
entero sobre $C$ y, \emph{a fortiori}, sobre $B$, es decir, $B[x]$ es una
$B$-álgebra finita~; además, $K$ se identifica con un subanillo de
$B\otimes_AK$, y el subanillo $B[x]$ de $B\otimes_AK$ se identifica con
$B\otimes_AA[x]$ por planitud~; se concluye que $A[x]$ es un $A$-módulo
de tipo finito (Bourbaki, \emph{Alg. comm.}, cap.~I\textsuperscript{er}, §~3,
n\textsuperscript{o}~6, prop.~11), y, por tanto, que $x\in A'$.
\end{proof}

Establecido este lema, vamos a aplicarlo a la inyección canónica de $A$
en su completado $\widehat A$, que es un $A$-módulo fielmente plano~;
$B=(\widehat A)_{\mathrm{red}}=\widehat A/\mathfrak N$ (donde $\mathfrak N$ es
el nilradical de $\widehat A$) es un anillo local noetheriano completo reducido,
cuyo anillo total de

\oldpage[0\textsubscript{IV-1}]{220}

fracciones $R$ es, por tanto, producto directo de un número finito de cuerpos $L_i$~;
la imagen canónica $B_i$ de $B$ en $L_i$ es un anillo local noetheriano
\emph{íntegro y completo}, cuyo cuerpo de fracciones es $L_i$, y la
clausura integral $B'$ de $B$ en $R$ es \emph{producto directo} de los
$B'_i$, donde $B'_i$ es la clausura integral de $B_i$. Pero en virtud de
(23.1.5), $B'_i$ es una $B_i$-álgebra \emph{finita}, y, por tanto, un anillo local
\emph{noetheriano} íntegramente cerrado. Se sabe (Bourbaki,
\emph{Alg. comm.}, cap.~VII, §~1, n\textsuperscript{o}~3, cor. del th.~2) que
$B'_i$ es un anillo de Krull. Ahora bien, para todo $i$, se tiene un homomorfismo
$\varphi_i:K\to L_i$, que es inyectivo, y por consiguiente
$\varphi_i^{-1}(B'_i)$ es un anillo de Krull en $K$. Pero como
$A'=B'\cap K$ en virtud del lema y
$A'=\bigcap_i\varphi_i^{-1}(B'_i)$\footnote{Lectura diplomática: el miembro izquierdo impreso es B prima; se corrige a A prima porque el miembro derecho es un subanillo de K.}, $A'$ es intersección de un número finito
de anillos de Krull y, por tanto, es un anillo de Krull.

Para demostrar la última afirmación de la proposición, observemos que existe una
$A$-álgebra finita $C\subset K$ tal que el morfismo
$\operatorname{Spec}(A')\to\operatorname{Spec}(C)$ es radicial y un
homeomorfismo (23.2.4)~; para todo ideal primo
$\mathfrak p'\in\operatorname{Spec}(A')$, $\mathfrak p'$ es, por tanto, el único
ideal primo de $A'$ situado sobre $\mathfrak p=\mathfrak p'\cap C$ y
$A'_{\mathfrak p'}$ es la clausura integral de $C_{\mathfrak p}$~; por otra parte,
el hecho de que $\operatorname{Spec}(A')\to\operatorname{Spec}(C)$ sea un
homeomorfismo implica que la aplicación
$\mathfrak p'\mapsto\mathfrak p'\cap C$ es una biyección del conjunto de los
ideales primos de altura $1$ de $A'$ sobre el conjunto de los ideales primos de
altura $1$ de $C$~; de donde resulta la conclusión, teniendo en cuenta Bourbaki,
\emph{Alg. comm.}, cap.~VII, §~1, n\textsuperscript{o}~6, th.~4.

\begin{remarks}[23.2.8]
\label{0.23.2.8-es}
\begin{enumerate}
  \item[(i)] Conviene observar, al aplicar (23.2.6) y
  (23.2.7), que el anillo $A'$ \emph{no es necesariamente un anillo
  noetheriano}, como muestra un ejemplo de Nagata con $\dim(A)=3$ [30].

  \item[(ii)] Las conclusiones de (23.2.7) siguen siendo válidas si $A'$ es
  la clausura integral de $A$ en una extensión finita $K'$ de $K$~; en efecto,
  basta considerar una $A$-álgebra finita $B$ cuyo cuerpo de fracciones sea
  $K'$~; $B$ es un anillo semilocal noetheriano, y, por tanto,
  intersección de un número finito de anillos locales noetherianos
  $B_{\mathfrak m_i}$ ($\mathfrak m_i$ ideales maximales de $B$), y su clausura
  integral (que es igual a $A'$) es intersección de los
  $A'_{\mathfrak m_i}$ (Bourbaki, \emph{Alg. comm.}, cap.~II, §~3,
  n\textsuperscript{o}~3, cor.~4 del th.~1), que son las clausuras integrales
  de los $B_{\mathfrak m_i}$, y, por tanto, anillos de Krull por (23.2.7)\footnote{Lectura diplomática: la fuente remite a la inexistente 23.1.13; el resultado utilizado es la Proposición 23.2.7.}~; por
  consiguiente, $A'$ es un anillo de Krull.

  \item[(iii)] Se puede demostrar ([30], 3.3.10) que para todo anillo noetheriano
  íntegro $A$ (no necesariamente local), la clausura integral de $A$ es un
  anillo de Krull.
\end{enumerate}
\end{remarks}

\begin{corollary}[23.2.9]
\label{0.23.2.9-es}
Sean $A$ un anillo noetheriano íntegro, $A'$ su clausura integral. Supongamos
que para todo anillo $C$ tal que $A\subset C\subset A'$ y tal que $C$ sea
una $A$-álgebra finita, y para todo ideal primo $\mathfrak q$ de altura $1$
en $C$, $\mathfrak q\cap A$ sea un ideal primo de altura $1$ en $A$.
Sea $P$ el conjunto de los ideales primos de altura $1$ de $A$~; se tiene entonces
\[
  A'=\bigcap_{\mathfrak p\in P}A'_{\mathfrak p}.
\]
\end{corollary}

\begin{proof}
Como $A'$ es un $A$-módulo sin torsión, se tiene
$A'=\bigcap A'_{\mathfrak m}$, donde $\mathfrak m$ recorre el conjunto de los ideales
maximales de $A$ (Bourbaki, \emph{Alg. comm.}, cap.~II, §~3,
n\textsuperscript{o}~3, cor.~4 del th.~1)~;

\oldpage[0\textsubscript{IV-1}]{221}

bastará, por tanto, demostrar que $A'_{\mathfrak m}$ es la intersección de los
$A'_{\mathfrak p}$ para los ideales primos $\mathfrak p\subset\mathfrak m$
de altura $1$. Como $A'_{\mathfrak m}$ es la clausura integral de
$A_{\mathfrak m}$, se ve que podemos limitarnos a demostrar el corolario para
$A_{\mathfrak m}$. Ahora bien, en virtud de (23.2.7), existe una
$A_{\mathfrak m}$-álgebra finita $B$ contenida en $A'_{\mathfrak m}$, tal
que $A'_{\mathfrak m}$ es la intersección de las clausuras integrales de los anillos
$B_{\mathfrak q}$, donde $\mathfrak q$ recorre el conjunto de los ideales primos de
altura $1$ de $B$. Pero $B$ es de la forma $C_{\mathfrak m}$, donde $C$ es
una $A$-álgebra finita~; por tanto, para todo ideal primo $\mathfrak q$ de altura
$1$ en $B$, $\mathfrak q\cap A$ es, por hipótesis, un ideal primo de
altura $1$ en $A$~; como la clausura integral de $B_{\mathfrak q}$
contiene la de $A_{\mathfrak q\cap A}$ y esta última contiene
$A'_{\mathfrak m}$, resulta que $A'_{\mathfrak m}$ es efectivamente
la intersección de los $A'_{\mathfrak p}$ para $\mathfrak p\in P$ y
$\mathfrak p\subset\mathfrak m$, lo que termina la demostración.
\end{proof}

\begin{remarks}[23.2.10]
\label{0.23.2.10-es}
\begin{enumerate}
  \item[(i)] Veremos (IV, 5.6.3 y 5.6.10) que la hipótesis de (23.2.9)
  se cumple cuando $A$ es \emph{universalmente catenario}, en particular
  (IV, 5.10.17 y 5.11.2) cuando $A$ es \emph{cociente de un anillo
  regular}~; por último, se cumple evidentemente si
  $\operatorname{Spec}(A')\to\operatorname{Spec}(A)$ es un homeomorfismo,
  en particular si este morfismo es radicial (23.2.2).

  \item[(ii)] La conclusión de (23.2.9) ya no es necesariamente cierta
  cuando se omite la hipótesis sobre las sub-$A$-álgebras finitas de $A'$. En
  efecto, se puede definir un anillo local noetheriano íntegro $A$, de dimensión
  $2$, cuya clausura integral $A'$ es una $A$-álgebra finita, pero tal que
  existe un ideal primo $\mathfrak p'$ de $A'$ de altura $1$ situado sobre
  el ideal maximal $\mathfrak m$ (de altura $2$) de $A$, y tal además que
  para todo ideal primo $\mathfrak p$ de altura $1$ de $A$,
  $A_{\mathfrak p}$ es íntegramente cerrado (IV, 5.6.11)~; la intersección de
  estos anillos $A_{\mathfrak p}$ no puede entonces ser igual a $A'$, puesto que $A'$
  es un anillo de Krull (Bourbaki, \emph{Alg. comm.}, cap.~VII, §~1,
  n\textsuperscript{o}~5, prop.~9).
\end{enumerate}
\end{remarks}

\begin{flushright}
\emph{(Continuará.)}
\end{flushright}

\end{document}
